\documentclass[11pt]{article}

\usepackage[a4paper,margin=1in]{geometry}
\usepackage[T1]{fontenc}
\DeclareUnicodeCharacter{0394}{\ensuremath{\Delta}}
\DeclareUnicodeCharacter{03B5}{\ensuremath{\varepsilon}}
\usepackage{microtype}
\usepackage{amsmath}
\usepackage{amssymb}
\usepackage{amsthm}
\usepackage{mathtools}
\usepackage{bm}
\usepackage{graphicx}
\usepackage{booktabs}
\usepackage{enumitem}
\usepackage{xcolor}
\usepackage[numbers,sort&compress]{natbib}
\usepackage[
  colorlinks=true,
  linkcolor=blue,
  citecolor=teal,
  urlcolor=blue
]{hyperref}
\usepackage[nameinlink,noabbrev]{cleveref}

% Common notation used throughout the paper.
\DeclareMathOperator{\tr}{tr}
\DeclareMathOperator{\rank}{rank}
\DeclareMathOperator{\diag}{diag}
\DeclareMathOperator{\supp}{supp}
\DeclareMathOperator{\nnz}{nnz}
\DeclareMathOperator{\poly}{poly}
\DeclareMathOperator{\polylog}{polylog}
\DeclareMathOperator{\cond}{cond}

\newcommand{\Otilde}{\widetilde{O}}
\newcommand{\Thetatilde}{\widetilde{\Theta}}
\newcommand{\Omegatilde}{\widetilde{\Omega}}

\DeclarePairedDelimiter{\norm}{\lVert}{\rVert}
\DeclarePairedDelimiterX{\inner}[2]{\langle}{\rangle}{#1, #2}

\newcommand{\R}{\mathbb{R}}
\newcommand{\C}{\mathbb{C}}
\newcommand{\Sym}{\mathbb{S}}
\newcommand{\Splus}{\mathbb{S}_{+}}
\newcommand{\Dtr}{D_{\mathrm{tr}}}
\newcommand{\ket}[1]{\lvert #1\rangle}
\newcommand{\bra}[1]{\langle #1\rvert}


\theoremstyle{plain}
\newtheorem{theorem}{Theorem}[section]
\newtheorem{proposition}[theorem]{Proposition}
\newtheorem{lemma}[theorem]{Lemma}
\newtheorem{corollary}[theorem]{Corollary}

\theoremstyle{definition}
\newtheorem{definition}[theorem]{Definition}
\newtheorem{example}[theorem]{Example}
\newtheorem{assumption}[theorem]{Assumption}
\newtheorem{openproblem}[theorem]{Open Problem}

\theoremstyle{remark}
\newtheorem{remark}[theorem]{Remark}

\crefname{openproblem}{open problem}{open problems}
\Crefname{openproblem}{Open Problem}{Open Problems}
\crefname{assumption}{assumption}{assumptions}
\Crefname{assumption}{Assumption}{Assumptions}

\title{Advances in Quantum Interior-Point Methods\\
for Sparse Problems: Conditioning Geometry, Query Lower Bounds, and
Access-Model Frontiers}
\author{}
\date{}

\begin{document}

\maketitle

\begin{abstract}
Can quantum interior-point methods (QIPMs) outperform classical
interior-point methods on sparse linear and semidefinite programs?  We give a
calibrated answer by separating geometry, data access, and output.  First, a
sublevel-chord law lower-bounds the reduced central-path Hessian condition
number for every self-concordant barrier; canonical log and log-det barriers
attain this floor up to instance constants, so barrier redesign cannot in
general remove the geometric obstruction.  Degenerate LP log-barrier tails
nevertheless have clustered spectra that are benign for classical Krylov
methods.  On the same clusters, plain quantum block-encoding access in
sufficiently large dimension has linear worst-case
condition dependence, while norm-tight factor access unconditionally supports a
square-root polynomial route in its model; only the corresponding end-to-end
solve needs explicit overlap assumptions.  Exact coupling and effective
spectral-dimension laws then identify which beyond-condition-number solver is
appropriate for a given Newton right-hand side.

For block log-det centers we prove a winner-take-all condensation transition,
finite mass--conditioning--visibility laws, and matching raw-query bounds for
sparse holonomy SDPs.  Verified-sample decoders and polynomial symmetrization
give accuracy-dependent state-conversion curves.  Separate LP and SDP
families show that loading or recovering optimization states can require
linear queries even when the relevant reduced Newton condition number is one.
Convex-mixture, electrical, Fourier, congruence, preconditioning, and
state-interface frontiers delimit the gadget and access classes covered by
these lower bounds.  The surviving positive results are correspondingly
conditional: residual-certified pathwise direction reuse, partition-free
equilibration and face repair, implicit dual coordinates with compressed
output, and block-angular border aggregation.

The results do not give an unconditional end-to-end quantum advantage or
disadvantage for sparse optimization.  They instead locate the remaining
bottlenecks in loading, recovery, dynamic oracle maintenance, and output.
Every result reported here was re-derived and re-verified; in that process we
correct the simulator-norm and clock analysis in the QCPM preprint and several
earlier claims from our own research program.
\end{abstract}

\tableofcontents

\section{Introduction}
\label{sec:introduction}

Sparse linear and semidefinite programs are an attractive test case for
quantum optimization.  Their data may be queried or applied without reading a
dense matrix, and an interior-point method reduces optimization to a sequence
of structured linear systems.  A quantum linear-system solver (QLS) can, in
favorable access models, prepare a state proportional to a Newton direction
without first writing every coordinate.  These observations motivate a simple
question:
\begin{quote}
Do quantum interior-point methods beat classical interior-point methods on
sparse problems?
\end{quote}
The word ``sparse'' does not by itself determine the answer.  An input matrix
can be sparse while its normal equations, null-space basis, factorization, or
inverse is dense.  Conversely, a matrix that looks unfavorable through its
condition number can have a clustered spectrum that a classical Krylov method
uses efficiently.  A quantum state proportional to a direction is not yet a
classical iterate, and a preconditioned matrix oracle is not free merely
because the matrix it represents is well conditioned.  The relevant
comparison must therefore name the Newton formulation, the input oracle, the
block-encoding normalization, the right-hand side, the requested output, and
the classical solver that receives the same structure.

This paper develops that comparison at three levels.  At the \emph{geometric}
level, we ask which central-path conditioning is inherent in the optimization
problem and which can be changed by selecting another barrier.  At the
\emph{access-model} level, we ask what a classical or quantum linear solver can
extract from the same spectral structure and what it costs to construct the
operator and right-hand side it receives.  At the \emph{interface} level, we
ask whether the output of one solve can be recovered, certified, and used to
form the next solve.  The resulting picture is neither a blanket quantum
speedup nor a blanket no-go theorem.  It is a set of sharp laws and scoped
frontiers that identify where either conclusion would still have to be
proved.

\subsection{State of play}
\label{sec:introduction-state-of-play}

The first generation of QIPMs inserted a QLS into a classical path-following
framework.  Kerenidis and Prakash treated LPs and SDPs; Casares and
Martin-Delgado developed a predictor--corrector LP method; Augustino et al.
gave conic and semidefinite variants; and the Terlaky-group line developed
inexact feasible, infeasible, orthogonal-subspace, and refinement-based
methods
\cite{kerenidis2020-quantum-interior-point-method-lps,
casares2020-quantum-interior-point-predictor-corrector,
augustino2023-quantum-interior-point-methods-semidefinite,
brandon2023-quantum-interior-point-methods-conic,
augustino2021-an-infeasible-inexact-quantum-interior,
augustino2022-quantum-interior-point-methods-for,
wu2023-inexact-feasible-quantum-interior-point,
mohammadisiahroudi2024-efficient-use-quantum-linear-system,
mohammadisiahroudi2025-inexact-feasible-interior-point-method}.
These methods made the potential role of QLS state preparation precise, but
their ledgers also exposed condition, normalization, tomography, QRAM, and
classical-update costs.  Dalzell et al. emphasize the same end-to-end
accounting issues in their QIPM review and resource analysis
\cite{dalzell2025-quantum-interior-point-methods,
dalzell2023-end-end-resource-analysis-quantum}.

Two newer lines change the architecture rather than merely improving the
linear solver.  The quantum central-path method of Augustino, Leng,
Nannicini, Terlaky, and Wu uses Hamiltonian evolution and potential
evaluation rather than repeated QLS tomography
\cite{augustino2024-quantum-central-path-algorithm-linear}.  The
self-concordant Schr\"odinger-operator program of Gribling et al. likewise
seeks optimization without a Newton condition-number factor
\cite{gribling2025-self-concordant-schrodinger-operators-spectral}.  These
approaches show why a conditioning theorem must be scoped to a named Newton
matrix: a lower bound on a barrier Hessian is not a lower bound on every
quantum optimization architecture.

Apers and Gribling take a different condition-number-free route.  Their tall
LP algorithms use quantum spectral approximation and local-norm estimation,
return an explicit solution, and prove a row-query separation when the number
of constraints is sufficiently larger than the variable dimension
\cite{apers2026-quantum-speedups-linear-programming-interior}.  Within the
QIPM evidence compared here, this is the only proved asymptotic
quantum--classical separation.  It is a row-query separation under a succinct
row oracle, not automatically a wall-clock separation from the best
\(\nnz\)-aware sparse solver.  Its requirements and high powers of the
variable dimension also make it a specific positive regime rather than a
generic consequence of sparsity.

The QLS primitive itself has evolved from HHL through the
Childs--Kothari--Somma precision improvement and QSVT pseudoinversion
\cite{harrow2009-quantum-algorithm-linear-systems-equations,
andrew2017-quantum-algorithm-systems-linear-equations,
andras2019-quantum-singular-value-transformation-beyond}.  Orsucci and Dunjko
proved access-sensitive upper and lower bounds for positive-definite systems.
Mori et al. proved an \(\Omega(\kappa\sqrt{s})\) sparse-access lower bound at
fixed constant error and restated an \(\Omega(\kappa\log(1/\epsilon))\)
constant-sparsity bound from earlier unpublished work of Harrow and Kothari.
Dalzell, Li, and Su introduced truncation and filtering parameters that can be
smaller than a worst-case condition number
\cite{orsucci2021-on-solving-classes-of-positive,
mori2026-sparsity-dependent-complexity-lower-bound,
dalzell2026-faster-quantum-linear-system-solver}.  Somma and de Wolf's related
plain-versus-sum-of-squares separation concerns guided ground-state energy,
not QLS, and we keep that distinction explicit in
\cref{sec:spectra-quantum}
\cite{somma2026-optimal-lower-bound-ground-state-energy}.  These results make
``the QLS cost'' an incomplete phrase: cost depends on whether one has plain
matrix access, factor access, sparse-entry access, or a stronger supplied
oracle, and on how the right-hand side overlaps the hard spectral subspace.

Tomography is the second standard bottleneck.  A path-following method usually
needs a classical direction to update all primal and dual coordinates, while
a QLS naturally returns only a normalized state.  State-preparation-unitary
tomography improves the dependence on readout accuracy, but still scales with
the direction dimension
\cite{mohammadisiahroudi2023-exponentially-more-precise-tomography-for,
joran2023-quantum-tomography-state-preparation-unitaries}.  Avoiding that cost
requires a different output contract and, crucially, a way to generate the
next changing diagonal or Newton operator from the compressed representation.
The interface cannot be removed by renaming a state as an iterate.

Classical sparsity exploitation sets a demanding baseline.  Practical IPMs
choose between normal equations and augmented systems, control fill, reuse
symbolic factorizations and preconditioners, and adapt inner tolerances
\cite{wright1997-primal-dual-interior-point-methods,
jacek2012-interior-point-methods-25-years,
gondzio2023-general-purpose-preconditioning-regularized-interior-point}.
For bounded-treewidth structure, supplied elimination orders and specialized
IPMs can be nearly linear in the problem size
\cite{dong2020-a-nearly-linear-time-algorithm,
gu2022-a-faster-small-treewidth-sdp}.  A sparse-QIPM claim must therefore be
compared with a structure-aware factorization or Krylov method, not with dense
Gaussian elimination by default.

\subsection{What a matched comparison counts}
\label{sec:introduction-ledgers}

The algorithmic literature above cannot be compared by retaining only the
dimension and the sparsity of the original constraint matrix.  Its QIPM lines
use different Newton formulations, state-preparation primitives, refinement
schemes, and outputs.  We use a raw-query/QRAM-write ledger throughout the
paper, counting calls after the oracle is precisely defined and writes needed
to build or refresh QRAM.  We pair it with qualitative end-to-end audits of
construction, reuse, recovery, and output.  Gate and arithmetic work are
populated only for selected results, so these audits are not a uniform
wall-clock ledger.  Neither accounting view substitutes for the other.

\paragraph{From sparse input to the solved operator.}
For an LP, a diagonal scaling of a sparse constraint matrix may produce a
normal equation \(AD^2A^\top\), an augmented KKT matrix, or an
orthogonal-subspace system.  Constant row sparsity of \(A\) need not imply
constant sparsity of any of these objects: one dense column creates a clique
in the normal-equation graph, and a null-space basis can be dense.  The
opposite can also occur, since a carefully chosen basis or augmented
formulation can preserve structure that normal equations lose.  Consequently
our notation distinguishes \(\nnz(A)\), the sparsity \(s_H\) of the actual
system \(H\), and the access description of a factor \(G\) when
\(H=G^\dagger G\); the formal definitions are in
\cref{sec:newton-formulations,sec:prelim-quantum-access}.  A claim that starts
from sparse \(A\) must prove the access properties of the operator it actually
solves.

\paragraph{Condition and normalization.}
A block encoding presents \(H/\alpha_H\), so the inversion scale is governed
by \(\alpha_H\lVert H^{-1}\rVert\), not by the spectral condition number in
isolation.  Rescaling \(H\) changes neither this product nor the underlying
problem.  Multiplying separately encoded factors can multiply their
normalizations, and forming a perfect preconditioned product can require the
same input information that the original solve was meant to discover.  The
factor, congruence, and parity results of \cref{sec:preconditioning} make these
exchanges exact on representative families.  On the classical side, an
effective preconditioner also has a construction and application cost, but it
can often be reused and can exploit entries already stored explicitly.  A
matched comparison prices both constructions rather than granting either side
an ideal inverse oracle.

\paragraph{Right-hand sides and solve accuracy.}
Worst-case QLS complexity is not always the right estimate for a Newton step.
The direction depends on a particular RHS, and truncation can ignore a hard
eigenspace only if the requested output error permits losing its solution
mass.  A residual guarantee, a normalized-state error, and a coordinatewise
direction error are different contracts.  They are connected by
condition-dependent inequalities, so comparing a quantum state-error bound
with a classical residual tolerance without conversion is not meaningful.  The
coupling and spectral-measure analysis in \cref{sec:beyond-kappa} evaluates
these distinctions on exact IPM systems, while the output contracts of
\cref{sec:prelim-output-contracts} state which error is being promised.

\paragraph{State preparation and readout.}
Under normalized operator and RHS access, a QLS call produces a quantum state
proportional to a solution.  A hybrid QIPM that updates a classical iterate
must then estimate a full vector, recover its norm and scale, and ensure that
the resulting step satisfies the outer neighborhood condition.  Repeating the
state preparation for tomography repeats every nonpersistent oracle and RHS
cost.  A method that requests only an observable or compact certificate can
avoid full readout, but then it must show how that object is sufficient for
optimization and for constructing the next iteration.  The distinction among
classical-vector, compact, state, density, coordinate-oracle, and heralded
outputs in \cref{sec:prelim-output-contracts} is designed to prevent a proof
for one interface from being silently reused for another.

\paragraph{The classical comparator.}
For a sparse direct method, the per-step cost is controlled by realized fill,
numerical factorization, and triangular solves, with a symbolic ordering
normally reused.  Multiple predictor or corrector RHSs can reuse the same
factorization.  For a Krylov method, the relevant costs are matrix--vector and
preconditioner applications times a spectrum- and tolerance-dependent
iteration count.  Clustered eigenvalues can make that count far smaller than
a worst-case \(\sqrt\kappa\) estimate, as
\cref{prop:two-cluster-cg} demonstrates.  Specialized graph, treewidth,
block-angular, and low-rank structure can improve the baseline further.  We
therefore do not compare a quantum query count directly with dense cubic
arithmetic or treat the act of writing a dense classical direction as the
only classical cost.

\paragraph{Setup, refresh, and trajectory reuse.}
One-step costs do not automatically multiply by the IPM iteration count.  A
fixed sparse input can be read once, a symbolic factorization can persist,
and a lazy oracle can combine stored data with a changing scalar path
parameter.  Conversely, a QRAM or block encoding of changing primal and slack
diagonals may require per-step writes, and tomography may be repeated at every
step.  The distinction is formalized by the amortized accounting contract in
\cref{def:amortized-accounting}.  Our online direct-sum lower bound applies
only when new oracle information is genuinely fresh
\cref{thm:prec-online-direct-sum}; the caching-ceiling argument in
\cref{sec:preconditioning-lazy} shows why no unqualified
\(T\)-times-one-step claim follows for a static LP.

\paragraph{Query separation versus time separation.}
Raw coefficient, sparse-entry, row, block, QRAM, and preconditioned-product
queries reveal different amounts of information.  A sublinear row-query
algorithm can be a genuine oracle separation even when generating a queried
row or performing local arithmetic prevents a time advantage.  Conversely,
an explicitly stored sparse matrix may have no interesting input-query cost,
while its repeated numerical solves remain expensive.  The Apers--Gribling
result is positive in its declared row-query model; our parity and holonomy
lower bounds are negative in their declared raw local-query models.  Neither
statement changes models for the sake of the comparison.  This discipline is
why the later theorems often conclude with a setup--solve--recovery inequality
rather than with a single condition-number exponent.

\subsection{Answer in outline}
\label{sec:introduction-answer}

Our first answer is that central-path conditioning is geometry before it is
barrier design.  For a compact feasible slice, the diameter of an objective
sublevel set and the length of its longest chord control the extreme
eigenvalues of every self-concordant barrier Hessian.  The sublevel chord law
in \cref{thm:sublevel-chord-law} gives a condition-number floor at a fixed
objective gap, and \cref{thm:canonical-achievability} shows that the canonical
log and log-det barriers attain that floor up to explicit instance constants.
The conclusion is minimax and gap matched: changing the barrier cannot in
general remove the divergence.  It does not say that every optimum is ill
conditioned, nor does it cover a method that avoids barrier-Newton solves.

Our second answer is that a large condition number does not determine the
relative classical and quantum costs.  Degenerate LP log-barrier Hessians
split into two intervals with bounded internal ratios
\cref{prop:two-cluster-hessian}.  A classical residual polynomial can exploit
the empty interval, and \cref{prop:two-cluster-cg} gives polylogarithmic
iteration dependence on the inter-cluster separation.  Quantum polynomial
transformations must also remain bounded on the domain imposed by their
access model.  Under plain block-encoding access the worst-case QLS cost is
\(\Thetatilde(\kappa)\), even for two-point clusters when
\(N_{\rm QLS}\gtrsim\kappa\).  A norm-tight factor presentation
unconditionally has a \(\Thetatilde(\sqrt\kappa)\) polynomial route in its
model; only the corresponding end-to-end solve needs the overlap hypotheses
in \cref{thm:quantum-cluster-obstruction}.  On-path right-hand sides can be
much easier still: \cref{prop:newton-rhs-alignment,cor:top-window-solve} show
that the exact central-path right-hand side has \(O(\mu^2)\) weak-cluster
fraction, so a top-window solve meets the fixed relative-residual contract
once the spectral gap is resolved.

This classical--quantum contrast is not paradoxical.  CG is charged for
matrix--vector products and can place roots in spectral gaps.  QSVT is charged
for a bounded transformation of a supplied normalized operator, and the
normalization and symmetric-domain constraints remain part of the problem.
The comparison reverses the naive view that a large \(\kappa\) is always bad
for the classical method and automatically favorable to a quantum method.
It also motivates the instance-selection laws of
\cref{sec:beyond-kappa}: one should measure right-hand-side coupling and
spectral mass, not select a solver from \(\kappa\) alone.

Our third answer is that the decisive lower bounds often live outside the
linear solve.  The LP constructions of \cref{sec:lp-hardness} can make the
reduced Newton geometry condition one while hiding a parity-dependent affine
translation or original-coordinate direction.  The SDP constructions of
\cref{sec:sdp-hardness} separate intrinsic value information, Schur solves,
KKT access, and recovery.  The preconditioning results of
\cref{sec:preconditioning} show product laws: improving spectral condition
can worsen factor normalization or physical-state sensitivity, and a state
from one solve need not provide a coordinate oracle for the next.  The right
unit of analysis is therefore a complete setup--load--solve--recover--output
ledger.

Condensation supplies a complementary positive and negative example.  A
shared trace budget drives winner-take-all mass onto the lowest-cost PSD
block.  That mass simultaneously controls reduced-Hessian conditioning,
trace-distance visibility, and the query cost of state preparation.  The
transition has matching lower and upper query bounds on a bounded-incidence
holonomy SDP, while small-success search gives a tight accuracy curve on its
proved domain; the \(\delta=o(k/G)\) sliver and the internal-state
\(\eta_{\rm int}\)-channel remain open
\cref{sec:state-conversion-open}.  Thus output hardness is not inferred from
condition alone: it is proved through a visible, decodable output contract.

Finally, our no-go results are frontier theorems, not universal
impossibilities.  Convex mixtures, cut resistance, bounded-arity Walsh rows,
symmetry work identities, scalar congruences, and state-interface products
rule out the natural gadget classes to which their hypotheses apply.  The
section records the five escape routes in \cref{sec:frontier-synthesis}.  The
positive modules of
\cref{sec:positive-modules} inhabit some of those routes, but remain
conditional and per-step: they require stable geometry, explicit access,
compressed output, or a low-dimensional border.  No section silently upgrades
one such module to an end-to-end sparse-QIPM advantage.

\subsection{Contributions}
\label{sec:introduction-contributions}

The technical contributions follow the paper's order.  Each item names the
result that supports its headline and the boundary under which it should be
read.

\begin{enumerate}[leftmargin=*]
 \item \textbf{Universal conditioning geometry (\cref{sec:conditioning}).}
 We prove minimum-eigenvalue rigidity and the sublevel chord law for every
 self-concordant barrier
 \cref{thm:lambda-min-rigidity,thm:sublevel-chord-law}, sharpen it to a
 width-by-width spectral sandwich \cref{thm:width-spectrum-rigidity}, and
 show canonical LP and product-SDP barriers attain the resulting floor up to
 an explicit gap-matched factor \cref{thm:canonical-achievability}.  The
 conditioning dictionary includes positive-dimensional optimal faces,
 degenerate vertices, curved optima, singularity-degree exponents, and
 approximate centrality
 \cref{cor:conditioning-regimes,cor:fractional-conditioning-laws,cor:near-degeneracy-plateau,cor:approximate-centrality-transfer}.

 \item \textbf{Clustered spectra and solver separation
 (\cref{sec:spectra}).}  We prove the two-cluster LP tail
 \cref{prop:two-cluster-hessian}, the cluster-sensitive CG bound
 \cref{prop:two-cluster-cg}, and the plain-versus-factor quantum polynomial
 frontier \cref{thm:quantum-cluster-obstruction}.  The exact central-path
 right-hand side obeys a weak-mode cancellation law, yielding a
 condition-number-free top-window solve under the fixed relative-residual
 contract once the spectral gap is resolved
 \cref{prop:newton-rhs-alignment,cor:top-window-solve}.  The generic
 \(\Thetatilde(\sqrt\kappa)\) factor-access QLS interpretation is explicitly
 rejected: the theorem gives a polynomial metric and an overlap-qualified
 solve, not such an unrestricted algorithm.

 \item \textbf{Instance laws beyond \(\kappa\)
 (\cref{sec:beyond-kappa}).}  Exact block inversion identifies when DLS
 filtering improves on worst-case conditioning
 \cref{thm:dls-coupling}; a truncation transition separates solution-state
 and residual contracts \cref{prop:dls-truncation-transition}; and the
 effective-spectral-dimension theorem places truncation and filtering
 thresholds at dimensions four and eight
 \cref{thm:spectral-dimension-laws}.  Stable active-range equilibration
 preserves the Newton state for \(f\in\operatorname{range}Q\) while removing
 the divergent tail when its projector is supplied
 \cref{thm:stable-split-equilibration}.

 \item \textbf{Block log-det condensation
 (\cref{sec:condensation}).}  We derive the exact capacity transition and
 ground-multiplicity-dependent conditioning phases
 \cref{thm:condensation-threshold,thm:cond-conditioning-phases}, prove
 finite mass--conditioning laws \cref{thm:cond-mass-conditioning}, and link
 them to trace-distance visibility and raw-query complexity
 \cref{thm:cond-visibility,thm:cond-unique-or}.  The transition extends to
 sublinear multiwinner sets and to a stated approximate-centrality model
 \cref{thm:cond-multiwinner,thm:cond-robust}.  These are statements about the
 equality-reduced Frobenius Hessian and trace-normalized output, not every
 preconditioned KKT representation.

 \item \textbf{A sparse winner-take-all SDP
 (\cref{sec:winner-take-all}).}  A bounded-incidence qutrit-holonomy family
 realizes the condensation law with tight \(\Theta(N\sqrt G)\) quantum and
 \(\Theta(NG)\) randomized value-query costs, a condition-one public root
 start, robust winner recovery \cref{thm:wta-robust-optimizer}, and a matching
 search-then-solve compact-output algorithm
 \cref{thm:wta-search-upper}.  The formulation tradeoff in
 \cref{sec:wta-scope} keeps the bounded-incidence claim separate from an
 isometric ambient conditioning claim.

 \item \textbf{Accuracy curves for hidden-block state conversion
 (\cref{sec:state-conversion}).}  A verified-sample decoder gives a
 mixed-output lower bound \cref{thm:sc-m1}; detuned rejection preparation
 supplies the upper bound \cref{thm:sc-m2}; and parity collapse yields the
 tight nontrivial conversion curve
 \cref{thm:sc-small-success,cor:sc-s2}.  The batch results separate
 nonamortizable vector materialization from scalar approximate counting
 \cref{thm:sc-batch-vector,thm:sc-batch-scalar}.  The theorem domains retain
 the loser-state discrepancy, the small-success sliver, and logarithmic
 fixed-point costs rather than extrapolating through them.

 \item \textbf{LP loading and recovery hardness
 (\cref{sec:lp-hardness}).}  The gain--plateau family gives a linear robust
 primal-state lower bound \cref{thm:lp-gain-plateau}.  Prefix rigidification
 strengthens this to condition-one reduced Newton geometry with linear
 original-coordinate recovery cost \cref{thm:lp-prefix-capstone}.  Complete
 KKT and Schur-preconditioned variants isolate public-right-hand-side,
 residual, bounded-data, and scale tradeoffs
 \cref{thm:lp-full-kkt,thm:lp-preconditioned-kkt,thm:lp-bounded-exact,thm:lp-unit-threshold,thm:lp-dual-homogeneous-full}.
 None is presented as a condition-one QLS
 lower bound with free transformed oracles.

 \item \textbf{Gadget-class frontiers
 (\cref{sec:gadget-frontiers}).}  A single-flip convex-mixture theorem gives
 the graph-free default obstruction \cref{thm:frontier-single-flip}.
 Cut--resistance and Walsh-row theorems quantify the volume or scale needed
 by signed-copy and bounded-arity constructions
 \cref{thm:frontier-cut-resistance,thm:frontier-walsh}; gauge-flow, work, and
 Fourier identities locate the multiplier or input-incidence cost
 \cref{thm:frontier-gauge-flow,thm:frontier-work,thm:frontier-fourier-source}.
 The escape routes in
 \cref{sec:frontier-synthesis} are part of the result.

 \item \textbf{SDP-native loading and recovery hardness
 (\cref{sec:sdp-hardness}).}  An intrinsic trace-normalized qutrit family has
 parity-hard value and a public, condition-one reduced Newton step whose
 primal and complete-KKT solution states are parity hard
 \cref{thm:sdp-intrinsic,thm:sdp-newton-state}.  Non-SOCP Schur families give
 an exact loading/accuracy frontier
 \cref{thm:sdp-schur,thm:sdp-accuracy-frontier}, and tree and group holonomy
 lift the mechanism to matrix-valued words
 \cref{thm:sdp-tree-loading,thm:sdp-holonomy-spectrum,thm:sdp-s3-hardness}.
 The proved query exponent for the \(S_3\) family is
 still obtained on an abelian restriction; intrinsically nonabelian word
 hardness remains open.

 \item \textbf{Preconditioning and state-interface products
 (\cref{sec:preconditioning}).}  Whitening and reusable inverse access obey
 normalization tradeoffs
 \cref{thm:prec-whitening-frontier,thm:prec-inverse-access}; scalar
 congruence trades transformed condition against physical recovery
 \cref{thm:prec-congruence-frontier,thm:prec-general-scalar-frontier}; and
 the parity ledger charges setup, RHS preparation, solve, and recovery
 \cref{thm:prec-factor-rhs-hardness,thm:prec-parity-ledger}.  State-to-next-solve
 and lazy-oracle theorems distinguish a one-use state from a reusable
 coordinate service
 \cref{thm:prec-state-diagonal,thm:prec-state-next-solve,thm:prec-lazy-lookup,thm:prec-reusable-endpoints}.

 \item \textbf{Conditional positive modules
 (\cref{sec:positive-modules}).}  We prove a residual-certified pathwise
 refresh bound \cref{thm:positive-lazy-refresh}, a constructive stable
 active-range equilibration and exact face repair
 \cref{thm:positive-partition-free-equilibration,thm:positive-face-repair}, an
 implicit slack oracle and conditional
 compressed dual QIPM
 \cref{thm:positive-implicit-slack,thm:positive-compressed-qipm}, and a
 block-angular border solver whose \(\sqrt N\) block-count exponent is
 optimal in the block-oracle model, as proved in
 \cref{sec:positive-block-angular}.  Every theorem keeps its tail, gap,
 forcing, normalization, and output assumptions visible.

 \item \textbf{Reusable structural tools
 (\cref{sec:structural-tools}).}  Involution and compact-group work
 identities \cref{thm:structural-involution-work,thm:structural-group-work},
 a noncommutative central-mixture sandwich
 \cref{thm:structural-sdp-kantorovich}, and rank and support barriers
 \cref{thm:structural-rank-progress,thm:structural-support-barrier} explain
 several no-go mechanisms abstractly.  Sparse OSS and treewidth-two KKT
 embeddings preserve known QLS restrictions
 \cref{thm:structural-oss-embedding,thm:structural-treewidth-two}, while
 exact-treewidth value hardness and a central-oracle light cone distinguish
 query, output, and physical-locality models
 \cref{thm:structural-exact-treewidth,thm:structural-light-cone}.

 \item \textbf{Corrections and dead ends (\cref{sec:corrections}).}  We show
 that the global potential supremum used in the QCPM simulator ledger cannot
 scale as asserted \cref{prop:corr-qcpm-supremum}, correct the independent
 clock/time-dilation error \cref{eq:corr-qcpm-clock}, derive the corrected norm
 scale and its speed tradeoff
 \cref{thm:corr-qcpm-norm,cor:corr-qcpm-tradeoff}, and give explicit
 densification and hidden-start counterexamples
 \cref{thm:corr-path-simplex,thm:corr-hidden-sign}.  The recorded dead ends
 prevent first-inverse norms, sparse metrics, exact Schur products, or
 one-step bounds from being reused as claims they do not prove.
\end{enumerate}

\subsection{Related work and positioning}
\label{sec:introduction-related-work}

\paragraph{Quantum interior-point methods.}
The QIPM literature spans QLS-based hybrid path following, iterative
refinement, orthogonal-subspace formulations, condition-number-free sampling,
and Hamiltonian central paths.  Representative lines include
Kerenidis--Prakash, Casares--Martin-Delgado, Augustino et al.,
Mohammadisiahroudi et al.,
Wu et al., Apers--Gribling, and the reviews and resource ledgers of Dalzell et
al.
\cite{kerenidis2020-quantum-interior-point-method-lps,
casares2020-quantum-interior-point-predictor-corrector,
augustino2023-quantum-interior-point-methods-semidefinite,
mohammadisiahroudi2025-inexact-feasible-interior-point-method,
wu2023-inexact-feasible-quantum-interior-point,
apers2026-quantum-speedups-linear-programming-interior,
dalzell2025-quantum-interior-point-methods}.  The QCPM and
self-concordant-Schr\"odinger lines avoid the usual repeated-QLS architecture
\cite{augustino2024-quantum-central-path-algorithm-linear,
gribling2025-self-concordant-schrodinger-operators-spectral}.  We do not
collapse these algorithms into a single complexity table: each uses a
different access, condition, precision, and output contract.  The per-section
novelty paragraphs compare the exact theorem at hand with its nearest line of
work.

\paragraph{Quantum linear systems.}
HHL introduced sparse QLS state preparation, CKS improved the precision
dependence, and QSVT placed inversion in a general block-encoding calculus
\cite{harrow2009-quantum-algorithm-linear-systems-equations,
andrew2017-quantum-algorithm-systems-linear-equations,
andras2019-quantum-singular-value-transformation-beyond}.  Orsucci--Dunjko's
positive-definite lower bounds ground our plain-access clustered comparison.
Mori et al. prove \(\Omega(\kappa\sqrt{s})\) at fixed constant error and
restate an \(\Omega(\kappa\log(1/\epsilon))\) constant-sparsity bound from
earlier unpublished work of Harrow and Kothari.  DLS provide the
beyond-condition-number parameters evaluated in
\cref{sec:beyond-kappa}
\cite{orsucci2021-on-solving-classes-of-positive,
mori2026-sparsity-dependent-complexity-lower-bound,
dalzell2026-faster-quantum-linear-system-solver}.  Somma--de Wolf provide an
analogous plain-versus-sum-of-squares resolution separation in a guided
Hamiltonian problem, not a QLS theorem
\cite{somma2026-optimal-lower-bound-ground-state-energy}.  This access-model
discipline is essential to \cref{sec:spectra,sec:preconditioning}.

\paragraph{Quantum optimization lower bounds.}
Van Apeldoorn et al. establish general quantum SDP upper and lower bounds and
Boolean-in-optimum constructions; Brand\~ao--Svore initiated the quantum SDP
solver line; and Apers--Gribling's Theorem~8.4 gives a sparse-LP value lower
bound already having a linear exponent in balanced dimensions
\cite{apeldoorn2020-quantum-sdp-solvers,
fernando2017-quantum-speed-ups-solving-semidefinite,
apers2026-quantum-speedups-linear-programming-interior}.  Our parity and
holonomy families do not claim parity, adversary composition, or the linear
exponent as new.  Their contribution is the conjunction with exact central
paths, condition-one or bounded-condition formulations, robust output
contracts, and complete loading/recovery ledgers described in
\cref{sec:winner-take-all,sec:lp-hardness,sec:sdp-hardness}.

\paragraph{Classical interior-point structure.}
Wright's work on logarithmic-barrier Hessians and numerical error, the
standard primal--dual theory, and the Forsgren--Gill--Wright survey provide
the conditioning and implementation background
\cite{wright1994-properties-hessian-logarithmic-barrier,
wright1998-ill-conditioning-computational-error-interior,
wright1997-primal-dual-interior-point-methods,
forsgren2002-interior-methods-nonlinear-optimization}.  Gondzio and
collaborators document preconditioning, structure exploitation, and practical
IPM evolution
\cite{jacek2012-interior-point-methods-25-years,
gondzio2023-general-purpose-preconditioning-regularized-interior-point}.
Small-treewidth LP and SDP algorithms show why bounded graph width is a
classical algorithmic asset rather than merely a sparse-oracle promise
\cite{dong2020-a-nearly-linear-time-algorithm,
gu2022-a-faster-small-treewidth-sdp}.  The sublevel chord law differs from the
classical log-barrier conditioning analyses by quantifying over all
self-concordant barriers, while \cref{sec:spectra} uses rather than ignores
the classical cluster-sensitive Krylov advantage.

\paragraph{Positioning discipline.}
Many ingredients used below are established: self-concordant containment,
weak-sharp and singularity-degree error bounds, matrix-Burg resolvents,
Rayleigh--Jeans condensation, parity and search, group synchronization,
Kantorovich inequalities, effective resistance, and sparse elimination.  We
credit them where each enters.  Because the nearest prior work changes from
one theorem package to the next, this introduction does not restate every
collision or priority qualification.  The detailed novelty and scope
statements in
\cref{sec:conditioning-overview,sec:condensation-discussion,sec:wta-scope,sec:sdp-synthesis,sec:preconditioning-synthesis,sec:positive-synthesis,sec:structural-synthesis}
are part of the claims.

\subsection{Conventions, verification, and corrections}
\label{sec:introduction-verification}

The distinctions formalized in \cref{sec:preliminaries} are standing
conventions, not optional caveats.  In particular, input sparsity differs from
Newton-system sparsity; \(\kappa_H\) differs from block-encoding normalization;
reduced and ambient KKT conditioning are different quantities; and value,
classical-vector, quantum-state, coordinate-oracle, and heralded outputs are
different contracts.  Query bounds count the setup and transformations named
by their access model.  An uncharged stronger oracle is treated as a stronger
problem statement, not as an implementation of raw sparse access.

The verification process changed several statements inherited from
preliminary research notes and identified two independent issues in a public
preprint.  The corrections relevant
to reading the paper are:
\begin{enumerate}[leftmargin=*]
 \item \cref{sec:corrections-qcpm} corrects the QCPM v2 simulator-norm
 estimate and its independent clock/time-dilation error: a global supremum on
 a fixed truncation box is not controlled by concentration near the
 instantaneous center, and the time-varying kinetic coefficient requires the
 corrected clock in \cref{eq:corr-qcpm-clock}.
 \Cref{thm:corr-qcpm-norm,cor:corr-qcpm-tradeoff} give the resulting scale.
 This is a correction to that derivation, not a general impossibility result
 for Hamiltonian optimization.
 \item \cref{sec:spectra-quantum} corrects an overbroad factor-access
 inference in our earlier analysis.  A square-root polynomial metric, or a
 two-stage solve with explicit overlap, does not imply a generic
 \(\Thetatilde(\sqrt\kappa)\) factor-access QLS; the Orsucci--Dunjko instance
 rules out that unrestricted interpretation.
 \item \cref{sec:spectra-two-cluster} replaces a coarse three-scale picture
 of the degree-two SDP witness by the verified four-window spectrum.  The
 correction changes the intermediate spectral description, not the
 condition-number exponent.
 \item \cref{sec:effective-spectral-dimension} anchors effective spectral
 dimension at the smallest visible eigenvalue.  This repairs the earlier
 shell formulation and is needed for the dimension-four and dimension-eight
 phase laws in \cref{thm:spectral-dimension-laws}.
 \item \cref{sec:state-conversion-model,sec:state-conversion-small-success,sec:state-conversion-open}
 replace the collision shortcut by a verified-sample decoder, restrict the
 classical linear-in-success law to its actual domain, keep the
 very-small-success sliver open, and attribute the \(k\)-marked small-success
 search bound to Dohotaru--H\o yer rather than to Zalka's one-marked theorem
 alone.
 \item \cref{sec:positive-face-repair} corrects the earlier claim that two
 convergent leakage ledgers alone bound projective variation.  Neither the
 ledgers nor an additional $O(\mu_t)$ tail bound suffices: an
 oscillating-face-displacement counterexample has convergent ledgers, the tail
 bound, and linearly growing realized variation, so finite realized variation
 stays a separate hypothesis.
\end{enumerate}
The list is deliberately narrow.  It records changes that affect the theorem
statements or their interpretation, rather than presenting ordinary proof
repairs as additional results.  Further failed shortcuts and their valid
replacements are collected in \cref{sec:corrections-dead-ends}.

\subsection{Organization and reading paths}
\label{sec:introduction-organization}

The conventions are fixed in \cref{sec:preliminaries}; thereafter we use the
section pointers rather than redefining those contracts.  Conditioning
geometry, clustered spectra, and beyond-condition-number instance laws occupy
\cref{sec:conditioning,sec:spectra,sec:beyond-kappa}.  Condensation, its sparse
winner-take-all realization, and the associated state-conversion accuracy
curve occupy \cref{sec:condensation,sec:winner-take-all,sec:state-conversion}.
The LP and SDP loading/recovery lower bounds, gadget frontiers, and
preconditioning/interface obstructions are in
\cref{sec:lp-hardness,sec:gadget-frontiers,sec:sdp-hardness,sec:preconditioning}.
The conditional positive modules are in
\cref{sec:positive-modules}, structural tools in \cref{sec:structural-tools},
corrections and dead ends in \cref{sec:corrections}, and the combined verdict
and open frontier in \cref{sec:discussion}.

Three shorter paths may be useful.  A reader interested only in conditioning
can read
\cref{sec:preliminaries,sec:conditioning,sec:spectra,sec:beyond-kappa}.  A
reader interested only in lower bounds can read
\cref{sec:preliminaries,sec:condensation,sec:winner-take-all,sec:state-conversion,sec:lp-hardness,sec:gadget-frontiers,sec:sdp-hardness,sec:preconditioning}.
A reader interested only in algorithmic
modules can read
\cref{sec:preliminaries,sec:beyond-kappa,sec:positive-modules}.  In every
path, the output and access definitions of \cref{sec:preliminaries} remain
part of the theorem statements.

\section{Preliminaries}
\label{sec:preliminaries}

This section fixes the optimization, access, and output models used throughout
the paper.  The distinctions are part of the problem specification: input
sparsity is not Newton-system sparsity, the spectral condition number is not a
block-encoding normalization, and a normalized quantum state is not an
explicit solution vector.

\subsection{Conic programs and central paths}
\label{sec:prelim-conic}

For a linear program (LP) in standard form, the primal--dual pair is
\begin{equation}
 \begin{aligned}
  \text{(P$_{\rm LP}$)}\qquad
  \min_{x\in\R^n}\quad &c^\top x
       &\text{subject to}\quad& Ax=b,\quad x\geq 0,\\
  \text{(D$_{\rm LP}$)}\qquad
  \max_{y\in\R^m,\,s\in\R^n}\quad &b^\top y
       &\text{subject to}\quad& A^\top y+s=c,\quad s\geq 0,
 \end{aligned}
 \label{eq:lp-pair}
\end{equation}
Throughout the LP discussion, $A\in\R^{m\times n}$ has full row rank unless
stated otherwise.  We write
$X=\diag(x)$, $S=\diag(s)$, and $e$ for the all-ones vector.  A feasible
primal--dual point has gap
\begin{equation}
 c^\top x-b^\top y=x^\top s.
 \label{eq:lp-pd-gap}
\end{equation}

The product-cone form of a semidefinite program (SDP) used below has
$q$ matrix blocks.  Put
$\mathcal K=\prod_{j=1}^q\Splus^{d_j}$ and equip it with
$\inner{X}{S}=\sum_j\tr(X_jS_j)$.  For symmetric data
$C_j,A_{ij}\in\Sym^{d_j}$, define
$\mathcal A(X)_i=\sum_j\tr(A_{ij}X_j)$ and
$(\mathcal A^*y)_j=\sum_i y_iA_{ij}$.  The pair is
\begin{equation}
 \begin{aligned}
  \text{(P$_{\rm SDP}$)}\qquad
  \min_{X\in\mathcal K}\quad &\inner{C}{X}
       &\text{subject to}\quad& \mathcal A(X)=b,\\
  \text{(D$_{\rm SDP}$)}\qquad
  \max_{y\in\R^m,\,S\in\mathcal K}\quad &b^\top y
       &\text{subject to}\quad& \mathcal A^*y+S=C.
 \end{aligned}
 \label{eq:sdp-pair}
\end{equation}
Its primal--dual gap is $\inner{C}{X}-b^\top y=\inner{X}{S}$.
The one-block formulation is the special case $q=1$; scalar nonnegative
variables can equivalently be included as $1\times1$ blocks.  These are the
standard conic dual pairs under the trace inner product
\cite{boyd2004-convex-optimization}.

\begin{definition}[Strict feasibility]
\label{def:strict-feasibility}
An LP is \emph{strictly primal--dual feasible} if there are points satisfying
\eqref{eq:lp-pair} with $x>0$ and $s>0$.  An SDP is strictly primal--dual
feasible if there are feasible points with $X_j\succ0$ and $S_j\succ0$ for
every block.  We use \emph{strict feasibility} for this two-sided condition
unless one side is named.  Because both sides are then feasible, their optimal
values are finite; the two-sided Slater condition gives zero duality gap and
attainment on both sides \cite{boyd2004-convex-optimization}.
\end{definition}

We next recall the barrier terminology in the normalization used by
Nesterov and Nemirovskii \cite{nesterov1994-interior-point-polynomial-algorithms-convex}.
For a $C^3$ convex function $F$ on an open convex domain, define the local
norms
\begin{equation}
 \norm{h}_x=\sqrt{D^2F(x)[h,h]}.
 \label{eq:local-norms}
\end{equation}
On an affine feasible slice, derivatives and Hessians in
\eqref{eq:local-norms} are restricted to its tangent space.

\begin{definition}[$\nu$-self-concordant barrier]
\label{def:self-concordant-barrier}
A function $F$ is a \emph{$\nu$-self-concordant barrier} for a convex set
$Q$ if $\operatorname{dom}F=\operatorname{relint}Q$, it tends to $+\infty$
at the relative boundary, its Hessian is positive definite on the tangent
space, and, for all $x\in\operatorname{relint}Q$ and tangent directions $h$,
\begin{equation}
 |D^3F(x)[h,h,h]|\leq2\norm{h}_x^3,
 \qquad
 |DF(x)[h]|\leq\sqrt\nu\,\norm{h}_x.
 \label{eq:sc-barrier}
\end{equation}
The declared value $\nu$ is called a barrier parameter; it need not be the
least admissible value.
\end{definition}

The logarithmic barrier for the nonnegative orthant and the log-determinant
barrier for the product PSD cone are, respectively,
\begin{equation}
 F_{\log}(x)=-\sum_{i=1}^n\log x_i,
 \qquad
 F_{\log\det}(X)=-\sum_{j=1}^q\log\det X_j.
 \label{eq:canonical-barriers}
\end{equation}
They have parameters $\nu=n$ and $\nu=\sum_jd_j$.  Both are logarithmically
homogeneous: $F(tx)=F(x)-\nu\log t$ in the appropriate cone notation.  The
symbol $\nu$ always denotes this cone barrier's parameter, which for these
canonical barriers equals the degree of logarithmic homogeneity.  Restricting
the cone barrier to an affine feasible slice preserves self-concordance with
parameter at most $\nu$, possibly with a smaller least parameter; we do not
redefine $\nu$ after that restriction.  The cone barrier parameter, rather
than the ambient number of stored coefficients, controls the standard
path-following iteration bound
\cite{nesterov1994-interior-point-polynomial-algorithms-convex}.

\begin{definition}[Central path and its two parametrizations]
\label{def:central-path}
Let $F$ be a self-concordant barrier on the relative interior of the primal
feasible region, suppose $\mathrm{OPT}$ is finite and attained, and suppose
the barrier subproblem has a unique minimizer.
The \emph{$\mu$-parametrized primal central path} is
\begin{equation}
 x_F(\mu)=\arg\min\{c^\top x+\mu F(x):Ax=b,\ x\in\operatorname{dom}F\},
 \qquad \mu>0.
 \label{eq:general-central-path}
\end{equation}
Its primal objective gap is
\begin{equation}
 g_F(\mu)=c^\top x_F(\mu)-\mathrm{OPT}.
 \label{eq:objective-gap-g}
\end{equation}
On any interval where $g_F$ is one-to-one, the same curve indexed by
$g>0$ is the \emph{$g$-parametrized central path}, denoted $x_F(g)$.
For the canonical LP barrier, the primal--dual central path additionally
satisfies
\begin{equation}
 Ax=b,\quad A^\top y+s=c,\quad x,s>0,\quad XSe=\mu e.
 \label{eq:lp-central-path}
\end{equation}
For the product-PSD log-det barrier it satisfies
\begin{equation}
 \mathcal A(X)=b,\quad \mathcal A^*y+S=C,\quad
 X_j,S_j\succ0,\quad S_j=\mu X_j^{-1}\quad(1\leq j\leq q).
 \label{eq:sdp-central-path}
\end{equation}
Thus the primal--dual gap on either canonical path is exactly $\nu\mu$, with
$\nu$ the cone barrier parameter fixed above.
\end{definition}

The first-order condition in \eqref{eq:general-central-path} is the source of
the dual slack relation $s=-\mu\nabla F(x)$ whenever that gradient is
identified with a cone dual variable.  For a general self-concordant barrier,
$g_F(\mu)\leq\nu\mu$; equality need not hold
\cite{nesterov1994-interior-point-polynomial-algorithms-convex}.

\begin{definition}[Trace-normalized log-det center]
\label{def:trace-logdet-center}
On the affine slice $\sum_j\tr X_j=1$, the \emph{trace-normalized log-det
center} at parameter $\mu$ is the minimizer of
\begin{equation}
 \inner{C}{X}-\mu\sum_j\log\det X_j,
 \qquad X_j\succ0,\quad \sum_j\tr X_j=1.
 \label{eq:trace-logdet-center}
\end{equation}
With the Lagrangian sign convention
$\inner{C}{X}-\mu\sum_j\log\det X_j-
\lambda(\sum_j\tr X_j-1)$, stationarity reads
$X_j=\mu(C_j-\lambda I)^{-1}$.  This expression is admissible only for
$\lambda<\min_j\lambda_{\min}(C_j)$, so every inverse is positive definite.
The corresponding log-det Hessian is
$D^2F(X)[U,V]=\sum_j\tr(X_j^{-1}U_jX_j^{-1}V_j)$
\cite{boyd2004-convex-optimization}.
\end{definition}

This is a specialization of the product-PSD central path, not a different
optimization model.  Trace normalization also makes
$\bigoplus_jX_j$ a density operator, connecting the conic and quantum output
conventions below.

\begin{remark}[The convention for $g$]
\label{rem:g-convention}
The conditioning sections use $g$ only for the primal objective gap
\eqref{eq:objective-gap-g}.  Although an objective gap is sometimes called a
``duality gap'' informally, here \emph{primal--dual gap} means
$c^\top x-b^\top y=x^\top s$, which equals $\nu\mu$ on the canonical LP
central path.  In general $g_F(\mu)\leq\nu\mu$, and equality need not hold.
\end{remark}

\begin{definition}[Analytic center]
\label{def:analytic-center}
If $F$ attains a minimum over the relative interior of a feasible convex set,
its minimizer is the \emph{$F$-analytic center}.  With the canonical barriers
we simply say \emph{analytic center}.  For a bounded feasible region, the
large-$\mu$ limit of \eqref{eq:general-central-path}, when it exists, is this
center.  For a face on the boundary of the original domain, ``analytic center
of the face'' means the minimizer of a barrier defined on that face's relative
interior; it does not mean evaluating the original barrier at a boundary
point \cite{nesterov1994-interior-point-polynomial-algorithms-convex}.
\end{definition}

For an LP point satisfying the two feasibility equations, put
$\bar\mu=x^\top s/n$.  The narrow neighborhood used by short-step analyses is
the following standard one \cite{wright1997-primal-dual-interior-point-methods}.

\begin{definition}[$N_2(\theta)$ neighborhood]
\label{def:n2-neighborhood}
For $0<\theta<1$, the \emph{$N_2(\theta)$ neighborhood} of the LP central path
is
\begin{equation}
 N_2(\theta)=\left\{(x,y,s):
 \begin{array}{l}
 Ax=b,\ A^\top y+s=c,\ x,s>0,\quad \bar\mu=x^\top s/n,\\[-2pt]
 \norm{XSe-\bar\mu e}_2\leq\theta\bar\mu
 \end{array}\right\}.
 \label{eq:n2-neighborhood}
\end{equation}
Unless an infeasible or matrix-valued analogue is explicitly named,
$N_2(\theta)$ always means \eqref{eq:n2-neighborhood}.
\end{definition}

A short-step path-following method keeps a fixed small $\theta$, changes the
target complementarity by a $1-\Theta(1/\sqrt\nu)$ factor, and takes a full Newton
step back into the neighborhood.  This gives
$O(\sqrt\nu\log(\nu/\epsilon_{\rm opt}))$ Newton steps under standard
initialization and scaling conventions; this statement is an iteration count,
not a linear-solver, data-loading, or readout cost
\cite{nesterov1994-interior-point-polynomial-algorithms-convex,wright1997-primal-dual-interior-point-methods}.

\subsubsection{Newton formulations and conditioning}
\label{sec:newton-formulations}

At a positive, possibly infeasible LP iterate, define residuals
\begin{equation}
 r_p=b-Ax,\qquad r_d=c-A^\top y-s,\qquad
 r_c=\widehat\mu e-XSe,
 \label{eq:newton-residuals}
\end{equation}
where frequently $\widehat\mu=\sigma\bar\mu$ for a centering parameter
$\sigma\in[0,1]$.  Linearizing primal feasibility, dual feasibility, and
complementarity gives the full primal--dual KKT system
\begin{equation}
 \begin{bmatrix}
  A&0&0\\
  0&A^\top&I\\
  S&0&X
 \end{bmatrix}
 \begin{bmatrix}\Delta x\\\Delta y\\\Delta s\end{bmatrix}
 =
 \begin{bmatrix}r_p\\r_d\\r_c\end{bmatrix}.
 \label{eq:full-kkt}
\end{equation}
This equation, with the residual signs in \eqref{eq:newton-residuals}, is what
we call the \emph{standard infeasible-start Newton correction}; a
feasible-start direction has $r_p=r_d=0$
\cite{wright1997-primal-dual-interior-point-methods}.  A full step eliminates
the affine primal and dual residuals exactly, but complementarity is targeted
only to first order.  After a full step satisfying the third block row exactly,
$(X+\Delta X)(S+\Delta S)e=\widehat\mu e+\Delta X\Delta S e$.
The quadratic term is uncontrolled by the linear system alone, so progress
also requires a neighborhood invariant and a step-size rule.

Eliminating $\Delta x$ and $\Delta s$ gives the normal-equation coefficient
matrix
\begin{equation}
 H_{\rm NE}=AXS^{-1}A^\top
 \label{eq:normal-matrix}
\end{equation}
and the equation
\begin{equation}
 H_{\rm NE}\Delta y
 =r_p-AS^{-1}r_c+AXS^{-1}r_d.
 \label{eq:normal-equations}
\end{equation}
Under the standing rank assumption, it is positive definite for $x,s>0$.  Even when
$A$ is sparse, $H_{\rm NE}$ can be dense because two rows are coupled whenever
they share a nonzero column.  Accordingly, $\nnz(A)$ and the maximum
row/column sparsity $s_H$ of the matrix actually sent to a solver are distinct
parameters.

Two reformulations used later move an inexact-solve residual to a more useful
part of the Newton equations.  At a feasible iterate, if
$V\in\R^{n\times(n-m)}$ has full column rank
and $\operatorname{range}V=\ker A$, a feasible direction can be represented as
$\Delta x=V\lambda$ and $\Delta s=-A^\top\Delta y$.  Substitution in the third
row of \eqref{eq:full-kkt} gives the \emph{orthogonal-subspaces system (OSS)}
\begin{equation}
 \underbrace{[-XA^\top\ \ SV]}_{H_{\rm OSS}}
 \begin{bmatrix}\Delta y\\\lambda\end{bmatrix}=r_c.
 \label{eq:oss}
\end{equation}
Every approximate coefficient vector still reconstructs directions satisfying
$A\Delta x=0$ and $A^\top\Delta y+\Delta s=0$ exactly.  This is the
\emph{OSS} introduced by Mohammadisiahroudi et al.
\cite{mohammadisiahroudi2025-inexact-feasible-interior-point-method}.

The \emph{modified normal-equation system (MNES)} instead chooses an index
set $B$ for which $A_B\in\R^{m\times m}$ is nonsingular.  With
$D=X^{1/2}S^{-1/2}$, its principal submatrix $D_B$, $\widehat A=A_B^{-1}A$, and
$E=D_B^{-1}\widehat A D$, its coefficient and right-hand side are
\begin{equation}
 H_{\rm MNES}=EE^\top
 =D_B^{-1}A_B^{-1}H_{\rm NE}(D_B^{-1}A_B^{-1})^\top,
 \qquad h_{\rm MNES}=D_B^{-1}A_B^{-1}h_{\rm NE},
 \label{eq:mnes}
\end{equation}
where $h_{\rm NE}$ denotes the right-hand side of
\eqref{eq:normal-equations}.  If an approximate solution has residual
$\widehat r$, the reconstruction uses the correction
$v_B=D_B\widehat r$, $v_{\bar B}=0$; the primal and dual feasibility rows are
then exact, while the residual $-Sv$ is confined to the complementarity row
\cite{mohammadisiahroudi2024-efficient-use-quantum-linear-system}.  The names
OSS and MNES always mean the formulations \eqref{eq:oss} and
\eqref{eq:mnes}.  Neither name implies that $V$, $\widehat A$, or the
resulting system is sparse; that property must be established for the
instance at hand.

\begin{definition}[OSS/MNES inexact Newton residual convention]
\label{def:inexact-newton-residual}
Let $\eta\in[0,1)$ be the named forcing tolerance.  An inexact OSS solve uses
the linear-system residual $r_{\rm OSS}=H_{\rm OSS}\widetilde z-r_c$ and the
convention $\norm{r_{\rm OSS}}_2\leq\eta\bar\mu$.  An inexact MNES solve uses
$\widehat r=H_{\rm MNES}\widetilde z-h_{\rm MNES}$ and the
$\sqrt{\bar\mu}$-scaled convention
\begin{equation}
 \norm{\widehat r}_2\leq\eta\sqrt{\bar\mu/n},
 \qquad\text{which ensures}\qquad
 \norm{Sv}_\infty\leq\eta\bar\mu.
 \label{eq:inexact-newton-residual}
\end{equation}
These are residuals of the named Newton formulation, not the global
RHS-relative feasibility residual of Definition~\ref{def:global-residual}.
The formulation-specific scaling is part of the contract
\cite{mohammadisiahroudi2025-inexact-feasible-interior-point-method,mohammadisiahroudi2024-efficient-use-quantum-linear-system}.
\end{definition}

For primal-barrier geometry, let $x_F(\mu)$ lie in the affine slice
$x^0+\ker A$, and let $W$ have orthonormal columns spanning $\ker A$.  Its
reduced Hessian is
\begin{equation}
 H_{\rm red}(\mu)=W^\top\nabla^2F(x_F(\mu))W.
 \label{eq:reduced-hessian}
\end{equation}

\begin{definition}[Reduced central-path Hessian condition number]
\label{def:reduced-kappa}
The \emph{reduced central-path Hessian condition number} is
\begin{equation}
 \kappa_{\rm red}(\mu;F)
 =\frac{\lambda_{\max}(H_{\rm red}(\mu))}
        {\lambda_{\min}(H_{\rm red}(\mu))}.
 \label{eq:reduced-kappa}
\end{equation}
The Euclidean metric and an orthonormal tangent basis are part of this
definition.  The value is invariant under replacing $W$ by another
orthonormal basis, but not under an arbitrary nonorthogonal reparametrization.
For the LP log barrier, $\nabla^2F(x)=X^{-2}$.
\end{definition}

\begin{definition}[Ambient and KKT condition numbers]
\label{def:ambient-kkt-kappa}
For the named ambient cone barrier $F$, the \emph{ambient barrier-Hessian
condition number} is
$\kappa_{\rm amb}(x;F)=\kappa_2(\nabla^2F(x))$ before restriction to the
feasible tangent space.  It is representation-dependent: changing the
ambient cone representation or its coordinates changes the named $F$ and may
change $\kappa_{\rm amb}$.  For a named nonsingular Newton or saddle matrix $K$
in a specified scaling, its \emph{KKT condition number} is
$\kappa_2(K)=\sigma_{\max}(K)/\sigma_{\min}(K)$.  We write $\kappa_H$ for the
condition number of the actual linear-system matrix $H$ being solved.
\end{definition}

These quantities are not interchangeable.  Eliminating variables, changing
row scales, choosing a nonorthogonal null-space basis, or passing from an
indefinite KKT matrix to normal equations can change singular values and can
square a condition number.  In particular, a statement about
$\kappa_{\rm red}$ says nothing by itself about
$\kappa_2(H_{\rm NE})$, $\kappa_2(H_{\rm OSS})$, or the full saddle system.

\subsection{Quantum computation and access models}
\label{sec:prelim-quantum-access}

Algorithms are uniform quantum circuits with arbitrary input-independent
gates, intermediate measurements, and classical control.  A query and a gate
are different resources.  Calls to an oracle, its inverse, or a controlled
version count as oracle calls unless a result explicitly states otherwise.

\begin{definition}[Raw sparse coefficient (sign) access]
\label{def:raw-sparse-access}
In the \emph{raw sparse coefficient (sign) access model}, the dimensions and a
fixed list of nonzero positions of the LP or SDP data are public.  The hidden
input is a string $\varsigma\in\{-1,+1\}^{n_{\rm bit}}$ that determines designated
coefficient signs (or designated entries from a fixed finite alphabet).
The charged oracle is
\begin{equation}
 O_\varsigma\ket{i,z}=\ket{i,z\mathbin\oplus\operatorname{enc}(\varsigma_i)}.
 \label{eq:sign-xor-oracle}
\end{equation}
A sparse row, sparse column, or coefficient query is admissible only when it
can be simulated coherently with the stated number of calls to
$O_\varsigma$; a query does not reveal several hidden signs merely because a
coefficient is repeated in a derived matrix.  The equivalent phase convention
$\ket{i}\mapsto\varsigma_i\ket{i}$ and the XOR convention
\eqref{eq:sign-xor-oracle} interconvert with one charged query in either
direction, using phase kickback and the assumed controlled form of the
oracle.
\end{definition}

\begin{remark}
Definition~\ref{def:raw-sparse-access} is a deliberately restrictive,
project-specific specialization of the standard quantum query model
\cite{beals2001-quantum-lower-bounds-by-polynomials}.  Public support and
hidden signs make reductions auditable.  Supplying a prefix-parity table, an
eliminated matrix, or an arbitrary input-dependent block-encoding completion
at unit cost is a stronger model and is never implicit.
\end{remark}

\begin{definition}[Block encoding]
\label{def:block-encoding}
Let $M$ act on $\ell$ qubits.  An $(\alpha,a,\epsilon)$ \emph{block encoding} of
$M$ is an $(\ell+a)$-qubit unitary $U$ such that
\begin{equation}
 \norm{M-\alpha(\bra{0^a}\otimes I)U(\ket{0^a}\otimes I)}_2\leq\epsilon,
 \qquad \alpha>0.
 \label{eq:block-encoding}
\end{equation}
The scalar $\alpha$ is the \emph{block-encoding normalization} (also called
subnormalization), $a$ is the ancilla count, and $\epsilon$ is absolute
operator-norm error \cite{andras2019-quantum-singular-value-transformation-beyond}.
\end{definition}

Normalizations multiply under products: composing encodings with
normalizations $\alpha$ and $\beta$ yields normalization $\alpha\beta$ and
the corresponding propagated error.  Hence the cost of polynomial inversion
is governed by the normalized spectral cutoff.  For invertible $H$, the
access-normalized condition parameter is
\begin{equation}
 \kappa_{\rm BE}(H;\alpha_H)=\alpha_H\norm{H^{-1}}_2,
 \label{eq:kappa-be}
\end{equation}
whereas $\kappa_H=\norm{H}_2\norm{H^{-1}}_2$.  They coincide only when the
encoding is normalized at the operator norm.  Sparse-access constructions can
have $\alpha_H$ controlled by the row/column sparsity of $H$, not by
$\nnz(A)$, and product encodings can have a still larger normalization
\cite{andras2019-quantum-singular-value-transformation-beyond}.

\begin{definition}[State-preparation access]
\label{def:state-preparation-access}
For a nonzero vector $v\in\C^N$, \emph{state-preparation access} supplies a
unitary $U_v$ with
\begin{equation}
 U_v\ket{0}=\ket{v}:=\frac{1}{\norm{v}_2}\sum_{i=1}^N v_i\ket{i},
 \label{eq:state-prep}
\end{equation}
together with controlled access to $U_v$ and $U_v^\dagger$.  If
$\norm{v}_2$ is needed by the task, it must be supplied or estimated
separately; when signed amplitudes matter, the oracle promise also fixes a
phase convention for the controlled circuit.  For a mixed state, the
analogous oracle prepares a specified purification.  This model is stronger
than receiving independent copies, because inverse access enables
amplitude-estimation procedures
\cite{joran2023-quantum-tomography-state-preparation-unitaries}.
\end{definition}

State preparation exposes one normalized ray.  It does not, merely by being
callable, provide coherent numerical access to individual $v_i$, a diagonal
operator $\diag(v)$, or a reusable update oracle.  Such interfaces are
separately defined output contracts below.

\begin{definition}[Plain block-encoding access and factor access]
\label{def:plain-factor-access}
\emph{Plain block-encoding access} to a Hermitian positive-definite matrix
$H$ supplies a controlled $(\alpha_H,a_H,\epsilon_H)$ block encoding of $H$
and its adjoint.  \emph{Factor access} supplies instead a controlled
$(\alpha_G,a_G,\epsilon_G)$ block encoding of a possibly rectangular $G$
(padded to square dimensions in the standard way) and its adjoint, together
with the promise
\begin{equation}
 H=G^\dagger G.
 \label{eq:factor-access}
\end{equation}
These are different input models.  The standard product construction using
two factor-oracle calls gives an
$(\alpha_G^2,2a_G,2\alpha_G\epsilon_G+\epsilon_G^2)$ block encoding of $H$
(the $\epsilon_G^2$ term is needed because \cref{def:block-encoding} does not
require $\alpha_G\ge\norm G$).  Thus a QLS
solver applied to this composed encoding uses normalization
$\alpha_H=\alpha_G^2$ and access-normalized inversion parameter
$\alpha_G^2\norm{H^{-1}}_2$.  A solver that acts on $G$ directly must instead
state its dependence on $\alpha_G$ and on the relevant singular-value cutoff
of $G$; it may not silently substitute plain access to $H$.
\end{definition}

Normal equations and reduced log-barrier Hessians often have natural factors,
so later results state explicitly which side of
Definition~\ref{def:plain-factor-access} they assume.  No factorization is
free merely because one exists algebraically.

\begin{definition}[Quantum linear-system solver]
\label{def:qls}
A \emph{quantum linear-system (QLS) solver} for a nonsingular Hermitian matrix
$H$ is a procedure that, given the stated matrix access and a preparation
oracle for $\ket b$, produces a state $\rho$ satisfying
\begin{equation}
 \Dtr\!\left(\rho,
 \frac{H^{-1}\ket b\!\bra bH^{-1}}
      {\norm{H^{-1}\ket b}_2^2}\right)\leq\epsilon_{\rm ls}.
 \label{eq:qls-contract}
\end{equation}
Its resource statement must name the linear-system condition number
$\kappa_H$, the block-encoding normalization $\alpha_H$ (or sparse-oracle
normalization), the preparation cost of $\ket b$, success probability, and
precision $\epsilon_{\rm ls}$.
\end{definition}

HHL introduced the state-output linear-systems primitive
\cite{harrow2009-quantum-algorithm-linear-systems-equations}.  The
Childs--Kothari--Somma algorithm achieves polylogarithmic precision dependence
in sparse access, variable-time techniques improve amplification overhead,
and QSVT implements inversion as a singular-value polynomial
\cite{andrew2017-quantum-algorithm-systems-linear-equations,andris2012-variable-time-amplitude-amplification-quantum,andras2019-quantum-singular-value-transformation-beyond}.
Discrete-adiabatic and adiabatic solvers provide other implementations
\cite{luigi2022-optimal-scaling-quantum-linear-systems,subasi2019-adiabatic-quantum-linear-system}.
These algorithm families have different constants and parameter dependences,
but all require an access-and-output ledger; none turns
\eqref{eq:qls-contract} into a classical vector for free.

\begin{definition}[Charged QRAM access]
\label{def:charged-qram}
A \emph{QRAM access assumption} supplies a coherent lookup unitary for an
explicit classical table, for example
$\ket{i,z}\mapsto\ket{i,z\mathbin\oplus\operatorname{enc}(d_i)}$
\cite{giovannetti2008-architectures-for-a-quantum-random}.
In \emph{charged QRAM access}, raw-input oracle calls and QRAM writes used to
build or update that table are charged as separate resources; subsequent
lookups have the advertised QRAM query cost.
\end{definition}

QRAM can reduce coherent lookup depth once data have been loaded, but it does
not erase the information cost of producing an input-dependent table.  The
lower bounds in this paper therefore charge setup.  A QRAM containing all hidden parities,
Newton entries, or solution coordinates has already been given information
that raw coefficient access was designed to hide.

Tomography illustrates the output bottleneck.  For an $N$-dimensional pure
state, a classical amplitude vector with $\ell_2$ error $\epsilon$ requires
$\Thetatilde(N/\epsilon)$ controlled calls in the state-preparation
unitary model, while sampling-based approaches have quadratic dependence on
$1/\epsilon$; mixed-state costs also depend on rank and the requested matrix
norm \cite{joran2023-quantum-tomography-state-preparation-unitaries}.  Thus a
polylogarithmic-dimensional QLS state-generation cost is compatible with a
linear-dimensional classical readout cost.

\subsection{Output and accounting contracts}
\label{sec:prelim-output-contracts}

An algorithm solves only the contract it is asked to solve.  The following
names are canonical throughout the paper.

\begin{definition}[Value-estimation output]
\label{def:value-output}
For a scalar target $v$, an \emph{additive $\epsilon$ value-estimation output}
is a classical $\widetilde v$ such that
$|\widetilde v-v|\leq\epsilon$.  A \emph{target-relative $\epsilon$
value-estimation output} satisfies
$|\widetilde v-v|\leq\epsilon|v|$ and is used only under a nonzero-target
promise.  An unqualified ``relative value estimate'' always means this
target-relative contract.  The required success probability is part of
either contract.
\end{definition}

\begin{definition}[Reference-scaled value-estimation output]
\label{def:reference-scaled-value-output}
For a stated scale $s_{\rm ref}>0$, a \emph{reference-scaled $\epsilon$
value-estimation output} satisfies
$|\widetilde v-v|\leq\epsilon s_{\rm ref}$.  This contract is used when zero
may be a valid target or when the problem supplies a natural external scale;
the scale is always named rather than called a relative error.
\end{definition}

\begin{definition}[Classical vector output]
\label{def:classical-vector-output}
A \emph{classical vector output} for $x\in\R^N$ is an explicit classical list
$\widetilde x\in\R^N$ satisfying a named error guarantee, including its norm
and scale.  Unless stated otherwise that guarantee is
$\norm{\widetilde x-x}_2\leq\epsilon$.  Producing only
$\widetilde x/\norm{\widetilde x}_2$, samples from its coordinates, or an
implicit circuit is a different contract.
\end{definition}

\begin{definition}[Compact optimizer output]
\label{def:compact-optimizer-output}
A \emph{compact (sparse-support) optimizer output} is a classical description
of an identified component or subset $I$ together with the nonzero
coordinates of a solution supported on $I$, satisfying a named optimality and
feasibility tolerance.  It is an explicit sparse solution description, not a
full dense vector and not merely a quantum state or sampling oracle.
\end{definition}

\begin{definition}[Normalized solution-state output]
\label{def:solution-state-output}
For a nonzero solution vector $x$ in the original problem coordinates, an
\emph{$\epsilon$ normalized solution-state output} (also called an
\emph{amplitude-encoded solution state}) is a density operator $\rho$ such that
\begin{equation}
 \Dtr\!\left(\rho,
 \ket{x/\norm{x}_2}\!\bra{x/\norm{x}_2}\right)\leq\epsilon,
 \qquad
 \Dtr(\rho,\sigma)=\tfrac12\norm{\rho-\sigma}_1.
 \label{eq:solution-state-output}
\end{equation}
The output does not include $\norm{x}_2$, coordinate values, or a preparation
unitary unless those items are explicitly requested.  When $x$ is an LP
primal candidate, we use the synonymous stable name \emph{amplitude-encoded
primal-state output}.  A state encoding coefficients in a reduced or
null-space basis is a different contract; the later original-coordinate
state lower bounds do not automatically apply to that reduced-coordinate
output.
\end{definition}

\begin{definition}[Trace-normalized density output]
\label{def:trace-density-output}
For a nonzero PSD solution $X$, an \emph{$\epsilon$ trace-normalized density
output} is $\rho$ satisfying
\begin{equation}
 \Dtr\!\left(\rho,X/\tr X\right)\leq\epsilon.
 \label{eq:trace-density-output}
\end{equation}
For a product-PSD solution this means the direct-sum density operator
$\bigoplus_jX_j/(\sum_j\tr X_j)$, including its block-label register.  This
contract deliberately discards the total trace, which must be output
separately when needed.
\end{definition}

\begin{definition}[Coordinate-oracle output]
\label{def:coordinate-oracle-output}
An \emph{$\epsilon$ coordinate-oracle output} for $x\in\R^N$ is a reusable,
controlled circuit $O_{\widetilde x}$ and its adjoint acting as
\begin{equation}
 O_{\widetilde x}\ket{i,z}
 =\ket{i,z\mathbin\oplus\operatorname{enc}(\widetilde x_i)},
 \qquad |\widetilde x_i-x_i|\leq\epsilon
 \label{eq:coordinate-oracle-output}
\end{equation}
for every $i$, with a stated fixed-point encoding and failure convention.
The cost of one invocation and any one-time compilation cost are both part of
the contract.  A normalized solution state is not a coordinate-oracle output.
\end{definition}

\begin{definition}[Unconditional and heralded preparation]
\label{def:heralded-output}
An \emph{unconditional preparation} is a channel whose output state itself
satisfies the promised error bound after all internal randomness and flags
are discarded.  A \emph{heralded preparation} outputs a classical or
computational-basis success flag; conditioned on that flag, the state obeys
the error bound, the success probability is at least a stated $p>0$ on every
allowed input, and repetitions or amplification needed to obtain success are
charged.  An unobserved or unspecified ``good branch'' is neither contract.
\end{definition}

\begin{definition}[Global RHS-relative residual]
\label{def:global-residual}
For $b\neq0$, the \emph{global RHS-relative $\ell_2$ primal residual} of an LP
candidate $x$ is
\begin{equation}
 \operatorname{res}_p(x)=\frac{\norm{Ax-b}_2}{\norm b_2}.
 \label{eq:global-residual}
\end{equation}
For an SDP candidate $X$, replace $Ax$ by $\mathcal A(X)$.  A statement that
the residual is at most $\eta$ always refers to this ratio unless it names a
different contract.  Per-row relative residuals, coefficientwise residuals,
normal-equation residuals, local-norm residuals, and backward error under
perturbed data are different and do not imply one another without additional
scale assumptions.
\end{definition}

\begin{definition}[First-preparation and amortized total-query accounting]
\label{def:amortized-accounting}
For one fixed input oracle and $R$ requested outputs, \emph{setup} is the
reusable input-dependent work performed before the first output is produced.
Let $Q_{\rm setup}$ and $W_{\rm setup}$ be, respectively, its raw-input query
count and QRAM-write count, and let $Q_t$ and $W_t$ be the corresponding
nonreusable costs attributable to output $t$.  Raw queries and QRAM writes are
distinct columns in the resource ledger and are never added together.  The
\emph{first-preparation cost} is the pair
$(Q_{\rm setup}+Q_1,W_{\rm setup}+W_1)$.  The total ledgers are
\begin{equation}
 Q_{\rm total}(R)=Q_{\rm setup}+\sum_{t=1}^R Q_t,
 \qquad
 W_{\rm total}(R)=W_{\rm setup}+\sum_{t=1}^R W_t,
 \label{eq:total-query-ledger}
\end{equation}
with amortized costs obtained by dividing each column by $R$.  If setup caches
the complete fixed instance, later $Q_t$ may be zero in a pure query model.
Conversely, every raw query or QRAM write used to prepare consumed advice or
additional state copies is charged in setup or in the consuming call.
\end{definition}

Definition~\ref{def:amortized-accounting} prevents two opposite accounting
errors: repeatedly charging the same static information as if it were fresh,
and declaring an input-dependent reusable data structure free.  In the pure
query model with a fixed, finite, cacheable hidden input, an additive
per-round raw-query lower bound therefore needs a genuine freshness, memory,
communication, or consumable-resource assumption.

\subsection{Lower-bound toolkit}
\label{sec:prelim-lower-bound-tools}

\paragraph{Parity.}
For $n_{\rm bit}$ oracle bits, bounded-error quantum computation of their
parity needs $\Omega(n_{\rm bit})$ queries.  This follows from the polynomial
method because a
$T$-query algorithm has an acceptance probability represented by a real
multilinear polynomial of degree at most $2T$, whereas parity has linear
approximate degree \cite{beals2001-quantum-lower-bounds-by-polynomials}.  Our
raw-sign reductions therefore verify that each optimization-data query can be
simulated by $O(1)$ sign queries and then decode parity from the contracted
output.

\paragraph{Adversary bounds and composition.}
Ambainis's adversary method assigns nonnegative weights to pairs of inputs
with different answers and bounds the progress one query can make
\cite{ambainis2002-quantum-lower-bounds-by-quantum}.  Its positive-weight
spectral formulation composes, and already covers an OR of independent
parity gadgets, with the product of the outer and inner adversary scales
\cite{hyer2005-tight-adversary-bounds-for-composite}.  The more general
negative-weight bound $\operatorname{Adv}^{\pm}$ composes for Boolean outer functions and
characterizes bounded-error Boolean function evaluation up to constants
\cite{lee2011-quantum-query-complexity-of-state,reichardt2011-reflections-for-quantum-query-algorithms}.
We do not invoke either composition statement when component queries overlap
or an aggregate predicate is supplied.

\paragraph{Polynomial method and symmetrization.}
More generally, every measurement probability after $T$ standard input
queries is a degree-$2T$ polynomial in the input bits.  Averaging over inputs
of equal Hamming weight turns it into a univariate polynomial without
increasing degree.  Approximation constraints at the promised weights can
then be combined with classical polynomial inequalities to lower-bound $T$
\cite{beals2001-quantum-lower-bounds-by-polynomials}.  Later applications
state the promise points and error interval explicitly; symmetrization does
not justify inserting unpromised inputs.

\paragraph{Trace distance and Helstrom decoding.}
Trace distance is contractive under channels, and the statistical distance
between the outcome distributions of any measurement on $\rho$ and $\sigma$
is at most $\Dtr(\rho,\sigma)$.  Conversely, for equal priors the optimal
binary discrimination success probability is
$\tfrac12(1+\Dtr(\rho,\sigma))$, the Helstrom bound
\cite{nielsen2010-quantum-computation-quantum-information-10th}.  Hence a fixed
measurement that decodes a hidden bit from an ideal solution state with bias
$\beta$ still has positive bias from any output within trace distance less
than $\beta$.

\paragraph{State conversion.}
In the state-conversion problem, an input oracle indexed by $x$ must transform
a prescribed family of initial states $\{\ket{\rho_x}\}$ into target states
$\{\ket{\sigma_x}\}$ to bounded error.  The filtered $\gamma_2$ query distance
compares the two Gram matrices
$[\langle\rho_x|\rho_y\rangle]$ and
$[\langle\sigma_x|\sigma_y\rangle]$ through the matrices that record which
input symbols a query can distinguish.  It gives lower bounds for general
state conversion, the direction used in this paper, and matching upper and
lower bounds up to a gap in their error parameters.  In the special case of
function evaluation, the resulting general adversary bound characterizes
bounded-error query complexity up to constants
\cite{lee2011-quantum-query-complexity-of-state}.  Later uses specify the full
oracle completion: prescribing only $U_x\ket0$ is insufficient when an
algorithm may query $U_x$ away from $\ket0$.

\subsection{Recurring notation}
\label{sec:prelim-notation}

\begin{table}[t]
 \centering
 \small
 \caption{Notation used across the technical sections.  Several sections
 locally rebind $G$, $N$, $P$, $\alpha$, and $\eta$ and state those
 conventions locally.}
 \label{tab:recurring-notation}
 \begin{tabular}{@{}lp{0.72\linewidth}@{}}
  \toprule
  Symbol & Meaning \\
  \midrule
  $m,n$ & Number of LP equality constraints and primal variables. \\
  $q,d_j$ & Number of PSD blocks and dimension of block $j$. \\
  $n_{\rm bit},\varsigma$ & Number of hidden signs and their string in a raw coefficient-oracle reduction. \\
  $N$ & Dimension of a vector requested by an output contract. \\
  $N_H$ & Dimension of the actual Newton, normal-equation, OSS, or MNES system sent to a linear solver. \\
  $A,b,c$; $\mathcal A,b,C$ & LP data; product-form SDP data. \\
  $X,S$ & Diagonal LP matrices $\diag(x),\diag(s)$, or primal and dual PSD variables when the context is semidefinite. \\
  $F,\nu$ & Named ambient cone barrier and its declared cone parameter. \\
  $\mu$ & Central-path parameter. \\
  $\bar\mu$ & Average LP complementarity $x^\top s/n$ at an iterate. \\
  $\sigma$; $\sigma_{\min},\sigma_{\max}$ & Newton centering parameter; minimum and maximum singular values. \\
  $g$ & Primal objective gap $c^\top x_F-\mathrm{OPT}$, not automatically the primal--dual gap. \\
  $N_2(\theta)$ & Feasible Euclidean complementarity neighborhood in \eqref{eq:n2-neighborhood}. \\
  $W,V$ & Orthonormal tangent basis used in $H_{\rm red}$; a general full-rank null-space basis used in OSS. \\
  $H_{\rm red},H_{\rm NE},H_{\rm OSS},H_{\rm MNES}$ & Reduced barrier Hessian and the named Newton-system matrices. \\
  $\kappa_{\rm red}$ & Condition number of $H_{\rm red}$ in an orthonormal Euclidean tangent basis. \\
  $\kappa_{\rm amb}$ & Condition number of the named ambient cone-barrier Hessian before tangent restriction. \\
  $\kappa_H$ & Spectral condition number of the actual named linear-system matrix $H$; not $\kappa_2(A)$ unless proved equal. \\
  $s_H$ & Maximum row/column sparsity of $H$; not $\nnz(A)$. \\
  $(\alpha_H,a_H,\epsilon_H)$ & Normalization, ancilla count, and operator error of a block encoding of $H$. \\
  $(\alpha_G,a_G,\epsilon_G)$ & Corresponding factor block-encoding data for $H=G^\dagger G$. \\
  $\kappa_{\rm BE}$ & Access-normalized inversion parameter $\alpha_H\norm{H^{-1}}_2$. \\
  $\epsilon_{\rm opt},\epsilon_{\rm ls}$ & Optimization accuracy and one-linear-solve state accuracy. \\
  $\eta$ & Named feasibility or inexact-Newton residual tolerance, as specified by the applicable contract. \\
  $\Dtr$ & Trace distance $\tfrac12\norm{\rho-\sigma}_1$. \\
  $Q_{\rm setup},Q_{\rm total}$ & Charged setup queries and total raw queries including setup. \\
  $W_{\rm setup},W_{\rm total}$ & Charged setup writes and total QRAM writes including setup. \\
  $\nnz(A)$ & Number of stored nonzeros of $A$; it does not encode fill or derived-system sparsity. \\
  \bottomrule
 \end{tabular}
\end{table}

\section{The Geometry of Central-Path Conditioning}
\label{sec:conditioning}

\subsection{Question, answer, and relation to prior work}
\label{sec:conditioning-overview}

Quantum interior-point methods that invoke a quantum linear-system solver, and
classical interior-point methods that solve Newton equations inexactly, pay for
the conditioning of a Newton matrix.  A familiar estimate is
$\kappa=\Theta(\mu^{-2})$, but the classical derivations concern the
logarithmic barrier and often a primal--dual Schur complement rather than the
reduced primal Hessian.  It is therefore natural to ask whether a different
self-concordant barrier can remove the ill-conditioning.

For reduced barrier Hessians the answer is no in the relevant minimax sense.
The obstruction is the Euclidean geometry of the objective sublevel set.  We
prove a lower bound for every self-concordant barrier, pin every eigenvalue
between directional chord widths, and show that the canonical log and log-det
barriers attain the condition-number floor up to an instance-dependent
constant.  Error-bound exponents then give a dictionary of conditioning laws.

The universal quantifier over barriers is important.  Earlier conditioning
analyses concern the log barrier, including Wright's analysis of logarithmic
barrier Hessians and the survey of Forsgren, Gill, and Wright
\cite{wright1994-properties-hessian-logarithmic-barrier,forsgren2002-interior-methods-nonlinear-optimization}.
The $\mu^{-2}$-type estimates in
\citet[Section~7.4]{augustino2023-quantum-interior-point-methods-semidefinite},
for example, are for primal--dual Schur systems.  Conversely,
\citet{allamigeon2022-no-self-concordant-barrier} prove a lower bound uniform
over self-concordant barriers, but for the number and geometry of central-path
iterations, not Hessian conditioning.  The error-bound consequences below
connect to weak sharp minima \cite{burke1993-weak-sharp-minima} and to Sturm's
singularity-degree bounds for semidefinite systems
\cite{sturm2000-error-bounds-linear-matrix-inequalities}.

\begin{remark}[Novelty and textbook ingredients]
\label{rem:conditioning-novelty}
The Dikin-ellipsoid and asymmetric-containment facts used in the proofs, as
well as the centrality identity
$\Pi_{\mathcal V}\nabla F(x_\mu)=-\Pi_{\mathcal V}c/\mu$, are standard
self-concordance tools
\cite{nesterov1994-interior-point-polynomial-algorithms-convex,renegar2001-mathematical-view-ipm}.
The contribution here is their universal-over-barriers conditioning statement
and its geometric consequences, not new containment machinery.  In
particular, the evidence for novelty is a prior-art audit, not a claim that
these proof ingredients are new.
\end{remark}

\subsection{Sublevel chords and two containment facts}
\label{sec:conditioning-setup}

Let $P\subset\R^N$ be compact and convex, with affine hull
$x^0+\mathcal V$ and nonempty relative interior $P^\circ$.  Let $W$ have
orthonormal columns spanning the $d$-dimensional tangent space $\mathcal V$.
We minimize the nonconstant linear functional $c^\top x$, so
$c_{\mathcal V}:=\Pi_{\mathcal V}c\ne0$.  Put
\begin{equation}
 \mathrm{OPT}=\min_{y\in P}c^\top y,
 \qquad
 \Phi=\{y\in P:c^\top y=\mathrm{OPT}\}.
 \label{eq:conditioning-optimal-face}
\end{equation}
For $t\geq0$, define the objective sublevel set, its diameter, and the longest
sublevel chord from an interior point by
\begin{equation}
 L(t)=\{y\in P:c^\top y\leq\mathrm{OPT}+t\},\qquad
 D(t)=\operatorname{diam} L(t),\qquad
 \ell(x)=\sup_{y\in L(c^\top x-\mathrm{OPT})}\norm{y-x}_2.
 \label{eq:sublevel-chord-definitions}
\end{equation}
Thus $D$ is nondecreasing.  If $x$ has objective gap $g$, then $x\in L(g)$
and the triangle inequality gives
\begin{equation}
 \frac{D(g)}2\leq\ell(x)\leq D(g).
 \label{eq:diameter-chord-comparison}
\end{equation}
Indeed, for any $y,z\in L(g)$, at least one of
$\norm{x-y}_2,\norm{x-z}_2$ is at least $\norm{y-z}_2/2$.

Let $F$ be a $\nu_F$-self-concordant barrier for $P$ in the sense of
\cref{def:self-concordant-barrier}.  We use two standard consequences of
self-concordance.  For $x\in P^\circ$,
\begin{align}
 y\in x^0+\mathcal V,\quad
 \norm{y-x}_x<1&\quad\Longrightarrow\quad y\in P^\circ,
 \tag{Dikin}\label{eq:dikin-containment}\\
 y\in P,\quad
 \inner{\nabla F(x)}{y-x}\geq0
 &\quad\Longrightarrow\quad
 \norm{y-x}_x\leq C_F,
 \qquad C_F:=\nu_F+2\sqrt{\nu_F}.
 \tag{asymmetric}\label{eq:asymmetric-containment}
\end{align}
The first implication is understood for points in the affine hull; its
contrapositive says that a boundary or exterior point has local distance at
least one.  We also use semiboundedness,
\begin{equation}
 \inner{\nabla F(x)}{y-x}\leq\nu_F\qquad(y\in P).
 \label{eq:barrier-semiboundedness}
\end{equation}
These facts follow from the standard self-concordant-barrier containment
theorems; affine restriction preserves them and does not increase the
declared parameter
\cite{nesterov1994-interior-point-polynomial-algorithms-convex,renegar2001-mathematical-view-ipm}.

At an exact central point $x=x_F(\mu)$, first-order optimality on the affine
slice is
\begin{equation}
 W^\top(c+\mu\nabla F(x))=0.
 \label{eq:centrality-tangent}
\end{equation}
Consequently, for every $y\in P$,
\begin{equation}
 c^\top(x-y)=\mu\inner{\nabla F(x)}{y-x}\leq\nu_F\mu.
 \label{eq:central-objective-bound-pointwise}
\end{equation}
Taking $y\in\Phi$ proves the convention required in
\cref{rem:g-convention}:
\begin{equation}
 0<g_F(\mu)\leq\nu_F\mu.
 \label{eq:central-gap-bound}
\end{equation}
The strict inequality holds because a relative-interior minimizer of a
nonconstant linear functional cannot exist.  Notice that
\eqref{eq:central-gap-bound} is generally an inequality.  The equality
$\nu\mu$ in \cref{def:central-path} is the canonical \emph{primal--dual}
gap, not necessarily the primal objective gap $g$ used here.

\subsection{Rigidity of the Hessian spectrum}
\label{sec:conditioning-rigidity}

We begin with the barrier-independent lower edge.  It does not require
centrality.

\begin{theorem}[Universal minimum-eigenvalue rigidity]
\label{thm:lambda-min-rigidity}
Under the setup of \cref{sec:conditioning-setup}, for every
$\nu_F$-self-concordant barrier $F$ for $P$ and every $x\in P^\circ$,
\begin{equation}
 \lambda_{\min}\!\left(W^\top\nabla^2F(x)W\right)
 \geq \frac1{\ell(x)^2}.
 \label{eq:lambda-min-rigidity}
\end{equation}
Here $\ell(x)$ is formed from the objective level of that same point, whether
or not $x$ is central.
\end{theorem}

\begin{proof}
Fix a Euclidean unit vector $h\in\mathcal V$ and replace it by $-h$ if
necessary so that $c^\top h\leq0$.  Compactness implies that the ray
$x+th$ leaves $P^\circ$ in finite time.  More precisely, if
\[
 \tau=\sup\{t\geq0:x+sh\in P^\circ\text{ for }0\leq s<t\},
\]
then $0<\tau<\infty$ and $y=x+\tau h\in P$.  The objective does not
increase along the ray, so $y\in L(c^\top x-\mathrm{OPT})$ and hence
$\tau\leq\ell(x)$.  Since $y\notin P^\circ$, the contrapositive of
\eqref{eq:dikin-containment} gives
\[
 1\leq\norm{\tau h}_x^2
   =\tau^2\,\inner{\nabla^2F(x)h}{h}.
\]
Thus the Rayleigh quotient is at least
$\tau^{-2}\geq\ell(x)^{-2}$.  The Hessian form is even in $h$, so the
initial sign choice loses no direction.  Minimizing over the Euclidean unit
sphere in $\mathcal V$ proves \eqref{eq:lambda-min-rigidity}.
\end{proof}

The opposite spectral edge is forced by the distance in objective value to
the optimum.  Combining it with \cref{thm:lambda-min-rigidity} gives the
sublevel chord law.

\begin{theorem}[Sublevel chord law]
\label{thm:sublevel-chord-law}
Under the setup of \cref{sec:conditioning-setup}, let
$x=x_F(\mu)$ be an exact central point of any
$\nu_F$-self-concordant barrier, and set $g=g_F(\mu)$.  Then
\begin{align}
 \lambda_{\max}\!\left(W^\top\nabla^2F(x)W\right)
 &\geq \frac{\norm{c_{\mathcal V}}_2^2}{g^2},
 \label{eq:chord-lmax}\\
 \frac1{\ell(x)^2}
 \leq\lambda_{\min}\!\left(W^\top\nabla^2F(x)W\right)
 &\leq\frac{C_F^2}{\ell(x)^2},
 \label{eq:chord-lmin}\\
 \kappa_{\rm red}(\mu;F)
 &\geq
 \left(\frac{\ell(x)\norm{c_{\mathcal V}}_2}{C_Fg}\right)^2
 \geq
 \left(\frac{D(g)\norm{c_{\mathcal V}}_2}{2C_Fg}\right)^2.
 \label{eq:sublevel-chord-floor}
\end{align}
In particular, using only the $\mu$ parametrization,
\begin{equation}
 \kappa_{\rm red}(\mu;F)
 \geq
 \left(
  \frac{D(g_F(\mu))\norm{c_{\mathcal V}}_2}
       {2C_F\nu_F\mu}
 \right)^2.
 \label{eq:sublevel-chord-floor-mu}
\end{equation}
There is no small-$\mu$ threshold and no polyhedrality, nondegeneracy, or
strict-complementarity hypothesis.
\end{theorem}

\begin{proof}
The gap statement \eqref{eq:central-gap-bound} was proved above.  For the
largest eigenvalue, let
\[
 \widehat h=-\frac{c_{\mathcal V}}{\norm{c_{\mathcal V}}_2},
 \qquad
 t_* =\frac{g}{\norm{c_{\mathcal V}}_2}.
\]
The affine point $z=x+t_*\widehat h$ has $c^\top z=\mathrm{OPT}$.  If it
were in $P^\circ$, it would be a relative-interior minimizer of a linear
functional, forcing $c_{\mathcal V}=0$.  Therefore $z\notin P^\circ$, and
\eqref{eq:dikin-containment} gives
\[
 \inner{\nabla^2F(x)\widehat h}{\widehat h}
 \geq t_*^{-2}=\frac{\norm{c_{\mathcal V}}_2^2}{g^2}.
\]
This proves \eqref{eq:chord-lmax}.

For the upper bound in \eqref{eq:chord-lmin}, take any $y\in L(g)$.  By
\eqref{eq:centrality-tangent},
\[
 \inner{\nabla F(x)}{y-x}
 =-\frac{c^\top(y-x)}\mu\geq0.
\]
Asymmetric containment \eqref{eq:asymmetric-containment} therefore gives
$\norm{y-x}_x\leq C_F$.  If $y\ne x$, the Euclidean unit vector
$v=(y-x)/\norm{y-x}_2$ satisfies
\[
 \inner{\nabla^2F(x)v}{v}
 \leq\frac{C_F^2}{\norm{y-x}_2^2}.
\]
The smallest eigenvalue is no larger than this Rayleigh quotient.  Taking a
sequence of $y$ whose distance from $x$ tends to $\ell(x)$ proves the upper
bound; the lower bound is \cref{thm:lambda-min-rigidity}.  Dividing
\eqref{eq:chord-lmax} by the upper bound in \eqref{eq:chord-lmin}, then using
\eqref{eq:diameter-chord-comparison}, proves
\eqref{eq:sublevel-chord-floor}.  Finally $g\leq\nu_F\mu$ proves
\eqref{eq:sublevel-chord-floor-mu}.
\end{proof}

The maximum-eigenvalue constant in \eqref{eq:chord-lmax} can be exact.  On
the simplex problem $\min x_3$ subject to $x_1+x_2+x_3=1$, numerical central
paths for the two tested barriers give
$\lambda_{\max}g^2/\norm{c_{\mathcal V}}_2^2\to1$; for the ordinary log
barrier, $g/\mu\to1$ and
$\lambda_{\max}\mu^2\to2/3=\norm{c_{\mathcal V}}_2^2$.

The chord idea controls more than the two extreme eigenvalues.  For an
unoriented one-dimensional subspace $[h]\subset\mathcal V$, choose its unit
representative $\widehat h$ with $c^\top\widehat h<0$.  If
$c^\top h=0$, choose the sign having the longer exit time, breaking a tie
arbitrarily.  Define
\begin{equation}
 \tau_x([h])=\sup\{t\geq0:x+s\widehat h\in P^\circ
                     \text{ for }0\leq s<t\}.
 \label{eq:directional-exit-time}
\end{equation}
The selected exit point is a sublevel point, so
$0<\tau_x([h])\leq\ell(x)$.  The equatorial convention makes $\tau_x$
well-defined on unoriented lines; continuity is not needed.

For $j=1,\ldots,d$, define the two width profiles
\begin{align}
 w_j^-(x)
 &=\inf_{\substack{U\subseteq\mathcal V\\\operatorname{codim}U=j-1}}
   \ \sup_{\substack{h\in U\\\norm h_2=1}}\tau_x([h]),
 \label{eq:lower-width-profile}\\
 w_j^+(x)
 &=\sup_{\substack{S\subseteq\mathcal V\\\dim S=j}}
   \ \inf_{\substack{h\in S\\\norm h_2=1}}\tau_x([h]).
 \label{eq:upper-width-profile}
\end{align}
We use infima and suprema because the sign choice can make the exit-time
function discontinuous on $c^\perp$.

\begin{theorem}[Width-spectrum rigidity]
\label{thm:width-spectrum-rigidity}
In the setting of \cref{thm:sublevel-chord-law}, list the eigenvalues of
$W^\top\nabla^2F(x)W$ in nondecreasing order.  For every
$j=1,\ldots,d$,
\begin{equation}
 \frac1{(w_j^-(x))^2}
 \leq\lambda_j\!\left(W^\top\nabla^2F(x)W\right)
 \leq\frac{C_F^2}{(w_j^+(x))^2}.
 \label{eq:width-spectrum-sandwich}
\end{equation}
Moreover $w_j^+(x)\leq w_j^-(x)$ for every $j$, and
$w_1^+(x)=\sup_{[h]}\tau_x([h])\geq\ell(x)$.
\end{theorem}

\begin{proof}
For every unit $h$, Dikin containment at the exit point gives
\begin{equation}
 \inner{\nabla^2F(x)h}{h}\geq\tau_x([h])^{-2}.
 \label{eq:width-pointwise-lower}
\end{equation}
For the selected orientation, centrality gives
\[
 \inner{\nabla F(x)}{\tau_x([h])\widehat h}
 =-\frac{\tau_x([h])c^\top\widehat h}{\mu}\geq0.
\]
Asymmetric containment \eqref{eq:asymmetric-containment} applies directly
to the exit point and gives
\begin{equation}
 \inner{\nabla^2F(x)h}{h}\leq C_F^2\tau_x([h])^{-2}.
 \label{eq:width-pointwise-upper}
\end{equation}
The Rayleigh quotient is unchanged by orientation.

The two Courant--Fischer formulas for eigenvalues in nondecreasing order are
\begin{align*}
 \lambda_j&=\sup_{\operatorname{codim}U=j-1}
             \ \inf_{h\in U,\norm h_2=1}R(h),\\
 \lambda_j&=\inf_{\dim S=j}
             \ \sup_{h\in S,\norm h_2=1}R(h),
\end{align*}
where $R(h)=\inner{\nabla^2F(x)h}{h}$.  Applying
\eqref{eq:width-pointwise-lower} in the first formula yields
\[
 \lambda_j\geq
 \sup_U\frac1{(\sup_{h\in U}\tau_x([h]))^2}
 =\frac1{(w_j^-(x))^2}.
\]
Applying \eqref{eq:width-pointwise-upper} in the second formula gives
\[
 \lambda_j\leq C_F^2
 \inf_S\frac1{(\inf_{h\in S}\tau_x([h]))^2}
 =\frac{C_F^2}{(w_j^+(x))^2}.
\]

For any $j$-dimensional $S$ and codimension-$(j-1)$ subspace $U$, the
dimension formula gives a nonzero line $[h]\subseteq S\cap U$.  Hence
\[
 \inf_{v\in S,\norm v=1}\tau_x([v])
 \leq\tau_x([h])
 \leq\sup_{v\in U,\norm v=1}\tau_x([v]).
\]
Taking the supremum over $S$ and infimum over $U$ proves
$w_j^+\leq w_j^-$.  Finally, if $y\in L(g)$, the segment $[x,y)$ lies in
$P^\circ$.  The line through $y-x$ is oriented toward $y$ when the objective
strictly decreases, and our longer-exit convention gives the same conclusion
when it is unchanged.  Thus
$\tau_x([y-x])\geq\norm{y-x}_2$.  Taking the supremum over $y$ proves
$w_1^+\geq\ell(x)$.
\end{proof}

Thus even the interior spectrum is constrained by directional sublevel
widths, up to the explicit $C_F^2$ window and the gap between the two minimax
widths.  The empirical consequences for spectral clustering are deferred to
the next section.

\subsection{Canonical barriers attain the floor}
\label{sec:conditioning-achievability}

We next specialize to compact affine slices of the nonnegative orthant and of
the product PSD cone in \cref{sec:prelim-conic}.  To distinguish parameters
when two barriers are compared, write $\nu_{\rm can}=n$ for the LP log
barrier and $\nu_{\rm can}=\sum_jd_j$ for the product log-det barrier.  This
is the cone parameter denoted by $\nu$ in \cref{sec:prelim-conic}.

\begin{theorem}[Canonical achievability and gap-matched comparison]
\label{thm:canonical-achievability}
Consider either of the following strictly primal--dual feasible problems with
a compact primal feasible region and nonconstant objective.
\begin{enumerate}[label=\textup{(\roman*)}]
 \item An LP in \eqref{eq:lp-pair}, with canonical path
 $(x(\mu),y(\mu),s(\mu))$ and $M(\mu)=\norm{s(\mu)}_\infty$.
 \item A product-form SDP in \eqref{eq:sdp-pair}.  Its canonical path is
 $(X(\mu),y(\mu),S(\mu))$, and
 $M(\mu)=\max_j\norm{S_j(\mu)}_2$.
\end{enumerate}
Fix $\bar\mu>0$ and suppose the tail constant
\begin{equation}
 M_{\rm tail}=\sup_{0<\mu\leq\bar\mu}M(\mu)
 \label{eq:dual-tail-constant}
\end{equation}
is finite.  Use the Euclidean metric for the LP and the product Frobenius
metric for the SDP.  Write
$g_{\rm can}(\bar\mu)=g_{F_{\rm can}}(\bar\mu)$ for the canonical primal
objective gap at the endpoint of the tail.  The notation
$\kappa_{\rm red}^{\rm can}(g)$ and $\kappa_{\rm red}^{F}(g)$ means the
reduced condition number at the central point of gap $g$ for the named
barrier, using the $g$-parametrization in \cref{def:central-path} on an
interval where the named objective-gap map $g_F$ is one-to-one.  For a
canonical central point with $0<\mu\leq\bar\mu$ and primal objective gap $g$,
\begin{align}
 \lambda_{\min}(H_{\rm red})&\geq\ell(x)^{-2},
 \label{eq:canonical-lmin}\\
 \lambda_{\max}(H_{\rm red})&\leq M_{\rm tail}^2\mu^{-2}
 \leq M_{\rm tail}^2\nu_{\rm can}^2g^{-2},
 \label{eq:canonical-lmax}\\
 \kappa_{\rm red}^{\rm can}(g)&\leq
 \left(\frac{M_{\rm tail}\nu_{\rm can}D(g)}g\right)^2.
 \label{eq:canonical-kappa-upper}
\end{align}
For every gap $0<g\leq g_{\rm can}(\bar\mu)$ attained on the canonical tail,
any $\nu_F$-self-concordant barrier $F$ on the same feasible region, and any
exact $F$-central point having that gap,
\begin{equation}
 \kappa_{\rm red}^{\rm can}(g)
 \leq
 \left(
  \frac{2M_{\rm tail}\nu_{\rm can}
              (\nu_F+2\sqrt{\nu_F})}
       {\norm{c_{\mathcal V}}_2}
 \right)^2
 \kappa_{\rm red}^{F}(g).
 \label{eq:gap-matched-barrier-comparison}
\end{equation}
The comparison is between points of equal objective gap; their values of
$\mu$ need not agree.
\end{theorem}

\begin{proof}
The first bound is \cref{thm:lambda-min-rigidity}.  For the LP, canonical
complementarity gives $x_i=\mu/s_i$, and therefore, for every Euclidean unit
tangent vector $h$,
\[
 \inner{\nabla^2F_{\log}(x)h}{h}
 =\sum_i\frac{h_i^2}{x_i^2}
 \leq\frac{M_{\rm tail}^2}{\mu^2}.
\]
For the SDP, $X_j^{-1}=S_j/\mu$, and for a unit product-Frobenius tangent
direction $H$,
\begin{align*}
 D^2F_{\log\det}(X)[H,H]
 &=\sum_j\norm{X_j^{-1/2}H_jX_j^{-1/2}}_F^2\\
 &\leq \max_j\norm{X_j^{-1}}_2^2\sum_j\norm{H_j}_F^2
 \leq\frac{M_{\rm tail}^2}{\mu^2}.
\end{align*}
This proves the first inequality in \eqref{eq:canonical-lmax}.  The canonical
primal--dual gap is $\nu_{\rm can}\mu$, whereas its primal objective gap
satisfies $g\leq\nu_{\rm can}\mu$ by \eqref{eq:central-gap-bound}; hence
$\mu^{-1}\leq\nu_{\rm can}/g$.  Combining the largest-eigenvalue bound with
\eqref{eq:canonical-lmin} and $\ell(x)\leq D(g)$ proves
\eqref{eq:canonical-kappa-upper}.  Finally divide that upper bound by the
universal lower bound \eqref{eq:sublevel-chord-floor} for $F$ to obtain
\eqref{eq:gap-matched-barrier-comparison}.
\end{proof}

The tail boundedness in \eqref{eq:dual-tail-constant} is an explicit
instance hypothesis; it holds in the standard convergent canonical-path
setting.  The restriction to a tail is necessary because the corresponding
supremum can diverge as $\mu\to\infty$.  More importantly,
\eqref{eq:gap-matched-barrier-comparison} says that no alternative barrier can
improve the canonical reduced-Hessian condition number by more than a
$g$-independent instance factor.  It does not say that an arbitrary barrier
cannot be worse.

\subsection{A dictionary from error bounds to conditioning}
\label{sec:conditioning-dictionary}

The diameter $D(g)$ is the relevant sharpness modulus.  If, for some
$\theta\in[0,1]$ and $\gamma>0$,
\begin{equation}
 \operatorname{dist}(y,\Phi)\leq\gamma^{-1}
       (c^\top y-\mathrm{OPT})^\theta\qquad(y\in P),
 \label{eq:holder-error-bound}
\end{equation}
then choosing nearest optimal points for two members of $L(g)$ gives
\begin{equation}
 D(g)\leq\operatorname{diam}\Phi+2\gamma^{-1}g^\theta.
 \label{eq:diameter-from-error-bound}
\end{equation}
Conversely, $\Phi\subseteq L(g)$ and a central point $x$ of gap $g$ belongs
to $L(g)$.  For every $x^*\in\Phi$,
$g=c_{\mathcal V}^\top(x-x^*)\leq
\norm{c_{\mathcal V}}_2\norm{x-x^*}_2$, so
\begin{equation}
 D(g)\geq\operatorname{diam}\Phi,
 \qquad
 D(g)\geq\frac g{\norm{c_{\mathcal V}}_2}
 \quad\text{at every attained central gap }g.
 \label{eq:diameter-lower-basics}
\end{equation}

\begin{corollary}[Conditioning regimes]
\label{cor:conditioning-regimes}
Assume the hypotheses of \cref{thm:sublevel-chord-law}; for canonical upper
bounds also assume \cref{thm:canonical-achievability}, and restrict those
conclusions to the canonical tail.
\begin{enumerate}[label=\textup{(\roman*)}]
 \item If $\dim\Phi\geq1$, then every fixed $\nu_F$-self-concordant barrier
 satisfies $\kappa_{\rm red}^{F}(g)=\Omega(g^{-2})$.  The canonical barrier
 satisfies $\kappa_{\rm red}^{\rm can}(g)=\Theta(g^{-2})$.
 \item If $P$ is a polytope with a unique optimum, including a degenerate
 vertex, then $D(g)=\Theta(g)$ as $g\downarrow0$.  Consequently the canonical
 log barrier has $\kappa_{\rm red}^{\log}(g)=\Theta(1)$.
 \item If the optimum is unique and $D(g)=\Theta(\sqrt g)$, as under a
 two-sided quadratic-growth geometry, every fixed barrier has
 $\kappa_{\rm red}^{F}(g)=\Omega(g^{-1})$ and the canonical barrier has
 $\kappa_{\rm red}^{\rm can}(g)=\Theta(g^{-1})$.
\end{enumerate}
The constants may depend on the instance and, for a noncanonical barrier, on
$\nu_F$.  Whenever the canonical objective gap is $\Theta(\mu)$, as in the
standard strictly complementary tail, the same three canonical rates are
$\Theta(\mu^{-2})$, $\Theta(1)$, and $\Theta(\mu^{-1})$, respectively.
\end{corollary}

\begin{proof}
If $\dim\Phi\geq1$, the face contains two distinct points and hence
$\operatorname{diam}\Phi>0$; compactness makes this diameter finite and
attained.
Substitution of $D(g)\geq\operatorname{diam}\Phi$ in
\eqref{eq:sublevel-chord-floor} proves the universal lower bound.  Since
$D(g)\leq\operatorname{diam}P$, \eqref{eq:canonical-kappa-upper} gives the
matching canonical upper bound.

For a polytope, the solution set of a feasible bounded linear program is a
weak sharp minimum
\cite{mangasarian1979-nonlinear-perturbation-linear-programs,burke1993-weak-sharp-minima}.
With a unique optimum $x^*$, this gives \eqref{eq:holder-error-bound} with
$\theta=1$, hence $D(g)\leq2g/\gamma$.  For the reverse bound, choose a
feasible $z\ne x^*$ and set
$\Delta=c^\top z-\mathrm{OPT}>0$.  For every $0<g\leq\Delta$, convexity makes
\[
 y_g=\left(1-\frac g\Delta\right)x^*+\frac g\Delta z
\]
feasible, and $c^\top y_g=\mathrm{OPT}+g$.  Thus
\[
 D(g)\geq\norm{y_g-x^*}_2
 =\frac{\norm{z-x^*}_2}{\Delta}\,g.
\]
Combining these two bounds with
\cref{thm:sublevel-chord-law,thm:canonical-achievability} proves (ii).
Part (iii) follows by substituting
$D(g)=\Theta(\sqrt g)$ in the same two theorems.  The final conversion is
immediate under the separately stated relation $g=\Theta(\mu)$.
\end{proof}

\begin{remark}
The statement uniform over barriers in \cref{cor:conditioning-regimes} is the
lower rate.  The matching upper rate is established for the canonical
barrier.  Thus the barrier-optimal rate in the positive-dimensional-face
case is $\Theta(g^{-2})$; the results do not supply an upper bound of that
order for every conceivable self-concordant barrier.
\end{remark}

Part (ii) includes degenerate vertices: nondegeneracy is not needed for weak
sharpness.  Part (iii) explains why the reduced primal log-det Hessian at a
generic curved, unique SDP optimum can have condition number
$\Theta(\mu^{-1})$, even though its largest eigenvalue and a primal--dual
Schur complement can live at the $\mu^{-2}$ scale.  This is the conic analogue
of the $\mu^{-1}$ structured logarithmic-barrier behavior identified for
nonlinear programs by
\citet{wright1994-properties-hessian-logarithmic-barrier}.

The same substitution produces fractional exponents.  Sturm proves an
$O(\epsilon^{2^{-d}})$ distance bound when the relevant semidefinite
feasibility system has singularity degree $d$
\cite{sturm2000-error-bounds-linear-matrix-inequalities}.  An upper error
bound alone does not force a conditioning lower bound: its exponent must also
be attained by sublevel chords.  When it is sharp, the result is immediate.

\begin{corollary}[Attained H\"older and singularity-degree laws]
\label{cor:fractional-conditioning-laws}
Suppose that, along central gaps $g\downarrow0$,
$D(g)\geq a g^\theta$ for constants $a>0$ and $\theta\in[0,1]$.  Then every
$\nu_F$-self-concordant barrier satisfies
\begin{equation}
 \kappa_{\rm red}^{F}(g)
 \geq
 \frac{a^2\norm{c_{\mathcal V}}_2^2}{4C_F^2}
 g^{2\theta-2}.
 \label{eq:holder-conditioning-law}
\end{equation}
If, in addition, the problem is a canonical LP or product-form SDP satisfying
all hypotheses of \cref{thm:canonical-achievability} and
$D(g)=O(g^\theta)$, then the canonical barrier attains
$\Theta(g^{2\theta-2})$ on its tail.  In particular, when the Sturm exponent
$\theta=2^{-d}$ is attained, the exponents are
\begin{equation}
 d=0:\ 0,\qquad d=1:\ -1,\qquad d=2:\ -\tfrac32,
 \qquad
 2\theta-2=2^{1-d}-2\longrightarrow-2.
 \label{eq:singularity-degree-exponents}
\end{equation}
\end{corollary}

\begin{proof}
Insert $D(g)\geq ag^\theta$ into
\eqref{eq:sublevel-chord-floor}.  If the matching upper bound on $D$ holds,
combine the resulting lower bound with
\eqref{eq:canonical-kappa-upper}.  Substituting $\theta=2^{-d}$ gives
\eqref{eq:singularity-degree-exponents}.
\end{proof}

The following $3\times3$ problem shows that the $d=2$ fractional law occurs
in a concrete SDP rather than only in the error-bound formalism.

\begin{proposition}[An explicit $g^{-3/2}$ witness]
\label{prop:degree-two-witness}
Consider
\begin{equation}
 \begin{aligned}
  \min_{X\in\Sym^3}\quad &X_{22}\\
  \text{subject to}\quad &\tr X=1,\qquad X_{33}=X_{12},
  \qquad X\succeq0.
 \end{aligned}
 \label{eq:degree-two-sdp}
\end{equation}
Its unique optimum is $X^*=e_1e_1^\top$, its optimality system reaches the
minimal face in two facial-reduction steps, and
\begin{equation}
 D(g)=\Theta(g^{1/4}).
 \label{eq:degree-two-diameter}
\end{equation}
Consequently every self-concordant barrier obeys
$\kappa_{\rm red}(g)=\Omega(g^{-3/2})$, while the log-det barrier, under the
hypotheses of \cref{thm:canonical-achievability}, has
$\kappa_{\rm red}(g)=\Theta(g^{-3/2})$ on its canonical tail.
\end{proposition}

\begin{proof}
Write a feasible matrix as
\[
 X=\begin{bmatrix}
      a&b&t\\ b&d&e\\ t&e&b
    \end{bmatrix},
 \qquad a+d+b=1,\qquad X\succeq0.
\]
Here $d=X_{22}$ is the objective and $b=X_{12}=X_{33}\geq0$.  If $d=0$,
positive semidefiniteness makes the second row and column zero, hence $b=e=0$.
Then the third diagonal entry is zero, so the third row and column vanish and
$t=0$; the trace constraint gives $a=1$.  The optimum is therefore unique.

For $X\in L(g)$, the $2\times2$ principal minors give
\[
 b^2\leq ad\leq d\leq g,
 \qquad
 t^2\leq ab\leq b\leq\sqrt g,
 \qquad
 e^2\leq db\leq g^{3/2}.
\]
Also $|a-1|=d+b\leq g+\sqrt g$.  Thus
$\norm{X-X^*}_F=O(g^{1/4})$, and the triangle inequality gives
$D(g)=O(g^{1/4})$.

For the reverse bound, let $0<g\leq1/16$, set
\[
 d=g,\qquad b=\frac{\sqrt g}{2},\qquad
 a=1-g-\frac{\sqrt g}{2},\qquad e=0,\qquad
 t=\pm\frac{g^{1/4}}{2\sqrt2}.
\]
Then $a\geq13/16$.  The one- and two-dimensional principal minors are
positive, and
\[
 \det X
 =ag b-b^3-gt^2
 =g^{3/2}\left(\frac a2-\frac14\right)>0.
\]
Both sign choices therefore give positive-definite feasible points in
$L(g)$.  Their Frobenius distance is
$2\sqrt2|t|=g^{1/4}$, proving the lower bound in
\eqref{eq:degree-two-diameter}.

For completeness, the two facial reductions can be read directly from the
optimality equations.  The first exposing equation $X_{22}=0$ removes the
$e_2$ row and column; then $X_{12}=0$ turns the affine equation
$X_{33}=X_{12}$ into $X_{33}=0$, which removes the $e_3$ row and column.
No one-step positive-semidefinite exposing combination can do both: a PSD
linear combination of the trace, $X_{33}-X_{12}$, and $X_{22}$ equations
whose $(1,1)$ entry is zero must have zero first row, forcing the coefficient
of $X_{33}-X_{12}$ to vanish.  Thus two steps are necessary and sufficient.
Finally apply \cref{cor:fractional-conditioning-laws} with $\theta=1/4$.
\end{proof}

For the log-det barrier, the leading constant can also be computed
exactly. For sufficiently small positive gap $g$, the center and central parameter are
\[
 X_{\log}(g)=\begin{pmatrix}a&b&0\\b&g&0\\0&0&b\end{pmatrix},\qquad
 b=\frac{-g+\sqrt{3g-2g^2}}3,\quad a=1-g-b,\quad
 \mu=\frac{ag-b^2}{1-b-2g}.
\]
Indeed, sign symmetry forces the two displayed zero entries, and
stationarity of $g/\mu-\log b-\log(g(1-g-b)-b^2)$ gives
$3b^2+2gb+g^2-g=0$ and the stated $\mu$. Strict convexity makes
this interior stationary point the unique center. The spectral calculation
in \cref{sec:spectra} gives
\[
 \kappa_{\rm red}\mu^{3/2}\longrightarrow\frac{2\sqrt2}{7},
 \qquad
 \kappa_{\rm red}g^{3/2}\longrightarrow\frac{3\sqrt3}{7}.
\]
The first constant is approximately $0.40406$, refining the earlier
numerical value $0.40$. The exact center, Frobenius-metric spectrum, and
these limits are verified in Lean; the accompanying report
\texttt{formal/FRACTIONAL\_SDP.md} states the verification scope.
The geometric and singularity-degree proof above is separate from that
formalization.

Near-optimal faces yield a related finite plateau rather than a divergent
tail.

\begin{corollary}[Near-degeneracy plateau]
\label{cor:near-degeneracy-plateau}
Suppose a family of compact instances has a scale $\vartheta>0$ and constants
independent of $g,\vartheta$ such that
\begin{equation}
 D(g)=
 \begin{cases}
  \Theta(g/\vartheta),&0<g\lesssim\vartheta,\\
  \Theta(1),&\vartheta\lesssim g\leq g_0.
 \end{cases}
 \label{eq:near-face-diameter-law}
\end{equation}
Assume also that $\norm{c_{\mathcal V}}_2$ is bounded below and that the
barrier parameters under comparison are bounded above uniformly in
$\vartheta$; for the canonical upper bound, assume the same of
$M_{\rm tail}\nu_{\rm can}$ and assume
$g_0\leq g_{\rm can}(\bar\mu)$ uniformly in $\vartheta$.
Then every self-concordant barrier has the corresponding floor
\begin{equation}
 \kappa_{\rm red}^{F}(g)=
 \begin{cases}
  \Omega(\vartheta^{-2}),&0<g\lesssim\vartheta,\\
  \Omega(g^{-2}),&\vartheta\lesssim g\leq g_0,
 \end{cases}
 \label{eq:near-face-floor}
\end{equation}
up to its explicit $C_F$ and instance constants.  The canonical barrier
matches both orders on the tail under
\cref{thm:canonical-achievability}.  Thus the
$g^{-2}$ growth stops near $g\asymp\vartheta$ at height
$\Theta(\vartheta^{-2})$.
\end{corollary}

\begin{proof}
In the first range, $D(g)/g=\Theta(\vartheta^{-1})$; in the second it is
$\Theta(g^{-1})$.  Apply the lower bound
\eqref{eq:sublevel-chord-floor} and the canonical upper bound
\eqref{eq:canonical-kappa-upper}.
\end{proof}

For the simplex family
$\min\{x_3+\vartheta x_2:x_1+x_2+x_3=1,\ x\geq0\}$, the log-barrier
plateau has the following precise two-scale limit.  For each fixed
$\vartheta>0$,
\begin{equation}
 \lim_{\mu\downarrow0}\kappa_{\rm red}(\mu)
 =
 \frac{1+\vartheta^2+\sqrt{1-\vartheta^2+\vartheta^4}}
      {1+\vartheta^2-\sqrt{1-\vartheta^2+\vartheta^4}},
 \qquad
 \lim_{\vartheta\downarrow0}\vartheta^2
       \left(\lim_{\mu\downarrow0}\kappa_{\rm red}(\mu)\right)=\frac43.
 \label{eq:simplex-plateau-double-limit}
\end{equation}
To verify the inner limit, write the equality multiplier as
$\alpha_\mu=\mu/x_1$.  Centrality gives
$x_2=\mu/(\vartheta+\alpha_\mu)$,
$x_3=\mu/(1+\alpha_\mu)$, and $\alpha_\mu\to0$.  On the tangent plane, put
$h=(-u-v,u,v)$.  After multiplication by $\mu^2$, the limiting Rayleigh
quotient is
\[
 \frac{\vartheta^2u^2+v^2}{2(u^2+uv+v^2)}.
\]
Its generalized eigenvalues are
$\bigl(1+\vartheta^2\pm\sqrt{1-\vartheta^2+\vartheta^4}\bigr)/3$, whose
ratio is the first expression in
\eqref{eq:simplex-plateau-double-limit}; expansion at
$\vartheta=0$ gives the second limit.  The computations at
$\vartheta=10^{-2}$ and $10^{-3}$ round
$\vartheta^2\lim_{\mu\downarrow0}\kappa_{\rm red}(\mu)$ to $4/3$.

\subsection{Transfer to approximate centrality}
\label{sec:conditioning-approximate}

The floor is stable in the short-step neighborhood.  All norms in the next
statement are restricted to the tangent space.

\begin{corollary}[Approximate-centrality transfer]
\label{cor:approximate-centrality-transfer}
Fix $\mu>0$, let $x_*=x_F(\mu)$, and suppose $x\in P^\circ$ satisfies
\begin{equation}
 \lambda_\mu(x):=\norm{c/\mu+\nabla F(x)}_x^*\leq\rho\leq\frac16.
 \label{eq:approximate-centrality-decrement}
\end{equation}
Then, with $H_*=W^\top\nabla^2F(x_*)W$ and
$H=W^\top\nabla^2F(x)W$,
\begin{equation}
 \norm{x-x_*}_{x_*}\leq\frac{\rho}{(1-\rho)^2}\leq\frac14,
 \qquad
 \kappa(H)\geq\left(\frac34\right)^4\kappa(H_*).
 \label{eq:approximate-centrality-transfer}
\end{equation}
Consequently $x$ inherits the right-hand sides of
\eqref{eq:sublevel-chord-floor} and
\eqref{eq:sublevel-chord-floor-mu}, evaluated at the exact center $x_*$,
with the additional factor $(3/4)^4$.
\end{corollary}

\begin{proof}
Restrict $f=F+c^\top(\cdot)/\mu$ to the affine slice.  It is standard
self-concordant, has Hessian $\nabla^2F$, and is minimized at $x_*$.  With
\[
 \omega(t)=t-\log(1+t),\qquad
 \omega_*(t)=-t-\log(1-t),
\]
the self-concordant Newton estimates give
\[
 \omega(\norm{x-x_*}_{x_*})
 \leq f(x)-f(x_*)
 \leq\omega_*(\lambda_\mu(x))
 \leq\omega_*(\rho).
\]
The elementary scalar inequality
$\omega^{-1}(\omega_*(\rho))\leq\rho/(1-\rho)^2$ follows by substituting
$t=\rho/(1-\rho)^2$ and differentiating the difference
$\omega(t)-\omega_*(\rho)$ from its value zero at $\rho=0$.  This proves the
distance bound, and $\rho\leq1/6$ gives
$\rho/(1-\rho)^2\leq6/25<1/4$.

If $r=\norm{x-x_*}_{x_*}<1$, self-concordant Hessian comparison gives
\[
 (1-r)^2H_*\preceq H\preceq(1-r)^{-2}H_*.
\]
It follows that
$\lambda_{\max}(H)\geq(1-r)^2\lambda_{\max}(H_*)$ and
$\lambda_{\min}(H)\leq(1-r)^{-2}\lambda_{\min}(H_*)$; hence
$\kappa(H)\geq(1-r)^4\kappa(H_*)\geq(3/4)^4\kappa(H_*)$.
\end{proof}

\subsection{Implications and open geometric questions}
\label{sec:conditioning-implications}

For QIPMs that solve barrier-Newton systems, changing the self-concordant
barrier cannot remove a sublevel-chord conditioning obstruction.  At target
primal objective gap $\epsilon$, every barrier pays at least
\[
 \left(
  \frac{D(\epsilon)\norm{c_{\mathcal V}}_2}
       {2(\nu_F+2\sqrt{\nu_F})\epsilon}
 \right)^2
\]
in the reduced Hessian, and, for canonical LPs or product-form SDPs with
finite $M_{\rm tail}$ and $0<\epsilon\leq g_{\rm can}(\bar\mu)$ attained on
the canonical tail, the canonical barrier is within the explicit instance
factor in \eqref{eq:gap-matched-barrier-comparison}.  In particular,
positive-dimensional optimal faces impose a quadratic inverse-gap floor,
curved unique optima can impose a linear inverse-gap floor, and sharp
singularity-degree exponents can impose fractional laws.

This conclusion is deliberately scoped to $\kappa_{\rm red}$ as defined in
\cref{def:reduced-kappa}.  It makes no claim about condition-number-free QIPM
approaches such as
\citet{apers2026-quantum-speedups-linear-programming-interior}, about
preconditioners outside barrier-Hessian form, or about scalar congruences.
Those access-model and system-representation questions are treated in later
sections.

\begin{openproblem}[Degenerate-vertex constants]
\label{open:degenerate-vertex-constants}
For a polytope with a unique degenerate optimal vertex, characterize the
finite limiting constant (or limiting spectrum) of the canonical reduced log
Hessian directly from the active-set geometry.
\end{openproblem}

\begin{openproblem}[Facial characterization of chord growth]
\label{open:facial-chord-growth}
Characterize the exact growth of $D(g)$ for semidefinite programs from facial
data beyond the generic strict-complementarity and attained
singularity-degree cases.  In particular, determine when Sturm's
$2^{-d}$ exponent is sharp for objective sublevel diameters rather than only
an upper error-bound exponent.
\end{openproblem}

\section{Clustered Spectra and Solver Cost}
\label{sec:spectra}

The chord law in \cref{thm:sublevel-chord-law} says when the condition
number must diverge, but not how the eigenvalues between its two edges are
distributed.  This section resolves that question for the logarithmic
barrier of a degenerate LP and then asks what the resulting spectrum costs.
The answer separates solver families.  The spectrum eventually consists of
two intervals with parameter-independent internal condition numbers, so a
classical Krylov method needs only polylogarithmically many matrix--vector
products.  Quantum polynomial solvers cannot use the empty spectral gap in
the same way: in sufficiently large dimension, plain block-encoding access
costs $\Thetatilde(\kappa)$ in the worst case, whereas the factor access already
defined in \cref{def:plain-factor-access} has a
$\Thetatilde(\sqrt\kappa)$ polynomial route under norm-tight encodings.
Clustering improves neither quantum model beyond the dependence its access
model already provides, while it exposes a classical--quantum
separation on the Newton systems themselves.  There is an important
qualification.  At exact central points the path's own Newton right-hand
sides are asymptotically orthogonal to the weak modes, so path-following
steps can use a condition-number-free window solve after the gap opens.

\subsection{The two-cluster tail}
\label{sec:spectra-two-cluster}

Consider the standard-form LP \eqref{eq:lp-pair} and suppose that its primal
feasible polytope $P$ is compact with $P^\circ\ne\varnothing$.  Let
$Z\in\R^{n\times p}$ have orthonormal columns spanning $\ker A$, assume as
in \cref{sec:conditioning} that the objective is nonconstant on the affine
hull, and follow
the exact log-barrier central path $(x(\mu),y(\mu),s(\mu))$.  Let $B$ be the
maximal primal optimal support and $N=[n]\setminus B$; equivalently, this is
the Goldman--Tucker optimal partition.  Write $P_B$ and $P_N$ for coordinate
projections, and put $D_P=\operatorname{diam}P$.

Choose a fixed central-path tail $0<\mu\leq\bar\mu$ on which
\begin{equation}
 S_*=\sup_{0<\mu\leq\bar\mu}\max_j s_j(\mu)<\infty,
 \qquad
 \beta=\inf_{0<\mu\leq\bar\mu}\min_{j\in B}x_j(\mu)>0.
 \label{eq:spectra-tail-constants}
\end{equation}
Strict complementarity of the central-path limit also gives
$\underline s_N:=\inf_{0<\mu\leq\bar\mu}\min_{i\in N}s_i(\mu)>0$.
These are instance constants, not asymptotic variables.

\begin{proposition}[Two-cluster log-barrier Hessian]
\label{prop:two-cluster-hessian}
Let
\begin{equation}
 H_\mu=Z^\top X(\mu)^{-2}Z=H_{\rm red}(\mu),
 \qquad r=\rank(P_NZ),
 \label{eq:two-cluster-hessian}
\end{equation}
where the equality uses \cref{sec:newton-formulations}'s reduced-Hessian notation
with $W=Z$.  If $r>0$, let $\sigma>0$ be the smallest positive singular
value of $P_NZ$.  The largest $r$ eigenvalues of $H_\mu$ then lie in
\begin{equation}
 \left[
  \frac{\sigma^2\min_{i\in N}s_i(\mu)^2}{\mu^2},
  \frac{\max_{i\in N}s_i(\mu)^2}{\mu^2}+\beta^{-2}
 \right].
 \label{eq:two-cluster-top}
\end{equation}
If $p-r>0$, the remaining $p-r$ eigenvalues lie in
\begin{equation}
 [D_P^{-2},\,\beta^{-2}].
 \label{eq:two-cluster-bottom}
\end{equation}
Each nonempty interval has a $\mu$-independent internal ratio.  When both
are nonempty, their separation grows as $\Theta(\mu^{-2})$; when $r=0$ the
top-cluster assertion is empty.
\end{proposition}

\begin{proof}
Complementarity gives $x_i(\mu)^{-2}=s_i(\mu)^2/\mu^2$.  Split
\begin{equation}
 H_\mu=H_{\mu,N}+H_{\mu,B},
 \quad
 H_{\mu,N}=Z^\top P_N^\top X_N^{-2}P_NZ,
 \quad
 H_{\mu,B}=Z^\top P_B^\top X_B^{-2}P_BZ.
 \label{eq:two-cluster-split}
\end{equation}
The first summand is positive semidefinite of rank $r$.  When $r>0$, its nonzero
eigenvalues are the squared nonzero singular values of $X_N^{-1}P_NZ$ and
therefore lie in
\[
 \left[
  \frac{\sigma^2\min_{i\in N}s_i(\mu)^2}{\mu^2},
  \frac{\max_{i\in N}s_i(\mu)^2}{\mu^2}
 \right].
\]
The tail bound on $x_B$ gives
$0\preceq H_{\mu,B}\preceq\beta^{-2}I$.
Weyl monotonicity now gives the lower edge in
\eqref{eq:two-cluster-top}, and Weyl's upper inequality gives its upper
edge.  Since $\ker H_{\mu,N}=\ker(P_NZ)$ has dimension $p-r$, the min--max
principle applied to this subspace gives
$\lambda_{p-r}^{\uparrow}(H_\mu)\leq\beta^{-2}$ when $p-r>0$, where
eigenvalues are ordered increasingly.  Hence every one of the lowest $p-r$
eigenvalues has the claimed upper bound.

It remains to bound their lower edge.  For every Euclidean unit tangent
direction, one of its two signs exits the compact polytope within distance
$D_P$.  The Dikin containment argument in
\cref{thm:lambda-min-rigidity} therefore gives
$H_\mu\succeq D_P^{-2}I$.  Finally,
\eqref{eq:spectra-tail-constants} and $\underline s_N>0$ bound each
nonempty interval's endpoint ratio independently of $\mu$.  If both clusters
are nonempty, the lower endpoint of \eqref{eq:two-cluster-top} is
$\Theta(\mu^{-2})$ while \eqref{eq:two-cluster-bottom} is fixed.
\end{proof}

The extremes agree with the geometry in \cref{sec:conditioning}.  At a
unique nondegenerate optimum---equivalently here, a singleton optimal face
with a nondegenerate optimal basis---$r=p$: if $h\in\ker A$ and $h_N=0$,
then the optimal basis gives $A_Bh_B=0$ and hence $h=0$.  The bottom cluster
is empty and the condition number remains bounded.  If the optimal face has positive
dimension, $\ker(P_NZ)$ contains its tangent space and the bottom cluster is
nonempty; its fixed scale and the $\mu^{-2}$ top scale recover the chord-law
divergence.

The two-window statement is specific to polyhedral log-barrier tails.
The log-det central path of \eqref{eq:degree-two-sdp}, the degree-two SDP
witness in \cref{prop:degree-two-witness}, has four increasingly ordered
reduced-Hessian eigenvalues in the inherited Frobenius metric:
\[
 \lambda_1\sim\sqrt2\,\mu^{-1/2},\qquad
 \lambda_2\sim\mu^{-1},\qquad
 \lambda_3\sim\sqrt2\,\mu^{-3/2},\qquad
 \lambda_4\sim\frac47\mu^{-2}.
\]
To see the metric normalization, use tangent coordinates $(b,g,t,e)$
for the symmetric matrix with diagonal $(1-g-b,g,b)$ and upper
entries $(b,t,e)$. At the center $t=e=0$, write $q=ag-b^2$ and
$A_2=\bigl(\begin{smallmatrix}a&b\\b&g\end{smallmatrix}\bigr)$.
The off-block operator is $A_2^{-1}/b$, giving the first and third
scales because $b\sim\sqrt{\mu/2}$ and $q\sim\mu$.
The $(b,g)$ coordinate Hessian $K$ has Gram matrix
$G=\bigl(\begin{smallmatrix}4&1\\1&2\end{smallmatrix}\bigr)$;
its operator is $G^{-1}K$. Direct differentiation gives
$\mu K_{bb}\to6$, $\mu^{3/2}K_{bg}\to-\sqrt2$, and
$\mu^2K_{gg}\to1$. Thus the larger eigenvalue has coefficient
$(G^{-1})_{22}=4/7$, while
$\mu^3\det K\to4$ and $\det G=7$ give the smaller coefficient one.
These four distinct powers fix the eventual order. In particular, there
is no $O(1)$ window. The complete calculation, including the actual
second derivative, the characteristic polynomial with multiplicities,
and the full central-path tail, is verified in Lean as documented in
\texttt{formal/FRACTIONAL\_SDP.md}. It replaces the earlier numerical
four-window observation with exact asymptotics.

\subsection{Why conjugate gradients see the clusters}
\label{sec:spectra-cg}

Cluster-sensitive CG estimates have a long classical lineage
\cite{axelsson1994-iterative-solution-methods}.  The following elementary
construction is enough here and makes the logarithmic dependence on the
inter-cluster gap explicit.

\begin{proposition}[Two intervals are Krylov-benign]
\label{prop:two-cluster-cg}
Let $H\succ0$ have
\[
 \operatorname{spec}(H)\subseteq[a_1,b_1]\cup[a_2,b_2],
 \qquad 0<a_1\leq b_1<a_2\leq b_2,
\]
and set $\rho_i=b_i/a_i$ and $\Lambda=b_2/a_1$.  Thus
$\kappa(H)\leq\Lambda$, with equality only when both enclosure endpoints
are attained.  For every $0<\epsilon\leq1/2$ there is a residual
polynomial $q$, with $q(0)=1$ and
$\max_{\lambda\in\operatorname{spec}(H)}|q(\lambda)|\leq\epsilon$, of
degree
\begin{equation}
 O\!\left(
  \sqrt{\rho_1\rho_2}\log\Lambda\log(1/\epsilon)
  +\sqrt{\rho_2}\log(1/\epsilon)
 \right).
 \label{eq:two-cluster-cg-degree}
\end{equation}
Consequently CG reaches relative $H$-norm error $\epsilon$ in
$O_\rho(\log(\Lambda/\epsilon)\log(1/\epsilon))$ iterations.
\end{proposition}

\begin{proof}
If $a_2<b_2$, let $p_2$ be the shifted-Chebyshev residual for $[a_2,b_2]$,
normalized by $p_2(0)=1$, of degree
$m_2=\lceil\sqrt{\rho_2}\rceil$.  The standard Chebyshev estimate gives
\[
 \max_{[a_2,b_2]}|p_2|
 \leq2e^{-2m_2/\sqrt{\rho_2}}\leq\frac12.
\]
On $[0,a_2]$ its magnitude is at most one: after the affine map from
$[a_2,b_2]$ to $[-1,1]$, the normalized Chebyshev quotient is monotone in
magnitude between the image of zero and the left endpoint.  At the edge
$a_2=b_2$, instead take $p_2(x)=1-x/a_2$ and $m_2=1$; the same two bounds
hold.

If $a_1=b_1$, take the exact factor $r_1(x)=1-x/a_1$ and $m_1=1$.
Otherwise let $r_1$ be the shifted-Chebyshev residual for $[a_1,b_1]$ of
degree $m_1=O(\sqrt{\rho_1}\log(1/\epsilon))$, with $r_1(0)=1$ and
$\max_{[a_1,b_1]}|r_1|\leq\epsilon$.  In its explicit quotient, the
denominator is $T_{m_1}(\eta_1)$, where
$\eta_1=(a_1+b_1)/(b_1-a_1)>1$, and
$T_{m_1}(\eta_1)\geq\eta_1^{m_1}$.  Combining this denominator bound with
$|T_m(t)|\leq(2|t|)^m$ for the numerator outside the approximation interval
gives, also for the exact-factor edge case,
\begin{equation}
 \max_{[a_2,b_2]}|r_1|\leq(4\Lambda)^{m_1}.
 \label{eq:lower-chebyshev-growth}
\end{equation}
Choose
\[
 j=\left\lceil m_1\log_2(4\Lambda)+\log_2(1/\epsilon)\right\rceil,
 \qquad q=r_1p_2^j.
\]
Then $q(0)=1$.  On the first interval, $|q|\leq\epsilon$ because
$|p_2|\leq1$; on the second, \eqref{eq:lower-chebyshev-growth} and the
choice of $j$ again give $|q|\leq\epsilon$.  Its degree is
$m_1+jm_2$, which is \eqref{eq:two-cluster-cg-degree}.  CG minimizes the
$H$-norm error over residual polynomials of its degree, so its error is no
larger.
\end{proof}

The numerical check used a 40-variable simplex-slice LP with a
two-coordinate active set and a 37-dimensional optimal face.  From
$\mu=10^{-2}$ to $10^{-8}$, $\kappa$ grew from about $10$ to
$7\times10^{12}$, the bottom-cluster ratio stayed below $1.006$, and the top
ratio approached $4.004$ (from $2.44$ at the shallowest point).  CG solved
random right-hand sides to a recursively updated residual tolerance of
$10^{-8}$ in 3--8 iterations.  At $\mu=10^{-8}$, the reported eight-step
recursive residual masks a postcomputed true relative residual
$\norm{b-Hx}_2/\norm b_2\approx10^{-4}$ and relative $H$-norm error about
$6\times10^{-6}$, a float64 loss-of-conjugacy effect.  A separate 80-bit
check, outside the accompanying scripts, used six iterations;
the benign iteration count therefore survives the precision check.  The
unstructured $\sqrt\kappa$ estimate is about $2.6\times10^6$ iterations.
The accompanying analysis scripts produce the float64 spectrum and
recursive-residual counts.

\subsection{Quantum polynomial boundedness and the access-model gap}
\label{sec:spectra-quantum}

CG evaluates its residual polynomial only on eigenvalues.  QSVT has a
different contract: after normalization, its polynomial must be bounded on
all of $[-1,1]$, including the empty spectral gap and the negative half of
the interval.  This difference rules out transplanting the composite
polynomial in \cref{prop:two-cluster-cg}.  An order-of-magnitude evaluation
of this construction on representative intervals reaches about $10^{85}$
inside the gap; this illustrative estimate is not used below.

Normalize $\norm{H}_2=1$ and write
\begin{equation}
 \operatorname{spec}(H)\subseteq
 [1/\Lambda,\rho_1/\Lambda]\cup[1/\rho_2,1],
 \qquad \Lambda=\kappa(H),
 \label{eq:quantum-two-cluster-normalization}
\end{equation}
where $\rho_1,\rho_2\geq1$ are fixed.  Cost statements below
assume norm-tight encodings,
$\alpha_H=\Theta(\norm H)$ and $\alpha_G=\Theta(\norm G)$.  Without this
hypothesis the plain-access parameter is
$\kappa_{\rm BE}=\alpha_H\norm{H^{-1}}$ from \cref{eq:kappa-be}, and the
factor parameter is $\kappa_G=\alpha_G\norm{H^{-1}}^{1/2}$.

\begin{theorem}[Theorem Q: clustered solves in the two access models]
\label{thm:quantum-cluster-obstruction}
For the two-cluster spectra \eqref{eq:quantum-two-cluster-normalization}:
\begin{enumerate}[label=\textup{(\roman*)}]
 \item In the one-sided polynomial model, a residual polynomial bounded by
 $M$ on $[0,1]$, with $q(0)=1$ and
 $|q(1/\Lambda)|\leq\epsilon\leq1/2$, has degree
 \begin{equation}
  d\geq\sqrt{\frac{(1-\epsilon)\Lambda}{2M}}.
  \label{eq:markov-one-sided-no-go}
 \end{equation}
 Thus its degree--amplitude-normalization metric satisfies
 $dM=\Omega(\sqrt\Lambda)$.
 \item In the fixed-polynomial plain-access model, any single polynomial
 that, for $\rho_1\geq2$, relatively inverts a bottom interval
 containing the points $y=1,2$ to accuracy at most $1/8$, and is bounded
 by $M$ on the QSVT-required symmetric interval
 $[-\Lambda,\Lambda]$, obeys
 \begin{equation}
  dM\geq\frac5{16}\sqrt{\Lambda^2-4}=\Omega(\Lambda).
  \label{eq:symmetric-bernstein-no-go}
 \end{equation}
 This is a lower bound on the polynomial metric for a right-hand side
 supported on the bottom cluster, not a general degree-to-QLS-cost theorem.
 \item With factor access $H=G^\dagger G$, the one-sided construction is an
 even singular-value polynomial with
 $dM=\Otilde(\sqrt\Lambda)$.  Item~\textup{(i)}, applied to the residual of
 this construction, gives the matching
 $\Omega(\sqrt\Lambda)$ degree lower bound.
 \item Let $N_{\rm QLS}$ denote the Hilbert-space dimension, distinct from
 the optimal-partition index set $N$ above.  In the full query model,
 plain access has $\Thetatilde(\kappa)$ worst-case QLS cost for
 $N_{\rm QLS}\gtrsim\kappa$, even for two-point clusters.  More precisely,
 \citet[Proposition~6]{orsucci2021-on-solving-classes-of-positive} prove an
 $\Omega(\min\{\kappa,N_{\rm QLS}\})$ normalized-block-encoding lower bound
 for all algorithms.  Their Proposition~7 supplies the sparse/Grover lower
 bound $\Omega(\min\{\kappa,\sqrt{N_{\rm QLS}}\})$, whose
 linear-in-$\kappa$ regime requires $N_{\rm QLS}\geq\kappa^2$.
\end{enumerate}
\end{theorem}

\begin{proof}
For (i), the mean-value theorem gives a point
$\xi\in(0,1/\Lambda)$ with
$|q'(\xi)|\geq(1-\epsilon)\Lambda$.  Markov's inequality, rescaled from
$[-1,1]$ to $[0,1]$, gives
$\max_{[0,1]}|q'|\leq2d^2M$.  Rearranging proves
\eqref{eq:markov-one-sided-no-go}.  Since $M\geq|q(0)|=1$, multiplying the
bound by $M$ gives $dM=\Omega(\sqrt\Lambda)$.

For (ii), use the rescaled eigenvalue $y=\Lambda\lambda$.  Relative error
$\delta'\leq1/8$ for the target $1/y$ implies
\[
 p(1)\geq1-\delta',
 \qquad
 p(2)\leq\frac{1+\delta'}2,
 \qquad
 p(1)-p(2)\geq\frac5{16}.
\]
The mean-value theorem therefore gives a $\xi\in(1,2)$ with
$|p'(\xi)|\geq5/16$.  Bernstein's inequality on the symmetric interval
states
\[
 |p'(y)|\leq\frac{dM}{\sqrt{\Lambda^2-y^2}}.
\]
Since $|\xi|<2$, this yields
\eqref{eq:symmetric-bernstein-no-go}.

The metric $dM$ has the stated operational meaning only after fixing the
input support.  For a right-hand side in the bottom cluster, normalizing a
polynomial bounded by $M$ reduces the useful transformed amplitude by
$\Theta(M)$, so amplitude amplification gives the $dM$ proxy.  For a general
right-hand side the success amplitude also depends on its spectral overlap.
Indeed, on a two-point spectrum the degree-one separator
$p(y)=1-(y-1)/\Lambda$ has $dM=O(1)$ even though the all-algorithms result in
item~\textup{(iv)} still gives an $\Omega(\kappa)$ worst case.  Thus
item~\textup{(ii)} is not the basis for the plain-access row of the cost
table; item~\textup{(iv)} is.

For (iii), the matching lower bound is item~\textup{(i)} applied to the
residual of the construction below.  Writing $q(y)=1-y\,p_y(y)$ in the
rescaled variable $y=\Lambda\lambda$, one has $q(0)=1$,
$|q(1)|\le e^{-T}+O(\epsilon)=O(\epsilon)$, and
$|q(y)|\le e^{-Ty}+y\,|p_y(y)-f_T(y)|\le1+O(\epsilon)$ on
$[0,\max\{L_G,\Lambda\}]\supseteq[0,\Lambda]$, the approximation interval of
the construction below.  In $\lambda$ units this is a residual
polynomial of degree $d+1$ bounded by $1+O(\epsilon)$ on $[0,1]$, so
item~\textup{(i)} forces $d=\Omega(\sqrt\Lambda)$; since
$M\ge|p_y(1)|\ge1-O(\epsilon)$, also $dM=\Omega(\sqrt\Lambda)$.  Here $M$ is
the amplitude bound of the inverse polynomial in $y$ units, whose target is
$1/y$; an inverse approximant in $\lambda$ units necessarily has bound at
least $(1-\epsilon)\Lambda$, so the two normalizations are not
interchangeable.

Now include the encoding normalization explicitly.  A block encoding of
$G$ presents singular values $s=\sigma/\alpha_G$ in $[-1,1]$.  Put
$L_G=\Lambda\alpha_G^2$.  If $p_y$ is bounded on $[0,L_G]$, then
\begin{equation}
 q_G(s)=p_y(\Lambda\alpha_G^2s^2)
 \label{eq:even-factor-polynomial}
\end{equation}
is even and bounded on $[-1,1]$.  On a singular value of $G$, its argument
is $\Lambda\sigma^2$, exactly the one-sided eigenvalue variable for
$G^\dagger G$.  Norm-tightness gives $\alpha_G=\Theta(\norm G)=\Theta(1)$
under \eqref{eq:quantum-two-cluster-normalization}, so
$L_G=\Theta(\Lambda)$; it does not set $\alpha_G$ equal to one.

For completeness, a $\Otilde(\sqrt\Lambda)$ one-sided polynomial follows from
the endpoint geometry.  In $y=\Lambda\lambda$ units define
\[
 f_T(y)=\frac{1-e^{-Ty}}y=\int_0^T e^{-ty}\,dt,
 \qquad T=\Theta(\log(1/\epsilon)).
\]
For $y\geq1$, its relative error as an approximation to $1/y$ is at most
$e^{-T}$, and $|f_T(y)|\leq T$ on $[0,L_G']$, where
$L_G'=\max\{L_G,\Lambda\}=\Theta(\Lambda)$.  The shifted-Chebyshev expansion
of $e^{-ty}$ on this interval, truncated uniformly for
$0\leq t\leq T$, has degree
\[
 O\!\left(\sqrt{L_G'T\log(L_G'T/\epsilon)}
       +\log(L_G'T/\epsilon)\right).
\]
Integrating its coefficients over $t$ produces a polynomial $p_y(y)$ approximating
$f_T$ to absolute error $O(\epsilon/\Lambda)$, hence to relative error
$O(\epsilon)$ for every $y\in[1,\Lambda]$.  Its bound is
$M=O(T)$, so degree times subnormalization is
$\Otilde(\sqrt\Lambda)$ under norm-tight access.  The even polynomial
$q_G$ in \eqref{eq:even-factor-polynomial} has only twice the degree.  This
makes the construction admissible in the factor-access
polynomial model and proves item~\textup{(iii)}.  Algebraically, the
corresponding factor route to the linear-system solution is also
transparent.  For a full-column-rank factor, first form
$w=(G^\dagger)^+b$ and then $x=G^+w$; an SVD verifies
$x=(G^\dagger G)^{-1}b$.  Define the two relative output amplitudes
\[
 \zeta_1=
 \frac{\norm{(G^\dagger)^+b}}{\norm{(G^\dagger)^+}\norm b},
 \qquad
 \zeta_2=
 \frac{\norm{G^+w}}{\norm{G^+}\norm w}.
\]
Each pseudoinverse application costs $\Otilde(\kappa_G)$ to constant
relative amplitude.  The end-to-end
$\Otilde(\kappa_G)=\Otilde(\sqrt\kappa)$ solve follows only under the
explicit overlap hypothesis $\zeta_1,\zeta_2=\Omega(1)$; otherwise amplitude
amplification pays the corresponding inverse overlaps.

There is no generic factor-access QLS guarantee at
$\Otilde(\sqrt\kappa)$.  On the Orsucci--Dunjko two-point instance, the
square-root factor is diagonal and its oracle is implementable with
$O(1)$ hidden-input queries.  A generic square-root-cost factor solver would
therefore contradict their $\Omega(\kappa)$ query lower bound.  The factor
statement here is a polynomial-metric result and an overlap-hypothesis solve,
not an unrestricted QLS upper bound.
Item~\textup{(iv)} is exactly the prior-art all-algorithms grounding for the
plain-access worst case.  Together with the standard
$\Otilde(\kappa_{\rm BE})$ block-encoding QLS upper bound, it establishes
the norm-tight $\Thetatilde(\kappa)$ worst-case row below.
\end{proof}

When $\rho_1=1$, item~\textup{(ii)} is vacuous because its two interpolation
points are absent; this does not make the case algorithmically easy.
Orsucci--Dunjko's all-algorithms lower bound in item~\textup{(iv)} already
covers exact two-point clusters.

The symmetric constraint is decisive.  The polynomial in the last proof is
well behaved on $[0,\Lambda]$ but grows on the negative half; it is not a
plain-access QSVT polynomial.  The value $f_T(0)=T\ne0$ also prevents the odd
extension used for the usual $1/x$ target.  A discretized Remez linear
program makes the distinction visible.  Requiring a polynomial to separate
$p(1/\Lambda)\leq0.1$ from $p(2/\Lambda)\geq0.9$ while staying bounded by one
gave the following minimum degrees; the computation is included in the
accompanying analysis scripts.
\begin{center}
\begin{tabular}{@{}rrrrr@{}}
\toprule
$\Lambda$ & $d_{[-1,1]}$ & $d_{[-1,1]}/\Lambda$
 & $d_{[0,1]}$ & $d_{[0,1]}/\sqrt\Lambda$ \\
\midrule
50  & 52  & 1.04 & 12 & 1.70 \\
100 & 102 & 1.02 & 17 & 1.70 \\
200 & 205 & 1.02 & 24 & 1.70 \\
400 & 376 & 0.94 & 33 & 1.65 \\
\bottomrule
\end{tabular}
\end{center}
Thus the symmetric calculation is $0.94$--$1.04$ times $\Lambda$, whereas
the one-sided relaxation is about $1.7\sqrt\Lambda$.

Suppressing logarithmic precision factors and keeping the access model
fixed, the result is the following solve-cost comparison.
\begin{center}
\begin{tabular}{@{}lll@{}}
\toprule
solver & matrix access & two-cluster cost \\
\midrule
classical CG & matrix--vector products & $\polylog(\kappa)$ \\
factor polynomial route\footnotemark & norm-tight $H=G^\dagger G$
 & $\Thetatilde(\sqrt\kappa)$ \\
QLS worst case, $N_{\rm QLS}\gtrsim\kappa$
 & norm-tight plain encoding of $H$ & $\Thetatilde(\kappa)$ \\
\bottomrule
\end{tabular}
\end{center}
\footnotetext{The factor row is a polynomial-metric statement and an
end-to-end solve cost only under the overlap hypothesis
$\zeta_1,\zeta_2=\Omega(1)$; it is not a generic factor-access QLS
guarantee.}
With \cref{sec:newton-formulations}'s convention
$D=X^{1/2}S^{-1/2}$, the normal-equations matrix is
$H=AD^2A^\top$, and the factor $G=DA^\top$ satisfies
$H=G^\dagger G$ and is naturally available from
the IPM data when the corresponding factor oracle can actually be
implemented.  One should act on that factor rather than
compose it into an encoding of $H$ and inadvertently square the singular
condition number.

The all-algorithms plain-access construction is prior art due to Orsucci and
Dunjko.  Somma and de Wolf's sum-of-squares separation is also prior art,
but it is a bound for guided ground-state energy estimation, under their
dimension condition
$\dim\geq\log(1/\epsilon)/\gamma_{\rm guide}^2$, where
$\gamma_{\rm guide}$ is their guide overlap, and their stated error assumptions
\cite{somma2026-optimal-lower-bound-ground-state-energy}.  It distinguishes
$\Omega(1/\delta)$ plain resolution from
$\Omega(1/\sqrt\delta)$ sum-of-squares resolution; we do not claim a transfer
of that result to QLS.  Within the polynomial model, the matching lower
bound for the even/one-sided factor class is instead item~\textup{(i)}.
Cluster-sensitive Krylov convergence is also classical.  What appears
to be new here is the observation that the hard two-cluster quantum instances
are solved by CG in $O(1)$ iterations when the clusters are points, together
with the resulting separation on central-path Newton spectra.  The novelty
claim is this classical-side observation and assembly, not the quantum lower
bounds themselves.

\begin{openproblem}[Exact-subnormalization shift]
\label{open:exact-subnormalization-shift}
The sharpest access-model question left by Theorem~Q is the following.
The shift $H\mapsto I-H$ moves the bottom cluster to an endpoint, where
$1/\Lambda$-scale discrimination would have square-root polynomial degree.
Can a subnormalization-one block encoding of $I-H$ be constructed from plain
access to $H$ without paying the apparent $\Theta(\Lambda)$ uniform
amplification cost near the new spectral edge?
\end{openproblem}

Orsucci and Dunjko raise the closely related normalization problem for
$B=I-\eta A$ (in their notation) in their Section~4.3
\cite{orsucci2021-on-solving-classes-of-positive}.  They observe that a
linear combination of $I$ and a normalized block encoding of $A$ yields only
the subnormalized $pB$, that any black-box amplification of $B/\beta$ to $B$
is in general inefficient by the lower bound of
\cite[Theorem~73]{andras2019-quantum-singular-value-transformation-beyond}
applied to $f(x)=\beta x$, and they leave open whether a normalized block
encoding of a matrix with operator norm at most one can be built from a
sparse-matrix oracle for it, answering it positively for diagonally dominant
$A$ and for sums of positive semidefinite local Hamiltonians.
\Cref{open:exact-subnormalization-shift} is the promise-sensitive form of
their black-box obstruction: the subnormalization to remove is the fixed
factor two, the spectrum of $H$ is promised to satisfy
\eqref{eq:quantum-two-cluster-normalization}, and the question is whether that
promise lowers the $\Theta(\Lambda)$ uniform-amplification degree at the
shifted edge, not whether a general normalization procedure exists.

\subsection{The central path does not use a worst-case right-hand side}
\label{sec:spectra-rhs}

The preceding obstruction concerns arbitrary right-hand sides.  The exact
log-barrier path generates a much more special vector.  Let
$P_{\rm bot}(\mu)$ denote the spectral projector of $H_\mu$ onto the bottom
cluster in \cref{prop:two-cluster-hessian}.

\begin{proposition}[Universal Newton right-hand side and weak-mode alignment]
\label{prop:newton-rhs-alignment}
In the setting of \cref{prop:two-cluster-hessian}, at an exact central point
define $v_\mu=Z^\top X^{-1}e$.  For a centering target
$\mu_2=\gamma\mu$ with $0<\gamma<1$, the exact reduced Newton equations are
\begin{equation}
 \mu_2H_\mu d_c=r_c,
 \qquad r_c=-(1-\gamma)\mu v_\mu,
 \qquad
 \mu H_\mu d_a=r_a,
 \qquad r_a=-\mu v_\mu.
 \label{eq:universal-newton-rhs}
\end{equation}
Here $d_a$ is the affine-scaling direction formed with the current local
barrier metric; $\gamma$ avoids collision with the singular value $\sigma$
in \cref{prop:two-cluster-hessian}.  If the right-hand side is nonzero,
then
\begin{equation}
 \frac{\norm{P_{\rm bot}(\mu)r}_2}{\norm r_2}=O(\mu^2)
 \qquad(r=r_c\text{ or }r_a).
 \label{eq:rhs-bottom-fraction}
\end{equation}
For either equation in \eqref{eq:universal-newton-rhs}, if the bottom
component of its exact solution is dropped, the resulting relative
Euclidean residual is exactly the fraction in \eqref{eq:rhs-bottom-fraction}.
\end{proposition}

\begin{proof}
At exact centrality,
$Z^\top(c-\mu X^{-1}e)=0$.  Negating the reduced gradient for the target
$\mu_2$ gives
\[
 -Z^\top(c-\mu_2X^{-1}e)
 =-(1-\gamma)\mu Z^\top X^{-1}e,
\]
and setting the target barrier term to zero gives
$-Z^\top c=-\mu Z^\top X^{-1}e$.  This proves
the right-hand sides in \eqref{eq:universal-newton-rhs}.  The reduced Hessian
of the target barrier objective is $\mu_2H_\mu$; the affine direction uses
the current local metric $\mu H_\mu$.  Thus the displayed system scalars are
also exact.

For the alignment claim, recall the splitting
$H_\mu=H_{\mu,N}+H_{\mu,B}$ in \eqref{eq:two-cluster-split}.  The bottom
eigenspace of $H_{\mu,N}$ is exactly
\begin{equation}
 \mathcal T_\Phi=\ker(P_NZ),
 \label{eq:face-tangent-space}
\end{equation}
the reduced tangent space of the primal optimal face.  The positive spectrum
of $H_{\mu,N}$ begins at $\Theta(\mu^{-2})$, whereas
$\norm{H_{\mu,B}}_2=O(1)$.  The Davis--Kahan theorem therefore gives
\begin{equation}
 \norm{P_{\rm bot}(\mu)-P_{\mathcal T_\Phi}}_2=O(\mu^2)
 \label{eq:bottom-projector-rotation}
\end{equation}
\cite{davis1970-rotation-eigenvectors-perturbation}.

Let $x^a$ be the analytic center of the primal optimal face.  The log-barrier
path converges to $x^a$, and strict complementarity gives the tail expansion
$x_B(\mu)=x_B^a+O(\mu)$
\cite{mclinden1980-analogue-moreau-proximation,wright1994-properties-hessian-logarithmic-barrier}.
For every reduced unit vector $u\in\mathcal T_\Phi$, the ambient direction
$Zu$ has $(Zu)_N=0$, and first-order optimality of $x^a$ on the face says
\[
 \inner{(Zu)_B}{(X_B^a)^{-1}e_B}=0.
\]
Since the coordinates of $x_B^a$ are positive, the tail expansion then gives
$\inner{u}{Z^\top X^{-1}e}=O(\mu)$.  On the other hand
$\norm{X^{-1}e}_2=O(1/\mu)$ because $x_N=\mu S_N^{-1}e_N$.
Combining these facts with \eqref{eq:bottom-projector-rotation} yields
\begin{equation}
 \norm{P_{\rm bot}(\mu)v_\mu}_2=O(\mu).
 \label{eq:bottom-rhs-amplitude}
\end{equation}
Centrality also gives $v_\mu=Z^\top c/\mu$.  The nonconstant-objective
hypothesis makes $Z^\top c\ne0$, so $\norm{v_\mu}_2=\Theta(1/\mu)$; dividing
\eqref{eq:bottom-rhs-amplitude} proves \eqref{eq:rhs-bottom-fraction}.

Finally, write $d=d_{\rm top}+d_{\rm bot}$ for either system
$\tau H_\mu d=r$, where $\tau=\mu_2$ for centering and $\tau=\mu$ for the
affine direction, and use $\widetilde d=d_{\rm top}$.  Then
\[
 \tau H_\mu\widetilde d-r
 =-\tau H_\mu d_{\rm bot}=-P_{\rm bot}(\mu)r.
\]
Taking Euclidean norms proves the exact residual identity.
\end{proof}

\begin{corollary}[Condition-number-free top-window solve]
\label{cor:top-window-solve}
Suppose an inexact Newton step only requires a relative linear-system
residual at most $\eta$, with $\eta$ independent of $\mu$.  On the exact
log-barrier path, once $\mu=O(\sqrt\eta)$, it suffices to invert the top
cluster.  Plain-access QSVT can do so with degree
\begin{equation}
 O\!\left(\rho_2\log(\rho_2/\eta)\right),
 \label{eq:top-window-qsvt-degree}
\end{equation}
independent of $\kappa(H_\mu)$.
\end{corollary}

\begin{proof}
Dropping the bottom solution component contributes $O(\mu^2)$ relative
residual by \cref{prop:newton-rhs-alignment}; choose the tail constant so
that this is at most $\eta/2$.  After scaling the top interval to
$[1/\rho_2,1]$, the standard bounded QSVT approximation to $1/x$ has degree
$O(\rho_2\log(\rho_2/\eta))$ and can make the remaining residual at most
$\eta/2$.  The interval stays a fixed distance from zero, so this polynomial
is admissible with plain access and has no dependence on the remote bottom
cluster.
\end{proof}

The contract in this corollary is the reduced-primal relative residual
$\norm{\tau H_\mu\widetilde d-r}/\norm r$, with $\tau=\mu_2$ for centering
and $\tau=\mu$ for the affine direction.  Here $r_c$ denotes the reduced
primal right-hand side in \eqref{eq:universal-newton-rhs}, not the KKT
complementarity residual denoted by $r_c$ in \cref{sec:newton-formulations}.
This contract is distinct from the absolute OSS and MNES conventions in
\cref{def:inexact-newton-residual}: OSS requires
an $O(\bar\mu)$ residual, while MNES requires an
$O(\sqrt{\bar\mu/n})$ system residual chosen to produce an
$O(\bar\mu)$ complementarity residual.  After identifying residual vectors,
a relative tolerance $\eta_{\rm rel}$ implies an absolute requirement
$\norm{\mathrm{res}}\leq\tau_{\rm abs}$ only when
$\eta_{\rm rel}\norm r\leq\tau_{\rm abs}$.  Here $r_a=-Z^\top c$ and
$r_c=-(1-\gamma)Z^\top c$ are generally constant in norm, so a constant
relative tolerance does not by itself imply either shrinking absolute
contract.

The measured law is substantially smaller than its asymptotic bound.  On
Netlib \texttt{afiro}, the dropped-bottom residual was
$6.7\times10^{-6}$, $6.4\times10^{-10}$, and $6.7\times10^{-14}$ at
$\mu=10^{-1},10^{-3},10^{-5}$, respectively: an instance constant of about
$7\times10^{-4}$ multiplying $\mu^2$.  Both centering and affine columns
agree, as \eqref{eq:universal-newton-rhs} predicts; the top-cluster RHS mass
rounded to $1.0000$, and the solution's top mass remained at least $0.9929$.
The symmetric 40-variable family had machine-zero weak-mode mass.  The
accompanying analysis scripts generate these measurements.

An end-to-end damped-Newton run tested more than exact central points.  It
used shrink factor $0.85$, up to four corrections per parameter value, and
fraction-to-boundary steps.  A window solver dropped the bottom modes only
after detecting a gap larger than two decades and otherwise used the full
solve.  The window and exact runs agreed in final gap and correction count:
the toy family finished at gap $1.758\times10^{-8}$ with 205 corrections in
both modes, and \texttt{afiro} finished at relative gap
$7.4\times10^{-10}$ with 393 corrections in both.  This does not establish
iterate-by-iterate identity; direct instrumentation of the two runs shows
that intermediate iterates can differ.  In a separate, independently reproduced experiment, a
naive geometric-mean threshold without gap detection stalled on
\texttt{afiro} at relative gap $0.98$ by splitting the early spectral
continuum arbitrarily.  That negative case is separate from the reported
adaptive experiment.  The operational rule is therefore: use full solves
until a genuine cluster gap opens, then drop the weak modes.  The accompanying
analysis scripts implement the adaptive experiment.

\subsection{A Netlib conditioning survey}
\label{sec:spectra-netlib}

To test how early the asymptotic picture becomes visible, we surveyed 14
small presolved standard-form Netlib LPs.  For each instance, a damped Newton
method followed the log-barrier path in $\ker A$.  We fitted
$\kappa_{\rm red}(g)\asymp g^{-\alpha}$ over the last four grid points,
spanning about three gap decades, counted clusters using gaps larger than
one decade, and ran CG to $10^{-8}$ tolerance with a 4000-iteration cap at
the deepest point.  A second computation solved directional LPs in 30 fixed
random directions over the objective sublevel polytope.  The resulting
$\widehat D(g)$ is a lower estimate of the diameter, not the diameter
itself; fitting $\widehat D(g)\asymp g^\theta$ and comparing with
$\alpha=2-2\theta$ is therefore a consistency check.  It is computationally
independent of the path Hessian eigensolves, but is not an exact second
measurement of the chord-law diameter.

\begin{table}[ht]
\centering
\small
\begin{tabular}{@{}lrrrrrr@{}}
\toprule
instance & $m$ & $n$ & $\alpha$ & $\kappa_{\rm end}$ & clusters & CG its \\
\midrule
adlittle  & 55  & 136 & 2.07 & $1.7\!\times\!10^{13}$ & 2 & 4000 cap \\
afiro     & 9   & 18  & 2.00 & $3.8\!\times\!10^{12}$ & 2 & 20 \\
beaconfd  & 11  & 30  & 0.02 & $1.3\!\times\!10^6$  & 3 & 28 \\
bore3d    & 36  & 65  & 0.03 & $3.4\!\times\!10^{11}$ & 3 & 521 \\
kb2       & 46  & 71  & 0.36 & $9.6\!\times\!10^9$  & 4 & 198 \\
sc105     & 34  & 65  & 0.03 & $9.9\!\times\!10^3$  & 1 & 55 \\
sc205     & 67  & 125 & 0.03 & $7.2\!\times\!10^5$  & 1 & 286 \\
sc50a     & 17  & 33  & 0.01 & $1.5\!\times\!10^3$  & 1 & 20 \\
sc50b     & 15  & 29  & 0.00 & $5.8\!\times\!10^2$  & 1 & 16 \\
scagr7    & 93  & 147 & 1.30 & $8.1\!\times\!10^{12}$ & 2 & 4000 cap \\
scorpion  & 77  & 136 & 0.12 & $1.1\!\times\!10^5$  & 1 & 160 \\
seba      & 10  & 17  & 0.00 & $1.3\!\times\!10^3$  & 2 & 7 \\
brandy    & 107 & 211 & \multicolumn{4}{l}{numerical failure} \\
recipelp  & 63  & 107 & \multicolumn{4}{l}{numerical failure} \\
\bottomrule
\end{tabular}
\caption{Tail conditioning fits and deepest-point solver diagnostics.  The
cluster count uses a one-decade heuristic and should not be read as an exact
spectral partition.}
\label{tab:netlib-conditioning}
\end{table}

Two endpoint classes dominate \cref{tab:netlib-conditioning}.  The
positive-dimensional-face instances \texttt{afiro} and \texttt{adlittle}
have $\alpha\approx2$, while the vertex-like \texttt{sc} instances,
\texttt{seba}, \texttt{scorpion}, \texttt{beaconfd}, and \texttt{bore3d}
have frozen exponents $\alpha\approx0$; their condition-number levels can
still be large.  For
\texttt{afiro}, $\widehat D(g)$ freezes at $126.84$, so $\theta\to0$ predicts
$\alpha=2.000$, matching the path fit $2.00$.  For \texttt{sc50b},
$\widehat D/g$ freezes at $126$, so $\theta\to1$ predicts $\alpha=0$, matching its
condition-number plateau near $582$.

The intermediate \texttt{kb2} slope is a crossover, not evidence for a
third asymptotic class.  Its random-direction estimate of $\theta$ climbs from
$0.08$ to $0.99$ across the available gap decades, the in-the-wild version
of the near-degeneracy plateau in
\cref{cor:near-degeneracy-plateau}; the measured $\alpha=0.36$ lies in the
knee and will tend to zero.  The \texttt{scagr7} local fitted $\theta$ slopes
range from $0.1$ to $0.6$ and may reflect a slower multiscale crossover.  We
claim no fractional asymptotic exponent for a Netlib instance; the proved
fractional example remains the constructed SDP in
\cref{prop:degree-two-witness}.

The instructive counterpoint is \texttt{adlittle}.  At the first float64-safe
probe, its 81 reduced eigenvalues form one cluster with spread ratio
$\kappa=4.3\times10^4$ (about $4.6$ decades).  The middle probe at
$\kappa=1.85\times10^7$ is also one segment, with spread ratio
$1.9\times10^7$ (about $7.3$ decades), and at the deepest safe probe the
spectrum spans nine decades at $\kappa=1.31\times10^9$.  No decade gap
appears at any of these three depths.  CG correspondingly takes 250, 985,
and 2935 iterations;
counts above 81 reflect finite-precision loss of conjugacy.  The deeper
survey row heuristically reports two clusters at
$\kappa=1.7\times10^{13}$, where float64 spectral diagnostics are already
fragile.  This sharpens rather than contradicts
\cref{prop:two-cluster-hessian}: the asymptotic split is real, but the safe
probe establishes only that it has not opened through
$\kappa\approx1.3\times10^9$.  The follow-up calculation is included in
the accompanying analysis scripts.

That wall is itself visible in \texttt{brandy} and \texttt{recipelp}: at
computed condition numbers above roughly $10^{15}$, roundoff produces
negative eigenvalues for a mathematically positive-definite reduced Hessian.
This is consistent with the classical finite-precision analysis of interior
methods \cite{wright1998-ill-conditioning-computational-error-interior}.
Presolve, rank reduction, and Phase-I choices affect the face approached by
the path, so every exponent in \cref{tab:netlib-conditioning} belongs to the
particular presolved representation.  The accompanying analysis scripts
implement the path survey and sublevel fits; the table reports their
deterministic cached-instance run.

\subsection{Scope of the escape}
\label{sec:spectra-scope}

The two-cluster theorem and the $O(\mu^2)$ right-hand-side fraction use the
LP logarithmic barrier, strict complementarity on a sufficiently small tail,
and exact centrality.  For a general barrier, it is plausible that the
barrier's face-center first-order condition produces an analogous alignment,
perhaps only an $O(\mu)$ fraction, but this is a conjecture: neither the
cluster identification nor the required path-limit expansion has been proved
in that generality.  Off-path iterates can acquire weak-mode components.
Nothing here establishes the same rate for predictor--corrector methods,
infeasible-start variants, or unreduced primal--dual systems.  Accordingly,
\cref{cor:top-window-solve} is an exact-path structural statement and the
trajectory experiment is evidence for the stated adaptive rule, not a
general convergence theorem for those variants.

\section{Beyond the Condition Number}
\label{sec:beyond-kappa}

The condition number is a worst-case spectral parameter: it charges for a
small eigenvalue even when the right-hand side barely sees its eigenspace.
Dalzell, Li, and Su (DLS) recently gave two quantum linear-system solvers that
replace this worst-case dependence by right-hand-side-dependent truncation and
filtering parameters \cite{dalzell2026-faster-quantum-linear-system-solver};
their filtering solver improves an earlier construction of
\citet{li2025-new-condition-number-application-multivariate-polynomial}.
The classical idea of effective conditioning is older
\cite{chan1988-effectively-well-conditioned-linear-systems}, but the DLS
parameters had not been evaluated on IPM Newton systems.  We do so here.

Sections~\ref{sec:conditioning}--\ref{sec:spectra} showed that a degenerate LP
produces two increasingly separated spectral clusters.  The question in this
section is finer: when does the Newton right-hand side allow a solver to exploit
that separation?  Exact block identities give a complete coupling criterion.
They also expose a truncation transition, while a basis-invariant spectral
dimension gives more general solver-selection laws.  Finally, when the active
range is supplied, a directly assembled equilibration removes the central-path
divergence while preserving the Newton state exactly.

Throughout the first part, $H\succ0$ is presented through a block encoding of
$H/\alpha_H$.  For a nonzero right-hand side $f$, set $x=H^{-1}f$ and
define
\begin{equation}
 \xi(H,f)=\alpha_H\frac{\norm{H^{-1}f}}{\norm f},
 \qquad
 \rho(H,f)=\alpha_H
 \frac{\norm{H^{-2}f}}{\norm{H^{-1}f}}.
 \label{eq:dls-two-parameters}
\end{equation}
The second quantity is the ideal DLS filtering parameter: it is
$\norm{(H/\alpha_H)^{-1}x}/\norm x$, not the first-inverse
quantity $\xi$.  We avoid $\nu$, which \cref{sec:prelim-conic} reserves for
the barrier parameter.  More generally, for a nonsingular, possibly
non-Hermitian matrix $A$ encoded as $A/\alpha_A$, filtering uses
\begin{equation}
 \rho(A,f)=\alpha_A\frac{\norm{A^{-\dagger}x}}{\norm x},
 \qquad x=A^{-1}f.
 \label{eq:dls-general-filtering-parameter}
\end{equation}
The DLS truncation parameter instead asks how far toward zero
one must invert in order to retain all but the prescribed normalized solution
mass.

\subsection{Exact-center coupling law}
\label{sec:beyond-kappa-coupling}

Consider a standard-form LP with full-row-rank $A$, and suppose its exact
central path converges to a strictly complementary solution.  Let
\[
 B=\{i:x_i^*>0\},\qquad N=\{i:s_i^*>0\},\qquad
 U=\operatorname{range}(A_B),\qquad K=U^\perp,
\]
and assume $0<\dim U<m$.  This is the mixed, primal-degenerate case; if
$U=\R^m$, the normalized normal matrix has no asymptotically slow cluster.
At an exact center, complementarity gives
\begin{equation}
 H_\mu=H_{\rm NE}(\mu):=AXS^{-1}A^\top
 =\mu^{-1}C_\mu+\mu D_\mu,
 \quad
 C_\mu=A_BX_B^2A_B^\top,
 \quad
 D_\mu=A_NS_N^{-2}A_N^\top.
 \label{eq:coupling-normal-split}
\end{equation}
Here $C_\mu\to C_0$, $D_\mu\to D_0$, and, in $U\oplus K$ blocks,
\[
 D_\mu=
 \begin{pmatrix}D_\mu^U&F_\mu^\top\\F_\mu&E_\mu\end{pmatrix}.
\]
The restrictions $C=C_0|_U$ and $E=E_0|_K$ are positive definite.  Indeed,
$C$ is a positive-weight Gram matrix with range $U$, and if $z\in K$ obeys
$A_N^\top z=0$, then also $A_B^\top z=0$; full row rank of $A$ forces
$z=0$, proving positivity of $E$ on $K$.  At the optimal limit,
$b=Ax^*=A_Bx_B^*$, so the exact-centering right-hand side lies in $U$.  A
short step to
$\sigma_\mu\mu$, $\sigma_\mu\ne1$, has
$H_\mu\Delta y=(1-\sigma_\mu)b$.  The scalar is immaterial to normalized
QLS parameters, so the next theorem uses $b\ne0$. This nonzero condition
also follows in the stated LP setting. Indeed $x^*\ne0$; if $b=0$,
optimality gives $c^Tx^*=0$, while a strictly feasible dual slack $s$
would give $s^Tx^*=c^Tx^*>0$, a contradiction.

\begin{theorem}[Exact-center coupling criterion]
\label{thm:dls-coupling}
Define
\begin{align}
 R_\mu&=D_\mu^U-F_\mu^\top E_\mu^{-1}F_\mu,\nonumber\\
 p_\mu&=(C_\mu|_U+\mu^2R_\mu)^{-1}b,
 &g_\mu&=E_\mu^{-1}F_\mu p_\mu,\label{eq:coupling-pg}\\
 q_\mu&=(C_\mu|_U+\mu^2R_\mu)^{-1}
 (p_\mu+F_\mu^\top E_\mu^{-1}g_\mu).\nonumber
\end{align}
Then the following identities are exact:
\begin{align}
 H_\mu^{-1}b
 &=\mu\begin{pmatrix}p_\mu\\-g_\mu\end{pmatrix},
 \label{eq:coupling-first-inverse}\\
 H_\mu^{-2}b
 &=\begin{pmatrix}
       \mu^2q_\mu\\
       -E_\mu^{-1}g_\mu-\mu^2E_\mu^{-1}F_\mu q_\mu
    \end{pmatrix}.
 \label{eq:coupling-second-inverse}
\end{align}
For a norm-tight encoding,
$\alpha_H=\Theta(\norm{H_\mu})=\Theta(\mu^{-1})$,
\begin{equation}
 \xi(H_\mu,b)=\Theta(1),\qquad
 \rho(H_\mu,b)=\Theta\!\left(1+\frac{\norm{g_\mu}}{\mu^2}\right).
 \label{eq:coupling-filtering-law}
\end{equation}
In particular, with
\begin{equation}
 g=E^{-1}FC^{-1}b,
 \label{eq:limiting-coupling}
\end{equation}
filtering improves asymptotically on the full condition-number scale if and
only if $g=0$.  If $g\ne0$, then
$\rho(H_\mu,b)=\Theta(\mu^{-2})=\Theta(\kappa(H_\mu))$; if $g=0$, then
$\rho=O(\mu^{-1})=O(\sqrt{\kappa(H_\mu)})$, with bounded cost precisely when
$g_\mu=O(\mu^2)$.
\end{theorem}

\begin{proof}
Write a vector in $U\oplus K$ coordinates.  Multiplying
$H_\mu(u,v)^\top=b$ by $\mu$, the lower block equation gives
$v=-E_\mu^{-1}F_\mu u$.  Substitution in the upper equation gives
$(C_\mu|_U+\mu^2R_\mu)u=\mu b$, proving
\eqref{eq:coupling-first-inverse}.

Apply the same elimination to
$H_\mu(u,v)^\top=H_\mu^{-1}b$.  Its lower equation is
$F_\mu u+E_\mu v=-g_\mu$, so
$v=-E_\mu^{-1}g_\mu-E_\mu^{-1}F_\mu u$.  The upper equation then becomes
\[
 (C_\mu|_U+\mu^2R_\mu)u
 =\mu^2(p_\mu+F_\mu^\top E_\mu^{-1}g_\mu),
\]
which proves \eqref{eq:coupling-second-inverse}.

Continuity and positive definiteness give
$p_\mu\to p=C^{-1}b\ne0$.  Moreover,
$q_\mu\to C^{-1}(p+F^\top E^{-1}g)$, and the vector in parentheses is
nonzero because its inner product with $p$ is
$\norm p^2+\norm{E^{-1}Fp}^2>0$.  Equations
\eqref{eq:coupling-first-inverse}--\eqref{eq:coupling-second-inverse}
therefore imply
\[
 \norm{H_\mu^{-1}b}=\Theta(\mu),\qquad
 \norm{H_\mu^{-2}b}=\Theta(\norm{g_\mu}+\mu^2).
\]
For the lower bound, write the second inverse as $(a_\mu,v_\mu)$,
where $a_\mu=\mu^2q_\mu$. Its lower block equation gives
$g_\mu=-E_\mu v_\mu-F_\mu a_\mu$. Boundedness of $E_\mu,F_\mu$
controls $\norm{g_\mu}$ by $\norm{(a_\mu,v_\mu)}$, and
$q_\mu\to q\ne0$ controls $\mu^2$ by the same norm. This also covers
cancellation between the two terms in the lower block.
Substitution in \eqref{eq:dls-two-parameters} proves
\eqref{eq:coupling-filtering-law}.  Finally $g_\mu\to g$.  Thus $g\ne0$
gives the full $\mu^{-2}$ scale.  Under strict complementarity the LP
central path, and hence $C_\mu,D_\mu,F_\mu,E_\mu$ and $g_\mu$, is analytic
at $\mu=0$.  Halick\'a proves this boundary analyticity for LP and the
corresponding result for SDP
\cite{halicka1999-analytical-properties-central-path-linear-programming,
halicka2002-analyticity-of-the-central-path}. If $g=0$, differentiability
alone gives $g_\mu=O(\mu)$, hence $\rho=O(\mu^{-1})$. More precisely,
let $g'_0$ denote the derivative of the coupling at zero. If $g'_0\ne0$,
then $\norm{g_\mu}=\Theta(\mu)$ and $\rho=\Theta(\mu^{-1})$.
If $g'_0=0$, the analytic Taylor remainder gives $g_\mu=O(\mu^2)$
and bounded $\rho$. Thus boundedness is equivalent to $g_0=g'_0=0$
for this analytic family. The norm law itself also proves boundedness
exactly when $g_\mu=O(\mu^2)$.
\end{proof}

Equivalently, the obstruction is
$\Pi_KD_0(C|_U)^{-1}b\ne0$.  The exceptional limiting right-hand-side space
is $C(\ker F)$, and the obstruction vanishes for every $b\in U$ exactly when
$F=0$, meaning that $D_0$ reduces $U\oplus K$.  Analyticity sharpens the
continuous-rate analysis: after $g=0$, the only divergent intermediate rate
is a first-order term $g_\mu=\Theta(\mu)$, which gives
$\rho=\Theta(\mu^{-1})$; every higher integer order gives bounded $\rho$.
For a non-tight encoding, the exact scale carried by the proof is
\begin{equation}
 \rho(H_\mu,b)=\Theta\!\left(
 \alpha_H\left(\mu+\frac{\norm{g_\mu}}\mu\right)\right).
 \label{eq:coupling-nontight}
\end{equation}

\begin{proposition}[Sparse sharpness witnesses]
\label{prop:coupling-witnesses}
There are two row- and column-two LPs with the same
$\kappa(H_\mu)\sim(2\mu^2)^{-1}$ and $\xi=\Theta(1)$, but respectively
$\rho=\Theta(\mu^{-2})$ and $\rho=1$.
\end{proposition}

\begin{proof}
For the coupled instance, take
\begin{equation}
 \min\ x_2+x_3
 \quad\text{subject to}\quad
 \begin{pmatrix}1&1&0\\0&1&-1\end{pmatrix}x
 =\begin{pmatrix}1\\0\end{pmatrix},\qquad x\ge0.
 \label{eq:coupled-small-lp}
\end{equation}
Its primal optimum and optimal partition are unique.  The optimal duals have
$y_1=0$, $y_2\in[-1,1]$; their strictly complementary members have
$y_2\in(-1,1)$.  The displayed slack
$s^*=(0,1,1)$ corresponds to $y^*=0$, the dual analytic center and hence the
central-path limit, while $x^*=e_1$.  The central point is
$x=(1-t,t,t)$, where
\[
 t=\frac{2+3\mu-\sqrt{4-4\mu+9\mu^2}}4=\mu+O(\mu^2).
\]
With $\theta_1=(1-t)^2/\mu$ and $\theta=t^2/\mu$,
\[
 H_\mu=
 \begin{pmatrix}\theta_1+\theta&\theta\\\theta&2\theta\end{pmatrix},
 \qquad \kappa(H_\mu)\sim\frac1{2\mu^2}.
\]
Here $C=1$, $E=2$, $F=1$, and the centering right-hand side is
$b=(1,0)^\top$, so $p=1$ and $g=1/2$.  Direct inversion gives
\[
 H_\mu^{-1}b=\frac{(2,-1)^\top}{2\theta_1+\theta},
 \qquad H_\mu^{-2}b\longrightarrow(0,-1/4)^\top.
\]
For $\alpha_H=\norm{H_\mu}$,
$\xi\to\sqrt5/2$ and $\mu^2\rho\to1/(2\sqrt5)$.  The normalized solution
is exactly $(2,-1)/\sqrt5$. Since the slow eigenspace tends to the
second coordinate axis, its slow-cluster amplitude tends to $1/\sqrt5$.

For the decoupled instance, replace the constraint matrix by
\[
 \begin{pmatrix}1&0&0\\0&1&-1\end{pmatrix}
\]
and keep the same objective and right-hand side.  It has the same unique
primal optimum, optimal partition, dual-optimal interval, central-path dual
limit, and it also satisfies the proposition's row/column-two sparsity bound,
but, for $0<\mu\le1/\sqrt2$ and $\alpha_H=\norm{H_\mu}$,
\[
 H_\mu=\diag(\mu^{-1},2\mu),\qquad
 \kappa(H_\mu)=\frac1{2\mu^2},\qquad \xi=\rho=1.
\]
Thus neither $\kappa$ nor the first-inverse norm distinguishes the two
instances; active-to-inactive coupling does.
\end{proof}

\paragraph{Formal verification and scope.}
The Lean development in \path{formal/QipmFormal/Coupling/} checks the LP
normal-matrix split and centering right-hand side, block elimination and
the actual inverse identities, noncancellation, and the exceptional
right-hand-side space. It derives the filtering law from continuous block
data and checks the differentiable and analytic rate consequences.
The sparse witnesses are checked from their primal and dual equations,
with Euclidean norms and the displayed normalization qualifications.
The claim map and reproduction command are in \path{formal/COUPLING.md}.
The audit permits only \texttt{propext}, \texttt{Classical.choice}, and
\texttt{Quot.sound}. The general existence and boundary analyticity of
the LP central path remain mathematical inputs; the development verifies
the block regularity consequences rather than formalizing those external
theorems. It does not verify the quantum solver, construct its block
encoding, or establish the later truncation and spectral-dimension claims.

\subsection{Truncation at the slow-shell boundary}
\label{sec:beyond-kappa-truncation}

Let $P_s(\mu)$ project onto the eigenvalues of $H_\mu$ of order
$\Theta(\mu)$, which become $\Theta(\mu^2)$ after norm normalization.
For a general nonzero $f$, define its slow right-hand-side amplitude and
slow solution mass by
\[
 q_s=\frac{\norm{P_sf}}{\norm f},\qquad
 p_s=\frac{\norm{P_sH_\mu^{-1}f}^2}{\norm{H_\mu^{-1}f}^2}.
\]

\begin{proposition}[Two-cluster truncation transition]
\label{prop:dls-truncation-transition}
If the fast-cluster component of $f$ is bounded below, then
\begin{equation}
 p_s=\Theta\!\left(\frac{q_s^2}{\mu^4+q_s^2}\right).
 \label{eq:slow-solution-mass}
\end{equation}
Consequently $q_s=\Theta(\mu^2)$ is the critical boundary.  There the slow
solution mass is constant and fixed-accuracy DLS truncation must retain the
full $\Theta(\mu^{-2})$ normalized inverse scale for every target error below
that mass.  If $q_s=o(\mu^2)$, truncation can eventually discard the slow
cluster at fixed accuracy; filtering need not yet be bounded.  Under a tight
encoding, $q_s=O(\mu^4)$ is sufficient for bounded ideal filtering cost.
At exact centering, generic coupling gives $q_s=\Theta(\mu^2)$.
\end{proposition}

\begin{proof}
On the unnormalized fast and slow clusters, $H_\mu^{-1}$ has scales
$\Theta(\mu)$ and $\Theta(\mu^{-1})$, respectively.  The lower bound on the
fast part of $f$ therefore gives fast solution norm $\Theta(\mu\norm f)$,
while the slow solution norm is $\Theta(\mu^{-1}q_s\norm f)$.  Dividing the
squared slow norm by the sum proves \eqref{eq:slow-solution-mass}.

For the exact-centering vector $b\in U$, standard separated-subspace
perturbation, or direct substitution in the block eigenvector equation,
gives in the leading $K$ component
\begin{equation}
 P_s(\mu)b=-\mu^2D_{0,KU}C^{-1}b+o(\mu^2).
 \label{eq:exact-center-slow-rhs}
\end{equation}
It is therefore $\Theta(\mu^2)$ when the limiting coupling is nonzero.
Equation \eqref{eq:slow-solution-mass} proves the truncation claims.  One
additional inverse weights the slow part by another $\mu^{-2}$ relative to
the fast part, so its contribution is bounded only at
$q_s=O(\mu^4)$, which is the filtering statement.
\end{proof}

The comparison with \cref{prop:newton-rhs-alignment,cor:top-window-solve}
requires naming the two formulations.  This section uses the $m$-dimensional
normal matrix
\[
 H_{\rm NE}(\mu)=AXS^{-1}A^\top,
 \qquad
 \operatorname{spec}(H_{\rm NE}(\mu))
 \subset\Theta(\mu^{-1})\cup\Theta(\mu),
\]
whereas \cref{sec:spectra-rhs} uses the $(n-m)$-dimensional reduced Hessian
\[
 H_{\rm red}(\mu)=Z^\top X^{-2}Z,
 \qquad
 \operatorname{spec}(H_{\rm red}(\mu))
 \subset\Theta(\mu^{-2})\cup\Theta(1).
\]
When the corresponding clusters are nonempty, both have fast normalized
scale $\Theta(1)$ and slow normalized scale $\Theta(\mu^2)$, but they act on
different variables and right-hand sides.  \Cref{sec:spectra} proves an
$O(\mu^2)$ slow right-hand-side fraction for $H_{\rm red}(\mu)$.  For
$H_{\rm NE}(\mu)$, the present coupling calculation
instead gives $q_s=\Theta(\mu^2)$ only when the limiting coupling is nonzero;
the decoupled witness has $P_sb=0$.

The residual-versus-state distinction can be proved directly for the normal
equation.  Under nonzero limiting coupling, let
$z=H_{\rm NE}(\mu)^{-1}b$ and discard
$z_s=P_sz$.  Because $P_s$ is a spectral projector,
\begin{equation}
 \frac{\norm{H_{\rm NE}(\mu)(z-z_s)-b}}{\norm b}
 =\frac{\norm{P_sb}}{\norm b}=q_s=\Theta(\mu^2),
 \label{eq:normal-window-residual}
\end{equation}
while \eqref{eq:slow-solution-mass} gives
$\norm{z_s}^2/\norm z^2=\Theta(1)$ under the same nonzero-coupling
hypothesis.  A state-preparation algorithm below that constant solution mass
must retain the slow cluster, but a fixed relative-residual contract may
discard it.  Thus the two formulations agree on their normalized spectral
separation, while both the formulation and the output contract matter; this
is the precise sense in which the result is consistent with
\cref{sec:spectra}.

\subsection{Effective spectral dimension}
\label{sec:effective-spectral-dimension}

Two clusters are not necessary for a sharp solver-selection law.  Normalize
$\norm H=1$ and $\norm b=1$, let
\begin{equation}
 \mathsf m_b(S)=\inner b{\mathbf1_S(H)b},
 \label{eq:spectral-measure}
\end{equation}
and let $\lambda_*$ be the smallest eigenvalue visible to $b$, with
$\kappa_*=\lambda_*^{-1}$.  We say that $(H,b)$ has effective spectral
dimension $D$ if constants $c,C>0$, independent of $\lambda_*$, satisfy
\begin{equation}
 c t^{D/2}\leq F_b(t):=\mathsf m_b([0,t])\leq C t^{D/2}
 \qquad(\lambda_*\leq t\leq1).
 \label{eq:effective-dimension-definition}
\end{equation}
This is invariant under simultaneous orthogonal changes
$(H,b)\mapsto(VHV^\top,Vb)$; it is not coordinate sparsity in disguise.
Unlike an unanchored per-shell condition, \eqref{eq:effective-dimension-definition}
controls the mass at the smallest visible eigenvalue and permits fixed gaps
between the first few eigenvalues.

\begin{theorem}[Spectral-dimension phase laws]
\label{thm:spectral-dimension-laws}
Assume the norm-tight unit-normalization convention: after $\norm H=1$, the
block-encoding normalization is $\alpha_H=\Theta(1)$, and truncation cutoffs
are measured in this same normalization.  Operationally,
\begin{equation}
 \kappa_{\rm eff}(\epsilon_{\rm ls})
 :=\inf\left\{K\in[1,\kappa_*]:
 \frac{\inner b{H^{-2}\mathbf1_{[0,1/K)}(H)b}}{\inner b{H^{-2}b}}
 \leq\epsilon_{\rm ls}^2\right\}.
 \label{eq:dls-effective-cutoff-definition}
\end{equation}
For an encoding of $H/\alpha_H$, the solver's encoded cutoff scale is
$\alpha_H\kappa_{\rm eff}$; the standing tight convention identifies these
up to constants.
Under \eqref{eq:effective-dimension-definition}, this quantity obeys
\begin{equation}
 \kappa_{\rm eff}(\epsilon_{\rm ls})=
 \begin{cases}
  \Theta_{\epsilon_{\rm ls}}(\kappa_*),&D<4,\ 0<\epsilon_{\rm ls}<1,\\
  \kappa_*^{\,1-\Theta(\epsilon_{\rm ls}^2)},&D=4,\\
  \Theta\!\left(\min\{\kappa_*,
                  \epsilon_{\rm ls}^{-4/(D-4)}\}\right),&D>4.
 \end{cases}
 \label{eq:dls-truncation-dimension-law}
\end{equation}
For $D<4$, the hidden constant may depend on the fixed error and on the
cumulative-law constants.  A sublinear power of $\kappa_*$ is possible only
when the allowed error tends to one as $\kappa_*\to\infty$.  At $D=4$, a
definite leading cumulative coefficient gives exponent
$1-\epsilon_{\rm ls}^2+o(1)$; under only
\eqref{eq:effective-dimension-definition},
$1-\Theta(\epsilon_{\rm ls}^2)$ is sharp, and the result saturates at
$\Theta(\kappa_*)$ when
$\epsilon_{\rm ls}^2\log\kappa_*=O(1)$.  The ideal filtering
parameter is
\begin{equation}
 \rho=
 \begin{cases}
  \Theta(\kappa_*),&D<4,\\
  \Theta(\kappa_*/\sqrt{\log\kappa_*}),&D=4,\\
  \Theta(\kappa_*^{(8-D)/4}),&4<D<8,\\
  \Theta(\sqrt{\log\kappa_*}),&D=8,\\
  \Theta(1),&D>8.
 \end{cases}
 \label{eq:dls-filter-dimension-law}
\end{equation}
Thus fixed-accuracy truncation becomes independent of mesh conditioning above
$D=4$, whereas filtering does so only above $D=8$.  For a non-tight encoding,
both the effective cutoff scale and every filtering branch acquire a factor
$\alpha_H$; in particular the $D>8$ filtering cost is
$\Theta(\alpha_H)$ rather than $\Theta(1)$.
\end{theorem}

\begin{proof}
For $p>0$, put $Z_p=\inner b{H^{-p}b}$ and write $s=D/2$.  Abel summation,
equivalently Stieltjes integration by parts, gives the exact layer-cake
identity
\begin{equation}
 Z_p=1+p\int_{\lambda_*}^1 t^{-p-1}F_b(t)\,dt.
 \label{eq:spectral-abel-identity}
\end{equation}
Using $F_b(t)=\Theta(t^s)$ in this integral gives
\begin{equation}
 Z_p=
 \begin{cases}
  \Theta(\kappa_*^{p-D/2}),&D<2p,\\
  \Theta(\log\kappa_*),&D=2p,\\
  \Theta(1),&D>2p.
 \end{cases}
 \label{eq:spectral-moment-law}
\end{equation}
The squared norm of $H^{-1}b$ is $Z_2$.  If a cutoff discards eigenvalues
strictly below $u=1/K$, as in \eqref{eq:dls-effective-cutoff-definition}, its
discarded squared solution norm is
\begin{equation}
 T_2(u)=u^{-2}F_b(u^-)
       +2\int_{\lambda_*}^{u}t^{-3}F_b(t)\,dt,
 \label{eq:spectral-partial-abel}
\end{equation}
where $F_b(u^-)=\mathsf m_b([0,u))$ is the left limit.
For $D<4$, fix a sufficiently small constant $u_0>0$.  Dyadically summing
the retained moment on $[u,1]$ and comparing it with $Z_2$ gives
\begin{equation}
 \frac{\inner b{H^{-2}\mathbf1_{[u,1]}(H)b}}{Z_2}
 =\Theta\!\left(
   \left(\frac{\lambda_*}{u}\right)^{2-D/2}\right)
 \qquad(\lambda_*\leq u\leq u_0).
 \label{eq:subcritical-retained-mass}
\end{equation}
Indeed, the upper bound follows by geometric summation, while a shell
$[u,\gamma_{\rm sh}u]$ with a fixed $\gamma_{\rm sh}>1$ chosen from the two
cumulative-law constants has mass $\Theta(u^{D/2})$ and gives the matching
lower bound.  Requiring the
retained fraction to be at least $1-\epsilon_{\rm ls}^2$ therefore gives
\begin{equation}
 u=\Theta\!\left(
 \lambda_*(1-\epsilon_{\rm ls}^2)^{-2/(4-D)}\right)
 =\Theta_{\epsilon_{\rm ls}}(\lambda_*),
 \label{eq:subcritical-cutoff}
\end{equation}
for every fixed $\epsilon_{\rm ls}<1$, proving the first line.  A cutoff
$u=\lambda_*^{1-\beta}$ for fixed $\beta>0$ has retained fraction tending to
zero, so it is
possible only when $\epsilon_{\rm ls}\to1$.  For $D=4$, both
\eqref{eq:spectral-abel-identity} and \eqref{eq:spectral-partial-abel} are
logarithmic.  Discarding an $\epsilon_{\rm ls}^2$ fraction yields
$K=\kappa_*^{1-\Theta(\epsilon_{\rm ls}^2)}$, with the stated refinement when the
leading cumulative coefficient exists.  For $D>4$, the discarded tail is
$\Theta(K^{-(D-4)/2})$ relative to the constant total.  Setting it to
$\epsilon_{\rm ls}^2$ and capping at $\kappa_*$ proves
\eqref{eq:dls-truncation-dimension-law}.

For filtering, $x=H^{-1}b$ and hence
\[
 \rho^2=\alpha_H^2\frac{\norm{H^{-1}x}^2}{\norm x^2}
 =\alpha_H^2\frac{Z_4}{Z_2}=\Theta\!\left(\frac{Z_4}{Z_2}\right).
\]
Applying \eqref{eq:spectral-moment-law} with $p=4$ in the numerator and
$p=2$ in the denominator gives all five cases in
\eqref{eq:dls-filter-dimension-law}.
\end{proof}

Local conventions in the remaining witnesses are as follows:
$e_0,e_{e_1},e_f$ are coordinate basis vectors, whereas $e_k$ in the
marked-diagonal construction is the $k$-dimensional all-ones vector; the
symbols $\delta_{\rm c},\delta_Q,\delta_{\mathcal G}$ denote, respectively,
the corrector perturbation, projector error, and expander contrast.

\begin{proposition}[Constant-degree torus witnesses]
\label{prop:torus-dimension-witnesses}
Let $L$ be even and $T_L^d$ the periodic $d$-dimensional $L$-grid.  With
its combinatorial Laplacian $\Delta_{T_L^d}$, set
\begin{equation}
 H_{L,d}=\frac{L^{-2}I+\Delta_{T_L^d}}{L^{-2}+4d}.
 \label{eq:massive-torus}
\end{equation}
Then $H_{L,d}$ is $(2d+1)$-sparse, has norm one, and has
$\kappa=1+4dL^2=\Theta(L^2)$.  A localized source $e_0$ realizes $D=d$;
the dipole $(e_0-e_{e_1})/\sqrt2$ realizes $D=d+2$.  For the dipole branch,
the resulting phases are:
\begin{center}
\begin{tabular}{@{}rrrr@{}}
\toprule
physical $d$ & effective $D$ & $\kappa_{\rm eff}$ & $\rho$ \\
\midrule
$3$ & $5$ & $\Theta(\min\{\kappa,\epsilon_{\rm ls}^{-4}\})$
            & $\Theta(\kappa^{3/4})$ \\
$4$ & $6$ & $\Theta(\min\{\kappa,\epsilon_{\rm ls}^{-2}\})$
            & $\Theta(\kappa^{1/2})$ \\
$5$ & $7$ & $\Theta(\min\{\kappa,\epsilon_{\rm ls}^{-4/3}\})$
            & $\Theta(\kappa^{1/4})$ \\
$6$ & $8$ & $\Theta(\min\{\kappa,\epsilon_{\rm ls}^{-1}\})$
            & $\Theta(\sqrt{\log\kappa})$ \\
$\geq7$ & $>8$ & mesh-independent at fixed $\epsilon_{\rm ls}$
                   & mesh-independent \\
\bottomrule
\end{tabular}
\end{center}
\end{proposition}

\begin{proof}
Fourier modes diagonalize the torus Laplacian, with
\[
 \lambda_k=\frac{L^{-2}+4\sum_{j=1}^d\sin^2(\pi k_j/L)}{L^{-2}+4d}.
\]
For the localized source, every squared Fourier overlap is $L^{-d}$.
Lattice-volume bounds, first for
$\lambda_k\asymp L^{-2}(1+\norm k^2)$ at low frequency and then by the
same product-counting bounds on the remaining fixed spectral range, give
\[
 \mathsf m_{e_0}([0,t])=\Theta(t^{d/2})
 \qquad(\lambda_*\leq t\leq1).
\]
Thus the cumulative hypothesis, including its bottom anchoring, holds with
$D=d$.  Equivalently, Abel summation gives
\[
 \sum_k\lambda_k^{-p}=
 \begin{cases}
  \Theta(L^{2p}),&d<2p,\\
  \Theta(L^{2p}\log L),&d=2p,\\
  \Theta(L^d),&d>2p,
 \end{cases}
\]
The dipole Fourier symbol is $1-e^{2\pi i k_1/L}$.  It vanishes on modes
with $k_1=0$, so no pointwise comparison with $\lambda_k$ is valid.  Instead
sum over permutation-symmetric spectral balls.  Symmetry gives
\[
 \sum_{\lambda_k\leq t}\sin^2(\pi k_1/L)
 =\frac1d\sum_{\lambda_k\leq t}\sum_{j=1}^d
       \sin^2(\pi k_j/L).
\]
Together with the lattice-volume estimate, this aggregated identity yields
\[
 \mathsf m_{(e_0-e_{e_1})/\sqrt2}([0,t])
 =\Theta(t^{(d+2)/2})
\]
from the smallest visible eigenvalue through $t=1$.  At the bottom, the
first visible modes have squared overlap $\Theta(L^{-(d+2)})$, exactly
$\Theta(\lambda_*^{(d+2)/2})$; the constant mode has zero overlap and
$\lambda_*=\Theta(L^{-2})$.  Hence the repaired cumulative definition is
satisfied with $D=d+2$.  The table follows by substituting $D=d+2$ in
\cref{thm:spectral-dimension-laws}.
\end{proof}

The two source types occur as actual QIPM correctors, not merely abstract
linear systems.

\begin{proposition}[Positive full-step feasible corrector]
\label{prop:torus-feasible-corrector}
Let $B_T$ be an oriented incidence matrix of $T_L^d$ and set
\begin{equation}
 \mathcal A=[B_T,I],\qquad
 W=\diag(I_E,L^{-2}I_V),\qquad
 \mathcal A W\mathcal A^\top=\Delta_{T_L^d}+L^{-2}I.
 \label{eq:torus-ipm-factor}
\end{equation}
For $d\geq2$, choose perpendicular torus edges $f,g$ meeting at a vertex,
fix $\tau>0$ and $0<\delta_{\rm c}\leq1/4$, and retain $x_i/s_i=w_i$, where
$w_i$
is the corresponding diagonal entry of $W$.  Set
\begin{equation}
 p_i=x_is_i=
 \begin{cases}
  \tau(1+\delta_{\rm c}),&i=f,\\
  \tau(1-\delta_{\rm c}),&i=g,\\
  \tau,&\text{otherwise},
 \end{cases}
 \quad
 x_i=\sqrt{p_iw_i},\quad s_i=\sqrt{p_i/w_i}.
 \label{eq:torus-perturbed-products}
\end{equation}
With $b_{\rm LP}=\mathcal Ax$, $c=s$, and $y=0$, this is primal and dual
feasible, has average complementarity $\tau$, and has centrality
$\norm{XSe-\tau e}/\tau=\sqrt2\delta_{\rm c}$.  Its full $\sigma=1$ feasible
centering step is positive, has a normal-equation right-hand side of
effective dimension $D=d+2$, preserves average complementarity, and satisfies
\begin{equation}
 \frac{\norm{X^+S^+e-\tau e}}\tau
 \leq\frac{\delta_{\rm c}^2}{1-\delta_{\rm c}^2}.
 \label{eq:torus-quadratic-centering}
\end{equation}
\end{proposition}

\begin{proof}
The definitions make feasibility and the centrality calculation immediate.
The feasible corrector equations are
\[
 \mathcal A\Delta x=0,
 \quad \mathcal A^\top\Delta y+\Delta s=0,
 \quad S\Delta x+X\Delta s=\tau e-XSe.
\]
Eliminating $\Delta x,\Delta s$ gives
\begin{equation}
 (\mathcal AW\mathcal A^\top)\Delta y
 =-\mathcal AS^{-1}(\tau e-XSe).
 \label{eq:torus-corrector-normal}
\end{equation}
After normalization, its right-hand side is proportional to
\[
 \frac{\mathcal A_f}{\sqrt{1+\delta_{\rm c}}}
 -\frac{\mathcal A_g}{\sqrt{1-\delta_{\rm c}}}.
\]
Its low-frequency Fourier symbol vanishes linearly and has nonzero gradient,
so its cumulative spectral weight is $\Theta(t^{(d+2)/2})$.  Explicitly, for
any nonzero
linear form $\ell$, lattice-ball comparison gives
$\sum_{\norm k\leq r}|\ell(k)|^2=\Theta(r^{d+2})$; applying this to the
linear part of the symbol gives the claimed cumulative spectral weight.

It remains to prove full-step safety.  Put
$u=X^{-1}\Delta x$, $v=S^{-1}\Delta s$, and
$P=\diag(p)$.  The Newton
equations imply
\begin{equation}
 P(u+v)=\tau e-p,\qquad
 u^\top Pv=\Delta x^\top\Delta s=0.
 \label{eq:torus-relative-step}
\end{equation}
Thus $P^{1/2}u$ and $P^{1/2}v$ are orthogonal and sum to
$P^{-1/2}(\tau e-p)$, whose squared norm is
$2\tau\delta_{\rm c}^2/(1-\delta_{\rm c}^2)$.  Since
$p_i\geq\tau(1-\delta_{\rm c})$, for either $z=u,v$,
\[
 |z_i|^2\leq
 \frac{2\delta_{\rm c}^2}
 {(1-\delta_{\rm c}^2)(1-\delta_{\rm c})}<1
 \qquad(\delta_{\rm c}\leq1/4).
\]
Hence $x+\Delta x>0$ and $s+\Delta s>0$.  Moreover
\[
 p_i^+=p_i(1+u_i)(1+v_i)=\tau+p_i u_iv_i.
\]
Equation \eqref{eq:torus-relative-step} gives
$\sum_i p_iu_iv_i=0$, so the average remains $\tau$.  Finally,
\[
 \norm{P(u\circ v)}
 \leq\norm{P^{1/2}u}\norm{P^{1/2}v}
 \leq\frac12\norm{P^{-1/2}(\tau e-p)}^2
 =\frac{\tau\delta_{\rm c}^2}{1-\delta_{\rm c}^2},
\]
which proves \eqref{eq:torus-quadratic-centering}.
\end{proof}

The localized branch is also QIPM-native, but through an infeasible
corrector.  At complementarity $\tau$, set
$x_i=\sqrt{\tau w_i}$, $s_i=\sqrt{\tau/w_i}$ and
$b_{\rm LP}=\mathcal Ax$.  Give the exactly centered
($x_is_i=\tau>0$), primal-feasible
point a one-coordinate dual residual $r_d=\eta e_f$.  A $\sigma=1$
infeasible feasibility corrector has
\begin{equation}
 (\mathcal AW\mathcal A^\top)\Delta y=\eta\mathcal AWe_f.
 \label{eq:torus-infeasible-corrector}
\end{equation}
A ground variable at vertex zero gives normalized right-hand side $e_0$ and
$D=d$; a torus edge gives a dipole and $D=d+2$.

These witnesses are solver-selection examples, not end-to-end speedups.  The
torus already admits FFT and multigrid solvers.  The constructions describe
a single corrector rather than an IPM trajectory, and do not charge block
encoding, norm estimation, state or observable readout, or classical iterate
recovery.

\subsection{Stable active-subspace equilibration}
\label{sec:stable-active-equilibration}

The negative coupling law has a positive structural counterpart.  Suppose
the active range is available through its orthogonal projector $Q$, and put
$R=I-Q$.  The next theorem directly assembles a left-equilibrated operator;
this direct assembly, rather than multiplication of normalized encodings, is
the source of the gain.

\begin{theorem}[Stable split equilibration]
\label{thm:stable-split-equilibration}
Suppose
\[
 H_\mu=\mu^{-1}C_\mu+\mu D_\mu,\qquad 0<\mu\leq1,
\]
where, for constants independent of $\mu$,
\[
 C_\mu=QC_\mu Q,\quad
 c_-Q\preceq C_\mu\preceq c_+Q,\quad
 D_\mu\succeq0,\quad \norm{D_\mu}\leq d_+,\quad
 RD_\mu R\succeq d_-R.
\]
Let $P_\mu=\mu Q+\mu^{-1}R$ and assemble
\begin{equation}
 M_\mu=C_\mu+RD_\mu+\mu^2QD_\mu
 =C_\mu+D_\mu-(1-\mu^2)QD_\mu.
 \label{eq:stable-split-operator}
\end{equation}
Then $M_\mu=P_\mu H_\mu$ exactly and
$\norm{M_\mu}+\norm{M_\mu^{-1}}=O(1)$.  For every
$f\in U=\operatorname{range}Q$,
\begin{equation}
 M_\mu^{-1}f=\mu^{-1}H_\mu^{-1}f,
 \label{eq:stable-state-preservation}
\end{equation}
so the normalized Newton state is preserved exactly, without recovery or
postselection.  If $C_\mu,D_\mu,Q$ have block-encoding normalizations
$\alpha_C,\alpha_D,1$, direct LCU assembly has normalization
\begin{equation}
 \Gamma_\mu\leq\alpha_C+(2-\mu^2)\alpha_D.
 \label{eq:stable-lcu-normalization}
\end{equation}
Thus bounded component normalizations give a uniformly conditioned,
uniformly normalized conventional QLS solve through the Hermitian dilation
below, and the ideal DLS filtering parameter is $O(\Gamma_\mu)$ as well.
\end{theorem}

The left-equilibrated $M_\mu$ need not be Hermitian.  To use the Hermitian
QLS contract of \cref{def:qls}, encode its standard dilation
\begin{equation}
 \mathcal H(M_\mu)=
 \begin{pmatrix}0&M_\mu\\M_\mu^\dagger&0\end{pmatrix}.
 \label{eq:stable-hermitian-dilation}
\end{equation}
The dilation has the same block-encoding normalization $\Gamma_\mu$ and
the same singular condition number as $M_\mu$.  Moreover,
$\mathcal H(M_\mu)^{-1}(f,0)^\top=(0,M_\mu^{-1}f)^\top$, so selecting the
second register preserves the solution normalization and requires no
postselection beyond the fixed register label.

\begin{proof}
Since $RC_\mu=0$,
\[
 P_\mu H_\mu=C_\mu+\mu^2QD_\mu+RD_\mu=M_\mu.
\]
In $U\oplus U^\perp$ blocks,
\[
 M_\mu=
 \begin{pmatrix}
 C_\mu+\mu^2D_{UU}&\mu^2D_{UK}\\
 D_{KU}&D_{KK}
 \end{pmatrix}.
\]
The Schur complement of $D_{KK}$ is
\[
 G_\mu=C_\mu+\mu^2
 (D_{UU}-D_{UK}D_{KK}^{-1}D_{KU})\succeq c_-I_U,
\]
because the parenthesized Schur complement is positive semidefinite.
Also $D_{KK}\succeq d_-I_{U^\perp}$.  The block inverse formula therefore
bounds every block of $M_\mu^{-1}$ using only $c_-,d_-,d_+$, while
$\norm{M_\mu}\leq c_++2d_+$.

For state preservation, let $x_\mu=H_\mu^{-1}f$.  Because $f\in U$,
$P_\mu f=\mu f$, and hence
$M_\mu x_\mu=P_\mu H_\mu x_\mu=\mu f$.  This proves
\eqref{eq:stable-state-preservation}.  Finally, applying LCU to the second
form in \eqref{eq:stable-split-operator} gives coefficient weight
$\alpha_C+\alpha_D+(1-\mu^2)\alpha_D$, proving
\eqref{eq:stable-lcu-normalization}.  Uniform conditioning bounds
$\Gamma_\mu\norm{M_\mu^{-1}}$ and its restriction to any solution state,
which proves the two solver claims.
\end{proof}

The split also supplies a projector-error firewall.  If
$\norm{\widehat Q-Q}\leq\delta_Q$ and one assembles
\[
 \widehat M_\mu=C_\mu+D_\mu-(1-\mu^2)\widehat QD_\mu,
\]
then $\norm{\widehat M_\mu-M_\mu}\leq\delta_Qd_+$.  For a sufficiently small
constant $\delta_Q$, the Neumann perturbation bound and
\cref{thm:stable-split-equilibration} give
\begin{equation}
 \frac{\norm{\widehat M_\mu^{-1}f-M_\mu^{-1}f}}
      {\norm{M_\mu^{-1}f}}=O(\delta_Q),
 \label{eq:projector-firewall}
\end{equation}
uniformly in $\mu$.  Thus $\delta_Q=O(\epsilon_{\rm ls})$ suffices.  In
contrast, forming $\widehat P_\mu H_\mu$ from the raw normal matrix loses
this stability.  For $H_\mu=\diag(\mu^{-1},\mu)$ and a projector rotated by
angle $\theta$, its $(2,1)$ entry is
$(1-\mu^{-2})\sin\theta\cos\theta$; uniform normalization requires
$\sin\theta=O(\mu^2)$.  Given a block encoding of $A_B/\alpha_B$, QSVT can
construct $Q$ to the needed accuracy in
\begin{equation}
 O\!\left(\frac{\alpha_B}{\sigma_{\min}^+(A_B)}
 \log(1/\epsilon_{\rm ls})\right)
 \label{eq:active-projector-cost}
\end{equation}
queries, charging active-basis conditioning but no polynomial power of
$\mu$.

\begin{proposition}[Separately normalized composition cannot help]
\label{prop:no-composed-preconditioning}
Let $H/\alpha_H$ and an invertible $P/\alpha_P$ be separately encoded, with
$\alpha_H\geq\norm H$ and $\alpha_P\geq\norm P$, and compose their encodings.
For $x=H^{-1}f$, the DLS parameter of the resulting system
$PHx=Pf$ satisfies
\begin{equation}
 \rho_{\rm prod}
 =\alpha_H\alpha_P
  \frac{\norm{P^{-\dagger}H^{-1}x}}{\norm x}
 \geq \alpha_H\frac{\norm{H^{-1}x}}{\norm x}
 =\rho_H.
 \label{eq:no-composition-law}
\end{equation}
\end{proposition}

\begin{proof}
Equation \eqref{eq:no-composition-law} uses the general non-Hermitian
parameter \eqref{eq:dls-general-filtering-parameter}, because
$(PH)^{-\dagger}=P^{-\dagger}H^{-1}$.  Since
$\alpha_P\geq\norm P$, every
vector $y$ satisfies
$\alpha_P\norm{P^{-\dagger}y}\geq\norm y$.  Apply this with
$y=H^{-1}x$.  The displayed expression is precisely the filtering
parameter for normalization $\alpha_H\alpha_P$, proving the result.
\end{proof}

This normalization obstruction is known in quantum preconditioning:
Lapworth and S\"underhauf show explicitly how multiplying separately
normalized block encodings can erase a classical gain, while Nie, Lai, and
An show how algebraic regrouping can restore it for Pauli structure
\cite{leigh2025-preconditioned-block-encodings-quantum-linear,hantao2026-pauli-structured-preconditioning-quantum-linear}.
Here separate composition has $\alpha_H,\alpha_P=\Theta(\mu^{-1})$, whereas the
direct split has $\Gamma_\mu=O(1)$, a $\Theta(\mu^{-2})$ normalization
compression.

\begin{proposition}[Constant-sparse expander occurrence]
\label{prop:stable-expander-witness}
There is a family of constant-sparse standard-form LPs for which the
unpreconditioned exact-centering solve has $\rho=\Theta(\mu^{-2})$, while
\eqref{eq:stable-split-operator} has $O(1)$ singular conditioning and
$O(1)$ block-encoding normalization and prepares exactly the same normalized
Newton state independently of $\mu$.  The uniform bounds hold on
$0<\mu\leq\mu_0$ with $\mu_0$ and all constants depending only on the degree
$d$, the Laplacian gap, and $\delta_{\mathcal G}$, not on $m$.
\end{proposition}

\begin{proof}
Let $\mathcal G=(\mathcal L,\mathcal R,\mathcal E)$ be a connected balanced
$d$-regular bipartite expander on $m$ vertices, with constant $d$ and a
constant Laplacian gap.  Orient each edge from $\mathcal L$ to $\mathcal R$,
let $B_{\mathcal G}$ be its signed incidence matrix, and put
\[
 r=\mathbf1_{\mathcal R},\qquad \zeta=2r-e,\qquad
 u_0=m^{-1/2}e.
\]
Here $r$ is the right-part indicator and $e$ is the all-ones vector of the
dimension required by context; these symbols are local to this witness and
are distinct from the coupling variables above.  Then
$B_{\mathcal G}^\top r=e$ and
$B_{\mathcal G}e=L_{\mathcal G}r=d\zeta$.  Fix
$0<\delta_{\mathcal G}<1$, set
$\omega_i=1+\delta_{\mathcal G}\zeta_i$ and
$a_i=\sqrt{2/\omega_i}$, and consider
\begin{equation}
 \begin{aligned}
 \min_{x,u,v\geq0}\quad&\sum_{i=1}^m a_i(u_i+v_i)\\
 \text{subject to}\quad&B_{\mathcal G}x+u-v=B_{\mathcal G}e.
 \end{aligned}
 \label{eq:expander-lp}
\end{equation}
Its constraint matrix $[B_{\mathcal G},I,-I]$ has row and column sparsity
$O(d)$.

The point $x=e,u=v=te$ is strictly primal feasible for $t>0$.  Taking
$y=-\varepsilon r$ for sufficiently small $\varepsilon>0$
gives strictly positive dual slacks.  Nonnegativity of the objective shows
that
\[
 x^*=e,\quad u^*=v^*=0,\quad y^*=0,
 \qquad s_x^*=0,\quad s_u^*=s_v^*=a
\]
is optimal and strictly complementary.  Its active face is bounded: if
$h\geq0$ and $B_{\mathcal G}h=0$, any left-vertex row forces all incident
components of $h$ to vanish.  Also $e$ is the primal face's analytic center
because $B_{\mathcal G}^\top r=e$.

For completeness, the dual limit is also unique even though the dual optimal
face is not a singleton.  Dual feasibility and optimality imply
$B_{\mathcal G}^\top y\leq0$ and
$e^\top B_{\mathcal G}^\top y=0$, hence
$B_{\mathcal G}^\top y=0$.  Connectedness gives $y=ce$, and the remaining
dual inequalities give
\[
 \{y:\ y\text{ is dual optimal}\}
 =\{ce:\ |c|\leq\min_i a_i\}.
\]
On its relative interior the dual analytic-center objective is
$\sum_i\log(a_i^2-c^2)$, which is strictly maximized at $c=0$.  Thus
$y(\mu)\to0$, and the displayed primal--dual point is the central-path
limit.

Along the path,
\begin{align}
 H_\mu&=\mu^{-1}C_\mu+\mu D_\mu,\nonumber\\
 C_\mu&=B_{\mathcal G}\diag(x_{\mathcal E}(\mu)^2)
          B_{\mathcal G}^\top\longrightarrow L_{\mathcal G},\nonumber\\
 D_\mu&=\diag\!\left(
 \frac1{(a_i-y_i(\mu))^2}+\frac1{(a_i+y_i(\mu))^2}
 \right)\longrightarrow\diag(\omega_i).
 \label{eq:expander-normal-split}
\end{align}
Here $U=\operatorname{range}B_{\mathcal G}=e^\perp$ and
$Q=I-\ket{u_0}\bra{u_0}$.  The constants are uniform in $m$.  Write
$s_x=-B_{\mathcal G}^\top y$, $s_u=a-y$, $s_v=a+y$; the central equations
are $x_es_{x,e}=\mu$, $u_is_{u,i}=\mu$, $v_is_{v,i}=\mu$, and
$B_{\mathcal G}x+u-v=B_{\mathcal G}e$.  Seek a solution constant on edges
and on each side: $x_e=\bar x$, $y_i=y_{\mathcal L}$ on $\mathcal L$,
$y_i=y_{\mathcal R}$ on $\mathcal R$, with $\bar x=1+\mu^2w$,
$y_{\mathcal L}=\mu z_{\mathcal L}$,
$y_{\mathcal R}=\mu z_{\mathcal R}$.  Since
$\omega=1\mp\delta_{\mathcal G}$ on $\mathcal L$ and $\mathcal R$, the
system reduces to three scalar equations independent of $m$:
\[
 dw+\frac{2z_{\mathcal R}}
 {2/(1+\delta_{\mathcal G})-\mu^2z_{\mathcal R}^2}=0,
 \qquad
 -dw+\frac{2z_{\mathcal L}}
 {2/(1-\delta_{\mathcal G})-\mu^2z_{\mathcal L}^2}=0,
 \qquad
 z_{\mathcal L}-z_{\mathcal R}=\frac1{1+\mu^2w}.
\]
At $\mu=0$ they have the solution
$w=(1-\delta_{\mathcal G}^2)/(2d)$,
$z_{\mathcal L}=(1+\delta_{\mathcal G})/2$,
$z_{\mathcal R}=-(1-\delta_{\mathcal G})/2$, with Jacobian determinant
$-2d\neq0$.  The implicit function theorem gives a branch on
$0<\mu\le\mu_0(d,\delta_{\mathcal G})$ with $\bar x>0$,
$y_{\mathcal L}-y_{\mathcal R}=\mu/\bar x>0$, and $|y_i|<a_i$; the ansatz
therefore satisfies the full central-path system, and uniqueness of the
central point identifies it as the central path.  Hence
$C_\mu=\bar x^2L_{\mathcal G}$ exactly, $D_\mu$ takes two side-constant
values, and the expander gap together with
$\omega_i\in[1-\delta_{\mathcal G},1+\delta_{\mathcal G}]$ gives the
hypotheses of \cref{thm:stable-split-equilibration} with constants depending
only on $d$, the Laplacian gap, and $\delta_{\mathcal G}$.

For the exact-centering right-hand side $f=b=B_{\mathcal G}e$,
$L_{\mathcal G}(Qr)=b$ and $Qr=\zeta/2$.  Its limiting active-to-null
response is nonzero:
\begin{equation}
 u_0^\top\diag(\omega)Qr=\frac{\delta_{\mathcal G}\sqrt m}{2}
 =\delta_{\mathcal G}\norm{Qr}.
 \label{eq:expander-coupling}
\end{equation}
The coupling theorem therefore gives $\rho(H_\mu,f)=\Theta(\mu^{-2})$.
In contrast, direct assembly gives
\begin{equation}
 M_\mu=C_\mu+\mu^2D_\mu
 +(1-\mu^2)\ket{u_0}\bra{u_0}D_\mu.
 \label{eq:expander-direct-split}
\end{equation}
It has uniform singular conditioning and $O(d)$ normalization: $C_\mu$ is
a bounded-weight degree-$d$ Laplacian, $D_\mu$ is diagonal, $u_0$ is the
uniform state, and the last term is rank one.  Since $b\in U$,
$M_\mu^{-1}b=\mu^{-1}H_\mu^{-1}b$ at every $\mu$, completing the witness.
\end{proof}

The theorem assumes a supplied active partition, coherent current weights,
and direct access to $C_\mu,D_\mu,Q$.  It also charges
$\sigma_{\min}^+(A_B)$, right-hand-side preparation, norm estimation, and
the requested output.  \Cref{sec:positive-modules} gives a partition-free
construction on
the strict-complementarity tail and thereby removes the supplied-partition
assumption; that construction is not needed here.  Classical application of
\eqref{eq:expander-direct-split} is already linear-time per Krylov iteration.
The quantum statement is therefore meaningful only for a state or a few
observables; tomography or a classical iterate restores dimension-dependent
cost.

\subsection{An exact-center marked-diagonal specialization}
\label{sec:marked-diagonal-specialization}

The sparse search lower bound of Orsucci and Dunjko
\cite[Proposition~7]{orsucci2021-on-solving-classes-of-positive} can also be
realized at an exact LP center.  This gives a useful QIPM specialization,
but its marked-search core is prior art.

Choose an integer $T\geq2$, set the local dimension
$N_{\rm md}=T^8$ and $\mu=T^{-2}=N_{\rm md}^{-1/4}$, and hide
$j\in[N_{\rm md}]$.  The subscript distinguishes this dimension from the
optimal-partition set $N$ above.  Let $e_k$ denote the $k$-dimensional
all-ones vector in this construction.  Row $i$ of
$A_j\in\R^{N_{\rm md}\times2N_{\rm md}}$ has disjoint support
$(2i-1,2i)$ and values
\begin{equation}
 (\mu/2,\mu/2)\quad(i=j),\qquad
 (1,\mu-1)\quad(i\ne j).
 \label{eq:marked-diagonal-A}
\end{equation}
Because $\mu$, $\sqrt\mu$, and $\mu^{3/2}$ are rational, all entries have
magnitude at most one and $O(\log N_{\rm md})$-bit rational descriptions;
row sparsity is two and column sparsity one.  Set
\[
 b=\mu^{3/2}e_{N_{\rm md}},\qquad
 c=\sqrt\mu e_{2N_{\rm md}},\qquad
 x^0=s^0=\sqrt\mu e_{2N_{\rm md}},\qquad y^0=0.
\]
Then $A_jx^0=b$, $(A_j)^\top y^0+s^0=c$, and
$X^0S^0e=\mu e$, so this is an exact $\mu$-center.  Since
$X^0(S^0)^{-1}=I$,
\begin{equation}
 H_j=A_jA_j^\top
 =\diag(h,\ldots,\mu^2/2,\ldots,h),
 \qquad h=1+(1-\mu)^2,
 \label{eq:marked-diagonal-H}
\end{equation}
and the exact short-step right-hand side is
$f=(1-\sigma)\mu^{3/2}e_{N_{\rm md}}$.  Its normalized preparation is uniform
and independent of $j$.  The squared solution weight on the marked
coordinate is at least $4/5$, while
\[
 \begin{aligned}
 \kappa(H_j)=\Theta(\mu^{-2}),\qquad
 \xi(H_j,f)=\Theta(1),\qquad
 \rho(H_j,f)&=\Theta(\mu^{-2}),\\
 \kappa_{\rm eff}(\epsilon_{\rm ls})
 &=\frac{2h}{\mu^2}=\Theta(\mu^{-2})
 \quad\text{for every fixed }\epsilon_{\rm ls}^2<4/5.
 \end{aligned}
\]
After norm-tight normalization the only eigenvalues are $1$ and
$\mu^2/(2h)$, so $\kappa_{\rm eff}$ drops to $1$ only once
$\epsilon_{\rm ls}^2$ reaches the marked weight, which is at least $4/5$.
Consequently, under the stated sparse-value oracle, or the canonical
diagonal block encoding reducible to it, preparing the conventional exact
normalized multiplier state to trace error $1/10$ (so that measuring the
marked coordinate succeeds with probability at least $4/5-1/10=7/10$)
requires
$\Omega(\sqrt{N_{\rm md}})=\Omega(\mu^{-2})=\Omega(\kappa)$ queries.

Taking $1-\sigma=\mu^2/4=\Theta(N_{\rm md}^{-1/2})$ gives a positive feasible
short step whose endpoint remains in a fixed two-norm central neighborhood.  This
makes the occurrence native to an IPM, subject to three essential caveats.
First, it is not representation invariant: row equilibration turns
\eqref{eq:marked-diagonal-H} into the identity and makes the scaled
multiplier state easy (the primal/slack direction itself is invariant).
Second, simply omitting the marked component leaves relative normal residual
only $\Theta(\mu^2)$, so ordinary residual-based inexact IPMs can bypass the
state-preparation contract.  Third, the hard scale occurs at one point of
each fixed instance's path; it does not persist as that fixed instance's
barrier parameter tends to zero.  The lower bound is therefore for the
original exact normalized multiplier-state contract, not for every QIPM or
every equivalent representation.  This distinction also accords with the
representation and finite-precision cautions in classical IPM analysis
\cite{wright1998-ill-conditioning-computational-error-interior}.
This $\Omega(\kappa)$ statement is a full-query-model bound and survives
factor access.  It is consistent with
\cref{thm:quantum-cluster-obstruction}\textup{(iii)}, which is a
polynomial-metric statement and becomes an end-to-end factor-access solve
only under its explicit overlap hypothesis.

\subsection{Novelty and scope}

The application of the 2026 DLS solvers to IPM Newton systems, the exact
coupling law, the $D=4$ versus $D=8$ selection rule and its QIPM correctors,
and the state-preserving stable split appear to be new.  This is deliberately
not an unconditional priority claim.  Effective conditioning, active-set
preconditioning, Green-function moment thresholds, and the generic
marked-diagonal search bound are prior art.  None of the results in this
section proves an end-to-end quantum LP speedup or a new universal QLS lower
bound.  Access construction, active-set discovery where needed, precision,
iterate maintenance, and output costs remain part of any such claim.

\section{Block Log-Det Condensation}
\label{sec:condensation}

This section studies the log-det center of an SDP whose variable is a tuple
of PSD blocks with one shared trace constraint $\sum_g\tr Z_g=1$.  Suppose
one block, the winner, has a cost whose least eigenvalue is smaller than
that of every other block.  As the path parameter decreases, the trace
multiplier approaches this eigenvalue, where the winner's resolvent has a
pole.  The losing blocks can hold only a bounded amount of trace near that
eigenvalue, so the remaining trace condenses on the winner.  We show that
the winner's trace mass controls three quantities: the condition number of
the equality-reduced Hessian, the trace distance
between the centers with and without a winner (the visibility of the
winner), and the query cost of preparing the center as a quantum state or
extracting the winner from it.

We work with an abstract two-cost block model.  \Cref{sec:winner-take-all}
gives a sparse SDP that realizes it, and \cref{sec:state-conversion} proves
a stronger lower bound for mixed-state output.  In this section only, $G$ is
the number of blocks, $d$ their common dimension, $m$ the multiplicity of
the least eigenvalue of the winning cost (not the number of equality
constraints), $\tau$ the path parameter (called $\mu$ in
\cref{def:central-path}), $\alpha=\tau G$, and $p$ the trace of the winning
block (the winner mass).  The $C$- and $\Gamma$-symbols and $d,b,c$ are also
local, and the resolvent trace $A(x)$ is unrelated to the constraint matrix
$A$ of \cref{sec:prelim-conic}.
\newcommand{\condtag}[1]{\tag{#1}}

\subsection{Two-cost model and exact center}
\label{sec:condensation-model}

Fix real symmetric costs $C_*,C_L\in\Sym^d$.  There is one winning block
with cost $C_*$ and $G-1$ identical losing blocks with cost $C_L$.  After
subtracting the winning spectral minimum $v_*$ from both costs, suppose
\begin{align}
 \operatorname{spec}(C_*-v_*I)
  &=\{\underbrace{0,\ldots,0}_{m},\gamma_{m+1},\ldots,\gamma_d\},
   &\gamma_i&>0,                                             \label{eq:cond-winner-spectrum}\\
 \operatorname{spec}(C_L-v_*I)
  &=\{\delta_1,\ldots,\delta_d\},
   &\delta_j&>0.                                             \label{eq:cond-loser-spectrum}
\end{align}
Here $1\le m\le d$, and the winning block is unique.  Unless otherwise
stated, $d$, $m$, and all positive gaps are fixed while $G$ grows.

At path parameter $\tau>0$, consider
\begin{equation}
 \min_{\substack{Z_g\succ0\\ \sum_g\tr Z_g=1}}
 \left\{\inner{C_*}{Z_*}+\sum_{g=2}^G\inner{C_L}{Z_g}
             -\tau\sum_{g=1}^G\log\det Z_g\right\}.
 \label{eq:cond-center-program}
\end{equation}
For $x>0$, define the winner and loser resolvent traces
\begin{equation}
 A(x)=\frac{m}{x}+\sum_{i=m+1}^d\frac1{x+\gamma_i},
 \qquad
 B(x)=\sum_{j=1}^d\frac1{x+\delta_j}.
 \label{eq:cond-resolvents}
\end{equation}
An empty regular sum is zero.  Also put
\begin{equation}
 B_0=B(0)=\sum_{j=1}^d\delta_j^{-1},
 \qquad
 \beta=-B'(0)=\sum_{j=1}^d\delta_j^{-2},
 \qquad
 \alpha_*=B_0^{-1}.
 \label{eq:cond-capacity}
\end{equation}

\begin{proposition}[Exact block center and finite concentration bounds]
\label{prop:cond-exact-center}
The unique minimizer of \eqref{eq:cond-center-program} is
\begin{equation}
 Z_*=\tau(C_*-\lambda I)^{-1},\qquad
 Z_g=\tau(C_L-\lambda I)^{-1}\quad(g>1),
 \label{eq:cond-exact-center}
\end{equation}
where $\lambda=v_*-x$ and $x$ is the unique positive solution of
\begin{equation}
 1=\tau\{A(x)+(G-1)B(x)\}.
 \label{eq:cond-secular}
\end{equation}
The winner mass $p_G=\tr Z_*=\tau A(x)$ satisfies, with
$\delta_{\max}=\max_j\delta_j$,
\begin{align}
 p_G&\ge1-\tau(G-1)B_0,                                  \label{eq:cond-finite-lower}\\
 p_G\ge p_0>0&\quad\Longrightarrow\quad
 x\le\frac{d\tau}{p_0},\qquad
 1-p_G\ge\frac{d\tau(G-1)}{\delta_{\max}+d\tau/p_0}.
                                                               \label{eq:cond-finite-converse}
\end{align}
Consequently, for $G\ge2$ and $0<\varepsilon<1$,
\begin{align}
 \tau\le\frac{\varepsilon}{B_0(G-1)}
 &\quad\Longrightarrow\quad p_G\ge1-\varepsilon,          \label{eq:cond-epsilon-sufficient}\\
 p_G\ge1-\varepsilon
 &\quad\Longrightarrow\quad
 \tau\le\frac{\varepsilon\delta_{\max}}
 {d\{G-1-\varepsilon/(1-\varepsilon)\}},                  \label{eq:cond-epsilon-necessary}
\end{align}
whenever the denominator in \eqref{eq:cond-epsilon-necessary} is positive.
For fixed $G$ and $\tau\downarrow0$, if
$A_{\rm reg}(0)=\sum_{i=m+1}^d\gamma_i^{-1}$, then
\begin{align}
 x&=m\tau+m\{A_{\rm reg}(0)+(G-1)B_0\}\tau^2+O_G(\tau^3),
                                                               \label{eq:cond-small-tau-x}\\
 p_G&=1-(G-1)B_0\tau+m(G-1)\beta\tau^2+O_G(\tau^3).
                                                               \label{eq:cond-small-tau-p}
\end{align}
\end{proposition}

\begin{proof}
The objective is strictly convex on the affine slice and diverges at its
relative boundary, so it has a unique minimizer.  With the multiplier sign
from \cref{def:trace-logdet-center}, stationarity gives
$C_g-\lambda I=\tau Z_g^{-1}$.  Positive definiteness requires
$\lambda<v_*$, and taking traces gives \eqref{eq:cond-secular}.  Its
right-hand side is continuous and strictly decreasing from infinity to zero,
which proves existence and uniqueness of $x$.

The loser mass is $1-p_G=\tau(G-1)B(x)$.  Since $B(x)\le B_0$, this proves
\eqref{eq:cond-finite-lower}.  Since $A(x)\le d/x$, $p_G\ge p_0$ implies
$x\le d\tau/p_0$.  The bound
$B(x)\ge d/(x+\delta_{\max})$ then proves
\eqref{eq:cond-finite-converse}.  Equations
\eqref{eq:cond-epsilon-sufficient}--\eqref{eq:cond-epsilon-necessary} follow
by substituting $p_0=1-\varepsilon$ and rearranging.

Finally, Taylor expansion gives
$A(x)=m/x+A_{\rm reg}(0)+O(x)$ and
$B(x)=B_0-\beta x+O(x^2)$.  Substitute the ansatz
$x=m\tau+c\tau^2+O_G(\tau^3)$ into \eqref{eq:cond-secular}.  Subtracting
one, dividing by $\tau$, and taking the constant term gives
$c=m\{A_{\rm reg}(0)+(G-1)B_0\}$.  Substitution into
$1-p_G=\tau(G-1)B(x)$ gives \eqref{eq:cond-small-tau-p}.
\end{proof}

Thus the winner mass stays bounded away from zero only when $\tau=O(1/G)$,
and it is close to one when $\tau$ is a small enough multiple of $1/G$.
The next result resolves this scale.

\subsection{Capacity transition}
\label{sec:condensation-threshold}

With $\tau=\alpha/G$, the losers hold total trace
$\tau(G-1)B(x)<\alpha B_0$.  We call $\alpha_*=1/B_0$ the capacity and the
cases $\alpha<\alpha_*$, $\alpha=\alpha_*$, and $\alpha>\alpha_*$
subcritical, critical, and supercritical.

\begin{theorem}[Block log-det condensation]
\label{thm:condensation-threshold}
Set $\tau=\alpha/G$ for fixed $\alpha>0$.  Under
\eqref{eq:cond-winner-spectrum}--\eqref{eq:cond-loser-spectrum}, the following
limits hold as $G\to\infty$:
\begin{alignat}{3}
 0<\alpha<\alpha_*:&\quad
 &\frac{x}{\tau}&\longrightarrow\frac{m}{1-\alpha B_0},
 &\qquad p_G&\longrightarrow1-\alpha B_0;                 \label{eq:cond-subcritical}\\
 \alpha=\alpha_*:&
 &\sqrt G\,x&\longrightarrow\sqrt{m/\beta},
 &\sqrt G\,p_G&\longrightarrow\alpha_*\sqrt{m\beta};    \label{eq:cond-critical}\\
 \alpha>\alpha_*:&
 &x&\longrightarrow x_\alpha,
 &G p_G&\longrightarrow\alpha A(x_\alpha),               \label{eq:cond-supercritical}
\end{alignat}
where $x_\alpha>0$ is the unique solution of
$B(x_\alpha)=1/\alpha$.
\end{theorem}

\begin{proof}
Write the secular equation as
\begin{equation}
 1=\frac{\alpha}{G}A(x)+\alpha(1-G^{-1})B(x).
 \label{eq:cond-scaled-secular}
\end{equation}
If $\alpha<\alpha_*$, $x$ must tend to zero: along any subsequence bounded
away from zero, the first term vanishes and the second has limit at most
$\alpha B_0<1$.  Using the expansions following
\eqref{eq:cond-small-tau-p} in \eqref{eq:cond-scaled-secular} gives
\[
 1=\frac{\alpha m}{Gx}+\alpha B_0+o(1),
\]
so $Gx\to\alpha m/(1-\alpha B_0)$ and
$p_G=(\alpha/G)A(x)\to1-\alpha B_0$.

At $\alpha=\alpha_*$ the same compactness argument gives $x\to0$.
After cancelling $\alpha_*B_0=1$, the Taylor expansions give
\begin{equation}
 0=\frac{\alpha_*m}{Gx}-\alpha_*\beta x
       +O(G^{-1})+O(x^2),
 \label{eq:cond-critical-balance}
\end{equation}
where the bounded regular part of $A$ is absorbed in $O(G^{-1})$.
Monotonicity first sandwiches $x$ between constant multiples of
$G^{-1/2}$; dividing \eqref{eq:cond-critical-balance} by $G^{-1/2}$ then
yields $\sqrt Gx\to\sqrt{m/\beta}$.  Since
$A(x)=m/x+O(1)$, the mass limit in \eqref{eq:cond-critical} follows.

If $\alpha>\alpha_*$, $x$ cannot tend to zero because the loser term alone
would tend to $\alpha B_0>1$.  It also cannot diverge.  Every convergent
subsequence therefore has a positive limit, at which the winner term in
\eqref{eq:cond-scaled-secular} vanishes and $B(x)=1/\alpha$.  Strict
monotonicity of $B$ makes that limit unique.  Multiplying
$p_G=(\alpha/G)A(x)$ by $G$ proves \eqref{eq:cond-supercritical}.
\end{proof}

Below capacity, the winner's pole term $m/x$ supplies the trace the losers
cannot hold; at capacity, it balances the first-order decrease $\beta x$ of
the loser trace; above capacity, the loser equation alone determines
$x_\alpha$.  The assumption $\delta_j>0$ is essential: a zero
loser gap makes $B_0$ infinite, and there is no finite capacity.

\subsection{Equality-reduced Hessian and conditioning}
\label{sec:condensation-conditioning}

All condition numbers below are those of \cref{def:reduced-kappa}: the
block Hessian restricted to the tangent space of the trace constraint, in
Frobenius coordinates.  They say nothing about preconditioned systems or the
saddle-point KKT matrix.  At the center put
$D_*=C_*-\lambda I$ and $D_L=C_L-\lambda I$.  On the tangent space
\begin{equation}
 \mathcal T=\left\{Y=(Y_1,\ldots,Y_G):
                    \sum_{g=1}^G\tr Y_g=0\right\},
 \label{eq:cond-tangent}
\end{equation}
the barrier-scaled log-det Hessian has quadratic form
\begin{equation}
\mathcal Q(Y)=\frac1\tau\left\{
  \tr(D_*Y_*D_*Y_*)+\sum_{g=2}^G\tr(D_LY_gD_LY_g)\right\}.
 \label{eq:cond-hessian-form}
\end{equation}
This \(\mathcal Q\) is the quadratic form of
\(\tau\nabla^2F_{\log\det}\) at the center, so
\(\kappa(\mathcal Q|_{\mathcal T})=\kappa_{\rm red}\); the factor \(\tau\)
does not change curvature ratios.

Let $\delta_{\min}=\min_j\delta_j$, let $r$ be its multiplicity, and, when
$m<d$, let $\gamma_{\min}=\min_{i>m}\gamma_i$.  Define
\begin{equation}
 b(x)=\min\bigl(\{x+\delta_{\min}\}\cup
                  \{x+\gamma_{\min}:m<d\}\bigr),
 \quad
 \Lambda=\max\bigl(\{\delta_j\}\cup
                  \{\gamma_i:i>m\}\bigr),
 \label{eq:cond-gap-extremes}
\end{equation}
with the second set omitted when $m=d$.

\begin{lemma}[Finite Hessian sandwich]
\label{lem:cond-hessian-sandwich}
For $m\ge2$,
\begin{equation}
 \lambda_{\min}(\mathcal Q|_{\mathcal T})=\frac{x^2}{\tau}.
 \label{eq:cond-min-degenerate}
\end{equation}
For $m=1$, let $\lambda_{\rm diag}$ be the least curvature on the diagonal
part of $\mathcal T$ (directions diagonal in the cost eigenbases).  Then
\begin{equation}
 \frac{(dG-1)x^2+b(x)^2}{dG\tau}
 \le\lambda_{\rm diag}\le
 \frac{r(G-1)x^2+(x+\delta_{\min})^2}
      {\tau\{r(G-1)+1\}}.                                  \label{eq:cond-diagonal-sandwich}
\end{equation}
If $m=1<d$, the ground--excited winner direction, which couples the ground
eigenvector with an eigenvector of gap $\gamma_{\min}$, has curvature
$x(x+\gamma_{\min})/\tau$.  Consequently, with constants depending only on
the fixed spectral data,
\begin{equation}
 \lambda_{\min}(\mathcal Q|_{\mathcal T})\asymp\frac1\tau
 \begin{cases}
 \min\{x(x+\gamma_{\min}),
        x^2+(x+\delta_{\min})^2/G\},&m=1<d,\\
 x^2+(x+\delta_1)^2/G,&m=d=1.
 \end{cases}                                                \label{eq:cond-min-simple}
\end{equation}
For $G\ge3$,
\begin{equation}
 \frac{(x+\delta_{\max})^2}{\tau}
 \le\lambda_{\max}(\mathcal Q|_{\mathcal T})
 \le\frac{(x+\Lambda)^2}{\tau}.                           \label{eq:cond-max-sandwich}
\end{equation}
\end{lemma}

\begin{proof}
Use eigenbases of the two costs.  Diagonal and symmetric off-diagonal
coordinates are orthogonal invariant subspaces.  An off-diagonal $(i,j)$
coordinate in a block with shifted eigenvalues $d_i,d_j$ has curvature
$d_id_j/\tau$, and a diagonal coordinate has curvature $d_i^2/\tau$.
Only diagonal coordinates are coupled by \eqref{eq:cond-tangent}.

If $m\ge2$, the difference of two winner ground diagonals, or their
off-diagonal coordinate, is tangent and has curvature $x^2/\tau$.  No
ambient coordinate has smaller curvature, proving \eqref{eq:cond-min-degenerate}.
For $m=1$, write the winner's ground diagonal coefficient as $a$ and all other
$dG-1$ diagonal coefficients as $z$.  The trace constraint gives
$\mathbf1^\top z=-a$, hence $\norm{z}_2^2\ge a^2/(dG-1)$; every corresponding
curvature is at least $b(x)^2/\tau$.  Minimizing the resulting two-term quotient
gives the lower bound in \eqref{eq:cond-diagonal-sandwich}.  For the upper
bound, cancel $a$ equally among the $r(G-1)$ loser coordinates with gap
$\delta_{\min}$ and calculate its Rayleigh quotient.

The stated winner off-diagonal coordinate supplies the other upper bound in
\eqref{eq:cond-min-simple}; all remaining off-diagonal curvatures are bounded
below by the same expression up to fixed spectral constants.  Together with
the diagonal sandwich this proves \eqref{eq:cond-min-simple}, including the
scalar case where no within-block off-diagonal exists.  Finally, opposite
maximum-gap diagonal directions in two loser blocks are tangent and give the
lower bound in \eqref{eq:cond-max-sandwich}; its upper bound is the largest
curvature of an ambient coordinate.
\end{proof}

\begin{theorem}[Multiplicity-dependent conditioning phases]
\label{thm:cond-conditioning-phases}
For fixed $\alpha>0$, $\tau=\alpha/G$, and fixed spectral data,
\begin{equation}
 \kappa_{\rm red}=\begin{cases}
 \Theta(G^2),&m\ge2,\ 0<\alpha<\alpha_*,\\
 \Theta(G),&m=1,\ 0<\alpha\le\alpha_*,\\
 \Theta(G),&m\ge2,\ \alpha=\alpha_*,\\
 \Theta(1),&\alpha>\alpha_*\quad\text{for every }m.
 \end{cases}
 \label{eq:cond-kappa-phases}
\end{equation}
\end{theorem}

\begin{proof}
In the subcritical phase, \cref{thm:condensation-threshold} gives
$x=\Theta(G^{-1})$.  By \eqref{eq:cond-max-sandwich}, the largest
curvature is $\Theta(G)$.  For $m\ge2$,
\eqref{eq:cond-min-degenerate} gives minimum curvature $\Theta(G^{-1})$;
for $m=1$, \eqref{eq:cond-min-simple} gives $\Theta(1)$.  At criticality
$x=\Theta(G^{-1/2})$, and either minimum formula gives $\Theta(1)$, while
the maximum remains $\Theta(G)$.  Above capacity $x$ is bounded away from
zero, so every curvature is $\Theta(G)$.
\end{proof}

The ground--excited direction needs $m<d$.  For $m=d=1$, the diagonal
direction that spreads the winner's trace across the losers plays its role,
and \eqref{eq:cond-min-simple} gives the same orders of $\kappa_{\rm red}$.

\subsection{Finite mass--conditioning laws}
\label{sec:condensation-mass-conditioning}

The next theorem bounds $\kappa_{\rm red}$ at finite $G$ in terms of the
winner mass $p$, without fixing the schedule $\tau=\alpha/G$ of
\cref{thm:cond-conditioning-phases}.

\begin{theorem}[Finite mass--conditioning law]
\label{thm:cond-mass-conditioning}
Let $G\ge3$ and let $p=p_G\in(0,1)$ be the winner mass at an exact center.
Define
\begin{equation*}
 R_G(p)=\frac{(G-1)p}{1-p}. \condtag{MC1}\label{eq:cond-mc1}
\end{equation*}
Then
\begin{align*}
 m\ge2:&\qquad \kappa_{\rm red}\ge R_G(p)^2,              \condtag{MC2}\label{eq:cond-mc2}\\
 m=1<d:&\qquad \kappa_{\rm red}\ge
 \min\left\{1,\frac{\delta_{\max}}{\gamma_{\min}}\right\}R_G(p).
                                                               \condtag{MC3}\label{eq:cond-mc3}
\end{align*}
If $m=d=1$, writing $\delta=\delta_1$, the exact closed form is
\begin{equation*}
 \boxed{\displaystyle
 \kappa_{\rm red}=\frac{G}
 {1+(1-p)^2/\{(G-1)p^2\}}.}                                \condtag{MC4}\label{eq:cond-mc4}
\end{equation*}
For scalar blocks with $G=2$, the tangent space is one-dimensional and
$\kappa_{\rm red}=1$, so \eqref{eq:cond-mc4} is asserted only for $G\ge3$.
\end{theorem}

\begin{proof}
The exact mass equations are
\begin{equation}
 p=\tau A(x),\qquad 1-p=\tau(G-1)B(x).                    \label{eq:cond-mass-equations}
\end{equation}
Because $A(x)\le d/x$ and $B(x)\ge d/(x+\delta_{\max})$,
\begin{equation}
 \frac{1-p}{p}=(G-1)\frac{B(x)}{A(x)}
 \ge(G-1)\frac{x}{x+\delta_{\max}},
 \qquad \frac{x+\delta_{\max}}x\ge R_G(p).              \label{eq:cond-mass-ratio}
\end{equation}
The tangent direction with opposite unit entries on the $\delta_{\max}$
diagonal coordinates of two loser blocks gives
$\lambda_{\max}\ge(x+\delta_{\max})^2/\tau$.  If $m\ge2$, divide this
by \eqref{eq:cond-min-degenerate} and use
\eqref{eq:cond-mass-ratio}.

If $m=1<d$, the winner ground--excited direction gives
$\lambda_{\min}\le x(x+\gamma_{\min})/\tau$, while
\[
 \frac{x+\delta_{\max}}{x+\gamma_{\min}}
 \ge\min\left\{1,\frac{\delta_{\max}}{\gamma_{\min}}\right\}.
\]
Factoring the curvature ratio proves \eqref{eq:cond-mc3}.

For $d=m=1$, put $a=x^2/\tau$ and $b=(x+\delta)^2/\tau$.  On the scalar
trace-zero space, $G-2$ loser-difference directions have curvature $b$, and
the direction balancing the winner against equal loser changes has curvature
$\varrho=\{(G-1)a+b\}/G$.  Since $b>a$, these are the largest and smallest
eigenvalues.  The mass equations reduce to
\[
 \frac{x}{x+\delta}=\frac{1-p}{(G-1)p}.
\]
Thus $\kappa_{\rm red}=b/\varrho=G/\{1+(G-1)a/b\}$, which is
\eqref{eq:cond-mc4}.
\end{proof}

In particular, $p\ge p_0>0$ forces $\kappa_{\rm red}=\Omega(G^2)$ when
$m\ge2$ and $\kappa_{\rm red}=\Omega(G)$ when $m=1$.  This also follows from
$x\le d\tau/p_0$, the bound $\tau=O(G^{-1})$ from the loser mass, and the
test directions above; the exact laws add the explicit dependence on $p$.

\subsection{Visibility, collision, and the conditioning dichotomy}
\label{sec:condensation-visibility}

The next result uses only the symmetry of the two-cost model, not the
spectral assumptions of \cref{sec:condensation-model}.  Let
$C_W,C_L\in\Sym^d$, $G\ge2$, and $\tau>0$.  We compare the trace-one center
$\rho_0$ where every block has the loser cost with the center $\rho_1$ where
block $1$ has cost $C_W$:
\begin{align}
 \rho_0&=\arg\min_{\substack{Z_g\succ0\\\sum_g\tr Z_g=1}}
 \left\{\sum_{g=1}^G\inner{C_L}{Z_g}-\tau\sum_g\log\det Z_g\right\},
                                                               \label{eq:cond-rho-zero}\\
 \rho_1&=\arg\min_{\substack{Z_g\succ0\\\sum_g\tr Z_g=1}}
 \left\{\inner{C_W}{Z_1}+\sum_{g=2}^G\inner{C_L}{Z_g}
       -\tau\sum_g\log\det Z_g\right\}.                 \label{eq:cond-rho-one}
\end{align}
Identify each tuple with its block-diagonal density matrix and write
\begin{equation}
 \rho_0=\bigoplus_{g=1}^G L_0,\qquad
 \rho_1=W_1\oplus\bigoplus_{g=2}^G L_1,\qquad p=\tr W_1.
 \label{eq:cond-rho-blocks}
\end{equation}
Permutation symmetry gives $\tr L_0=1/G$ and
$\tr L_1=(1-p)/(G-1)$.

\begin{theorem}[Visibility and label collision]
\label{thm:cond-visibility}
For the centers above,
\begin{equation*}
 \boxed{\left|p-\frac1G\right|
 \le\Dtr(\rho_0,\rho_1)\le\max\left\{p,\frac1G\right\}.}
 \condtag{O1}\label{eq:cond-visibility}
\end{equation*}
If the block labels of two independent copies are measured and compared,
the probabilities $c_0,c_1$ that they agree are
\begin{equation*}
 c_0=\frac1G,\qquad
 c_1=p^2+\frac{(1-p)^2}{G-1},\qquad
 \boxed{c_1-c_0=\frac{(Gp-1)^2}{G(G-1)}.}
 \condtag{O2}\label{eq:cond-collision}
\end{equation*}
The measurement is public (independent of the input) and acts only on the
two label registers, so the identity holds regardless of degeneracy inside
the blocks and is unchanged by purifying and then discarding auxiliary
registers.
\end{theorem}

\begin{proof}
Measuring whether the label is $1$ gives the lower bound in
\eqref{eq:cond-visibility}.  Let $\lambda_0,\lambda_1$ be the scalar trace
multipliers of the two centers.  Their loser blocks satisfy
\[
 L_b=\tau(C_L-\lambda_bI)^{-1}\quad(b=0,1),
\]
and are therefore Loewner ordered.  If $\lambda_1\le\lambda_0$, then
$L_1\preceq L_0$ and the trace identities imply $p\ge1/G$.  Hence
\[
 (G-1)\norm{L_1-L_0}_1=p-\frac1G,
 \qquad \norm{W_1-L_0}_1\le p+\frac1G.
\]
Adding the trace norms of the direct-sum blocks and dividing by two gives
$\Dtr(\rho_0,\rho_1)\le p$.  If $\lambda_1\ge\lambda_0$, the order and
inequality reverse, and the same calculation gives the upper bound $1/G$.
Finally the two label distributions are
\[
 (1/G,\ldots,1/G),\qquad
 \left(p,\frac{1-p}{G-1},\ldots,\frac{1-p}{G-1}\right).
\]
Summing their squared entries and subtracting proves
\eqref{eq:cond-collision}.
\end{proof}

Only the upper trace-distance bound uses the resolvent form of the centers;
the lower bound and the collision identity need only a uniform label
distribution for $\rho_0$ and exchangeable losers.  The condition $p\ge1/G$
holds whenever $C_W\preceq C_L$, since $X\mapsto(X-\lambda I)^{-1}$ is
operator antitone on $X\succ\lambda I$, and for the family of
\cref{sec:wta-central-path} by \eqref{eq:wta-resolvent-order}.  Then the
trace-distance bounds reduce to $p-1/G\le\Dtr\le p$.

For $C_W=C_*$, under the spectral assumptions of
\cref{sec:condensation-model} and $G\ge3$,
\cref{thm:cond-visibility,thm:cond-mass-conditioning} give the following
finite laws relating conditioning and visibility.  With
$c=\min\{1,\delta_{\max}/\gamma_{\min}\}$,
\begin{align*}
 p&\le\frac{\sqrt{\kappa_{\rm red}}}
 {G-1+\sqrt{\kappa_{\rm red}}} &&(m\ge2),                 \condtag{CV1}\label{eq:cond-cv1}\\
 p&\le\frac{\kappa_{\rm red}}
 {c(G-1)+\kappa_{\rm red}} &&(m=1<d),                    \condtag{CV2}\label{eq:cond-cv2}\\
 p&=\frac1{1+\sqrt{G-1}\sqrt{G/\kappa_{\rm red}-1}}
 &&(m=d=1).                                                \condtag{CV3}\label{eq:cond-cv3}
\end{align*}
By \eqref{eq:cond-visibility}, the maximum of each right-hand side and $1/G$
bounds $\Dtr(\rho_0,\rho_1)$.  Conversely, if
$\Dtr(\rho_0,\rho_1)\ge\upsilon>1/G$, then $p\ge\upsilon$ and
\begin{equation*}
 \kappa_{\rm red}\ge
 \left(\frac{(G-1)\upsilon}{1-\upsilon}\right)^2\ (m\ge2),
 \qquad
 \kappa_{\rm red}\ge
 c\frac{(G-1)\upsilon}{1-\upsilon}\ (m=1<d).             \condtag{CV4}\label{eq:cond-cv4}
\end{equation*}
For $m=d=1$, \eqref{eq:cond-cv3} gives the same linear order.  Thus
constant visibility forces $\kappa_{\rm red}=\Omega(G^2)$ when $m\ge2$ and
$\kappa_{\rm red}=\Omega(G)$ when $m=1$.  A constant collision gap also
forces $p=\Omega(1)$ and hence the same dichotomy.

\paragraph{Approximate outputs and reduced symmetry.}
Let $\widetilde\rho_b$ satisfy
$\Dtr(\widetilde\rho_b,\rho_b)\le\epsilon$.  Contractivity and the
tensor-product triangle inequality give
$\Dtr(\widetilde\rho_b^{\otimes2},\rho_b^{\otimes2})\le2\epsilon$.
Thus the observed two-copy collision gap is at least
\begin{equation}
 \Delta_{\rm out}\ge d(p,G)-4\epsilon,
 \qquad d(p,G):=\frac{(Gp-1)^2}{G(G-1)}.                  \label{eq:cond-output-gap}
\end{equation}
By standard concentration bounds, $O(\Delta_{\rm out}^{-2})$ repetitions
distinguish the two cases with constant success probability, and an extra
$\log(1/\delta)$ factor gives failure probability $\delta$.

The collision lower bound does not need equal loser masses.  Suppose a
public label measurement is uniform on the no-winner input and, on every
unique-winner input, gives label masses $w_g$ with winner mass
$q\ge p\ge1/G$.  Cauchy--Schwarz gives
\begin{equation}
 \sum_{g=1}^G w_g^2\ge q^2+\frac{(1-q)^2}{G-1}
 \ge p^2+\frac{(1-p)^2}{G-1},                            \label{eq:cond-symmetry-light}
\end{equation}
where the last inequality holds because
$q\mapsto q^2+(1-q)^2/(G-1)$ is increasing for $q\ge1/G$.  The collision gap is therefore at least $d(p,G)$, with
equality exactly at winner mass $p$ and uniform loser mass.  If the
no-winner collision probability is only bounded above by $b$, replace the
gap by $p^2+(1-p)^2/(G-1)-b$.  In particular, if $p\ge p_0>0$ and
$G\ge2/p_0$, then $d(p,G)\ge p_0^2/4$.

\subsection{Unique-OR query lower bounds}
\label{sec:condensation-query-lower}

To turn \cref{thm:cond-visibility} into a query lower bound, we use a
composed problem.  Let $G\ge2$ and let a possibly partial Boolean predicate
$f:\mathcal X_0\sqcup\mathcal X_1\to\{0,1\}$ have general adversary value
$\mathsf{Adv}_f=\operatorname{Adv}^{\pm}(f)$.  For $G$ independently queried
blocks, set
\begin{equation}
 \mathcal P_0=\mathcal X_0^G,
 \qquad
 \mathcal P_1=\bigsqcup_{a=1}^G
 \mathcal X_0^{a-1}\times\mathcal X_1\times\mathcal X_0^{G-a},
 \label{eq:cond-unique-or-promise}
\end{equation}
and let $F$ be the problem of distinguishing $\mathcal P_0$ (no winner) from
$\mathcal P_1$ (exactly one winner).

\begin{theorem}[Unique-OR adversary and preparation]
\label{thm:cond-unique-or}
In the raw-query model where one query addresses one coordinate of one block,
\begin{equation*}
 \operatorname{Adv}^{\pm}(F)\ge \mathsf{Adv}_f\sqrt G.     \condtag{O3}\label{eq:cond-unique-or-adv}
\end{equation*}
Consequently, suppose that fixed public projectors onto the blocks give
the uniform label distribution on every no-winner target state and winner
mass at least $p\ge1/G$ on every unique-winner target state; loser masses
may be arbitrary.  Any circuit that prepares, on every promised input, a state
within trace distance $\epsilon$ of such a target uses
$\Omega(\mathsf{Adv}_f\sqrt G)$ raw queries whenever
$d(p,G)-4\epsilon=\Omega(1)$.  By tightness of the general adversary bound,
this is equivalently $\Omega(Q(f)\sqrt G)$ up to universal constants.
\end{theorem}

\begin{proof}
Take the rectangular zero-to-one block $B$, with rows indexed by
$\mathcal X_0$ and columns by $\mathcal X_1$, of a feasible adversary matrix
for $f$, normalized so that
\begin{equation}
 \norm{B}=\mathsf{Adv}_f,\qquad
 \max_j\norm{B\circ\Delta_j}\le1,                         \label{eq:cond-inner-adversary}
\end{equation}
where $\Delta_j[x,y]=1$ exactly when query $j$ has different answers on
$x,y$, and let $I_0$ be the identity on $\ell_2(\mathcal X_0)$.  For the
sector of $\mathcal P_1$ whose winner is block $a$, define
\[
 C_a=I_0^{\otimes(a-1)}\otimes B\otimes I_0^{\otimes(G-a)},
 \qquad C=[C_1\ C_2\ \cdots\ C_G],
\]
and use the Hermitian adversary matrix
\[
 \Gamma=\begin{pmatrix}0&C\\C^*&0\end{pmatrix}.
\]
The winner sectors are disjoint, and
\[
 CC^*=\sum_{a=1}^G I_0^{\otimes(a-1)}\otimes BB^*\otimes
                          I_0^{\otimes(G-a)}.
\]
The summands commute.  Tensoring a top eigenvector of $BB^*$ across all
$G$ blocks gives their common top eigenvector, so
$\norm{\Gamma}=\mathsf{Adv}_f\sqrt G$.  Filtering by raw query $(a,j)$ kills every
winner sector except $a$, since row and column inputs agree in block $a$ in
all other sectors.  The remaining block is $B\circ\Delta_j$ and has norm at
most one by \eqref{eq:cond-inner-adversary}.  Thus $\Gamma$ is feasible and
proves \eqref{eq:cond-unique-or-adv}.

If $d(p,G)-4\epsilon$ is a positive constant,
\eqref{eq:cond-output-gap} turns a constant number of fresh preparations and
public collision measurements into a bounded-error algorithm for $F$.
The adversary lower bound therefore applies to the preparation circuit.
Tightness and composition of the negative-weight adversary bound give the
equivalent $Q(f)$ statement
\cite{hyer2005-tight-adversary-bounds-for-composite,lee2011-quantum-query-complexity-of-state,reichardt2011-reflections-for-quantum-query-algorithms}.
\end{proof}

The theorem assumes that all product inputs in
\eqref{eq:cond-unique-or-promise} are valid, both fibers of $f$ are
nonempty, and the block projectors are public.  It needs no clean evaluator, candidate
verification, inverse preparation, or restriction on the internal states of
the blocks.  If one coefficient-oracle call can be
simulated with at most $s$ raw queries, the theorem gives a lower bound of
\begin{equation}
 \Omega\!\left(\frac{\mathsf{Adv}_f\sqrt G}{s}\right)     \label{eq:cond-coefficient-transfer}
\end{equation}
coefficient-oracle calls.  An oracle that returns a winner bit, an
aggregate flag, or an input-dependent projector is stronger and may remove
the inner factor $\mathsf{Adv}_f$.

No such lower bound holds when the output error is large compared with the
visibility: by \cref{thm:cond-visibility}, outputting the public all-loser
center $\rho_0$ uses no queries and is valid on every input with zero or
one winner when
\begin{equation}
 \epsilon\ge\max\{p,1/G\}.                                \label{eq:cond-zero-query-boundary}
\end{equation}
Hence a constant primal objective gap \eqref{eq:objective-gap-g} need not imply a lower bound for
trace-normalized primal-density output at constant error.  This concerns
the direct-sum density of \cref{def:trace-density-output}, not a normalized
vectorization, whose label probabilities are Frobenius rather than trace
weights.

\subsection{Preparation, extraction, and the conservation law}
\label{sec:condensation-algorithms}

For the upper bounds we assume exactly one winner $w$ and a clean coherent
evaluator of the predicate
\begin{equation*}
 V_f\ket{g,0}=\ket{g,f_g},\qquad f_g=1\Longleftrightarrow g=w,          \condtag{A1}\label{eq:cond-clean-evaluator}
\end{equation*}
that costs $T_{\rm eval}$ raw queries, as does its inverse.  The two cost
matrices, $G$, and $\tau$ are public, so the center has the form
\begin{equation*}
 \rho_p=p\ket w\!\bra w\otimes\rho_W
 +\frac{1-p}{G-1}\sum_{g\ne w}\ket g\!\bra g\otimes\rho_L,            \condtag{A2}\label{eq:cond-central-state}
\end{equation*}
where the normalized block resolvents $\rho_W,\rho_L$ are public.  Put
\begin{equation*}
 a_W=p,\qquad a_L=\frac{1-p}{G-1},\qquad
 a_{\max}=\max\{a_W,a_L\}.                                            \condtag{A3}\label{eq:cond-component-weights}
\end{equation*}

\begin{proposition}[Coherent rejection preparation]
\label{prop:cond-rejection-preparation}
There is a purification circuit for \eqref{eq:cond-central-state} using
\begin{equation*}
 \Otilde\!\left(T_{\rm eval}\sqrt{Ga_{\max}}
              +T_{\rm loc}\right)                                    \condtag{A4}\label{eq:cond-preparation-upper}
\end{equation*}
operations, where $T_{\rm loc}$ is the gate cost (not a raw-query cost) of
the one public preparation of $\rho_W$ or $\rho_L$ conditioned on the
predicate.  When $p\ge1/G$, the raw-query part is
\begin{equation*}
 \Otilde(T_{\rm eval}\sqrt{Gp}).                                     \condtag{A5}\label{eq:cond-preparation-upper-p}
\end{equation*}
\end{proposition}

\begin{proof}
Prepare $G^{-1/2}\sum_g\ket g$, compute $f_g$, and rotate an acceptance
qubit with amplitude $\sqrt{a_{f_g}/a_{\max}}$, but do not yet copy the
label or prepare the internal state.  This trial succeeds with probability
\begin{equation*}
 P_{\rm acc}=\frac1G\left(\frac{a_W}{a_{\max}}+
                    (G-1)\frac{a_L}{a_{\max}}\right)
             =\frac1{Ga_{\max}}.                                      \condtag{A6}\label{eq:cond-acceptance}
\end{equation*}
This probability is known because the secular equation is solved
classically, so exact amplitude amplification with matched phases applies
to this trial, acting on the label, predicate, and acceptance registers
\cite{brassard2002-quantum-amplitude-amplification-and-estimation}.  After
it succeeds, the label has probability $a_{f_g}$; only then coherently copy
the label to an environment register and use the retained predicate once to
prepare a purification of $\rho_W$ or $\rho_L$.  Tracing out the copied
label, predicate, acceptance, and local purification registers gives
exactly \eqref{eq:cond-central-state}.  Only the evaluator calls are
repeated by the amplification, and the local preparation runs once, proving
both bounds; the tilde hides only polylogarithmic factors in the requested
implementation accuracy.
\end{proof}

The secular equation and the fixed-size resolvents are computed
classically; if $d$ varies, one must add $\poly(d,\log(1/\epsilon))$ costs
for arithmetic and state synthesis.  A bounded-error evaluator must first
have its error reduced coherently below the final state error.

Suppose independent copies, each costing $Q_\rho$ raw queries, satisfy
$\Dtr(\widetilde\rho,\rho_p)\le\epsilon<p$.  Measuring the label proposes
the winner with probability at least $p-\epsilon$, so verifying and
repeating finds the winner with failure probability $\delta$ using
\begin{equation*}
 O\!\left(\frac{Q_\rho+T_{\rm eval}}{p-\epsilon}
          \log\frac1\delta\right)                                     \condtag{A7}\label{eq:cond-copy-extraction}
\end{equation*}
raw queries.  For exact copies take $\epsilon=0$; if
$\epsilon\le p/2$, this has the usual $O(1/p)$ repetition factor.  Given
instead a coherent preparation unitary $U_\rho$ costing $Q_U$ raw queries,
its inverse, and a clean winner reflection, amplitude amplification gives
\begin{equation*}
 O\!\left(\frac{Q_U+T_{\rm eval}}{\sqrt p}\log\frac1\delta\right).     \condtag{A8}\label{eq:cond-coherent-extraction}
\end{equation*}
This bound needs $U_\rho$ and its inverse as operators; an accurate output
density matrix says nothing about how the unitary acts on inputs other than
$\ket0$.

Combining \eqref{eq:cond-preparation-upper-p} and
\eqref{eq:cond-coherent-extraction} gives the preparation--extraction
conservation law
\begin{equation*}
 \Otilde(T_{\rm eval}\sqrt{Gp})\,O(p^{-1/2})
   =\Otilde(T_{\rm eval}\sqrt G).                                     \condtag{A9}\label{eq:cond-conservation}
\end{equation*}
Stronger condensation moves work from extraction to preparation without
changing the end-to-end number of raw queries up to logarithmic factors.
The three regimes give the following costs for exact output
($\epsilon\le p/2$ suffices in the copy-extraction column):
\begin{equation*}
\begin{array}{@{}lccc@{}}
 \text{regime} & \text{central preparation} & \text{copy extraction}
   & \text{coherent extraction}\\ \hline
 \text{subcritical} & \Otilde(T_{\rm eval}G^{1/2}) & O(1) & O(1)\\
 \text{critical} & \Otilde(T_{\rm eval}G^{1/4}) & O(G^{1/2}) & O(G^{1/4})\\
 \text{supercritical} & \Otilde(T_{\rm eval}) & O(G) & O(G^{1/2})
\end{array}                                                            \condtag{A10}\label{eq:cond-algorithm-phases}
\end{equation*}
The supercritical preparation bound matters only for exact output or for
accuracy fine enough to resolve a winner mass of order $1/G$; at fixed
error, the zero-query output of \eqref{eq:cond-zero-query-boundary} may
suffice.

Alternatively, quantum search finds the winner at cost
\begin{equation*}
 \Otilde(T_{\rm eval}\sqrt G),                                      \condtag{A11}\label{eq:cond-direct-search}
\end{equation*}
and solving and recovering the selected block at costs $T_{\rm solve}$ and
$T_{\rm recover}$ gives the end-to-end bound
\begin{equation*}
 \Otilde(T_{\rm eval}\sqrt G)+T_{\rm solve}+T_{\rm recover}.          \condtag{A12}\label{eq:cond-search-crossover}
\end{equation*}
These bounds use standard quantum search and minimum finding
\cite{boyer1998-tight-bounds-on-quantum-searching,durr1996-a-quantum-algorithm-for-finding}.
This is search followed by crossover, not a new interior-point method.  If
the local predicate has bounded-error raw-query complexity $q$ and the
blocks are independent, composition of the general adversary bound
\cite{hyer2005-tight-adversary-bounds-for-composite,lee2011-quantum-query-complexity-of-state,reichardt2011-reflections-for-quantum-query-algorithms}
gives the total bounded-error raw-query cost
\begin{equation*}
 Q(\operatorname{UniqueOR}_G\circ f)=\Theta(q\sqrt G).                \condtag{A13}\label{eq:cond-search-composition}
\end{equation*}
Combining this lower bound with the coherent extraction bound
\eqref{eq:cond-coherent-extraction} gives
\begin{equation*}
 Q_U+T_{\rm eval}=\Omega(q\sqrt{Gp}),                                  \condtag{A14}\label{eq:cond-prep-tradeoff}
\end{equation*}
and, for $\epsilon$-accurate independent copies with $\epsilon<p$, the copy
extraction bound gives the weaker
\begin{equation*}
 Q_\rho+T_{\rm eval}=\Omega((p-\epsilon)q\sqrt G).                    \condtag{A15}\label{eq:cond-copy-tradeoff}
\end{equation*}
For exact output this is $\Omega(pq\sqrt G)$, and the same asymptotic bound
holds when $\epsilon\le p/2$.
These statements require independent block oracles; if one query can touch
several blocks, or an aggregate winner flag is available, the composition
argument fails.

For comparison, the block product barrier has parameter $dG$, so short-step
path following from constant $\tau$ to $\tau=\Theta(1/G)$ takes
$O(\sqrt{dG}\log G)$ Newton iterations.  This count excludes the linear
solves, whose cost depends on a condition number, and state preparation and
recovery.  It must not be multiplied by \eqref{eq:cond-kappa-phases} unless
an actual solver and preconditioner are chosen and analyzed.

\subsection{Sublinear multiwinner capacity}
\label{sec:condensation-multiwinner}

Now let exactly $k$ blocks have cost $C_W$, with
$v=\lambda_{\min}(C_W)$ of multiplicity $m$, and suppose
$C_L-vI\succ0$.  In this subsection, redefine
\begin{equation}
 A(x)=\tr(C_W-vI+xI)^{-1},\qquad
 B(x)=\tr(C_L-vI+xI)^{-1}.                              \label{eq:cond-multi-resolvents}
\end{equation}
Then $A(x)=m/x+O(1)$,
$B_0=B(0)<\infty$, and $\beta=-B'(0)>0$.  The exact equations are
\begin{equation}
 1=\tau\{kA(x)+(G-k)B(x)\},\qquad
 p_{G,k}=\tau kA(x),\quad1-p_{G,k}=\tau(G-k)B(x).         \label{eq:cond-multi-secular}
\end{equation}

\begin{theorem}[Multiwinner capacity law]
\label{thm:cond-multiwinner}
Let $1\le k=k_G=o(G)$, put $\tau=\alpha/G$, and keep $\alpha>0$ fixed.
With $\alpha_*=1/B_0$,
\begin{align}
 p_{G,k}&\longrightarrow1-\alpha B_0,
 &x&\sim\frac{\alpha m k}{G(1-\alpha B_0)}
 &&(0<\alpha<\alpha_*),                                  \label{eq:cond-multi-subcritical}\\
 p_{G,k}&\sim\alpha_*\sqrt{m\beta k/G},
 &x&\sim\sqrt{mk/(\beta G)}
 &&(\alpha=\alpha_*),                                    \label{eq:cond-multi-critical}\\
 p_{G,k}&\sim\alpha A(x_\alpha)k/G,
 &B(x_\alpha)&=1/\alpha
 &&(\alpha>\alpha_*).                                    \label{eq:cond-multi-supercritical}
\end{align}
\end{theorem}

\begin{proof}
Substitute $\tau=\alpha/G$ into the first equation in
\eqref{eq:cond-multi-secular}.  Below criticality $x\to0$: otherwise the
winner term vanishes because $k/G\to0$, while the loser term is at most
$\alpha B_0<1$.  The loser mass then tends to $\alpha B_0$, proving the
first mass law, and the pole $A(x)\sim m/x$ gives the formula for $x$.

At criticality put $q=k/G$.  Taylor expansion of the loser equation and the
winner pole give two expressions for the total winner mass:
\begin{equation}
 p_{G,k}=q+\frac{\beta}{B_0}x+o(q+x)
       =\frac{qm}{B_0x}+O(q).                            \label{eq:cond-multi-balance}
\end{equation}
As in the one-winner proof, monotonicity first bounds $x$ by constant
multiples of $\sqrt q$; balancing the leading terms then gives
$x\sim\sqrt{mq/\beta}$ and the critical mass.  Above criticality the same
compactness and strict-monotonicity argument as in
\cref{thm:condensation-threshold} gives $x\to x_\alpha>0$; the winner
equation gives the final law.
\end{proof}

The assumption $k=o(G)$ is needed.  If $k/G\to q\in(0,1]$, the limit
instead solves $1=\alpha\{qA(x)+(1-q)B(x)\}$ with $x>0$, and there is no
singular capacity transition at fixed $\alpha$.

\paragraph{Nonidentical losers.}
Let every losing cost satisfy $C_{\ell,G}-vI\succ0$, and define
\[
 B_{\ell,G}(x)=\tr(C_{\ell,G}-vI+xI)^{-1}
\]
and their average
\[
 \overline B_G(x)=\frac1{G-k}
       \sum_{\ell\ {\rm losing}}B_{\ell,G}(x).
\]
If $k=o(G)$ and $\overline B_G\to B$ locally uniformly on $x\ge0$, with
$B(0)=B_0\in(0,\infty)$, the subcritical and supercritical conclusions and
the subcritical formula for $x$ remain valid.  The critical law also holds
if, in addition,
\begin{equation}
 \overline B_G(0)-B_0=o(\sqrt{k/G}),\qquad
 \overline B_G(x)=\overline B_G(0)-\beta x+r_G(x),         \label{eq:cond-empirical-resolvent}
\end{equation}
where $\beta>0$ and, for every fixed $C>0$,
\[
 \sup_{0<x\le C\sqrt{k/G}}\frac{|r_G(x)|}{x}\longrightarrow0.
\]
Indeed, replacing $B$ by $\overline B_G$ in
\eqref{eq:cond-multi-secular} leaves the subcritical and supercritical limit
arguments unchanged.  Under \eqref{eq:cond-empirical-resolvent}, the
finite-size error is of lower order than both leading terms in
\eqref{eq:cond-multi-balance}, which proves the critical law.  These
hypotheses control only the winner mass; they do not give exchangeable
losers, a common public internal state, or the exact collision identity.

For identical blocks and a fixed winner set $S$ of size $k<G$, let
$q=k/G$.  The Loewner-order argument in the proof of
\cref{thm:cond-visibility} gives
\begin{equation}
 |p_{G,k}-q|\le\Dtr(\rho_0,\rho_S)\le\max\{p_{G,k},q\},     \label{eq:cond-multi-visibility}
\end{equation}
and direct calculation gives the two-copy label-collision gap
\begin{equation}
 c_S-c_0=\frac{p_{G,k}^2}{k}+\frac{(1-p_{G,k})^2}{G-k}-\frac1G
 =\boxed{\frac{(Gp_{G,k}-k)^2}{Gk(G-k)}}.                 \label{eq:cond-multi-collision}
\end{equation}
Thus a constant total winner mass gives a collision gap of only $O(1/k)$
when $k$ grows, and a trace-distance bound for each fixed $S$ does not give
one measurement for an unknown $S$.  \Cref{sec:state-conversion}
avoids this $1/k$ loss for its structured family by the composed
small-success bound (\cref{thm:sc-small-success}); that stronger conclusion
is not available in the abstract two-cost model.

\subsection{Exactly-\texorpdfstring{$k$}{k} coherent preparation}
\label{sec:condensation-exact-k-preparation}

Assume now that $1\le k<G$ is known, that there are exactly $k$ winners, and
that a clean coherent membership marker costs $T_{\rm mark}$ raw queries.
The two costs, $G$, $k$, $\tau$, the trace multiplier, and the normalized
internal resolvents are public, and controlled marker calls and the inverse
preparation circuit are available.  Put, for $p=p_{G,k}$,
\begin{equation*}
 a_W=\frac pk,\qquad a_L=\frac{1-p}{G-k},\qquad
 a_{\max}=\max\{a_W,a_L\}.                               \condtag{AP0}\label{eq:cond-ap0}
\end{equation*}
The construction in \cref{prop:cond-rejection-preparation}, with the
membership bit in place of the unique-winner bit, succeeds in one trial with
probability $1/(Ga_{\max})$.  This probability is known exactly, so exact
amplitude amplification with matched phases uses
\begin{equation*}
 O(\sqrt{Ga_{\max}}+1)                                    \condtag{AP1}\label{eq:cond-ap1}
\end{equation*}
marker calls
\cite{brassard2002-quantum-amplitude-amplification-and-estimation}.  After
it succeeds, we copy the label to an environment register and, controlled
on the membership bit, prepare once a public purification of the normalized
winner or loser block, where $W$ and $L$ are the common unnormalized winner
and loser blocks of the center and
\begin{equation*}
 \omega_W=\frac{W}{p/k},\qquad
 \omega_L=\frac{L}{(1-p)/(G-k)}.                                     \condtag{AP2}\label{eq:cond-ap2}
\end{equation*}
The resulting circuit satisfies
\begin{equation*}
 \boxed{Q_U=O\!\left(T_{\rm mark}[\sqrt{Ga_{\max}}+1]\right).}         \condtag{AP3}\label{eq:cond-ap3}
\end{equation*}
When $p\ge k/G$, this is $O(T_{\rm mark}[\sqrt{Gp/k}+1])$; when
$p<k/G$, it is $O(T_{\rm mark}[\sqrt{G(1-p)/(G-k)}+1])$.  Preparing to
error $\epsilon>0$ rather than exactly adds the standard factor logarithmic
in $1/\epsilon$.  Costs of loading nonpublic eigenvectors, resolvents, or
multipliers must be added.

Conversely, amplification from the prepared purification finds a winner in
\begin{equation*}
 O\!\left(\frac{Q_U+T_{\rm mark}}{\sqrt p}\right)          \condtag{AP4}\label{eq:cond-ap4}
\end{equation*}
raw queries.  If the local predicate has bounded-error raw-query complexity
$q_f$ and independent composition gives the lower bound
$\Omega(q_f\sqrt{G/k})$ for finding a winner, then \eqref{eq:cond-ap4}
implies
\begin{equation*}
 \boxed{Q_U+T_{\rm mark}=\Omega\!\left(q_f\sqrt{Gp/k}\right).}         \condtag{AP5}\label{eq:cond-ap5}
\end{equation*}
This requires the preparation unitary and its inverse as operators; for
mixed-state output, independent copies give only the weaker tradeoff from
repetition.

If $f$ is the parity of $N$ input bits, then $T_{\rm mark}=\Theta(N)$ and
adversary composition gives winner-finding complexity
$\Theta(N\sqrt{G/k})$.  For $k=o(G)$, \cref{thm:cond-multiwinner} gives the
parity profile
\begin{equation*}
 Q_U=\begin{cases}
 \Theta(N\sqrt{G/k}),&0<\alpha<\alpha_*,\\
 \Theta(N(G/k)^{1/4}),&\alpha=\alpha_*,\\
 O(N),&\alpha>\alpha_*.
 \end{cases}                                                \condtag{AP6}\label{eq:cond-ap6}
\end{equation*}
The first two lines match \eqref{eq:cond-ap5} when the marker is optimal and
the amplification factor diverges.  The supercritical upper bound follows
from $a_{\max}=\Theta(1/G)$ even if $p<k/G$.  At criticality $p\to0$, so
this operator-level lower bound need not hold for density output at fixed
error.

This circuit does not solve the problem under the promise of zero or $k$
winners: with no winner, the center has a different trace multiplier and
loser resolvent.  Under that promise, exact existence search followed by
conditional preparation costs $O(T_{\rm mark}\sqrt{G/k})$, while returning
the public all-loser center costs zero queries and has error at most
$\max\{p_{G,k},k/G\}$ by \eqref{eq:cond-multi-visibility}.  The two-copy
collision measurement gives a matching lower bound only when the gap
\eqref{eq:cond-multi-collision}, reduced by the output error, stays
constant.

\subsection{Robust approximate centrality}
\label{sec:condensation-robustness}

The mass and conditioning bounds also hold, with modified constants, at
approximately central points whose blocks need not commute with the costs.
Keep the one-winner notation and, for positive definite $\widetilde Z_g$
and $\widetilde\lambda<v_*$, put
\begin{equation}
 x=v_*-\widetilde\lambda,\quad
 D_*=C_*-\widetilde\lambda I,\quad D_L=C_L-\widetilde\lambda I,\quad
 \widetilde T=\sum_g\tr\widetilde Z_g,\quad
 \widetilde p=\frac{\tr\widetilde Z_*}{\widetilde T}.       \label{eq:cond-robust-defs}
\end{equation}
Write $D_g=D_*$ for the winner and $D_g=D_L$ for every loser.
Assume a dimensionless stationarity residual in the barrier norm and a
trace residual:
\begin{equation}
 E_g=\tau^{-1}\widetilde Z_g^{1/2}D_g\widetilde Z_g^{1/2}-I,\qquad
 \max_g\norm{E_g}_{\rm op}\le\eta<1,\qquad
 |\widetilde T-1|\le\zeta<1.                              \label{eq:cond-robust-residual}
\end{equation}

\begin{theorem}[Robust capacity, mass, and conditioning]
\label{thm:cond-robust}
Every point satisfying \eqref{eq:cond-robust-defs}--
\eqref{eq:cond-robust-residual} satisfies the noncommutative Loewner
sandwich
\begin{equation*}
 (1-\eta)\tau D_g^{-1}\preceq\widetilde Z_g
 \preceq(1+\eta)\tau D_g^{-1}.                            \condtag{R1}\label{eq:cond-r3}
\end{equation*}
Consequently,
\begin{equation*}
 \boxed{\widetilde p\ge
 1-\frac{(1+\eta)\tau(G-1)B_0}{1-\zeta}
 \ge1-\frac{(1+\eta)\alpha B_0}{1-\zeta}.}               \condtag{R2}\label{eq:cond-r4}
\end{equation*}
Hence a uniform margin $(1+\eta)\alpha B_0\le1-\zeta-c$ with a constant
$c>0$ forces a positive, $G$-independent lower bound on the winner mass.  Conversely, if $G\to\infty$,
$\alpha\to\bar\alpha$, and $x\to0$, then
\begin{equation*}
 (1-\eta)\bar\alpha B_0\le1+\zeta.                       \condtag{R3}\label{eq:cond-r5}
\end{equation*}
Thus $(1+\eta)\alpha B_0=1-\zeta$ and
$(1-\eta)\alpha B_0=1+\zeta$ bracket the critical surface allowed by these
residual bounds.

For the conditioning statement, suppose $G\ge2$, $d\ge2$, and
$\widetilde p\ge p_0\in(0,1)$.  Set
\begin{equation}
 q_0=p_0(1-\zeta),\qquad
 \Gamma_*=\lambda_{\max}(C_*-v_*I),\qquad
 \delta_{\max}=\lambda_{\max}(C_L-v_*I).                 \label{eq:cond-robust-constants}
\end{equation}
Then
\begin{equation*}
 x\le\frac{d(1+\eta)\tau}{q_0},                           \condtag{R4}\label{eq:cond-r8}
\end{equation*}
and, whenever the denominator is positive,
\begin{equation*}
 \boxed{\tau\le
 \frac{(1-p_0)(1+\zeta)\delta_{\max}}
 {d\{(1-\eta)(G-1)-(1-p_0)(1+\zeta)(1+\eta)/q_0\}}.}     \condtag{R5}\label{eq:cond-r9}
\end{equation*}
For the log-det Hessian at the approximate point,
\begin{equation}
 \widetilde{\mathcal H}[Y]_g
 =\tau\widetilde Z_g^{-1}Y_g\widetilde Z_g^{-1}
 \quad\text{on }\mathcal T,                               \label{eq:cond-robust-hessian}
\end{equation}
we have
\begin{equation*}
 \boxed{\kappa(\widetilde{\mathcal H}|_{\mathcal T})\ge
 \frac{\delta_{\min}^2(1-\eta)q_0}
 {d(1+\eta)^2\tau(x+\Gamma_*)}=\Omega(G).}                \condtag{R6}\label{eq:cond-r13}
\end{equation*}
If $m\ge2$, the stronger bound is
\begin{equation*}
 \boxed{\kappa(\widetilde{\mathcal H}|_{\mathcal T})\ge
 \frac{\delta_{\min}^2(1-\eta)^2}{(1+\eta)^2x^2}
 =\Omega(G^2).}                                            \condtag{R7}\label{eq:cond-r14}
\end{equation*}
For fixed $p_0,\eta,\zeta,C_*,C_L$, all asymptotic constants are uniform
over such points.
\end{theorem}

\begin{proof}
The residual bound is equivalent to
\[
 (1-\eta)I\preceq\tau^{-1}\widetilde Z_g^{1/2}D_g
 \widetilde Z_g^{1/2}\preceq(1+\eta)I.
\]
Congruence by $\widetilde Z_g^{-1/2}$, inversion, and multiplication by
$\tau$ prove \eqref{eq:cond-r3}.  Since
$D_L\succeq C_L-v_*I$,
\[
 \sum_{g=2}^G\tr\widetilde Z_g
 \le(1+\eta)\tau(G-1)\tr D_L^{-1}
 \le(1+\eta)\tau(G-1)B_0.
\]
Divide by $\widetilde T\ge1-\zeta$ to prove
\eqref{eq:cond-r4}.  The lower sandwich similarly gives
$\widetilde T\ge(1-\eta)\tau(G-1)\tr D_L^{-1}$; taking the limit in the
theorem proves \eqref{eq:cond-r5}.

Winner mass and the upper sandwich give
\[
 q_0\le\tr\widetilde Z_*
 \le(1+\eta)\tau\tr D_*^{-1}
 \le\frac{d(1+\eta)\tau}{x},
\]
which is \eqref{eq:cond-r8}.  Every eigenvalue of $D_L$ is at most
$x+\delta_{\max}$.  The lower sandwich on loser blocks and
$(1-\widetilde p)\widetilde T\le(1-p_0)(1+\zeta)$ give
\[
 (1-p_0)(1+\zeta)\ge
 \frac{(1-\eta)\tau(G-1)d}{x+\delta_{\max}}.
\]
Substitute \eqref{eq:cond-r8} and rearrange to obtain
\eqref{eq:cond-r9}, hence $\tau=O(G^{-1})$.

No commutation is needed for the curvature tests.  In an eigenbasis of any
loser block $\widetilde Z_g$, a Frobenius-unit symmetric off-diagonal matrix
is tangent.  The upper sandwich and $D_L\succeq\delta_{\min}I$ bound every
eigenvalue of this block by $(1+\eta)\tau/\delta_{\min}$, so this direction
has curvature at least
$\delta_{\min}^2/\{(1+\eta)^2\tau\}$.  The largest eigenvalue of the winner
block is at least $q_0/d$, while its other eigenvalues are at least
$(1-\eta)\tau/(x+\Gamma_*)$.  An off-diagonal direction between a largest
eigenvector and an orthogonal eigenvector therefore has curvature at most
$d(x+\Gamma_*)/\{(1-\eta)q_0\}$.  Taking the quotient proves the finite
bound in \eqref{eq:cond-r13}; \eqref{eq:cond-r8}--\eqref{eq:cond-r9} give
its asymptotic form.

If $m\ge2$, the min--max principle and the lower sandwich show that the two
largest winner eigenvalues are at least $(1-\eta)\tau/x$.  Their
off-diagonal direction has curvature at most
$x^2/\{(1-\eta)^2\tau\}$.  Comparing with the loser off-diagonal direction
above proves \eqref{eq:cond-r14}.
\end{proof}

Scalar blocks, excluded above by $d\ge2$, also have a robust linear lower
bound.  When $d=1$ and $G\ge3$, \eqref{eq:cond-r8}--\eqref{eq:cond-r9}
remain valid.
A difference of two loser coordinates has curvature at least
$\delta_1^2/\{(1+\eta)^2\tau\}$.  Balancing a unit winner coordinate by
$-1/(G-1)$ in every loser yields a tangent Rayleigh quotient at most
\begin{equation}
 \frac{\tau}{q_0^2}+
 \frac{(x+\delta_1)^2}{(1-\eta)^2\tau(G-1)}.              \label{eq:cond-robust-scalar-soft}
\end{equation}
The ratio of these two curvature bounds, together with
\eqref{eq:cond-r8}--\eqref{eq:cond-r9}, is $\Omega(G)$.  The residual in
\eqref{eq:cond-robust-residual} must be measured relative to the barrier: a
small unscaled stationarity residual need not imply a Loewner sandwich near
the spectral edge, where $D_*$ becomes singular.

\subsection{Relation to prior work and open questions}
\label{sec:condensation-discussion}

The resolvent form \eqref{eq:cond-exact-center} of the center is classical:
it is the density associated with the matrix Burg entropy, and Ishihara
derives it from a maximum-entropy problem with trace and energy constraints
\cite{ishihara2019-density-operators-generalized-entropies}.  For individual
eigenmodes, the precise statistical-mechanics analogue is Rayleigh--Jeans
condensation, where inverse-gap occupancies, saturation of excited modes,
and a macroscopic ground-mode mass are established
\cite{baudin2020-rayleigh-jeans-condensation,zanaglia2024-bridging-rayleigh-jeans}.
Inverse-gap capacity and condensation onto several states also appear in
spherical models
\cite{lukkarinen2018-multi-state-condensation,crisanti2019-condensation-ordering}.
We therefore call our setting a matrix-Burg / Rayleigh--Jeans analogue and
claim neither a Bose--Einstein distribution nor a new form of the center.

Square-root rates on degenerate SDP central paths and generic ill-conditioning
of Newton systems near low-rank optima are also known
\cite{cruzneto2005-asymptotic-central-path,alizadeh1998-primal-dual-sdp,zhang2017-modified-interior-point-method-for},
but those works do not study the reduced Frobenius Hessian as the number $G$
of blocks grows.  Our contributions are the critical coefficients as $G$
grows, the dependence of the conditioning order ($G$ versus $G^2$) on $m$,
the finite mass--conditioning laws, their combination with visibility and
raw-query bounds, and the conservation law.  The quantum search,
amplification, and adversary tools are standard.

Identical losers are needed for a sharp common capacity, the exact collision
identity, and coherent preparation from public resolvents.  The
average-resolvent hypotheses \eqref{eq:cond-empirical-resolvent} extend the
mass law, but visibility for nonidentical losers requires a separate public
statistic such as \eqref{eq:cond-symmetry-light}.  It is open how the
bound $\eta$ on the barrier-norm residual in \eqref{eq:cond-robust-residual}
relates to the algorithm-specific inexact Newton residuals of
\cref{def:inexact-newton-residual}, whose forcing tolerance is a separate
$\eta$; we claim no conversion.

\section{Winner-Take-All Global-Trace Holonomy SDP}
\label{sec:winner-take-all}

This section gives an explicit SDP family with bounded scalar incidence to
which the condensation theory of \cref{sec:condensation} applies.  It couples
$G$ groups, each a copy of the qutrit-holonomy component of
\cref{sec:sdp-intrinsic}.  Each group has \(N=n_{\rm bit}\) private signs,
and its holonomy is their parity.  A single trace constraint shared by all
groups makes the optimal value the least minimum eigenvalue of the group
costs, and a chain of free scalar accumulators keeps every row and column of
the equality matrix sparse.

\begin{center}
\begin{tabular}{@{}ll@{}}
\toprule
quantity & conclusion \\ \midrule
value queries & $Q=\Theta(N\sqrt G)$, $R=\Theta(NG)$ for error $<1/20$ \\
one-odd central path & $p=\Theta(1),\Theta(G^{-1/2}),\Theta(G^{-1})$ \\
 & below, at, and above $\alpha_*=2/45$ \\
reduced-Hessian conditioning & $\Theta(G),\Theta(G),\Theta(1)$ in the same phases \\
coherent central-state preparation & $\Theta(N\sqrt G),\Theta(NG^{1/4}),O(N)$ \\
end-to-end algorithm & search for the winner, then solve only its group \\
\bottomrule
\end{tabular}
\end{center}
Here $Q$ and $R$ are the quantum and randomized query complexities of the
optimal value, $p$ is the central-path trace mass of the odd group when
exactly one group is odd, and the phases refer to $\alpha=\tau G$
(\cref{sec:wta-central-path}).  The preparation costs assume exactly one
winner.  At the public starting point the root Hessian has condition number
one, but the Newton-direction state does not reveal the winner
(\cref{rem:wta-dilution}).

\subsection{A bounded-incidence product SDP}
\label{sec:wta-construction}

Set
\begin{equation}
 K=16N,\qquad P=2K+N=33N,
 \qquad u=\frac1{10},\qquad q=\frac KP=\frac{16}{33},
 \qquad \gamma=\frac uq=\frac{33}{160}.
 \label{eq:wta-parameters}
\end{equation}
For $a<b$, write $H_{ab}=e_ae_b^\top+e_be_a^\top$, and put
$D_s=\diag(s,1,1)$ for $s\in\{-1,+1\}$.  Group $g$ has private signs
$\sigma_{g,1},\ldots,\sigma_{g,N}$ and holonomy
\begin{equation}
 h_g=\prod_{j=1}^N\sigma_{g,j}.
 \label{eq:wta-holonomy}
\end{equation}
Group $g$ has $P$ blocks $X_{g,0},\ldots,X_{g,P-1}\in\Splus^3$, joined in a
path.  For the edge from $i-1$ to $i$, let $G_{g,i}=I$ except that
$G_{g,K+j-1}=D_{\sigma_{g,j}}$ for $1\le j\le N$, and impose the copy
equations, each as six isometric-svec scalar equations:
\begin{equation}
 X_{g,i}=G_{g,i}X_{g,i-1}G_{g,i}^\top
 \qquad(1\le i<P).
 \label{eq:wta-copy-rows}
\end{equation}
If $R_{g,0}=I$ and $R_{g,i}=G_{g,i}\cdots G_{g,1}$, these equations state that
\begin{equation}
 X_{g,i}=R_{g,i}Z_gR_{g,i}^\top,\qquad Z_g=X_{g,0}\succeq0.
 \label{eq:wta-root-lift}
\end{equation}
The first $K$ matrices $R_{g,i}$ equal $I$, and the last $K$ equal
$D_{h_g}$.

The public block costs are
\begin{equation}
 C_{g,i}=\begin{cases}
 3I+uH_{23}+\gamma H_{12},&0\le i<K,\\
 3I+uH_{23},&K\le i<K+N,\\
 3I+uH_{23}+\gamma H_{13},&K+N\le i<P.
 \end{cases}
 \label{eq:wta-local-costs}
\end{equation}
Because $\sqrt{u^2+\gamma^2}<1/4$, every $C_{g,i}$ has spectrum in
$(11/4,13/4)$.  Introduce free public scalars $t_1,\ldots,t_G$ and add
\begin{align}
 t_1-\tr X_{1,0}&=0,\nonumber\\
 t_g-t_{g-1}-\tr X_{g,0}&=0 &&(2\le g\le G),\nonumber\\
 t_G&=1. \label{eq:wta-accumulator}
\end{align}
For $G=1$, the first and last rows impose $t_1=\tr X_{1,0}=1$.
The SDP minimizes
\begin{equation}
 f(X)=\frac1P\sum_{g=1}^G\sum_{i=0}^{P-1}
                  \inner{C_{g,i}}{X_{g,i}}
 \label{eq:wta-objective}
\end{equation}
subject to \eqref{eq:wta-copy-rows}--\eqref{eq:wta-accumulator} and
$X_{g,i}\succeq0$.

There are
\begin{equation}
 L=GP=33GN
 \label{eq:wta-block-count}
\end{equation}
blocks in $\Splus^3$.  A copy row has two scalar nonzeros.  The first
accumulator row has four, an interior accumulator row has five, and the
final row $t_G=1$ has one.
Every matrix coordinate and every free-scalar column occurs in at most two
rows.  Thus the equality matrix has scalar row sparsity at most five and
scalar column sparsity at most two.  Its support, right-hand side, and all
costs are public; only the $12$ and $13$ copy coefficients on the $N$ signed
edges of each path depend on the input.

Eliminating the path and accumulator variables leaves the roots $Z_g$,
whose feasible set is exactly
\begin{equation}
 \mathcal B_G=\left\{(Z_1,\ldots,Z_G):Z_g\succeq0,
                       \ \sum_{g=1}^G\tr Z_g=1\right\}.
 \label{eq:wta-root-feasible-set}
\end{equation}
The public point
\begin{equation}
 X_{g,i}^0=\frac1{3G}I,\qquad t_g^0=\frac gG
 \label{eq:wta-public-point}
\end{equation}
is strictly primal feasible.  With all equality multipliers zero, the dual
PSD slacks are $C_{g,i}/P\succ0$ and stationarity in every free scalar is
zero.  Hence the formulation is strictly primal--dual feasible in the sense
of \cref{def:strict-feasibility}.

\subsection{Exact spectral minimum}
\label{sec:wta-optimum}

In terms of the roots, the objective \eqref{eq:wta-objective} becomes
\begin{equation}
 \overline C_h=3I+u(H_{12}+H_{23}+hH_{13}),
 \qquad
 \min_{Z\in\mathcal B_G}\sum_g\inner{\overline C_{h_g}}{Z_g}.
 \label{eq:wta-effective-cost}
\end{equation}
Indeed, the first and last $K$ blocks of each path contribute
$K\gamma/P=u$ to $H_{12}$ and $H_{13}$, respectively, and
conjugation by $D_h$ changes $H_{13}$ to $hH_{13}$.  The two effective costs
have spectra
\begin{equation}
 \operatorname{spec}(\overline C_+)
   =\left\{\frac{16}5,\frac{29}{10},\frac{29}{10}\right\},
 \qquad
 \operatorname{spec}(\overline C_-)
   =\left\{\frac{31}{10},\frac{31}{10},\frac{14}5\right\}.
 \label{eq:wta-cost-spectra}
\end{equation}
These follow either by diagonalizing the two signed triangle adjacency
matrices, whose eigenvalues are $(2,-1,-1)$ and $(1,1,-2)$, or from the
intrinsic component analysis in \cref{sec:sdp-intrinsic}.

For any feasible roots, Rayleigh--Ritz gives
\begin{equation}
 \sum_g\inner{\overline C_{h_g}}{Z_g}
 \ge\sum_g\lambda_{\min}(\overline C_{h_g})\tr Z_g
 \ge\min_g\lambda_{\min}(\overline C_{h_g}).
 \label{eq:wta-rayleigh-ritz}
\end{equation}
Equality holds when all trace is placed on a minimum eigenvector of a
minimizing group.  Therefore
\begin{equation}
 \boxed{\operatorname{OPT}=
 \begin{cases}
 29/10,&h_g=+1\text{ for every }g,\\
 14/5,&h_g=-1\text{ for at least one }g.
 \end{cases}}
 \label{eq:wta-optimum}
\end{equation}
The gap is $1/10$, so the optimal value encodes
$\operatorname{OR}_G\circ\operatorname{PARITY}_N$: whether some $h_g=-1$.

The root feasible set \eqref{eq:wta-root-feasible-set} has no lift over a
finite product of second-order cones.  Its conic hull is $(\Splus^3)^G$.  If
\eqref{eq:wta-root-feasible-set} had such a lift, homogenizing the compact lift and
projecting to one factor would give such a lift of $\Splus^3$, contrary to
Fawzi's theorem \cite{fawzi2018-on-representing-the-positive-semidefinite}.
The path-and-accumulator formulation projects linearly onto the root set, so
it has no finite SOC lift either.

\subsection{Sharp value-query complexity}
\label{sec:wta-query-complexity}

We use the raw sparse coefficient access of \cref{def:raw-sparse-access}.
In the canonical oracle, which returns isometric-svec coefficients at fixed
positions, a public reversible address map sends a requested
input-dependent position to $(g,j)$ and a public position to a dummy zero.
The value oracle XORs the sign of one of the two hidden coefficients
$\pm\sigma_{g,j}$ directly into its target register.  Thus one raw-sign query, which returns one $\sigma_{g,j}$, simulates
one coefficient query, coherently and also under inverse
or external control.  Conversely, one fixed input-dependent copy coefficient
reveals $\sigma_{g,j}$.  The two models
therefore have the same query counts.  A clean value oracle can
instead be used as a phase oracle by compute--phase--uncompute, at two calls
per query.

We instantiate \cref{thm:cond-unique-or} for parity blocks.
Let $E,O\subseteq\{0,1\}^N$ be the even- and odd-parity strings and let
\begin{equation}
 B_{xy}=\mathbf1[|x-y|_1=1]\qquad(x\in E,\ y\in O).
 \label{eq:wta-hypercube}
\end{equation}
The constant vectors show that $\norm B=N$, while each
$B\circ\Delta_j$ is a permutation matrix and has norm one.  Hence the inner
adversary value in \eqref{eq:cond-inner-adversary} is exactly $N$.
Applying \eqref{eq:cond-unique-or-adv} on the all-even versus
exactly-one-odd promise gives
\begin{equation}
 Q_{1/3}=\Omega(N\sqrt G).
 \label{eq:wta-quantum-lower}
\end{equation}
The hypercube matrix \eqref{eq:wta-hypercube} is thus an explicit
adversary witness for \cref{sec:condensation-query-lower}.

If $|\widetilde v-\operatorname{OPT}|\le\epsilon<1/20$, thresholding at
$57/20$ distinguishes the two cases in \eqref{eq:wta-optimum}.  The same
threshold works for target-relative error $1/100$, since
\begin{equation}
 \frac{101}{100}\frac{14}5<\frac{57}{20}
 <\frac{99}{100}\frac{29}{10}.
 \label{eq:wta-relative-threshold}
\end{equation}
At $\epsilon=1/20$, the midpoint $57/20$ is a valid answer for both values
and uses no query, so the additive threshold is sharp.

For the randomized lower bound, let $\mathcal D_0$ make all groups
independent and uniform conditional on even parity; let $\mathcal D_1$
choose a uniform group $J$, make it uniform conditional on odd parity, and
leave the others as in $\mathcal D_0$.  Couple adaptive decision trees by returning common fair answers to the first $N-1$ distinct positions in each
group and using the final unseen sign to enforce the promised parity.  The
transcripts can differ only if the algorithm queries all $N$ signs of
group $J$.  A depth-$r$ tree queries all signs of at most $r/N$ groups, so
\begin{equation}
 \operatorname{TV}(\mathsf{Tr}_{\mathcal D_0},
                    \mathsf{Tr}_{\mathcal D_1})\le\frac r{NG}.
 \label{eq:wta-yao-coupling}
\end{equation}
Yao's principle therefore gives $r\ge NG/3$ for worst-case success $2/3$;
reading every sign is a matching upper bound.  Consequently, for every
$\epsilon<1/20$,
\begin{equation}
 \boxed{Q_\epsilon(\operatorname{OPT})=\Theta(N\sqrt G),
 \qquad R_\epsilon(\operatorname{OPT})=\Theta(NG).}
 \label{eq:wta-value-complexity}
\end{equation}

\subsection{The Newton step at the public point}
\label{sec:wta-public-start}

Use the log-det barrier on the PSD blocks, but no barrier on the free
accumulators:
\begin{equation}
 \Phi_\mu(X)=f(X)-\mu\sum_{g,i}\log\det X_{g,i},
 \qquad \mu=\frac1{GP}.
 \label{eq:wta-start-barrier}
\end{equation}
At \eqref{eq:wta-public-point}, the PSD-block gradient and Hessian in
isometric-svec coordinates are
\begin{equation}
 g_{g,i}=\frac1P(C_{g,i}-3I),
 \qquad \nabla_X^2\Phi_\mu(X^0)=\frac{9G}{P}I.
 \label{eq:wta-start-gradient-hessian}
\end{equation}
These and all equality residuals are public; the input enters the unreduced
Newton system only through the signed copy constraints in its matrix.

After exact elimination of the copies and accumulators, a tangent
direction in root coordinates is $Y=(Y_1,\ldots,Y_G)$ with
$\sum_g\tr Y_g=0$.  The copy maps are orthogonal congruences and each root
has $P$ path copies, so
\begin{equation}
 \boxed{\nabla^2_{\rm root}\Phi_\mu(X^0)=9G I_{6G-1}.}
 \label{eq:wta-root-hessian}
\end{equation}
This root-coordinate Hessian has condition number one, including in
directions that move trace between groups.  With
\begin{equation}
 A_h=H_{12}+H_{23}+hH_{13},
 \label{eq:wta-Ah}
\end{equation}
the pulled-back gradient is $uA_{h_g}$, which is trace zero, and the exact
Newton directions are
\begin{equation}
 \Delta Z_g=-\frac{u}{9G}A_{h_g},\qquad \Delta t_g=0.
 \label{eq:wta-newton-direction}
\end{equation}
Since $\norm{A_h}_F^2=6$, the decrement is
\begin{equation}
 \lambda^2=\sum_g\inner{uA_{h_g}}{(9G)^{-1}uA_{h_g}}
 =\frac1{150},\qquad \lambda=\frac1{\sqrt{150}}.
 \label{eq:wta-decrement}
\end{equation}
The full step remains strictly feasible.  The eigenvalues inside the
parentheses below are $14/45,31/90,31/90$ for $h=+1$ and
$29/90,29/90,16/45$ for $h=-1$.  Consequently,
\begin{equation}
 Z_g^0+\Delta Z_g=\frac1G\left(\frac13I-\frac u9A_{h_g}\right)
 \succeq\frac{14}{45G}I.
 \label{eq:wta-full-step-margin}
\end{equation}
On the root slice the barrier weight is $\tau=\mu P=1/G$ (see
\eqref{eq:wta-root-center-problem}), which is below one for every $G\ge2$;
then $\Phi_\mu$ is not standard self-concordant there, and the usual
full-step decrement estimate does not apply to $\lambda$.  We instead bound the next
decrement directly.  With
$B_h=\tfrac13I-\tfrac u9A_h$, the full step is $\widehat Z_g=B_{h_g}/G$ and
its slice gradient is $\overline C_{h_g}-B_{h_g}^{-1}$.  On the trace slice
the Newton decrement is the minimum over the trace multiplier of the
Hessian-inverse norm of the shifted gradient, so taking that multiplier to
be zero gives an upper bound:
\begin{equation}
 \lambda_{\rm next}^2
 \le\frac1G\sum_g\tr\bigl[(B_{h_g}\overline C_{h_g}-I)^2\bigr]
 =\frac{u^4}{81}\tr A_h^4=\frac{2u^4}{9},
 \qquad
 \lambda_{\rm next}\le\frac{\sqrt2}{300},
 \label{eq:wta-next-decrement}
\end{equation}
because $\overline C_h=3I+uA_h$ gives
$B_h\overline C_h-I=-\tfrac{u^2}9A_h^2$ and $\tr A_h^4=18$ for both
parities.  This bound is uniform in $G$ and in the parity pattern.

\begin{remark}[The Newton-direction state hides the winner]
\label{rem:wta-dilution}
The normalized root-direction state is
\begin{equation}
 \ket{\Delta(h)}=\frac1{\sqrt G}\sum_{g=1}^G\ket g\ket{a_{h_g}},
 \qquad
 \ket{a_h}=\frac{\ket{\operatorname{svec}(A_h)}}{\sqrt6}.
 \label{eq:wta-newton-state}
\end{equation}
Looping once over the $N$ raw signs and applying phase kickback only to the
$13$ coordinate prepares this state exactly with $N$ queries.  Conversely,
the density-matrix entry between group~1's $12$ and $13$ coordinates is
$h_1/(3G)$.  Since each density-matrix entry of a $q$-query
output has degree at most $2q$ in the input signs, exact preparation needs $q\ge N/2$; the same holds at trace error
$\delta<1/(6G)$.  Hence exact preparation costs $\Theta(N)$.

Moreover, the all-even and one-odd states have overlap $1-2/(3G)$ and trace
distance
\begin{equation}
 \frac{2\sqrt{3G-1}}{3G}=\Theta(G^{-1/2}).
 \label{eq:wta-newton-distance}
\end{equation}
Therefore a constant-error Newton state cannot decide the unique-winner
promise with constant bias.  The $\Omega(N\sqrt G)$ lower bound applies to
the value and to optimizer outputs that reveal the winner, not to this
state.
\end{remark}

The condition number one in \eqref{eq:wta-root-hessian} holds in root
coordinates only.  A feasible root direction lifts to accumulator changes
equal to prefix sums of traces, which is not an isometry, and the free
scalars have zero barrier Hessian.  We therefore do not claim condition
number one for the ambient augmented Hessian or the saddle-point KKT
matrix.

\subsection{Approximate optimizers}
\label{sec:wta-approximate-optimizers}

For exactly feasible roots on an instance with at least one odd group, put
$m_-=\sum_{g:h_g=-1}\tr Z_g$.  If the objective is at most
$14/5+\delta_{\rm obj}$, \eqref{eq:wta-cost-spectra} and
\eqref{eq:wta-rayleigh-ritz} give
\begin{equation}
 \frac{14}5+\frac{1-m_-}{10}\le \sum_g
 \inner{\overline C_{h_g}}{Z_g}
 \le\frac{14}5+\delta_{\rm obj},
 \qquad
 \boxed{m_-\ge1-10\delta_{\rm obj}.}
 \label{eq:wta-exact-optimizer-mass}
\end{equation}
When the instance contains both an odd and an even group and
$0\le\delta_{\rm obj}\le1/10$, the constant is sharp: put trace $1-10\delta_{\rm obj}$ in an odd
and $10\delta_{\rm obj}$ in an even minimum eigenspace.

A similar bound holds for approximate solutions of the sparse formulation,
which may violate its equalities.  Define the copy residuals
\begin{equation}
 E_{g,i}=X_{g,i}-G_{g,i}X_{g,i-1}G_{g,i}^\top
 \label{eq:wta-copy-residual}
\end{equation}
and let $r_1,\ldots,r_G,r_{\rm top}$ be the residuals in
\eqref{eq:wta-accumulator}.  Suppose $X_{g,i}\succeq0$ and
\begin{equation}
 \sum_{g,i\ge1}\norm{E_{g,i}}_F^2+
 \sum_{g=1}^G r_g^2+r_{\rm top}^2\le\epsilon^2,
 \qquad f(X)\le\operatorname{OPT}+\delta_{\rm obj}.
 \label{eq:wta-approx-contract}
\end{equation}
Put $Z_{g,i}=R_{g,i}^\top X_{g,i}R_{g,i}$ (each block in its root frame) and
telescope along the paths and the accumulator chain.
Cauchy--Schwarz and the spectral bound on $C_{g,i}$ give
\begin{align}
 \left|\sum_g\tr Z_{g,0}-1\right|&\le\sqrt{G+1}\,\epsilon,
 \label{eq:wta-root-trace-error}\\
 \left|\frac1P\sum_{g,i}\tr X_{g,i}-1\right|&\le\sqrt{GP+1}\,\epsilon,
 \label{eq:wta-physical-trace-error}\\
 \left|f(X)-\sum_g\inner{\overline C_{h_g}}{Z_{g,0}}\right|
 &<\frac{13}{4}\sqrt{GP}\,\epsilon.
 \label{eq:wta-objective-lift-error}
\end{align}
Combining these estimates with \eqref{eq:wta-cost-spectra} gives, for the
even root mass $m_+$ on a yes-instance,
\begin{equation}
 m_+\le10\delta_{\rm obj}+28\sqrt{G+1}\,\epsilon
              +\frac{65}{2}\sqrt{GP}\,\epsilon.
 \label{eq:wta-robust-even-mass}
\end{equation}

\begin{theorem}[Robust winner recovery]
\label{thm:wta-robust-optimizer}
Suppose an algorithm prepares a state within trace distance $1/100$ of the
trace-normalized physical density
\begin{equation}
 \rho(X)=\frac{\bigoplus_{g,i}X_{g,i}}
                 {\sum_{g,i}\tr X_{g,i}}
 \label{eq:wta-optimizer-density}
\end{equation}
for PSD blocks and free accumulators satisfying
\eqref{eq:wta-approx-contract}, where
\begin{equation}
 \delta_{\rm obj}\le\frac1{100},
 \qquad \epsilon\le\frac1{1000\sqrt{GP}}.
 \label{eq:wta-optimizer-tolerances}
\end{equation}
Then a fixed input-independent measurement decides whether any
group is odd with constant bias.  On every yes-instance, measuring the group register returns
an actual odd group with probability greater than $5/6$.  Consequently any
such first preparation (in the sense of \cref{def:amortized-accounting})
uses $\Omega(N\sqrt G)$ raw or canonical coefficient
queries.
\end{theorem}

\begin{proof}
Set $s=\sqrt{GP}\epsilon\le1/1000$.  Since $P\ge33$,
$\sqrt{G+1}\epsilon<s/4$.  Equations
\eqref{eq:wta-physical-trace-error} and
\eqref{eq:wta-robust-even-mass} put the physical trace of the even groups
below
$1/10+(81/2)s$ and the total trace above $1-(33/32)s$.  Their ratio is less
than $1/7$; trace error $1/100$ leaves odd-group probability above $5/6$.
For existence alone, use the input-independent observable
\begin{equation}
 M=\frac1{2/5}\left(\bigoplus_{g,i}C_{g,i}-\frac{57}{20}I\right)
 \label{eq:wta-existence-observable}
\end{equation}
which has norm at most one.  The same three telescoping estimates put
$f(X)/\overline T<2.82$ on yes-instances and $>2.89$ on no-instances,
where $\overline T=P^{-1}\sum_{g,i}\tr X_{g,i}$.  Equivalently,
$\tr(M\rho(X))<-3/40$ and $>1/10$, respectively.
The POVM $(I\pm M)/2$, repeated a constant number of times, therefore
decides the promise.  Now apply \eqref{eq:wta-quantum-lower}.
\end{proof}

The order $1/\sqrt{GP}$ of the residual tolerance in
\eqref{eq:wta-optimizer-tolerances} is necessary.  On an all-even base input,
flip each of the $GN$ raw signs in turn, take the neighboring instance's
exact odd optimizer, and average the resulting physical solutions.  This
PSD mixture has objective $14/5$, exact accumulator rows, and only the $GN$
flipped-edge copy residuals.  With
$r_-=(1,-1,1)^\top/\sqrt3$ and $Z_-^*=r_-r_-^\top$, each nonzero residual before
averaging has Frobenius norm
$\norm{D_-Z_-^*D_--Z_-^*}_F=4/3$.  The total Euclidean residual is therefore
\begin{equation}
 \frac4{3\sqrt{GN}}=\frac{4\sqrt{33}}{3\sqrt{GP}}.
 \label{eq:wta-residual-sharpness}
\end{equation}
Yet on the all-even input its density makes
\eqref{eq:wta-existence-observable} report a winner.  Feasibility alone is
also insufficient: the public point \eqref{eq:wta-public-point} is exactly
feasible for every input.  Hence both conditions in
\cref{thm:wta-robust-optimizer} are needed.

\subsection{Exact central path and condensation}
\label{sec:wta-central-path}

Eliminating a length-$P$ path multiplies the physical barrier coefficient by
$P$.  Thus, with $\tau=\mu P$, the root problem is
\begin{equation}
 \min_{Z_g\succ0,\ \sum_g\tr Z_g=1}
 \left\{\sum_g\inner{\overline C_{h_g}}{Z_g}
                    -\tau\sum_g\log\det Z_g\right\}.
 \label{eq:wta-root-center-problem}
\end{equation}
Its exact center is
\begin{equation}
 Z_g(\tau)=\tau(\overline C_{h_g}-\lambda I)^{-1},
 \qquad
 1=\tau\sum_g\tr(\overline C_{h_g}-\lambda I)^{-1},
 \label{eq:wta-exact-center}
\end{equation}
with the unique $\lambda<\min_g\lambda_{\min}(\overline C_{h_g})$.
This instantiates \cref{prop:cond-exact-center}.

On the all-even input, every group has trace exactly $1/G$.  With
$y=29/10-\lambda$ and $\alpha=\tau G$, the scalar equation and its positive
solution are
\begin{align}
 1&=\alpha\left(\frac2y+\frac1{y+3/10}\right),
 \label{eq:wta-even-secular}\\
 y_\alpha&=\frac{3\alpha-3/10+
 \sqrt{(3/10-3\alpha)^2+(12/5)\alpha}}2.
 \label{eq:wta-even-root}
\end{align}

Now suppose group~1 is the unique odd group and put $x=14/5-\lambda$.  The
winner and loser resolvents are
\begin{equation}
 A(x)=\frac1x+\frac2{x+3/10},
 \qquad
 B(x)=\frac1{x+2/5}+\frac2{x+1/10}.
 \label{eq:wta-resolvents}
\end{equation}
The exact central path is determined by
\begin{equation}
 \boxed{1=\tau\{A(x)+(G-1)B(x)\}},
 \qquad
 \boxed{p_G=\tau A(x),\quad1-p_G=\tau(G-1)B(x)}.
 \label{eq:wta-one-odd-center}
\end{equation}
Direct algebra gives
\begin{equation}
 A(x)-B(x)=
 \frac{3/250}{x(x+1/10)(x+3/10)(x+2/5)}>0
 \qquad(x>0).
 \label{eq:wta-resolvent-order}
\end{equation}
Together with \eqref{eq:wta-one-odd-center}, this yields
$p_G/(1-p_G)>1/(G-1)$, and hence $p_G\ge1/G$.  This implication is special
to the present spectra and is not part of the general two-cost model of
\cref{sec:condensation-model}.
Since $B$ is decreasing on $x>0$ and $B(0)=45/2$,
\eqref{eq:wta-one-odd-center} also gives
\begin{equation}
 p_G\ge1-\frac{45}{2}\tau(G-1).
 \label{eq:wta-finite-mass}
\end{equation}

The constants of \cref{thm:condensation-threshold} are
\begin{equation}
 B_0=B(0)=\frac{45}{2},\qquad
 \beta=-B'(0)=\frac{825}{4},\qquad
 \boxed{\alpha_*=B_0^{-1}=\frac2{45}}.
 \label{eq:wta-critical-constants}
\end{equation}
Applying \cref{thm:condensation-threshold} gives, for fixed $\alpha>0$ and
$G\to\infty$,
\begin{align}
 0<\alpha<2/45:&\quad
 p_G\longrightarrow1-\frac{45}{2}\alpha,
 &\frac{x}{\tau}&\longrightarrow
       \frac1{1-(45/2)\alpha},
 \label{eq:wta-subcritical}\\
 \alpha=2/45:&\quad
 \sqrt G\,p_G\longrightarrow\frac{\sqrt{33}}9,
 &\sqrt G\,x&\longrightarrow\frac2{\sqrt{825}},
 \label{eq:wta-critical}\\
 \alpha>2/45:&\quad
 Gp_G\longrightarrow\alpha A(x_\alpha),
 &B(x_\alpha)&=\frac1\alpha.
 \label{eq:wta-supercritical}
\end{align}
In the supercritical case, \eqref{eq:wta-resolvent-order} implies the strict
coefficient inequality $\alpha A(x_\alpha)>\alpha B(x_\alpha)=1$.
The coefficient $\sqrt{33}/9$ follows from
$\alpha_*\sqrt\beta=(2/45)(\sqrt{825}/2)$.

At the center, the reduced Hessian in the Frobenius inner product is
\begin{equation}
 \mathcal H_\tau[Y]_g=\frac1\tau D_gY_gD_g,
 \qquad D_g=\overline C_{h_g}-\lambda I,
 \qquad \sum_g\tr Y_g=0.
 \label{eq:wta-central-hessian}
\end{equation}
For one odd group, the spectral gaps are
$(x,x+3/10,x+3/10)$ in the winner and
$(x+1/10,x+1/10,x+2/5)$ in every loser.  For $G\ge3$,
\cref{lem:cond-hessian-sandwich} specializes to
\begin{equation}
 \kappa_{\rm red}\asymp
 \frac{(x+2/5)^2}{
  \min\{x(x+3/10),\ x^2+(x+1/10)^2/G\}}.
 \label{eq:wta-kappa-finite}
\end{equation}
Therefore \cref{thm:cond-conditioning-phases} gives
\begin{equation}
 \kappa_{\rm red}=\begin{cases}
 \Theta(G),&0<\alpha\le2/45,\\
 \Theta(1),&\alpha>2/45.
 \end{cases}
 \label{eq:wta-conditioning-transition}
\end{equation}
On the all-even input and for $G\ge2$, by contrast,
\begin{equation}
 \kappa_{\rm even}=\left(\frac{y_\alpha+3/10}{y_\alpha}\right)^2
 =\Theta_\alpha(1).
 \label{eq:wta-even-conditioning}
\end{equation}
The bound \eqref{eq:cond-mc3} simplifies here because $\min\{1,(2/5)/(3/10)\}=1$: for $G\ge3$,
\begin{equation}
 \kappa_{\rm red}\ge\frac{(G-1)p_G}{1-p_G}.
 \label{eq:wta-mass-conditioning}
\end{equation}
Thus a constant winner mass forces $\kappa_{\rm red}=\Omega(G)$ without
preconditioning.  This concerns the central path only and gives no lower
bound for methods with an input-dependent preconditioner.  The corresponding
visibility bounds are \cref{eq:cond-cv2,eq:cond-cv4}.

For approximate centers, we measure centrality by the dimensionless local
residual of \cref{thm:cond-robust}, not an unscaled stationarity error.
If positive definite
$\widetilde Z_g$ and a multiplier $\widetilde\lambda<14/5$ satisfy
\begin{equation}
 \left\|\tau^{-1}\widetilde Z_g^{1/2}
 (\overline C_{h_g}-\widetilde\lambda I)
 \widetilde Z_g^{1/2}-I\right\|_{\rm op}\le\eta<1,
 \qquad
 \left|\sum_g\tr\widetilde Z_g-1\right|\le\zeta<1,
 \label{eq:wta-robust-centrality}
\end{equation}
put
$\widetilde p_G=\tr(\widetilde Z_1)/\sum_g\tr(\widetilde Z_g)$.
The denominator need not equal one under
\eqref{eq:wta-robust-centrality}.  Then
\eqref{eq:cond-r3}--\eqref{eq:cond-r5} specialize to
\begin{equation}
 \widetilde p_G\ge
 1-\frac{(1+\eta)(45/2)\alpha}{1-\zeta}.
 \label{eq:wta-robust-mass}
\end{equation}
The approximate transition lies between
$(1+\eta)(45/2)\alpha=1-\zeta$ and
$(1-\eta)(45/2)\alpha=1+\zeta$.
Equation \eqref{eq:cond-r13} also gives
$\kappa(\widetilde{\mathcal H}|_{\mathcal T})=\Omega(G)$ whenever
$\widetilde p_G$ is bounded below.  Trace error $\delta$ lowers the probability of observing the winner label
by at most $\delta$.

\subsection{Preparation, visibility, and extraction}
\label{sec:wta-central-state}

The root and physical central densities are
\begin{equation}
 \rho_\tau=\bigoplus_g Z_g(\tau),
 \qquad
 \rho_\tau^{\rm phys}=\frac1P\bigoplus_{g,i}
 R_{g,i}Z_g(\tau)R_{g,i}^\top.
 \label{eq:wta-central-densities}
\end{equation}
Both return the unique winner with probability $p_G$ when the group label is
measured.  The ordering
\eqref{eq:wta-resolvent-order} proves $p_G\ge1/G$, so the general upper
bound $\max\{p_G,1/G\}$ in \eqref{eq:cond-visibility} becomes $p_G$.  Thus
\cref{thm:cond-visibility} gives
\begin{align}
 p_G-\frac1G&\le\Dtr(\rho_0,\rho_1)\le p_G,
 \label{eq:wta-visibility}\\
 c_1-c_0&=\frac{(Gp_G-1)^2}{G(G-1)},\qquad c_0=\frac1G.
 \label{eq:wta-collision}
\end{align}

These formulas bound the cost of the first preparation through a collision
test on two independent preparations, with no query to verify a candidate
winner.  If $\tau\le1/(90G)$, \eqref{eq:wta-finite-mass} gives $p_G\ge3/4$,
and \eqref{eq:wta-collision} gives $c_1-c_0\ge1/8$ for every $G\ge2$.
Two independently prepared states of trace error $\epsilon$ change the
collision gap $c_1-c_0$ by at most $4\epsilon$.  Thus every fixed $\epsilon<1/32$ leaves
a constant gap.  At $\tau\le1/(135G)$, one has $p_G\ge5/6$; trace error
$1/20$ leaves gap at least $2/9-1/5=1/45$.  \Cref{thm:cond-unique-or} then implies
\begin{equation}
 \boxed{Q_\rho=\Omega(N\sqrt G)}
 \label{eq:wta-first-use-lower}
\end{equation}
for a first root or physical central-density preparation at either stated
schedule and tolerance.  At the first schedule, the physical duality gap is $3G\tau\le1/30$.  Under
\cref{def:amortized-accounting}, \eqref{eq:wta-first-use-lower} includes
setup queries; it does not cover an algorithm given a parity table or a
central-state oracle, which already contains the aggregated input.

Under the promise of exactly one winner, there is a matching coherent
upper bound.  A clean parity evaluator costs $T_{\rm eval}=N$ raw queries.
Instantiating
\eqref{eq:cond-preparation-upper}--\eqref{eq:cond-preparation-upper-p} in
\cref{prop:cond-rejection-preparation} gives
\begin{equation}
 Q_U=O\!\left(N(\sqrt{Gp_G}+1)\right),
 \label{eq:wta-preparation-upper}
\end{equation}
because the normalized $3$-dimensional winner and loser resolvents are public
functions of $(G,\tau)$.  The acceptance probability is classically known,
so exact phase-matched amplitude amplification removes any logarithmic
overhead from the raw-query count; the conditional local preparation cost
$T_{\rm loc}$ is gate cost only.  With the lower bound
\eqref{eq:cond-prep-tradeoff}, this is tight when the amplification
factor $\sqrt{Gp_G}$ diverges and the error resolves the winner mass $p_G$.  Hence
\begin{equation}
 Q_U=\begin{cases}
 \Theta(N\sqrt G),&0<\alpha<2/45,\\
 \Theta(NG^{1/4}),&\alpha=2/45,\\
 O(N),&\alpha>2/45.
 \end{cases}
 \label{eq:wta-preparation-phases}
\end{equation}
The extension to $k$ promised winners,
\eqref{eq:cond-ap3}--\eqref{eq:cond-ap6} in
\cref{sec:condensation-exact-k-preparation}, is not needed here.

Given an exact central-state preparation,
\eqref{eq:cond-algorithm-phases} gives these costs for extracting the
winner label:
\begin{center}
\begin{tabular}{@{}lccc@{}}
\toprule
regime & $p_G$ & verified copies & coherent calls \\ \midrule
$0<\alpha<2/45$ & $\Theta(1)$ & $\Theta(1)$ & $\Theta(1)$ \\
$\alpha=2/45$ & $\Theta(G^{-1/2})$ & $\Theta(G^{1/2})$ & $\Theta(G^{1/4})$ \\
$\alpha>2/45$ & $\Theta(G^{-1})$ & $\Theta(G)$ & $\Theta(G^{1/2})$ \\
\bottomrule
\end{tabular}
\end{center}
Verifying a label costs another $N$ raw queries.  Coherent
amplification needs a preparation unitary, its inverse, and a clean marker;
one trace-close output state does not suffice
(\cref{sec:state-conversion} treats that case).  The table counts uses of an
existing preparation, whereas \eqref{eq:wta-first-use-lower} bounds the
first preparation.

Visibility also shows when a fixed error allows a zero-query answer.  If
the allowed root trace error is at least $p_G$, then, under the promise of
zero or one winner, the public all-even root center is a valid output, by \eqref{eq:wta-resolvent-order} and
\eqref{eq:cond-visibility}.  At criticality this
happens for $G=\Omega(\epsilon^{-2})$, and for fixed supercritical $\alpha$
it happens for $G=\Omega_\alpha(\epsilon^{-1})$.  In physical coordinates,
lifting that state through the input's prefix products $R_{g,i}$ costs
exactly $N$ phase queries and preserves trace distance, so
\begin{equation}
 p_G\le\epsilon\quad\Longrightarrow\quad
 Q_\epsilon(\rho_{\rm central}^{\rm root})=0,
 \qquad Q_\epsilon(\rho_{\rm central}^{\rm phys})\le N.
 \label{eq:wta-zero-query-boundary}
\end{equation}
For sufficiently large fixed $\alpha$, even the public maximally mixed
physical state is zero-query: with $y_\alpha$ from
\eqref{eq:wta-even-root}, it suffices that
\begin{equation}
 \frac13-\frac{\alpha}{y_\alpha+3/10}+p_G\le\epsilon.
 \label{eq:wta-public-physical-boundary}
\end{equation}
No $\Omega(N)$ physical-state lower bound is claimed above the transition.

For algorithm design, condensation reveals the winner only in central
states whose preparation already costs as much as search, and a constant
winner mass forces $\kappa_{\rm red}$ linear in $G$.  The conservation law
\eqref{eq:cond-conservation} and the direct-search bounds
\cref{eq:cond-direct-search,eq:cond-search-crossover} therefore favor
search-then-crossover (find the winner, then solve only its group), at
cost
\begin{equation}
 O(N\sqrt G)+T_{\rm component\ solve}+O(N)_{\rm recovery}.
 \label{eq:wta-architecture}
\end{equation}
This family does not exhibit a QIPM speedup.

\subsection{Search then solve}
\label{sec:wta-search}

Write $\sigma_{g,j}=(-1)^{z_{g,j}}$.  With the XOR-oracle answer register in
$\ket-= (\ket0-\ket1)/\sqrt2$, query the public indices $j=1,\ldots,N$ in
order while retaining the group address.  Phase kickback implements
\begin{equation}
 \ket g\longmapsto\prod_{j=1}^N(-1)^{z_{g,j}}\ket g=h_g\ket g
 \label{eq:wta-clean-marker}
\end{equation}
with exactly $N$ raw queries and no parity or prefix garbage.

\begin{theorem}[Search-then-solve upper bound]
\label{thm:wta-search-upper}
For failure probability $\delta$, the value of the SDP can be returned using
$O(N\sqrt G\log(1/\delta))$ raw queries.  The same procedure returns a verified winning label, or `NONE' on an
all-even input, with a compact description of the corresponding exact root
optimizer; another $O(N)$ queries prepare that optimizer's physical state
instead.  It needs no QRAM.
\end{theorem}

\begin{proof}
Apply bounded-error quantum search to the clean marker
\eqref{eq:wta-clean-marker}, and verify a returned label by evaluating its
parity.  Output $14/5$ after successful verification and $29/10$ when the
search reports `NONE'.  Standard reversible addressing uses
$O(N\sqrt G\log(GP)\log(1/\delta))$ elementary routing gates and
$O(\log G+\log P)$ work qubits, apart from raw-query implementation.

If $g_*$ is odd, take $Z_{g_*}=r_-r_-^\top$ with
$r_-=(1,-1,1)^\top/\sqrt3$; if no odd group exists, take $g_*=1$ and
$Z_{g_*}=(I-r_+r_+^\top)/2$ with $r_+=(1,1,1)^\top/\sqrt3$.  Set every other
root to zero.  This is a compact optimizer output in the sense of
\cref{def:compact-optimizer-output}.  Streaming the $N$ prefix signs of
$g_*$ lifts it through \eqref{eq:wta-root-lift} and prepares the normalized
physical direct-sum state.  Listing its $P=33N$ nonzero blocks instead costs
$\Theta(P)$ output work.
\end{proof}

Together, \cref{thm:wta-search-upper} and
\eqref{eq:wta-quantum-lower} show that bounded-error value
estimation and winner recovery cost $\Theta(N\sqrt G)$, using standard search
\cite{boyer1998-tight-bounds-on-quantum-searching,
durr1996-a-quantum-algorithm-for-finding}.

Exact quantum query complexity depends on the promise.  Under
the promise of zero or one winner,
phase-matched exact amplification followed by exact parity verification gives
\begin{equation}
 Q_E(\operatorname{OR}^{0,1}_G\circ\operatorname{PARITY}_N)
 =\Theta(N\sqrt G).
 \label{eq:wta-exact-promised}
\end{equation}
Without that promise, the unique multilinear output polynomial has a
nonzero degree-$NG$ monomial, so the polynomial method gives an $NG/2$
lower bound; computing every parity gives
\begin{equation}
 Q_E(\operatorname{OR}_G\circ\operatorname{PARITY}_N)=\Theta(NG).
 \label{eq:wta-exact-unrestricted}
\end{equation}
If exactly one winner is promised, the decision is trivial, but finding
the winner's label still costs
$\Theta(N\sqrt G)$ for $G\ge2$.

\begin{remark}[Stronger access and setup costs]
\label{rem:wta-access-scope}
An oracle for $h_g$ removes the factor $N$ only because it is a stronger
input.  Parallel query ports may reduce depth but not the total number of
queries.  The accumulator chain is public.  The end-to-end bound requires
compact optimizer output; writing out $G$ arrays costs their output size.  If preprocessing uses $S$ raw queries and
later calls use $q_t$, any procedure that outputs the value or a decodable
winner obeys
\begin{equation}
 S+\sum_tq_t=\Omega(N\sqrt G).
 \label{eq:wta-setup-charge}
\end{equation}
In particular, the queries needed to build an input-dependent block
encoding, QRAM parity table, prefix table, null-space basis, or the cost
matrix obtained by eliminating the path variables count toward $S$.
\end{remark}

\subsection{Formulation tradeoff and novelty}
\label{sec:wta-scope}

A single trace row on the roots has $3G$ nonzeros, but its lift from root
directions to the matrix variables is an isometry.  The accumulator
chain \eqref{eq:wta-accumulator} reduces row sparsity to five and column
sparsity to two without changing the root SDP or the root-coordinate
Hessian, but its lift is not an isometry (\cref{sec:wta-public-start}).
Replacing the free scalars by differences of nonnegative variables, or
adding scalar log barriers, would change the geometry.  Thus we do not claim
a formulation with both bounded incidence and ambient condition number
one.

The query tools are prior work: adversary composition is due to
Høyer--Lee--Špalek, Belovs--Lee, and Reichardt; the randomized composition
framework includes Ben-David--Kothari; search and minimum finding are due to
BBHT and Dürr--Høyer; and fixed-point search is standard
\cite{hyer2005-tight-adversary-bounds-for-composite,
belovs2020-the-quantum-query-complexity-of,
reichardt2011-reflections-for-quantum-query-algorithms,
bendavid2018-randomized-query-complexity,
boyer1998-tight-bounds-on-quantum-searching,
durr1996-a-quantum-algorithm-for-finding,
yoder2014-fixed-point-quantum-search}.  Embeddings of Boolean functions into
conic optimal values, and sparsity-dependent coefficient-query lower bounds,
already appear in van
Apeldoorn et al., Theorem~29 and Corollary~30, and Apers--Gribling,
Theorem~8.4
\cite{apeldoorn2020-quantum-sdp-solvers,
apers2026-quantum-speedups-linear-programming-interior}.  Fawzi's non-SOC
theorem and the Rayleigh--Ritz trace identity are likewise established
tools.

To our knowledge, what is new is only their combination: a
bounded-incidence public accumulator chain for the trace normalization, an
exact spectral minimum that encodes
$\operatorname{OR}\circ\operatorname{PARITY}$ with a constant gap, a
public start whose root Hessian has condition number one, the exact
constants in the condensation transition, and the preparation costs in each
phase.  None of Grover search, parity, composition, the Rayleigh--Ritz trace
identity, generic SDP ill-conditioning, or log-det resolvents is claimed as
new by itself.

\section{Hidden-Block State Conversion}
\label{sec:state-conversion}

The central paths of \cref{sec:wta-central-path,sec:condensation-multiwinner}
concentrate trace on a hidden set of winning blocks.  This section studies the query cost of
preparing such a central state to trace-distance accuracy $\varepsilon$ when
the number of winners is promised to be either zero or $k$.  The two
extremes already appear in \cref{sec:condensation-exact-k-preparation}: a
sufficiently large error allows the public all-loser output, whereas a
sufficiently accurate coherent preparation costs as much as search.  In
between, the collision statistic in
\eqref{eq:cond-multi-collision} loses a factor $1/k$, and a bound on
pairwise trace distance gives no single measurement that works for every
unknown winner set.  Instead, we measure the label of the prepared state and
verify it with fresh queries.  This verified-sample decoder accepts a
zero-winner input only through verifier error and applies directly to
trace-close mixed outputs.

For parity blocks, the polynomial method then gives the query cost at every
nontrivial accuracy.  The key step is that averaging over inputs with fixed
block parities removes every character supported on a proper nonempty part
of a block, which divides the polynomial degree by the block length.  We
also compare the coupled centers with a batch of uncoupled qutrit-holonomy
SDPs.  Returning all component values admits no amortization, whereas
estimating their average has the usual approximate-counting query
complexity.

\subsection{The zero-versus-\texorpdfstring{$k$}{k} conversion problem}
\label{sec:state-conversion-model}

Let $f:\mathcal X_0\sqcup\mathcal X_1\to\{0,1\}$ be a possibly partial
Boolean predicate on one block.  Its bounded-error query complexity and
negative-weight adversary value satisfy
\[
 Q(f)=\Theta(A),\qquad A=\operatorname{Adv}^{\pm}(f).
\]
There are $G$ independently queried blocks with inputs $x_1,\ldots,x_G$; in
\cref{sec:winner-take-all}, block $g$ is group $g$, whose input is its sign
string $(\sigma_{g,1},\ldots,\sigma_{g,N})$, not one of its PSD blocks
$X_{g,i}$.  The winner set
$S=\{g:f(x_g)=1\}$ is promised to have size zero or $k$, where
$1\le k<G$.  One raw query addresses one coordinate of one block.  The
target on a $k$-winner input is
\begin{equation}
 \rho_S=\frac pk\sum_{g\in S}\ket g\!\bra g\otimes\rho_W
 +\frac{1-p}{G-k}\sum_{g\notin S}\ket g\!\bra g\otimes\rho_L,
 \qquad
 \rho_0=\frac1G\sum_{g=1}^G\ket g\!\bra g\otimes\rho_L^0,
 \label{eq:sc-targets}
\end{equation}
where $p=p_{G,k}$ is the winner mass and the constant-dimensional internal
states $\rho_W$, $\rho_L$, and $\rho_L^0$ are public.
Write
\begin{equation}
 \eta_{\rm int}=\Dtr(\rho_L,\rho_L^0),\qquad \delta=p-\varepsilon.
 \label{eq:sc-parameters}
\end{equation}
Here \(\eta_{\rm int}\) is the local internal-state discrepancy, not the
residual tolerance \(\eta\) used elsewhere.
The task is to output, on every promised input, a state within trace distance
$\varepsilon$ of its target.  The upper bounds below use a clean coherent
evaluator as in \eqref{eq:cond-clean-evaluator}, which computes the
membership bit for $S$, at cost $T_{\rm eval}$ per call.  An exact evaluator
costs $\Theta(Q_E(f))$ raw queries, where $Q_E$ denotes exact quantum query
complexity.  A bounded-error evaluator may be used instead after coherent
error reduction to operator-norm error below the per-call budget of the
amplification, at cost
$O(Q(f)\log(G/\min\{\delta,\varepsilon\}))$ per call; the tildes absorb
this factor.  For $f=\operatorname{PARITY}_N$, $Q_E(f)=\lceil N/2\rceil$, and
$Q(f)$ and $\operatorname{Adv}^{\pm}(f)$ are both $\Theta(N)$.  As in the
condensation applications, we assume $p\ge k/G$ when invoking rejection
preparation.
For classical algorithms we count only queries; preparing the output state
from the classical transcript is free.

Consider the test that measures the label of an $\varepsilon$-accurate
output and then evaluates $f$ on the sampled block.  On a $k$-winner input,
it accepts with probability at least $p-\varepsilon=\delta$, ignoring
verifier error.  On a zero-winner input, every proposed block is a loser, so
only verifier error can cause acceptance, whatever state the circuit
produces.  Unlike the collision statistic, this test is a single
measurement that works for every winner set $S$.

\begin{theorem}[Verified-sample lower bound]
\label{thm:sc-m1}
Fix $\delta\in(0,1)$.  Suppose a $T$-query circuit solves
\eqref{eq:sc-targets} with $\varepsilon\le p-\delta$.  If
\begin{equation}
 \sqrt{G/k}\ge C\delta^{-1}\log(e/\delta)
 \label{eq:sc-m1-domain}
\end{equation}
for a sufficiently large universal constant $C$, then
\[
 T=\Omega\!\left(\delta A\sqrt{G/k}\right).
\]
For $f=\operatorname{PARITY}_N$, this is
$T=\Omega(\delta N\sqrt{G/k})$.
\end{theorem}

\begin{proof}
First reduce the zero-versus-one problem of \cref{thm:cond-unique-or} to
zero-versus-$k$.  Put $G'=\lfloor G/k\rfloor$.  Replace each of the $G'$
input blocks by $k$ addressable copies and pin the remaining
$G-kG'<k$ blocks to a fixed public element of $\mathcal X_0$.  A query to a
copy is simulated by one query to its original.  Zero winners remain zero,
and a unique winner produces exactly $k$ winners.  The Cartesian-product
promise is preserved.  Since \eqref{eq:sc-m1-domain} forces $G/k$ to be at least a large
constant, $G'=\Theta(G/k)$, and
\eqref{eq:cond-unique-or-adv} gives a decision lower bound
$\Omega(A\sqrt{G/k})$.  This construction also works when $k$ does not divide $G$.

Now repeat the following trial $R=\Theta(1/\delta)$ times.  Run the
preparation circuit afresh, measure its label $g$, and evaluate $f(x_g)$
with error at most $\beta=\delta/8$.  Standard error reduction costs
$O(Q(f)\log(e/\delta))$ queries per verification.  On a $k$-winner input,
contractivity of trace distance under the two-outcome measurement
$\{\sum_{g\in S}\ket g\!\bra g\otimes I,I-\cdots\}$ gives
\[
 \Pr[g\in S]\ge p-\varepsilon\ge\delta,
 \qquad \Pr[\text{accept}]\ge\delta(1-\beta)\ge\delta/2.
\]
On a zero-winner input, every label fails the true predicate, regardless of
the prepared state, and $\Pr[\text{accept}]\le\beta=\delta/8$.  A
multiplicative Chernoff bound distinguishes these two rates with constant
error using $R=\Theta(1/\delta)$ trials.  Its total cost
\[
 O\!\left(\delta^{-1}
   [T+Q(f)\log(e/\delta)]\right)
\]
is therefore $\Omega(A\sqrt{G/k})$.  Since $Q(f)=\Theta(A)$,
\eqref{eq:sc-m1-domain} absorbs the verification term and proves the claim.
\end{proof}

The domain condition \eqref{eq:sc-m1-domain} is restrictive.  At
constant $\delta$ it requires only $G/k=\Omega(1)$; \cref{cor:sc-m3}
assumes the stronger uniform condition $k=O(G/\log^2G)$.  In general it
requires $G/k=\Omega(\delta^{-2})$ up to logarithms.  In exchange,
\cref{thm:sc-m1} holds for a general inner predicate $f$ and assumes no
access to a purification or to the inverse of the preparation.

\begin{theorem}[Detuned rejection preparation]
\label{thm:sc-m2}
Assume $p\ge k/G$, $1-p\ge c_0>0$, and $0<\varepsilon<p$.  With
$\delta=p-\varepsilon$, if $4\eta_{\rm int}\le\varepsilon$, then
\begin{equation}
 Q_\varepsilon=O\!\left(
 T_{\rm eval}\left[1+\sqrt{\delta G/k}\right]c_0^{-1/2}
 \log\frac{e}{\min\{\delta,\varepsilon\}}\right).
 \label{eq:sc-m2-upper}
\end{equation}
If $0<\varepsilon<4\eta_{\rm int}$, then
\begin{equation}
 Q_\varepsilon=O\!\left(
 T_{\rm eval}\sqrt{G/k}\log(e/\varepsilon)\right).
 \label{eq:sc-m2-fallback}
\end{equation}
When $\varepsilon=\Omega(\delta)$, the logarithm
in \eqref{eq:sc-m2-upper} is simply $O(\log(e/\delta))$.
\end{theorem}

\begin{proof}
We prepare with a detuned winner mass $m$, chosen in two regimes:
\begin{equation}
 \begin{aligned}
 m&=\begin{cases}
 \max\{3\delta/2,k/G\},&\delta\le2p/3,\\
 \max\{p-\varepsilon/2,k/G\},&\delta>2p/3,
 \end{cases}\\[-2pt]
 m&\le p:\quad
 \begin{cases}
 3\delta/2\le p\ \text{and}\ k/G\le p,&\delta\le2p/3,\\
 p-\varepsilon/2\le p\ \text{and}\ k/G\le p,&\delta>2p/3.
 \end{cases}
 \end{aligned}
 \label{eq:sc-detuning}
\end{equation}
Then $k/G\le m\le p$, $m-\delta\ge\delta/2$ in the first
regime, and $m-\delta\ge\varepsilon/2$ in the second.  Moreover
$\sqrt{Gm/k}=O(1+\sqrt{\delta G/k})$: this is immediate in the first
regime, while $m=\Theta(\delta)$ in the second.  The cap $m\le p$
is essential: with $1-p\ge c_0$, it keeps the zero-winner acceptance
probability from adding a factor $(1-m)^{-1/2}$ to the amplification cost.

Run the rejection construction of \cref{prop:cond-rejection-preparation}
with the label copied before amplification and the copy included in every
reflection, and with detuned weights
\[
 a'_W=m/k,\qquad a'_L=(1-m)/(G-k),\qquad a'_{\max}=m/k.
\]
Starting from the uniform label superposition, copy the label to an
environment register, compute the membership bit, rotate an acceptance
qubit with amplitude $\sqrt{a'_b/a'_{\max}}$, and conditionally attach
public purifications of $\rho_W$ and $\rho_L$.  A trial has acceptance
probability
\begin{equation}
 P_k=\frac{k}{Gm},\qquad
 P_0=\frac{k(1-m)}{m(G-k)}
 \label{eq:sc-branch-acceptance}
\end{equation}
on $k$-winner and zero-winner inputs, respectively.  Both are at least a
constant times $c_0k/(Gm)$, because $m\le p$ implies
$1-m\ge1-p\ge c_0$.

Apply fixed-point amplification simultaneously to the two possible
acceptance angles \cite{yoder2014-fixed-point-quantum-search}.  The
fixed-point reflections act on every register the trial touches, including
the environment copy of the label, which is uncomputed and recomputed inside
each reflection.  With
\[
 r=\frac18\min\{\delta,\varepsilon\},
\]
the acceptance-subspace residual can be made at most $r^2$ on both kinds of input
using
\[
 O\!\left(\sqrt{Gm/k}\,c_0^{-1/2}\log(e/r)\right)
\]
evaluator calls.  The quadratic residual is necessary: overlap error of a
purification becomes its square root in trace distance.  Writing $q=k/G$,
\eqref{eq:sc-branch-acceptance} gives
$P_0/P_k=(1-m)/(1-q)$.  Thus $m=q$ gives equal acceptance probabilities.  When
$q=o(m)$, however, an exact schedule tuned to the larger probability $P_k$
under-rotates a zero-winner input and can leave $\Theta(m^2)$ junk probability,
which exceeds the error budget for constant $p$.  Fixed-point amplification avoids
this mismatch.

Conditioned on acceptance, the output on a $k$-winner input is exactly
\eqref{eq:sc-targets} with winner mass $m$ in place of $p$; the zero-winner
output is exactly $G^{-1}\sum_g\ket g\!\bra g\otimes\rho_L$.  Hence the
respective errors are at most
\[
 (p-m)+r,\qquad \eta_{\rm int}+r.
\]
If $\delta\le2p/3$, then $\varepsilon=p-\delta\ge\delta/2$ and
$p-m\le\varepsilon-\delta/2$, so both errors are below
$\varepsilon$.  If $\delta>2p/3$, then $p-m\le\varepsilon/2$ and
$r\le\varepsilon/8$, so again both errors are at most $\varepsilon$.  Substituting the bound on
$m$ gives \eqref{eq:sc-m2-upper}.

For \eqref{eq:sc-m2-fallback}, first decide whether there are zero or $k$
winners by fixed-point search, reducing its error to $O(\varepsilon)$.  If
the decision is $k$, use the exactly-$k$ circuit \eqref{eq:cond-ap3}; if it
is zero, prepare $\rho_0$ directly.  The decision error contributes at most
$O(\varepsilon)$ in trace distance, and the stated cost dominates both
steps.  Deciding first avoids circular use of the exactly-$k$ promise.
\end{proof}

The precision logarithm in \eqref{eq:sc-m2-upper} depends on
$\min\{\delta,\varepsilon\}$, not on $\delta$ alone.  In the second regime of \eqref{eq:sc-detuning}, a
fixed-point residual of order $\delta$ would not suffice when
$\varepsilon\ll\delta$; the residual must be $O(\varepsilon)$.

\begin{corollary}[Constant-visibility regime]
\label{cor:sc-m3}
Let $f=\operatorname{PARITY}_N$, let $p$ be bounded away from zero and one,
and suppose
\[
 k=O(G/\log^2G),\qquad
 4\eta_{\rm int}\le\varepsilon\le(1-c)p
\]
for a constant $c>0$.  Then, up to fixed-point precision logarithms,
\[
 Q_\varepsilon=\Thetatilde\!\left(N\sqrt{G/k}\right).
\]
By \eqref{eq:sc-m2-fallback}, the same scale holds for
$\varepsilon<4\eta_{\rm int}$ when $\delta=\Omega(1)$.  If $\varepsilon$
exceeds the minimax radius of the target family, zero queries suffice; when
the internal loser states agree, \eqref{eq:cond-multi-visibility} bounds
the public output's error by $\max\{p,k/G\}$.
\end{corollary}

\begin{proof}
Here $\delta=p-\varepsilon=\Omega(1)$ and $A=\Theta(N)$.
The domain assumption makes \eqref{eq:sc-m1-domain} valid, so
\cref{thm:sc-m1} gives the lower bound.  The upper bound is
\cref{thm:sc-m2}.  The zero-query statement follows by returning the public
all-loser target and using \eqref{eq:cond-multi-visibility}.
\end{proof}

\begin{theorem}[Classical query complexity]
\label{thm:sc-m4}
For $f=\operatorname{PARITY}_N$, suppose $p\ge k/G$,
$1-p=\Omega(1)$, and
\begin{equation}
 2\eta_{\rm int}\le\varepsilon,
 \qquad Ck/G\le\delta=p-\varepsilon\le1/2.
 \label{eq:sc-m4-domain}
\end{equation}
Here $C>0$ is a sufficiently large universal constant.
Then the randomized classical raw-query complexity is
\[
 R_\varepsilon=\Theta(N\delta G/k).
\]
\end{theorem}

\begin{proof}
For the upper bound, inspect
$t=\Theta(\delta G/k)$ uniformly random blocks without replacement and
compute each parity exactly using $N$ queries.  Put $q=k/G$.  The probability
of finding at least one winner is
\[
 P_{\rm find}=1-\frac{\binom{G-k}{t}}{\binom Gt}.
\]
The constant in $t$ can be chosen so that
$P_{\rm find}\ge\delta$ throughout
\eqref{eq:sc-m4-domain}.  This use of a linear budget relies on
$\delta\le1/2$; as $\delta\to1$, the correct sampling budget contains
$(G/k)\log(1/(1-\delta))$.

Let $E_{\rm all}$ be the event that every sampled block is a winner.  Its
public probability is at most $q^t\le q\le\delta$.  On $E_{\rm all}$ output
a uniformly selected sampled winner.  On
$E_{\rm find}\setminus E_{\rm all}$, do the same with a public probability
chosen so that the total probability of winner output is exactly $\delta$;
this is possible because
$\Pr(E_{\rm all})\le\delta\le\Pr(E_{\rm find})$.  In all other cases select
uniformly among sampled losers.  A loser is available whenever the latter
rule is used.  Exchangeability makes the selected winner uniform on $S$ and
the selected loser uniform on its complement.  Attaching the corresponding
public internal state therefore prepares exactly the target shape with
winner mass $\delta$ in place of $p$.  Its trace error on a $k$-winner input is
$p-\delta=\varepsilon$.  On a zero-winner input the output is the uniform-label
$\rho_L$ state and has error $\eta_{\rm int}\le\varepsilon/2$.  Thus
$O(N\delta G/k)$ queries suffice.

For the lower bound, append exact $N$-query parity verification to a sampled
output label and repeat the preparation $\Theta(1/\delta)$ times.  This
decides zero versus $k$ winners with one-sided error and cost
$O(\delta^{-1}(T+N))$.  For the decision lower bound, use the $k$-copy
embedding from \cref{thm:sc-m1} and draw each original block uniformly
conditioned on its parity.  Every proper subset of a block is uniform, so a
classical transcript is independent of the planted parity until all $N$
positions of that block have been queried.  A $q$-query tree completes at
most $q/N$ original blocks, and it hits the uniformly planted winner with
probability at most $O(qk/(NG))$.  Yao's principle therefore gives
$\Omega(NG/k)$ decision queries.  Consequently
\[
 \delta^{-1}(T+N)=\Omega(NG/k).
\]
The lower restriction $\delta\ge Ck/G$ absorbs the additive $N$ term.
\end{proof}

We claim no lower bound when $\delta=o(k/G)$: there, a public output with
uniform label may already meet the accuracy requirement.  The upper
restriction $\delta\le1/2$ in \eqref{eq:sc-m4-domain} is also needed:
near one, the probability of finding a winner is no longer linear in the
number of sampled blocks.

\subsection{Composed small-success search}
\label{sec:state-conversion-small-success}

The verified-sample argument loses a factor $\sqrt\delta$ by repeating the
preparation.  For parity blocks, the polynomial method recovers it.  We
state the degree-division step as a lemma because
\cref{sec:state-conversion-batch} also uses it.

\begin{lemma}[Parity collapse]
\label{lem:sc-parity-collapse}
Let a $T$-query quantum algorithm access
$x\in\{\pm1\}^{G\times N}$, and let $P(x)$ be the probability of a fixed
final measurement event.  For block parities
$z_g=\prod_{j=1}^N x_{g,j}$, define
\[
 q(z)=\mathbb E[P(x)\mid \textstyle\prod_jx_{g,j}=z_g
                  \text{ for every }g].
\]
Here the conditional expectation is uniform over each parity fiber,
independently across blocks.
Then $q$ is a real multilinear polynomial in $z\in\{\pm1\}^G$ of degree
at most $\lfloor2T/N\rfloor$.  If $0\le P\le1$, then $0\le q\le1$ on
the Boolean cube.
\end{lemma}

\begin{proof}
The polynomial method writes $P$ as a real multilinear polynomial of degree
at most $2T$ \cite{beals2001-quantum-lower-bounds-by-polynomials}.  Consider
a monomial $x_R$.  Conditional expectation factors over blocks.  In one
block, the average of a character on a parity coset is one when its share of
$R$ is empty, is $z_g$ when that share is the full block, and is zero for
every proper nonempty share.  The last assertion follows by pairing inputs
that differ in two coordinates, one inside and one outside the proper
share.  Thus a raw monomial with nonzero average becomes a monomial in $z$ whose
degree is the number of full blocks it contains, at most
$|R|/N\le2T/N$.  Averaging all monomials proves the claim.
\end{proof}

\begin{remark}[Computational check]
We checked the character calculation in \cref{lem:sc-parity-collapse} by
exhaustive enumeration for $N=2,3,4,5$, and the full identity on random
multilinear polynomials.
\end{remark}

\begin{theorem}[One-sided composed small-success bound]
\label{thm:sc-small-success}
Suppose a $T$-query quantum algorithm accepts with probability at most
$\beta\le\delta/2$ on every all-even input and at least $\delta$ on every
input with exactly $k$ odd-parity blocks.  Then, for absolute constants
$c_1,c_2>0$,
\begin{equation}
 T\ge c_1N\sqrt{\delta G/k}-c_2N.
 \label{eq:sc-small-success}
\end{equation}
\end{theorem}

\begin{proof}
Let $P(x)$ be the acceptance probability.  By
\cref{lem:sc-parity-collapse}, conditioning on all block parities gives a
polynomial $q(z)$ of degree at most
$D=\lfloor2T/N\rfloor$, bounded between zero and one at every Boolean
input.  It is at most $\beta$ at the all-even point and at least $\delta$
at every parity vector of odd-block weight $k$.

Average $q$ over all block permutations.  Minsky--Papert symmetrization
\cite{minsky1988-perceptrons} gives a univariate polynomial $h$ of degree at
most $D$ such that
\[
 0\le h(w)\le1\quad(w=0,1,\ldots,G),\qquad
 h(0)\le\beta,\qquad h(k)\ge\delta.
\]
Fix a sufficiently small absolute $\gamma>0$.  First suppose
$D\le\gamma\sqrt G$.  The Coppersmith--Rivlin inequality bounds a
degree-$D$ polynomial bounded at the $G+1$ equally spaced integer points by
\[
 \sup_{t\in[0,G]}|h(t)|\le C(\gamma),
\]
where $C(\gamma)$ is an absolute constant in this degree range
\cite{coppersmith1992-growth-polynomials,klauck2007-quantum-classical-sdpt}.
We use only the existence of the universal Coppersmith--Rivlin constants,
which the original paper does not give numerically; this is the form quoted
in Theorem~8 of Klauck--\v Spalek--de Wolf.  Erd\'elyi proves a related
explicit-constant inequality in a different, single-point-versus-supremum
parametrization
\cite{erdelyi2016-coppersmith-rivlin}.

Markov's inequality, rescaled from $[-1,1]$ to $[0,G]$, now gives
\[
 \sup_{[0,G]}|h'|\le\frac{2D^2C(\gamma)}G.
\]
The mean value theorem on $[0,k]$ gives a point $\xi$ with
\[
 h'(\xi)\ge\frac{h(k)-h(0)}k
 \ge\frac{\delta-\beta}k\ge\frac{\delta}{2k}.
\]
Therefore $D=\Omega(\sqrt{\delta G/k})$.  In the complementary case
$D>\gamma\sqrt G$, the same conclusion follows directly from
$\delta\le1\le k$.  Finally, $D\le2T/N$, with an additive $O(N)$ allowance for integer
rounding, proves
\eqref{eq:sc-small-success}.
\end{proof}

To our knowledge, \cref{thm:sc-small-success} is new as stated, but its
ingredients are standard and experts could plausibly derive it.  It is not
a new general composition theorem.

\begin{corollary}[Composed search]
\label{cor:sc-s1}
If a quantum algorithm outputs an odd-parity block with probability at least
$\delta$ whenever exactly $k$ blocks are odd, then
\[
 T=\Omega\!\left(\sqrt\delta\,N\sqrt{G/k}\right)-O(N).
\]
This is tight up to constants and the additive $N$ term.
\end{corollary}

\begin{proof}
Verify the output block by computing its parity exactly with $N$ fresh
queries.  Acceptance is then zero on an all-even input and at least
$\delta$ on an exactly-$k$ input, so \cref{thm:sc-small-success} applies to
$T+N$ queries.  Conversely, if $\delta\le4k/G$ a uniformly random label is
odd with probability $k/G\ge\delta/4$ at zero queries.  Otherwise $k/G<1/4$,
so the initial Grover angle $\theta=\arcsin\sqrt{k/G}$ is below $\pi/6$, and
$t=\lceil c\sqrt{\delta G/k}\rceil\ge1$ Grover iterations with $c=1/6$ keep
$(2t+1)\theta\le\pi/2$; each uses a $\Theta(N)$-query parity marker, and
together they produce a winner with probability
$\sin^2((2t+1)\theta)\ge\sin^2(2c\sqrt\delta)=\Theta(\delta)$.  Direct exact
verification costs another $N$ queries.
For one marked item, Zalka proved exact optimality
\cite{zalka1999-grovers-quantum-searching-algorithm-is}.  The small-success
scaling for $k$ marked items follows from the angle bound in Theorem~2 of
the published Dohotaru--H\o yer paper (Theorem~9 of the arXiv version) plus
the standard $k$-copy embedding.
Carolan--Poremba Corollary~2.4 records the packaged inequality
$\delta\le8(T+1)^2k/G$; our attribution uses the angle theorem and embedding,
not their intermediate distance-based derivation
\cite{dohotaru2009-exact-grover,carolan2024-quantum-one-wayness}.  Zalka's
result alone does not establish the $k$-marked statement.  Aaronson's
direct-product Theorem~4.6 instead concerns the find-all-$k$ analogue
\cite{aaronson2005-limitations-quantum-advice}.
\end{proof}

\begin{corollary}[Full conversion curve]
\label{cor:sc-s2}
Let $f=\operatorname{PARITY}_N$, $p\ge k/G$, $1-p=\Omega(1)$,
$4\eta_{\rm int}\le\varepsilon<p$, and $\delta=p-\varepsilon\ge Ck/G$ for a
sufficiently large constant $C$.  Then
\begin{equation}
 \boxed{Q_\varepsilon=\Thetatilde\!\left(
 N\sqrt{(p-\varepsilon)G/k}\right).}
 \label{eq:sc-full-curve}
\end{equation}
The tilde hides only fixed-point precision logarithms: the hidden factor is
$\log(e/\min\{\delta,\varepsilon\})=\log(1/\varepsilon)+O(1)$ when
$\varepsilon\ll\delta$, and can be unbounded relative to the main term for
superpolynomially small $\varepsilon$.  The lower bound is nonvacuous once
$\delta G/k$ is a sufficiently large constant.
\end{corollary}

\begin{proof}
One preparation followed by label measurement and exact parity verification
uses $T+N$ queries.  It accepts with probability zero on a zero-winner input
and at least $p-\varepsilon=\delta$ on a $k$-winner input.  Applying
\cref{thm:sc-small-success} and absorbing its additive $O(N)$ term under the
domain assumption gives the lower bound.  The upper bound is
\cref{thm:sc-m2}; its additive one is absorbed on the same domain.
\end{proof}

Unlike \cref{thm:sc-m1}, this conclusion needs neither
$k=O(G/\log^2G)$ nor $G/k=\widetilde\Omega(\delta^{-2})$.  It is
parity-specific: the gain comes from \cref{lem:sc-parity-collapse}.

\begin{corollary}[Mixed-output preparation profile]
\label{cor:sc-s3}
Suppose a $T$-query circuit, on every exactly-$k$ parity input, outputs a
state within trace distance $\varepsilon\le p/2$ of the winner-mass-$p$
target.  Then
\begin{equation}
 T=\Omega\!\left(\sqrt p\,N\sqrt{G/k}\right)-O(N).
 \label{eq:sc-mixed-profile}
\end{equation}
Thus the first two lines of the preparation profile \eqref{eq:cond-ap6}
are tight even for trace-close mixed outputs: the lower bounds are
$\Omega(N\sqrt{G/k})$ in the subcritical phase and
$\Omega(N(G/k)^{1/4})$ at criticality.  In the supercritical phase the
displayed lower bound is vacuous, so it does not show that the $O(N)$ upper
bound is tight.
\end{corollary}

\begin{proof}
Run the circuit also on an all-even input, which is outside its preparation
promise but on which the physical circuit remains defined.  Measure a label
and verify its parity exactly.  Acceptance is zero on this input and is at
least $p-\varepsilon\ge p/2$ on every exactly-$k$ input.  Apply
\cref{thm:sc-small-success} with $\delta=p/2$.  The three phase statements
follow by substituting the winner masses from
\cref{thm:cond-multiwinner}: $p=\Theta(1)$,
$p=\Theta(\sqrt{k/G})$, and $p=\Theta(k/G)$.
\end{proof}

\Cref{cor:sc-s3} weakens the hypothesis of the preparation lower bounds in
\cref{sec:condensation-exact-k-preparation}.  The lower bound \eqref{eq:cond-ap5} was obtained by coherent
extraction through \eqref{eq:cond-ap4}, and therefore assumed access to a
preparation unitary and its inverse.  For parity blocks and
$\varepsilon\le p/2$, \cref{cor:sc-s3} replaces this access by the weaker
requirement that the circuit output a trace-close mixed state.  The upper bound \eqref{eq:cond-ap3} and the caveat on the
supercritical phase in \eqref{eq:cond-ap6} are unchanged.  The family
of \cref{sec:wta-central-state} has parity blocks and $k=1$, so this
shows which of its condensed central states reveal the hidden holonomy; the
search-then-crossover approach \eqref{eq:wta-architecture} is unchanged.

Combining \cref{cor:sc-s2,thm:sc-m4} gives, on their common domain,
\begin{center}
\begin{tabular}{@{}lcc@{}}
\toprule
model & query complexity & dependence on $\delta=p-\varepsilon$ \\ \midrule
quantum & $\Thetatilde(N\sqrt{\delta G/k})$ & square root \\
randomized classical & $\Theta(N\delta G/k)$ & linear \\
\bottomrule
\end{tabular}
\end{center}
Thus at every nontrivial accuracy, from $\delta=\Theta(k/G)$ to
constant $\delta$, the classical-to-quantum ratio is
$\widetilde\Theta(\sqrt{\delta G/k})$, as for search.  It is
$\widetilde\Theta(1)$ near the public-output boundary $\delta=\Theta(k/G)$
and $\widetilde\Theta(\sqrt{G/k})$ at constant $\delta$.

\begin{openproblem}[General-inner composition]
\label{open:sc-general-composition}
Does every predicate $f$ satisfy a composed small-success lower bound
\[
 \Omega\!\left(\sqrt\delta\,Q(f)\sqrt{G/k}\right)
\]
for finding a winner with probability $\delta$?  The adversary reduction in
\cref{thm:sc-m1} gives only linear dependence on $\delta$, while the parity
proof uses the fact that every character on a proper nonempty part of a
block averages to zero.
\end{openproblem}

\subsection{Uncoupled holonomy batches}
\label{sec:state-conversion-batch}

The global trace row in \cref{sec:wta-construction} couples all components.
For contrast, take $G$ disjoint copies of the qutrit-holonomy component,
each with its own public trace-one row, and no equality that joins two
components.  The path length and local costs are those of
\cref{eq:wta-parameters,eq:wta-local-costs}.  Consequently the product has
$L=33GN$ constant-size PSD blocks, scalar row sparsity at most three, and
scalar column sparsity at most two.  If
\[
 h_g=\prod_{j=1}^N\sigma_{g,j},\qquad y_g=(1+h_g)/2,
\]
then \eqref{eq:wta-cost-spectra} gives component optima
$v_+=29/10$ and $v_-=14/5$.

\begin{theorem}[No amortization for vector output]
\label{thm:sc-batch-vector}
Returning estimates $\widehat v_1,\ldots,\widehat v_G$ satisfying
$|\widehat v_g-v_{h_g}|\le1/25$ for every $g$, with worst-case joint success
probability at least $2/3$, has quantum raw-query complexity
\[
 \Theta(GN)=\Theta(L).
\]
Moreover, there are constants $a,b>0$ such that every algorithm with at
most $aGN$ queries has worst-case joint success probability at most
$2^{-bG}$.
\end{theorem}

\begin{proof}
Threshold every estimate at $57/20$ to recover all $h_g$.  Their product is
the parity of all $GN$ raw signs, which has quantum query complexity
$\Omega(GN)$; reading all signs is a matching upper bound.  For the stronger
claim, specialize Lee--Roland's strong direct-product theorem to
$G$ copies of $\operatorname{PARITY}_N$
\cite{lee2012-a-strong-direct-product-theorem}.  Its Theorem~1.1 with
success parameter $3/4$ gives
\[
 Q_{1-(3/4)^{G/2}}(f^{(G)})
 \ge\frac{G\ln(9/8)}{8000}Q_{1/4}(f).
\]
Since $Q_{1/4}(\operatorname{PARITY}_N)=\Theta(N)$, with fewer than a
fixed constant times $GN$ queries, simultaneous decoding succeeds with
worst-case probability at most $(3/4)^{G/2}=2^{-\Omega(G)}$.
\end{proof}

Now average the $G$ component objectives into one SDP.  Since the components
share no constraint, its optimal value is
\begin{equation}
 V(\sigma)=\frac1G\sum_{g=1}^Gv_{h_g}
 =\frac{14}{5}+\frac1{10G}\sum_{g=1}^Gy_g.
 \label{eq:sc-batch-value}
\end{equation}

\begin{theorem}[Query complexity of the averaged value]
\label{thm:sc-batch-scalar}
For $0<\varepsilon\le1/100$,
\begin{align}
 Q_\varepsilon(V)&=\Theta\!\left(
 N\min\{G,1/\varepsilon\}\right),
 \label{eq:sc-batch-quantum}\\
 R_\varepsilon(V)&=\Theta\!\left(
 N\min\{G,1/\varepsilon^2\}\right).
 \label{eq:sc-batch-classical}
\end{align}
\end{theorem}

\begin{proof}
An exact coherent query to $y_g$ costs $\Theta(N)$ raw sign queries.
Amplitude estimation estimates its mean to error $10\varepsilon$ using
$O(1/\varepsilon)$ such calls
\cite{brassard2002-quantum-amplitude-amplification-and-estimation}; when
$1/\varepsilon>G$, compute all parities.  This proves the quantum upper
bound.  Classically, sample
$O(\min\{G,1/\varepsilon^2\})$ groups without replacement, compute their
parities, and use the empirical mean, scanning all groups when the minimum
is $G$.

For every $G$, restricting to inputs in which all $G$ groups carry the same
$N$-bit string $x$ makes $V$ change by $1/10>2\varepsilon$ when the parity of
$x$ flips, so an $\varepsilon$-estimator computes
$\operatorname{PARITY}_N$ and $T\ge N/2$; this settles $G<6$, where
$\min\{G,1/\varepsilon\}=O(1)$.  For $G\ge6$, choose Hamming weights
$\ell,\ell'\in[G/4,3G/4]$ separated by
$\Delta=\max\{1,\lceil21\varepsilon G\rceil\}$.  Their values in
\eqref{eq:sc-batch-value} differ by more than $2\varepsilon$, so thresholding
an estimator gives a bounded-error polynomial distinguishing the two
weights.  Nayak--Wu's symmetric approximate-counting lower bound gives
degree
\[
 \Omega\!\left(\frac{\sqrt{m(G-m)}}{\Delta}\right)
 =\Omega(\min\{G,1/\varepsilon\}),
\]
where $m$ is the selected weight farther from $G/2$
\cite{nayak1999-the-quantum-query-complexity-of}.  Applying the already
proved degree division in \cref{lem:sc-parity-collapse} forces
$\lfloor2T/N\rfloor=\Omega(\min\{G,1/\varepsilon\})$.

The same lower bound also follows from perfect negative-adversary
composition: compose the partial symmetric outer gap-counting problem with
$G$ independent copies of $\operatorname{PARITY}_N$ and combine Nayak--Wu's outer bound with the
Belovs--Lee composition theorem
\cite{belovs2020-the-quantum-query-complexity-of}.  Thus the parity-collapse
argument here is a construction-specific proof, not a new composition
theorem.

For the classical bound, we show that even an adaptive algorithm must read
every sign of a group to learn its parity.  Fix $h\in\{\pm1\}^G$ and draw the raw
signs of each group uniformly conditional on their product being $h_g$.
While at least two signs of a group remain unseen, a simulator answers
each new query to that group with an independent fair sign.  Only when the
final unseen sign in a group is queried does the simulator query $h_g$ and
force the required product.
This exactly reproduces the conditional distribution under arbitrary
interleaving.  Each parity-oracle query consumes $N$ distinct raw queries,
so a depth-$q$ raw tree induces a parity-bit tree of depth at most
$\lfloor q/N\rfloor$.  The classical bounded-error complexity of estimating
a $G$-bit mean to error $10\varepsilon$ is
$\Theta(\min\{G,1/\varepsilon^2\})$, by the standard Yao distribution on
two nearby central Hamming weights.  Multiplication by $N$ proves
\eqref{eq:sc-batch-classical}.
\end{proof}

The two bounds give three regimes:
\begin{center}
\begin{tabular}{@{}lcc@{}}
\toprule
accuracy & $Q_\varepsilon(V)$ & $R_\varepsilon(V)$ \\ \midrule
$\varepsilon=\Omega(G^{-1/2})$
 & $\Theta(N/\varepsilon)$ & $\Theta(N/\varepsilon^2)$ \\
$G^{-1}=O(\varepsilon)=O(G^{-1/2})$
 & $\Theta(N/\varepsilon)$ & $\Theta(NG)$ \\
$\varepsilon=O(G^{-1})$
 & $\Theta(NG)$ & $\Theta(NG)$ \\
\bottomrule
\end{tabular}
\end{center}
At $\varepsilon=\Theta(G^{-1/2})$, the costs are
$\Theta(N\sqrt G)$ and $\Theta(NG)$, exactly the quantum and classical
scales of the coupled winner-take-all value bound
\eqref{eq:wta-value-complexity}.  The quantum advantage therefore depends on how the accuracy scales with
$G$; it is not an exponential advantage.

\begin{corollary}[Classical output of all Newton directions]
\label{cor:sc-parallel-newton}
At the public point $X_{g,i}=I/3$ and path parameter
$\mu=1/(GP)$, the uncoupled batch has reduced root Hessian $(9/G)I$ and
condition number one.  Nevertheless, producing classical descriptions of
all $G$ root Newton directions to entrywise error less than $1/200$ costs
$\Theta(GN)$ queries, and query budgets below a suitable constant times $GN$ have
$2^{-\Omega(G)}$ worst-case joint success probability.
\end{corollary}

\begin{proof}
Using $A_h$ from \eqref{eq:wta-Ah}, direct elimination gives
\[
 \Delta Z_g=-\frac1{90}(H_{12}+H_{23}+h_gH_{13}).
\]
Every full physical step remains positive definite, with minimum eigenvalue
at least $14/45$.  The hidden $13$ entry has magnitude $1/90>1/200$, so its
sign reveals $h_g$.  Apply \cref{thm:sc-batch-vector}.
\end{proof}

These batch results concern raw queries to the signed copy rows.  An
oracle that supplies group parities, aggregate gradients, or component
optima has already done the hard part: combining the raw signs into
parities or aggregates of them.
\Cref{thm:sc-batch-vector} concerns writing out all components classically,
not preparing one normalized quantum state of the global direction.  Such a
state gives each group amplitude $G^{-1/2}$, so a preparation with constant
error need not reveal any single entry, as \cref{rem:wta-dilution} shows.  Finally, the claims assume exact
equalities; the path operator has small singular values, so they imply
nothing about solutions with constant residuals.

\subsection{Open questions}
\label{sec:state-conversion-open}

Three questions remain open.  First, the decoder based on winner mass does
not detect the internal-state discrepancy $\eta_{\rm int}$.  The minimax
radius and query complexity of the conversion problem when $\eta_{\rm int}$
matters, especially when $p$ is below the condensation threshold, are
open.  Second,
\cref{open:sc-general-composition} asks for the $\sqrt\delta$ composition
bound beyond parity.  Third, our tight bounds exclude the range
$\delta=o(k/G)$: there, public outputs with uniform label can already meet
some accuracy requirements, and the exact boundary depends on the internal
states, not only on the winner mass.

\section{LP Loading and Recovery Lower Bounds}
\label{sec:lp-hardness}

This section gives the LP thread of the paper.  Its endpoint is a
linear-size, bounded-degree standard-form LP whose reduced primal Newton
matrix is a scalar identity at a public infeasible start, whose reduced
log-barrier Hessian is a scalar identity along the entire central path, and
whose original-primal Newton-direction state nevertheless needs a linear
number of raw coefficient queries.  The obstruction is not the scalar solve:
it is construction of the hidden affine translation and recovery in the
original coordinates.

All state, access, residual, and success conventions are those of
\cref{sec:prelim-quantum-access,sec:prelim-output-contracts}.  Throughout
this section \(N\geq2\).  In particular, a residual without another
qualifier is the global RHS-relative residual
\(\norm{Ax-b}_2/\norm b_2\), and an output state is unconditional.  The same
results hold for a heralded branch of input-independent constant success
probability when repetitions are charged.  A reduced-coordinate state is not
an original-coordinate output, and an unobserved ``good branch'' is not an
output contract.  We count all raw queries used for setup, preconditioning,
right-hand-side formation, solution, and recovery.

The underlying amplification motifs are established.  Parity and the
polynomial method are classical query-complexity tools
\cite{beals2001-quantum-lower-bounds-by-polynomials}; endpoint padding and
biased clocks occur in HHL and Feynman--Kitaev constructions
\cite{harrow2009-quantum-algorithm-linear-systems-equations,caha2018-feynman-kitaev-clock,bausch2018-analysis-limitations-modified-circuit};
and sparse QLS states carrying an endpoint bit occur in Wang--Zhang and Mori
et al.\ \cite{wang2024-tight-quantum-depth-lower-bound,mori2026-sparsity-dependent-complexity-lower-bound}.
Apers--Gribling Theorem~8.4 already gives a linear exponent for sparse-LP
\emph{value} estimation in balanced dimensions
\cite{apers2026-quantum-speedups-linear-programming-interior}.  None of these
results is a constant-condition QLS lower bound of the kind sometimes
informally inferred from the statements below.  Our contribution is the
LP-native conjunction of regularity, reduced geometry, original-coordinate
output contracts, sharp residual thresholds, and complete access accounting.

\subsection{Parity amplification and feasible-state hardness}
\label{sec:lp-foundation}

Let \(\sigma_1,\ldots,\sigma_N\in\{\pm1\}\), and put
\[
 p_0=1,\qquad p_i=\prod_{j=1}^i\sigma_j.
\]
Represent a signed scalar at node \(j\) by a nonnegative pair
\((u_j,v_j)\), with \(d_j=u_j-v_j\) and \(q_j=u_j+v_j\).  A chain fixes
\begin{equation}
 d_0=1,\qquad d_i-\sigma_i d_{i-1}=0\quad(1\leq i\leq N).
\end{equation}
At node \(N\), attach a complete binary copy tree with
\begin{equation}
 M=2^{\lceil\log_2(8(N+1))\rceil},\quad K=2M-2,
 \quad P=N+1+K,
\end{equation}
and impose \(d_j-d_{\pi(j)}=0\) on every new node.  Thus all \(K\) copy
nodes have difference \(p_N\), while row and column sparsity remain at most
four and three.

We use the following elementary decoder throughout the section.

\begin{lemma}[Fixed-observable parity decoder]
\label{lem:lp-decoder}
Let \(O\) be an input-independent Hermitian observable with \(\norm O_2\leq1\).
If a family of normalized ideal states obeys
\[
 p_N\langle\psi_\sigma|O|\psi_\sigma\rangle\geq\delta>0,
\]
then the POVM \((I\pm O)/2\) guesses \(p_N\) with bias at least
\(\delta/2\).  From any \(\rho_\sigma\) with
\(\Dtr(\rho_\sigma,\ket{\psi_\sigma}\!\bra{\psi_\sigma})\leq\epsilon\),
the bias is at least \((\delta-2\epsilon)/2\).  If this is positive by an
absolute constant, preparing \(\rho_\sigma\) requires \(\Omega(N)\) raw
coefficient queries.
\end{lemma}

\begin{proof}
Trace duality changes the expectation of \(O\) by at most \(2\epsilon\).
The stated POVM has signed expectation \(\operatorname{Tr}(O\rho)\).  Constant
repetition turns any constant positive bias into bounded error.  Every sparse
row, column, position, or value query in the families below exposes at most one
\(\sigma_i\), even coherently, and is simulated by one call to the sign oracle.
The parity lower bound completes the reduction.
\end{proof}

First minimize \(\sum_jq_j\) subject to the chain and copy equations.  Exact
feasibility fixes every \(d_j\in\{\pm1\}\), and \(q_j\geq1\).  Hence the
unique optimum has \(q_j=1\) and value \(P\).  The nullspace has the public
orthonormal basis
\begin{equation}
 V_j=\frac{e_{u_j}+e_{v_j}}{\sqrt2}.
\end{equation}
On the feasible affine space the barrier objective is
\begin{equation}
 \Phi_\mu(q)=\sum_j\left[q_j-
       \mu\log\frac{q_j^2-1}{4}\right].
\end{equation}
Thus every central coordinate is
\(q(\mu)=\mu+\sqrt{\mu^2+1}\), and
\begin{equation}
 V^T\nabla^2\Phi_\mu V
 =2\mu\bigl[(q+1)^{-2}+(q-1)^{-2}\bigr]I.
\end{equation}
At \(\mu=3/4\), \(q=2\), so every oriented pair is
\((3/2,1/2)\).  If instead an exactly feasible point has objective gap at
most one, writing \(q_j=1+\delta_j\) gives
\(\sum_j\delta_j\leq1\) and
\begin{equation}
 \norm x_2^2=\sum_j\frac{q_j^2+1}{2}\leq P+\frac32.
\end{equation}
Since \(K\geq16(N+1)-2\), more than \(9/10\) of this mass lies in the
copy tree, and within an oriented copy pair the correct coordinate has at
least \(9/10\) of the pair mass.  The computational-basis decoder therefore
succeeds with probability greater than \(43/50\); trace error \(1/10\)
leaves success greater than \(19/25\).  This proves an \(\Omega(N)\) lower
bound both for every such near-optimal primal state and for the exact
\(\mu=3/4\) central state, despite \(\kappa_{\rm red}=1\).

There is a stronger any-feasible-point contract.  Add \(t_j\geq0\) and
\begin{equation}
 2u_j+2v_j+2t_j=3
\end{equation}
at every node, and retain objective \(\sum_jq_j\).  Exact feasibility now
gives
\begin{equation}
 1\leq q_j\leq\frac32,\qquad t_j=\frac32-q_j.
\end{equation}
The equality matrix has rank \(2P\), and its public orthonormal null basis is
\begin{equation}
 \bar V_j=\frac{e_{u_j}+e_{v_j}-2e_{t_j}}{\sqrt6}.
\end{equation}
The positive optimal columns form a basis; the opposite pair coordinate has
dual slack two.  Thus the LP is bounded and strictly primal--dual feasible,
with a unique nondegenerate strictly complementary optimum.  Its reduced
central Hessian is again a scalar identity, because the reduced barrier is a
sum of identical functions of \(q_j\).

For \(q\in[1,3/2]\), one triple has squared norm
\begin{equation}
 \bar w(q)=\frac{q^2+1}{2}+\left(\frac32-q\right)^2
 \in\left[\frac54,\frac{13}{8}\right].
\end{equation}
The copy triples therefore carry more than \(9/10\) of the total mass, and
the plus/minus decoder, randomized on \(t\), succeeds conditionally with
probability at least \(9/10\).  Consequently \emph{every} exact feasible
amplitude state, without an objective promise, has the same
\(\Omega(N)\) query lower bound at trace error \(1/10\).  The proof also
survives \(x\geq0\) and the absolute global \(\ell_1\) residual
\(\norm{\bar Ax-\bar b}_1\leq1/50\): telescoping gives
\(p_Nd_j\geq49/50\) at every copy node, while the cap gives
\(q_j+t_j\leq151/100\).  The resulting copy mass is greater than \(43/50\),
and the signed conditional expectation exceeds \(63/100\).  Before trace
error the total signed expectation is therefore
\begin{equation}
 \frac{43}{50}\frac{63}{100}=\frac{2709}{5000},
\end{equation}
so trace error \(1/10\) leaves signed expectation at least
\[
 \frac{2709}{5000}-\frac15=\frac{1709}{5000}>\frac13.
\]

\subsection{The robustness frontier}
\label{sec:lp-robustness-frontier}

The preceding \(\ell_1\) theorem cannot be promoted to a constant
\(\ell_2\) theorem on one unit-height path.  In gauged variables set
\begin{equation}
 p_id_i=1-\frac{i}{N},\qquad d_j=0\quad(j\text{ in the copy tree}),
\end{equation}
and take \(q_j=1,t_j=1/2\).  All cap and copy equations are exact, the
\(N\) chain residuals equal \(1/N\), and the total residual is
\(N^{-1/2}\), yet the amplified coordinates have no parity signal.  Using
\(1-2i/N\) instead gives an entirely wrong copy plateau at residual
\(2/\sqrt N\).  This is the basic drift obstruction.

A bounded-height repair, for power-of-two \(N\), uses \(R=N\) length-\(N\)
signed paths in parallel.
A binary fanout tree feeds the paths, and a binary tree with \(L=16N\)
leaves is attached to every endpoint.  The resulting tree has
\begin{equation}
 P=(2N-1)+N^2+N(2L-2)=33N^2-1,
 \qquad |S|=16N^2
\end{equation}
vertices and designated output leaves \(S\).  At every vertex use
\((u,v,h,t)\geq0\), impose pinned labelled-copy equations on
\(d=u-v\), pinned unsigned-copy equations on \(h\), and
\(q+t-2h=0\).  Minimize \(\sum_j(2u_j+2v_j+h_j+t_j)\).

\begin{theorem}[Parallel-path residual tube]
\label{thm:lp-parallel-path}
This LP has \(4P\) variables, \(3P\) full-rank equalities, coefficients of
magnitude at most two, and row and column sparsity at most four.  It has the
same regularity properties as the capped family and a scalar reduced central
Hessian for every \(\mu>0\).  If \(x\geq0\), \(x\ne0\), and
\begin{equation}
 \frac{\norm{Ax-b}_2}{\norm b_2}\leq\frac1{200},
 \qquad \norm b_2=\sqrt2,
\end{equation}
then the output-coordinate decoder has ideal bias greater than \(1/80\).
Consequently, preparing an amplitude state within trace distance \(1/100\)
of any such point needs \(\Omega(N)=\Omega(\sqrt P)\) raw queries.  Any
extra two-sided objective or centrality promise is immaterial.
\end{theorem}

\begin{proof}
Let \(\lambda_j\) be the exact propagated tree label.  Gauging removes all signs.
Descendant counting in each binary tree and Cauchy--Schwarz on the \(N\)
parallel paths give
\begin{equation}
 \frac{\norm{\operatorname{Diag}(\lambda)z-\mathbf1}_2}{\sqrt P}
 \leq12\norm{(z_r-1,Dz)}_2,
\quad
 \frac{\norm{\operatorname{Diag}(\lambda_S)z_S-\mathbf1}_2}{\sqrt{|S|}}
 \leq12\norm{(z_r-1,Dz)}_2.
\label{eq:lp-tag-9-15}
\end{equation}
The binary-level operator norms sum to
\(\sum_{\ell\geq1}2^{-\ell/2}=1+\sqrt2<5/2\), while the path term is
\(\sqrt{N/R}=1\).  With \(E=\norm{Ax-b}_2\leq\sqrt2/200\),
\eqref{eq:lp-tag-9-15} for
\(d\) and \(h\), plus the cap and nonnegativity, gives
\begin{equation}
 \norm x_2^2<11P,
 \qquad
 \sum_{j\in S}\lambda_jd_jq_j>\frac35|S|.
\end{equation}
The two loss terms in the second inequality, divided by \(|S|\), are at most
\[
 12E\left(1+2.18\sqrt{P/|S|}\right)
 <\frac{12\sqrt2}{200}\left(1+2.18\sqrt{33/16}\right)<\frac25.
\]
Hence the measurement bias is larger than
\[
 \frac{(3/5)16}{2\cdot11\cdot33}=\frac8{605}>\frac1{80}.
\]
Trace error \(1/100\) leaves positive constant bias, and
Lemma~\ref{lem:lp-decoder} applies.  Full rank and regularity follow by
pivoting successively on the pinned \(d\), pinned \(h\), and private
\(t\) blocks.  The public null vectors are the local vectors in
\eqref{eq:lp-gain-null}; identical feasible intervals make the central
curvature scalar.
\end{proof}

The quadratic volume is necessary in the signed-edge gadget class, not in all
bounded-coefficient real-linear gadgets.

\begin{proposition}[Effective-resistance volume bound]
\label{prop:lp-resistance-volume}
Let a connected graph have maximum degree \(\Delta\), a pinned root, and
outputs \(S\).  Suppose every consistent edge constraint
\(d_v-a_{uv}(\sigma)d_u=0\) has \(a_{uv}\in\{\pm1\}\) depending on at most
\(k\) input bits, and every output label is parity.  Then there is a vector
with root value one and wrong-parity value at every output whose residual is
at most
\begin{equation}
 \frac{k\sqrt{2\Delta|V|}}{N}.
\label{eq:lp-tag-9-17}
\end{equation}
Thus excluding all wrong-output vectors below an absolute residual \(\eta\)
requires \(|V|>\eta^2N^2/(2\Delta k^2)\).
\end{proposition}

\begin{proof}
The product of edge labels on a root-to-output path is parity.  Fourier degree
is subadditive under products, so every such path has length at least \(N/k\).
Gauge the consistent labels away and consider a unit root-to-\(S\) flow.
Removing circulations and decomposing into paths gives
\(\sum_e|f_e|\geq N/k\).  Cauchy--Schwarz and Thomson's principle yield
\[
 R_{\rm eff}(r,S)\geq\frac{N^2}{k^2|E|},\qquad
 C_{\rm eff}(r,S)\leq\frac{\Delta k^2|V|}{2N^2}.
\]
The Dirichlet energy for voltage difference two is \(4C_{\rm eff}\); its
harmonic minimizer lies in \([-1,1]\).  Taking the square root and undoing
the gauge proves \eqref{eq:lp-tag-9-17}.  The next section develops the general mixture and
cut frontiers; only this signed-edge specialization is needed here.
\end{proof}

There is a load-bearing escape witness.  On one path replace the signed
recurrence by
\(d_i-2\sigma_i d_{i-1}=0\), add a public reference
\(h_i-2h_{i-1}=0\), and retain \(q_i+t_i-2h_i=0\).  The exact scale is
\(h_i=|d_i|=2^i\).  With absolute residual
\(E=\norm{Ax-b}_2\), rescaling by \(2^{-i}\) gives
\begin{equation}
 \left|2^{-i}p_id_i-1\right|
 \leq\left(1+\frac1{\sqrt3}\right)E.
\end{equation}
Put \(K_0=1+1/\sqrt3\), and assume \(E\leq10^{-2}\).  Applying the same
estimate to the reference rows gives
\[
 p_id_i\geq(1-K_0E)2^i,\qquad
 0\leq h_i\leq(1+K_0E)2^i.
\]
If \(c_i=q_i+t_i-2h_i\) is the cap residual, nonnegativity gives
\(q_i,t_i\leq2h_i+|c_i|\).  Therefore
\begin{equation}
 \norm x_2^2
 \leq\sum_i\bigl[(q_i+t_i)^2+h_i^2\bigr]
 \leq12(1+K_0E)^2\,4^N+2E^2,
\qquad
 x_{{\rm correct},N}\geq(1-K_0E)2^N.
\end{equation}
The endpoint consequently has constant normalized mass (and ideal decoder
bias greater than \(3/100\) at \(E=10^{-2}\)), giving a linear-size robust
family.  It has exponential range and poor conditioning, showing that bounded
local coefficients alone do not imply the quadratic volume law.  The next
construction replaces the long gain chain by a short gain prefix and a
polynomial plateau.

\subsection{The linear-size gain--plateau family}
\label{sec:lp-gain-plateau}

Set
\begin{equation}
 T=\left\lceil\frac12\log_2N\right\rceil,
 \quad K=16N,\quad P=17N+1,
 \quad H_i=2^{\min\{i,T\}},\quad H=2^T\in[\sqrt N,2\sqrt N).
\end{equation}
There are nodes \(i=0,\ldots,N+K\), gains
\(g_i=H_i/H_{i-1}\in\{1,2\}\), and labels \(a_i=\sigma_i\) for
\(i\leq N\) and \(a_i=1\) afterward.  Define
\[
 \tau_0=1,\qquad \tau_i=\prod_{j=1}^i a_j,\qquad
 R_\tau=\diag(\tau_0,\ldots,\tau_{P-1}).
\]
At every node introduce
\((u_i,v_i,h_i,t_i)\geq0\), put \(d_i=u_i-v_i\), \(q_i=u_i+v_i\), and
impose
\begin{equation}
\begin{aligned}
 d_0&=1,&d_i-g_ia_id_{i-1}&=0,\\
 h_0&=1,&h_i-g_ih_{i-1}&=0,\\
 &&q_i+t_i-2h_i&=0.
\end{aligned}
\label{eq:lp-tag-9-20}
\end{equation}
The base objective is \(\sum_i(2u_i+2v_i+h_i+t_i)\).  Its public
orthonormal null basis is
\begin{equation}
 W_i=\sqrt{\frac23}\left(\frac12e_{u_i}+\frac12e_{v_i}-e_{t_i}\right).
 \label{eq:lp-gain-null}
\end{equation}

\begin{theorem}[Linear robust primal-state lower bound]
\label{thm:lp-gain-plateau}
The LP \eqref{eq:lp-tag-9-20} has \(4P\) variables, \(3P\) full-rank equalities,
coefficients of magnitude at most two, and row and column sparsity at most
four and three.  It is bounded and strictly primal--dual feasible, with a
unique nondegenerate strictly complementary optimum.  The optimum is
\(\Theta(N^{3/2})\), its largest coordinate is \(\Theta(\sqrt N)\), and
its Euclidean norm is \(\Theta(N)\).

For every nonzero \(x\geq0\) with
\begin{equation}
 \operatorname{res}_p(x)\leq\frac1{100},
\label{eq:lp-tag-9-21}
\end{equation}
the fixed plateau decoder has ideal bias greater than \(1/30\).  Preparing
its original-primal state to trace error \(1/100\) therefore needs
\(\Omega(N)=\Omega(P)\) raw queries.  Moreover,
\begin{equation}
 \kappa_{\rm red}(\mu)<\frac76\qquad(0<\mu\leq1/16).
\label{eq:lp-tag-9-22}
\end{equation}
\end{theorem}

\begin{proof}
The pinned signed and unsigned bidiagonal blocks fix
\(d_i=H_i\tau_i\) and \(h_i=H_i\); private \(t_i\) columns finish the rank
proof.  Exact feasibility is
\begin{equation}
 H_i\leq q_i\leq2H_i,\qquad t_i=2H_i-q_i.
\end{equation}
The local objective is \(q_i+3H_i\), so the unique optimum has
\(q_i=t_i=H_i\).  The sign-selected pair member together with \(h_i,t_i\)
forms a triangular positive basis, and the opposite coordinate has slack two.
This proves all regularity claims.  Since there are \(\Theta(N)\) nodes of
height \(H=\Theta(\sqrt N)\), the three scale estimates follow.

Write a central coordinate as \(q_i=H_i+z_i\).  Stationarity and curvature
are
\begin{equation}
 \frac1\mu=\frac1{2H_i+z_i}+\frac1{z_i}-\frac1{H_i-z_i},
\quad
 f_i''=\mu\left[(2H_i+z_i)^{-2}+z_i^{-2}+(H_i-z_i)^{-2}\right].
\end{equation}
For \(\mu\leq1/16\),
\(15\mu/16\leq z_i<\mu\), whence
\(1/\mu<f_i''<7/(6\mu)\).  Changing to the orthonormal coordinates
\eqref{eq:lp-gain-null} multiplies every eigenvalue by \(2/3\), proving
\eqref{eq:lp-tag-9-22}.

For robustness put \(E=\norm{Ax-b}_2\leq\sqrt2/100\).  The weighted
resistance is
\[
 \sum_{i\geq1}H_i^{-2}<\frac13+\frac{17N}{H^2}\leq\frac{52}{3}.
\]
Weighted telescoping gives, simultaneously at every node,
\begin{equation}
 \left|\frac{\tau_id_i}{H_i}-1\right|<\frac3{40},
 \qquad \left|\frac{h_i}{H_i}-1\right|<\frac3{40}.
\end{equation}
The cap then gives \(q_i,t_i<2.165H_i\), and
\begin{equation}
 \norm x_2^2<11\sum_iH_i^2
 <\frac{583}{3}NH^2.
\end{equation}
On the \(K=16N\) output nodes,
\(p_Nd_iq_i>(37H/40)^2\).  The computational-basis decoder has bias
\begin{equation}
 \frac{16N(37/40)^2H^2}{2(583/3)NH^2}>\frac1{30}.
\end{equation}
Trace error \(1/100\) leaves bias greater than \(7/300\), and
Lemma~\ref{lem:lp-decoder} completes the proof.
\end{proof}

The early-path conditioning claim cannot be strengthened: for \(\mu/H\to
\infty\), local curvature scales as \(\mu/H_i^2\), so the unrigidified
condition ratio can reach \(\Theta(H^2)=\Theta(N)\).  The gain is also the
precise escape from the quadratic path-bundle law.  If \(R\) internally
disjoint length-\(N\) paths have reference height at most \(H\), a harmonic
wrong-endpoint interpolation has squared residual at most \(4RH^2/N\).
Therefore exclusion below absolute residual \(\eta\) requires
\begin{equation}
 RH^2>\frac{\eta^2N}{4},\qquad
 P=\Theta(RN)\Longrightarrow PH^2=\Omega(N^2).
\label{eq:lp-tag-9-28}
\end{equation}
Both the bounded-height parallel construction and the linear-volume
gain--plateau construction attain this frontier.  The general gadget-class
mixture and cut theorems are given in \cref{sec:gadget-frontiers};
\eqref{eq:lp-tag-9-28} is
the only
path-bundle statement used here.  Naive balanced-XOR/PCP encodings do not
improve it: boundary truth-table points obstruct strict feasibility, and
homogenized shallow circuits still admit fractional drift across a small cut.

For one static input, Theorem~\ref{thm:lp-gain-plateau} gives the tight total
ledger
\begin{equation}
 Q_{\rm prep}+\sum_{t=1}^R
 (Q_{{\rm form},t}+Q_{{\rm solve},t}+Q_{{\rm recover},t})+Q_{\rm out}
 =\Omega(P).
\label{eq:lp-tag-9-29}
\end{equation}
The bound is total, not per iteration.  Querying all \(N\) signs and caching
their prefixes costs \(O(P)\) and makes every later raw query free, so no
\(\Omega(RP)\) fixed-instance theorem is possible.

\subsection{Reduced geometry versus the affine offset}
\label{sec:lp-affine-separation}

All reduced geometry in the gain--plateau family is public.  Every exact
feasible point has the form
\begin{equation}
 x_\sigma(q)_i=
 \left(\frac{q_i+H_i\tau_i}{2},
       \frac{q_i-H_i\tau_i}{2},H_i,2H_i-q_i\right),
\label{eq:lp-tag-9-30}
\end{equation}
and
\begin{equation}
 \{x:A_\sigma x=b\}=x_\sigma^{\rm off}+\operatorname{range}(W),
\quad
 (x_\sigma^{\rm off})_i
 =H_i\left(\frac{1+\tau_i}{2},\frac{1-\tau_i}{2},1,1\right).
\end{equation}
Even the minimum-norm offset is hidden: minimizing the local squared norm in
\eqref{eq:lp-tag-9-30} gives \(A_\sigma^\dagger b=x_\sigma(4H_i/3)\).  By contrast, the
full reduced barrier, all derivatives, central scalars, \(W\), both
projectors, and the diagonal reduced Hessian are independent of \(\sigma\).
For \(\mu\leq1/16\), its eigenvalues lie in
\((2/(3\mu^2),7/(9\mu^2))\) in the unscaled log-barrier convention of
\cref{def:reduced-kappa}; the \(\mu\)-weighted Hessian
\(W^\top\nabla^2\Phi_\mu W\) has eigenvalues in
\((2/(3\mu),7/(9\mu))\).

\begin{corollary}[Affine-offset oracle dichotomy]
\label{cor:lp-affine-dichotomy}
Grant arbitrary exact input-independent reduced-geometry oracles for free.
For the gain--plateau family, loading any original-primal state satisfying
\eqref{eq:lp-tag-9-21} costs \(\Omega(P)\), while the free reduced Hessian has condition
below \(7/6\) on the late tail.  For the parallel-path family, the analogous
cost is \(\Omega(\sqrt P)\), while the reduced Hessian is an exact scalar
identity for the entire central path.  In both cases the hidden object is the
affine offset, not the reduced solution.
\end{corollary}

\begin{proof}
Every granted channel is the same for all \(\sigma\) and can be hardwired in
the parity reduction.  Theorem~\ref{thm:lp-gain-plateau} and
Theorem~\ref{thm:lp-parallel-path} then apply unchanged.  A free feasible
offset, prefix table, or loading isometry would not be input-independent: one
known output pair reveals \(p_N\), and a normalized offset state carries a
constant amount of output-pair mass.
\end{proof}

\subsection{Exact Newton secants}
\label{sec:lp-newton-secants}

The base objective in the gain--plateau family gives a uniformly conditioned
late central tail, but its exact public-start construction initially used
different complementarity parameters on the \(O(\log N)\) gain nodes.  A
public reweighting removes that qualification.  At height \(H_i\), replace
the local objective by
\begin{equation}
 (1+\omega_i)(u_i+v_i)+h_i+t_i,
 \qquad \omega_i=\frac7{36H_i}.
\end{equation}
On the feasible line its nonconstant part is \(\omega_iq_i\), so the unique
optimum and all regularity properties survive; the nonbasic reduced cost is
\(2\omega_i>0\).

Set
\begin{equation}
 \mu_1=\frac1{16},\qquad
 x_i^1=H_i\left(\frac98,\frac18,1,\frac34\right)
\label{eq:lp-tag-9-33}
\end{equation}
up to the hidden pair swap, and put \(s_j^1=\mu_1/x_j^1\).  The reduced
stationarity equation at \(z_i=H_i/4\) is exact because
\begin{equation}
 \frac{\omega_i}{\mu_1}=\frac{28}{9H_i}
 =\frac1{2H_i+H_i/4}+\frac1{H_i/4}-\frac1{H_i-H_i/4}.
\end{equation}
Thus \((x^1,y^1,s^1)\) is the genuine global \(\mu_1\)-central point.

Use the public start
\begin{equation}
\begin{split}
 x_i^0&=H_i(5/6,5/6,20/27,5/9),\\
 s_i^0&=H_i^{-1}(10/27,10/27,5/12,5/9),\qquad y^0=0.
\end{split}
\label{eq:lp-tag-9-35}
\end{equation}
Every coordinate product is
\begin{equation}
 \mu_0=\frac{25}{81},\qquad
 \widehat\mu=\frac{25}{162}=\frac{\mu_0}{2}.
\end{equation}
Every Newton right-hand side is public, because the pair differences at the
start vanish.  For either hidden pair orientation,
\begin{equation}
 \frac{10}{27H_i}\frac{9H_i}{8}
 +\frac{5H_i}{6}\frac1{18H_i}
 =\frac{10}{27H_i}\frac{H_i}{8}
 +\frac{5H_i}{6}\frac1{2H_i}
 =\frac{25}{54}.
\label{eq:lp-tag-9-37}
\end{equation}
The \(h,t\) coordinates give the same \(25/54\).  Hence
\begin{equation}
 S^0(x^1-x^0)+X^0(s^1-s^0)
 =\left(\frac{25}{54}-2\frac{25}{81}\right)e
 =(\widehat\mu-\mu_0)e.
\label{eq:lp-tag-9-38}
\end{equation}
Primal and dual endpoint feasibility provide the other two rows of the
standard system \eqref{eq:full-kkt}; full row rank makes the correction
unique.  One full algebraic Newton step therefore reaches the genuine global
\(\mu_1=1/16\) center.  The endpoint parameter differs from
\(\widehat\mu\) because the Newton equation omits
\(\Delta X\Delta s\).

Up to the hidden swap, the primal direction block is
\begin{equation}
 \Delta x_i=H_i(7/24,-17/24,7/27,7/36),
 \qquad \norm{\Delta x_i}_2^2
 =D H_i^2,\quad D=\frac{16139}{23328}.
\label{eq:lp-tag-9-39}
\end{equation}
The unselected-minus-selected squared pair gap is \(5H_i^2/12\) (the
unselected coordinate carries the larger direction mass, \(17H_i/24\)
against \(7H_i/24\)).  Since
\(\sum_iH_i^2<(17N+4/3)H^2\), the fixed output-copy decoder has bias
\begin{equation}
 \frac{16N(5/12)H^2}{2D(17N+4/3)H^2}>\frac14.
\label{eq:lp-tag-9-40}
\end{equation}
Thus the exact original-primal direction state at trace error \(1/100\)
needs \(\Omega(P)\) raw queries.

In the unscaled log-barrier convention of
\cref{def:reduced-kappa}, the \(\mu_1\)-scaled endpoint reduced Hessian and
the start reduced Newton matrix are
\begin{equation}
 \mu_1 W^T\nabla^2F(x^1)W
 =\frac{182}{243}\diag(H_i^{-2}),
\qquad
 W^T(X^0)^{-1}S^0W=\frac{22}{27}\diag(H_i^{-2}).
\end{equation}
Thus all but the \(T=O(\log N)\) public gain modes form an exact scalar
block.  The public congruence \(D_H=\diag(H_i)\) makes both matrices scalar;
as a standard-form diagonal change of variables it raises the coefficient
scale to \(\Theta(\sqrt N)\).  More generally, scaling by \(H_i^\eta\)
gives
\begin{equation}
 B_\eta=\Theta(H^\eta),\quad
 \kappa_\eta=H^{2(1-\eta)},\quad
 B_\eta^2\kappa_\eta=\Theta(N).
\end{equation}
This is a property of this diagonal family, not a universal
normalization--condition theorem.

The remaining outliers are forced inside this local ansatz.

\begin{proposition}[Four-variable-node incompatibility]
\label{prop:lp-local-incompatibility}
Consider hidden-swappable endpoint pairs with values \(L_i>S_i>0\) at one
common central parameter \(\mu_1\), and suppose their differences are the
prescribed heights \(L_i-S_i=H_i\).  Within the separable four-variable
node gadget, the following three requirements are incompatible: a public
start centered coordinatewise at one \(\mu_0\); one full standard Newton
correction with a common target that works for both hidden orientations; and
an \(O(1)\)-conditioned endpoint Hessian in the raw orthonormal null basis.
This statement does not cover added rows, weighted barriers, or other local
gadgets.
\end{proposition}

\begin{proof}
Write the two public start values as \(a_i,b_i\), with slacks
\(\mu_0/a_i,\mu_0/b_i\).  Equality of the first-coordinate secant cross
terms for the two hidden orientations requires
\[
 \frac{\mu_0L_i}{a_i}+\frac{a_i\mu_1}{L_i}
 =\frac{\mu_0S_i}{a_i}+\frac{a_i\mu_1}{S_i}.
\]
Since \(L_i\ne S_i\), this gives
\(a_i^2=\mu_0L_iS_i/\mu_1\); the second coordinate similarly gives
\(b_i=a_i\).  The common cross term is therefore
\[
 C_i=\sqrt{\mu_0\mu_1}\,
       \frac{L_i+S_i}{\sqrt{L_iS_i}}.
\]
One global Newton target makes \(C_i\) independent of \(i\), hence the
aspect ratio \(L_i/S_i\) is constant.  The cap row
\(t_i=2H_i-(L_i+S_i)>0\) makes \(L_i+S_i\) proportional to \(H_i\) and
forces that constant aspect ratio to exceed three.  Together with
\(L_i-S_i=H_i\), this makes both endpoint values proportional to \(H_i\).
The log-barrier curvature in the local orthonormal null direction
consequently scales as \(H_i^{-2}\), giving condition ratio
\(H^2=\Theta(N)\).
\end{proof}

The reweighted construction also exposes an important negative result.  A
fixed relative residual for the full Newton system alone does not force the
primal direction ray.  On the \(K\) output nodes let
\(w=\sum_{i\in S}W_i\), and perturb the exact solution by
\begin{equation}
 e_x=Mw,\qquad e_y=0,\qquad
 e_s=-(X^0)^{-1}S^0e_x.
\end{equation}
The primal and complementarity residual perturbations vanish.  On the
plateau, direct substitution in \eqref{eq:lp-tag-9-35} gives the sharp bound
\[
 \norm{(X^0)^{-1}S^0w}_2\leq\frac{\norm w_2}{H^2}
\]
(the exact coefficient is below \(0.856/H^2\)), whereas the public cap RHS
has norm \(\Omega(H\sqrt K)\).  With \(M=H^2\), the relative full-KKT
residual is \(O(H^{-1})\) and the normalized primal direction converges to
the public zero-query state \(\ket{w/\norm w}\).  The perturbed update has
negative \(t\)-coordinates.  Therefore the robust direction theorem below
must retain the geometric condition \(x^0+\widetilde{\Delta x}\geq0\); a
bare fixed full-KKT residual promise is insufficient.

\subsubsection{Prefix rigidification: the condition-one capstone}
\label{sec:lp-prefix-rigidified}

Return to the constraints \eqref{eq:lp-tag-9-20}.  At each of the \(T\) small-height nodes
\(i<T\), add the public row
\begin{equation}
 q_i-\frac54h_i=0,
\label{eq:lp-tag-9-44}
\end{equation}
and use the uniform objective
\begin{equation}
 \sum_i[(1+\omega)(u_i+v_i)+h_i+t_i],
 \qquad \omega=\frac7{36H}.
\end{equation}
Write the augmented matrix and RHS as \(A'_\sigma,b'\).

\begin{theorem}[Condition-one original-primal Newton hardness]
\label{thm:lp-prefix-capstone}
The prefix-rigidified family has \(4P\) variables and \(3P+T\) full-rank
equalities.  Coefficients have magnitude at most two, and row and column
sparsity are at most four.  It is bounded and strictly primal--dual feasible,
with a unique nondegenerate strictly complementary optimum.  In the raw
orthonormal reduced coordinates,
\begin{equation}
 \kappa_{\rm red}(\mu)=1\quad\text{for every }\mu>0,
\end{equation}
and the reduced Newton matrix at the public start \eqref{eq:lp-tag-9-35} is also exactly
condition one.

The unique standard Newton correction from that start with target
\(\widehat\mu=25/162\) lands after one full step at the genuine global
\(\mu_1=1/16\) center \eqref{eq:lp-tag-9-33}.  Its original-primal direction state has
decoder bias greater than \(1/4\), so trace error \(1/100\) requires
\(\Omega(N)=\Omega(P)\) raw queries.

The same lower bound holds for a state of every nonzero \(x\geq0\) satisfying
\(\operatorname{res}_p(x)\leq1/100\), and for every nonzero direction
\(\widetilde{\Delta x}\) whose update is nonnegative and obeys that residual,
at state trace error \(1/100\).  In the latter case a fixed interference
decoder has bias greater than \(1/54\).
\end{theorem}

\begin{proof}
For the row \(r_i\) in \eqref{eq:lp-tag-9-44} and the old null vectors,
\begin{equation}
 r_i^TW_j=\sqrt{2/3}\,\delta_{ij}.
\end{equation}
Hence the new rows are independent modulo the old row space and
\begin{equation}
 \ker A'_\sigma=\operatorname{span}\{W_i:i\geq T\}.
\label{eq:lp-tag-9-48}
\end{equation}
They fix \(q_i=5H_i/4\) on the prefix, leaving only equal-height \(H\)
coordinates free.  At the optimum all four prefix variables and three
variables at every free node are positive, exactly \(3P+T\) in total.  A
null vector supported on them must vanish because every nonzero free \(W_i\)
touches that node's zero pair coordinate.  The positive columns form a basis,
and each nonbasic reduced cost is \(2\omega>0\).  This proves rank and
regularity.

All surviving nodes have the same height, objective slope, central scalar,
and local curvature.  Equation~\eqref{eq:lp-tag-9-48} therefore gives
\[
 W^T\nabla^2F(x(\mu))W=\lambda(\mu)I_{P-T}.
\]
At \(\mu_1=1/16\), the free coordinate is \(q=5H/4\), while the fixed
prefix coordinates already have the same local template.  Positive slacks
\(s^1=\mu_1(X^1)^{-1}e\) are orthogonal to the surviving kernel after
subtracting the objective, so full row rank supplies the central multiplier.
The secant identities \eqref{eq:lp-tag-9-37}--\eqref{eq:lp-tag-9-38} are unchanged, and the endpoint satisfies
the new rows.  This proves the exact Newton statement.  At the start,
\begin{equation}
 W^T(X^0)^{-1}S^0W=\frac{22}{27H^2}I_{P-T}.
\end{equation}
The fraction-to-boundary maximum is exactly
\(\alpha_{\max}=20/17>1\), so the full step preserves strict positivity;
any safety factor at least \(17/20\) permits it.  The direction formula
\eqref{eq:lp-tag-9-39} survives and \eqref{eq:lp-tag-9-40} proves the exact-state lower bound.

For a robust feasible state, deleting the public new residual coordinates
gives the base residual bound with the same RHS norm
\(\norm{b'}_2=\sqrt2\); Theorem~\ref{thm:lp-gain-plateau} applies.  For a
robust direction, set \(\widetilde x=x^0+\widetilde{\Delta x}\geq0\).  On
every plateau node the same telescoping estimates give
\begin{equation}
 \tau_i\widetilde d_i>\frac{37}{40}H,\quad
 \frac{37}{40}H<\widetilde h_i<\frac{43}{40}H,\quad
 \widetilde q_i,\widetilde t_i<\frac{433}{200}H.
\label{eq:lp-tag-9-50}
\end{equation}
Together with the public start, these imply
\begin{equation}
 \norm{\widetilde{\Delta x}}_2^2<100NH^2,\qquad
 \widetilde{\Delta h}_i>\frac{199}{1080}H
 \quad(i\in S).
\label{eq:lp-tag-9-51}
\end{equation}
Let \(\ket{g_i}=(\ket{u_i}-\ket{v_i})/\sqrt2\), and use the public norm-one
observable
\begin{equation}
 O=\bigoplus_{i\in S}
 (\ket{h_i}\!\bra{g_i}+\ket{g_i}\!\bra{h_i}).
\end{equation}
Equations \eqref{eq:lp-tag-9-50}--\eqref{eq:lp-tag-9-51} give signed expectation greater than
\begin{equation}
 \frac{7363\sqrt2}{270000}>\frac1{27},
\end{equation}
so its associated binary POVM has bias greater than \(1/54\).  Trace error
\(1/100\) leaves positive constant bias.  Lemma~\ref{lem:lp-decoder}
finishes the proof.
\end{proof}

The condition-one claim is intentionally reduced-primal.  At the same start,
the raw OSS matrix \({\cal M}_0=[-X^0(A')^T\ \ S^0W]\) and the eliminated
augmented KKT matrix satisfy
\begin{equation}
 \kappa_2({\cal M}_0)\geq\sqrt{681/44}\,H^2,
 \qquad
 \kappa_2({\cal K}_0)\geq\sqrt{1701/178}\,H^2,
\end{equation}
so both are \(\Omega(P)\).  Exact block preconditioning can reduce the
second condition to \(\varphi^2\), but its difference-row Schur block has
the gauge form
\begin{equation}
 G_\sigma^{(d)}=R_\tau G_+^{(d)}R_\tau,\qquad
 (G_\sigma^{(d)})^{-1/2}
 =R_\tau(G_+^{(d)})^{-1/2}R_\tau.
\end{equation}
The all-positive block is an irreducible tridiagonal Stieltjes matrix, whose
inverse square root is entrywise strictly positive by its resolvent integral.
Thus the sign of a sign-resolving preconditioner entry reveals a prefix
product.  A factor block encoding without entry access requires a separate
implementation analysis; \cref{sec:preconditioning} treats that
factor-access frontier.

The access localization can be exact.  In the Moore--Penrose gauge write
\begin{equation}
 \Delta x=p_\sigma+Wz,\qquad
 p_\sigma=(A'_\sigma)^\dagger(b'-A'_\sigma x^0).
\end{equation}
On every free node,
\begin{equation}
 (p_d,p_q,p_h,p_t)=H(\tau_i,-4/27,7/27,-2/27),
\end{equation}
with pair coordinates recovered from \(p_d,p_q\).  The complete reduced
system is public:
\begin{equation}
 M=\frac{22}{27H^2}I,\qquad
 z_i=-\frac{29}{108}\sqrt{\frac32}\,H,\qquad Mz=g.
\end{equation}
Normalized at \(\alpha_M=22/(27H^2)\), it has
\(M/\alpha_M=I\), \(\kappa=1\), and
\begin{equation}
 \alpha_M\frac{\norm{M^{-1}g}_2}{\norm g_2}=1.
\end{equation}
The states of \(g,z,Wz\) are public and free.  For the fixed diagonal-pair
observable
\[
 O_S=\sum_{i\in S}
 \bigl(\ket{u_i}\!\bra{u_i}-\ket{v_i}\!\bra{v_i}\bigr),
\]
each output node contributes
\(u_i^2-v_i^2=p_d p_q=-(4/27)p_NH^2\).  A free-node block has squared norm
\((851/1458)H^2\), a prefix block has squared norm
\(DH_i^2\) with \(D=16139/23328<7/10\), and
\(\sum_iH_i^2<(53/3)NH^2\).  Hence
\begin{equation}
 \left|\left\langle
 \frac{p_\sigma}{\norm{p_\sigma}_2},O_S
 \frac{p_\sigma}{\norm{p_\sigma}_2}
 \right\rangle\right|
 >\frac{(64/27)NH^2}{(7/10)(53/3)NH^2}
 =\frac{1920}{10017}>\frac16.
\end{equation}
Its sign is \(-p_N\), so
preparing \(\ket{p_\sigma/\norm{p_\sigma}}\) to trace error \(1/100\)
costs \(\Theta(P)\) raw queries.  Because
\(p_\sigma\perp Wz\), their coherent combination succeeds with probability
at least \(1/2\), provided one has phase-referenced coherent preparation
access to both normalized summands.  A density-state contract alone does not
supply that access, as \cref{sec:prelim-output-contracts} records.  The
localization is gauge-specific: an arbitrary
particular solution can trade a null component with \(z\); the invariant
claim is the end-to-end original-coordinate lower bound.

The capstone has the same tight static ledger \eqref{eq:lp-tag-9-29}.  Its matching upper
bound caches all signs in \(N\) queries.  A genuine iteration multiplier is
valid only for a streaming service in which round \(t\) releases a fresh
independent sign block after output \(t-1\) is committed.  Before the next
block is released, round \(t\) must return an unconditional state within
trace distance \(1/100\) of that round's hard original-coordinate Newton
direction (or of a point under the robust contract in
\cref{thm:lp-prefix-capstone}).  Let \(Q_t\) count every call to the current
block's coefficient oracle or inverse, all setup after that block is
released, and all queries used to prepare classical or quantum advice
consumed in that round.  Conditioned on all past memory, the new block is
independent, so \cref{thm:lp-prefix-capstone} applies in every round and
\begin{equation}
 \sum_{t=1}^RQ_t=\Omega(RN)=\Omega(RP).
\end{equation}
This is a freshness direct sum, not a per-iteration theorem for one fixed LP.

\subsection{Full-KKT and preconditioned formulations}
\label{sec:lp-full-kkt}

The preceding results request only the original-primal block.  A related
construction makes the complete primal--dual direction hard while keeping its
unpreconditioned right-hand side public.  Use the unrigidified gain--plateau
constraints \eqref{eq:lp-tag-9-20}, the base objective \(c_i=(2,2,1,1)\), and the public start
\begin{equation}
 x^0=s^0=e,\qquad y^0=0,\qquad \mu^0=1.
\end{equation}
For a public centering target \(\theta\in[0,1]\), the Newton equations are
\begin{equation}
 A_\sigma\Delta x=b-A_\sigma e,\quad
 A_\sigma^T\Delta y+\Delta s=c-e,\quad
 \Delta x+\Delta s=(\theta-1)e.
\label{eq:lp-tag-9-62}
\end{equation}
The RHS is input-independent: every hidden coefficient multiplies the zero
pair difference at the start.

Let \(R_{r,j}=\sum_{i=j}^{N+K}H_i^r\).  Direct elimination and backward
substitution give
\begin{equation}
\begin{aligned}
 \Delta x_{u_i}&=(4H_i-8+3H_i\tau_i)/6,&
 \Delta x_{v_i}&=(4H_i-8-3H_i\tau_i)/6,\\
 \Delta x_{h_i}&=H_i-1,&
 \Delta x_{t_i}&=(2H_i-1)/3,\\
 \alpha_j&=\tau_jR_{2,j}/(2H_j),&
 \beta_j&=[7R_{2,j}+(4-9\theta)R_{1,j}]/(3H_j),\\
 \gamma_j&=(2H_j+2)/3-\theta .
\end{aligned}
\label{eq:lp-tag-9-63}
\end{equation}
Here \((\alpha,\beta,\gamma)=\Delta y\) correspond to the signed,
reference, and cap rows.

\begin{theorem}[Public-RHS full-KKT direction]
\label{thm:lp-full-kkt}
Let \(w_\sigma=(\Delta x,\Delta y,\Delta s)\) solve \eqref{eq:lp-tag-9-62}.
Preparing its
normalized state to trace distance \(1/100\) takes at least \(N/2\) raw
coefficient queries.  The same holds if an approximate vector obeys
\(\norm{\widetilde w-w_\sigma}_2\leq\norm{w_\sigma}_2/300\) and its state
has trace error at most \(1/300\).

For \(\theta=1\), the two-block elimination has an explicit canonical sparse
block encoding of normalization six and a public RHS state.  The same
\(N/2\) lower bound holds for calls to that fully specified block encoding,
not merely to its principal matrix block.
\end{theorem}

\begin{proof}
On the \(K\) output rows define
\(A_S=H^2\sum_{\ell=1}^K\ell^2/4\).  The height sums and
\eqref{eq:lp-tag-9-63} imply
\begin{equation}
 \norm\alpha_2^2<3A_S,\qquad
 \norm\beta_2^2<\frac{484}{3}A_S,\qquad
 \norm{w_\sigma}_2^2<165A_S.
\end{equation}
Moreover \(\alpha_j\) has sign \(p_N\) and \(\beta_j>0\) on the output
suffix, with \(\beta_j/|\alpha_j|\geq3\).  The public norm-one observable
that swaps the corresponding \(\alpha_j,\beta_j\) coordinates therefore has
parity-signed expectation
\begin{equation}
 p_N\langle O_{\rm KKT}\rangle>\frac6{165}=\frac2{55}>\frac1{30}.
\label{eq:lp-tag-9-65}
\end{equation}
Trace error \(1/100\) changes this by at most \(1/50\).  The resulting
expectation is a real polynomial of degree at most twice the query count; its
nonzero parity Fourier coefficient forces degree at least \(N\).  Thus the
query count is at least \(N/2\).  The triangle inequality for normalized
rays gives the approximate-vector statement.

At \(\theta=1\), eliminate \(\Delta s=-\Delta x\), put \(y'=-\Delta y\),
and write
\begin{equation}
 {\cal K}_\sigma=\begin{bmatrix}I&A_\sigma^T\\A_\sigma&0\end{bmatrix},
 \qquad
 r_\sigma=\binom{e-c}{b-A_\sigma e}.
\end{equation}
The maximum absolute row and column sums are six.  The standard
magnitude-weighted row/column state-pair construction gives an exact normalization-six
block encoding.  Its support, magnitudes, failure arms, and left preparation
are public; the right preparation uses one local sign phase, hence exactly
one sign query per call, adjoint, or controlled call.  Also
\(\norm{r_\sigma}_2^2=3P+T+1\), and its preparation is public.  Removing
\(\Delta s\) can only strengthen \eqref{eq:lp-tag-9-65}, so the same degree argument applies
to this complete oracle implementation.
\end{proof}

Ideal Schur preconditioning makes the matrix well conditioned but does not
make the end-to-end raw-input task easy.  Put
\[
 {\cal S}_\sigma=A_\sigma A_\sigma^T,\quad
 Q_\sigma={\cal S}_\sigma^{-1/2}A_\sigma,\quad
 P_\sigma=\diag(I,{\cal S}_\sigma^{-1/2}).
\]
Define
\[
 \overline w_\sigma=P_\sigma^{-1}\binom{\Delta x}{y'},\qquad
 \overline r_\sigma=P_\sigma r_\sigma,\qquad
 \overline K_\sigma=P_\sigma{\cal K}_\sigma P_\sigma.
\]
Then
\begin{equation}
 \overline K_\sigma=
 \begin{bmatrix}I&Q_\sigma^T\\Q_\sigma&0\end{bmatrix},
 \qquad
 \overline K_\sigma\overline w_\sigma=\overline r_\sigma,
\end{equation}
and the primal block of \(\overline w_\sigma\) is exactly \(\Delta x\).
Since \(Q_\sigma Q_\sigma^T=I\), the spectrum is \(1\) on the nullspace and
\begin{equation}
 \frac{1+\sqrt5}{2},\qquad \frac{1-\sqrt5}{2}
\end{equation}
on every coupled two-plane.  Therefore
\begin{equation}
 \norm{\overline K_\sigma}_2
 =\norm{\overline K_\sigma^{-1}}_2=\varphi,\qquad
 \kappa_{\rm abs}(\overline K_\sigma)=\varphi^2
 =\frac{3+\sqrt5}{2}.
\label{eq:lp-tag-9-69}
\end{equation}
This is the classical ideal-Schur spectrum of Murphy--Golub--Wathen
\cite{murphy2000-a-note-on-preconditioning-for}.

\begin{theorem}[Residual-robust Schur-preconditioned KKT state]
\label{thm:lp-preconditioned-kkt}
Suppose \(\widetilde{\overline w}\ne0\) obeys
\begin{equation}
 \norm{\overline K_\sigma\widetilde{\overline w}
       -\overline r_\sigma}_2
 \leq10^{-3}\norm{\overline r_\sigma}_2.
\label{eq:lp-tag-9-70}
\end{equation}
Preparing its normalized state to trace error \(1/200\) takes at least
\(N/2\) total raw coefficient queries, counting preconditioner, transformed
RHS, solve, and recovery access.  The result transfers to the
prefix-rigidified family.
\end{theorem}

\begin{proof}
Let
\[
 O_x=\sum_{i\in S}
 \bigl(\ket{D_i}\!\bra{h_i}+\ket{h_i}\!\bra{D_i}\bigr),
 \qquad
 \ket{D_i}=\frac{\ket{u_i}-\ket{v_i}}{\sqrt2}.
\]
This norm-one, primal-supported observable has exact parity-signed
expectation greater than \(2/87\) on the preconditioned solution.  Indeed,
\begin{equation}
 \norm{\Delta x}_2^2\leq\frac{11}{3}PH^2,\qquad
 p_N\frac{\langle\Delta x|O_x|\Delta x\rangle}
              {\norm{\Delta x}_2^2}>\frac16,\qquad
 \norm{\overline w}_2^2<\frac{29}{4}\norm{\Delta x}_2^2.
\end{equation}
Equations \eqref{eq:lp-tag-9-69}--\eqref{eq:lp-tag-9-70} give
\(\norm{\widetilde{\overline w}-\overline w}_2
 \leq10^{-3}\varphi^2\norm{\overline w}_2\).  The total target-state
distance is at most
\begin{equation}
 \frac1{200}+\frac{2\varphi^2}{1000}<\frac1{87},
\label{eq:lp-tag-9-72}
\end{equation}
which preserves the sign of the observable.  The polynomial argument gives
\(N/2\).

For the prefix-rigidified transfer, use the same all-ones start.  Every
added row \(q_i-\frac54h_i=0\) has the public primal residual \(-3/4\), and
the uniform objective is
\[
 c'_i=(1+\omega,1+\omega,1,1),\qquad
 \omega=\frac7{36H}.
\]
The hidden difference coefficients still multiply \(u_i-v_i=0\), so the
complete Newton RHS remains public.  On every free node, null projection and
exact feasibility give
\[
 \Delta q_i=\frac{4H-6-2\omega}{3},\qquad
 \Delta h_i=H-1,\qquad
 \Delta u_i-\Delta v_i=H\tau_i.
\]
The prefix variables are fixed by feasibility.  These formulas yield
\[
 \norm{\Delta x}_2^2<5\sum_iH_i^2,\qquad
 p_N\frac{\langle\Delta x|O_x|\Delta x\rangle}
              {\norm{\Delta x}_2^2}>\frac18,
\]
because the numerator is \(K\sqrt2H(H-1)\), each prefix node contributes
less than \(5H_i^2\), and each free node less than \(4H^2\).
For the exact block preconditioner the transformed primal RHS is
\(e-c'=(-\omega,-\omega,0,0)\) per node, and
\[
 \frac{\norm{e-c'}_2}{\norm{\Delta x}_2}<\frac{21}{100},
 \qquad
 \norm{\overline w_\sigma}_2^2
 <\frac52\norm{\Delta x}_2^2.
\]
Thus the parity-signed expectation on the complete preconditioned state
exceeds \(1/20\).  Its condition remains \(\varphi^2\), while the solve and
output errors in \eqref{eq:lp-tag-9-72} change the expectation by less than
\(2/87<1/20\).  The same \(N/2\) conclusion follows.
\end{proof}

There is no matrix-only normalization/condition tradeoff behind this theorem.
The row-space projector is public because \(\ker A_\sigma\) is public.  If
\(Q_0\) is a public row isometry with \(Q_0^TQ_0=I-WW^T\), then
\(O_\sigma=Q_\sigma Q_0^T\) is orthogonal and
\begin{equation}
 \diag(I,O_\sigma)^T\overline K_\sigma\diag(I,O_\sigma)
 =\begin{bmatrix}I&Q_0^T\\Q_0&0\end{bmatrix},
\end{equation}
which is wholly public.  Input dependence has moved to the transformed RHS
and recovery map.  Accordingly the invariant accounting statement is
\begin{equation}
 Q_{\rm pre}+Q_{\rm on}\geq N/2.
\end{equation}
It is tight: after \(N\) cached sign queries, later raw-query cost can vanish.
The separate factor-access frontier is given in
\cref{sec:preconditioning}; it is not proved here.

\subsection{Bounded data and the sharp residual scale}
\label{sec:lp-bounded-data}

Remove the gain prefix by setting every height and gain to one, retain a
length-\(N\) signed chain followed by \(K=16N\) public copy edges, and put
\(P=17N+1\).  Use the same four nonnegative variables at every node and the
equations
\begin{equation}
\begin{aligned}
 d_0&=1,&d_i-a_id_{i-1}&=0,\\
 h_0&=1,&h_i-h_{i-1}&=0,\\
 &&q_i+t_i-2h_i&=0,
\end{aligned}
\qquad
 a_i=\begin{cases}\sigma_i&i\leq N,\\1&i>N.\end{cases}
\end{equation}
Use objective
\begin{equation}
 c^Tx=\sum_i\left[\frac{43}{36}(u_i+v_i)+h_i+t_i\right].
\end{equation}
All matrix, RHS, objective, feasible-primal, and relevant slack coordinates
are bounded by absolute constants; \(\norm b_2=\sqrt2\).  Exact feasibility
is \(d_i=\tau_i,h_i=1,t_i=2-q_i,1\leq q_i\leq2\), so the same basis
\eqref{eq:lp-gain-null} spans the nullspace.

\begin{theorem}[Bounded exact-state separation]
\label{thm:lp-bounded-exact}
The all-unit family is bounded and strictly primal--dual feasible, with a
unique nondegenerate strictly complementary optimum.  Its raw orthonormal
reduced central Hessian has condition one for every \(\mu>0\), and its
reduced Newton matrix at the public start below is \((22/27)I\).

At \(\mu_1=1/16\), one exact Newton correction from
\begin{equation}
 x_i^0=(5/6,5/6,20/27,5/9),\quad
 s_i^0=(10/27,10/27,5/12,5/9),\quad y^0=0
\label{eq:lp-tag-9-77}
\end{equation}
with \(\mu_0=25/81\) and target \(25/162\) lands at the exact center
\begin{equation}
 x_i^1=(9/8,1/8,1,3/4),\qquad
 s_i^1=(1/18,1/2,1/16,1/12),
\end{equation}
up to the hidden pair swap.  Preparing either the exact primal direction
state or exact central-primal state to trace error \(1/100\) requires
\(\Omega(N)=\Omega(P)\) raw queries.
\end{theorem}

\begin{proof}
On the feasible line the local objective is \(3+(7/36)q_i\), so the unique
optimum has \(q_i=1\).  The positive columns form a basis and the zero pair
coordinate has reduced cost \(7/18\).  All central coordinates solve the
same scalar equation
\[
 \frac{7}{36\mu}=\frac1{q+1}+\frac1{q-1}-\frac1{2-q},
\]
so the reduced Hessian is scalar.  At \(q=5/4,\mu=1/16\), both sides are
\(28/9\), and direct evaluation gives the \(\mu_1\)-scaled reduced
curvature \(182/243\).  Equations
\eqref{eq:lp-tag-9-37}--\eqref{eq:lp-tag-9-38} with \(H_i=1\) prove the same exact Newton secant and
\((22/27)I\) start matrix.

The primal direction is the constant block in \eqref{eq:lp-tag-9-39}.  Its plateau decoder
has bias
\[
 \frac{K(5/12)}{2DP}>\frac14.
\]
Each central block has squared norm \(91/32\) and oriented pair gap \(5/4\),
so the central-state decoder has bias
\[
 \frac{K(5/4)}{2(91/32)P}=\frac{20K}{91P}>\frac15.
\]
Lemma~\ref{lem:lp-decoder} completes both lower bounds.
\end{proof}

The all-unit family is robust, but only at the sharp vanishing scale.

\begin{theorem}[Sharp all-unit residual threshold]
\label{thm:lp-unit-threshold}
For every nonzero \(x\geq0\) satisfying the absolute residual bound
\begin{equation}
 \norm{A_\sigma x-b}_2\leq\frac1{16\sqrt N},
\label{eq:lp-tag-9-79}
\end{equation}
the normalized primal state has constant correct-parity bias.  Preparing it
to trace error \(1/100\) needs \(\Omega(N)=\Omega(P)\) raw queries.  In
relative form, \eqref{eq:lp-tag-9-79} is \(1/(16\sqrt{2N})\); in particular it contains
the \(1/(100\sqrt N)\) absolute residual tube.

The order is sharp: a nonnegative zero-signal plateau exists at absolute
residual \(1/\sqrt N\), and a nonnegative entirely wrong-parity plateau
exists at absolute residual \(2/\sqrt N\), even with exact optimal objective.
After public homogeneous rigidification, the wrong witness can also have
exact dual feasibility and exact coordinatewise complementarity.
\end{theorem}

\begin{proof}
Let \(E=\norm{Ax-b}_2\), \(z_i=\tau_id_i\), and
\(\delta=\sqrt P E\).  Telescoping the signed and reference recurrences gives
\begin{equation}
 |z_i-1|\leq\sqrt P E,\qquad |h_i-1|\leq\sqrt P E.
\end{equation}
The factor \(\sqrt P\) is exact by equality in Cauchy--Schwarz.  Under
\eqref{eq:lp-tag-9-79}, \(P/N\leq35/2\) gives \(\delta<4/15\), so on every output node
\(p_Nd_i>11/15\).  The cap and nonnegativity give
\(q_i,t_i<31/12\), \(h_i<19/15\), and
\(\norm x_2^2<15P\).  Thus the output decoder has bias
\begin{equation}
 \frac{K(11/15)^2}{30P}>\frac1{70}.
\end{equation}
Trace error \(1/100\) leaves bias at least
\(1/70-1/100=3/700>0\).

For the zero-signal witness take
\begin{equation}
 z_i=1-i/N\ (i\leq N),\qquad z_i=0\ (i>N),
\end{equation}
and for the wrong-parity witness replace this by
\begin{equation}
 z_i=1-2i/N\ (i\leq N),\qquad z_i=-1\ (i>N).
\end{equation}
Set \(d_i=\tau_i z_i\) and \(h_i=q_i=t_i=1\).  All reference and cap rows
are exact.  The \(N\) signed residuals are respectively \(1/N\) and
\(2/N\), proving the claimed norms.  Both points are nonnegative and have
the exact optimal objective.

For the stronger failure, add \(q_i-5h_i/4=0\) on all non-output nodes and
average the exact centers of the \(N\) inputs obtained by flipping one sign.
The average has the wrong common output template and the same
\(2/\sqrt N\) primal residual.  Every coordinate is positive.  Setting
\(\bar s=\mu_1\bar X^{-1}e\) gives exact complementarity.  The remaining
nullspace is supported on output nodes, where \(\bar x\) is itself an exact
central template; hence \(c-\bar s\perp\ker A\), which supplies an exact
dual multiplier.
\end{proof}

The recurrence-only cutoff is exactly \(E<1/\sqrt P\): taking
\(z_i=1-(i+1)/P\) makes all gauged residuals \(-1/P\), reaches zero at the
last node, and has norm \(1/\sqrt P\).  Inverting the robust scale over all
hidden lengths \(L\) with \(17L+1\leq P\) gives the worst-case accuracy law
\begin{equation}
 Q(P,\epsilon)=\Omega\bigl(\min\{P,\epsilon^{-2}\}\bigr)
\end{equation}
for absolute or RHS-relative residual, up to constants.  This is a
reparameterized family bound, not accumulating hardness at several accuracies
on one instance.

The public copy plateau is necessary for constant one-copy distinguishability
at unit height.  Flipping only \(\sigma_N\) swaps exactly \(K+1\) target
blocks.  For the central and direction states the normalized overlaps are
\begin{equation}
 1-\frac{32(K+1)}{91P},\qquad
 1-\frac{K+1}{DP},
\end{equation}
respectively.  If \(K=o(N)\), their trace distance is \(o(1)\) even for
opposite parities.  Hence the correct bounded-data resource is
\(K=\Theta(N)\) copies, not \(\Theta(\sqrt N)\) dynamic range.

\subsection{Dual homogeneity and the bounded full direction}
\label{sec:lp-dual-homogeneity}

There is one exact way to suppress the multiplier accumulation of a unit-height
chain.  Positive dual scaling is an elementary symmetry of the central and
Newton equations, consistent with the general scale-invariance viewpoint in
interior-point theory \cite{tuncel1998-primal-dual-symmetry-scale-invariance}.

\begin{lemma}[Dual homogeneity]
\label{lem:lp-dual-homogeneity}
In the standard Newton system, leave \((A,b,x^0)\) fixed and replace
\begin{equation}
 (c,y^0,s^0,\widehat\mu)
 \longmapsto
 (\lambda c,\lambda y^0,\lambda s^0,\lambda\widehat\mu),
 \qquad\lambda>0.
\end{equation}
Then \((\Delta x,\Delta y,\Delta s)\) is replaced by
\((\Delta x,\lambda\Delta y,\lambda\Delta s)\).  The primal central
trajectory is unchanged, with \(\mu\) reparameterized as
\(\mu'=\lambda\mu\).
\end{lemma}

\begin{proof}
The primal equation is unchanged, while substitution multiplies the dual and
complementarity equations by \(\lambda\).  Full row rank and positivity make
the Newton solution unique.  The same substitution in
\(Ax=b,A^Ty+s=c,Xs=\mu e\) proves the central-path statement.
\end{proof}

For the full-state theorem use the all-unit signed chain and copy plateau, but
pin each reference coordinate independently by \(h_i=1\).  This changes the
row representation, not the primal feasible set.  Starting from
\eqref{eq:lp-tag-9-77}, scale
the objective, initial slacks, and complementarity data by
\begin{equation}
 \lambda=\frac1P.
\label{eq:lp-tag-9-87}
\end{equation}
The exact primal correction remains the block \eqref{eq:lp-tag-9-39}.  If the first coordinate
is paired with the larger endpoint coordinate, the unscaled slack correction
is
\begin{equation}
 \Delta\bar s_i=(-17/54,7/54,-17/48,-17/36),
\label{eq:lp-tag-9-88}
\end{equation}
and, with difference, height-pin, and cap multipliers,
\begin{equation}
 \Delta y=\frac1P(\bar\alpha,\bar\beta,\bar\gamma),\quad
 \bar\alpha_j=\frac29\tau_j(P-j),\quad
 \bar\beta_j=\frac{133}{48},\quad
 \bar\gamma_j=\frac{11}{12}.
\label{eq:lp-tag-9-89}
\end{equation}

\begin{theorem}[Bounded full-KKT direction]
\label{thm:lp-dual-homogeneous-full}
Every coordinate of the complete \(11P\)-dimensional direction
\(w=(\Delta x,\Delta y,\Delta s)\) in
\eqref{eq:lp-tag-9-87}--\eqref{eq:lp-tag-9-89} is bounded by an
absolute constant, and \(\norm w_2=\Theta(\sqrt P)\).  Preparing its
normalized state to trace error \(1/100\) requires
\(\Omega(N)=\Omega(P)\) raw queries.

The explicit canonical sparse block encoding of either the three-block Newton
matrix or its two-block elimination has normalization five, a public RHS, and
the same linear call lower bound.  The reduced Newton condition remains one,
but the central endpoint is at \(\mu_1'=1/(16P)\), and both raw full-KKT
matrices have condition \(\Theta(P)\).
\end{theorem}

\begin{proof}
Equations \eqref{eq:lp-tag-9-39}, \eqref{eq:lp-tag-9-88}, and
\eqref{eq:lp-tag-9-89} give
\[
 \norm{\Delta x}_\infty\leq17/24,
 \quad \norm{\Delta y}_\infty\leq2/9,
 \quad \norm{\Delta s}_\infty\leq17/(36P).
\]
Writing \(D=16139/23328\), direct summation gives
\begin{equation}
 DP\leq\norm w_2^2<\left(D+\frac1{40}\right)P<\frac34P.
\end{equation}
Use the diagonal decoder convention \(+1\) on every \(v\) coordinate and
\(-1\) on every \(u\) coordinate.  The underlying \(u\)-minus-\(v\)
squared-mass difference on each output node is \(-p_N(5/12)\), so the
decoder expectation contributes \(+p_N(5/12)\) per node and has magnitude
\begin{equation}
 \frac{(5/12)K}{(3/4)P}=\frac{5K}{9P}
 \geq\frac{32}{63}>\frac12.
\end{equation}
Trace error preserves constant bias.  The polynomial method gives at least
\(N/2\) sign queries.

The maximum absolute row and column sums of both canonical matrices are five.
The same magnitude-weighted state-pair construction used in
Theorem~\ref{thm:lp-full-kkt} therefore gives normalization five and one
sign query per full oracle call.  The RHS remains public because the start
has zero pair difference.  This proves the canonical-oracle claim.

Dual homogeneity gives \(\mu_1'=1/(16P)\) and the reduced start matrix
\((22/(27P))I\).  If \(z\in\ker A\) is unit, applying either raw KKT matrix
to the corresponding nullspace test vector has norm at most \(1/P\), while
the largest singular value is bounded below by a constant.  Conversely, the
pinned signed incidence has smallest singular value \(\Theta(P^{-1})\);
bounded local changes reduce the remaining blocks to identities.  The paired
saddle eigenvalues then show that both conditions are \(\Theta(P)\).
\end{proof}

The price is real: objective coefficients, positive slacks, reduced costs, and
complementarity are \(\Theta(P^{-1})\).  The theorem is not a simultaneous
constant-scale and constant-full-condition result.  It also does not inherit
the sharp relative-residual theorem, because independent height pins make
\(\norm b_2=\Theta(\sqrt P)\).

\subsection{Why the inverse-scale price remains open}
\label{sec:lp-beyond-pairs}

The obstruction can be stated without signed pairs.  At any positive start,
eliminate \(\Delta s\) from the Newton equations and put
\[
 H=(X^0)^{-1}S^0,\qquad g=(X^0)^{-1}r_c-r_d.
\]
Then
\begin{equation}
 H\Delta x-A^T\Delta y=g,\qquad
 \boxed{r_p^T\Delta y=\Delta x^TH\Delta x-g^T\Delta x}.
\end{equation}
This exact work identity is representation-independent.  If \(g=0\),
\(H\succeq\rho I\), and \(\norm{r_p}_1\leq R\), it gives
\begin{equation}
 \norm{\Delta y}_\infty\geq
 \frac{\rho\norm{\Delta x}_2^2}{R}.
\end{equation}
A constant-scale correction with \(\Theta(P)\) squared norm behind a sparse
public source therefore forces an \(\Omega(P)\) multiplier, independently of
redundant rows or branching.

There is also a two-input version for constant-dimensional copy blocks.  Let a
flat output region \(S\) carry \(z_v\in\R^k\) and orthogonal copy equations
\(z_v-U_ez_u=0\).  After public transports \(G_v\), suppose stationarity is
\begin{equation}
 \operatorname{div}_U(f)_v=D_v\Delta z_v+g_v,
 \qquad D_v\succeq\rho I,
\label{eq:lp-tag-9-94}
\end{equation}
with the same numerical \(D_v,g_v\) for two opposite-parity inputs, and their
transported codewords differ by a common \(w\), \(\norm w_2\geq2a\).  If a
boundary row scale is at most \(B\), summing the difference of
\eqref{eq:lp-tag-9-94} over
\(S\) cancels every public source and internal force.  Taking the inner
product with \(w\) yields
\begin{equation}
 \max_{p\in\{+,-\}}\norm{\Delta y_{\partial S}^p}_{2,\infty}
 \geq\frac{\rho a|S|}{B|\partial S|}.
\end{equation}
Thus swaps, simplex permutations, and public flat rotations do not remove
dual accumulation behind a bounded interface.

The most tempting counterexample is a real quarter-turn
\(J=\left[\begin{smallmatrix}0&-1\\1&0\end{smallmatrix}\right]\), with
\(z_i=\sigma_iJz_{i-1}\) and closure \(z_0-Jz_N=e_1\).  For \(4\mid N\),
\begin{equation}
 z_0(p)=\frac12(1,p)^T,
\end{equation}
so all blocks are bounded and locally reveal parity after a public inverse
rotation.  Nevertheless the work identity gives the closure multiplier
\begin{equation}
 e_1^T\Delta y_{\rm close}=\sum_{i=0}^N\norm{z_i}_2^2
 =\frac{N+1}{2}.
\label{eq:lp-tag-9-97}
\end{equation}
After gauging, the constraint is a twisted cycle whose singular values are
\begin{equation}
 2\left|\sin\frac{2\pi\ell\pm\pi/2}{2(N+1)}\right|,
 \qquad \ell=0,\ldots,N,
\end{equation}
so its least singular value is
\(2\sin(\pi/[4(N+1)])=\Theta(N^{-1})\).  A positive/negative standard-form
LP realization preserves bounded primal and slack directions and even allows
a positive half-centering step, but \eqref{eq:lp-tag-9-97} survives exactly.  Circulations,
public antisymmetric sources, and dense public residuals fail for the same
work or two-parity cancellation reason unless they introduce linearly many
genuinely independent hard interfaces.

\begin{openproblem}[Bounded-scale full-KKT parity amplifier]
\label{open:lp-bounded-full-kkt}
Does a sparse standard-form LP family exist that simultaneously has:
\begin{enumerate}[label=(\roman*)]
 \item linear dimension and bounded row and column degree;
 \item bounded coefficients with one-bit-local input dependence;
 \item public, simply preparable Newton right-hand sides;
 \item primal and slack starts bounded above and below by positive constants;
 \item a coordinatewise bounded complete Newton direction; and
 \item a parity-hard primal component with constant mass in that complete
 direction state?
\end{enumerate}
The work identity and orthogonal-copy cut law rule out the natural signed,
permutation, rotation, and flat-copy constructions.  They do not cover
arbitrary nonorthogonal mixed-mode constraints, input-dependent sources whose
loading is separately charged, or output interfaces that quotient the equality
multipliers.
\end{openproblem}

The distinction closing the section is therefore exact.  The
prefix-rigidified construction proves condition-one \emph{reduced Newton
geometry} together with a linear original-coordinate loading/recovery cost.
The preconditioned full-KKT result proves a constant-condition residual theorem
only after charging construction of its data-dependent access.  The bounded
full direction obtained by dual homogeneity instead has canonical full-system
condition \(\Theta(P)\).  None is a constant-\(\kappa\) QLS lower bound with
free matrix and RHS preparation oracles; each is an end-to-end LP access and
output theorem.

\section{Limits of Gadget Classes}
\label{sec:gadget-frontiers}

\Cref{sec:lp-hardness} made a local parity signal robust by
replicating constraints, scaling rows, or increasing the height of the
propagated signal.  This section asks whether some other sparse LP gadget
can avoid all three costs.  For several broad, precisely stated classes the
answer is no.  The most general result is a convex-mixture theorem that
assumes no graph structure.  For signed-copy graphs and bounded-arity Walsh
rows, electrical resistance gives a geometric refinement.  Conservation of
the dual multipliers shows that complete central and Newton outputs force a
large row scale, inverse dual scale, multiplier size, or input incidence.

Throughout this section the hidden input is
\(\varsigma\in\{\pm1\}^N\), with \(N\geq2\), and
\(p(\varsigma)=\prod_{i=1}^N\varsigma_i\).  Input locality always refers to
the raw sparse coefficient access of \cref{def:raw-sparse-access}.
A residual written as \(\norm{Ax-b}_2\) is absolute; when a relative
threshold is named, it is the global RHS-relative residual of
\cref{def:global-residual}.  The results are structural limits on
output requirements.  They become query lower bounds only after a separate
reduction shows that the oracle can be simulated locally and that parity can
be decoded.

\subsection{Convex mixtures of neighboring witnesses}
\label{sec:frontier-mixture}

Consider a fixed-support family
\begin{equation}
 \mathcal F_{\varsigma}
 =\{x\in\R_{\geq0}^n:A_{\varsigma}x=b\},
 \label{eq:frontier-local-family}
\end{equation}
where \(b\) is public.  Every row has at most \(s\) nonzeros, every
coefficient has magnitude at most \(B\), and each input-dependent position
\((r,j)\) is a function of one designated bit
\(\varsigma_{\ell(r,j)}\).  Let \(M_{\rm dep}\) be the number of these
positions.  A bit may label many positions, including several in one row.
Fix a base input, let \(\varsigma^{(i)}\) flip bit \(i\), and suppose that
each such neighbor has an exact witness
\begin{equation}
 x^{(i)}\in\mathcal F_{\varsigma^{(i)}},
 \qquad \norm{x^{(i)}}_\infty\leq H
 \quad(1\leq i\leq N).
 \label{eq:frontier-neighbor-witnesses}
\end{equation}
The magnitude bounds \(B,H\) are nonnegative.  The residual estimates
allow \(s=0\); the resource inequalities below that divide by \(s\) assume
\(s>0\).

\begin{theorem}[Single-flip convex mixture]
\label{thm:frontier-single-flip}
The average \(\bar x=N^{-1}\sum_i x^{(i)}\) is nonnegative and satisfies
\begin{equation}
 \boxed{\norm{A_{\varsigma}\bar x-b}_2
 \leq \frac{2BH\sqrt{sM_{\rm dep}}}{N}.}
 \label{eq:frontier-single-flip}
\end{equation}
If a public objective \(c\) has the same optimum \(v_*\) on every instance
and the witnesses in \eqref{eq:frontier-neighbor-witnesses} are optimal, then
\(c^\top\bar x=v_*\): the mixture is objective-exact, although it is
generally infeasible for the base instance.
\end{theorem}

\begin{proof}
Feasibility for the neighboring instances gives
\[
 N(A_{\varsigma}\bar x-b)
 =\sum_{i=1}^N
   (A_{\varsigma}-A_{\varsigma^{(i)}})x^{(i)}.
\]
Let \(D_r\) be the dependent positions in row \(r\), and set
\(m_r=|D_r|\).  In the \(r\)-th coordinate, position \((r,j)\) contributes
only to the summand \(i=\ell(r,j)\), so this coordinate is a sum of at most
\(m_r\) terms, each of magnitude at most \(2BH\).
Since \(m_r\leq s\),
\[
 \left|N(A_{\varsigma}\bar x-b)_r\right|^2
 \leq4B^2H^2m_r^2
 \leq4sB^2H^2m_r.
\]
Summing over rows and using \(\sum_rm_r=M_{\rm dep}\) proves
\eqref{eq:frontier-single-flip}.  Convexity proves nonnegativity, and
linearity proves the objective statement.
\end{proof}

The theorem also covers a right-hand side whose entries each depend on one
bit: apply it to \([A_{\varsigma},-b_{\varsigma}](x,1)=0\).  This uses row
sparsity \(s+1\) and height \(\max\{H,1\}\), replaces \(B\) by a bound on
every entry of \([A_{\varsigma},-b_{\varsigma}]\), and counts the dependent
RHS positions.  Inequalities are covered after slack conversion only when
the selected slack coordinates satisfy the same height bound.

The output hypothesis cannot be omitted.  For an affine public score
\(\phi(x)=a^\top x-a_0\), suppose every selected witness for input \(\xi\)
obeys \(p(\xi)\phi(x^\xi)\geq\gamma>0\).  All single-flip neighbors have
parity \(-p(\varsigma)\), so
\(p(\varsigma)\phi(\bar x)\leq-\gamma\).  Assume \(b\neq0\), as in
\cref{def:global-residual}, and call an output requirement
\emph{sound at threshold \(\eta\geq0\)} if it rejects every nonnegative,
objective-exact point with the wrong output and relative residual at most
\(\eta\).  A sound output requirement needs
\begin{equation}
 M_{\rm dep}B^2H^2
 >\frac{\eta^2\norm b_2^2}{4s}N^2.
 \label{eq:frontier-mixture-product}
\end{equation}
In particular, for fixed \(s,B,\eta,\norm b_2\), linear incidence and
bounded height are incompatible with soundness at a fixed threshold
\(\eta\) (a fixed residual tube).  Here and in later asymptotic statements
about a fixed tube, \(\eta>0\).

The same conclusion holds for an output requirement on the normalized primal
state \(x/\norm x_2\) under a concrete mixture-stability condition.  Let
\(\mathcal S\) be \(L\geq1\) pairwise disjoint public output pairs
\((u_\ell,v_\ell)\), and suppose
\begin{equation}
 p(\xi)(u_\ell^\xi-v_\ell^\xi)\geq a>0
 \quad(\ell\in\mathcal S)
 \label{eq:frontier-pair-margin}
\end{equation}
for every selected witness.  Averaging gives
\(-p(\varsigma)(\bar u_\ell-\bar v_\ell)\geq a\).  Nonnegativity gives
\(\bar u_\ell+\bar v_\ell\geq a\), hence
\[
 -p(\varsigma)(\bar u_\ell^2-\bar v_\ell^2)\geq a^2.
\]
For
\(Q_{\mathcal S}=\sum_{\ell\in\mathcal S}
 (\ket{u_\ell}\!\bra{u_\ell}-\ket{v_\ell}\!\bra{v_\ell})\),
\begin{equation}
 -p(\varsigma)
 \left\langle\frac{\bar x}{\norm{\bar x}_2}\middle|
 Q_{\mathcal S}\middle|\frac{\bar x}{\norm{\bar x}_2}\right\rangle
 \geq\frac{La^2}{nH^2}.
 \label{eq:frontier-pair-decoder}
\end{equation}
Measure a coordinate, output its pair sign on the selected pairs, and
output an independent fair sign otherwise.  The wrong-parity bias of this
binary measurement is half the signed expectation in
\eqref{eq:frontier-pair-decoder}.
Thus \(La^2=\Omega(nH^2)\) gives a constant wrong-parity bias.
Without an affine score, a common output template, the pair margin
\eqref{eq:frontier-pair-margin}, or another explicit mixture-stability
property, \cref{thm:frontier-single-flip} makes no claim about a
normalized quadratic decoder.

Two refinements show the roles of sensitivity and uneven incidence.  For a
Boolean function \(f\), fix a base input whose sensitive
set \(S_f\) has size \(k>0\), and average only over its sensitive neighbors.
If \(M_S\) counts the positions labelled by bits in \(S_f\), the same proof
gives
\begin{equation}
 \norm{A_{\varsigma}\bar x_S-b}_2
 \leq\frac{2BH\sqrt{sM_S}}{k}.
 \label{eq:frontier-sensitivity}
\end{equation}
Under the common-objective and mixture-stable output hypotheses, soundness
at threshold \(\eta\) therefore requires
\[
 M_SB^2H^2>
 \frac{\eta^2\norm b_2^2k^2}{4s}.
\]
More sharply, let \(M_i\) be the number of positions labelled by sensitive
bit \(i\).  For weights \(w_i\geq0\) with \(\sum_iw_i=1\), set
\(x_w=\sum_iw_ix^{(i)}\).  Cauchy--Schwarz in each row gives
\begin{equation}
 \norm{A_{\varsigma}x_w-b}_2^2
 \leq4sB^2H^2\sum_iM_iw_i^2.
 \label{eq:frontier-weighted-mixture}
\end{equation}
When every \(M_i>0\), the right side is minimized by
\(w_i=M_i^{-1}/\sum_hM_h^{-1}\), which gives
\[
 \norm{A_{\varsigma}x_w-b}_2
 \leq2BH\sqrt{\frac{s}{\sum_iM_i^{-1}}}.
\]
The corresponding soundness condition is
\[
 B^2H^2>
 \frac{\eta^2\norm b_2^2}{4s}\sum_iM_i^{-1}.
\]
If some \(M_i=0\), that neighbor's witness is exactly feasible for the base
matrix and already makes the output requirement unsound.  These mixtures preserve
every uniform affine or pair margin and the common objective.  The resource
that matters is therefore \(\sum_iM_i^{-1}\), not merely the total
incidence.

For a fixed relative tube and fixed \(s,\eta,\norm b_2\),
Table~\ref{tab:frontier-currencies} lists three constructions that attain
\eqref{eq:frontier-mixture-product} up to constants.
\begin{table}[t]
 \centering
 \small
 \caption{Constructions attaining the one-bit-local mixture bound
 \eqref{eq:frontier-mixture-product} up to constants, each by increasing one
 resource.}
 \label{tab:frontier-currencies}
 \begin{tabular}{@{}lccc@{}}
  \toprule
  construction & \(M_{\rm dep}\) & \(B\) & \(H\) \\
  \midrule
  bounded-height parallel replication
    & \(\Theta(N^2)\) & \(\Theta(1)\) & \(\Theta(1)\) \\
  direct row scaling
    & \(\Theta(N)\) & \(\Theta(\sqrt N)\) & \(\Theta(1)\) \\
  gain--plateau scaling
    & \(\Theta(N)\) & \(\Theta(1)\) & \(\Theta(\sqrt N)\) \\
  \bottomrule
 \end{tabular}
\end{table}
Every row has \(M_{\rm dep}B^2H^2=\Theta(N^2)\).  The parallel and
gain--plateau families of
\cref{thm:lp-parallel-path,thm:lp-gain-plateau} attain the
first and third rows.  Scaling the propagation rows by \(\sqrt N\)
attains the second without changing the feasible set; the cost appears in
the matrix normalization.  This sharpness concerns only these scales, not
the conditioning or QIPM cost of the formulations.

This is not an extension-complexity lower bound.  Carr--Konjevod give an
\(O(N)\)-size dynamic-programming extended formulation for the parity
polytope, and Ermel--Walter develop further flow formulations
\cite{carr2005-polyhedral-combinatorics,ermel2018-parity-polytopes-and-binarization}.
Those formulations do not guarantee the correct parity output at every point
of a fixed equality-residual tube.  Sparse LP decoding also has classical
fractional pseudocodewords and graph-cover descriptions
\cite{feldman2005-using-linear-programming-to-decode,
koetter2005-characterizations-of-pseudo-codewords-of,
vontobel2005-graph-cover-decoding}; these are conceptually related but do
not give the quantitative bound \eqref{eq:frontier-single-flip}.  In
particular, expander LP decoding may use an input-dependent objective,
whereas exact objective preservation above requires one common objective
and optimum.

\subsection{Extensions: multi-bit coefficients, KKT systems, cones, and centrality}
\label{sec:frontier-mixture-supplements}

If coefficient position \((r,j)\) may depend on a set
\(I_{rj}\subseteq[N]\) of input bits, define
\[
 d=\max_r\left|\bigcup_jI_{rj}\right|,
 \qquad d_{\rm coeff}=\max_{r,j}|I_{rj}|\leq d,
 \qquad M=\sum_{(r,j)}|I_{rj}|,
\]
and keep the row sparsity \(s\), coefficient bound \(B\), and witness
height \(H\).

\begin{proposition}[Multi-bit coefficients]
\label{prop:frontier-multibit}
The single-flip average obeys
\begin{equation}
 \norm{A_{\varsigma}\bar x-b}_2
 \leq\frac{2BH\sqrt{d_{\rm coeff}sM}}{N}
 \leq\frac{2BH\sqrt{dsM}}{N}.
 \label{eq:frontier-multibit}
\end{equation}
For fixed \(d_{\rm coeff},s\), the product bound
\(MB^2H^2=\Omega(N^2)\) still holds under the objective and output
assumptions above.  The first bound recovers the one-bit constant when every
coefficient depends on at most one bit, however many distinct bits appear in
a row.
\end{proposition}

\begin{proof}
Write \(f_{rji}=(A_{\varsigma,rj}-A_{\varsigma^{(i)},rj})x_j^{(i)}\).
It vanishes unless \(i\in I_{rj}\), and \(|f_{rji}|\leq2BH\).
Applying Cauchy--Schwarz first over the at most \(s\) dependent positions in
a row, and then over each coefficient's input set, gives
\[
 \left(\sum_j\sum_{i\in I_{rj}}f_{rji}\right)^2
 \leq s\sum_j |I_{rj}|\sum_{i\in I_{rj}}f_{rji}^2
 \leq4s d_{\rm coeff} B^2H^2\sum_j|I_{rj}|.
\]
Sum over rows and divide by \(N^2\).  The row-union estimate follows
from \(d_{\rm coeff}\leq d\).
\end{proof}

A simultaneous primal--dual version also holds.  Suppose \(A_\varsigma\)
has maximum row and column sparsities \(s_r,s_c\) and \(M\)
one-bit-dependent positions, and that the selected exact optimal triples
\((x^\xi,y^\xi,s^\xi)\) have all coordinates bounded in magnitude by \(H\)
and a common optimal value \(v\).  Split \(y=y^+-y^-\), stack the primal
feasibility and dual stationarity equations, and put
\[
 s_{\rm KKT}=\max\{s_r,2s_c+1\},
 \qquad B_{\rm KKT}=\max\{B,1\}.
\]
Each dependent position occurs once in the primal block and twice in the
split dual block.  \Cref{thm:frontier-single-flip} therefore gives
\begin{equation}
 \left\|
 \begin{bmatrix}
 A_\varsigma\bar x-b\\
 A_\varsigma^\top\bar y+\bar s-c
 \end{bmatrix}\right\|_2
 \leq\frac{2B_{\rm KKT}H\sqrt{3s_{\rm KKT}M}}{N}.
 \label{eq:frontier-kkt-stack}
\end{equation}
By linearity \(c^\top\bar x=b^\top\bar y=v\), so the algebraic duality gap
is exactly zero.  This covers output requirements combining feasibility, stationarity,
and gap requirements; the mixture does not preserve coordinatewise
complementarity.

The mixture proof uses only convexity of the domain.  Let \(\mathcal K\) be
a fixed convex set or cone, and suppose that for every selected input
\(\xi\in\{\varsigma,\varsigma^{(1)},\ldots,\varsigma^{(N)}\}\) there is a
witness \(x^\xi\in\mathcal K\) with
\(A_\xi x^\xi=b\), \(\norm{x^\xi}_\infty\leq H\) in one fixed basis, and
\(c^\top x^\xi=v\) for one public objective and common value.  If, in
addition, a fixed parity-neutral linear map satisfies
\(Lx^\xi=p(\xi)y\) for a fixed nonzero \(y\), then the half-base mixture
\[
 \widehat x=\frac12x^\varsigma+
              \frac1{2N}\sum_i x^{\varsigma^{(i)}}
\]
belongs to \(\mathcal K\), has the common objective, satisfies
\(L\widehat x=0\), and has half the residual in
\eqref{eq:frontier-multibit}.  For a sparse SDP this statement applies after
fixing one \(\operatorname{svec}\) identification, with all sparsity,
locality, coefficient, and coordinate bounds measured in that basis.  By
itself it gives no conclusion about normalized states, rank, determinant
neighborhoods, or complementarity.

A mixture of central points stays near centrality only under an explicit
stability hypothesis.  Keep the preceding row, column, one-bit-locality, and
height bounds, and suppose the selected neighboring triples are exact
central points for a common \(\mu>0\), so \(x_j^{(i)}s_j^{(i)}=\mu\).  Let
\(\bar x,\bar s\) be their averages.  For a fixed coordinate \(j\), write
\(t_i=x_j^{(i)}\).  Symmetrizing over pairs \(i,k\) gives the exact
identity
\begin{equation}
 \boxed{
 \frac{\bar x_j\bar s_j}{\mu}-1
 =\frac1{2N^2}\sum_{i,k=1}^N
   \frac{(t_i-t_k)^2}{t_it_k}.}
 \label{eq:frontier-multiplicative-variance}
\end{equation}
In particular, the excess is nonnegative and vanishes exactly when the
neighboring \(j\)-th primal coordinates agree.  If, for every coordinate
\(j\),
\(\max_i x_j^{(i)}/\min_i x_j^{(i)}\leq R\), where \(R\geq1\), then
\begin{equation}
 \mu\leq\bar x_j\bar s_j\leq K(R)\mu,
 \qquad K(R)=\frac{(R+1)^2}{4R}
 =1+\frac{(R-1)^2}{4R}.
 \label{eq:frontier-kantorovich}
\end{equation}
Indeed, for \(\alpha\leq t_i\leq\beta\),
\(t_i+\alpha\beta/t_i\leq\alpha+\beta\).  Averaging and maximizing the
product under this linear constraint gives
\((N^{-1}\sum_it_i)(N^{-1}\sum_it_i^{-1})
\leq(\alpha+\beta)^2/(4\alpha\beta)\), the scalar Kantorovich inequality
\cite{kantorovich1948-functional-analysis-applied-mathematics}.

\begin{proposition}[Central-mixture soundness dichotomy]
\label{prop:frontier-central-dichotomy}
Assume separately that the output is wrong for the base input uniformly
over mixtures, in one of two forms.  Write
\(Z^{(i)}=(x^{(i)},y^{(i)},s^{(i)})\).  Either a fixed affine
decoder \(\Phi\) satisfies
\(p(\varsigma)\Phi(Z^{(i)})\leq-\gamma<0\) for every \(i\), or the
primal components have the public pair decoder of
\eqref{eq:frontier-pair-decoder}, with
\(-p(\varsigma)(u_\ell^{(i)}-v_\ell^{(i)})\geq a\) for every \(i,\ell\)
and \(La^2=\Omega(nH^2)\).  In the first case the decoded object is the
stacked triple \(Z_w\); in the second it is the normalized primal state
\(\ket{x_w}=x_w/\norm{x_w}_2\).

Let \(w\) be any probability vector and set
\(Z_w=(x_w,y_w,s_w)=\sum_iw_i(x^{(i)},y^{(i)},s^{(i)})\).  With \(M_i\) the
one-bit incidence of bit \(i\), and with \(t_{ij}=x_j^{(i)}\), the stacked
residual and the centrality defect obey
\begin{align}
 \left\|
 \begin{bmatrix}A_\varsigma x_w-b\\
 A_\varsigma^\top y_w+s_w-c\end{bmatrix}\right\|_2^2
 &\leq12s_{\rm KKT}B_{\rm KKT}^2H^2\sum_iM_iw_i^2,
 \label{eq:frontier-weighted-kkt}\\
 \delta_j(w):=\frac{(x_w)_j(s_w)_j}{\mu}-1
 &=\frac12\sum_{i,k}w_iw_k
   \frac{(t_{ij}-t_{kj})^2}{t_{ij}t_{kj}}.
 \label{eq:frontier-weighted-centrality}
\end{align}
The weighted residual estimate admits the sharper constant
\begin{equation}
 \left\|
 \begin{bmatrix}A_\varsigma x_w-b\\
 A_\varsigma^\top y_w+s_w-c\end{bmatrix}\right\|_2^2
 \leq4(s_r+s_c)B^2H^2\sum_iM_iw_i^2.
 \label{eq:frontier-direct-kkt}
\end{equation}
Apply the weighted matrix-change estimate separately to the primal block
and to the transposed coefficient changes in stationarity.  The latter
allows signed multipliers and has the same incidence profile,
with column sparsity \(s_c\).  Adding the squared bounds proves
\eqref{eq:frontier-direct-kkt}; the earlier bound follows by increasing
the constant.  Thus the soundness inequalities below also hold with
\(12s_{\rm KKT}B_{\rm KKT}^2\) replaced by \(4(s_r+s_c)B^2\).
The algebraic objective difference, primal minus dual objective, remains
exactly \(n\mu\).  This follows from feasibility, stationarity, and
complementarity of each neighboring triple; the neighbors need not share
their individual primal or dual objective values.  An output requirement asking for
absolute stacked residual at most \(\varepsilon\) and an \(\ell_\infty\)
neighborhood of width \(\theta\) around the common \(\mu\) is therefore
sound only if, for every \(w\),
\begin{equation}
 12s_{\rm KKT}B_{\rm KKT}^2H^2\sum_iM_iw_i^2>\varepsilon^2
 \quad\text{or}\quad
 \max_j\delta_j(w)>\theta.
 \label{eq:frontier-central-dichotomy}
\end{equation}
Here \(\varepsilon,\theta\geq0\).  For the point-centered convention,
\(\mu_w=x_w^\top s_w/n\) with \(n>0\), replace the second term by the exact
width
\[
 \frac{\max_j|\delta_j-\bar\delta|}{1+\bar\delta},
 \qquad \bar\delta=\frac1n\sum_j\delta_j.
\]
For uniform weights, suppose for every coordinate \(j\) that
\[
 \frac{\max_i t_{ij}}{\min_i t_{ij}}\leq R.
\]
Then the width is at most \(K(R)-1\) in either convention, so soundness
requires the simpler dichotomy
\begin{equation}
 \frac{12s_{\rm KKT}B_{\rm KKT}^2H^2M}{N^2}>\varepsilon^2
 \quad\text{or}\quad
 \frac{(R-1)^2}{4R}>\theta.
 \label{eq:frontier-central-uniform}
\end{equation}
\end{proposition}

\begin{proof}
The stacked mixture estimate is the single-flip argument applied to primal
feasibility and stationarity.  Pairwise symmetrization gives the displayed
formula for \(\delta_j(w)\), and linearity preserves the algebraic objective
difference.  The affine or pair-margin hypothesis makes the decoded object wrong for the
base input.  If neither strict inequality in
\eqref{eq:frontier-central-dichotomy} held, this mixture would satisfy
every condition of the output requirement.  The coordinatewise Kantorovich
bound gives \eqref{eq:frontier-central-uniform} for uniform weights.
\end{proof}

The stability hypothesis is necessary.  Fix \(0<\mu<1\), let
\(b=\mu e,c=0\), and let bit \(j\) choose
\(a_j\in\{\mu,1\}\) in the diagonal LP
\(\operatorname{Diag}(a)x=b\).  Its exact \(\mu\)-central point is
\[
 x_j=\mu/a_j,\qquad s_j=a_j,\qquad y_j=-1.
\]
At the all-\(\mu\) base input, each single-flip neighbor changes one
coordinate from \((x,s)=(1,\mu)\) to \((\mu,1)\).  Thus every coordinate of
the average has product
\[
 \bar x_j\bar s_j
 =\frac{N-1+\mu}{N}\frac{1+(N-1)\mu}{N},
\]
which is \(\Theta(1/N)\), rather than \(\Theta(\mu)\), when
\(\mu\ll1/N\), although the primal and stationarity residuals vanish as
\(O(N^{-1/2})\).  Bounded coefficients and bounded central coordinates alone
therefore do not keep central mixtures near the common \(\mu\).  This
particular uniform mixture has the same complementarity
product in every coordinate, so its point-centered width is exactly zero.
Its algebraic objective difference is still \(N\mu\), whereas its
complementarity sum is \(N\bar x_j\bar s_j\); the two need not agree at
an infeasible mixture.

A bounded-coordinate counterexample also exists for the point-centered
convention.  Append \(N\) public diagonal coordinates with
coefficient \(\mu\) and central coordinates \((x,s,y)=(1,\mu,-1)\).
The moving coordinates share the product
\(q=(N-1+\mu)(1+(N-1)\mu)/N^2\), while the public coordinates have
product \(\mu\).  The point-centered mean is \((q+\mu)/2\), and
every coordinate has relative deviation \((q-\mu)/(q+\mu)\).
For \(N\geq8\) and \(\mu=N^{-2}\), the exact values satisfy
\[
 \left\|\begin{bmatrix}A\bar x-b\\A^\top\bar y+\bar s-c\end{bmatrix}\right\|_2^2
 =\frac{(1+\mu^2)(1-\mu)^2}{N}\leq\frac2N,
 \qquad \frac{q-\mu}{q+\mu}\geq\frac13.
\]
The unpadded example has the same residual and
\(\max_j\delta_j\geq N/8\) for \(\mu=N^{-2}\), \(N\geq2\).
Both constructions have row and column sparsity one, coefficient and
coordinate magnitudes at most one, and exactly \(N\) dependent positions.
Thus these bounds and a vanishing residual do not imply a point-centered
width below any fixed \(\theta<1/3\).

\paragraph{Formal verification.}
The Lean development in \path{formal/QipmFormal/Mixture/} verifies the
residual and centrality results of these two subsections from the actual
finite matrices and neighboring witnesses.  It covers weighted and
sensitive-subset mixtures, the stronger coefficientwise multibit bound,
RHS and slack augmentation, the split KKT construction, the direct KKT
bound, harmonic incidence weights, affine and normalized pair decoders,
the scalar variance and Kantorovich identities, both neighborhood
conventions, and the explicit counterexamples.  The soundness hypotheses
are stated as conditions on candidate outputs; the proofs derive the
residual, objective, and decoder properties of the mixture.  The development
uses Euclidean sums of squares and relies only on Lean's standard axioms
\texttt{propext}, \texttt{Classical.choice}, and \texttt{Quot.sound}.
The claim map and reproduction command are in \path{formal/MIXTURE.md}.
The sharpness constructions in Table~\ref{tab:frontier-currencies},
the electrical bounds below, and the quantum query reductions are not
covered by this verification.  The noncommutative extension has a separate
Lean development and claim map in \path{formal/SDP_MIXTURE.md}; see
\cref{sec:structural-sdp-mixtures}.

\subsection{Electrical-resistance bounds}
\label{sec:frontier-electrical}

The graph argument constructs, within the base instance, a point that
erases the outputs and a point that flips them, and shows the geometry behind the mixture bound: disjoint input
cuts force a large effective resistance.  Let
\(G=(V,E)\) be a connected oriented multigraph with root \(r\), nonempty
output set \(S\), public heights \(h_v>0\) with \(h_r=1\), and propagation
rows
\begin{equation}
 \lambda_e(d_v-g_ea_e(\varsigma)d_u)=0,
 \qquad g_e=h_v/h_u,
 \label{eq:frontier-signed-edge}
\end{equation}
where \(e=(u,v)\), \(\lambda_e\neq0\) is a public scale, \(a_e\) is either a
fixed sign or a fixed sign times one input bit, and the rows are
cycle-consistent.  Thus there are signs \(\tau_v\) with
\(\tau_r=1\) and \(\tau_v=a_e\tau_u\).  Assume separately that the
designated outputs carry the parity:
\begin{equation}
 \tau_s(\varsigma)=\epsilon_s p(\varsigma)
 \quad(s\in S),\qquad \epsilon_s\in\{\pm1\}\text{ public}.
 \label{eq:frontier-cut-output-parity}
\end{equation}
Embed each \(d_v\) using the capped pair of
\cref{sec:lp-foundation}: write \(d_v=u_v-v_v\) and
\(q_v=u_v+v_v\), fix a reference variable \(h_v^{\rm var}=h_v\), impose
\(q_v+t_v-2h_v^{\rm var}=0\), and use the local objective
\(2u_v+2v_v+h_v^{\rm var}+t_v\).  Put
\[
 L=\max_{e=(u,v)}|\lambda_e|h_v\leq BH,
 \qquad M=|E|,
\]
where the inequality assumes \(0<|\lambda_e|\leq B\) and \(h_v\leq H\).

\begin{theorem}[Cut--resistance bound]
\label{thm:frontier-cut-resistance}
There is a nonnegative, objective-exact point that satisfies every root,
reference, and cap row exactly, erases all output differences, and obeys
\begin{equation}
 \norm{Ax-b}_2^2\leq\frac{ML^2}{N^2}
 \leq\frac{MB^2H^2}{N^2}.
 \label{eq:frontier-cut-residual}
\end{equation}
There is also a point with every output sign flipped, whose squared residual
is at most four times this bound.  If the two unit anchors are the only
nonzero RHS entries, so \(\norm b_2=\sqrt2\), any relative-residual output requirement
with threshold \(\eta\) that rules out the erased point requires
\begin{equation}
 MB^2H^2>2\eta^2N^2.
 \label{eq:frontier-cut-product}
\end{equation}
If each edge label contains at most \(k\) raw signs, the residual bound
\eqref{eq:frontier-cut-residual} for the erased point acquires a factor
\(k^2\), and soundness requires \(Mk^2B^2H^2=\Omega(N^2)\).
\end{theorem}

\begin{proof}
For bit \(i\), let \(E_i\) be the edges whose label changes when bit
\(i\) flips.  By cycle consistency every cycle meets \(E_i\) evenly, so signed-graph switching identifies \(E_i\) with a cut
separating the root from every output
\cite{harary1953-balance-signed-graph,
zaslavsky2013-matrices-signed-simple-graphs}.  One-bit locality makes the
\(N\) cuts edge-disjoint.

Give edge \(e=(u,v)\) conductance
\(c_e=(\lambda_eh_v)^2\leq L^2\), and merge all outputs into one node.  For any
unit root-to-output flow \(f\), one unit crosses every cut.  Weighted
Cauchy--Schwarz gives
\[
 \sum_{e\in E_i}\frac{f_e^2}{c_e}\geq\frac1{C_i},
 \qquad C_i=\sum_{e\in E_i}c_e.
\]
Because the cuts are disjoint, Thomson's principle and a second application
of Cauchy--Schwarz give the weighted Nash--Williams bound
\begin{equation}
 \mathcal R_{r,S}\geq\sum_i\frac1{C_i}
 \geq\frac{N^2}{\sum_iC_i}
 \geq\frac{N^2}{ML^2}.
 \label{eq:frontier-nash-williams}
\end{equation}
We include this derivation to show where the oracle's disjoint input cuts
enter; the resistance inequality itself is classical
\cite{nashwilliams1959-random-walk-electric-currents}.

Let \(w\in[0,1]^V\) be the harmonic voltage with \(w_r=1\) and \(w_s=0\)
on every output.  By Dirichlet's principle and
\eqref{eq:frontier-nash-williams}, its energy is at most \(ML^2/N^2\).  Set
\[
 d_v=\tau_vh_vw_v,\quad q_v=h_v^{\rm var}=t_v=h_v,
 \quad u_v=(h_v+d_v)/2,\quad v_v=(h_v-d_v)/2.
\]
The maximum principle gives nonnegativity.  Since every $\lambda_e$ is
nonzero and the graph is connected, every feasible point has $|d_v|=h_v$,
hence $q_v=u_v+v_v\ge|d_v|=h_v$; the local objective reduces to $q_v+3h_v$,
so this point, which has $q_v=h_v$, is objective-exact.  All nonpropagation
rows hold exactly, while edge \(e\) has residual
\(\lambda_e\tau_vh_v(w_v-w_u)\), whose squared sum is the Dirichlet
energy.  Every output has \(d_s=0\).
Replacing the output boundary value zero by \(-1\) doubles every voltage
difference and proves the statement about the sign-flipped point.

For \(k\)-local labels, an edge lies in at most \(k\) of the bit cuts.  Hence
\(\mathcal R_{r,S}\geq k^{-1}\sum_iC_i^{-1}\), while
\(\sum_iC_i\leq kML^2\).  This gives
\(\mathcal R_{r,S}\geq N^2/(k^2ML^2)\) and the stated factor.
\end{proof}

By \cref{thm:frontier-cut-resistance}, the gain--plateau family is
optimal among cycle-consistent signed graphs: linear size and bounded row
scales force \(H=\Omega(\sqrt N)\).  The other two rows of
Table~\ref{tab:frontier-currencies} attain the remaining extremes.
Adaptivity does not rescue the pointwise output requirement, because the
harmonic point exists at every requested stage.  However, the theorem gives
no unconditional lower bound on fresh queries per iteration:
after all \(N\) fixed signs are cached, later raw-query cost may be zero
under \cref{def:amortized-accounting}.

The resistance argument extends to rows with several wires.
Give wire \(v\) a public height \(h_v\in(0,H]\) and intended value
\(h_v\chi_{J_v}(\varsigma)\), where
\(\chi_J=\prod_{i\in J}\varsigma_i\).  Let each of \(m\) homogeneous rows
have arity at most \(q\):
\begin{equation}
 \sum_{v\in V_a}\alpha_{av}\chi_{I_{av}}(\varsigma)d_v=0,
 \qquad |\alpha_{av}|\leq A,\qquad |I_{av}|\leq k.
 \label{eq:frontier-walsh-row}
\end{equation}
Roots have \(J_v=\varnothing\), outputs have \(J_v=[N]\), and every row
holds on the intended wires for every input.  We may assume that the
comparison graph in the proof connects the roots to the outputs; otherwise the output character is unconstrained and can be erased with zero
propagation residual.
For the objective-exact conclusion, assume also that the intended
full-height assignment is optimal for the capped LP.  It suffices that
every exactly feasible assignment has \(|d_v|\geq h_v\) on every wire,
since nonnegativity then forces \(q_v\geq h_v\).

\begin{theorem}[Bounded-arity Walsh-row bound]
\label{thm:frontier-walsh}
The capped-pair embedding has a nonnegative, objective-exact point that
satisfies the anchors, references, and caps, erases every output, and has
\begin{equation}
 \norm r_2^2\leq
 \frac{4k^2q^3mA^2H^2}{N^2}.
 \label{eq:frontier-walsh-residual}
\end{equation}
Consequently, for fixed \(k,q\), soundness at a constant absolute residual
threshold requires \(mA^2H^2=\Omega(N^2)\).
\end{theorem}

\begin{proof}
Substituting the intended wires into row \(a\) gives
\(\sum_v\beta_{av}\chi_{K_{av}}=0\), where
\(\beta_{av}=\alpha_{av}h_v\) and
\(K_{av}=I_{av}\mathbin\triangle J_v\).  Since distinct Walsh characters
are linearly independent, \(\sum_{v\in G_{a,K}}\beta_{av}=0\) within each
equal-character group.  Choose one pivot per group and join it to the other
group members by comparison edges of conductance
\(q^2\beta_{av}^2\).  There are at most
\(qm\) edges, each with conductance at most \(q^2A^2H^2\).  For arbitrary
interpolation values \(w_v\), a character group with pivot \(p\) has
residual
\[
 \sum_{v\in G}\beta_{av}w_v
 =\sum_{v\in G\setminus\{p\}}\beta_{av}(w_v-w_p).
\]
Applying Cauchy--Schwarz first over the at most \(q\) terms of a group and
then over the at most \(q\) groups of a row gives
\[
 \norm r_2^2\leq
 \sum_ec_e(w_{h(e)}-w_{t(e)})^2.
\]

The endpoints of a comparison edge satisfy
\(J_u\triangle J_v=I_{au}\triangle I_{av}\), so the edge crosses at most
\(2k\) of the vertex cuts \(\{v:i\in J_v\}\).  The bounded-overlap form of
the Nash--Williams calculation proved above gives
\[
 \mathcal R_{\rm eff}
 \geq\frac{N^2}{4k^2q^3mA^2H^2}.
\]
Take the harmonic interpolation equal to one at the roots and zero at the
outputs.  Its energy gives \eqref{eq:frontier-walsh-residual}; the maximum
principle and the capped-pair construction from the preceding proof give
nonnegativity and exact auxiliary rows.  The erased point has
\(q_v=h_v\) on every wire, hence the local objective of the intended
assignment, which is optimal by hypothesis.
\end{proof}

Formulations built from exact local XOR hulls have three further concrete
defects.  If a set \(U\) of unfixed variables meets every parity check in zero or at
least two positions, then the point that sets \(x_j=1/2\) on \(U\) and keeps
a satisfying Boolean codeword elsewhere lies in every local parity polytope.
Indeed, for a check meeting \(U\) in \(r\geq2\) variables, this point is the
uniform average of the \(2^{r-1}\) assignments with the required residual
parity.  This stopping-set pseudocodeword can make an output exactly
unbiased.  If the checks instead peel along a dependency chain of length
\(D\), linearly interpolating the gauged wire error from zero to one
violates one opposing facet at each gate by \(1/D\), for total Euclidean
violation \(1/\sqrt D\).  Finally, every satisfying truth-table assignment
is a vertex of its exact gate hull, so some standard-form slacks vanish; a
Boolean circuit forced by its inputs is not strictly primal feasible.
Thickening the hulls restores positive slacks only by admitting fractional
drift.  These facts explain the LP-decoding analogy but are not a universal
lower bound for all inequality extended formulations.

\subsection{Multiplier conservation}
\label{sec:frontier-multiplier}

Residual robustness constrains primal witnesses.  Complete central and
Newton outputs face a second constraint: conservation of the dual
multipliers.  In a signed-copy graph \(G=(V\cup T,E)\), vertices in \(V\) carry variables with
\(d_v=x_v^+-x_v^-\) and \(q_v=x_v^++x_v^-\), and vertices in \(T\) are fixed
terminals used for anchor rows.  A variable--variable edge \(e=(u,v)\)
represents \(d_v-a_ed_u=0\).  A terminal edge is oriented from its terminal
to its variable endpoint, with the fixed terminal term moved to the RHS.
Let \(B_a\) and \(B\) be the signed and the ordinary oriented incidence
matrices, with all terminal columns omitted.  For
\(S\subseteq V\), the boundary \(\partial S\) consists of the edges to
\(V\setminus S\) and all terminal edges incident to \(S\).  Each normalized
edge equation has a multiplier \(y_e\); an actual LP row may be multiplied
by a public nonzero scale \(r_e\ne0\), in which case its scale-invariant row
force is \(f_e=r_ey_e\).
Assume exact primal--dual centrality, \(d_v\ne0\), a pair-symmetric objective,
and pair-symmetric columns for all equality constraints outside the graph.

\begin{theorem}[Gauge flow and isoperimetric law]
\label{thm:frontier-gauge-flow}
Let \(\tau_v=\operatorname{sgn}(d_v)\).  After the edge gauge
\(z_e=\tau_{h(e)}y_e\), signed stationarity becomes the positive-divergence
law
\begin{equation}
 B^\top z=g,
 \qquad g_v=\frac{2\mu|d_v|}{q_v^2-d_v^2}>0.
 \label{eq:frontier-gauge-flow}
\end{equation}
For every vertex set \(S\) with nonempty boundary,
\begin{equation}
 \sum_{e\in\partial S}|y_e|\geq G_S,
 \quad
 \norm{y_{\partial S}}_\infty\geq\frac{G_S}{|\partial S|},
 \quad
 \norm{y_{\partial S}}_2\geq\frac{G_S}{\sqrt{|\partial S|}},
 \quad G_S=\sum_{v\in S}g_v.
 \label{eq:frontier-cut-flow}
\end{equation}
If \(S\ne\varnothing\) but \(\partial S=\varnothing\), no such positive
central point exists.
If \(|d_v|\geq a>0\) and \(q_v\leq Q\) on \(S\), then
\begin{equation}
 \norm{y_{\partial S}}_\infty
 \geq\frac{2\mu a}{Q^2-a^2}\frac{|S|}{|\partial S|},
 \qquad
 \norm{y_{\partial S}}_2
 \geq\frac{2\mu a}{Q^2-a^2}\frac{|S|}{\sqrt{|\partial S|}}.
 \label{eq:frontier-isoperimetric}
\end{equation}
For a row written as \(r_e(d_v-a_ed_u)=0\), all statements hold for the
row force \(f_e=r_ey_e\); hence, if \(|r_e|\leq B_{\rm row}\), the lower
bounds apply to \(B_{\rm row}\norm y_\infty\), not to the multiplier
alone, which depends on the row normalization.
\end{theorem}

\begin{proof}
Subtract the two stationarity equations of each pair:
\[
 2(B_a^\top y)_v+s_v^+-s_v^-=0.
\]
Since \(x_v^\pm=(q_v\pm d_v)/2\) and
\(s_v^\pm=\mu/x_v^\pm\), this becomes
\((B_a^\top y)_v=2\mu d_v/(q_v^2-d_v^2)\).  Feasibility implies
\(a_e=\tau_u\tau_v\) on variable--variable edges; terminal rows have the
same single variable coefficient in \(B_a\) and \(B\).  Switching the signed
incidence matrix therefore gives
\eqref{eq:frontier-gauge-flow}.  Summing over \(S\) cancels the internal
edges; the triangle inequality and Cauchy--Schwarz give
\eqref{eq:frontier-cut-flow}.  Positivity gives \(q_v>|d_v|\), and the
margin hypotheses bound every \(g_v\) below by
\(2\mu a/(Q^2-a^2)\).  A diagonal row scaling enters stationarity only
through \(r_ey_e\), which proves the final statement.
\end{proof}

Pair symmetry is not needed to compare two inputs if the pair-antisymmetric
source \(h_v\) below is the same vector, in the original pair coordinates,
for both.  Let two exact central points, denoted \(+\) and \(-\), have
row forces \(f_e^\pm=r_ey_e^\pm\), and suppose their subtracted stationarity
equations are
\[
 (B_{a^\pm}^\top f^\pm)_v
 =h_v+\frac{2\mu_\pm d_v^\pm}
              {(q_v^\pm)^2-(d_v^\pm)^2}.
\]
If \(d_v^+>0>d_v^-\) on \(S\), then switching the two inputs, adding their
equations, and summing over \(S\) cancels both the common source and every
internal edge.  Hence
\begin{equation}
 \sum_{e\in\partial S}(|f_e^+|+|f_e^-|)
 \geq\sum_{v\in S}(g_v^++g_v^-),
 \qquad
 g_v^\pm=\frac{2\mu_\pm|d_v^\pm|}
 {(q_v^\pm)^2-(d_v^\pm)^2}.
 \label{eq:frontier-two-input}
\end{equation}
If \(\mu_\pm\geq\mu_0\), \(|d_v^\pm|\geq a\), \(q_v^\pm\leq Q\), and
\(|r_e|\leq B_{\rm row}\), this gives
\begin{equation}
 \max_{p\in\{+,-\}}\norm{y_{\partial S}^p}_\infty
 \geq\frac{2\mu_0a}{Q^2-a^2}
       \frac{|S|}{B_{\rm row}|\partial S|}.
 \label{eq:frontier-two-input-norm}
\end{equation}
Thus a public asymmetric objective can cancel the demand \(g_v\) for one
sign pattern of \(d\) on \(S\), but not for both when \(\partial S\) is
small.  Coefficients of another public asymmetric constraint do not suffice
unless their contribution through the multipliers is also the same vector
\(h\) for both inputs.

The identity for the work \(b_a^\top y_a\) does not need a graph.  Let
\(Cd=b_a\) be an arbitrary pair-antisymmetric equality block, with columns
\(C,-C\), and let all other local data be pair-symmetric.

\begin{theorem}[Central and Newton work identities]
\label{thm:frontier-work}
At an exact central point,
\begin{equation}
 \boxed{b_a^\top y_a
 =\sum_v\frac{2\mu d_v^2}{q_v^2-d_v^2}\geq0.}
 \label{eq:frontier-central-work}
\end{equation}
The work is strictly positive exactly when \(d\ne0\).  At a pair-symmetric
Newton start with pair-symmetric dual and complementarity right-hand sides,
pair-symmetric other equality columns,
\(x_v^{0,+}=x_v^{0,-}=\xi_v\), and
\(s_v^{0,+}=s_v^{0,-}=\zeta_v\), the antisymmetric block satisfies
\begin{equation}
 C_\varsigma\Delta d=r_a,
 \qquad C_\varsigma^\top\alpha=D\Delta d,
 \qquad D=\operatorname{Diag}\!\left(\frac{\zeta_v}{2\xi_v}\right),
 \quad
 \boxed{r_a^\top\alpha=\Delta d^\top D\Delta d.}
 \label{eq:frontier-newton-work}
\end{equation}
Both pairings are invariant under
\((C,b_a,y_a)\mapsto(UC,Ub_a,U^{-\top}y_a)\), and likewise for
\((C_\varsigma,r_a,\alpha)\).  If \(D\succeq\rho I\),
\(|\Delta d_v|\geq a\) on \(L\) coordinates, and \(\norm{r_a}_1\leq R\),
then \(\norm\alpha_\infty\geq\rho a^2L/R\).
Likewise, if \(|d_v|\geq a\) and \(q_v\leq Q\) on \(L\) coordinates,
then
\begin{equation}
 \norm{y_a}_\infty\geq
 \frac{2\mu a^2L}{(Q^2-a^2)\norm{b_a}_1},
 \qquad
 \norm{y_a}_2\geq
 \frac{2\mu a^2L}{(Q^2-a^2)\norm{b_a}_2}.
 \label{eq:frontier-central-work-norms}
\end{equation}
A zero denominator with nonzero work means that no such point exists.
\end{theorem}

\begin{proof}
Pair subtraction gives
\(C^\top y_a=2\mu d/(q^2-d^2)\) coordinatewise.  Taking the inner product
with \(d\) and using \(Cd=b_a\) proves
\eqref{eq:frontier-central-work}.  At the Newton start, subtracting the dual
and complementarity equations and eliminating the slack difference gives
\(C_\varsigma^\top\alpha=D\Delta d\); primal antisymmetric feasibility gives
\(C_\varsigma\Delta d=r_a\).  Taking inner products proves
\eqref{eq:frontier-newton-work}, and H\"older's inequality gives the norm
bounds.  The contragredient transformations preserve both scalar pairings.
\end{proof}

For the bounded secant start of \cref{sec:lp-dual-homogeneity},
\(D=(2/9)I\), \(\Delta d_v=\tau_v\), and the antisymmetric RHS at the root
is a unit vector.  Equation~\eqref{eq:frontier-newton-work} therefore gives
the root multiplier \(2P/9\).  Scaling all dual data by \(\lambda\)
changes this to \(2\lambda P/9\); bounded full Newton directions require
\(\lambda=O(1/P)\), exactly the inverse dual scale used in
\cref{thm:lp-dual-homogeneous-full}.  \Cref{sec:lp-beyond-pairs}
states the corresponding unrestricted Newton virtual-work identity.  We use
the paired forms here because they expose the gadget resources;
\cref{sec:structural-tools} places them in the general
involution/symmetry framework.

The cut laws also hold directly for Newton directions.  Under a
pair-symmetric start, pair-symmetric Newton right-hand sides, and homogeneous
signed-copy primal equations, switching by
\(\operatorname{sgn}(\Delta d_v)\) gives
\begin{equation}
 B^\top z=g^{\rm N},
 \qquad g_v^{\rm N}=\frac{\zeta_v}{2\xi_v}|\Delta d_v|.
 \label{eq:frontier-newton-flow}
\end{equation}
The proof is the Newton half of \cref{thm:frontier-work} followed by
the cut summation in \cref{thm:frontier-gauge-flow}.  In particular,
if \(\zeta_v/\xi_v\geq r\) and \(|\Delta d_v|\geq a\) on \(S\), then
\[
 \norm{\Delta y_{\partial S}}_\infty
 \geq\frac{ra}{2}\frac{|S|}{|\partial S|}.
\]

Now assume that terminal values, labels, and terminal Newton residuals are
public and fixed, and fix a base input such that it and all its
single-flip neighbors have exact central points with nonzero \(d_v\).  Let
\(O\) be a fixed set of \(K\) output vertices on which
\(\operatorname{sgn}d_v\) flips under every single-bit flip, and define
\[
 S_j=\{v:\operatorname{sgn}d_v(\varsigma^{(j)})
             =-\operatorname{sgn}d_v(\varsigma)\}.
\]
Let \(M_j\) count the signed rows that change under flip \(j\), and put
\(M_{\rm dep}=\sum_jM_j\).  Feasibility on an internal edge gives
\(a_e=\operatorname{sgn}(d_u)\operatorname{sgn}(d_v)\), so its label changes
exactly when one endpoint is in \(S_j\).  Public terminal signs act as fixed
outside vertices, so the same boundary count includes terminal edges.
Hence
\(|\partial S_j|\leq M_j\), with equality when rows rather than expanded
coefficient positions are counted, and \(O\subseteq S_j\).  The base
multiplier must satisfy all \(N\) cut inequalities simultaneously.  Assume
explicitly that the base central point has
\(|d_v|\geq a=\Omega(1)\) and \(q_v\leq Q=O(1)\) on \(O\).  Applying this
central margin only on \(O\) and summing over \(j\) proves
\begin{equation}
 M_{\rm dep}B_{\rm row}\norm{y}_\infty
 \geq\frac{2\mu a}{Q^2-a^2}KN.
 \label{eq:frontier-central-incidence}
\end{equation}
For Newton directions, suppose instead that the base and every single-flip
system satisfy the hypotheses of \eqref{eq:frontier-newton-flow}, define
\(S_j\) from the signs of \(\Delta d\), and assume that every output
direction flips with \(|\Delta d_v|\geq a=\Omega(1)\) on \(O\).  The same
argument gives
\begin{equation}
 M_{\rm dep}B_{\rm row}\norm{\Delta y}_\infty
 \geq\rho aKN,
 \qquad \rho=\min_{v\in O}\frac{s_v^0}{2x_v^0}.
 \label{eq:frontier-newton-incidence}
\end{equation}
Thus \(M_{\rm dep}=\Theta(N)\), \(K=\Theta(N)\), bounded row scale, and
bounded multipliers, together with the explicit constant margin
\(a=\Omega(1)\) and central range \(Q=O(1)\), force
\(\mu=O(1/N)\) in the central system and \(\rho=O(1/N)\) in the Newton
system.  Equation
\eqref{eq:frontier-newton-incidence} is the general conservation law behind
the multiplier lower bound beyond signed pairs in \cref{sec:lp-beyond-pairs}.
Every signed-copy gadget must make the input incidence, row scale,
multiplier size, or inverse dual scale large.  Branching and redundant copy
rows only move the cost into the relevant boundary count.

Pair-asymmetric local data do not provide a cheap universal cancellation.
Suppose a prescribed direct source \(h_v(\varsigma)\) changes central
stationarity to
\[
 (B_a^\top f)_v=
 \frac{2\mu d_v}{q_v^2-d_v^2}+h_v(\varsigma),
 \qquad f_e=r_ey_e.
\]
On a fixed set \(O\) of \(K\) output vertices, assume
\(\operatorname{sgn}(d_v)=\kappa_vp(\varsigma)\),
\(|d_v|\geq a\), and \(q_v\leq Q\) for every input.  With the uniform
Fourier convention
\(\widehat h_v([N])=2^{-N}\sum_\varsigma p(\varsigma)h_v(\varsigma)\), put
\(\varepsilon_O=K^{-1}\sum_{v\in O}|\widehat h_v([N])|\).

\begin{theorem}[Fourier obstruction to direct cancellation]
\label{thm:frontier-fourier-source}
For some input \(\varsigma^*\),
\begin{equation}
 \sum_{e\in\partial O}|r_ey_e(\varsigma^*)|
 \geq K\left(\frac{2\mu a}{Q^2-a^2}-\varepsilon_O\right).
 \label{eq:frontier-fourier-source}
\end{equation}
In particular, bounded-bit juntas that depend on strictly fewer than \(N\)
input bits, and all sources of Fourier degree less than \(N\), have zero
full-parity coefficient and cannot cancel the demand for all inputs.  The
bound is tight: the direct source
\(h_v=-\kappa_vp(\varsigma)g_0\) cancels a constant demand \(g_0\) exactly.
\end{theorem}

\begin{proof}
After switching, sum the stationarity equations over \(O\).  The direct
contribution is
\(H(\varsigma)=\sum_{v\in O}\kappa_vp(\varsigma)h_v(\varsigma)\), whose
average over inputs is at least \(-K\varepsilon_O\).  Some \(\varsigma^*\) has
\(H(\varsigma^*)\) at least this average.  The positive central demand is at
least \(2\mu aK/(Q^2-a^2)\); internal edge forces cancel, and the triangle
inequality proves \eqref{eq:frontier-fourier-source}.
\end{proof}

There is an exact Newton analogue.  At a pair-symmetric start put
\(x_v^{0,\pm}=\xi_v\), \(s_v^{0,\pm}=\zeta_v\), and allow asymmetric dual
and complementarity right-hand sides.  Pair subtraction gives
\[
 (B_a^\top\Delta f)_v
 =\frac{\zeta_v}{2\xi_v}\Delta d_v+h_v^{\rm N}(\varsigma),
 \qquad
 h_v^{\rm N}=\frac{r_{d,v}^+-r_{d,v}^-}{2}
 -\frac{r_{c,v}^+-r_{c,v}^-}{2\xi_v}.
\]
If \(\operatorname{sgn}(\Delta d_v)=\kappa_vp(\varsigma)\),
\(|\Delta d_v|\geq a\), and \(\zeta_v/(2\xi_v)\geq\rho\) on \(O\), then
the same proof gives some input \(\varsigma^*\) with
\begin{equation}
 \sum_{e\in\partial O}|r_e\Delta y_e(\varsigma^*)|
 \geq K(\rho a-\varepsilon_O^{\rm N}),
 \qquad
 \varepsilon_O^{\rm N}=K^{-1}\sum_{v\in O}
 |\widehat h_v^{\rm N}([N])|.
 \label{eq:frontier-newton-fourier}
\end{equation}
This Fourier coefficient must be computed after division by the start
coordinates: input dependence in \(1/\xi_v\) can create full parity from
raw data that are otherwise local.  The word \emph{direct} also matters.
An unknown multiplier of another antisymmetric constraint may depend on all
input bits even when that row is local, so
\cref{thm:frontier-fourier-source} does not assume that such a
multiplier has low Fourier degree.

Switching and conservation of incidence flows are classical
\cite{harary1953-balance-signed-graph,zaslavsky2013-matrices-signed-simple-graphs};
the work identities are special cases of standard KKT virtual-work and
Schur-energy identities.  Electrical flows on interior-point central paths
are also established algorithmic tools \cite{madry2013-navigating-central-path}.
The only contribution claimed here is that the local input cuts, the
parity-sized output region, multiplier conservation, row scaling, and
incidence accounting must all be satisfied at once.  Effective resistance has a separate established role in
quantum graph algorithms
\cite{belovs2012-span-programs-st-connectivity,
jarret2018-quantum-connectivity,apers2022-quantum-speedup-for-graph-sparsification};
those works do not imply the LP residual or multiplier bounds of this
section.

\subsection{Relative-value estimation and long shortcut rows}
\label{sec:frontier-remnants}

The path family also gives a simple lower bound for value outputs.  Put
\(p_0=1\) and \(p_i=\prod_{k=1}^i\varsigma_k\).  On a single signed path,
minimize
\begin{equation}
 \frac1{N+1}\sum_{i=0}^N(u_i+v_i)
 +\frac14(u_N-v_N)
 \label{eq:frontier-relative-value-lp}
\end{equation}
subject to \(d_0=1\) and
\(d_i-\varsigma_id_{i-1}=0\).

\begin{theorem}[Relative-value gap]
\label{thm:frontier-relative-value}
In raw sparse coefficient access, estimating the optimum of
\eqref{eq:frontier-relative-value-lp} to target-relative error
\(\epsilon<1/4\) requires \(\Omega(N)\) queries.  The LP has a finite,
attained optimum, is strictly primal--dual feasible, and has a unique,
nondegenerate, strictly complementary optimal solution.  Its objective norm
is bounded by an absolute constant, and its primal log-barrier Hessian,
restricted to the feasible affine space, has condition number one along the
entire central path.
\end{theorem}

\begin{proof}
Exact feasibility fixes \(d_i=p_i\) and \(q_i=u_i+v_i\geq1\).  On this
affine space the objective is
\((N+1)^{-1}\sum_iq_i+p_N/4\), so the unique optimum has every \(q_i=1\)
and
\begin{equation}
 \operatorname{OPT}=1+\frac14p(\varsigma)
 \in\left\{\frac34,\frac54\right\}.
 \label{eq:frontier-relative-value-gap}
\end{equation}
The two target-relative error intervals are disjoint for every
\(\epsilon<1/4\); hence such value estimation requires \(\Omega(N)\) raw
queries by the parity lower bound.  Even an exactly feasible point with
objective at most \(3\operatorname{OPT}/2\) distinguishes the two cases by
a threshold between \(9/8\) and \(5/4\).  Each
sparse coefficient query is simulated by one raw sign query.  Also
\(\norm c_2^2=2/(N+1)+1/8=O(1)\).

The equality matrix is full row rank, and its public orthonormal null basis
consists of \((e_{u_i}+e_{v_i})/\sqrt2\).  Taking every \(q_i=2\) proves
strict primal feasibility.  Taking \(s_{u_i}=s_{v_i}=1/(N+1)\) gives a
positive slack vector for which \(c-s\) is the endpoint difference term; it is
orthogonal to the nullspace and hence lies in \(\operatorname{range}(A^\top)\),
proving strict dual feasibility.  At the optimum, one positive coordinate
from each pair gives a triangular basis, and every opposite coordinate has
positive reduced cost \(2/(N+1)\).  This proves nondegeneracy and strict
complementarity.

Finally, on the feasible slice the barrier objective is
\[
 \sum_{i=0}^N\left[\frac{q_i}{N+1}
 -\mu\log\frac{q_i^2-1}{4}\right]+\frac14p_N.
\]
Every central coordinate solves the same scalar equation
\((N+1)^{-1}=2\mu q/(q^2-1)\), so the Hessian in the displayed orthonormal
null basis is a positive scalar multiple of the identity.
\end{proof}

This is a value reduction with the normalization guarantees stated above,
not a query lower bound for generic LPs.

One tempting repair is a long two-coordinate shortcut row.  Suppose
\(0\leq i<j\leq N\) and a homogeneous row
\(a(\varsigma)d_i+b(\varsigma)d_j=0\) holds on every exact signed path,
with both \(a(\varsigma)\) and \(b(\varsigma)\) nonzero for every input
covered by the conclusion.  Substitution then gives
\begin{equation}
 -\frac{a(\varsigma)}{b(\varsigma)}
 =\frac{p_j}{p_i}=\prod_{k=i+1}^j\varsigma_k.
 \label{eq:frontier-shortcut-parity}
\end{equation}
Thus the row itself carries the parity of bits \(i+1,\ldots,j\),
invariant under common row rescaling, and generally costs
\(\Omega(j-i)\) raw queries to simulate.  Otherwise, replicated-path and
expander constructions are covered by the mixture, resistance, and
Walsh-row theorems above.

\subsection{Scope of the conclusions}
\label{sec:frontier-synthesis}

The parallel-path and gain--plateau constructions of \cref{sec:lp-hardness}
are optimal in their stated bounded-height and linear-size classes,
respectively.  The \(\Theta(N^{-1/2})\) residual threshold of
the all-unit family in \cref{thm:lp-unit-threshold} is the matching
threshold without amplification, and dual homogeneity attains the scaling
\(\rho=\Theta(1/N)\) that \eqref{eq:frontier-newton-incidence} forces.
None of these conclusions says that parity needs a quadratic-size LP, that
every bounded-range LP fails, or that an adaptive QIPM pays a fresh linear
query cost at every iteration.

A bounded-range construction with linear incidence must violate at least
one hypothesis used here.  The escape routes are:
\begin{enumerate}[label=(\roman*)]
 \item use an output promise that is nonlinear or otherwise not
       mixture-stable, and prove its robustness directly;
 \item require only central or Newton outputs, and allow the neighboring
       central coordinates to leave every fixed Kantorovich neighborhood
       \eqref{eq:frontier-kantorovich};
 \item use nonlocal input coefficients and charge their full raw-query
       cost;
 \item give up the common objective or common optimum, so neighbor
       mixtures need not be objective-exact; or
 \item impose integrality or another nonconvex promise, which leaves linear
       programming.
\end{enumerate}
The second route is not automatic:
\eqref{eq:frontier-central-dichotomy}--\eqref{eq:frontier-central-uniform}
show exactly how much central instability it requires.

Two problems remain open within the present scope.  First, can an arbitrary
sparse bounded-range equality LP with a nonlinear normalized-state decoder
be robust in a fixed global residual tube without nonlocal coefficient
access?  Second, can the bounded-scale full-KKT problem of
\cref{open:lp-bounded-full-kkt} be solved by a nonorthogonal
mixed-mode gadget when all input-dependent sources and recovery maps are
charged?  Together with the orthogonal-copy cut law in
\cref{sec:lp-beyond-pairs}, the theorems here rule out the natural
mixture-stable, signed-copy, constant-arity Walsh, permutation, rotation,
and flat-copy approaches; they do not settle either question.

\section{SDP Hardness Families}
\label{sec:sdp-hardness}

This section gives SDP-native lower bounds rather than embedding the LP
families of \cref{sec:lp-hardness}.  Three output levels must be kept
separate.  The first family has an intrinsic, constant gap in the optimal
value.  The second has a public Schur-complement solve but makes the requested
original-coordinate state expensive to load.  At the third level, even a
constant-normalization encoding of a canonical KKT matrix, or the first
synthesis of a tangent projector, may require linear raw-input work.  All
claims use the access and output contracts of
\cref{sec:prelim-quantum-access,sec:prelim-output-contracts}.

The geometric boundary is equally important.  The first trace-normalized
family exposes the density slice of $\Splus^3$.  In the Schur family the
$3\times3$ blocks leave only $\Splus^2$ and hence are SOCP-representable;
$4\times4$ blocks are the first members whose free Schur complement is
$\Splus^3$ and has no finite second-order-cone lift.

\subsection{An intrinsic-spectrum trace-normalized family}
\label{sec:sdp-intrinsic}

Let $n_{\rm bit}\geq1$, let
\begin{equation}
 K=16n_{\rm bit},\quad P=2K+n_{\rm bit}=33n_{\rm bit},
 \quad h=\prod_{j=1}^{n_{\rm bit}}\varsigma_j,
 \label{eq:sdp-intrinsic-parameters}
\end{equation}
and put $D_s=\diag(s,1,1)$.  There are $P$ blocks
$X_0,\ldots,X_{P-1}\in\Splus^3$.  Consecutive blocks obey
\begin{equation}
 X_i=G_iX_{i-1}G_i^T.                         \label{eq:sdp-copy}
\end{equation}
The first $K-1$ and last $K$ edges have $G_i=I$; the intervening
$n_{\rm bit}$ edges have $G_i=D_{\varsigma_j}$ in input order.  Add
$\tr X_0=1$.  If $R_0=I$ and $R_i=G_i\cdots G_1$, elimination gives exactly
\begin{equation}
 X_i=R_iZR_i^T,\qquad Z\succeq0,\qquad\tr Z=1.              \label{eq:sdp-root-lift}
\end{equation}
In the Frobenius-isometric coordinates
\begin{equation}
 \operatorname{svec}(X)=(X_{11},\sqrt2X_{12},\sqrt2X_{13},
                          X_{22},\sqrt2X_{23},X_{33}),       \label{eq:sdp-svec}
\end{equation}
congruence by $D_s$ is $\diag(1,s,s,1,1,1)$.  Each copy row
has two nonzeros, the trace row has three, and each scalar column occurs in
at most two rows.  One input sign changes only the two predecessor
coefficients in the $12$ and $13$ rows.  Ordering the copy rows along the
path proves that they are independent; their kernel is the six-dimensional
root family in \eqref{eq:sdp-root-lift}, and the trace row removes one
further dimension.

For $a<b$, write $H_{ab}=e_ae_b^T+e_be_a^T$, and set
\begin{equation}
 u=\frac1{10},\quad q=\frac KP=\frac{16}{33},\quad
 \gamma=\frac uq=\frac{33}{160}.                           \label{eq:sdp-anisotropy}
\end{equation}
Use the public costs
\begin{equation}
 C_i=\begin{cases}
  3I+uH_{23}+\gamma H_{12},&0\leq i<K,\\
  3I+uH_{23},&K\leq i<K+n_{\rm bit},\\
  3I+uH_{23}+\gamma H_{13},&K+n_{\rm bit}\leq i<P,
 \end{cases}                                                \label{eq:sdp-local-costs}
\end{equation}
and consider
\begin{equation}
 \min\left\{\frac1P\sum_{i=0}^{P-1}\inner{C_i}{X_i}:
       \eqref{eq:sdp-copy},\ \tr X_0=1,\ X_i\succeq0\right\}.
 \label{eq:sdp-intrinsic-program}
\end{equation}
Because $\sqrt{u^2+\gamma^2}<1/4$, all local costs have spectrum in
$(11/4,13/4)$.  The point $X_i=I/3$ is strictly primal feasible, and zero
equality multipliers with slacks $C_i/P$ prove strict dual feasibility.

\begin{theorem}[Intrinsic qutrit parity]
\label{thm:sdp-intrinsic}
The program \eqref{eq:sdp-intrinsic-program} has optimal values
\begin{equation}
 v_+=\frac{29}{10},\qquad v_-=\frac{14}{5}.                 \label{eq:sdp-intrinsic-values}
\end{equation}
Consequently additive-$1/25$ or target-relative-$1/100$ value estimation
requires $\Omega(P)$ raw coefficient queries.  More exactly,
\begin{equation}
 Q_{\rm add}(\epsilon)=
 \begin{cases}\lceil n_{\rm bit}/2\rceil,&\epsilon<1/20,\\0,&\epsilon\geq1/20,
 \end{cases}
 \quad
 Q_{\rm rel}(\eta)=
 \begin{cases}\lceil n_{\rm bit}/2\rceil,&\eta<1/57,\\0,&\eta\geq1/57.
 \end{cases}                                                \label{eq:sdp-value-regimes}
\end{equation}
The feasible spectrahedron has no lift over any finite product of
second-order cones.
\end{theorem}

\begin{proof}
The $H_{23}$ term is invariant under every $R_i$.  The first anchor remains
in the root frame and the last acquires the final sign $h$.  Since
$q\gamma=u$, pulling the objective through \eqref{eq:sdp-root-lift} gives
\begin{equation}
 \overline C_h=3I+uA_h,\qquad
 A_h=H_{12}+H_{23}+hH_{13}.                                \label{eq:sdp-effective-cost}
\end{equation}
For $h=+1$, $A_h$ has eigenvalues $2,-1,-1$; for $h=-1$, its signed cycle
product is negative and its eigenvalues are $1,1,-2$.  Hence
\begin{equation}
 \operatorname{spec}(\overline C_+)
 =\left\{\frac{16}5,\frac{29}{10},\frac{29}{10}\right\},
 \quad
 \operatorname{spec}(\overline C_-)
 =\left\{\frac{31}{10},\frac{31}{10},\frac{14}5\right\}.  \label{eq:sdp-intrinsic-spectra}
\end{equation}
Rayleigh--Ritz on $Z\succeq0$, $\tr Z=1$, proves
\eqref{eq:sdp-intrinsic-values}.  The common trace nine rules out an
orthogonal congruence followed by nontrivial positive rescaling.  Even after
adding a scalar multiple of $I$, the simple-largest/repeated-smallest and
repeated-largest/simple-smallest multiplicities cannot be exchanged.  The
gap is therefore intrinsic to the optimization problem.

The additive intervals are disjoint exactly when $\epsilon<1/20$ and meet
at $57/20$.  The relative intervals meet exactly at
\begin{equation}
 \eta_*=\frac{v_+-v_-}{v_++v_-}=\frac1{57},
 \qquad v_+(1-\eta_*)=v_-(1+\eta_*)=\frac{812}{285}.
\end{equation}
Below either threshold, comparison with a public separating point computes
parity with bounded error.  One coefficient query uses at most one sign query,
and the approximate-degree bound recalled in \cref{sec:prelim-lower-bound-tools}
gives the bounded-error lower bound $\lceil n_{\rm bit}/2\rceil$.  Conversely, apply
Deutsch's one-query circuit to each pair and query an unpaired last sign;
this computes $h$ in exactly that many queries.  At and above the thresholds
the displayed intersection points are public zero-query answers.

Projection onto $X_0$ identifies the feasible set with
$\mathcal D_3=\{Z\succeq0:\tr Z=1\}$.  If it had a finite SOC lift, replace
each affine lift equation $Ay=b$ by $Ay=tb$ and impose $t\geq0$.  For $t>0$
the projection is $t\mathcal D_3$.  At $t=0$, a homogeneous lift direction
could be added at arbitrary scale to a fixed feasible lift; compactness of
$\mathcal D_3$ forces its matrix projection to be zero.  The homogenized
projection is exactly $\Splus^3$, contradicting Fawzi's theorem
\cite{fawzi2018-on-representing-the-positive-semidefinite}.
\end{proof}

With $r_+=(1,1,1)^T/\sqrt3$ and $r_-=(1,-1,1)^T/\sqrt3$, strictly
complementary reduced pairs are
\begin{equation}
 Z_+^*=\tfrac12(I-r_+r_+^T),\quad Z_-^*=r_-r_-^T,
 \quad \overline S_+^*=\tfrac3{10}r_+r_+^T,
 \quad \overline S_-^*=\tfrac3{10}(I-r_-r_-^T).            \label{eq:sdp-strict-complementarity}
\end{equation}
Lifting by $X_i^*=R_iZ_h^*R_i^T$ and
$S_i^*=P^{-1}R_i\overline S_h^*R_i^T$ gives a strictly complementary
product-cone optimum.  Stationarity follows because the remaining
coefficient tuple is orthogonal to every trace-zero feasible tangent and
hence lies in the equality adjoint's range.

\subsubsection{A public, condition-one Newton step}
\label{sec:sdp-intrinsic-newton}

Use
\begin{equation}
 \Phi_\mu(X)=\frac1P\sum_i\inner{C_i}{X_i}
              -\mu\sum_i\log\det X_i,\qquad
 X_i^0=I/3,\quad\mu=1/P.                                  \label{eq:sdp-public-start}
\end{equation}
At this point, in \eqref{eq:sdp-svec},
\begin{equation}
 g_i=(C_i-3I)/P,\qquad \nabla^2\Phi_\mu(X^0)=(9/P)I.       \label{eq:sdp-public-hessian}
\end{equation}
Thus both right-hand sides of
\begin{equation}
 \widetilde A_\varsigma\Delta x=0,\qquad
 (9/P)\Delta x-\widetilde A_\varsigma^T\Delta y=-g       \label{eq:sdp-kkt}
\end{equation}
are public.  For trace-zero $Y$, define the isometry
$\mathcal W_\varsigma(Y)=P^{-1/2}(R_0YR_0^T,\ldots,R_{P-1}YR_{P-1}^T)$.
Then
\begin{equation}
 \mathcal W_\varsigma^*\nabla^2\Phi_\mu(X^0)
 \mathcal W_\varsigma=(9/P)I_5.                            \label{eq:sdp-reduced-nine}
\end{equation}
In the physical root coordinate the Hessian is $9I_5$ and
\begin{equation}
 \Delta Z_h=-\frac u9A_h.                                  \label{eq:sdp-root-direction}
\end{equation}
Its decrement and the standard full-step bound are
\begin{equation}
 \lambda=\frac1{\sqrt{150}},\qquad
 \lambda_{\rm next}\leq\frac1{(\sqrt{150}-1)^2}<\frac1{126}.
 \label{eq:sdp-decrement}
\end{equation}
Indeed $\lambda^2=\inner{uA_h}{(9I)^{-1}uA_h}=1/150$.
The post-step root eigenvalues are $(14/45,31/90,31/90)$ for $h=+1$ and
$(29/90,29/90,16/45)$ for $h=-1$, so the undamped step remains positive in
every block.  The next-decrement estimate is the universal self-concordant
bound, not a claim about a paper-specific primal--dual neighborhood.

\begin{theorem}[Newton and complete-KKT state hardness]
\label{thm:sdp-newton-state}
At \eqref{eq:sdp-public-start}, preparing the normalized root or global
primal Newton-direction state to trace distance $1/100$ requires
$\Omega(P)$ raw queries.  This remains true for every exact feasible lift of
a reduced solve $Y$ satisfying
\begin{equation}
 \frac{\norm{9Y+uA_h}_F}{\norm{uA_h}_F}\leq\frac1{10}.       \label{eq:sdp-reduced-residual}
\end{equation}
It also holds for the normalized complete primal-plus-multiplier solution of
\eqref{eq:sdp-kkt} in the displayed unit row scaling and multiplier
convention.  The latter has a public parity observable of magnitude at least
$1489/2097>7/10$.
\end{theorem}

\begin{proof}
Up to global sign,
\begin{equation}
 \ket{\delta z_h}=\frac{\ket{12}+\ket{23}+h\ket{13}}{\sqrt3}.
\end{equation}
The public contraction $T=\ket{12}\!\bra{13}+\ket{13}\!\bra{12}$ has
expectation $2h/3$.  At block $i$, both swapped coordinates acquire the same
prefix sign, so the same expectation holds for
$T_{\rm glob}=\bigoplus_iT$.  The POVM $(I\pm T)/2$ succeeds with probability
$5/6$, and trace error $1/100$ leaves constant bias.

For the residual claim, \eqref{eq:sdp-reduced-residual} is exactly
$\norm{Y-\Delta Z_h}_F\leq\eta\norm{\Delta Z_h}_F$ with
$\eta\leq1/10$.  The projective sine angle from the exact ray is at most
$\eta$.  Direct minimization gives the sharp bound
\begin{equation}
 h\langle T\rangle\geq m(\eta):=\frac23-\frac{\eta^2}{3}
 -\frac{2\sqrt2}{3}\eta\sqrt{1-\eta^2},\qquad0\leq\eta\leq\frac13.
 \label{eq:sdp-sharp-decoder}
\end{equation}
To verify it, at sine angle $s$ write the only useful tangent coefficient as
$x\in[-1,1]$.  The expectation is
$\frac23(1-s^2)-s^2+\frac{2\sqrt2}{3}s\sqrt{1-s^2}x+
\frac43s^2x^2$.  For $s\leq1/3$ its unconstrained minimizer is below $-1$,
so $x=-1$, and the result decreases with $s$.  At $\eta=1/10$ it is
$199/300-\sqrt{198}/150>17/30$.  Also
$\lambda_{\min}(I/3+Y)>3/10$.

For the complete vector, put
\begin{equation}
 \alpha=16/33,\quad c=\sqrt2\gamma/P,\quad d=\sqrt2u/9,
 \quad v_e=\begin{cases}c(1-\alpha)e,&e\leq K,\\
                          c\alpha(P-e),&K<e<P.
             \end{cases}                                  \label{eq:sdp-tent}
\end{equation}
If $\tau_i$ is the prefix sign, the three primal chains are
$-d\tau_i,-dh\tau_i,-d$ in coordinates $12,13,23$.  Switching signed path
incidence to ordinary incidence shows that the only nonzero copy multipliers
are $-\tau_ev_e$ on $12$ and $h\tau_ev_{P-e}$ on $13$.  The public
within-sector swap then has expectation $2h(D-C)/(3D+2V)$, where
\begin{align}
 D&=\frac{11n_{\rm bit}}{1350},\nonumber\\
 V&=\frac{289n_{\rm bit}}{4950}+\frac{17}{158400n_{\rm bit}},\nonumber\\
 C&=\frac{577n_{\rm bit}}{9900}+\frac1{9900n_{\rm bit}}.  \label{eq:sdp-kkt-sums}
\end{align}
Its sign is $-h$.  The leading ratio is $1489/2097$; the ratio of positive
$1/n_{\rm bit}$ corrections is $16/17>1489/2097$, so finite size only
increases the magnitude.  Reversing outcome labels decodes parity.  Each
preparer would therefore compute parity and costs
$\Omega(n_{\rm bit})=\Omega(P)$ queries.
\end{proof}

The complete-state claim is coordinate-specific: rescaling equality rows
changes multiplier amplitudes.  Also $D A_-D=-A_+$ for
$D=\diag(1,-1,1)$, so direction rays are related by an input-dependent
orthogonal gauge.  The value gap is invariant; the state theorem requests
fixed public coordinates, consistent with \cref{sec:condensation}.

\subsubsection{Ambient robustness and access objects}
\label{sec:sdp-intrinsic-access}

For arbitrary $X_i\succeq0$, define
\begin{equation}
 e_0=\tr X_0-1,\quad E_i=X_i-G_iX_{i-1}G_i^T,\quad
 e_0^2+\sum_{i=1}^{P-1}\norm{E_i}_F^2\leq\epsilon^2.       \label{eq:sdp-ambient-residual}
\end{equation}
Gauging to $Z_i=R_i^TX_iR_i$ and telescoping gives
\begin{equation}
 |P^{-1}\sum_i\tr X_i-1|<\sqrt P\epsilon,\quad
 \left|P^{-1}\sum_i\inner{C_i}{X_i}-\inner{\overline C_h}{Z_0}\right|
 <\frac{13}{4}\sqrt P\epsilon.                            \label{eq:sdp-ambient-telescope}
\end{equation}
The second bound writes the difference as suffix costs against
$Z_i-Z_{i-1}$ and uses $\norm{C_i}_{\rm op}<13/4$.  If the objective is at
most $v_h+1/100$ and
$\epsilon\leq1/(400\sqrt P)$, define
\begin{equation}
 \rho(X):=\frac{\bigoplus_iX_i}{\sum_i\tr X_i}.             \label{eq:sdp-ambient-density}
\end{equation}
The trace bounds in \eqref{eq:sdp-ambient-telescope} make this density
operator well-defined.  The public contraction
\begin{equation}
 M=\frac1{2/5}\left(\bigoplus_iC_i-\frac{57}{20}I\right)
\end{equation}
satisfies $h\tr(M\rho(X))>1/16$.  Specifically the normalized cost is
$>23/8$ for $h=+1$ and $<113/40$ for $h=-1$.  Thus any procedure which,
for every input, returns a density operator $\widetilde\rho$ satisfying
$\Dtr(\widetilde\rho,\rho(X))\leq1/100$ for some PSD tuple $X$ meeting the
displayed residual and objective conditions requires $\Omega(P)$ raw
queries.  The tuple may be selected adversarially by the procedure.  The
residual order is necessary: averaging the
$n_{\rm bit}$ neighboring negative-parity optimizers against a fixed
positive-parity input gives wrong value and residual
\begin{equation}
 \frac4{3\sqrt{n_{\rm bit}}}=\Theta(P^{-1/2}).              \label{eq:sdp-wrong-parity-witness}
\end{equation}

For access objects, use the following state-pair construction.  If a real
matrix $B$ has public support and magnitudes and maximum absolute row and
column sums at most $a$, give each supported pair $(r,c)$ its own token.
Row and column states with amplitudes
$\sqrt{|B_{rc}|/a}$ and
$\operatorname{sgn}(B_{rc})\sqrt{|B_{rc}|/a}$, completed on disjoint failure
tokens, are orthonormal and have overlap $B_{rc}/a$.  Unitary completions
therefore give an exact normalization-$a$ block encoding.

Put $y'=-\Delta y$ in \eqref{eq:sdp-kkt} and apply this to
\begin{equation}
 M_\varsigma=\begin{pmatrix}(9/P)I&\widetilde A_\varsigma^T\\
                 \widetilde A_\varsigma&0\end{pmatrix}.    \label{eq:sdp-canonical-kkt}
\end{equation}
Its maximum absolute row and column sums equal three; the trace row attains
three and a primal row has sum at most $2+9/P<3$.  One query phases the four
symmetric hidden KKT tokens.  This gives a fully specified exact
normalization-three encoding, and substitution into the decoders proves
\begin{equation}
 q_{\rm BE}\geq\left\lceil\frac{n_{\rm bit}}2\right\rceil  \label{eq:sdp-kkt-call-lower}
\end{equation}
for root, global, or complete-KKT state preparation.  The claim is relative
to this completion; input-dependent junk may contain parity and must be
charged.

Let $\Pi_\varsigma$ project onto $\ker\widetilde A_\varsigma$.  Projecting
\eqref{eq:sdp-kkt} and calculating norms gives
\begin{equation}
 \norm{\Pi_\varsigma g}_2^2=\frac3{50P},\quad
 \norm g_2^2=\frac{41}{400P},\quad
 \norm{\Pi_\varsigma\ket{\widehat g}}_2^2=\frac{24}{41}. \label{eq:sdp-projector-signal}
\end{equation}
One use of an operator-error-$1/32$ normalization-one projector encoding
then leaves signed expectation at least
$16/41-\sqrt{24/41}/16-1/1024>0.341$.  Hence first synthesis costs at least
$\lceil n_{\rm bit}/2\rceil$ raw queries.  Conversely, applying all prefix
signs to the $12,13$ chains constructs the tangent isometry; conjugating the
public trace-zero root projector gives an exact encoding with
$2n_{\rm bit}$ queries and no QRAM.  Projector synthesis is therefore
$\Theta(n_{\rm bit})$; a supplied projector already performs the aggregation.

Condition one in \eqref{eq:sdp-reduced-nine} does not describe the unscaled
saddle matrix.  Switching removes all path signs without changing singular
values.  On each individual off-diagonal coordinate, the path-difference
operator has singular values $2\sin(j\pi/(2P))$, $1\leq j<P$; in particular,
this gives $s_{\min}^+(\widetilde A_\varsigma)\leq\pi/P$.

The trace-coupled diagonal subsystem supplies the missing lower bound.  Let
the three-coordinate vector $x=(x^{(1)},x^{(2)},x^{(3)})$ be orthogonal to
its kernel, and write $m_j=P^{-1}{\bf1}^Tx^{(j)}$.
Orthogonality to every constant triple whose
coefficients sum to zero implies $m_1=m_2=m_3=:m$.  Put
$x^{(j)}=m{\bf1}+z^{(j)}$, with each $z^{(j)}$ mean zero, and define
\begin{equation}
 E^2=\sum_{j=1}^3\norm{Bz^{(j)}}_2^2,\qquad
 t=\sum_{j=1}^3x^{(j)}_0,
\end{equation}
where $B$ is ordinary path incidence.  The path Poincar\'e inequality and
telescoping give
\begin{equation}
 \sum_j\norm{z^{(j)}}_2^2\leq\frac{P^2}{4}E^2,\qquad
 |z^{(j)}_0|\leq\sqrt P\norm{Bz^{(j)}}_2.
\end{equation}
Since $3m=t-\sum_jz^{(j)}_0$, it follows that
$\norm x_2^2\leq(9/4)P^2(t^2+E^2)$.  The last parenthesis is exactly the
squared norm of the diagonal subsystem applied to $x$.  Combining diagonal
and off-diagonal coordinates, and using the maximum absolute row and column
sums three and two, yields
\begin{equation}
 \frac{2}{3P}\leq s_{\min}^+(\widetilde A_\varsigma)
 \leq\frac\pi P,\qquad
 s_{\max}(\widetilde A_\varsigma)\leq\sqrt6.               \label{eq:sdp-equality-singular-bounds}
\end{equation}
The kernel gives the saddle eigenvalue $9/P$, while every positive equality
singular value $s$ gives eigenvalues
$\{9/P\pm\sqrt{(9/P)^2+4s^2}\}/2$.  The two-sided bounds in
\eqref{eq:sdp-equality-singular-bounds} therefore prove
\begin{equation}
 \kappa_2(M_\varsigma)=\Theta(P).                          \label{eq:sdp-saddle-kappa}
\end{equation}
Its scalar off-diagonal graph is nevertheless a forest: each coordinate
forms a subdivided path and the trace multiplier joins three distinct
diagonal paths at one new vertex.  Thus its treewidth is one.

\subsection{Non-SOCP Schur-block families}
\label{sec:sdp-schur}

Fix $k\geq2$, let $K=16n_{\rm bit}$ and $L=17n_{\rm bit}+1$, and set
$a_i=\varsigma_i$ for $i\leq n_{\rm bit}$ and $a_i=1$ thereafter, with
$\tau_i=\prod_{j=1}^ia_j$.  For $X_i\in\Splus^{k+1}$, put
$F_j=(e_je_{k+1}^T+e_{k+1}e_j^T)/\sqrt2$ and
$E=e_{k+1}e_{k+1}^T$.  Impose
\begin{align}
 \inner{F_1}{X_0}&=\sqrt2,&
 \inner{F_1}{X_i}-a_i\inner{F_1}{X_{i-1}}&=0,\nonumber\\
 \inner{F_j}{X_0}&=0,&
 \inner{F_j}{X_i}-\inner{F_j}{X_{i-1}}&=0\quad(2\leq j\leq k),\nonumber\\
 \inner E{X_0}&=1,&
 \inner E{X_i}-\inner E{X_{i-1}}&=0                       \label{eq:sdp-schur-chains}
\end{align}
for $1\leq i<L$.  Rows and scalar columns have incidence at most two.  Every
feasible block is
\begin{equation}
 X_i=\begin{pmatrix}A_i&\tau_i e_1\\\tau_i e_1^T&1\end{pmatrix},
 \qquad Z_i=A_i-e_1e_1^T\succeq0.                          \label{eq:sdp-schur-map}
\end{equation}
Put $\mu_0=1/L$ and use $C_0=\mu_0\diag(I_k,2)$ on every block.

\begin{theorem}[Exact Schur path and SOCP boundary]
\label{thm:sdp-schur}
The unique, strictly complementary optimum has $Z_i=0$.  For
$t=\mu/\mu_0$, the complete central path is
\begin{equation}
 X_{\tau_i}^{(k)}(t)=
 \begin{pmatrix}tI_k+e_1e_1^T&\tau_i e_1\\\tau_i e_1^T&1\end{pmatrix},
 \quad S_i(\mu)=
 \begin{pmatrix}\mu_0I_k&-\mu_0\tau_i e_1\\
 -\mu_0\tau_i e_1^T&\mu+\mu_0\end{pmatrix}.              \label{eq:sdp-schur-path}
\end{equation}
In every public Frobenius-orthonormal tangent basis,
$H_{\rm red}(\mu)=(\mu_0/\mu)^2I=t^{-2}I$.  All primal, slack, and multiplier
coordinates are bounded for $0<\mu\leq\mu_0$.  The feasible block slice has
a finite SOCP lift exactly when $k\leq2$.
\end{theorem}

\begin{proof}
Schur complementation gives \eqref{eq:sdp-schur-map}, $\det X_i=\det Z_i$,
and $\inner{C_0}{X_i}=\mu_0(\tr Z_i+3)$.  Thus $Z_i=0$ is the unique
optimum, and $\mu_0\tr Z_i-\mu\log\det Z_i$ is uniquely minimized at
$tI_k$, proving the primal formula.  Direct inversion gives the slack.
With $S=C-\mathcal A^*y$, the nonzero chain multipliers are
\begin{equation}
 \alpha_i=\sqrt2\mu_0(L-i)\tau_i,\qquad
 \beta_i=(\mu_0-\mu)(L-i).                                \label{eq:sdp-schur-multipliers}
\end{equation}
Their divergences reproduce $C-S$.  On $0<t\leq1$, primal entries are at
most two, slack entries at most $2/L$, and
$\norm\alpha_\infty\leq\sqrt2$, $\norm\beta_\infty\leq1$.  At $t=0$ the
rank-one primal and rank-$k$ displayed slack are strictly complementary.
For dual uniqueness, put $w_i=(\tau_i e_1,1)^T$, so
$X_i^{\rm opt}=w_iw_i^T$.  Zero gap for any dual optimum gives
$S_iw_i=0$.  Stationarity in the unconstrained upper-left coordinates fixes
$(S_i)_{1:k,1:k}=\mu_0I_k$; annihilating $w_i$ then forces the last column
and row to be $(-\mu_0\tau_i e_1,\mu_0)^T$, exactly the $t=0$ limit in
\eqref{eq:sdp-schur-path}.  Finally, full row rank of $\mathcal A$ makes
$\mathcal A^*$ injective, so the equality multipliers are unique as well.

Every feasible tangent is
$\Delta X_i=\left(\begin{smallmatrix}H_i&0\\0&0\end{smallmatrix}\right)$.
Using the inverse of \eqref{eq:sdp-schur-path}, whose upper-left block is
$t^{-1}I_k$ because $\tau_i^2=1$, the reduced log-det Hessian is
\[
 \sum_i\tr(X_i^{-1}\Delta X_iX_i^{-1}\Delta Y_i)
 =\frac{\mu_0^2}{\mu^2}\sum_i\inner{H_i}{K_i},
\]
proving scalarity.  The two local centers do not commute:
$[X_+^{(k)}(t),X_-^{(k)}(t)]=2t(e_{k+1}e_1^T-e_1e_{k+1}^T)$.
Finally, \eqref{eq:sdp-schur-map} is an affine bijection from $\Splus^k$.
$\Splus^2$ is a Lorentz cone, while $\Splus^3$ has no finite SOC lift; larger
$k$ contain it as a face.  This proves the exact boundary.
\end{proof}

The noncommutativity is removable by an input-dependent congruence; in the
$Z_i$ coordinates the optimization is public.  The product-cone
representation is essential: replacing it by one unrestricted
$(k+1)L$-dimensional PSD block and killing off-block entries takes
quadratically many scalar equalities.

\subsubsection{Tight loading and the accuracy frontier}
\label{sec:sdp-schur-frontier}

For $k=3$ and $0<t\leq1$, define
\begin{equation}
 \rho_\varsigma(t)=\frac{\bigoplus_iX_{\tau_i}^{(3)}(t)}{L(3t+2)}.
\end{equation}
On the suffix $T=\{n_{\rm bit}+1,\ldots,n_{\rm bit}+K\}$, every
$\tau_i=h$.  Let $O$ equal $e_1e_4^T+e_4e_1^T$ on these blocks and zero
elsewhere.  Then
\begin{equation}
 h\tr(O\rho_\varsigma(t))=\frac{2K}{L(3t+2)}
 \geq\frac{32n_{\rm bit}}{5(17n_{\rm bit}+1)}>\frac13.    \label{eq:sdp-schur-density-bias}
\end{equation}
Thus trace-distance-$1/100$ density preparation costs $\Omega(L)$ queries.

For the complete central triple
$g_\varsigma(t)=(\operatorname{svec}X,y,\operatorname{svec}S)$,
\begin{equation}
 \norm{g_\varsigma(t)}_2^2
 =L(3t^2+2t+4)+[2+(1-t)^2]\frac{(L+1)(2L+1)}{6L}
 +\frac{5+(1+t)^2}{L}<11L.                                \label{eq:sdp-schur-triple-norm}
\end{equation}
Swapping the suffix primal coordinates $1$ and $\sqrt2\tau_i$ gives a
public contraction $W$ with
$h\langle W\rangle=2\sqrt2K/\norm{g_\varsigma(t)}_2^2>1/5$.
Hence complete-triple amplitude preparation is also $\Omega(L)$ throughout
the late tail.

Here is the promised canonical normalization-two access object.  The equality
matrix has public unit-magnitude entries and maximum absolute row and column
sums two.  Give each row and column two public slots.  A degree-two row or
column prepares the equal superposition of its two entry tokens; a degree-one
row or column sends the unused slot to its own failure label, and a degree-zero
free column goes entirely to its failure label.  Distinct rows, columns, and
failure arms use disjoint labels.  Public two-slot Hadamards, reversible token
maps, and a fixed lexicographic bijection on invalid addresses give fixed
unitary completions $L_A$ and $R_A^{(0)}$, with every hidden sign set to
$+1$ in $R_A^{(0)}$ while retaining the public predecessor minus signs.  A
phase unitary $S_\varsigma$ recognizes the predecessor token in
signed transition $i$ and phases it by $\varsigma_i$, while acting as the
identity on all other entry, failure, and invalid labels.  Extending the sign
oracle by the public dummy value $\varsigma_0=1$ makes this one coherent raw
query.  Therefore
\begin{equation}
 U_A=L_A^\dagger S_\varsigma R_A^{(0)},\qquad
 \bra{0,r}U_A\ket{0,c}=\frac{(\mathcal A_\varsigma)_{rc}}2.              \label{eq:sdp-schur-canonical-be}
\end{equation}
After padding the row and column address spaces to a common public dimension,
this is an exact, fully specified normalization-two encoding, and every ordinary,
adjoint, or controlled call uses one raw query.  Thus
$q_{\rm BE}\geq n_{\rm bit}/2$ for either output.  An alternative
input-dependent completion, junk action, prefix table, or preprocessing stage
is admissible only when all queries used to construct it are charged.

For arbitrary $X_i\succeq0$, let $d_i=\sqrt2(X_i)_{14}$ and suppose
\begin{equation}
 \norm{\mathcal A_\varsigma(X)-b}_2\leq\epsilon,\qquad
 \sum_i\inner{C_0}{X_i}\leq6,\qquad
 \epsilon\leq\frac1{4\sqrt L}.                            \label{eq:sdp-schur-robust-contract}
\end{equation}
Define the trace-normalized density output of such a point by
\begin{equation}
 \rho(X):=\frac{\bigoplus_iX_i}{\sum_i\tr X_i}.             \label{eq:sdp-schur-robust-density}
\end{equation}
It is well-defined because the root residual forces a nonzero off-diagonal
entry and a PSD matrix with such an entry has positive trace.  If
$r_0=d_0-\sqrt2$ and $r_i=d_i-a_id_{i-1}$, switching the recurrence gives
$\tau_id_i=\sqrt2+\sum_{j\leq i}\tau_jr_j$.  Summing over $T$ and applying
Cauchy--Schwarz gives $h\sum_{i\in T}d_i>K$.  PSD and the objective cap give
$\sum_i\tr X_i\leq6L$.  Since $\tr(O_iX_i)=\sqrt2d_i$ on the suffix,
\begin{equation}
 h\tr(O\rho(X))
 =\frac{\sqrt2h\sum_{i\in T}d_i}{\sum_i\tr X_i}
 >\frac{\sqrt2K}{6L}>\frac15.                              \label{eq:sdp-schur-robust-bias}
\end{equation}
Therefore preparing any $\widetilde\rho$ with
$\Dtr(\widetilde\rho,\rho(X))\leq1/100$ for some blockwise PSD point meeting
\eqref{eq:sdp-schur-robust-contract} costs $\Omega(L)$ raw queries.  The
point may depend adversarially on the input; no centrality or matrix-norm
proximity is needed.  Conversely, averaging the exact centers
obtained by flipping each sign once gives the wrong suffix parity and
\begin{equation}
 \norm{\mathcal A_\varsigma(\overline X)-b}_2
 =2\sqrt{2/n_{\rm bit}}=\Theta(L^{-1/2}).                  \label{eq:sdp-schur-wrong-mixture}
\end{equation}

For precision parameters $\epsilon,\eta>0$, let
$Q_{\rm abs}(P,\epsilon)$ be the worst-case first-preparation raw-query cost,
including setup, of producing an unconditional density output
$\widetilde\rho$ within trace distance $1/100$ of $\rho(X)$ for some
blockwise PSD $X$ with at most $P$ blocks, absolute equality residual at most
$\epsilon$, and the displayed objective cap.  Define
$Q_{\rm rel}(P,\eta)$ identically with
$\norm{\mathcal A_\varsigma(X)-b}_2/\norm b_2\leq\eta$.
For a heralded procedure, its success probability and all repetition or
amplification costs are charged as in \cref{def:heralded-output}; reusable
preprocessing is charged as in \cref{def:amortized-accounting}.

\begin{theorem}[Accuracy--query frontier]
\label{thm:sdp-accuracy-frontier}
For at most $P$ blocks and the relational output contract above,
\begin{equation}
 Q_{\rm abs}(P,\epsilon)=\Omega(\min\{P,\epsilon^{-2}\}),
 \qquad Q_{\rm rel}(P,\eta)=\Omega(\min\{P,\eta^{-2}\})  \label{eq:sdp-accuracy-frontier}
\end{equation}
for $0<\epsilon,\eta\leq1/8$.  The relative residual is the global
RHS-relative residual of \cref{def:global-residual}.  Exact-$P$ public
padding preserves the theorem, and the bounds are tight for the tuned family.
\end{theorem}

\begin{proof}
For absolute error choose
$n_{\rm bit}=\lfloor\min\{(P-1)/17,1/(2048\epsilon^2)\}\rfloor$.
Apart from a constant-size range, this is
$\Theta(\min\{P,\epsilon^{-2}\})$ and ensures the last inequality in
\eqref{eq:sdp-schur-robust-contract}.  The decoder reduces parity to every
valid output.  Since $\norm b_2=\sqrt3$, replacing $2048$ by $6144$ proves
the relative claim.

For exact $P$, add $P-L$ public $4\times4$ dummy blocks fixed by ten
one-coordinate rows to $P^{-1}I_4$, and set the common cost scale and
parameter to $1/P$.  The cap becomes $6L/P+5(P-L)/P^2$, so total trace is at
most $6L+5$ and active bias remains above $1/5$.  Dummy blocks have no
tangent coordinates, have bounded multipliers and a strictly complementary
optimum, and do not change the active non-SOCP slice.  Their RHS makes
$\norm{b_{\rm pad}}_2^2<4$, so $8192$ replaces $6144$ for exact-$P$ relative
padding.  Reading all active signs gives an exact valid output in
$O(n_{\rm bit})$ queries, while \eqref{eq:sdp-schur-wrong-mixture} matches
the residual exponent.
\end{proof}

A hidden coordinate in block $i$ is the prefix $\tau_i$ and has complexity
$\Theta(\min\{i,n_{\rm bit}\})$.  A classical scan streams all blocks with
exactly $n_{\rm bit}$ queries and constant extra workspace.  Quantumly, use
\begin{equation}
 X_\tau^{(k)}(t)=t\sum_{j=1}^ke_je_j^T+w_\tau w_\tau^T,
 \qquad w_\tau=\tau e_1+e_{k+1},                           \label{eq:sdp-schur-gram}
\end{equation}
to prepare a public local purification.  A coherent prefix-phase loop queries
and unqueries each sign, using at most $2n_{\rm bit}$ queries, logarithmic
workspace, and no QRAM.  Public coordinate predicates extend it to the
complete-triple state.  Both first-preparation complexities are therefore
$\Theta(L)$.

If $q_k=(k+1)(k+2)/2$, a path bag containing two adjacent block cliques and
their $k+1$ transition multipliers proves
\begin{equation}
 \operatorname{tw}(K_{\rm KKT})\leq2q_k+k.                 \label{eq:sdp-schur-treewidth}
\end{equation}
Thus the free central solve uses zero input queries while original-coordinate
loading uses $\Theta(L)$.  This complements bounded-treewidth chordal
conversion \cite{zhang2024-complexity-of-chordal-conversion-for}.  The bound
is for first preparation: charged cached prefixes can make later centers
cheap.  It gives no lower bound for a reduced $Z$ output or the public value.

\subsection{Abstract tree-holonomy loading}
\label{sec:sdp-tree-loading}

\begin{theorem}[Tree-holonomy loading]
\label{thm:sdp-tree-loading}
Let $\mathcal T=(V,E)$ be a public rooted tree with edge signs
$\varsigma_e$ and transported labels
$\tau_v=\prod_{e\in\operatorname{path}(r,v)}\varsigma_e$.  For public
constant-dimensional $B_+,B_-\succeq0$ of common trace $T_0$, put
\begin{equation}
 \rho_\varsigma=\frac1{|V|T_0}\bigoplus_{v\in V}B_{\tau_v},
 \qquad d=\frac12\norm{B_+/T_0-B_-/T_0}_1.                 \label{eq:sdp-tree-state}
\end{equation}
If public constant-size circuits prepare local purifications, a purification
of \eqref{eq:sdp-tree-state} needs at most $4|E|$ sign queries and logarithmic
workspace, without QRAM.  If an unknown path of length $n_{\rm bit}$ is
followed by $K=\Theta(n_{\rm bit})$ public descendants, $|V|=\Theta(n_{\rm bit})$,
and $d>0$, preparation to trace distance below a constant depending only on
$d$ and $K/|V|$ needs at least $n_{\rm bit}/2$ queries on the first
preparation.
\end{theorem}

\begin{proof}
In a public depth-first order, an edge is on the root-to-$v$ path exactly when
$v$ lies in the child subtree's preorder interval.  Prepare a uniform
entangled vertex address, query every edge, toggle a label on that interval
when the sign is negative, and unquery.  Select $B_{\tau_v}/T_0$, then reverse
the loop to erase the label.  The two compute--uncompute loops use four
queries per edge and logarithmic workspace.

For the lower bound, project onto the $K$ descendants.  This has public
probability $\alpha=K/|V|$.  Conditioned on it, the local state is $B_h/T_0$,
whose Helstrom success is $(1+d)/2$.  Guessing randomly off the tail gives
$1/2+\alpha d/2$.  Preparation error below $\alpha d/4$ leaves constant
advantage, so parity costs $n_{\rm bit}/2$ queries.  A classical tree scan
matches the linear order and may restrict itself to the union of requested
root paths.
\end{proof}

\subsection{Matrix-valued group holonomy}
\label{sec:sdp-group-holonomy}

Let $U:G\to O(d)$, put $k=3d$, and define
$W=U_{g_{n_{\rm bit}}}\cdots U_{g_1}$.  Use the copied construction with
$P=33n_{\rm bit}$ and hidden transports
$\mathcal D_g=\diag(U_g,I_d,I_d)$.  Replace $3I$ in
\eqref{eq:sdp-local-costs} by $kI_k$ and replace scalar $H_{ab}$ by
$H_{ab}(I_d)$.  Elimination gives
\begin{equation}
 \overline C_W=kI_k+uA_W,\qquad
 A_W=\begin{pmatrix}0&I_d&W^T\\I_d&0&I_d\\W&I_d&0\end{pmatrix}.          \label{eq:sdp-group-cost}
\end{equation}
Conjugating $W$ by $Q$ conjugates $A_W$ by $\diag(Q,Q,Q)$.

\begin{theorem}[Holonomy spectrum]
\label{thm:sdp-holonomy-spectrum}
\begin{equation}
 \operatorname{spec}(A_W)=
 \bigcup_{e^{i\theta}\in\operatorname{spec}(W)}
 \left\{2\cos\frac{\theta+2\pi\ell}{3}:\ell=0,1,2\right\}.              \label{eq:sdp-holonomy-spectrum}
\end{equation}
Thus $v(W)=k+u\lambda_{\min}(A_W)$ is a conjugacy-class function.  For the
natural permutation representation of $S_3$, the identity, transposition,
and three-cycle values are
\begin{equation}
 \frac{89}{10},\qquad \frac{44}{5},\qquad
 9+\frac15\cos\frac{8\pi}{9}.                             \label{eq:sdp-s3-values}
\end{equation}
\end{theorem}

\begin{proof}
On a $W$-eigenvector with eigenvalue $e^{i\theta}$, the complexified operator
is $\left(\begin{smallmatrix}0&1&e^{-i\theta}\\1&0&1\\e^{i\theta}&1&0\end{smallmatrix}\right)$,
whose characteristic polynomial is $\lambda^3-3\lambda-2\cos\theta$.
Substituting $\lambda=2\cos\phi$ proves the spectrum.  The natural $S_3$
eigenvalue lists are $(1,1,1)$, $(1,1,-1)$, and
$(1,e^{2\pi i/3},e^{-2\pi i/3})$, with least connection eigenvalues
$-1,-2,2\cos(8\pi/9)$.  Identity and three-cycles agree in the abelianization
but have different values, so the unrestricted observable is genuinely
nonabelian.
\end{proof}

Signed-permutation representations keep each svec copy row at two
unit-magnitude nonzeros, though the hidden symbol may move the predecessor
position.  The natural $S_3$ example uses $9\times9$ blocks, public costs in
$(9-1/4,9+1/4)$, and is strictly primal--dual feasible.

\begin{theorem}[$S_3$ value and Newton hardness]
\label{thm:sdp-s3-hardness}
Restrict every symbol to $\{e,t\}$ for one fixed transposition.  Additive
$1/25$ or target-relative $1/200$ value estimation takes $\Omega(P)$ raw
queries.  At $X_i=I_9/9$, $\mu=1/P$, the ambient and reduced primal Hessians
have condition number one, the KKT right-hand side is public, and
\begin{equation}
 \Delta Z_W=-\frac{u}{81}A_W,\qquad
 \lambda=\frac{\sqrt2}{30},\qquad
 \lambda_{\min}(I_9/9+\Delta Z_W)\geq\frac{44}{405}.       \label{eq:sdp-s3-newton-constants}
\end{equation}
Preparing the normalized root or global primal direction, or the normalized
complete primal-plus-multiplier KKT state in the stated row scaling, to trace
distance $1/100$ requires $\Omega(P)$ raw queries.  The complete-state
observable has magnitude greater than $33/100$.
\end{theorem}

\begin{proof}
The restricted word is identity or $t$ according to parity.  The first two
values in \eqref{eq:sdp-s3-values} differ by $1/10$ and are separated at
$177/20$; their relative intervals are disjoint below $1/177$.  A sparse
position query is simulated by a constant number of bit queries, so parity
proves value hardness.

At the public start the gradient is $(C_i-9I)/P$ and the ambient Hessian is
$(81/P)I$.  Orthogonal transport makes the normalized tangent map isometric,
so the physical root Hessian is $81I$.  Since $\tr A_W=0$, the direction is
as displayed.  Also $\norm{A_W}_F^2=18$ and
$\norm{A_W}_{\rm op}\leq2$, so
$\lambda^2=2u^2/9=1/450$ and
$1/9-2u/81=44/405$.  The next decrement is
$2/(30-\sqrt2)^2<1/400$.

For identity and a transposition, the normalized $A_W$ overlap is $7/9$.
The root-direction states have trace distance $4\sqrt2/9$, so their fixed
Helstrom measurement succeeds with probability
$\tfrac12(1+4\sqrt2/9)>0.81$.  Let $J$ swap corresponding
Frobenius-isometric svec coordinates in the $(1,2)$ and $(1,3)$ connection
blocks and act as zero on every other coordinate.  It has expectation $2/3$
for identity and $2/9$ for a transposition; centering as
$T_G=(9J-4I)/13$ gives $2h/13$ on every physical block and hence on the
global direction.

For the complete state, a public basis makes $U_t=\diag(-1,1,1)$.  Two modes
are public and one is the signed qutrit mode.  With
$d_0=\sqrt2u/81$, its primal contribution is
$D=Pd_0^2=11n_{\rm bit}/109350$ and its multiplier tent has the $V,C$ in
\eqref{eq:sdp-kkt-sums}.  Swapping $12,13$ only in the signed mode gives
\begin{equation}
 |\langle T_{\rm sign}\rangle|
 =\frac{2239504n_{\rm bit}^2+3888}
        {6759216n_{\rm bit}^2+12393}>\frac{33}{100},        \label{eq:sdp-s3-complete-bias}
\end{equation}
by cross multiplication.  These fixed decoders retain positive bias at
trace error $1/100$, and parity proves all state lower bounds.
\end{proof}

For canonical KKT access, $45$ scalar coordinates lie in each block.  A
copy-row state uses amplitude $1/3$ on two tokens and a trace-row state uses
amplitude $1/3$ on nine.  The column preparation routes predecessor tokens
through the svec involution induced by $t$, which one symbol query implements
by toggling its public orbit bit.  This exactly encodes the equality dilation
with normalization nine.  Adding $(81/P)$ times the public primal-sector
projector $P_x$ by a two-term LCU gives
\begin{equation}
 \alpha_M=9+\frac{81}{P}<12,                               \label{eq:sdp-s3-alpha}
\end{equation}
and $q_{\rm BE}\geq\lceil n_{\rm bit}/2\rceil$ for the three outputs above.
The completion is fixed; other input-dependent junk must be charged.

The tangent projector again has signal probability $24/41$.  Combining it
with the global decoder gives $48/533$.  At operator error $1/64$, the signed
expectation remains
$48/533-2\sqrt{24/41}/64-1/4096>0.065$.  First synthesis is therefore
$\Omega(n_{\rm bit})$.  Accumulating each prefix svec involution and
conjugating the public trace-zero projector gives an exact $2n_{\rm bit}$
query construction, so the bound is tight.  The scalar KKT graph is a forest:
a layer permutation makes disjoint paths, and the trace multiplier joins
distinct diagonal paths without a cycle.

If setup uses $S$ queries and later iterations use $q_t$, the canonical XOR
position convention yields
\begin{equation}
 S+\sum_tq_t\geq n_{\rm bit}/4,                            \label{eq:sdp-group-ledger}
\end{equation}
improved to $n_{\rm bit}/2$ for the clean one-query convention.  Linear
setup may store the word or prefixes and make later iterations cheap.

The unrestricted objective detects nonabelian information, but every proved
$\Omega(P)$ query theorem above restricts inputs to the abelian subgroup
$\{e,t\}\cong\mathbb Z_2$.  This is parity hardness, not a bound for a
genuinely noncommuting ordered-word promise.

The scalar value is invariant under simultaneous orthogonal gauge by
construction.  The direction decoders instead use the fixed natural
permutation basis and root frame.  A common public gauge merely conjugates
both public measurements, whereas an input-dependent output gauge must be
learned or supplied and is outside the stated output contract.

\begin{openproblem}[Intrinsically nonabelian word hardness]
\label{open:sdp-nonabelian}
Prove a local-symbol query lower bound for distinguishing conjugacy classes
of a genuinely noncommuting word promise, and transfer it through
\eqref{eq:sdp-group-cost}, without restricting the alphabet to one cyclic
subgroup.
\end{openproblem}

\subsection{Historical boundary and synthesis}
\label{sec:sdp-synthesis}

Signed switching and cycle products are classical graph invariants
\cite{kadyan2023-switching-equivalence-hermitian}; group synchronization
organizes the mechanism as gauge transport and holonomy
\cite{gao2019-the-geometry-of-synchronization-problems}, and orthogonal
synchronization is routinely relaxed by SDPs
\cite{ling2022-solving-orthogonal-group-synchronization-via}.  General group
oracle classification and symmetric-oracle problems use different promises
\cite{zhandry2015-quantum-oracle-classification,
copeland2021-quantum-query-complexity-of-symmetric}.  Fawzi's theorem gives
the exact lift boundary, not a coefficient-query lower bound.

The early $2\times2$ sparse-block examples reduce to paired LPs in their
exact-scalar form; their shared-trace variant is genuinely PSD.  The first
noncommuting $3\times3$ Schur example is still SOCP-representable, which is
why $4\times4$ blocks are the honest headline.  The two-anchor precursor used
asymmetric anchors: its two effective costs, and hence its central matrices,
have different spectra and are not related by any orthogonal similarity.

Separately, the free-$\Splus^3$ parity-central-Hessian precursor gives one
small operator result.  For
\begin{equation}
 \overline C_h=I+a(H_{12}+hH_{13}),\quad a=8/33,
 \qquad\mathcal L_h(Y)=\overline C_hY\overline C_h,         \label{eq:sdp-hessian-precursor}
\end{equation}
the reduced Hessian has condition number below five, but synthesizing a
public-constant-normalization block encoding to unscaled error $1/8$ costs
at least $n_{\rm bit}/2$ raw queries.  Indeed
$\operatorname{svec}(\mathcal L_h(E_{33}))=(a^2,0,\sqrt2ha,0,0,1)^T$;
swapping its $13,33$ coordinates has signed expectation
$2\sqrt2a/(1+a^2)^2>0.61$, and error $1/8$ leaves constant bias.  This is a
first-synthesis statement: a supplied Hessian oracle reveals parity in one
call.  In this free-$\Splus^3$ family---unlike the two-anchor family---the
two Hessians are related by an input-dependent orthogonal gauge.

Linear sparse optimal-value lower bounds already follow from LP constructions,
including Theorem~8.4 of Apers--Gribling
\cite{apers2026-quantum-speedups-linear-programming-interior}, and general SDP
value bounds predate these families \cite{apeldoorn2020-quantum-sdp-solvers}.
Mori et al.'s sparse QLS bound also uses parity propagation
\cite{mori2026-sparsity-dependent-complexity-lower-bound}.  The contribution
here is the optimization-native conjunction, not parity, propagation, or the
linear exponent.

The qutrit component imported by \cref{sec:winner-take-all} is exactly
\eqref{eq:sdp-intrinsic-parameters} with the costs
\eqref{eq:sdp-local-costs}; under that import \(N=n_{\rm bit}\).  Its spectra are
\eqref{eq:sdp-intrinsic-spectra} and its value gap is $1/10$.  Nothing here
makes that section's public-start Newton state winner-revealing; trace
allocation is governed by \cref{sec:condensation}.  As in the LP families of
\cref{sec:lp-hardness}, reduced conditioning, KKT sparsity, data access, and
requested coordinates are separate resources.  On these parity promises the
only quantum improvement is Deutsch's factor of two; classical reading takes
$n_{\rm bit}$ signs, so there is no asymptotic quantum speedup.

\section{Preconditioning and Interface Obstructions}
\label{sec:preconditioning}

``Just precondition'' and ``just keep the state'' are incomplete complexity
arguments.  Preconditioning may exchange spectral condition for factor
normalization, right-hand-side preparation, or recovery sensitivity.  A
normalized Newton state may in turn be too weak to supply the coordinatewise
diagonal needed by the next Newton system.  This section prices those exchanges
in the access models of \cref{sec:prelim-quantum-access}.  The conclusions are
products and complete ledgers, never lower bounds on one spectral factor in
isolation.

The analytic condition/normalization statements, including the fixed-factor
minimax result, are model free.  Only their compilation/accounting
corollaries, the parity reductions, and the lazy-oracle results use the raw
sparse coefficient model of
\cref{def:raw-sparse-access}.  The state-interface results instead use the
distinct state-preparation access of \cref{def:state-preparation-access}; they
are not raw-coefficient query lower bounds.  The lazy-oracle results finally
explain why the caching ceiling in \cref{sec:lp-gain-plateau} prevents an
unconditional per-iteration multiplier on a fixed instance.

\subsection{Factor normalization and whitening}
\label{sec:preconditioning-whitening}

We first isolate the signed gain-propagation block used by
\cref{thm:lp-gain-plateau}.  Let $N\geq2$, $M=17N$, and
\[
 T=\left\lceil\tfrac12\log_2N\right\rceil,
 \qquad H=2^T,\qquad H_i=2^{\min\{i,T\}},\qquad g_i=H_i/H_{i-1},
\]
and take $a_i=\varsigma_i$ for $i\leq N$ and $a_i=1$ afterward.  Thus
$\varsigma_i$ here is the hidden sign denoted $\sigma_i$ in
\cref{thm:lp-gain-plateau}.  Here \(H\) is the plateau height; the normal
matrix \(H_\mu\) used later is unrelated.  On
$d=(d_0,\ldots,d_M)^T$, define the pinned lower-bidiagonal factor
\begin{equation}
 (B_\varsigma d)_0=d_0,
 \qquad (B_\varsigma d)_i=d_i-g_ia_id_{i-1}.
 \label{eq:prec-gain-factor}
\end{equation}
Put $\tau_0=1$, $\tau_i=\prod_{j=1}^ia_j$, and
$D_\tau=\diag(\tau_0,\ldots,\tau_M)$.  Let $B_+$ be
\eqref{eq:prec-gain-factor} with every sign positive and set
$S_+=B_+B_+^T$.  Then
\begin{equation}
 B_\varsigma=D_\tau B_+D_\tau,
 \qquad S_\varsigma:=B_\varsigma B_\varsigma^T
 =D_\tau S_+D_\tau.
 \label{eq:prec-gauge}
\end{equation}
Thus the singular values are independent of the signs, although implementing
the gauge requires their prefix products.  This switching identity and the
corresponding inverse formulas are standard signed-graph facts
\cite{alazemi2024-signed-graphs-and-inverses-of}.

\begin{lemma}[Plateau singular scale]
\label{lem:prec-small-singular}
If $L=M-T+1$ is the number of plateau nodes, then
\begin{equation}
 \sigma_{\min}(B_\varsigma)
 \leq2\sin\frac{\pi}{2(L+1)}
 \leq\frac{\pi}{L+1}<\frac{\pi}{16N},
 \qquad \norm{B_\varsigma^{-1}}_2=\Theta(N).
 \label{eq:prec-small-singular}
\end{equation}
\end{lemma}

\begin{proof}
Take a Dirichlet sine vector supported on nodes $T,\ldots,M$, with virtual
zeros at $T-1$ and $M+1$.  On this support $B_+$ is the first-difference
operator.  Adding the omitted final difference can only increase its energy,
and the full Dirichlet path quotient is
$4\sin^2(\pi/(2(L+1)))$.  Equation~\eqref{eq:prec-gauge} gives the signed
case, while $\sin x\leq x$ and $L+1>16N$ give the two remaining inequalities.

For $i\geq j$, direct forward substitution gives
\begin{equation}
 (B_+^{-1})_{ij}=H_i/H_j.                    \label{eq:prec-factor-inverse}
\end{equation}
The column $B_+^{-1}e_0=(H_i)_i$ has norm $\Theta(N)$, whereas summing the
squares in \eqref{eq:prec-factor-inverse} gives
$\norm{B_+^{-1}}_F=O(N)$.  Orthogonal equivalence in
\eqref{eq:prec-gauge} proves the final claim.
\end{proof}

\begin{theorem}[Sharp plateau whitening frontier]
\label{thm:prec-whitening-frontier}
Let $P_\varsigma$ be any square factor such that, for $K\geq1$,
\begin{equation}
 \frac1K I\preceq P_\varsigma S_\varsigma P_\varsigma^*
 \preceq I.                                      \label{eq:prec-whitened-spectrum}
\end{equation}
Then
\begin{equation}
 \norm{P_\varsigma}_2\sqrt K
 \geq\frac1{\sigma_{\min}(B_\varsigma)}
 >\frac{16N}{\pi}.                              \label{eq:prec-whitening-frontier}
\end{equation}
Consequently every exact $(\alpha_P,a_{\rm anc},0)$ block encoding of $P_\varsigma$
obeys
\begin{equation}
 \boxed{\alpha_P\sqrt K>\frac{16N}{\pi}}.       \label{eq:prec-alpha-frontier}
\end{equation}
For $1\leq a\leq\Theta(N)$, the regularized family
\begin{equation}
 P_\varsigma^{(a)}
 =D_\tau(S_++a^{-2}I)^{-1/2}D_\tau             \label{eq:prec-regularized-family}
\end{equation}
has $\norm{P_\varsigma^{(a)}}_2=\Theta(a)$ and, after a constant
input-independent normalization of the largest eigenvalue,
\begin{equation}
 K=\Theta((N/a)^2),\qquad
 \norm{P_\varsigma^{(a)}}_2\sqrt K=\Theta(N).  \label{eq:prec-tight-frontier}
\end{equation}
Thus the order of the frontier is attained throughout, not only at its
endpoints.
\end{theorem}

\begin{proof}
Equation~\eqref{eq:prec-whitened-spectrum} implies
$\sigma_{\min}(P_\varsigma B_\varsigma)\geq K^{-1/2}$.  On a right
singular vector of $B_\varsigma$ for its least singular value,
\[
 \sigma_{\min}(P_\varsigma B_\varsigma)
 \leq\norm{P_\varsigma}_2\sigma_{\min}(B_\varsigma).
\]
This and \cref{lem:prec-small-singular} prove
\eqref{eq:prec-whitening-frontier}; the normalization of an exact block
encoding is at least the encoded norm.

For tightness, the eigenvalues of the congruence in
\eqref{eq:prec-regularized-family} are
\[
 \frac{\lambda}{\lambda+a^{-2}},\qquad \lambda\in\sigma(S_+).
\]
The bidiagonal factor has bounded norm, a diagonal entry of one, and inverse
norm $\Theta(N)$ by \cref{lem:prec-small-singular}.  Hence
$\lambda_{\max}(S_+)=\Theta(1)$ and
$\lambda_{\min}(S_+)=\Theta(N^{-2})$.  It follows that the factor norm is
$\Theta(a)$ and that the ratio of the displayed largest and smallest
eigenvalues is $\Theta((N/a)^2)$ in the stated range.  This proves
\eqref{eq:prec-tight-frontier}.
\end{proof}

The theorem is a specialization, not a new general quantum-preconditioning
obstruction.  Theorem~3 and Corollary~4 of Nie, Lai, and An already prove the
separate-block-encoding limitation for arbitrary preconditioners
\cite{hantao2026-pauli-structured-preconditioning-quantum-linear}.  What is
specific here is the constant $16/\pi$, the tight interpolation
\eqref{eq:prec-regularized-family}--\eqref{eq:prec-tight-frontier}, and the
preprocessing-aware raw-query theorem below.  Directly encoding a classically
formed or algebraically regrouped product is a different access model, as both
that work and \cref{prop:no-composed-preconditioning} emphasize.

Two useful approximate operator consequences are immediate.  If the factor is
supplied through an exact $(\alpha_P,a_{\rm anc},0)$ block encoding, then, for fixed
$\delta<1$,
\begin{align}
 \norm{P_\varsigma S_\varsigma P_\varsigma^*-I}_2\leq\delta
 &\Longrightarrow
 \alpha_P\geq
 \frac{\sqrt{1-\delta}}{\sigma_{\min}(B_\varsigma)},\label{eq:prec-approx-white}\\
 \norm{P_\varsigma B_\varsigma-I}_2\leq\delta
 &\Longrightarrow
 \alpha_P\geq
 \frac{1-\delta}{\sigma_{\min}(B_\varsigma)}.\label{eq:prec-approx-inverse}
\end{align}
Indeed, the first inequality gives
$\sigma_{\min}(P_\varsigma B_\varsigma)\geq\sqrt{1-\delta}$, while the
second gives $\sigma_{\min}(P_\varsigma B_\varsigma)\geq1-\delta$; the same
singular-vector argument then applies.
Constant-error approximate whitening or inversion supplied through such an
exact factor encoding therefore has $\alpha_P=\Omega(N)$.  More generally, if
$\beta=\norm{P_\varsigma S_\varsigma P_\varsigma^*}_2$ and the congruence
has condition at most $K$, the scale-invariant statement is
\begin{equation}
 \norm{P_\varsigma}_2\sqrt{K/\beta}
 \geq\sigma_{\min}(B_\varsigma)^{-1}.          \label{eq:prec-scale-invariant}
\end{equation}
This normalization is essential: scalar rescaling defeats an unnormalized
statement.
For a left factor normalized by
$\sigma(P_\varsigma B_\varsigma)\subseteq[1/K,1]$ and supplied through an
exact $(\alpha_P,a_{\rm anc},0)$ block encoding, the same argument gives
\begin{equation}
 \alpha_PK\geq\norm{P_\varsigma}_2K>\frac{16N}{\pi}.       \label{eq:prec-left-frontier}
\end{equation}

The factor norm still does not price compilation from raw signs.  Define
\begin{equation}
 F_\varsigma=\diag(B_\varsigma^{-1},B_+^{-1}),\quad
 G_\varsigma=\diag(B_\varsigma,B_+),\quad
 \ket r=\frac{\ket{0,e_0}+\ket{1,e_0}}{\sqrt2}.            \label{eq:prec-FG}
\end{equation}
By \eqref{eq:prec-factor-inverse},
\begin{equation}
 F_\varsigma\ket r
 =\frac1{\sqrt2}\bigl((H_i\tau_i)_i\oplus(H_i)_i\bigr),
 \qquad \norm{F_\varsigma r}_2=\left(\sum_iH_i^2\right)^{1/2}=\Theta(N).
 \label{eq:prec-inverse-state}
\end{equation}
Pair the signed and unsigned coordinates and measure their side qubit in the
$X$ basis on the $16N$ plateau-copy nodes.  Its parity-signed expectation is
\begin{equation}
 \frac{16NH^2}{\sum_iH_i^2}
 >\frac{16N}{17N+4/3}>\frac9{10}.             \label{eq:prec-inverse-bias}
\end{equation}

\begin{theorem}[Reusable inverse access--normalization tradeoff]
\label{thm:prec-inverse-access}
Suppose an offline stage makes $s$ raw sign queries and produces a reusable
exact $(\alpha,a,0)$ block encoding of $\widetilde F_\varsigma$.  One call,
its inverse, or a controlled call makes at most $q$ further raw queries.
If, for fixed $\delta\leq1/20$,
\begin{equation}
 \norm{\widetilde F_\varsigma G_\varsigma-I}_2\leq\delta, \label{eq:prec-inverse-promise}
\end{equation}
then
\begin{equation}
 \boxed{s+O(q\alpha/N)=\Omega(N)},\qquad
 \boxed{Ns+q\alpha=\Omega(N^2)}.             \label{eq:prec-inverse-product}
\end{equation}
The constants hidden in $O$ and $\Omega$ are absolute.
\end{theorem}

\begin{proof}
Since $G_\varsigma F_\varsigma=I$,
\[
 \norm{(\widetilde F_\varsigma-F_\varsigma)r}_2
 \leq\delta\norm{F_\varsigma r}_2.
\]
Thus the normalized approximate-inverse state retains a constant parity bias
under the decoder in \eqref{eq:prec-inverse-bias}, and
$\norm{\widetilde F_\varsigma r}_2\geq(1-\delta)\Theta(N)$.
Fixed-point amplitude amplification prepares it with constant success using
$O(\alpha/N)$ block-encoding calls.  Including setup, the resulting parity
algorithm uses $s+O(q\alpha/N)$ raw queries.  Quantum parity needs
$\Omega(N)$ queries \cite{beals2001-quantum-lower-bounds-by-polynomials},
proving the first inequality; multiplication by $N$
gives the second.  Equation~\eqref{eq:prec-approx-inverse} also gives
$\alpha=\Omega(N)$, so the amplification count is never subconstant.
\end{proof}

Three no-go qualifications delimit these statements.  First, a
preconditioned product supplied directly can be the identity:
$P_\varsigma=B_\varsigma^{-1}$ gives $P_\varsigma B_\varsigma=I$, and exact
Schur whitening gives $P_\varsigma S_\varsigma P_\varsigma^*=I$.  No lower
bound on the product's normalization and condition alone is possible.
Second, applying $D_\tau$ is not free.  Any operator-norm-$<1/3$
implementation needs $\Omega(N)$ raw queries: apply it to
$(\ket0+\ket M)/\sqrt2$ and measure in the
$(\ket0\pm\ket M)/\sqrt2$ basis to recover endpoint parity.  Third, the
factor-independent conservation law needs the hard input/output contract.
Consider specifically the plateau solve
$G_\varsigma x=r$ from \eqref{eq:prec-FG}, and require an output state within
trace distance $1/20$ of
$F_\varsigma\ket r/\norm{F_\varsigma r}_2$.  For any such end-to-end
preconditioned implementation, compose its output with the fixed $X$-basis
decoder of \eqref{eq:prec-inverse-bias}.  The reduction gives
\begin{equation}
 q_{\rm set}+q_{\rm rhs}+q_{\rm rec}+m q_{\rm BE}=\Omega(N),
 \label{eq:prec-end-to-end-conservation}
\end{equation}
where $m$ is the number of solver calls, $q_{\rm BE}$ is the raw cost per
call, the solve phase makes no other raw queries, and the other terms charge
setup, transformed-RHS preparation, and recovery.  Indeed,
composing all phases with that decoder gives a bounded-error parity algorithm,
so their total raw-query cost is $\Omega(N)$.  Without this particular
RHS/output guarantee, no such law is claimed; a public solution may carry no
hidden information.  An orthogonal factor rotation can move parity among the
charged phases.  The statement is for the pinned plateau block, not for every
unrestricted full-KKT preconditioner.

\subsection{Congruence condition versus physical-state recovery}
\label{sec:preconditioning-congruence}

A second frontier arises even when the preconditioner is an elementary
active--inactive scaling.  Let the fixed Euclidean space split orthogonally as
$\R^r=U\oplus K$, with projectors $Q$ and $R$, and suppose
\begin{equation}
 H_\mu=\mu^{-1}C_\mu+\mu D_\mu,\qquad 0<\mu\leq\mu_0.       \label{eq:prec-congruence-H}
\end{equation}
Assume $C_\mu=QC_\mu Q$ and $C_\mu|_U\to C\succ0$;
$D_\mu\succeq0$ is uniformly bounded;
$D_{KK,\mu}\to D_{KK}\succ0$; and all other blocks converge.  For a fixed
nonzero $f\in U$, put
\begin{equation}
 p=C^{-1}f,\qquad g=D_{KK}^{-1}D_{KU}p,                    \label{eq:prec-pg}
\end{equation}
and assume the nontrivial coupling $g\ne0$.  These are fixed-instance
strict-complementarity asymptotics, not uniform claims in a growing
dimension.

For $0\leq\gamma\leq1$, define
\begin{equation}
 T_{\mu,\gamma}=\mu^{\gamma/2}Q+\mu^{-\gamma/2}R,\quad
 G_{\mu,\gamma}=T_{\mu,\gamma}H_\mu T_{\mu,\gamma},\quad
 \eta_\mu=H_\mu^{-1}f,\quad z_{\mu,\gamma}=T_{\mu,\gamma}^{-1}\eta_\mu.
 \label{eq:prec-power-congruence}
\end{equation}
Because $f\in U$, $T_{\mu,\gamma}f=\mu^{\gamma/2}f$; the normalized
transformed RHS is exactly the original RHS state.  The tradeoff below is
therefore a condition--recovery effect, not a loading effect.
The normalized contraction that recovers physical coordinates is
\begin{equation}
 \mathsf R_{\mu,\gamma}=T_{\mu,\gamma}/\norm{T_{\mu,\gamma}}_2
 =\mu^\gamma Q+R,\qquad
 \Phi_{\mu,\gamma}(w)=
 \frac{\mathsf R_{\mu,\gamma}w}{\norm{\mathsf R_{\mu,\gamma}w}_2}.
 \label{eq:prec-recovery-map}
\end{equation}
For a unit $w$, measure its local normalized-state sensitivity by
\begin{equation}
 \chi_{\mu,\gamma}(w)=
 \left\|D\Phi_{\mu,\gamma}(w)|_{w^\perp}\right\|_{2\to2}.
 \label{eq:prec-recovery-sensitivity}
\end{equation}

\begin{theorem}[Power congruence frontier]
\label{thm:prec-congruence-frontier}
Under the preceding assumptions, for each fixed $\gamma\in[0,1]$ as
$\mu\downarrow0$,
\begin{align}
 \kappa_2(G_{\mu,\gamma})
 &=\Theta(\mu^{-2(1-\gamma)}),                         \label{eq:prec-congruence-kappa}\\
 \eta_\mu&=\mu\binom{p}{-g}+o(\mu),                    \label{eq:prec-physical-expansion}\\
 z_{\mu,\gamma}&=
 \binom{\mu^{1-\gamma/2}(p+o(1))}
       {-\mu^{1+\gamma/2}(g+o(1))}.                    \label{eq:prec-scaled-expansion}
\end{align}
Writing $\widehat z=z/\norm z_2$, exact normalized recovery holds and direct
postselection has
\begin{equation}
 \Phi_{\mu,\gamma}(\widehat z)
 =\frac{\eta_\mu}{\norm{\eta_\mu}_2},\qquad
 P_{\rm rec}=\norm{\mathsf R_{\mu,\gamma}\widehat z}_2^2
 =\Theta(\mu^{2\gamma}),                                \label{eq:prec-recovery-probability}
\end{equation}
while
\begin{equation}
 \chi_{\mu,\gamma}(\widehat z)=\Theta(\mu^{-\gamma}).   \label{eq:prec-recovery-chi}
\end{equation}
Consequently
\begin{equation}
 \boxed{\sqrt{\kappa_2(G_{\mu,\gamma})}\,
 \chi_{\mu,\gamma}(\widehat z)=\Theta(1/\mu)},
 \qquad
 \sqrt{\kappa_2(G_{\mu,\gamma})}P_{\rm rec}^{-1/2}
 =\Theta(1/\mu).                                         \label{eq:prec-congruence-product}
\end{equation}
\end{theorem}

\begin{proof}
Set $a_\mu=\mu^{-(1-\gamma)}$ and
$b_\mu=\mu^{1-\gamma}=a_\mu^{-1}$.  In $U\oplus K$ blocks,
\begin{equation}
 G_{\mu,\gamma}=
 \begin{pmatrix}
 a_\mu C_\mu+\mu^{1+\gamma}D_{UU,\mu}&\mu D_{UK,\mu}\\
 \mu D_{KU,\mu}&b_\mu D_{KK,\mu}
 \end{pmatrix}.                                         \label{eq:prec-congruence-blocks}
\end{equation}
The upper-left block $A_\mu$ is $\Theta(a_\mu)I_U$, and its Schur
complement is
\[
 b_\mu D_{KK,\mu}-\mu^2D_{KU,\mu}A_\mu^{-1}D_{UK,\mu}
 =\Theta(b_\mu)I_K.
\]
The elimination factor has off-diagonal norm
$O(\mu b_\mu)=o(1)$.  A uniformly conditioned near-identity congruence
therefore reduces \eqref{eq:prec-congruence-blocks} to blocks of sizes
$\Theta(a_\mu)$ and $\Theta(b_\mu)$, proving
\eqref{eq:prec-congruence-kappa}.

The $K$ block of $H_\mu\eta_\mu=f$ gives exactly
\[
 \eta_{K,\mu}=-D_{KK,\mu}^{-1}D_{KU,\mu}\eta_{U,\mu}.
\]
Substitution in the $U$ block gives
\[
 [\mu^{-1}C_\mu+\mu E_\mu]\eta_{U,\mu}=f,
 \quad
 E_\mu=D_{UU,\mu}-D_{UK,\mu}D_{KK,\mu}^{-1}D_{KU,\mu},
\]
where $E_\mu$ is bounded.  Hence
$\eta_{U,\mu}=\mu(C_\mu+\mu^2E_\mu)^{-1}f=\mu(p+o(1))$,
which proves \eqref{eq:prec-physical-expansion}; applying $T^{-1}$ proves
\eqref{eq:prec-scaled-expansion}.

Because $\mathsf Rz=\mu^{\gamma/2}\eta$, normalized recovery is exact.
The two expansions give
\[
 \norm{\mathsf R\widehat z}_2
 =\mu^\gamma
 \frac{(\norm p_2^2+\norm g_2^2+o(1))^{1/2}}
      {(\norm p_2^2+\mu^{2\gamma}\norm g_2^2+o(1))^{1/2}}
 =\Theta(\mu^\gamma),
\]
proving \eqref{eq:prec-recovery-probability}.  Differentiation gives, with
$\widehat\eta=\Phi(w)$,
\begin{equation}
 \chi_{\mu,\gamma}(w)=
 \frac{\norm{(I-\widehat\eta\widehat\eta^T)
 \mathsf R_{\mu,\gamma}(I-ww^T)}_2}
      {\norm{\mathsf R_{\mu,\gamma}w}_2}.                \label{eq:prec-chi-formula}
\end{equation}
This is $O(\mu^{-\gamma})$ at $w=\widehat z$.  For $\gamma>0$, project
$g/\norm g$ onto $\widehat z^\perp$.  The resulting unit tangent is
$g/\norm g+o(1)$ in $K$, while the recovered state tends to
$(p,-g)/(\norm p^2+\norm g^2)^{1/2}$.  Since $p\ne0$, the numerator in
\eqref{eq:prec-chi-formula} is bounded below by
$\norm p/(\norm p^2+\norm g^2)^{1/2}>0$.  This proves the matching lower
bound.  For $\gamma=0$, $\mathsf R=I$ and $\chi=1$.  Multiplying the
estimates proves \eqref{eq:prec-congruence-product}.
\end{proof}

The same two-sparse LP witnesses every exponent.  Consider
\begin{equation}
 \min_{x\geq0}x_2+x_3,\qquad
 x_1+x_2=1,\quad x_2-x_3=0,                               \label{eq:prec-small-witness}
\end{equation}
so
$A=\left(\begin{smallmatrix}1&1&0\\0&1&-1\end{smallmatrix}\right)$,
$b=e_1$, and $c=(0,1,1)^T$.  Its unique optimum is $(1,0,0)$.  Along the
central path $x_2=x_3=t_\mu$, $x_1=1-t_\mu$, with
$t_\mu=\mu+O(\mu^2)$, and
\begin{equation}
 H_\mu=A(XS^{-1})A^T
 =\frac{(1-t_\mu)^2}{\mu}
 \begin{pmatrix}1&0\\0&0\end{pmatrix}
 +\frac{t_\mu^2}{\mu}
 \begin{pmatrix}1&1\\1&2\end{pmatrix}.                 \label{eq:prec-witness-normal}
\end{equation}
Indeed, dual feasibility and complementarity give equal inactive slacks
$q_\mu=1+\mu/[2(1-t_\mu)]$ and $t_\mu q_\mu=\mu$, which yields the
expansion.  For $U=\operatorname{span}\{e_1\}$ and
$K=\operatorname{span}\{e_2\}$, the limiting blocks in
\eqref{eq:prec-congruence-H} are $C=e_1e_1^T$ and
$D=\left(\begin{smallmatrix}1&1\\1&2\end{smallmatrix}\right)$.  Thus
$p=1$, $g=1/2$, and $H_\mu^{-1}b=\mu(1,-1/2)^T+o(\mu)$.
This fixed two-row, three-variable LP therefore realizes all parts of
\cref{thm:prec-congruence-frontier}; at $\gamma=1$ its scaled condition has
a finite positive limit and its recovery amplitude is asymptotic to
$\sqrt{5/4}\,\mu$.

The power family is also the complete order-optimal interpolation within all
scalar block-diagonal SPD congruences.  Let
\begin{equation}
 T_\mu=t_U(\mu)Q+t_K(\mu)R,\qquad
 r_\mu=t_U/t_K,\quad
 M_\mu=\max\{r_\mu/\mu,\mu/r_\mu\},\quad
 K_\mu=\max\{r_\mu,r_\mu^{-1}\}.                         \label{eq:prec-general-scalars}
\end{equation}
Put
\begin{equation}
 \widetilde G_\mu=T_\mu H_\mu T_\mu,\qquad
 \widetilde z_\mu=T_\mu^{-1}\eta_\mu,\qquad
 \widetilde{\mathsf R}_\mu=T_\mu/\norm{T_\mu}_2,\qquad
 \widetilde\Phi_\mu(w)=
 \frac{\widetilde{\mathsf R}_\mu w}
      {\norm{\widetilde{\mathsf R}_\mu w}_2}.            \label{eq:prec-general-recovery}
\end{equation}
Define $\widetilde P_{\rm rec}$ by applying the squared norm in
\eqref{eq:prec-recovery-probability} to
$\widetilde{\mathsf R}_\mu$ and $\widehat{\widetilde z}_\mu$, and define
$\widetilde\chi_\mu$ by \eqref{eq:prec-recovery-sensitivity} with
$\widetilde\Phi_\mu$.

\begin{theorem}[General scalar-block frontier]
\label{thm:prec-general-scalar-frontier}
Uniformly over all positive scalar functions $t_U,t_K$,
\begin{equation}
 \sqrt{\kappa_2(\widetilde G_\mu)}=\Theta(M_\mu),\quad
 \widetilde P_{\rm rec}^{-1/2}=\Theta(K_\mu),\quad
 \widetilde\chi_\mu=\Theta(K_\mu).                      \label{eq:prec-general-frontier}
\end{equation}
Consequently
\begin{equation}
 \sqrt{\kappa_2(\widetilde G_\mu)}\,\widetilde\chi_\mu
 =\Theta(M_\mu K_\mu)\geq c/\mu.                       \label{eq:prec-general-product}
\end{equation}
The same holds with $\widetilde P_{\rm rec}^{-1/2}$ in place of
$\widetilde\chi_\mu$.  Equality in order,
$M_\mu K_\mu=\Theta(1/\mu)$, holds exactly for the maximal class
\begin{equation}
 r_\mu=\Omega(\mu),\qquad r_\mu=O(1).                    \label{eq:prec-equality-class}
\end{equation}
More explicitly,
\begin{equation}
 M_\mu K_\mu=
 \begin{cases}
  \mu/r_\mu^2,&r_\mu<\mu,\\
  1/\mu,&\mu\leq r_\mu\leq1,\\
  r_\mu^2/\mu,&r_\mu>1.
 \end{cases}                                             \label{eq:prec-three-cases}
\end{equation}
\end{theorem}

\begin{proof}
The assumptions imply uniform spectral equivalence
\begin{equation}
 c_0(\mu^{-1}Q+\mu R)\preceq H_\mu
 \preceq C_0(\mu^{-1}Q+\mu R).                           \label{eq:prec-spectral-equivalence}
\end{equation}
For the lower bound, write a vector as $u+k$.  Fixed $c,d,M>0$ give
\[
 (u+k)^TH_\mu(u+k)
 \geq(c/\mu-2M^2\mu/d)\norm u^2+(d\mu/2)\norm k^2,
\]
after Young's inequality; for small $\mu$ this proves the lower half of
\eqref{eq:prec-spectral-equivalence}, and boundedness proves the upper half.
Congruence by $T_\mu$ shows that the two spectral scales are
$t_U^2/\mu$ and $\mu t_K^2$, proving the condition estimate in
\eqref{eq:prec-general-frontier}.

Equation~\eqref{eq:prec-physical-expansion} gives
$\norm{\eta_{U,\mu}}=\Theta(\mu)$ and
$\norm{\eta_{K,\mu}}=\Theta(\mu)$.  Therefore
\[
 \norm{\widetilde{\mathsf R}_\mu\widehat{\widetilde z}_\mu}
 =\frac{\norm{\eta_\mu}}
 {\max\{t_U,t_K\}
  \sqrt{\norm{\eta_{U,\mu}}^2/t_U^2+
        \norm{\eta_{K,\mu}}^2/t_K^2}}
 =\Theta(K_\mu^{-1}).
\]
This proves the recovery-probability claim.  In the plane spanned by the
normalized $U$ and $K$ components, differentiating the output angle of
$\diag(t_U,t_K)$ gives tangent gain
\[
 \frac{r_\mu^{-1}\norm{\eta_{U,\mu}}^2+
       r_\mu\norm{\eta_{K,\mu}}^2}
      {\norm{\eta_{U,\mu}}^2+\norm{\eta_{K,\mu}}^2}
 =\Theta(K_\mu).
\]
This is a lower bound on $\widetilde\chi_\mu$; differentiation as in
\eqref{eq:prec-chi-formula} gives the matching upper bound.  The elementary
three-case calculation \eqref{eq:prec-three-cases} proves the product and
the equality class.
\end{proof}

This is not a universal preconditioning lower bound.  It excludes arbitrary
left preconditioners, input-dependent changes mixing $U$ and $K$, direct
preparation of the physical state, variable-time recovery, and algorithms
that request only an observable.  It also charges neither construction of
$Q$ nor block-encoding normalization.  If $g=0$, the recovery conclusion can
fail.  In particular, the state-preserving left equilibration in
\cref{thm:stable-split-equilibration} lies outside this congruence-recovery
contract and is a genuine positive escape.  The theorem is therefore a sharp
structural lemma for a specified family, realized by
\eqref{eq:prec-small-witness}, rather than a claim against all
preconditioners.

\subsection{The parity preconditioner dichotomy}
\label{sec:preconditioning-parity}

The parity LP gives a complementary dichotomy.  Its reduced primal Hessian is
condition one, but its equality normal matrix is a signed grounded-tree
Laplacian.  Fixed preconditioners cannot serve all signings, whereas a perfect
data-dependent sparse factor is locally trivial.  The global information then
appears in factor application, transformed-RHS preparation, or recovery.

The minimax claim already holds on a path.  Let
\begin{equation}
 B=\begin{pmatrix}
 1\\-1&1\\&-1&1\\&&\ddots&\ddots
 \end{pmatrix},\qquad C=BB^T,\qquad
 R=\diag(1,-1,1,-1,\ldots)                               \label{eq:prec-path-pair}
\end{equation}
for $m\geq1$.  Both $C$ and $RCR$ are legal sign inputs.

\begin{theorem}[Minimax fixed-preconditioner bound]
\label{thm:prec-fixed-minimax}
For SPD matrices write $\kappa=\lambda_{\max}/\lambda_{\min}$.
For every input-independent SPD matrix $M$,
\begin{equation}
 \max\left\{
 \kappa(M^{-1/2}CM^{-1/2}),
 \kappa(M^{-1/2}RCRM^{-1/2})
 \right\}\geq\frac{4m^2-1}{3}\geq m^2.                  \label{eq:prec-minimax}
\end{equation}
The identity gives condition number at most $4m^2$ for either signing, so it
is minimax optimal in order.  The statement allows dense $M$ and holds more
generally for every fixed invertible factor $P$, with the two matrices
$P^TCP$ and $P^TRCRP$.  The bound $(4m^2-1)/3$ is a witness lower bound,
not a claim of the exact minimax value.  It strengthens the former
$m^2/4$ bound for $m\geq4$.
\end{theorem}

\begin{proof}
For SPD $A_0,A_1,M$, generalized Rayleigh quotients give
\begin{equation}
 \kappa(A_0^{-1/2}A_1A_0^{-1/2})
 \leq\kappa(M^{-1/2}A_0M^{-1/2})
       \kappa(M^{-1/2}A_1M^{-1/2}).                       \label{eq:prec-simultaneous}
\end{equation}
Indeed, bound the numerator and denominator of
$x^TA_1x/(x^TA_0x)$ by the extremal eigenvalues in $M$ coordinates; the
independent scales cancel in the quotient of extrema.

For $r=(m,m-1,\ldots,1)^T$,
\[
 x^TCx=\sum_{j=0}^{m-2}(x_j-x_{j+1})^2+x_{m-1}^2,
 \qquad r^TCr=m.
\]
Summing the alternating witness exactly gives
\[
 r^TRCRr=\sum_{k=1}^m(2k-1)^2=\frac{m(4m^2-1)}3.
\]
Put $t=(4m^2-1)/3$.  The generalized quotient is $t$ at $r$ and
$1/t$ at $Rr$, so
$\kappa(C^{-1/2}RCRC^{-1/2})\geq t^2$.
Apply \eqref{eq:prec-simultaneous} with $A_0=C$ and $A_1=RCR$; one of its
two factors must be at least $t$.  The same argument in coordinates
$P^{-1}r$ and $P^{-1}Rr$ proves the assertion for any invertible $P$,
without a polar-decomposition assumption.

For the upper bound, write $d_i=x_i-x_{i+1}$ for $i<m-1$ and
$d_{m-1}=x_{m-1}$.  Then $x_i=\sum_{j=i}^{m-1}d_j$ and
$x^TCx=\sum_jd_j^2$.  Cauchy--Schwarz gives
$\norm{x}_2^2\leq m^2x^TCx$, while
$(a-b)^2\leq2a^2+2b^2$ gives $x^TCx\leq4\norm{x}_2^2$.
Thus $\kappa(C)\leq4m^2$.  Since $R$ is orthogonal, the same bound
holds for $RCR$; more generally every signed path is an orthogonal
diagonal conjugate of $C$.
\end{proof}

The fixed-preconditioner result is formalized in Lean, including the
actual incidence matrices, their positive definiteness, extremal-eigenvalue
condition numbers, the simultaneous-preconditioning inequality, the exact
witness ratio, and the identity upper bound for every edge signing.
The development also proves the arbitrary-factor statement directly and
identifies the SPD specialization with the positive inverse square root.
It does not formalize the data-dependent factor or quantum-query results
that follow.  The repository's \texttt{formal/PRECONDITIONER.md} maps these
claims to the proofs and records the verification commands.

For data-dependent preconditioning, let $D_\varsigma$ be the square grounded
signed incidence matrix of the chain-plus-copy tree: its root row is $e_0^T$
and the row of a nonroot node $j$ is
$e_j^T-\tau_je_{\pi(j)}^T$, where $\tau_j$ is its edge sign.  With pair
variables the equality matrix is
\begin{equation}
 A_\varsigma=D_\varsigma[I,-I],\qquad
 C_\varsigma=D_\varsigma D_\varsigma^T.                  \label{eq:prec-tree-normal}
\end{equation}
Every row of $D_\varsigma$ has at most two nonzeros, each determined by one
local sign, yet
\begin{equation}
 D_\varsigma^{-1}C_\varsigma D_\varsigma^{-T}=I,\qquad
 D_\varsigma^{-1}e_0=(p_j)_j,                             \label{eq:prec-perfect-local}
\end{equation}
where $p_j$ is the root-to-$j$ sign product.  Set
$R_\varsigma=\diag(p_j)$, and let $D_+$ be the unsigned grounded incidence
matrix on the same rooted tree.  The switching identity is
$D_\varsigma=R_\varsigma D_+R_\varsigma$.  Thus construction of a sparse
perfect preconditioner is not hard; applying its inverse exposes all prefixes.
This distinction is the classical requirement that a useful preconditioner be
cheap to apply, expressed here in support-theory language
\cite{boman2003-support-theory-for-preconditioning}.

The right-hand side $e_0$ in \eqref{eq:prec-perfect-local} occurs in an actual
Newton equation.  At the infeasible-start point
$x=s=\mathbf1$, $y=0$, $\mu=1$ for the pair LP with $c=\mathbf1$ and
$b=e_0$, the point is centered and dual feasible, and $A_\varsigma\mathbf1=0$.
With the residual convention of \eqref{eq:newton-residuals}, any centering
target $0<\sigma_c\leq1$ gives
\begin{equation}
 A_\varsigma\Delta x=e_0,\quad
 A_\varsigma^T\Delta y+\Delta s=0,\quad
 \Delta x+\Delta s=(\sigma_c-1)\mathbf1,\quad
 A_\varsigma A_\varsigma^T\Delta y=e_0.                  \label{eq:prec-actual-newton}
\end{equation}
Since $A_\varsigma A_\varsigma^T=2D_\varsigma D_\varsigma^T$, the signed
incidence congruence transforms this exact Newton RHS into
$D_\varsigma^{-1}e_0$ up to a public scalar.

For the state lower bound, take $N\geq2$ and attach two isomorphic complete
copy trees, one to the root and one to the endpoint of the $N$-edge hidden
spine.  Each has $M$ leaves and $J=2M-2$ new nodes, where
\begin{equation}
 M=2^{\lceil\log_2(16(N+1))\rceil}.                       \label{eq:prec-two-tree-size}
\end{equation}
Here $M$ and $m=N+1$ below are local to this amplified-tree construction and
are unrelated to the plateau length $M$ and path dimension $m$ used earlier.
All new edges have positive sign; write $\widehat D_\varsigma$ for this
grounded incidence factor and define
\begin{equation}
 \widehat A_\varsigma=\widehat D_\varsigma[I,-I],\qquad
 \widehat C_\varsigma=\widehat D_\varsigma\widehat D_\varsigma^T.
 \label{eq:prec-amplified-normal}
\end{equation}
At the same infeasible-start point $x=s=\mathbf1$, $y=0$, $\mu=1$, now for
the amplified pair LP with $c=\mathbf1$ and $b=e_0$, the Newton equations are
\begin{equation}
 \widehat A_\varsigma\Delta x=e_0,\quad
 \widehat A_\varsigma^T\Delta y+\Delta s=0,\quad
 \Delta x+\Delta s=(\sigma_c-1)\mathbf1,\quad
 \widehat A_\varsigma\widehat A_\varsigma^T\Delta y=e_0.
 \label{eq:prec-amplified-newton}
\end{equation}
Thus the signed-incidence congruence maps this actual amplified Newton RHS to
$\widehat D_\varsigma^{-1}e_0$ up to a public scalar.

The weighted factor below is also the exact local factor away from the start.
At any positive interior point let
$\Theta=XS^{-1}=\diag(\Theta_+,\Theta_-)$, where
$\Theta_\pm=\diag(\theta_j^\pm)_j$, and set
\begin{equation}
 \Omega=\diag(\omega_j)_j,\qquad
 \omega_j=\theta^+_j+\theta^-_j>0.                       \label{eq:prec-local-omega}
\end{equation}
Then
\begin{equation}
 \widehat A_\varsigma\Theta\widehat A_\varsigma^T
 =\widehat D_\varsigma\Omega\widehat D_\varsigma^T
 =L_\varsigma L_\varsigma^T,\qquad
 L_\varsigma:=\widehat D_\varsigma\Omega^{1/2}.         \label{eq:prec-local-factor}
\end{equation}
Here $g$ below is the corresponding eliminated Newton RHS.  In particular,
$4/5\leq x_i,s_i\leq6/5$ implies
$4/3\leq\omega_j\leq3$, hence
$\Gamma:=\max_j\omega_j/\min_j\omega_j\leq9/4<4$.

\begin{theorem}[Factor-specific Newton-RHS hardness]
\label{thm:prec-factor-rhs-hardness}
Every unconditional algorithm that, on each sign input, outputs a state within
trace distance $1/20$ of
\begin{equation}
 \ket{z_\varsigma}=
 \frac{\widehat D_\varsigma^{-1}e_0}
      {\norm{\widehat D_\varsigma^{-1}e_0}_2}             \label{eq:prec-prefix-state}
\end{equation}
makes $\Omega(N)$ raw coefficient queries.

The conclusion remains true for the weighted nearby factor $L_\varsigma$ in
\eqref{eq:prec-local-factor} and eliminated RHS $g$ whenever
\begin{equation}
 \frac{\max_j\omega_j}{\min_j\omega_j}\leq4,\qquad
 \norm{g-e_0}_1\leq\frac1{100};                           \label{eq:prec-l1-robust}
\end{equation}
the target is the normalized state proportional to
$\Omega^{-1/2}\widehat D_\varsigma^{-1}g$.  If $g$ or $\Omega$ is
input-dependent, all raw queries used to construct and access it are included.
Thus the theorem applies at any nearby central iterate satisfying these two
promises.  Freely supplied side oracles whose numerical values themselves
encode parity are outside the theorem.
\end{theorem}

\begin{proof}
In the exact case, the prefix vector is $+1$ on every node of the root copy
tree and endpoint parity $p_N$ on every node of the endpoint copy tree.  Pair
corresponding nodes and measure their side qubit in the $X$ basis.  If
$m=N+1$, then $J\geq31m$ and the probability of a paired node is
\begin{equation}
 \frac{2J}{m+2J}>\frac{62}{63}.                           \label{eq:prec-copy-mass}
\end{equation}
Conditional on that event the measurement returns $p_N$ exactly; outputting a
fair bit otherwise succeeds with probability greater than
$1/2+31/63=125/126$.  Trace error $1/20$ leaves success above $2/3$, so the
fixed decoder computes parity and requires $\Omega(N)$ queries.

For \eqref{eq:prec-l1-robust}, first take $g=e_0$ and put
$\Gamma=\max\omega/\min\omega\leq4$.  Paired nodes have squared mass at
least
\[
 \alpha\geq\frac{2J}{2J+\Gamma m}\geq\frac{31}{33}.
\]
The amplitudes of a corresponding pair are $a$ and $p_Nb$ with
$a/b\in[1/2,2]$.  The $X$ measurement is correct with probability
$1/2+ab/(a^2+b^2)\geq9/10$.  The full ideal success is therefore at least
$1/2+(2/5)(31/33)=289/330$.

For $\delta=g-e_0$, every coordinate of
$\widehat D_\varsigma^{-1}\delta$ is a signed ancestor sum, so, if $P$ is
the number of tree nodes,
\[
 \norm{\widehat D_\varsigma^{-1}\delta}_2
 \leq\sqrt P\norm\delta_1,
 \qquad
 \frac{\norm{\Omega^{-1/2}\widehat D_\varsigma^{-1}\delta}_2}
 {\norm{\Omega^{-1/2}\widehat D_\varsigma^{-1}e_0}_2}
 \leq\sqrt\Gamma\norm\delta_1\leq\frac1{50}.
\]
The corresponding normalized pure states are at Euclidean distance at most
$1/25$.  Including trace error $1/20$ changes outcome probabilities by at
most $9/100$; subtracting this from $289/330$ leaves success above $2/3$.
Parity again proves the claim.
\end{proof}

The $\ell_1$ hypothesis is essential to this proof.  On the signed spine set
$\widetilde d_j=p_j(1-j/N)$, keep the root copy tree at value one, and set the
endpoint copy tree to zero.  Then
\begin{equation}
 \norm{\widehat D_\varsigma\widetilde d-e_0}_2=N^{-1/2},  \label{eq:prec-l2-failure}
\end{equation}
although the amplified endpoint has vanished and $\widetilde d$ stays a
constant relative distance from the prefix vector.  Thus even a vanishing
$\ell_2$ feasibility residual does not imply the state closeness used above;
\eqref{eq:prec-l2-failure} is a failure of this proof route, not a no-go for
every possible $\ell_2$-robust reduction.

Nor is transformed-RHS hardness invariant under the factor.  On the amplified
tree, let $\widehat R_\varsigma=\diag(p_j)$ be the prefix gauge and let
$\widehat D_+$ be the unsigned grounded incidence matrix.  Then
$\widehat D_\varsigma=\widehat R_\varsigma\widehat D_+
\widehat R_\varsigma$, and the equally perfect factor
\begin{equation}
 \widehat P_\varsigma=\widehat D_+^{-1}\widehat R_\varsigma
 =\widehat R_\varsigma\widehat D_\varsigma^{-1}
 \label{eq:prec-orthogonal-loophole}
\end{equation}
satisfies
\begin{equation}
 \widehat P_\varsigma\widehat C_\varsigma\widehat P_\varsigma^T=I,
 \qquad
 \widehat P_\varsigma e_0=\widehat D_+^{-1}e_0=\mathbf1.
 \label{eq:prec-easy-rhs}
\end{equation}
The RHS is now input independent.  Recovery uses
$\widehat P_\varsigma^T=\widehat R_\varsigma\widehat D_+^{-T}$, so the
parity cost has moved to factor application or original-coordinate recovery.
This orthogonal-factor
loophole is why \cref{thm:prec-factor-rhs-hardness} is deliberately
factor-specific.

The invariant statement is end to end.  Specialize
\eqref{eq:prec-amplified-newton} to $\sigma_c=1$ and take the full step.  With
$p=\widehat D_\varsigma^{-1}e_0$, the unique minimum-norm direction and its
update are
\begin{equation}
 \Delta x=\tfrac12(p,-p),\qquad \Delta s=-\Delta x,\qquad
 x^+=\mathbf1+\Delta x.                                  \label{eq:prec-hard-update}
\end{equation}
Every pair of $x^+$ is $(3/2,1/2)$ or its swap, and
$X^+s^+=(3/4)\mathbf1$; this is the exact feasible center at $\mu=3/4$.
The endpoint copy tree occupies more than $15/32$ of all pairs.  Conditional
on measuring such a pair, the computational-basis coordinate identifies
$p_N$ with probability $9/10$.  Returning a fair bit otherwise gives ideal
success greater than $1/2+(2/5)(15/32)=11/16$, and trace error $1/100$
leaves success above $2/3$.

\begin{theorem}[Preconditioner-oracle accounting]
\label{thm:prec-parity-ledger}
Let $P_\varsigma$ be any possibly data-dependent invertible factor for the
normal equation in \eqref{eq:prec-amplified-newton}, with no sparsity or locality
assumption, and let an end-to-end Newton procedure output a state within trace
distance $1/100$ of the normalized $x^+$ in
\eqref{eq:prec-hard-update}.  Partition all raw coefficient queries into
setup, transformed-RHS preparation, iterative solve, and original-coordinate
recovery/update.  Then
\begin{equation}
 \boxed{Q_{\rm set}+Q_{\rm rhs}+Q_{\rm solve}+Q_{\rm rec}=\Omega(N).}
 \label{eq:prec-parity-ledger}
\end{equation}
If the preconditioned system has condition at most $K$ and the solve phase's
only raw queries occur inside at most $T(K,\epsilon)$ calls to its
preconditioned oracles, each costing at most $q_{\rm on}$ raw queries, then
\begin{equation}
 Q_{\rm set}+Q_{\rm rhs}+T(K,\epsilon)q_{\rm on}+Q_{\rm rec}=\Omega(N).
 \label{eq:prec-condition-ledger}
\end{equation}
In particular, for
$T(K,\epsilon)=\widetilde O(K\log(1/\epsilon))$ and hidden polylogarithmic
factor $L(N,\epsilon)$,
\begin{equation}
 Q+O(Kq_{\rm on}L(N,\epsilon))=\Omega(N),\qquad
 Q:=Q_{\rm set}+Q_{\rm rhs}+Q_{\rm rec}.                 \label{eq:prec-condition-corollary}
\end{equation}
\end{theorem}

\begin{proof}
All four phases together form one coefficient-query algorithm.  Compose its
output with the fixed decoder following \eqref{eq:prec-hard-update}; this
computes parity of the $N$ hidden spine signs with bounded error.  Each LP
coefficient query is simulated by $O(1)$ sign queries, proving
\eqref{eq:prec-parity-ledger}.  Under the oracle-only solve hypothesis,
$Q_{\rm solve}\leq T(K,\epsilon)q_{\rm on}$, which gives
\eqref{eq:prec-condition-ledger}, and substitution of the stated solver call
bound gives \eqref{eq:prec-condition-corollary}.
\end{proof}

A hybrid solver that also queries raw data directly remains covered by
\eqref{eq:prec-parity-ledger}, where such queries stay in $Q_{\rm solve}$.

The ledger is tight and is not a quantum--classical separation.  A classical
algorithm reads the $N$ signs, computes every prefix, and stores the gauge in
$\Theta(N)$ queries and work.  Ladner--Fischer prefix scan reduces parallel
depth to $O(\log N)$ with $O(N)$ bounded-fan-in work
\cite{ladner1980-parallel-prefix-computation}.  Caching then removes further
hidden-sign queries, but it does not make the unsigned triangular operations
free: applying $\widehat D_+^{-1}$ or $\widehat D_+^{-T}$ still takes
$\Theta(N)$ work and $O(\log N)$ parallel depth by prefix or subtree
aggregation.  Multiplication by the cached diagonal gauge is local.  Quantum
parity changes only the constant in the query count.  The right referee-level
interpretation is therefore a
lemma/counterexample suite: \cref{thm:prec-fixed-minimax} is a finite-dimensional
minimax lower bound, \cref{thm:prec-factor-rhs-hardness} locates cost for a canonical factor,
and \cref{thm:prec-parity-ledger} records how another factor can relocate it.
It is not a standalone universal preconditioning lower bound.  Contemporary
preconditioned QIPM analyses that assume prepared transformed data are
consistent with this conclusion; reducing those oracles to raw LP access is a
separate charged step \cite{zeguan2026-preconditioned-inexact-infeasible-quantum-interi}.

\subsection{Newton-state output is not an iterate oracle}
\label{sec:preconditioning-state-interface}

We now change oracle models.  Let $v\in\R_{>0}^n$, let its known norm be
$r_v=\norm v_2$, and grant controlled state-preparation access
\begin{equation}
 U_v\ket0=\ket{\widehat v}=\frac1{r_v}\sum_{j=1}^nv_j\ket j,
 \label{eq:prec-state-oracle}
\end{equation}
together with $U_v^\dagger$ and a fixed real phase.  This is precisely
\cref{def:state-preparation-access}; it is stronger than receiving copies but
distinct from coordinate-value, raw-coefficient, QRAM, or already compiled
diagonal access.  Those interfaces are not generally ordered by strength.
The next results count calls to $U_v,U_v^\dagger$, not raw LP queries.

\begin{theorem}[State-to-diagonal normalization--query product]
\label{thm:prec-state-diagonal}
Suppose a controlled circuit $W_v$, using $q$ calls to
$U_v,U_v^\dagger$, is an $(\alpha,a,1/8)$ block encoding of
$\diag(v)$.  If the reported normalization satisfies the uniform known bound
$\alpha\leq A$, where $A\geq1$, then in the worst case
\begin{equation}
 \boxed{qA=\Omega(\sqrt n)}.                              \label{eq:prec-qA}
\end{equation}
The product is tight in order on the hard near-uniform family.
\end{theorem}

\begin{proof}
For $n\geq10$, take the positive equal-norm pair
\begin{equation}
 v^{(0)}=(1,\ldots,1),\qquad
 v^{(1)}=(2,\beta,\ldots,\beta),\qquad
 \beta=\sqrt{\frac{n-4}{n-1}}.                           \label{eq:prec-state-pair}
\end{equation}
Both norms are $\sqrt n$, and
\begin{equation}
 \frac1{\sqrt n}\leq
 \norm{v^{(0)}/\sqrt n-v^{(1)}/\sqrt n}_2
 \leq\sqrt{2/n}.                                        \label{eq:prec-close-states}
\end{equation}
Choose unitary extensions $U_0,U_1$ related by the minimum planar rotation
between the prepared states.  Then
$\norm{U_0-U_1}_2=\Theta(n^{-1/2})$, so the hybrid argument requires
$\Omega(\sqrt n)$ oracle calls for constant-bias distinction.

On the other hand, $O(A)$ controlled calls to $W_v$ estimate
$\operatorname{Re}\bra{0^a,1}W_v\ket{0^a,1}$ to additive error $1/(16A)$.
Multiplication by the reported $\alpha$ contributes error at most $1/16$,
and block-encoding error contributes at most $1/8$.  The first diagonal entry
is therefore known within $1/4$, distinguishing one from two.  The composed
distinguisher uses $O(qA)$ original oracle calls, proving
\eqref{eq:prec-qA}.

Rattew and Rebentrost give an exact constant-query block encoding of
$\diag(v/r_v)$ from $U_v$ \cite{rattew2023-non-linear-transformations-of-quantum}.
As an encoding of $\diag(v)$ it has normalization $r_v=\sqrt n$ on
\eqref{eq:prec-state-pair}, attaining the product in order.  This is
worst-case tightness, not a per-vector lower bound for every structured
preparation circuit.
\end{proof}

With copy-only access, the pair in \eqref{eq:prec-state-pair} requires
$\Omega(n)$ copies: one-copy overlap is $1-\Theta(1/n)$, so only
$k=\Omega(n)$ copies have constant trace distance.  Thus inverse access in
\cref{def:state-preparation-access} changes the exponent but does not erase
the interface cost.

The same obstruction sharpens near a strictly complementary endpoint.  Let
$n=2m$, $m\geq5$, $0<\mu\leq1$, and define
\begin{equation}
 \beta=\sqrt{\frac{m-4}{m-1}},\quad
 v=(2,\beta\mathbf1_{m-1}),\quad
 r=(m-1)(1-\beta)=\frac3{1+\beta},\quad
 z=2(r,\mathbf1_{m-1}).                                  \label{eq:prec-central-pair}
\end{equation}
For the one-sparse LP
\begin{equation}
 A=[I_m\;0],\qquad b=\mathbf1_m,\qquad c=(z,\mathbf1_m), \label{eq:prec-central-lp}
\end{equation}
take the common $x=(\mathbf1_m,\mu\mathbf1_m)$ and
\begin{align}
 s^{(0)}&=(\mu\mathbf1_m,\mathbf1_m),&
 y^{(0)}&=z-\mu\mathbf1_m,\nonumber\\
 s^{(1)}&=(\mu v,\mathbf1_m),&
 y^{(1)}&=z-\mu v.                                      \label{eq:prec-central-triples}
\end{align}
Both triples are feasible.  Since $\norm v=\norm{\mathbf1_m}$ and
$z^T(v-\mathbf1_m)=0$, the corresponding slack norms and dual norms agree.

\begin{theorem}[Strict-complementarity state-access barrier]
\label{thm:prec-centrality-state}
In controlled state-preparation access to the slack and dual states in
\eqref{eq:prec-central-triples}, distinguishing
\begin{equation}
 \norm{Xs-\bar\mu\mathbf1}_2\leq\bar\mu/4
 \quad\text{from}\quad
 \norm{Xs-\bar\mu\mathbf1}_2\geq3\bar\mu/4              \label{eq:prec-centrality-promise}
\end{equation}
requires
\begin{equation}
 \boxed{\Omega(\sqrt n/\mu)}                             \label{eq:prec-centrality-lb}
\end{equation}
total controlled preparation/inverse calls.  With copy-only access the bound
is $\Omega(n/\mu^2)$.  The result holds although both triples approach the
same strictly complementary optimum.  The fixed LP data, the common $x$, all
vector norms, positivity, and the real phase conventions are free.
\end{theorem}

\begin{proof}
The zero triple is exactly centered with $\bar\mu_0=\mu$.  For the other,
$\bar\mu_1<\mu$ while the first complementarity product is $2\mu$, so
\begin{equation}
 \norm{Xs^{(1)}-\bar\mu_1\mathbf1}_2>\mu>\bar\mu_1.       \label{eq:prec-central-gap}
\end{equation}
The normalized slack-state distance is $\Theta(\mu/\sqrt n)$.  The duals
differ by the same unnormalized vector $\mu(v-\mathbf1_m)$ and have common
norm $\Omega(\sqrt m)$, so their normalized-state distance is also
$O(\mu/\sqrt n)$.  Choose close minimum-rotation extensions for both
oracles.  A hybrid over queries to either oracle gives
\eqref{eq:prec-centrality-lb}.  For independent copies, fidelity
multiplicativity gives the square of this dependence.  Finally both triples
converge to
$((\mathbf1_m,0),z,(0,\mathbf1_m))$, proving strict complementarity.
\end{proof}

The triples can also be generated by a sparse LP's objective-state oracle.
Put $d=v-\mathbf1_m$,
$z_0=2(r,\mathbf1_{m-1})$, and $z_\mu=z_0-(\mu/2)d$.  Then
\begin{equation}
 z_0^Td=0,\qquad \mathbf1^Td=1-r,\qquad 1\leq\norm d_2^2\leq2.
 \label{eq:prec-input-identities}
\end{equation}
\begin{theorem}[Input-generated state-only barrier]
\label{thm:prec-input-generated}
Use the common one-sparse $A=[I_m\;0]$, $b=\mathbf1_m$, and objectives
\begin{equation}
 c^{(0)}=(z_\mu,\mathbf1_m),\qquad
 c^{(1)}=(z_\mu+\mu d,\mathbf1_m).                       \label{eq:prec-input-objectives}
\end{equation}
Let the only input-dependent access be a controlled clean preparation of the
normalized objective and its inverse; the common norm, $A$, and $b$ are free.
Distinguishing the inputs, certifying which one is centered at the common pair
$x=(\mathbf1_m,\mu\mathbf1_m)$,
$y=z_\mu-\mu\mathbf1_m$, or estimating the optimum to additive error
$\mu/8$, needs $\Omega(\sqrt n/\mu)$ controlled oracle calls.  Copy-only
access needs $\Omega(n/\mu^2)$ copies.

Moreover, multiply both objectives symmetrically by $n$, take
$b=\mathbf1_m/6$, and use
\[
 \widetilde c^{(0)}=n(z_0-\tfrac\mu2d,\mathbf1_m),\qquad
 \widetilde c^{(1)}=n(z_0+\tfrac\mu2d,\mathbf1_m).
\]
The normalized objective oracles do not change, the value gap obeys
\begin{equation}
 |\operatorname{OPT}_1-\operatorname{OPT}_0|
 =\frac{n\mu}{6}(r-1)\geq\frac{n\mu}{12},                \label{eq:prec-value-gap}
\end{equation}
and additive value accuracy $n\mu/48$ has the same query lower bound.  This
is a fixed fraction of the centered total duality gap $n\mu$; here
$n=\nu$ is the canonical LP barrier parameter of \cref{sec:prelim-conic}.
\end{theorem}

\begin{proof}
The objective norms agree because
$2\mu z_\mu^Td+\mu^2\norm d^2=0$, and their normalized-state distance is
$\Theta(\mu/\sqrt n)$.  Minimum-rotation objective-state oracles therefore
require $\Omega(\sqrt n/\mu)$ calls to distinguish; fidelity
multiplicativity gives $\Omega(n/\mu^2)$ copies.  The displayed common pair
is centered for input zero and violates \eqref{eq:prec-centrality-promise}
for input one.  Every feasible point has first block $\mathbf1_m$ and its
positive-cost second block vanishes at optimum, so the original value gap is
$\mu|\mathbf1^Td|\geq\mu/2$.  Additive error $\mu/8$ distinguishes the
inputs.

For the scaled claim, the common point
\[
 \widetilde x=(\mathbf1_m/6,(\mu/n)\mathbf1_m),\qquad
 \widetilde y=n(z_0-\tfrac\mu2d)-6\mu\mathbf1_m
\]
has every complementarity product equal to $\mu$ for input zero; input one
remains positive because $n(1-\beta)\leq5<6$ and lies outside every fixed
$N_2$ neighborhood of radius below one.  Its centered total duality gap is
$n\mu$.  Feasibility fixes the first primal block to $\mathbf1_m/6$, which
gives \eqref{eq:prec-value-gap}; accuracy $n\mu/48$ distinguishes the inputs.
The inverse-$\mu$ factor is therefore not caused by demanding accuracy finer
than the centered duality-gap scale.
\end{proof}

This input-generated theorem has a sharp boundary.  A binary coefficient-value
oracle exposes the known exceptional coordinate, whose values differ by
$\Theta(\mu)$, in one query at sufficient requested precision.  Distinct
value-register strings are not close unitaries merely because the encoded
reals are close.  Likewise, the naturally normalized Newton diagonals for the
two inputs differ by a constant in their first entry, so a freely supplied
Newton block encoding distinguishes them in $O(1)$ calls.  The
$\Omega(\sqrt n/\mu)$ theorem is therefore an end-to-end result only for the
state-only objective/update interface; it does not lift to raw exact-value or
precompiled-Newton access.

State access also obstructs fraction-to-boundary line search.

\begin{corollary}[Fraction-to-boundary]
\label{cor:prec-fraction-boundary}
For $n\geq4$, let $x=\mathbf1_n$ and compare the equal-norm directions
\[
 \delta^{(0)}=-\tfrac12\mathbf1_n,\qquad
 \delta^{(1)}=(-1,-\gamma,\ldots,-\gamma),\qquad
 \gamma=\tfrac12\sqrt{\frac{n-4}{n-1}}.
\]
Their maximal positive steps are two and one, while their normalized states
are $\Theta(n^{-1/2})$ apart.  A state-only rule that is never larger than one
on the second input but is at least $3/2$ on the first needs
$\Omega(\sqrt n)$ controlled state-preparation calls.
\end{corollary}

\begin{proof}
Choose the two preparation unitaries as minimum rotations.  Any valid output
distinguishes them with constant bias, while a $Q$-query hybrid has distance
$O(Q/\sqrt n)$.  Thus $Q=\Omega(\sqrt n)$.
\end{proof}

The next theorem rules out bypassing diagonal compilation by directly solving
the following easy system.  For $n\geq4$, fix $0<\eta\leq1/4$ and define
\begin{equation}
 v^{(0)}=\mathbf1_n,\qquad
 v^{(1)}=(1+\eta,\beta_\eta,\ldots,\beta_\eta),\qquad
 \beta_\eta=\sqrt{1-\frac{2\eta+\eta^2}{n-1}}.           \label{eq:prec-transition-pair}
\end{equation}
Both norms equal $\sqrt n$, and their normalized distance $d_\eta$ satisfies
\begin{equation}
 \frac\eta{\sqrt n}\leq d_\eta\leq\frac{2\eta}{\sqrt n}. \label{eq:prec-transition-distance}
\end{equation}
Let $D_b=\diag(v^{(b)})$,
$r=(e_1+e_2)/\sqrt2$, and
$\ket{z_b}=D_b^{-1}r/\norm{D_b^{-1}r}_2$.  These systems are one-sparse and
$\kappa(D_b)<2$.

\begin{theorem}[State to next solve]
\label{thm:prec-state-next-solve}
Give an algorithm controlled access to the minimum-rotation preparations
$U_b,U_b^\dagger$ of the vectors in \eqref{eq:prec-transition-pair}.  If its
output $\rho_b$ satisfies
\begin{equation}
 \Dtr(\rho_b,\ket{z_b}\!\bra{z_b})\leq\eta/12            \label{eq:prec-next-solve-error}
\end{equation}
for both inputs, then it uses $\Omega(\sqrt n)$ oracle calls.  With copy-only
iterate access it needs $\Omega(n)$ copies.  In particular, fixed
$\eta=1/4$ is a constant-accuracy lower bound for a condition-below-two
system.
\end{theorem}

\begin{proof}
For $b=0$, measuring coordinate one in $\ket{z_b}$ has probability $1/2$.
For $b=1$, put $q=\beta_\eta/(1+\eta)$.  The probability is
$q^2/(1+q^2)$, and
\begin{equation}
 \frac12-\frac{q^2}{1+q^2}
 \geq\frac12-\frac1{1+(1+\eta)^2}\geq\frac\eta3.         \label{eq:prec-solution-separation}
\end{equation}
Contractivity and \eqref{eq:prec-next-solve-error} therefore make the two
algorithm outputs at least $\eta/6$ apart in trace distance.  A $Q$-query
hybrid is at most $Qd_\eta\leq2Q\eta/\sqrt n$ apart, whence
$Q\geq\sqrt n/12$.  For $k$ copies the input trace distance is at most
$\sqrt{k}d_\eta=O(\sqrt{k}\eta/\sqrt n)$, forcing $k=\Omega(n)$.
\end{proof}

The Hermitian dilation
$\left(\begin{smallmatrix}0&D_b\\D_b&0\end{smallmatrix}\right)$ has the
same sparsity and condition number, so this is also a standard Hermitian QLS
instance.  It does not contradict a polylogarithmic-dimensional QLS
algorithm: such an algorithm assumes the matrix oracle that this theorem
requires one to construct or bypass from the iterate state.

Per-round costs add only under an explicit fresh-oracle contract.  At round
$t=1,\ldots,T$, after earlier outputs are committed, an adversary chooses a
fresh bit $b_t$ independent of the current workspace and supplies only a new
minimum-rotation pair $U_{t,b_t},U_{t,b_t}^\dagger$ for
\eqref{eq:prec-transition-pair}; future oracles are unavailable.  Before the
next round, the algorithm must either satisfy
\eqref{eq:prec-next-solve-error}, or, after $S_t$ setup queries, expose a
reusable controlled $(\alpha_t,a,\eta/8)$ block encoding of $D_{b_t}$ with
$\alpha_t\leq A_t$.  One invocation may use another $q_t$ state-oracle calls.

\begin{theorem}[Online fresh-oracle direct sum]
\label{thm:prec-online-direct-sum}
In the preceding interface, every fixed worst-case per-round budget satisfies
\begin{equation}
 \sum_{t=1}^TQ_t=\Omega(T\sqrt n)                         \label{eq:prec-online-solves}
\end{equation}
for next-solution-state outputs, and
\begin{equation}
 \boxed{\sum_{t=1}^T(\eta S_t+q_tA_t)=\Omega(T\sqrt n)}  \label{eq:prec-online-compilers}
\end{equation}
for reusable diagonal compilers.  For fixed short-step $\eta$, the latter is
equivalently $\sum_t(S_t+q_tA_t)=\Omega(T\sqrt n)$.
\end{theorem}

\begin{proof}
Condition on the complete workspace and transcript at the start of round
$t$.  Freshness makes that side information independent of $b_t$, so it is a
fixed ancilla in \cref{thm:prec-state-next-solve}, which gives
$Q_t=\Omega(\sqrt n)$.  For a compiler, run setup once and estimate the first
encoded entry to $O(\eta/A_t)$ using $O(A_t/\eta)$ invocations.  This
distinguisher makes $S_t+O(q_tA_t/\eta)$ state-oracle calls, but the close
input oracles require $\Omega(\sqrt n/\eta)$ calls.  Hence
$\eta S_t+q_tA_t=\Omega(\sqrt n)$ in each conditioned round.  Summation
proves both claims.
\end{proof}

The full oracle ledger in \cref{thm:prec-online-direct-sum} is load-bearing:
\begin{itemize}[leftmargin=*]
 \item The atomic charged operations are controlled calls to the current
 $U_{t,b_t}$ and its inverse.  Gate complexity is not lower-bounded.
 \item The common norm $\sqrt n$, the two vector formulas, $r$, and the round
 number are free; only $b_t$ is hidden.
 \item No coordinate-value, QRAM, objective-coefficient, or compiled matrix
 oracle is available.  Each could reveal the exceptional coordinate in one
 query.
 \item The complete off-zero action is promised: the adversary may choose the
 close minimum-rotation completions.  Specifying only $U\ket0$ would not
 define a query lower bound.
 \item One-time setup and reusable advice are charged through $S_t$, while
 per-invocation use is charged through $q_t$.  Advice consumed by calls must
 be prepared in sufficient quantity and cannot be cloned for free.
 \item The required output is committed before the next fresh oracle arrives;
 this causality makes the conditional bounds additive.
\end{itemize}
A heralded compiler with success probability bounded below obeys the same
ledger after charging repetitions.  An unheralded constant failure
probability is not covered when $\eta\to0$, because it can exceed the target
states' $\Theta(\eta)$ separation.

The hard stream remains positive and obeys a coordinatewise short-update
promise: $|v_i^{(b_{t+1})}-v_i^{(b_t)}|\leq\eta v_i^{(b_t)}$.  Thus
positivity and short relative motion do not defeat the online theorem.  They
also do not turn it into a fixed-instance QIPM iteration bound.  The exact
obstructions are:
\begin{enumerate}[leftmargin=*]
 \item \emph{Freshness is stronger than a fixed optimization instance.}
 Every QIPM iterate is determined by one static input and prior computation;
 no independent oracle normally arrives after each output.
 \item \emph{One-vector worst-case hardness is not automatically a direct
 sum.}  If every iterate depends on one bit, learning it once makes every
 later diagonal easy.
 \item \emph{Circuit generation must be expanded to atomic queries.}  A
 circuit for $U_{x_{t+1}}$ may call $U_{x_t}$, which contains the prior
 history.  Unit-cost treatment hides inlining, while charging both as
 independent primitives double-counts it.
 \item \emph{Off-zero leakage is model dependent.}  The hybrid chooses close
 extensions, whereas a concrete solver circuit may reveal coordinate or
 trajectory data away from $\ket0$.  Its entire oracle action must be fixed.
 \item \emph{All-round coherent access changes the direct-sum problem.}  If a
 single multiplexed query reaches every time-indexed oracle, the causal proof
 fails and a separate coherent direct-sum theorem is needed.
 \item \emph{The QIPM equations may permit a bypass.}  The diagonal systems
 above have not been embedded as consecutive Newton systems of one fixed
 sparse LP; cancellation, scale-free formulations, or history-state methods
 might avoid separate compilation.
 \item \emph{The output contract matters.}  An architecture that exposes a
 reusable diagonal or a Newton state every round meets the theorem's
 contract.  An algorithm that proves its step analytically and emits only one
 final observable need not.
\end{enumerate}
Therefore a fixed-LP $\Omega(T\sqrt n)$ result would require $T$ independent
hard pieces in one sparse instance, a trajectory that reveals them one by
one, and a direct-sum theorem against the static input oracle.  None is proved
here.

Finally, coherent update errors accumulate.  If
\begin{equation}
 \norm{\ket{\widetilde v}-\ket{v/r_v}}_2\leq\xi,\qquad
 |\widetilde r_v-r_v|\leq\delta_r,                        \label{eq:prec-update-input-error}
\end{equation}
amplitude-to-diagonal lifting has absolute operator error at most
$\delta_r+\widetilde r_v\xi$.  For coherent LCU updates
$X_{t+1}=X_t+h_t\Delta_t$, let
$E_t:=\norm{\widetilde X_t-X_t}_2$ be the accumulated absolute operator error
before update $t$.  If the step lengths and LCU arithmetic are exact and the
only new error is the lifted direction error, the operator-norm triangle
inequality gives
\begin{equation}
 E_{t+1}\leq E_t+|h_t|
 (\delta_{r,t}+\widetilde r_t\xi_t),\qquad
 E_T\leq E_0+\sum_{t<T}|h_t|
 (\delta_{r,t}+\widetilde r_t\xi_t).                     \label{eq:prec-update-error-ledger}
\end{equation}
Rounding, imperfect LCU coefficients, or postselection errors must be added as
separate terms.
Late-path centrality therefore constrains the cumulative weighted error, not
only each local preparation.  Literal circuit reuse nests as well.  Let $C_t$
be the cost of the current iterate-preparation circuit, let $Q_t$ be the
number of its calls made by the next direction circuit, and let $G_t$ be all
nonrecursive work in that update.  With literal inlining, additive uniform
gate/query costs, and no sharing or memoization, the exact accounting is
\begin{equation}
 C_{t+1}=G_t+(1+Q_t)C_t.                                 \label{eq:prec-circuit-accounting}
\end{equation}
This identity is for literal composition, not a universal lower bound against
direct structured or history-state oracles.

\subsection{Lazy central oracles and the caching ceiling}
\label{sec:preconditioning-lazy}

Changing all diagonal weights does not imply that they must all be rewritten.
The following fixed sparse LP makes this failure explicit.  For $K$ disjoint
blocks of length $L$, let
$d_{a,i}=u_{a,i}-v_{a,i}$ and impose
\begin{equation}
 d_{a,0}=1,\qquad
 d_{a,i}-\varsigma_{a,i}d_{a,i-1}=0\quad(1\leq i\leq L),  \label{eq:prec-lazy-paths}
\end{equation}
with all $u,v\geq0$.  Minimize
\begin{equation}
 \sum_{a=1}^K\left[
  \sum_{i=0}^L(u_{a,i}+v_{a,i})
  +\alpha(u_{a,L}-v_{a,L})\right],\qquad 0<\alpha<1.       \label{eq:prec-lazy-objective}
\end{equation}
Rows and columns have constant sparsity and the instance size is
$\Theta(KL)$.  With
\begin{equation}
 p_{a,0}=1,\qquad p_{a,i}=\prod_{j=1}^i\varsigma_{a,j},\qquad
 q(\mu)=\mu+\sqrt{\mu^2+1},                              \label{eq:prec-lazy-prefix}
\end{equation}
the primal central path is exactly
\begin{equation}
 (u_{a,i}(\mu),v_{a,i}(\mu))
 =\left(\frac{q(\mu)+p_{a,i}}2,
         \frac{q(\mu)-p_{a,i}}2\right).                  \label{eq:prec-lazy-factorization}
\end{equation}
Indeed feasibility fixes the difference $p_{a,i}$; differentiating the
one-dimensional pair barrier gives the common sum $q(\mu)$.  In the fixed
orthonormal null basis
$(e_{u_{a,i}}+e_{v_{a,i}})/\sqrt2$, the reduced log-barrier Hessian
$W^\top X^{-2}W$ is the scalar
\[
 2[(q(\mu)+1)^{-2}+(q(\mu)-1)^{-2}]I,
\]
so it has condition one throughout the trajectory.  All changing numerical
information is the public scalar $q(\mu)$; all hidden structural information
is the static prefix table.

Give reversible raw access to bits
$\varsigma_{a,i}=(-1)^{b_{a,i}}$.  An endpoint central lookup at block $a$
returns \eqref{eq:prec-lazy-factorization} for $i=L$ and therefore reveals
$P_a=p_{a,L}$, even with coordinate error strictly below $1/2$.

\begin{theorem}[Lazy endpoint lookup]
\label{thm:prec-lazy-lookup}
One exact coherent endpoint lookup, or one whose decoded sign is correct with
probability at least $2/3$, needs $\Omega(L)$ raw queries in the worst case,
even when the block index is in superposition.  The bound is tight.
\end{theorem}

\begin{proof}
Fix the block register to any basis value $a$ and fix all other blocks.
Decoding the returned pair computes parity of the remaining $L$ bits, which
has bounded-error quantum query complexity $\Omega(L)$.  Conversely, the
exact quantum parity algorithm computes the selected block's parity in
$O(L)$ queries while carrying $a$ coherently as a spectator; reversible
arithmetic then combines it with $q(\mu)$.  The same circuit works in
superposition over $a$.
\end{proof}

Call a resource an exact reusable endpoint oracle if, once constructed, a
fixed processor can obtain every one of the $K$ endpoint values without
further raw queries, and the resource survives or otherwise supports those
$K$ invocations.

\begin{theorem}[Reusable endpoint preprocessing]
\label{thm:prec-reusable-endpoints}
Constructing an exact reusable endpoint oracle needs at least
\begin{equation}
 \boxed{\left\lceil\frac{KL}{2}\right\rceil}             \label{eq:prec-reusable-cost}
\end{equation}
raw queries.  A bounded-error resource also needs $\Omega(KL)$ queries when
all $K$ successive answers are jointly correct with probability at least
$2/3$ on every input.
\end{theorem}

\begin{proof}
After preprocessing, invoke the resource on every block and multiply the
decoded signs.  The result is
\[
 \prod_{a=1}^KP_a=(-1)^{\bigoplus_{a=1}^K\bigoplus_{i=1}^Lb_{a,i}},
\]
the parity of all $KL$ raw bits.  Exact quantum parity requires
$\lceil KL/2\rceil$ queries, and bounded-error parity requires $\Omega(KL)$.
The proof permits entangled program states and implicit formulas; reusability,
not classical table materialization, is the decisive promise.
\end{proof}

\begin{theorem}[Explicit-coordinate trajectory service]
\label{thm:prec-lazy-trajectory}
Fix arbitrary parameters $\mu_1,\ldots,\mu_K>0$.  A service that at stage
$a$ returns the classical endpoint pair of block $a$ at $\mu_a$, with both
coordinate errors below $1/2$ and all answers jointly correct with probability
at least $2/3$ on every input, uses
\begin{equation}
 \boxed{\Omega(KL)}                                      \label{eq:prec-trajectory-cost}
\end{equation}
total raw queries across preprocessing and all stages.  Thus its amortized
cost is $\Omega(L)$.
\end{theorem}

\begin{proof}
Each returned pair reveals the independent $P_a$.  Multiplying the $K$
answers computes parity of all $KL$ input bits, so bounded-error quantum
parity proves \eqref{eq:prec-trajectory-cost}.
\end{proof}

Theorem~\ref{thm:prec-lazy-trajectory} is a lower bound for an explicit-cell
service contract, not for every QIPM.  Quantum cell-probe lower bounds likewise
require a specified storage-and-probe model and do not transfer automatically
to an analytic central path \cite{sen2001-lower-bounds-quantum-cell-probe}.
There is no unconditional per-IPM-iteration refresh theorem here, for five
separate reasons:
\begin{enumerate}[leftmargin=*]
 \item A fixed LP has no fresh per-iteration information.  Once all $KL$
 bits have been processed, unrestricted persistent memory can amortize that
 setup over arbitrarily many iterations.
 \item The path parameter may be low dimensional.  In
 \eqref{eq:prec-lazy-factorization} every changing coordinate is generated by
 one scalar $q(\mu)$, so a stored structural oracle needs only $O(1)$ extra
 arithmetic when $\mu$ changes.
 \item Cell-probe lower bounds need a query contract.  An IPM follows an
 analytic, algorithm-chosen trajectory rather than an adversarial sequence of
 cell updates and requests; \cref{thm:prec-lazy-trajectory} applies only
 after the distinct endpoint outputs are required.
 \item Coherent lookup defeats a naive block direct sum.  One $O(L)$ circuit
 evaluates the selected block in superposition.  A factor $K$ requires a
 classical-output contract or a suitable direct-product theorem.
 \item A normalized state is weaker than a data structure.  It need not
 support reusable random access to every coordinate; assuming that it does
 silently adds tomography or QRAM construction.
\end{enumerate}
This is the same caching ceiling as \eqref{eq:total-query-ledger} and the
static ledger following \cref{thm:lp-gain-plateau}: caching all signs can make
later raw-query cost vanish.  A genuine refresh lower bound would require a
memory cap, mandatory coordinate monitoring, nonshareable independent
invocations, or a proven online-hard fixed-LP trajectory.  Sparsity and
changing diagonal weights alone provide none of these.

\subsection{Synthesis}
\label{sec:preconditioning-synthesis}

Across all five interfaces, the invariant is a product or a ledger.  Separate
exact factor access gives $\alpha_P\sqrt K=\Omega(N)$ and, with reusable
compilation, $Ns+q\alpha=\Omega(N^2)$.  Scalar congruence obeys
$\sqrt\kappa\,\chi\geq c/\mu$; equality in order holds exactly for the
maximal class $r_\mu=\Omega(\mu)$ and $r_\mu=O(1)$.  The parity instance
preserves the total setup--RHS--solve--recovery cost even when a perfect factor
makes one phase trivial.  State preparation obeys $qA=\Omega(\sqrt n)$, and its iterative
version becomes additive only under a fresh-oracle contract.  Lazy evaluation
then shows why static preprocessing can defeat any unqualified per-iteration
claim.

These are access-model obstructions, not a verdict against preconditioning or
state reuse.  The positive modules in \cref{sec:positive-modules} are designed
to respect the same ledgers: they directly assemble favorable operators,
identify which projectors or cached structures are charged, and state what
output interface is actually produced.

\section{Positive Algorithmic Modules}
\label{sec:positive-modules}

The lower bounds in \cref{sec:lp-hardness,sec:preconditioning} do not rule out
useful quantum subroutines.  They say what a positive result must count.  This
section gives such results at their natural interfaces: a pathwise reuse
theorem, a partition-free tail solver and face repair, an implicit-coordinate
dual method with compressed output, and an explicit border solver for
block-angular systems.  Some compose to conditional trajectory theorems; none
is an unconditional end-to-end speedup for arbitrary sparse LPs.

We use the following rules throughout.  We refer to these as the referee
rules; they are also a checklist for assessing any positive QIPM claim.
\begin{enumerate}[leftmargin=*]
 \item A small spectral condition number does not make loading, normalization,
 RHS preparation, or recovery free.  Every block encoding and output compiler
 must be expanded into the access model of
 \cref{sec:prelim-quantum-access,sec:prelim-output-contracts}.
 \item A strict-complementarity construction applies only after its stated,
 data-dependent tail threshold; the separation margin and active singular gaps
 are parameters, not consequences of sparsity.
 \item Compressed dual output is dimensionally useful only when the retained
 dimension satisfies $m=o(n)$, and only if the classical comparator is
 charged in the same model.
 \item A coherent primal-coordinate oracle is substantive output.  When its
 values are well spread, constant-overhead value-to-amplitude conversion makes
 it at least as informative as a primal-state preparer.
 \item A genuinely dual-only certificate can fall outside a primal-output lower
 bound.  A claimed contradiction requires an explicit low-query reduction to
 the hard output contract.
\end{enumerate}
These rules instantiate rather than repeat the ledgers of
\cref{sec:preconditioning}.  In particular, all costs below retain operator
normalization, RHS amplification, residual conversion, tomography, and static
input construction.

\subsection{Projective lazy refresh and its exact boundary}
\label{sec:positive-lazy}

Consider a realized sequence of nonsingular real systems, indexed by
$t=0,\ldots,T$, with $f_t\ne0$:
\begin{equation}
 H_tz_t=f_t,\qquad q_t=\norm{z_t},\qquad \psi_t=z_t/q_t,
 \qquad M_t=\frac{q_t}{\norm{f_t}}H_t,
 \label{eq:positive-projective-definitions}
\end{equation}
and define
\begin{equation}
 \delta_t=\norm{M_{t+1}(\psi_{t+1}-\psi_t)},
 \qquad V_{\rm proj}=\sum_{t<T}\delta_t.
 \label{eq:positive-projective-variation}
\end{equation}
Here $z_t$ is the coordinate vector of the solution in whatever fixed basis
the module stores it in (for the OSS systems below, the fixed $V$ spanning
$\ker A$), so $\delta_t$ and $V_{\rm proj}$ depend on that basis.
Fix $\eta>0$ and $\epsilon\geq0$.  At the initial checkpoint $s=0$ and
every subsequent refresh, store a classical $\widehat z_s$ satisfying
$\norm{H_s\widehat z_s-f_s}/\norm{f_s}\leq\epsilon$.
At each later time $t$, let $s$ be the most recent checkpoint and choose
the scalar minimizing $\norm{H_t(a\widehat z_s)-f_t}$.
For $v=H_t\widehat z_s\ne0$ this scalar is
$\inner{v}{f_t}/\norm v^2$; for $v=0$ choose zero.
Reuse the vector if this minimum relative residual is at most $\eta$;
otherwise refresh at $t$.  The minimization and threshold comparison here
are exact.  Write $R$ for the total number of checkpoint solves, including
the initial solve.

\begin{theorem}[Pathwise residual-certified refresh]
\label{thm:positive-lazy-refresh}
Suppose $G>0$ and, for every test at time $t$ with most recent checkpoint
$s$, including a test that triggers a refresh,
\begin{equation}
 \norm{M_tM_j^{-1}}\leq G\qquad(s\leq j\leq t).       \label{eq:positive-cross-gain}
\end{equation}
If $\epsilon\leq\eta/(2G)$, then the number $R$ of checkpoint solves satisfies
\begin{equation}
 \boxed{R\leq1+\frac{2GV_{\rm proj}}{\eta}.}          \label{eq:positive-refresh-count}
\end{equation}
The statement is pathwise: the systems may depend adaptively on earlier
inexact steps.
\end{theorem}

\begin{proof}
Put $u_s=\widehat z_s/q_s$ and
$e_s=M_s(u_s-\psi_s)=(H_s\widehat z_s-f_s)/\norm{f_s}$, so
$\norm{e_s}\leq\epsilon$.  The nonzero right-hand sides imply $q_s>0$.
For a test at $t$, the least-squares scalar is no worse than $a=q_t/q_s$.
The candidate's normalized residual vector is
\[
 M_t(u_s-\psi_t)=M_tM_s^{-1}e_s-
 \sum_{j=s}^{t-1}M_tM_{j+1}^{-1}
       M_{j+1}(\psi_{j+1}-\psi_j).
\]
The triangle inequality and \eqref{eq:positive-cross-gain} give
\[
 \frac{\norm{H_t(a\widehat z_s-z_t)}}{\norm{f_t}}
 \leq G\left(\epsilon+\sum_{j=s}^{t-1}\delta_j\right).
\]
Consequently every failed test after a checkpoint consumes more than
$\eta/(2G)$ of projective variation.  The consumed intervals are disjoint,
and are contained in $\{0,\ldots,T-1\}$, which proves
\eqref{eq:positive-refresh-count}, including when no test fails.
\end{proof}

\paragraph{Formal verification and scope.}
The Lean development in \path{formal/QipmFormal/Refresh/} verifies the
normalization identities, scalar least-squares minimization, residual
transport, and the bound for the sequential reuse/refresh procedure.
It checks the initial solve, equality at the acceptance threshold, and
the final interval with no further refresh.  The claim map and reproduction
command are in \path{formal/REFRESH.md}.  The audit permits only
\texttt{propext}, \texttt{Classical.choice}, and \texttt{Quot.sound}.
This verifies the pathwise theorem under its stated gain and checkpoint
accuracy hypotheses; it does not establish those hypotheses for an IPM
trajectory, nor verify the tail, OSS, cycling, or quantum-cost results below.
A test with numerical error needs an explicit rejection guarantee as well
as safe acceptance: rejecting a truly acceptable vector can add checkpoint
solves without consuming the variation used by this proof.

For $H_t=AD_tA^T$, the test uses $A^T\widehat z_s$, the current diagonal,
and $A$; it costs $O(\nnz(A))$ classical arithmetic without forming $H_t$.
This per-iteration term remains in the ledger.  The theorem amortizes QLS and
readout, not all outer-loop work.  It is a residual-screened classical cache,
distinct from coherent eigenpath traversal \cite{positive-boixo2010-eigenpath}
and from classical Krylov recycling or inverse maintenance
\cite{parks2006-recycling-krylov-subspaces-sequences-linear-systems,
lee2015-efficient-inverse-maintenance-and-faster}.
Classical warm starts and reoptimization instead transfer information between
perturbed problem instances
\cite{positive-yildirim2002-warm-start,
r4-n062-reoptimization-primal-dual-interior-point-method}; they do not certify
reuse of a direction along one fixed central tail.

The projective parameters are finite on the exact central tail even when the
unscaled normal matrix becomes ill conditioned.

\begin{theorem}[Directed strict-complementarity tail]
\label{thm:positive-directed-tail}
Let $Q$ project onto $U$, let $Q^\perp=I-Q$, and suppose
\begin{equation}
 H_\mu=\mu^{-1}C_\mu+\mu D_\mu,\qquad
 C_\mu=QC_\mu Q,                                       \label{eq:positive-tail-split}
\end{equation}
where, uniformly on $0\leq\mu\leq\mu_0$,
\[
 c_-Q\preceq C_\mu\preceq c_+Q,\quad D_\mu\succeq0,\quad
 \norm{D_\mu}\leq d_+,\quad Q^\perp D_\mu Q^\perp\succeq d_-Q^\perp.
\]
Assume $C_\mu,D_\mu$ are $C^1$ with bounded derivatives.  For nonzero
$f\in U$, set $z_\mu=H_\mu^{-1}f$, $\psi_\mu=z_\mu/\norm{z_\mu}$, and
$N_\mu=(\norm{z_\mu}/\norm f)H_\mu$.  There are finite constants $B,L,G$
depending only on the displayed bounds and $\mu_0$ such that
\begin{equation}
 \sup_\mu\norm{N_\mu}\leq B,\quad
 \sup_{0<\mu\leq\nu\leq\mu_0}\norm{N_\mu N_\nu^{-1}}\leq G,\quad
 \sup_\mu\norm{\psi_\mu'}\leq L.                       \label{eq:positive-directed-bounds}
\end{equation}
Hence every decreasing discretization, using the refresh policy above with
$\epsilon\leq\eta/(2G)$, satisfies
\begin{equation}
 V_{\rm proj}\leq BL\mu_0,\qquad
 R\leq1+\frac{2GBL\mu_0}{\eta},                        \label{eq:positive-tail-refresh}
\end{equation}
independently of its endpoint and number of steps.
\end{theorem}

\begin{proof}
Put $\mathcal A_\mu=C_\mu+\mu^2D_\mu$ and
$w_\mu=\mathcal A_\mu^{-1}f$.
Then $z_\mu=\mu w_\mu$ and
$N_\mu=a_\mu\mathcal A_\mu$, where $a_\mu=\norm{w_\mu}/\norm f$.
In $U\oplus U^\perp$ blocks, Schur elimination gives
\[
 E_\mu=D_{UU}-D_{U\perp}D_{\perp\perp}^{-1}D_{\perp U}\succeq0,\quad
 S_\mu=C_{UU}+\mu^2E_\mu,
\]
\[
 w_U=S_\mu^{-1}f,\qquad
 w_\perp=-D_{\perp\perp}^{-1}D_{\perp U}w_U.
\]
The hypotheses bound these factors and their derivatives and bound
$\norm{w_\mu}$ above and away from zero.  Thus $a_\mu$ and the normalized ray
are $C^1$ at zero, proving the first and third bounds.

The only unbounded block of $\mathcal A_\nu^{-1}$ is
$\nu^{-2}D_{\perp\perp,\nu}^{-1}$ in its lower-right block.  Every block of $\mathcal A_\mu$
that can multiply it contains $\mu^2$.  For $\mu\leq\nu$ this ratio is at most
one, so $\norm{\mathcal A_\mu\mathcal A_\nu^{-1}}=O(1)$.  The bounded ratio
$a_\mu/a_\nu$ proves the directed cross-gain.  Finally the fundamental
theorem of calculus gives
\[
 \sum_t\norm{N_{\mu_{t+1}}(\psi_{\mu_{t+1}}-\psi_{\mu_t})}
 \leq B\int_0^{\mu_0}\norm{\psi_\mu'}\,d\mu\leq BL\mu_0,
\]
and \cref{thm:positive-lazy-refresh} gives the refresh count.
\end{proof}

The order in \eqref{eq:positive-directed-bounds} is essential.  For
$C=\diag(1,0)$, $D=\diag(0,1)$, and $f=e_1$, the ray and forward gain are
constant, whereas the reverse gain is $(\nu/\mu)^2$.  On an exact LP central
path, $C_\mu=A_B\diag(x_B^2)A_B^T$, $D_\mu=A_N\diag(s_N^{-2})A_N^T$, and
$f=b\in\operatorname{range}(A_B)$.  Boundary analyticity supplies the stated
regularity even for degenerate optimal faces
\cite{halicka1999-analytical-properties-central-path-linear-programming,
positive-monteiro1996-limiting-derivatives}.  Its constants are instance
dependent and need not be bounded by sparsity or dimension alone.

At checkpoint $s$, let $\Gamma_{f,s}$ be the RHS-state amplification,
$\mathcal K_s=\alpha_s\norm{H_s^{-1}}$ the effective QLS parameter, and
$\rho_s=\alpha_s\norm{z_s}/\norm{f_s}$ the normalized-state-to-residual
conversion.  Suppose these are bounded by $\Gamma_f,\mathcal K,\rho$ on the
tail.  A refresh must have residual accuracy $O(\eta/G)$, as required by
\cref{thm:positive-lazy-refresh}; hence its QLS precision contributes
$\log(G/\eta)$, and tight coherent tomography uses
$\Otilde(mG\rho/\eta)$ controlled state preparations.  The resulting ledger is
\begin{equation}
 \Otilde\!\left[
 \left(1+\frac{GBL\mu_0}{\eta}\right)
 \Gamma_f\mathcal K\frac{mG\rho}{\eta}
 \log\frac{G}{\eta}\right],                           \label{eq:positive-lazy-quantum-ledger}
\end{equation}
where RHS preparation and amplification are explicit rather than hidden in
the operator cost.  Copy tomography instead has quadratic inverse dependence
on the required $O(\eta/(G\rho))$ normalized-state accuracy
\cite{joran2023-quantum-tomography-state-preparation-unitaries}.  The ordinary
$T=O(\sqrt n\log(\mu_0/\epsilon_{\rm opt}))$ sparse updates and residual
tests remain.

\begin{corollary}[Exact-central feasible OSS]
\label{cor:positive-exact-oss}
With a fixed full-rank $V$ spanning $\ker A$, let
\begin{equation}
 \mathcal M_\mu=[-X_\mu A^T\ \ S_\mu V],\qquad
 f_\mu=-\tau\mu e.                                     \label{eq:positive-oss}
\end{equation}
On the exact central path, its normalized solution ray has finite tail
variation and bounded directed cross-gain.  Every coefficient vector reconstructs
$\Delta x=V\lambda$, $\Delta s=-A^T\Delta y$ and hence preserves both
feasibility equations exactly.  Consequently residual-certified reuse with
checkpoint accuracy $\epsilon\leq\eta/(2G)$ has a
final-precision-independent refresh count on the exact-central tail.
\end{corollary}

\begin{proof}
At the center,
\[
 \mathcal M_\mu^T\mathcal M_\mu
 =\diag(AX_\mu^2A^T,V^TS_\mu^2V),
\]
so $\norm{\mathcal M_\mu^{-1}}=O(1/\mu)$.  Analyticity gives
$\norm{\mathcal M_\mu-\mathcal M_\nu}=O(|\mu-\nu|)$, and therefore
$\norm{\mathcal M_\mu\mathcal M_\nu^{-1}}=O(1)$ for $\mu\leq\nu$.
Differentiating centrality shows that the OSS solution divided by $\mu$ is an
analytic nonzero vector.  Thus the normalized exact-central OSS ray has finite
variation.  The reconstruction identities prove exact feasibility, and
\cref{thm:positive-lazy-refresh} gives the refresh conclusion.
\end{proof}

This does not extend to every standard inexact trajectory.

\begin{proposition}[Cycling and proximal no-go]
\label{prop:positive-cycling-proximal}
There is a row-$3$-sparse LP and a feasible three-phase trajectory in a fixed
$N_2(\theta)$ neighborhood, satisfying the standard fixed relative OSS
residual contract, for which
\begin{equation}
 \sum_{t<T}\delta_t=\Omega\!\left(
 \log\frac{\mu_0}{\epsilon_{\rm opt}}\right).          \label{eq:positive-cycling}
\end{equation}
Moreover, no Tikhonov rule inferred only from that neighborhood and residual
contract can simultaneously guarantee finite projective variation, admissible
OSS residual, and bounded reconstructed steps.  A residual-safe ridge scale
does not improve worst-case conditioning.
\end{proposition}

\begin{proof}
Use
\[
 \min\sum_{j=4}^n x_j\quad\text{subject to}\quad
 x_1+x_2+x_3=1,\quad x\geq0.
\]
Set $\tau=a_0/\sqrt n$, $\mu_t=(1-\tau)^t\mu_0$, and use
\[
 u_0=\frac{1}{\sqrt2}(1,-1,0),\qquad
 u_1=\frac{1}{\sqrt2}(0,1,-1),\qquad
 u_2=\frac{1}{\sqrt2}(-1,0,1).
\]
Cycle
\[
 x_{1:3}^{(t)}=e/3+d u_{t\bmod3},\quad x_{4:n}^{(t)}=\mu_te,
 \quad y^{(t)}=-3\mu_t,
 \quad s_{1:3}^{(t)}=3\mu_te,\quad s_{4:n}^{(t)}=e.
\]
Then $\norm{X_ts_t-\mu_te}=3d\mu_t$, and the feasible step to the next
phase has OSS residual $3d\mu_t(u_{j+1}-\tau u_j)$.  Taking
$d=\Theta(1/\sqrt n)$ meets the usual relative contract
$\eta_n=\Theta(1/\sqrt n)$.
With the fixed sparse basis
$V=[e_1-e_3,e_2-e_3,e_4,\ldots,e_n]$, the exact OSS coefficient rays converge
cyclically to the three distinct lines $(1,-1)$, $(0,1)$, and $(-1,0)$.
The normalized destination operator acts on their differences at constant
scale, so $\delta_t=\Omega(d)$ eventually.  Since
$T=\Theta(\sqrt n\log(\mu_0/\epsilon_{\rm opt}))$, this proves
\eqref{eq:positive-cycling}.

The preceding cycling half uses both $d=\Theta(1/\sqrt n)$ and
$\eta_n=\Theta(1/\sqrt n)$.  The proximal no-go switches both regimes: fix
constants $d>0$ and $0<\eta<\theta<1$ such that
\begin{equation}
 3d\leq\theta,\qquad
 \frac{c_1d}{\sqrt{a_0^2+9d^2}}>\eta,                  \label{eq:positive-separated-rays}
\end{equation}
where $c_1>0$ is the constant separation of the three projective lines.
For the ridge selector
\[
 z_\lambda=\arg\min_z\{\norm{Mz-f}^2+\lambda\norm{z-c}^2\},
\]
let $z_*=M^{-1}f$, $v$ be a right singular vector of singular value $\sigma$,
and $e=c-z_*$.  Direct diagonalization gives the exact frontier
\begin{align*}
 \langle v,z_\lambda-c\rangle
 &=-\frac{\sigma^2}{\sigma^2+\lambda}\langle v,e\rangle,\\
 \left|\left\langle \frac{Mv}{\sigma},Mz_\lambda-f\right\rangle\right|
 &=\frac{\lambda\sigma}{\sigma^2+\lambda}|\langle v,e\rangle|.
\end{align*}
On the cycling face $\sigma=\Theta(\mu)$.  A ray-locked sequence of bounded
steps and finite projective variation is Cauchy, so on at least one of every
three phases it stays a fixed distance from that phase's exact face target.
If the ridge moves a nonvanishing fraction toward those targets infinitely
often, its projective variation diverges.  Otherwise its residual is at least
$c_1d\mu_t$ on an infinite separated subsequence.  Since
$\norm{f_t}=\mu_t\sqrt{a_0^2+9d^2}$,
\eqref{eq:positive-separated-rays} makes this exceed the admissible
$\eta\norm{f_t}$.  This is the fixed-constant regime, not the
$d=\Theta(1/\sqrt n)$ regime used above.

Conversely the universal
bound
\begin{equation}
 \norm{Mz_\lambda-f}\leq \frac{\sqrt\lambda}{2}\norm e          \label{eq:positive-prox-residual}
\end{equation}
is sharp, since $\sigma\lambda/(\sigma^2+\lambda)$ is maximized at
$\sigma=\sqrt\lambda$.  With $\norm f=\Theta(\mu)$ and $\norm e=\Theta(1)$,
the uniform worst-case residual-safe scale is
$\lambda=O(\eta^2\mu^2)$.  The augmented system
$[M^T\ \sqrt\lambda I]^T$ then retains condition $\Theta(1/\mu)$; its normal
matrix retains $\Theta(1/\mu^2)$.  Thus regularization cannot establish the
missing trajectory property or improve the safe worst-case scale.
\end{proof}

Residual-certified reuse nevertheless gives an honest conditional QIPM.  If
the feasible short-step theorem accepts every coefficient vector satisfying
$\norm{\mathcal M_t\widehat z_t-f_t}\leq\eta\norm{f_t}$, then exact OSS
reconstruction preserves feasibility and the published neighborhood, gap, and
iteration guarantees apply unchanged
\cite{mohammadisiahroudi2025-inexact-feasible-interior-point-method}.  Under
finite realized $V_{\rm proj}$, the directed cross-gain bound, and checkpoint
accuracy $\epsilon\leq\eta/(2G)$, \cref{thm:positive-lazy-refresh} also applies.
Assume block encodings of $A^T/\alpha_A$ and $V/\alpha_V$, coherent value
access to the current $x_t,s_t$, and known bounds
$R_x\geq\norm{x_t}_\infty$ and $R_s\geq\norm{s_t}_\infty$.  If
$F_{\infty,t}\geq\norm{f_t}_\infty$, define
\[
 \alpha_{\mathcal M,t}=R_x\alpha_A+R_s\alpha_V,\quad
 \mathcal K_t=\alpha_{\mathcal M,t}\norm{\mathcal M_t^{-1}},\quad
 \Gamma_{f,t}=\frac{\sqrt n F_{\infty,t}}{\norm{f_t}},
\]
and the conversion factor
$\rho_t=\alpha_{\mathcal M,t}\norm{z_t}/\norm{f_t}$, its dense-output ledger is
\begin{equation}
 C_{\rm iter}+O(T\nnz\mathcal M)+
 \sum_{s\in\mathcal R}\Otilde\!\left(
  \Gamma_{f,s}\mathcal K_s\frac{n\rho_s}{\epsilon}\right),     \label{eq:positive-certified-ledger}
\end{equation}
plus $O(nR)$ checkpoint writes.  Here $C_{\rm iter}$ is the full cost of
maintaining the promised coherent iterate oracles.  In particular, rebuilding
dense $x_t,s_t$ memories each round gives
$C_{\rm iter}=\Theta(nTB_{\rm bit})$; the ledger omits that term only when the
outer implementation already supplies coherent reads without such rebuilds.
A quantum residual test costs
$\Otilde((\Lambda_{t,s}+\Gamma_{f,t})/\eta)$, where
$\Lambda_{t,s}=\alpha_{\mathcal M,t}\norm{\widehat z_s}/
\norm{\mathcal M_t\widehat z_s}$; on the cycling witness
$\Lambda_{t,s}=\Omega(1/\mu_t)$.  Generic verification is therefore not free
\cite{rolando2021-complexity-quantum-state-verification-quantum}.  Proximal
QLS warm starts remain useful for a certified instance, but do not change this
sequence-level conclusion \cite{kim2026-catalyst-framework-problem-proximal-point}.
Iterative-refinement tomography can change how one solves a single high-accuracy
system, but it does not remove the per-refresh/output factors in
\eqref{eq:positive-certified-ledger}
\cite{mohammadisiahroudi2023-exponentially-more-precise-tomography-for}.

\subsection{Partition-free equilibration and exact face repair}
\label{sec:positive-face-repair}

The stable solver in \cref{sec:stable-active-equilibration} assumes a supplied
active range.  On a strict-complementarity tail it can instead be constructed
from current coordinates.  Suppose
\begin{equation}
 \underline x\leq x_i\leq\overline x\ (i\in B),\qquad
 \underline s\leq s_i\leq\overline s\ (i\in N),        \label{eq:positive-tail-margins}
\end{equation}
and $(1-\theta)\mu\leq x_is_i\leq(1+\theta)\mu$.

\begin{lemma}[Constant-margin tail selector]
\label{lem:positive-tail-selector}
If
\begin{equation}
 \mu\leq \frac{1}{2(1+\theta)}
 \min\{\underline x^2,\underline s^2\},                 \label{eq:positive-selector-threshold}
\end{equation}
then $E=\diag({\bf1}\{x_i\geq s_i\})$ equals $P_B$, and every comparison has
absolute margin at least $\frac12\min\{\underline x,\underline s\}$.
\end{lemma}

\begin{proof}
For $i\in B$, $s_i\leq(1+\theta)\mu/\underline x\leq
\underline x/2\leq x_i/2$.  The symmetric calculation gives
$x_i\leq s_i/2$ for $i\in N$.
\end{proof}

This is the classical tail rule of Mehrotra--Ye, now used coherently
\cite{positive-mehrotra1993-optimal-face,positive-cartis2016-active-set}.  It
requires current value oracles and enough bits to resolve the promised margin;
it does not convert amplitude states into coordinate values.

\begin{theorem}[Constructive stable active-range equilibration]
\label{thm:positive-partition-free-equilibration}
Under \cref{lem:positive-tail-selector}, define $F=I-E$,
\begin{align}
 C_\mu&=\mu AEXS^{-1}EA^T,&
 D_\mu&=\mu^{-1}AFXS^{-1}FA^T,\nonumber\\
 Q_U&=\Pi_{\operatorname{range}(AE)},&
 \widehat M_\mu&=C_\mu+D_\mu-(1-\mu^2)Q_UD_\mu.        \label{eq:positive-constructive-split}
\end{align}
Then $H_\mu=AXS^{-1}A^T=\mu^{-1}C_\mu+\mu D_\mu$, and direct encodings have
\begin{equation}
 \alpha_C\leq\alpha_A^2\frac{\overline x^2}{1-\theta},
 \qquad
 \alpha_D\leq\alpha_A^2\frac{1+\theta}{\underline s^2}. \label{eq:positive-split-normalization}
\end{equation}
If $\sigma_B=\sigma_{\min}^+(A_B)>0$, QSVT constructs $Q_U$ to error $\delta$
in
\begin{equation}
 O\!\left(\frac{\alpha_A}{\sigma_B}\log\frac{1}{\delta}\right) \label{eq:positive-active-filter-cost}
\end{equation}
queries.  If the fixed tail block bounds of
\cref{thm:stable-split-equilibration} hold, $\widehat M_\mu$ has
$\mu$-independent normalization and inverse norm.  For $f\in\operatorname{range}(AE)$,
\begin{equation}
 \widehat M_\mu^{-1}f=\mu^{-1}H_\mu^{-1}f                 \label{eq:positive-state-identity}
\end{equation}
for the exact projector, and constant projector error gives constant
solution-state error.  With RHS factor $\Gamma_f$, state preparation costs
\begin{equation}
 \Otilde\!\left(
 \Gamma_f\Gamma_M\norm{\widehat M_\mu^{-1}}
 \frac{\alpha_A}{\sigma_B}\right),
 \quad \Gamma_M=\alpha_C+(2-\mu^2)\alpha_D,              \label{eq:positive-equilibration-ledger}
\end{equation}
with no polynomial $1/\mu$ factor.
\end{theorem}

\begin{proof}
The selector splits the columns exactly.  On $B$,
$\mu x_i/s_i=\mu x_i^2/(x_is_i)\leq\overline x^2/(1-\theta)$;
on $N$, $\mu^{-1}x_i/s_i=(x_is_i)/(\mu s_i^2)\leq
(1+\theta)/\underline s^2$.  This proves
\eqref{eq:positive-split-normalization}.  QSVT thresholding of $AE/\alpha_A$
at its smallest nonzero singular value gives
\eqref{eq:positive-active-filter-cost}
\cite{andras2019-quantum-singular-value-transformation-beyond}.
With $P_\mu=\mu Q_U+\mu^{-1}(I-Q_U)$, multiplication gives
$\widehat M_\mu=P_\mu H_\mu$.  Its $U\oplus U^\perp$ diagonal blocks are
uniformly invertible and the upper off-diagonal block is $O(\mu^2)$; the Schur
complement bounds its inverse.  Since $P_\mu f=\mu f$, the solution identity
follows.  Direct LCU gives $\Gamma_M$, and a Neumann bound makes constant
projector error sufficient.  Multiplying projector, QLS, RHS, and direct
encoding factors yields \eqref{eq:positive-equilibration-ledger}.
\end{proof}

Thus \cref{thm:positive-partition-free-equilibration} removes exactly the
supplied-partition assumption left open in \cref{sec:stable-active-equilibration}.
It is narrower than the published partition-estimated infeasible-QIPM
preconditioner \cite{zeguan2026-preconditioned-inexact-infeasible-quantum-interi}:
it gives constant tail conditioning for active-range RHSs by direct assembly,
not a full infeasible-QIPM theorem.  The input-dependent gaps and output remain
charged, consistent with classical active-partition preconditioning
\cite{positive-chai2007-preconditioning}.

The dangerous face components of an OSS residual can also be isolated exactly.
Let $Z$ span $\ker A_B$ and $Y$ span $\ker A_B^T$ orthonormally, and define
\begin{align}
 G_p&=Z^TX_B^{-1}S_BZ,&g_p(r)&=Z^TX_B^{-1}r_B,\nonumber\\
 G_d&=Y^TA_N S_N^{-1}X_NA_N^TY,&
 g_d(r)&=Y^TA_NS_N^{-1}r_N.                              \label{eq:positive-leakage-data}
\end{align}
Write
$\sigma_N=\sigma_{\min}(A_N^T|_{\ker A_B^T})$ when
$\ker A_B^T\ne\{0\}$.

\begin{theorem}[Face leakage, repair, and restricted stability]
\label{thm:positive-face-repair}
For any feasible OSS direction error with residual $r$, the primal and dual
face errors are
\begin{equation}
 p_F=ZG_p^{-1}g_p(r),\qquad q_F=-YG_d^{-1}g_d(r).         \label{eq:positive-face-errors}
\end{equation}
Under \eqref{eq:positive-tail-margins}, the primal map has sharp two-sided
$\Theta(1/\mu)$ amplification when $\ker A_B\ne\{0\}$, and the dual map has
the same amplification when $\ker A_B^T\ne\{0\}$; otherwise the corresponding
map is zero.  The correction
\begin{equation}
 \widehat{\Delta x}^{+}_B=\widehat{\Delta x}_B-p_F,\quad
 \widehat{\Delta y}^{+}=\widehat{\Delta y}+YG_d^{-1}g_d(r),\quad
 \widehat{\Delta s}^{+}=\widehat{\Delta s}-A^TYG_d^{-1}g_d(r)  \label{eq:positive-face-correction}
\end{equation}
preserves primal and dual feasibility and makes $g_p(r^+)=g_d(r^+)=0$ exactly.
It also obeys $\norm{r^+}\leq C_{\rm rep}\norm r$, and the inverse of the OSS
map restricted by these two equations is uniformly bounded as $\mu\downarrow0$.
The sharp worst-case amortized condition is
\begin{equation}
 \sum_t\left(\frac{\norm{g_p(r_t)}}{\mu_t}
             +\frac{\norm{g_d(r_t)}}{\mu_t}\right)<\infty.    \label{eq:positive-leakage-ledger}
\end{equation}
\end{theorem}

\begin{proof}
Writing the direction error as $(e_x,e_y,e_s)$ gives
$Ae_x=0$, $e_s=-A^Te_y$, and $Se_x-XA^Te_y=r$.  Multiply the $B$ rows by
$Z^TX_B^{-1}$; the dual term vanishes and
$Z^TX_B^{-1}S_B(e_x)_B=g_p(r)$.  This is the normal equation for the
$X_B^{-1}S_B$-weighted projection of $(e_x)_B$ onto $\ker A_B$.  Eliminating
$e_x$ gives $AXS^{-1}A^Te_y=-AS^{-1}r$; left-multiplying by $Y^T$ and using
$Y^TA_B=0$ gives
$Y^T(A_NX_NS_N^{-1}A_N^T)e_y=-g_d(r)$, the normal equation for the
$H_\mu$-weighted projection of $e_y$ onto $\ker A_B^T$, whose coordinate is
$-G_d^{-1}g_d(r)$ because $Y^TH_\mu Y=G_d$.  These are
\eqref{eq:positive-face-errors}.

Complementarity and the tail margins imply
\[
 c_p\mu I\preceq G_p\preceq C_p\mu I,\qquad
 c_d\mu I\preceq G_d\preceq C_d\mu I,
\]
where, when $\ker A_B^T\ne\{0\}$, the dual constants also contain
$\sigma_N^2>0$.  This proves both sides of the $1/\mu$ amplification on each
nonzero face.  The primal correction lies in $\ker A$ and is
supported on $B$; the dual correction lies in $\ker A_B^T$.  They therefore
preserve feasibility and affect disjoint projected residual components, giving
$g_p(r^+)=g_d(r^+)=0$.  The displayed spectral bounds and the bounded tail
coordinates give $\norm{r^+}\leq C_{\rm rep}\norm r$.

At $\mu=0$, the kernel of the limiting OSS map consists precisely of primal
motions in $\ker A_B$ and dual motions in $\ker A_B^T$.  The two corrected
equations remove those components.  Compactness of the unit sphere on the
remaining complement and continuity give a uniform smallest restricted
singular value.  Finally each step can realize face displacement comparable
to $\norm{g_p}/\mu+\norm{g_d}/\mu$; aligned successive residuals attain that
order.  Summability is therefore sufficient and worst-case sharp.
\end{proof}

The object controlled by \eqref{eq:positive-leakage-ledger} is the
residual-induced total variation inside the primal and dual optimal faces.
It is not by itself a bound on the realized projective variation.  A local
strict-complementarity chart also needs the transverse-forcing ledger
\begin{equation}
 W_r=\sum_t\left\|
 \frac{r_{t+1}}{\mu_{t+1}}-\frac{r_t}{\mu_t}\right\|<\infty.             \label{eq:positive-transverse-ledger}
\end{equation}
Define the total face leakage
\begin{equation}
 \mathcal L_\infty=
 \sum_t\bigl(\norm{p_F(r_t)}+\norm{q_F(r_t)}\bigr).       \label{eq:positive-total-face-leakage}
\end{equation}
These ledgers bound the residual-induced face motion and the transverse
forcing.  They do not bound realized $V_{\rm proj}$, even together with the
tail bound
$\sum_{k\geq t}(\norm{p_F(r_k)}+\norm{q_F(r_k)})=O(\mu_t)$: the
row-$3$-sparse family of \cref{prop:positive-cycling-proximal} with face
displacement $d_t=\varrho_{t\bmod3}\mu_t$ and $\varrho=(1,3,1)$ has convergent
$\mathcal L_\infty$ and $W_r$, satisfies that tail bound, and stays inside
the neighborhood and residual contract, yet its realized $V_{\rm proj}$
grows linearly in the number of steps.  An oscillating relative face
displacement $d_t/\mu_t$ moves the normalized ray by $\Theta(1)$ per step,
while both ledgers are absolute sums that do not see it; with the constant
ratio $\varrho=(1,1,1)$ the same family has bounded $V_{\rm proj}$ when
$\delta_t$ is measured in an orthonormal basis of $\ker A$, but not in the
skew stored basis $V=[e_1-e_3,e_2-e_3,e_4,\ldots,e_n]$ of
\cref{prop:positive-cycling-proximal}: there the three cycling directions
have coefficient norms $1,1/\sqrt2,1/\sqrt2$, so $\norm{z_t}$ is phase
dependent, $\delta_t=\Theta(1)$ on two of every three steps, and
$V_{\rm proj}$ still grows linearly.  The constant-ratio family therefore
does not certify finite $V_{\rm proj}$ for
\cref{thm:positive-lazy-refresh} in the basis the module stores.  Neither
$o(\mu_t)$ leakage nor convergence of $d_t/\mu_t$ suffices either:
$d_t=\varrho_{t\bmod3}\mu_t/\sqrt t$ gives convergent ledgers and an
$o(\mu_t)$ tail while $V_{\rm proj}$ diverges logarithmically.  Finite
realized $V_{\rm proj}$ therefore remains a separate trajectory hypothesis,
as in \cref{thm:positive-lazy-refresh}; it is the hypothesis under which that
theorem bounds the refresh count, and its failure is not shown here to
preclude reuse.  If
repaired directions are used, $r_t$ here is the corrected and retested
residual.

The repair is basis free.  With $T=E\sqrt{XS^{-1}}E$,
$B_r=E(XS)^{-1/2}E$, $K_F=AT$, $Q=\Pi_{\ker(AE)^T}$,
$D_N=FXS^{-1}F$, and $H_N=AD_NA^T$,
\begin{equation}
 p_F(r)=T\Pi_{\ker K_F}B_rr,\qquad
 q_F(r)=-(QH_NQ)^+QAFS^{-1}Fr.                          \label{eq:positive-basis-free-repair}
\end{equation}
Indeed the weighted kernel projector equals
$W^{1/2}Z(Z^TWZ)^{-1}Z^TW^{1/2}$ on $B$, with $W=X_B^{-1}S_B$;
sandwiching it yields the first expression.  The second is the inverse of the
$H_N$ form on $\ker A_B^T$.  Given known positive lower bounds on the relevant
gaps $\sigma_B$ and $\sigma_N$, QSVT implements the projectors with
\begin{equation}
 \chi_p=O\!\left(\frac{\alpha_A}{\sigma_B}
 \frac{\overline x}{\underline x}\sqrt{\frac{1+\theta}{1-\theta}}\right),
 \quad
 \chi_d=O\!\left(\frac{\alpha_A^2}{\sigma_N^2}
 \frac{1+\theta}{1-\theta}\frac{\overline s^2}{\underline s^2}\right), \label{eq:positive-face-condition-ledger}
\end{equation}
both independent of $\mu$.  With
\[
 R_T=\frac{\overline x}{\sqrt{(1-\theta)\mu}},\qquad
 R_B=\frac{1}{\sqrt{(1-\theta)\mu}},\qquad
 \alpha_p=R_TR_B=\Theta(1/\mu),
\]
the primal correction map has RHS factor
$\Gamma_p=\alpha_p\norm{Er}/\norm{p_F(r)}$.  If
$c_N(r)=QAFS^{-1}Fr$, the dual factor is
\[
 \Gamma_d(r)=\frac{\alpha_A\underline s^{-1}\norm{Fr}}
                   {\norm{c_N(r)}}.
\]
Writing $\Gamma_{r,B},\Gamma_{r,N}$ for selected-residual preparation, the
conservative normalized correction-state ledgers are
\begin{equation}
 Q_p=\Otilde(\Gamma_{r,B}\Gamma_p\chi_p),\qquad
 Q_d=\Otilde(\Gamma_{r,N}\Gamma_d\chi_p\chi_d),          \label{eq:positive-repair-query-ledger}
\end{equation}
with norm recovery and the RHS-specific inverse success amplitude charged
separately.  These are operator modules, not free corrections.

Replacing the fixed face radius in the row-$3$-sparse example by $d_t>0$
gives exactly
\[
 r_t=3\mu_t(d_{t+1}u_{j+1}-\tau d_tu_j),
\]
and, for adjacent radii within a fixed factor, its leakage is
$\Theta(d_t+d_{t+1})$.  This verified identity establishes the leakage ledger
as a necessary worst-case condition.  It does not establish the projective
threshold claimed in a preliminary version of this analysis: exact solves show that when
$d_t\gg\mu_t$, the phase-dependent transverse coefficient changes by
$\Theta(\mu_t/d_t)$, while the normalized OSS operator amplifies that component
by $\Theta(d_t/\mu_t)$, so $\delta_t=\Theta(1)$ rather than
$\Theta(d_t+d_{t+1})$.

For example, on a geometric schedule with
$d_t=\Theta(\mu_t^\gamma)$ and $0<\gamma<1$, both
$\mathcal L_\infty$ and $W_r$ converge, but realized $V_{\rm proj}$ grows
linearly with the number of steps.  This family also violates the tail
requirement
$\sum_{k\geq t}(\norm{p_F(r_k)}+\norm{q_F(r_k)})=O(\mu_t)$; the
oscillating-displacement family above shows that even that tail bound
together with both convergent ledgers does not bound projective variation.
Thus the sharpness-family
claim from that preliminary version is corrected here.  A blanket $O(\mu^2)$ residual rule
is sufficient but is not the sharp leakage boundary.

There is a sharp internal dichotomy.  If the selector is omitted and one uses
$T_{\rm all}=(X^{-1}S)^{-1/2}$ and $K_{\rm all}=AT_{\rm all}$, the weighted
kernel formula reconstructs the entire Newton error, not merely face leakage.
For
\[
 A=\begin{pmatrix}1&1&0&0\\0&0&1&-1\end{pmatrix},\quad
 x=(1/2,1/2,\mu,\mu),\quad s=(2\mu,2\mu,1,1),
\]
$\kappa(K_{\rm all})=1/(2\mu)$ and
$\kappa(AXS^{-1}A^T)=1/(4\mu^2)$.  A Tikhonov kernel filter
$\lambda(K_{\rm all}^TK_{\rm all}+\lambda I)^{-1}$ accurate to $\delta<1/2$
forces $\lambda\leq2\delta\mu/(1-\delta)$ and hence condition
$\Omega(1/(\delta\mu^2))$.  Cluster resolution, by the coordinate selector or
an equivalent spectral cutoff, is necessary.

Finally, if a materialized repair must meet the pointwise budget
$O(\tau\mu_t)$, an arbitrary constant-norm correction in dimension $d$ needs
absolute tomography error $O(\tau\mu_t)$.  Generic coherent tomography costs
\begin{equation}
 \Thetatilde\!\left(\frac{d}{\tau\mu_t}\right)          \label{eq:positive-face-tomography}
\end{equation}
state-unitary calls.  This is a boundary for that generic extraction route,
not an LP-specific output lower bound.  A concise formula, sparse output, or
the compressed representation below can evade it.

\subsection{Implicit dual coordinates and compressed output}
\label{sec:positive-compressed-dual}

The state-to-diagonal obstruction of
\cref{sec:preconditioning-state-interface} is avoided when the changing
diagonal is computed from a smaller explicit iterate instead of from an
amplitude state.  Let the columns of $A$ have at most $d_c$ nonzeros and store
only $y_t\in\R^m$.  A static block encoding of $A/\alpha_A$ costs $C_A$ per
use, a coherent read of one retained-iterate entry costs $C_{y,{\rm read}}$,
and $B_{\rm bit}$ is the common arithmetic word width.  The one-time cost
$C_{\rm build}(A)$ is separate.  Coherent RAM may make
$C_{y,{\rm read}}$ polylogarithmic; otherwise the actual multiplexed-read
cost must be substituted below.

\begin{theorem}[Implicit slack oracle and affine closure]
\label{thm:positive-implicit-slack}
Given static coherent sparse-column access, coherent value access to $c_i$,
and coherent reads of $y_t$, a
reversible circuit computes
\begin{equation}
 s_i(y_t)=c_i-a_i^Ty_t                                  \label{eq:positive-slack-oracle}
\end{equation}
to absolute error $\epsilon_s$ using $d_c$ column accesses and
\[
 \Otilde\!\left(d_cB_{\rm bit}+d_cC_{y,{\rm read}}
                  +\log(1/\epsilon_s)\right)
\]
gates.  If
$s_i\geq\underline s_t>0$, it gives a unit-normalized encoding of
$\underline s_tS_t^{-1}$ with operator error
\begin{equation}
 O\!\left(\frac{\epsilon_s}{\underline s_t}+\epsilon_{\rm arith}\right), \label{eq:positive-slack-operator-error}
\end{equation}
and therefore an encoding of $AS_t^{-2}A^T$ with normalization
$\alpha_A^2/\underline s_t^2$.  It suffices to make the error in
\eqref{eq:positive-slack-operator-error} at most
\[
 O\!\left(\frac{\eta}{J_t\rho_t\mathcal K_{M,t}}\right),
\]
with $\rho_t$ as in \eqref{eq:positive-forcing-factor}; the factor
$1/\mathcal K_{M,t}$ converts operator error into relative vector error, and
$1/(J_t\rho_t)$ is the vector accuracy target of
\cref{thm:positive-compressed-qipm}.  On active coordinates this follows,
for example, from absolute slack error
$O(\eta\mu_t/(J_t\rho_t\mathcal K_{M,t}))$.  Thus finite precision is charged
against the transformed residual budget.  Updating
$y_{t+1}=y_t+\lambda_t\Delta y_t$ updates every implicit slack exactly with no
history factor and no $n$-coordinate rewrite.
\end{theorem}

\begin{proof}
Query $c_i$ and the support and values of $a_i$, then read the matching entries
of $y_t$ and reversibly accumulate \eqref{eq:positive-slack-oracle}.  Approximate
$\underline s_t/z$ on the promised interval and use it as a controlled-rotation
amplitude; uncomputation yields the diagonal encoding.  Product encodings give
the normal matrix.  The affine identity
$c_i-a_i^Ty_{t+1}=s_i(y_t)-\lambda_ta_i^T\Delta y_t$ proves exact closure.
No state-preparation oracle for $s$ is queried, so the state-to-diagonal
premise is absent.
\end{proof}

The dual method does not store a primal iterate.  Its selector instead uses
the implicit exact-barrier coordinate
\begin{equation}
 x_i^b=\frac{\mu_t}{s_i(y_t)},\qquad
 E_t=\diag\!\left({\bf1}\{x_i^b\geq s_i(y_t)\}\right)
    =\diag\!\left({\bf1}\{s_i(y_t)^2\leq\mu_t\}\right). \label{eq:positive-dual-selector}
\end{equation}
This is \cref{lem:positive-tail-selector} with exact complementarity
($\theta=0$), applied to $(x^b,s(y_t))$; the tail-margin hypotheses in this
application are explicitly assumed for this implicit pair.  Below
$\mu_t\leq\frac12\min\{\underline x^2,\underline s^2\}$, an active coordinate
has $x_i^b\geq2s_i$ and an inactive coordinate has $s_i\geq2x_i^b$.
Consequently the $x^b$--$s$ comparison has the constant margin of that lemma;
equivalently, the directly queryable squared test is at least $\mu_t/2$ from
its threshold.  The squared test is resolved with $O(\log(1/\mu_t))$ bits and
does not require a stored $x_t$.

The partition-free stable dual equation is
\begin{equation}
 G_t\Delta y_t=g_t,\qquad
 G_t=\mu_tAS_t^{-2}A^T,\qquad g_t=b-\mu_tAS_t^{-1}e.      \label{eq:positive-dual-newton}
\end{equation}
Using \eqref{eq:positive-dual-selector}, define
$C_t=\mu_t^2AES_t^{-2}EA^T$, $D_t=AFS_t^{-2}FA^T$,
$Q_t=\Pi_{\operatorname{range}(AE)}$, and
\begin{equation}
 M_t=C_t+D_t-(1-\mu_t^2)Q_tD_t=P_tG_t,\qquad
 P_t=\mu_tQ_t+\mu_t^{-1}(I-Q_t).                        \label{eq:positive-dual-stable}
\end{equation}
With direct component normalizations $\alpha_{C,t}$ and $\alpha_{D,t}$,
\begin{equation}
 \alpha_{M,t}\leq\alpha_{C,t}+(2-\mu_t^2)\alpha_{D,t},\qquad
 \mathcal K_{M,t}=\alpha_{M,t}\norm{M_t^{-1}}.           \label{eq:positive-dual-normalization}
\end{equation}
The assembled block encoding of $M_t$ is an
$(\alpha_{M,t},\cdot,\epsilon_{M,t})$ encoding in the sense of
\cref{def:block-encoding}; its precision is fixed so that
$\epsilon_{M,t}\le\alpha_{M,t}$, a convention used below.
Both component normalizations and $\mathcal K_{M,t}$
are independent of $\mu_t$ under the explicit tail gaps.  Since
$b\in\operatorname{range}(A_B)$, the transformed RHS can be assembled without
encoding the badly normalized $P_t$:
\begin{align}
 \widetilde g_t=P_tg_t
 &=\mu_tb-\mu_t^2AES_t^{-1}Ee-\mu_t^2Q_tAFS_t^{-1}Fe\nonumber\\
 &\quad-(I-Q_t)AFS_t^{-1}Fe.                            \label{eq:positive-transformed-rhs}
\end{align}
Let $\Gamma_{g,t}$ be its actual LCU numerator divided by
$\norm{\widetilde g_t}$, so cancellation is charged.
Thus the transformed system actually solved is
\begin{equation}
 M_t\Delta y_t=\widetilde g_t.                            \label{eq:positive-transformed-system}
\end{equation}

\begin{theorem}[Conditional compressed dual QIPM]
\label{thm:positive-compressed-qipm}
Assume the implicit-access model of \cref{thm:positive-implicit-slack}, the
strict-complementarity margins and known gaps used above, and the inexact
dual-barrier theorem of
\cite{zeguan2026-quantum-dual-logarithmic-barrier-method} (their
Theorem~3.3), which with $B_t=S_t^{-1}A^\top$ and
$\delta^{\rm N}_t=\norm{B_t\Delta y_t}\le1/2$ accepts any step $\widehat z_t$
satisfying $\norm{B_t(\widehat z_t-\Delta y_t)}\le\eta\delta^{\rm N}_t$ for a fixed
$\eta\le0.1$, in
$T=O(\sqrt n\log(n\mu_0/\epsilon_{\rm opt}))$ iterations.  Assume throughout
the tail that $0<\mu_t\leq1$ and $\operatorname{range}Q_t\ne\{0\}$.  For
$g_t\ne0$ set
\begin{equation}
 J_t=\frac{\norm{P_t^{-1}}\norm{P_tg_t}}{\norm{g_t}},\qquad
 \rho_t=\frac{\alpha_{M,t}\norm{\Delta y_t}}{\norm{\widetilde g_t}}.     \label{eq:positive-forcing-factor}
\end{equation}
Solving \eqref{eq:positive-transformed-system} and coherently reconstructing
its $m$-vector direction $\widehat z_t$ with
\[
 \norm{\widehat z_t-\Delta y_t}
 \le\frac{\eta}{2J_t\rho_t}\norm{\Delta y_t}
\]
produces
\[
 (y_T,\mu_T,O_{s,T},O_{x,T}),\qquad
 x_i=\frac{\mu_T}{s_i}+\frac{\mu_T}{s_i^2}a_i^T\Delta y_T,
\]
with the positivity, neighborhood, and convergence guarantees of the outer
method.  The $T$ outer updates use solves $t<T$; one additional terminal solve
at $(y_T,\mu_T)$ supplies the displayed $\Delta y_T$ and final coordinate
oracle.  Its query count is
\begin{equation}
 \boxed{Q_{\rm path}=\Otilde\!\left(
 \sum_{t\leq T}\frac{mJ_t\rho_t}{\eta}
 \Gamma_{g,t}\mathcal K_{M,t}\chi_{Q,t}\right),}       \label{eq:positive-compressed-query}
\end{equation}
where $\chi_{Q,t}=O((\alpha_A/\sigma_B)\log(1/\delta_Q))$.  Its data and gate
ledger is
\begin{equation}
 C_{\rm build}(A)+\Otilde\sum_{t\leq T}\left[
 \frac{mJ_t\rho_t}{\eta}\Gamma_{g,t}\mathcal K_{M,t}\chi_{Q,t}
 \bigl(C_A+d_c(B_{\rm bit}+C_{y,{\rm read}})\bigr)
 +mB_{\rm bit}\right].                                  \label{eq:positive-compressed-gates}
\end{equation}
There is no $nT$ iterate-write term.  If every displayed factor and $\eta^{-1}$
is $O(1)$, the favorable tail regime is $Q_{\rm path}=\Otilde(Tm)$.
\end{theorem}

\begin{proof}
Let $e_t=\widehat z_t-\Delta y_t$ with
$\norm{e_t}\le(\eta/(2J_t\rho_t))\norm{\Delta y_t}$.  Since
$G_t=P_t^{-1}M_t$ and
$\norm{M_t}\le\alpha_{M,t}+\epsilon_{M,t}\le2\alpha_{M,t}$ by the
block-encoding contract and the precision convention
$\epsilon_{M,t}\le\alpha_{M,t}$ fixed after
\eqref{eq:positive-dual-normalization},
$\norm{G_t}\le2\norm{P_t^{-1}}\alpha_{M,t}$, while
\eqref{eq:positive-forcing-factor} gives
$J_t\rho_t=\alpha_{M,t}\norm{P_t^{-1}}\norm{\Delta y_t}/\norm{g_t}$.
Using $g_t=G_t\Delta y_t$ and $G_t^2\preceq\norm{G_t}G_t$,
\[
 J_t\rho_t
 \ge\frac{\norm{G_t}\norm{\Delta y_t}}
          {2\norm{G_t\Delta y_t}}
 \ge\frac{\norm{G_t}^{1/2}\norm{\Delta y_t}}
          {2\norm{\Delta y_t}_{G_t}}.
\]
With $G_t=\mu_tB_t^\top B_t$ we have
$\norm{B_te_t}=\mu_t^{-1/2}\norm{e_t}_{G_t}
\le \mu_t^{-1/2}\norm{G_t}^{1/2}\norm{e_t}$ and
$\delta^{\rm N}_t=\norm{B_t\Delta y_t}
=\mu_t^{-1/2}\norm{\Delta y_t}_{G_t}$, so
$\norm{B_te_t}\le\eta\delta^{\rm N}_t$.  Thus the reconstructed step
meets the weighted contract of the outer theorem.

The operator-error budget of \cref{thm:positive-implicit-slack} supplies this
vector accuracy; its factor $1/\rho_t$ enters only logarithmic precision
terms.  Normalized-state
and norm errors $O(\eta/(J_t\rho_t))$ suffice; tight coherent tomography uses
$\Otilde(mJ_t\rho_t/\eta)$ controlled preparations.  Multiplication by RHS,
QLS, and projector costs gives \eqref{eq:positive-compressed-query}.  The
slack and selector queries depend only on the current explicit $y_t$, so each
accepted direction costs $\Otilde(mB_{\rm bit})$ to store and never materializes an
$n$-vector.  Expanding one state preparation yields
\eqref{eq:positive-compressed-gates}.  The transferred step accuracy invokes
the outer theorem, completing the composition.
\end{proof}

Stable conditioning alone is insufficient.  On the normalized tail
$0<\mu_t\leq1$, assuming $\operatorname{range}Q_t\ne\{0\}$ and $g_t\ne0$, if
$u_t=\norm{Q_tg_t}$ and $v_t=\norm{(I-Q_t)g_t}$, direct substitution gives
\begin{equation}
 J_t=\frac{\sqrt{u_t^2+\mu_t^{-4}v_t^2}}{\sqrt{u_t^2+v_t^2}}.    \label{eq:positive-forcing-dichotomy}
\end{equation}
Thus $v_t\leq C\mu_t^2u_t$ implies $J_t\leq\sqrt{1+C^2}$, while purely
inactive forcing gives $J_t=\mu_t^{-2}$.  The favorable theorem explicitly
requires stable forcing, bounded RHS cancellation, and stable output
conversion, not only a stable matrix.  Full face-repair tomography still obeys
\eqref{eq:positive-face-tomography}; classically, the same uniformly
conditioned sparse repair can be computed by Krylov iteration with only
logarithmic precision dependence.

The matched classical baseline applies the same stable operator matrix-free.
It costs $\Theta(nd_c)$ arithmetic per multiplication and
\begin{equation}
 \Otilde\!\left(nd_c\log\frac{J_t}{\eta}\right)           \label{eq:positive-classical-krylov}
\end{equation}
per direction under the uniform condition promise.  Ignoring the different
query and gate units, a necessary regime for the compressed quantum route to
improve this comparator is
\begin{equation}
 \frac{mJ_t\rho_t\Gamma_{g,t}\mathcal K_{M,t}\chi_{Q,t}}{\eta}
 =o(nd_c).                                                \label{eq:positive-compressed-regime}
\end{equation}
Randomized column sketches may improve the classical side further.  Moreover,
$C_{\rm build}(A)=\Omega(nd_cB_{\rm bit})$ for an explicit sparse input file,
so no sublinear end-to-end wall-clock claim follows unless a reusable oracle
or concise input circuit is supplied.

The condition-one prefix family of \cref{sec:lp-hardness} neither contradicts
nor validates \cref{thm:positive-compressed-qipm}.  Its hard public start is
primal--dual infeasible, its selector tail begins only at $\mu=O(1/P)$, and
its retained dimension is $m=\Theta(P)$, already giving a
$\Omegatilde(P)$ generic tomography budget.  The compatibility is exact:

\begin{proposition}[Compressed-coordinate compiler bound]
\label{prop:positive-compiler-bound}
On the prefix family, suppose a compiler uses $S$ raw coefficient queries and
returns data for a coherent coordinate oracle of either the exact central
primal point or a nonzero $x\geq0$ satisfying the robust contract
$\operatorname{res}_p(x)\leq1/100$ of
\cref{thm:lp-prefix-capstone}.  If one oracle use costs $q$ further raw
queries, then
\begin{equation}
 S+q=\Omega(P).                                         \label{eq:positive-compiler-bound}
\end{equation}
The same conclusion holds for a direction oracle whose update is nonnegative
and satisfies that residual contract.  In particular, the affine compressed
output above, including its Newton correction, has $q=O(1)$; whenever it meets
this robust approximately feasible contract, $S=\Omega(P)$.
\end{proposition}

\begin{proof}
Here the plateau height is $H=\Theta(\sqrt P)$.  Both the central vector and
every output admitted by the robust contract have $\Theta(P)$ plateau
coordinates with
$\norm x=\Theta(H\sqrt P)$ and $\norm x_\infty=\Theta(H)$.  Query the value
oracle in a uniform superposition and rotate by $x_j/(CH)$.  The success
probability is constant, so constant-round amplification prepares the hard
primal state with $S+O(q)$ raw queries.  The central and robust primal-state
clauses of \cref{thm:lp-prefix-capstone} give $\Omega(P)$, proving the claim.
For a direction oracle, add the public start before applying the same argument.
\end{proof}

This is where hardness belongs: loading or retained-dimension tomography, not
the equilibrated condition number.  A dual-only certificate without a primal
coordinate oracle remains outside that reduction.

\subsection{Checkpoint spans and the classical sketching boundary}
\label{sec:positive-checkpoint-span}

Lazy refresh strengthens the implicit representation.  If the checkpoint
directions are $\widehat z_{s_1},\ldots,\widehat z_{s_R}$, then scalar reuse
gives the exact affine representation
\begin{equation}
 y_t=y_0+\sum_{r\leq R_t}\theta_{r,t}\widehat z_{s_r},\qquad R_t\leq R.     \label{eq:positive-checkpoint-span}
\end{equation}

\begin{theorem}[Checkpoint-span trajectory oracle]
\label{thm:positive-checkpoint-span}
Under the hypotheses of \cref{thm:positive-lazy-refresh}, if the stored
vectors have $B_{\rm bit}$ bits per entry and columns of $A$
have sparsity $d_c$, a $y_t$ coordinate uses $R_t+1$ coherent reads and a
slack coordinate uses $O(d_cR_t)$ reads.  Including reversible arithmetic, a
$y_t$ query costs $\Otilde(R_tB_{\rm bit})$ gates after the public $y_0$ term
is supplied, and a slack query costs
$\Otilde(d_c(R_t+1)B_{\rm bit})$ gates.  All vector loads and scalar updates
over the path cost
\begin{equation}
 \Otilde\!\left(mB_{\rm bit}
 \left(1+\frac{GV_{\rm proj}}{\eta}\right)+TB_{\rm bit}\right). \label{eq:positive-span-ledger}
\end{equation}
\end{theorem}

\begin{proof}
Collect consecutive scalar steps using one checkpoint to obtain
\eqref{eq:positive-checkpoint-span}.  A coordinate query reads that coordinate
from each stored vector and evaluates the affine sum.  A slack query repeats
this for the at most $d_c$ entries of its column.  Each checkpoint loads one
$m$-vector and each ordinary iteration changes one scalar.  Sum these costs and
apply \cref{thm:positive-lazy-refresh}.
\end{proof}

This representation assumes accepted steps meet the outer residual contract.
A scheduled exact-tail version can avoid per-iteration residual tests; an
online version must charge them.

For a checkable tall example, let $B$ be the oriented incidence matrix of a
fixed-degree balanced bipartite expander on $m$ vertices.  For each vertex
$i$ and copy $k\in[K]$, introduce $u_{ik},v_{ik}$ with respective columns
$K^{-1/2}e_i$ and $-K^{-1/2}e_i$ and costs $a_{ik}\in[1,2]$:
\begin{equation}
 \begin{aligned}
  \min\quad&\sum_{i,k}a_{ik}(u_{ik}+v_{ik})\\
  \text{subject to}\quad&Bx+K^{-1/2}\sum_{k=1}^K
       (u_{\cdot k}-v_{\cdot k})=B\mathbf1,\qquad x,u,v\geq0.
 \end{aligned}                                           \label{eq:positive-scaled-copy-family}
\end{equation}
Here $n=|E(B)|+2mK=\Theta(mK)$; $K$ is the copy multiplicity, unrelated to
the effective QLS symbols $\mathcal K$ above.  On a sufficiently small tail,
$|y_i|/\sqrt K\leq1/2$.  The quantum algorithm and its comparator receive the
same random/coherent-query access to $(a_{ik})$ and the formula-defined sparse
columns.  If the costs arrive as an explicit file, loading already costs
$\Omega(nB_{\rm bit})$.  This family has constant direct normalization and
condition and projective constants independent of $K$, but it is not a
quantum-advantage example.

\begin{proposition}[Bounded-average sketch obstruction]
\label{prop:positive-classical-sketch}
For the scaled-copy family, the inactive contribution at vertex $i$ is
\begin{equation}
 D_{ii}(y)=\frac{1}{K}\sum_{k=1}^K\left[
 \frac{1}{(a_{ik}-y_i/\sqrt K)^2}+\frac{1}{(a_{ik}+y_i/\sqrt K)^2}\right], \label{eq:positive-copy-average}
\end{equation}
where all summands lie in a fixed bounded interval on the promised tail.
Sampling $O(\varepsilon^{-2}\log(m/\delta))$ terms independently for every
vertex produces $\norm{\widehat D-D}\leq\varepsilon$ with probability at least
$1-\delta$, in $O(m\varepsilon^{-2}\log(m/\delta))$ classical queries,
independently of $K$.
\end{proposition}

\begin{proof}
Scalar Hoeffding at each vertex and a union bound give the diagonal operator
bound.  Uniform invertibility of the equilibrated matrix and the resolvent
identity convert it to $O(\varepsilon)$ relative direction error.
\end{proof}

The column scale is essential.  The inactive block has row sparsity $2K$,
column sparsity one, and entries $K^{-1/2}$, so its standard sparse-access
normalization is $O(1)$.  Making one hidden column contribute $\Theta(1)$
requires increasing its magnitude to $\Theta(1)$, which raises the standard
sparse-access factor normalization and its two-factor product Gram
normalization to
\begin{equation}
 \alpha_{A_N}^{\rm sparse}=\Theta(\sqrt K),\qquad
 \alpha_{H_N}^{\rm sparse}=\Theta(K).                    \label{eq:positive-rare-normalization}
\end{equation}
Since $H_N$ remains diagonal with $\norm{H_N}=\Theta(1)$, a direct diagonal
encoding can keep constant normalization, but assembling its value oracle
from raw column access requires locating the exceptional columns, which is
the Grover preprocessing charged below.

At constant IPM accuracy this $\Otilde(m)$ sketch matches the unavoidable
$m$-coordinate checkpoint readout.  More generally, well-spread bounded
contributions admit classical leverage or matrix-concentration sampling.  If a
rare column is made influential, a standard sparse block encoding acquires
large normalization; locating the exceptions separately yields only Grover
preprocessing.  The open target is therefore a structured access model that
supplies a constant-normalized quantum operator without also supplying a
bounded-leverage classical sampling distribution.  No explicit sparse LP
family with that conjunction is known here.

\subsection{A block-angular border module}
\label{sec:positive-block-angular}

Consider the arrowhead Newton system
\begin{equation}
 \begin{pmatrix}
 D_1&&&B_1\\&\ddots&&\vdots\\&&D_N&B_N\\
 B_1^T&\cdots&B_N^T&-D_0
 \end{pmatrix}\binom uv=\binom fg,                      \label{eq:positive-arrowhead}
\end{equation}
where $D_i\succ0$, $D_0\succeq0$, local blocks have size at most $p$, and
the border has size $k$.  Classical block-angular IPMs use the same Schur
decomposition \cite{birge1988-computing-block-angular-karmarkar-projections,
grigoriadis1996-an-interior-point-method-for,
kempke2026-massively-parallel-interior-point-method-arrowhead}.

\begin{theorem}[Explicit border and coherent local directions]
\label{thm:positive-border-module}
Let
\[
 S=D_0+\sum_{i=1}^NB_i^TD_i^{-1}B_i\succ0,\qquad
 h=\sum_{i=1}^NB_i^TD_i^{-1}f_i-g,\qquad v=S^{-1}h.
\]
Suppose one coherent block query returns $(D_i,B_i,f_i)$ and local
factorization costs $\poly(p,k)$.  Put $M=Np+k$,
$\alpha_z=\max_i\norm{D_i^{-1/2}f_i}_\infty$, and
$\chi_h=\alpha_z/\norm v_S$ for $v\ne0$, and suppose a known upper bound
$\overline\chi_h\geq\chi_h$ is supplied.  There is a bounded-error algorithm
using
\begin{equation}
 \Otilde\!\left(\sqrt{Mk}\left(\frac{1}{\epsilon}
                   +\frac{\overline\chi_h}{\eta}\right)\right)    \label{eq:positive-border-queries}
\end{equation}
block queries that returns an explicit $\widetilde v$ satisfying
\begin{equation}
 \frac{\norm{\widetilde v-v}_S}{\norm v_S}
 \leq\frac{\epsilon+\eta}{1-\epsilon}.                 \label{eq:positive-border-error}
\end{equation}
It also implements, with one additional block query, the coherent local oracle
$|i,0\rangle\mapsto|i,D_i^{-1}(f_i-B_i\widetilde v)\rangle$.  No tomography
of the $Np$ local coordinates and no condition-number factor is required.
Without a known relative bound, one must instead request absolute inverse-energy
error $\delta_h$, replacing the second term by $\alpha_z/\delta_h$, or charge
adaptive estimation of $\norm v_S$.
\end{theorem}

\begin{proof}
Stack $C_i=D_i^{-1/2}B_i$ and $D_0^{1/2}$ into $C$, so $S=C^TC$, and stack
$z_i=D_i^{-1/2}f_i$.  The tall-matrix spectral approximation and
inverse-local-norm mean-estimation primitives used by Apers--Gribling
\cite{apers2026-quantum-speedups-linear-programming-interior} return explicit
$\widetilde S,\widetilde h$ with
$(1-\epsilon)S\preceq\widetilde S\preceq(1+\epsilon)S$ and
$\norm{\widetilde h-h}_{S^{-1}}\leq\eta\norm v_S$ at the cost in
\eqref{eq:positive-border-queries}.  Cancellation with $g$ is exactly the
$\overline\chi_h$ factor.  With $\widetilde v=\widetilde S^{-1}\widetilde h$,
\[
 \widetilde v-v=\widetilde S^{-1}
 [\widetilde h-h+(S-\widetilde S)v].
\]
The spectral sandwich bounds the first operator by $(1-\epsilon)^{-1}$ in
$S$ norm and the bracket by $(\epsilon+\eta)\norm v_S$, proving
\eqref{eq:positive-border-error}.  Block elimination gives each local formula.
\end{proof}

The $\sqrt N$ exponent is optimal in the block-oracle model.  Set the border
data $D_0=0$ and $g=0$.  For scalar blocks take a known base scenario with
$D_{\rm base}=B_{\rm base}=f_{\rm base}=1$ and hidden
$z_i\in\{0,1\}$ of Hamming weight zero or one with
$D_i=1$, $B_i=z_i$, and $f_i=0$.  Then $S=1+\sum_i z_i$, $h=1$, and the border
value is $1$ or $1/2$, while the augmented system remains uniformly
conditioned.  Constant additive accuracy decides OR, proving
$\Omega(\sqrt N)$ quantum and $\Omega(N)$ randomized block-query lower bounds.

An exact LP realization exists at the public center $x=s=e$, $y=0$ using one
dense equality row: hidden coefficient pairs are $(\delta,\delta)$ or
$(1,2\delta-1)$ with $\delta=1/(2N)$.  Its scalar normal matrix has condition
one and its centering multiplier has a constant marked/unmarked gap.  This is
a column/coefficient-query lower bound, not a constant-row-sparsity lower bound.
Indeed, for normalized uniformly conditioned bounded-degree matrices,
polynomial approximation of $1/x$ implies exponential inverse decay with graph
distance; a bounded local perturbation far from a fixed border cannot have
constant effect.  Growing condition, RHS norm, output support, or graph degree
can evade that qualified obstruction.

The module is not yet an end-to-end QIPM.  First, the product log barrier still
has parameter $\Theta(Np)$: partial minimization preserves that certificate and
can give $-N\log x$ even when the projected set admits $-\log x$.  Second,
static blocks can be cached classically.  Third, later iterations need a
nonrecursive representation of updated local variables.  The concrete open
target is a projected barrier of parameter
$\Otilde(k+\poly(p))$ whose local center and derivatives are recomputable from
one scenario block and the explicit border.  Only with such a barrier would
\cref{thm:positive-border-module} yield a genuine
$\Otilde(\sqrt N\poly(p,k))$ border-output QIPM against the $\Omega(N)$
classical block-query bound.

\subsection{Supporting exact crossover}
\label{sec:positive-crossover}

For a tall integer LP $Az\geq b$ in dimension $d$, with coefficient magnitude
at most $2^\ell$, row sparsity $r$, and a unique nondegenerate strictly
complementary optimum, let
\[
 D=(\sqrt r\,2^\ell)^d,\qquad
 M=\max\{1,\sqrt r\,2^\ell\sqrt d
                 (\sqrt r\,2^\ell)^{d-1}\}.
\]
A feasible point of objective error
\begin{equation}
 E=\frac{1}{8D^2M}                                      \label{eq:positive-crossover-accuracy}
\end{equation}
gives true row slacks at most $1/(8D)$ on active rows and at least $7/(8D)$ on
inactive rows.  Evaluate each slack with error at most $1/(8D)$ and mark it
when the computed value is at most $1/(2D)$; these inequalities make the
predicate exact.  Its reversible arithmetic uses
$O(d(\ell+\log d))$-bit precision (and the corresponding polynomial gate
cost).  Thus $\log(1/E)=O(d(\ell+\log d))$, independent of the number $n$ of
rows.  Coherent enumeration identifies all binding rows in
$\Otilde(\sqrt{nd})$ row queries; exact rational solves for the primal and dual
variables and a Grover feasibility check then return an exact KKT certificate.
The enumeration cost is tight for this sparse-certificate wrapper by $d$
independent search blocks.

This is only a supporting crossover lemma for the tall-row-oracle QIPM of
Apers--Gribling.  It assumes the structural promise and coherent row access,
and explicit input loading can erase sublinear wall time.  More importantly,
Objois--Vladu give an exact quantum Clarkson algorithm with
$\Otilde(\sqrt n\,d^3)$ row queries under a comparable nondegeneracy promise
(a unique optimal basis of $d$ constraints) but without the integrality,
bit-size, or strict-complementarity structure, and
Dadush--V\'egh--Zambelli give an exact primal--dual oracle framework
\cite{positive-objois2025-adaptive-sparsification,
positive-dadush2026-exact-oracle-lp}.  The wrapper is therefore not a
best-known exact-LP result.

\subsection{Synthesis}
\label{sec:positive-synthesis}

The positive modules identify three real escape routes from the obstructions
of \cref{sec:beyond-kappa,sec:lp-hardness,sec:preconditioning}: amortize
classical readout by projective variation, generate changing diagonals from a
smaller affine iterate, or request an explicit low-dimensional border rather
than the full Newton vector.  Each route has a matching condition that cannot
be dropped.  Lazy reuse, as bounded here, needs finite realized $V_{\rm proj}$
and a fixed residual tolerance; partition-free equilibration needs a certified
tail, active gaps, stable forcing, and bounded RHS/output factors; compressed output
needs $m=o(n)$ and charges its coordinate compiler; block aggregation needs a
small-parameter projected barrier and a nonrecursive local-state
representation.

The first-class negative results explain why simply combining the modules is
not enough.  A standard feasible inexact neighborhood permits logarithmically
divergent face cycling.  Proximal locking trades motion for residual and gives
no residual-safe condition improvement.  Unfiltered global repair recreates
the ill-conditioned Newton inverse.  Generic full-vector tomography recreates
inverse-precision cost.  Benign tall averages admit equally benign classical
sketches, while rare influential columns move the cost into normalization or
Grover search.

Accordingly, this section proves reusable operators and conditional
end-to-end compressed-output theorems, not an unconditional quantum advantage.
A genuine advantage still requires one explicit problem family and access
model for which all of the following hold simultaneously: stable direct
normalization and inverse, stable forcing and RHS preparation, finite
projective variation or a low-parameter barrier, cheap static construction,
a useful compressed output, and a classical lower bound surviving randomized
sketching and preprocessing.  None of the modules above silently supplies a
missing item from that list.

\section{Structural Identities and Obstruction Tools}
\label{sec:structural-tools}

The preceding lower bounds use several small pieces of structure that are
useful beyond any one hard family.  This section collects them in one place.
The first two tools turn symmetry and cross-instance averaging into invariant
certificates.  The third says how much matrix rank an SDP Newton step must
carry to make complementarity progress.  The remaining tools delimit what
sparse support, bounded treewidth, and geometric locality can and cannot buy.
They are supporting results: most of their ingredients are classical, and we
separate the new obstruction formulations from those ingredients throughout.

\subsection{Symmetry as virtual work}
\label{sec:structural-symmetry}

Work in a finite-dimensional Euclidean space.  Let $R$ be an orthogonal
involution, let $\Omega$ be an open convex set with $R\Omega=\Omega$, and let
$F:\Omega\to\R$ be differentiable, convex, and $R$-invariant.  Assume also
that $R^Tc=c$ and use the stationarity convention
\begin{equation}
 A^Ty=c+\mu\nabla F(x),\qquad \mu>0.                 \label{eq:structural-stationarity}
\end{equation}
Thus the equality term in the Lagrangian is $-y^T(Ax-b)$.  Put
\begin{equation}
 x_-:=\frac{I-R}{2}x,\qquad d:=x-Rx=2x_-,\qquad r_-:=Ax_-.
                                                               \label{eq:structural-antisymmetric-source}
\end{equation}
The definition $r_-=Ax_-$, rather than a selected block of $b$, is what makes
the following statement insensitive to mixed equality rows.

\begin{theorem}[Involution work identity]
\label{thm:structural-involution-work}
Under the preceding assumptions,
\begin{align}
 r_-^Ty
 &=\mu\inner{x_-}{\nabla F(x)}                                      \notag\\
 &=\frac{\mu}{4}\inner{x-Rx}{\nabla F(x)-\nabla F(Rx)}             \label{eq:structural-work-chain}\\
 &=\frac{\mu}{2}D_F(Rx,x)
  =\frac{\mu}{2}D_F(x,Rx)\geq0,                                    \notag
\end{align}
where $D_F(u,v)=F(u)-F(v)-\inner{\nabla F(v)}{u-v}$.
If $F\in C^2(\Omega)$ and
\begin{equation}
 \overline H_x:=\int_0^1\nabla^2F(Rx+td)\,dt,
\end{equation}
then, exactly,
\begin{equation}
 r_-^Ty=\frac{\mu}{4}d^T\overline H_xd
       =\mu x_-^T\overline H_xx_-.                         \label{eq:structural-integrated-hessian}
\end{equation}
For a standard self-concordant barrier, set
$x_+=(x+Rx)/2$ and $\delta=\norm{x_-}_{x_+}<1$.  Then the
integrated chord comparison gives sharper local constants than a worst-point
comparison, and the universal lower bound can be strengthened further:
\begin{equation}
 \boxed{
 \mu\left(1-\delta+\frac{\delta^2}{3}\right)\delta^2
 \leq\frac{\mu\delta^2}{1+\delta}
 \leq r_-^Ty\leq\frac{\mu\delta^2}{1-\delta}.}
                                                               \label{eq:structural-sharp-chord}
\end{equation}
\end{theorem}

\begin{proof}
The invariant objective does no antisymmetric work:
\[
 \inner{x_-}{c}=\tfrac12\inner{x}{c-R^Tc}=0.
\]
Pairing \eqref{eq:structural-stationarity} with $x_-$ gives the first
equality in \eqref{eq:structural-work-chain}.  Differentiating
$F(Rx)=F(x)$ gives $\nabla F(Rx)=R\nabla F(x)$.  Since $Rd=-d$,
\[
 \inner{d}{\nabla F(Rx)}
 =\inner{Rd}{\nabla F(x)}=-\inner{d}{\nabla F(x)},
\]
which proves the second equality.  Gradient monotonicity proves
nonnegativity.  Invariance also gives
\[
 D_F(Rx,x)=D_F(x,Rx)=\inner{\nabla F(x)}{x-Rx},
\]
so the Bregman equalities follow.  The fundamental theorem of calculus on
the chord gives \eqref{eq:structural-integrated-hessian}.

For the last claim, parameterize the chord as $x_++ux_-$,
$-1\leq u\leq1$.  Self-concordant Hessian comparison bounds its directional
energy by
\[
 (1-|u|\delta)^2\delta^2
 \leq x_-^T\nabla^2F(x_++ux_-)x_-
 \leq(1-|u|\delta)^{-2}\delta^2.
\]
The relevant chord averages are
\[
 \int_0^1(1-u\delta)^2\,du
 =1-\delta+\delta^2/3,
 \qquad
 \int_0^1(1-u\delta)^{-2}\,du=(1-\delta)^{-1},
\]
so these constants are the exact integrals of the comparison envelopes, not
generally attained bounds.  For the stronger lower bound, put
\[
 \phi(u)=x_-^T\nabla^2F(x_++ux_-)x_-.
\]
Self-concordance gives
$|(\phi^{-1/2})'|\leq1$ and $\phi(0)=\delta^2$, hence
$\phi(u)\geq\delta^2/(1+u\delta)^2$ for $0\leq u\leq1$.
Invariance makes $\phi$ even and
$r_-^Ty=\mu\int_0^1\phi(u)\,du$, which proves the middle lower bound.
It strictly dominates the chord-comparison lower constant for
$0<\delta<1$, since
\[
 \frac1{1+\delta}-\left(1-\delta+\frac{\delta^2}{3}\right)
 =\frac{\delta^2(2-\delta)}{3(1+\delta)}>0.
\]
\end{proof}

No rank assumption on $A$ is hidden here.  If $\widetilde y$ is another
multiplier, then $A^T(\widetilde y-y)=0$ and hence
$(Ax_-)^T(\widetilde y-y)=0$.  Moreover, under any invertible row change
\begin{equation}
 A'=UA,\qquad y'=U^{-T}y,\qquad r_-'=Ur_-,
 \qquad (r_-')^Ty'=r_-^Ty.                              \label{eq:structural-row-covariance}
\end{equation}
Only the source--multiplier pairing is invariant; an ill-conditioned $U$
can shrink a multiplier norm while enlarging the source or the input
normalization.  If rows genuinely split as $A_+R=A_+$ and $A_-R=-A_-$
and $Ax=b=(b_+,b_-)$, then $r_-=(0,b_-)$ and the pairing is the familiar
$b_-^Ty_-$.  Arbitrary
row mixing preserves \eqref{eq:structural-row-covariance}, but not that
coordinate shortcut.

Two specializations fix the constants used earlier.  For paired positive
variables $x=(u,v)$, $R(u,v)=(v,u)$, and
$F=-\sum_j(\log u_j+\log v_j)$, put $d=u-v$ and $q=u+v$.  Then
\begin{equation}
 r_-^Ty=\frac{\mu}{2}\sum_j\frac{d_j^2}{u_jv_j}
       =\sum_j\frac{2\mu d_j^2}{q_j^2-d_j^2}.             \label{eq:structural-paired-log}
\end{equation}
This is exactly \eqref{eq:frontier-central-work}; the proof in
\cref{thm:frontier-work} is the paired-coordinate specialization, not a
second general identity.

For the log-determinant barrier on $\Sym_{++}^r$, let
$\mathcal R(X)=QXQ$ for a symmetric orthogonal $Q$, and set
$X_-=(X-QXQ)/2$.  With the Frobenius pairing,
\begin{align}
 (\mathcal A X_-)^Ty
 &=\frac{\mu}{2}\bigl[\tr(QXQX^{-1})-r\bigr]              \label{eq:structural-stein}\\
 &=\mu\int_0^1
 \norm{Z_t^{-1/2}X_-Z_t^{-1/2}}_F^2\,dt,
 \quad Z_t=QXQ+t(X-QXQ).                                  \label{eq:structural-logdet-integral}
\end{align}
The bracketed trace quantity in the first line is half the symmetrized Stein
loss, so the work is $\mu/4$ times that symmetrization.  This is not the
distinct Jensen--Bregman LogDet, or S-divergence.  The second line is the
exact noncommutative integrated-Hessian formula; no commuting assumption is
used.  If $\ell I\preceq X\preceq UI$, it lies between
$\mu\norm{X_-}_F^2/U^2$ and $\mu\norm{X_-}_F^2/\ell^2$.
Here $U$ is a spectral endpoint, unrelated to the row-change matrix in
\eqref{eq:structural-row-covariance}.

At a fixed start $x^0=Rx^0$, the Newton version is equally direct.  If
$G_0\succ0$ commutes with $R$ and
\begin{equation}
 A\Delta x=r_p,\qquad G_0\Delta x-A^T\Delta y=h_+,
 \qquad Rh_+=h_+,
\end{equation}
then, for $\Delta x_-=(I-R)\Delta x/2$,
\begin{equation}
 \boxed{(A\Delta x_-)^T\Delta y
 =\Delta x_-^TG_0\Delta x_-\geq0.}                       \label{eq:structural-newton-work}
\end{equation}
This recovers \eqref{eq:frontier-newton-work}.  A nonzero antisymmetric
forcing $h_-$ adds $-\inner{\Delta x_-}{h_-}$ and destroys the sign.

The involution is only the two-element case.  Let a finite or compact group
$\mathcal G$ act orthogonally, and assume $g\Omega=\Omega$ for every
$g\in\mathcal G$.  Let
\begin{equation}
 P_0=\int_{\mathcal G}g\,d\nu(g),\qquad
 x_0=P_0x,\qquad x_\perp=(I-P_0)x                       \label{eq:structural-reynolds}
\end{equation}
be the Reynolds projection.  The orbit is compact, and Carath\'eodory's
theorem shows that its convex hull is a compact subset of the convex domain;
thus $x_0\in\Omega$ and the segment from $x_0$ to $x$ stays in $\Omega$.

\begin{theorem}[Compact-group work]
\label{thm:structural-group-work}
If $F$ and $c$ are $\mathcal G$-invariant and
\eqref{eq:structural-stationarity} holds, then
\begin{equation}
 (Ax_\perp)^Ty
 =\mu\inner{x_\perp}{\nabla F(x)-\nabla F(x_0)}\geq0.
                                                               \label{eq:structural-group-central}
\end{equation}
When the Hessian exists, the complete chain is
\begin{align}
 (Ax_\perp)^Ty
 &=\mu\inner{x_\perp}{\nabla F(x)-\nabla F(x_0)}          \notag\\
 &=\mu\int_0^1\inner{x_\perp}
 {\nabla^2F(x_0+tx_\perp)x_\perp}\,dt                     \notag\\
 &\geq0.                                                  \label{eq:structural-group-hessian}
\end{align}
At a group-fixed Newton start, suppose $G_0\succeq0$ commutes with the action and
\begin{equation}
 A\Delta x=r_p,\qquad
 G_0\Delta x-A^T\Delta y=h_0,\qquad gh_0=h_0
 \quad(g\in\mathcal G).                                   \label{eq:structural-group-newton-system}
\end{equation}
Define $\Delta x_\perp=(I-P_0)\Delta x$.  Then
\begin{equation}
 (A\Delta x_\perp)^T\Delta y
 =\Delta x_\perp^TG_0\Delta x_\perp\geq0.               \label{eq:structural-group-newton}
\end{equation}
\end{theorem}

\begin{proof}
Both $c$ and $\nabla F(x_0)$ are fixed by the group and hence orthogonal to
$x_\perp$.  Pair stationarity with $x_\perp$, then apply gradient
monotonicity between $x_0$ and $x$.  Integration gives
\eqref{eq:structural-group-hessian}.  The Newton identity follows by pairing
the stationarity equation with the nontrivial-representation component;
$G_0$ preserves the orthogonal representation decomposition.
\end{proof}

For coordinate permutations, $P_0$ replaces each orbit by its average; for
the logarithmic barrier the work on an orbit $O$ is
$\mu(\bar x_O\sum_{i\in O}x_i^{-1}-|O|)$.  The simplex gives the same
formula with $\bar x=1/n$.  For the standard $SO(k)$ action and the unit-ball
barrier, it is $2\mu\norm z^2/(1-\norm z^2)$.  These examples explain the
permutation and rotation gadgets excluded in \cref{sec:frontier-synthesis}.

\begin{remark}[Scope and provenance]
\label{rem:structural-work-scope}
The sign reverses with the opposite multiplier convention.  An asymmetric
objective or Newton forcing contributes an uncontrolled linear term.  Mixed
rows require the canonical source $Ax_-$; convexity without strict curvature
gives only nonnegativity.  The scalar identity extends to a nonorthogonal
linear involution by using $\nabla F(Rx)=R^{-T}\nabla F(x)$, but the Euclidean
orthogonal decomposition does not.  Finally, arbitrary edge transports need
not arise from one simultaneous public symmetry action and may have no public
Reynolds projector.  Even when they do come from a common representation,
the Reynolds identity supplies no cut-local conservation law; the
signed/gain incidence and holonomy analysis of \cref{sec:sdp-hardness} is
separate.

The divergence in \eqref{eq:structural-work-chain} is the classical
symmetrized Bregman divergence
\cite{structural-bregman1967,structural-nielsen2009-bregman}; Reynolds
reduction and symmetric criticality are likewise classical
\cite{structural-palais1979-symmetric-criticality}.  The contribution used
in this paper is narrower: the row-mixing-invariant source
\eqref{eq:structural-antisymmetric-source} is an obstruction template for
sparse QIPM gadgets.
\end{remark}

\subsection{Central mixtures in semidefinite programming}
\label{sec:structural-sdp-mixtures}

The scalar Kantorovich argument in \cref{sec:frontier-mixture-supplements}
has a fully noncommutative analogue.  Consider SDPs with common $C,b$ and
input-dependent measurement maps $\mathcal A_\xi$, with matrix order $r\geq1$
and $N\geq1$ neighboring inputs.  Fix one Frobenius-isometric
$\operatorname{svec}:\Sym^r\to\R^{r(r+1)/2}$.  All sparsity and coefficient
bounds below refer to the resulting matrix $A_\xi$, not to the unvectorized
data.  At a base input $\varsigma$, suppose every one-bit neighbor
$\varsigma^{(i)}$ has an exact center at the same $\mu>0$:
\begin{equation}
 \mathcal A_{\varsigma^{(i)}}(X_i)=b,\qquad
 \mathcal A_{\varsigma^{(i)}}^*(y_i)+S_i=C,\qquad
 S_i=\mu X_i^{-1},\qquad X_i\succ0.                       \label{eq:structural-sdp-centers}
\end{equation}
For weights $w_i\geq0$, $\sum_iw_i=1$, write
\begin{equation}
 \begin{aligned}
 X_w&=\sum_iw_iX_i,& y_w&=\sum_iw_iy_i,&
 S_w&=\sum_iw_iS_i,\\
 D_i&=X_i-X_w,&&&
 Z_w&=X_w^{1/2}(S_w/\mu)X_w^{1/2}.
 \end{aligned}                                            \label{eq:structural-sdp-average}
\end{equation}

\begin{lemma}[Matrix multiplicative variance]
\label{lem:structural-matrix-variance}
The following identity is exact:
\begin{equation}
 \boxed{
 Z_w-I=X_w^{-1/2}\left(\sum_iw_iD_iX_i^{-1}D_i\right)X_w^{-1/2}.}
                                                               \label{eq:structural-matrix-variance}
\end{equation}
Consequently $Z_w\succeq I$, with equality exactly when all positively
weighted $X_i$ equal $X_w$.
\end{lemma}

\begin{proof}
Expanding the middle term and using $\sum_iw_iX_i=X_w$ gives
\[
 \sum_iw_iD_iX_i^{-1}D_i
 =X_w\left(\sum_iw_iX_i^{-1}\right)X_w-X_w.
\]
Congruence by $X_w^{-1/2}$ proves the identity.  Every un-congrued summand is
positive semidefinite, and a positive weighted sum of them vanishes only
when each corresponding $D_i$ vanishes.
\end{proof}

Identity \eqref{eq:structural-matrix-variance} is an elementary exact
expansion of the arithmetic--harmonic defect; it is not claimed as a new
matrix inequality.

\begin{theorem}[Noncommutative central sandwich]
\label{thm:structural-sdp-kantorovich}
If $0<m\leq M$, $mI\preceq X_i\preceq MI$ for every $i$, and $R=M/m$,
then
\begin{equation}
 \boxed{I\preceq Z_w\preceq K(R)I,\qquad
 K(R)=\frac{(R+1)^2}{4R}.}                                \label{eq:structural-sdp-sandwich}
\end{equation}
The constant is sharp, already for an equal scalar mixture of $m$ and $M$.
No two matrices in the mixture need commute.
\end{theorem}

Here and below $R=M/m$ is a spectral ratio, unrelated to the involution $R$
in \cref{sec:structural-symmetry}.

\begin{proof}
The lower bound is \cref{lem:structural-matrix-variance}.  Functional
calculus, applied to one $X_i$ at a time, gives
\begin{equation}
 X_i+mM X_i^{-1}\preceq(m+M)I.
\end{equation}
Averaging and congruencing by $X_w^{1/2}$ yields
\[
 Z_w\preceq\frac{(m+M)X_w-X_w^2}{mM}.
\]
Every eigenvalue $\lambda$ of $X_w$ is in $[m,M]$, and
\[
 \frac{\lambda(m+M-\lambda)}{mM}
 \leq\frac{(m+M)^2}{4mM}=K(R).
\]
Only $X_i$ and $X_w$ individually entered functional calculus, so no
commutation was assumed.
\end{proof}

The variance identity supplies a dimension-free refinement of the spectral
ratio bound.  If only $X_i\succeq mI$ with $m>0$ is assumed, then
\begin{align}
 \norm{Z_w-I}_F
 &\leq\frac1{m^2}\sum_iw_i\norm{D_i}_F^2,                 \label{eq:structural-sdp-variance-frobenius}\\
 \norm{Z_w-I}_{\mathrm{op}}
 &\leq\frac1{m^2}\sum_iw_i\norm{D_i}_{\mathrm{op}}^2.   \label{eq:structural-sdp-variance-operator}
\end{align}
Indeed, the two inverse factors come separately from the outer congruence
and $X_i^{-1}$ in \eqref{eq:structural-matrix-variance}.  Multiplying by
$\mu$ gives the corresponding complementarity-residual bounds.  They survive
standard recentering: if
\begin{equation}
 \mu_w=\frac{\inner{X_w}{S_w}}r=\mu\alpha,\qquad
 \alpha=\frac{\tr Z_w}{r}\geq1,
\end{equation}
then
\[
 \mu_w^{-1}X_w^{1/2}S_wX_w^{1/2}-I
 =\frac{(Z_w-I)-[\tr(Z_w-I)/r]I}{\alpha}.
\]
Removing the trace component is an orthogonal projection in Frobenius norm,
and it does not enlarge the operator range because $Z_w-I\succeq0$.

The objective identity uses no feasibility of the average.  Every point in
\eqref{eq:structural-sdp-centers} has algebraic gap $r\mu$, so linearity gives
\begin{equation}
 \boxed{\inner C{X_w}-b^Ty_w=r\mu.}                       \label{eq:structural-sdp-gap}
\end{equation}
By contrast,
$\inner{X_w}{S_w}-[\inner C{X_w}-b^Ty_w]
=\mu\tr(Z_w-I)$ is nonnegative.  Under the common sandwich
$0<m\leq M$ and $mI\preceq X_i\preceq MI$, it is at most
$r[K(M/m)-1]\mu$.  This is an arithmetic--harmonic Jensen gap at a
generally base-infeasible point, not an application of weak duality.

We next combine the operator statement with the sparse single-flip residual.
Let $s_r,s_c$ bound the row and column sizes of one fixed union support
of the matrices $A_\xi$, let
$M_{\rm dep}$ count input-dependent positions, each of which is a function
of one designated bit $\varsigma_{\ell(r,c)}$ (a bit may designate several
positions), let their magnitudes be at most $B\geq0$, and assume that every
coordinate of
$\operatorname{svec}(X_i),y_i,\operatorname{svec}(S_i)$ has magnitude at
most $H\geq0$.  Set
\begin{equation}
 s_{\rm KKT}=\max\{s_r,s_c+1\},\qquad
 B_{\rm KKT}=\max\{B,1\}.
\end{equation}
For the uniform average of all $N$ one-bit neighbors, the same incidence
count as \cref{thm:frontier-single-flip}, now once in the primal block and
once in the dual block, gives
\begin{equation}
 \left\|
 \begin{bmatrix}
  \mathcal A_\varsigma(X_w)-b\\
  \operatorname{svec}(\mathcal A_\varsigma^*(y_w)+S_w-C)
 \end{bmatrix}\right\|_2
 \leq\frac{2B_{\rm KKT}H\sqrt{2s_{\rm KKT}M_{\rm dep}}}{N}.
                                                               \label{eq:structural-sdp-single-flip}
\end{equation}
Keeping the dual multipliers signed gives the stronger weighted estimate
\begin{equation}
 \left\|
 \begin{bmatrix}
  \mathcal A_\varsigma(X_w)-b\\
  \operatorname{svec}(\mathcal A_\varsigma^*(y_w)+S_w-C)
 \end{bmatrix}\right\|_2^2
 \leq4(s_r+s_c)B^2H^2\sum_i M_iw_i^2,
 \label{eq:structural-sdp-weighted-residual}
\end{equation}
where $M_i$ counts dependent coefficient positions assigned to bit $i$.
For uniform weights this is
$4(s_r+s_c)B^2H^2M_{\rm dep}/N^2$, which implies
\eqref{eq:structural-sdp-single-flip}.
The PSD variables need no coordinatewise sign split; convexity preserves the
cone, and the identity block on $S$ is not input dependent.
For multibit coefficients the degree factor of
\cref{prop:frontier-multibit} enters, as in
\cref{sec:frontier-mixture-supplements}.

\begin{proposition}[SDP mixture phase boundary]
\label{prop:structural-sdp-phase}
Suppose $0<m\leq M$, every neighboring center obeys the common sandwich
$mI\preceq X_i\preceq MI$, and $R=M/m$.  Assume the one-bit-local dependence
stated before \eqref{eq:structural-sdp-single-flip}, that all $N$ bits are
sensitive, and that the selected neighboring centers have a common
mixture-stable wrong
output.  This may be either a fixed affine
score with a uniform signed margin, or a fixed symmetric contraction $O$
such that
\begin{equation}
 f(\xi)\inner O{X^\xi}\geq\gamma\tr X^\xi               \label{eq:structural-sdp-povm-margin}
\end{equation}
for an output $f\in\{-1,+1\}$ and a positive margin $0<\gamma\leq1$.
The contraction condition means $-I\preceq O\preceq I$; the binary
measurement has effects $(I+O)/2$ and $(I-O)/2$ on $X/\tr X$.
If a base-instance contract demands the correct
output from every positive-definite triple with combined residual at most
$\varepsilon\geq0$, algebraic gap $r\mu$, and either common-$\mu$ or point-centered
operator width at most $\theta$, then soundness requires
\begin{equation}
 \boxed{
 \frac{8s_{\rm KKT}B_{\rm KKT}^2H^2M_{\rm dep}}{N^2}>\varepsilon^2
 \quad\text{or}\quad
 \frac{(R-1)^2}{4R}>\theta.}                             \label{eq:structural-sdp-phase}
\end{equation}
For a Frobenius neighborhood of width $\eta_F$, the second alternative is
\begin{equation}
 \sqrt r\,\frac{(R-1)^2}{4R}>\eta_F,                     \label{eq:structural-sdp-frobenius-phase}
\end{equation}
unless \eqref{eq:structural-sdp-variance-frobenius} is smaller.
The residual alternative can be sharpened to
$4(s_r+s_c)B^2H^2M_{\rm dep}/N^2>\varepsilon^2$.
Under point centering the algebraic gap remains $r\mu$, not $r\mu_w$.
\end{proposition}

\begin{proof}
Equation \eqref{eq:structural-sdp-single-flip} gives the first certificate,
\eqref{eq:structural-sdp-gap} gives the second, and
\cref{thm:structural-sdp-kantorovich} gives operator width at most
$K(R)-1=(R-1)^2/(4R)$ under either centering convention.  The affine score is
preserved by averaging.  Under \eqref{eq:structural-sdp-povm-margin}, trace
weighting gives
\[
 -f(\varsigma)\tr\left(O\frac{X_w}{\tr X_w}\right)\geq\gamma,
\]
so the normalized density matrix has constant bias for the wrong sign.  If
neither strict inequality in \eqref{eq:structural-sdp-phase} held, this wrong
output would satisfy the contract.  Finally,
$\norm{Z_w-I}_F\leq\sqrt r\norm{Z_w-I}_{\rm op}$ gives
\eqref{eq:structural-sdp-frobenius-phase}; the same projection argument used
after \eqref{eq:structural-sdp-variance-operator} covers point centering.
\end{proof}

Since
\begin{equation}
 K(R)-1=\sinh^2\!\left(\frac{\log R}{2}\right),
\end{equation}
under a common sandwich $0<m\leq M$ and $mI\preceq X_i\preceq MI$, with
$R=M/m\geq1$, the worst-case common-$\mu$ Frobenius guarantee of width
$\eta_F\geq0$ based only on this global ratio is equivalent to
\begin{equation}
 \log R\leq2\operatorname{arsinh}(\sqrt{\eta_F}\,r^{-1/4})
 =O(r^{-1/4}).                                             \label{eq:structural-sdp-high-dim-law}
\end{equation}
Indeed, the equal mixture $X_1=mI$, $X_2=MI$ attains
$Z_w=K(R)I$ in every dimension.  The same ratio criterion is sufficient
for point centering, but this example has zero point-centered width, so
it does not establish necessity for that convention.  A finite consequence
of the criterion is $\log R\leq2\sqrt{\eta_F}\,r^{-1/4}$.
Actual Frobenius variance can be small even when the global spectral
ratio is not.

The common sandwich is load-bearing.  Individual condition numbers cannot
replace its ratio: the scalars $X_1=1,X_2=9$ each have condition number one,
but their equal mixture has $\norm{Z_w-I}_{\rm op}=16/9$.

\begin{remark}[Seven scope limits]
\label{rem:structural-sdp-mixture-scope}
The common-basis bounds $mI\preceq X_i\preceq MI$ are strong.  Operator
control incurs a generic $\sqrt r$ loss in Frobenius norm.  The factors
$m^{-2}$ in the additive-variance bounds are real.  Every sparsity and height
parameter depends on the chosen Frobenius-isometric $\operatorname{svec}$.
Output stability is a separate hypothesis, and
\eqref{eq:structural-sdp-povm-margin} says nothing about normalized
$\operatorname{svec}$ amplitudes.  The mixture is existential, so the theorem
is not a query lower bound without a reduction.  Finally,
\eqref{eq:structural-sdp-gap} is an algebraic identity at an infeasible point.
These qualifications are part of the theorem's usable boundary.

The arithmetic--harmonic operator inequality and its Kantorovich reverse are
classical
\cite{structural-kubo1980-means,structural-ando1983-arithmetic,
structural-marshall1990-kantorovich,lin2013-on-an-operator-kantorovich-inequality}.
The nearest slogan-level collision is conic-IPM warm starting
\cite{structural-chen2025-conic-warmstart}, which moves one old solution to a
centered start for a perturbed problem.  The conjunction here is different:
cross-instance averaging, a sparse incidence budget, simultaneous primal and
dual residuals, and a wrong-output boundary.  The scalar LP corollary was
already proved in \cref{prop:frontier-central-dichotomy}; the present result is
the SDP operator theory behind it.
\end{remark}

\paragraph{Formal verification.}
The Lean development in \path{formal/QipmFormal/SDPMixture/} verifies the
finite matrix-mixture argument in this subsection.  It includes the
variance identity and equality case, the sharp noncommutative sandwich,
operator and Frobenius variance bounds, trace recentering, and the
algebraic-gap discrepancy.  A constructed Frobenius-isometric symmetric
vectorization connects actual SDP measurement maps and their adjoints to
weighted, uniform, and multibit residual bounds.  The soundness results
combine these properties with explicit affine or trace-normalized
observable output contracts under both centering conventions.  The
scalar ratio criterion, endpoint sharpness, and the distinction between
the two centering conventions are also checked.  The claim map,
qualifications, and reproduction command are in \path{formal/SDP_MIXTURE.md}.
The development permits only \texttt{propext}, \texttt{Classical.choice},
and \texttt{Quot.sound}.  It does not assert existence of the neighboring
centers for arbitrary inputs or formalize an IPM trajectory or quantum
query reduction.

\subsection{Rank required for SDP complementarity progress}
\label{sec:structural-rank-progress}

Low rank at an SDP optimum does not imply low rank of an accurate Newton
direction.  The clean statement is in Nesterov--Todd (NT) coordinates.  At
an exact center $XS=\mu I$, let $W$ be the NT scaling, so $WSW=X$, put
$P=W^{-1/2}$, and define
\begin{equation}
 U=P\Delta XP^T,\qquad V=P^{-T}\Delta SP^{-1}.
\end{equation}
Here $U$ denotes the NT-scaled primal direction and is unrelated to the
row-change matrix and spectral endpoint denoted by $U$ above.
Write the symmetric scaled complementarity equation operationally as
\begin{equation}
 \sqrt\mu(U+V)=-(1-\sigma)\mu I+\widetilde R.             \label{eq:structural-scaled-complementarity}
\end{equation}
This definition avoids dependence on competing residual conventions.

\begin{theorem}[Sharp rank--progress barrier]
\label{thm:structural-rank-progress}
Let $0\leq\sigma<1$.  For $1\leq p<\infty$, if
\begin{equation}
 \norm{\widetilde R}_{S_p}
 \leq\theta(1-\sigma)\mu n^{1/p},\qquad 0\leq\theta<1,
                                                               \label{eq:structural-rank-residual}
\end{equation}
then
\begin{equation}
 \boxed{\rank(\Delta X)+\rank(\Delta S)
 \geq\left\lceil(1-\theta^p)n\right\rceil.}             \label{eq:structural-rank-bound}
\end{equation}
For spectral error,
\begin{equation}
 \norm{\widetilde R}_{\mathrm{op}}<(1-\sigma)\mu
 \quad\Longrightarrow\quad
 \rank(\Delta X)+\rank(\Delta S)\geq n.                 \label{eq:structural-rank-spectral}
\end{equation}
The result is invariant under primal congruence and holds blockwise on a
product PSD cone.
\end{theorem}

\begin{proof}
Congruence preserves rank, so, with
$k=\rank(\Delta X)+\rank(\Delta S)$,
\[
 \rank(U+V)\leq\rank(U)+\rank(V)=k.
\]
Equation \eqref{eq:structural-scaled-complementarity} makes
$\widetilde R$ the error of approximating $(1-\sigma)\mu I$ by a
rank-at-most-$k$ matrix.  Eckart--Young--Mirsky gives
\[
 \norm{\widetilde R}_{S_p}
 \geq(1-\sigma)\mu(n-k)_+^{1/p}.
\]
Combining this with \eqref{eq:structural-rank-residual} proves
\eqref{eq:structural-rank-bound}.  If $k<n$, a kernel vector of $U+V$
sees residual exactly $(1-\sigma)\mu$, proving
\eqref{eq:structural-rank-spectral}.
\end{proof}

The bound is sharp as a matrix statement: take
$U=-(1-\sigma)\sqrt\mu\diag(I_k,0)$ and $V=0$.  The residual is precisely
$(1-\sigma)\mu\diag(0,I_{n-k})$.  The approximation theorem used in the
proof is classical \cite{structural-eckart1936-low-rank,
structural-mirsky1960-unitarily-invariant}; the rank--progress interpretation
is the structural content here.

\subsubsection{A path witness and the norm audit}

Let $L$ be the Laplacian of the $n$-vertex path and consider
\begin{equation}
 \min\ \inner L X\qquad\text{subject to}\qquad
 \tr X=1,\quad X\succeq0.                                \label{eq:structural-path-sdp}
\end{equation}
For $u=\mathbf1/\sqrt n$, the unique optimum is $X^*=uu^T$, while the
dual optimum has slack $S^*=L$; hence strict complementarity holds.  Its
exact central path is
\begin{equation}
 y(t)=-t,\quad S(t)=L+tI,\quad Z(t)=\tr(S(t)^{-1}),
 \qquad \mu(t)=Z(t)^{-1},\quad X(t)=\mu(t)S(t)^{-1}.      \label{eq:structural-path-center}
\end{equation}
With $T(t)=\tr(S(t)^{-2})$, differentiation gives
\begin{equation}
 \dot S=I,\qquad \dot\mu=\mu^2T,\qquad
 \dot X=\dot\mu S^{-1}-\mu S^{-2}.                       \label{eq:structural-path-tangent}
\end{equation}
If $0=\lambda_0<\lambda_1\leq\cdots$ are the eigenvalues of $L$, the
eigenvalues of $\dot X$ are
\begin{equation}
 \frac{\mu}{(t+\lambda_i)^2}
 \left[\frac TZ(t+\lambda_i)-1\right].                   \label{eq:structural-path-raw-spectrum}
\end{equation}
The zero-mode entry is negative.  Every other entry is positive whenever
$t\leq\lambda_1/(2\sqrt n)$, because
$T/Z>1/(t+\lambda_1)$ follows from
$\lambda_1^2+\lambda_1t-(n-1)t^2>0$.  Thus both tangent matrices are full
rank arbitrarily close to a rank-one optimum.  The exact Newton direction
targeting $\sigma\mu$ is the scalar multiple
\begin{equation}
 (\Delta X,\Delta S,\Delta y)
 =-\frac{(1-\sigma)\mu}{\dot\mu}(\dot X,\dot S,\dot y).
                                                               \label{eq:structural-path-newton}
\end{equation}

Full algebraic rank alone is weak, so the scaled spectrum matters.  Put
$d_i=t+\lambda_i$ and $\beta=T/Z$.  Since
$W=\sqrt\mu S^{-1}$ on \eqref{eq:structural-path-center}, the NT-scaled
tangent eigenvalues are
\begin{equation}
 u_i=\sqrt\mu\left(\beta-d_i^{-1}\right),\qquad
 v_i=\sqrt\mu d_i^{-1},\qquad u_i+v_i=\sqrt\mu\beta.
                                                               \label{eq:structural-path-nt-spectrum}
\end{equation}
If $t\leq\lambda_1/(4\sqrt n)$, then $\beta d_i\geq3$ for $i\geq1$:
\[
 \beta\lambda_1
 \geq\frac{\lambda_1/t}{1+(n-1)t/\lambda_1}>3.
\]
Hence
\begin{equation}
 \frac23\sqrt\mu\beta<u_i<\sqrt\mu\beta\qquad(i\geq1),
                                                               \label{eq:structural-path-flat}
\end{equation}
and every rank-$k<n-1$ approximation obeys
\begin{equation}
 \norm{U-U_k}_F\geq\frac23\sqrt\mu\beta\sqrt{n-1-k}.
                                                               \label{eq:structural-path-approx-rank}
\end{equation}
Using the path trace identities below to control the exceptional zero mode,
any fixed relative Frobenius target error below $2/3$ therefore requires
rank $\Omega(n)$ in the NT metric.  More generally, for any fixed target
error below one, shrinking the constant in
$t\leq\lambda_1/(4\sqrt n)$ makes all $n-1$ nonexceptional singular values
an arbitrarily large fixed fraction of $\sqrt\mu\beta$ and again forces
rank $\Omega(n)$.

The coordinate system is essential.  For the path,
\begin{equation}
 \lambda_k=4\sin^2\frac{\pi k}{2n},\quad
 \tr L^+=\frac{n^2-1}{6},\quad
 \tr(L^+)^2=\frac{2n^4+5n^2-7}{180}.                    \label{eq:structural-path-traces}
\end{equation}
As $t\downarrow0$, the raw tangent spectrum tends to
$\{-\tr L^+,\lambda_1^{-1},\ldots,\lambda_{n-1}^{-1}\}$, whose stable rank
tends to $7/5$ and whose best rank-one relative Frobenius error tends to
$\sqrt{2/7}$.  At $t=\lambda_1/(2\sqrt n)$, by contrast, the $n-1$
nonzero-mode NT singular values are asymptotically flat.  This regime is
late and ill conditioned:
\begin{equation}
 t=\frac{\pi^2}{2n^{5/2}}(1+o(1)),\qquad
 \cond(L+tI)=\frac8{\pi^2}n^{5/2}(1+o(1)).               \label{eq:structural-path-late}
\end{equation}
For smaller $t$ only the lower bound $\Omega(n^{5/2})$ is uniform.

\subsubsection{Why the barrier is self-limiting}
\label{sec:structural-rank-escapes}

The rank theorem obstructs low-rank \emph{factor tomography} of the
one-block NT direction.  It does not prove an $\Omega(n^2)$ storage bound,
and the path witness makes the escapes explicit.

First, chordal conversion replaces the single block by edge blocks
\begin{equation}
 Y_i=\begin{pmatrix}x_{ii}&x_{i,i+1}\\x_{i,i+1}&x_{i+1,i+1}\end{pmatrix}
 \succeq0\qquad(1\leq i<n),                              \label{eq:structural-path-blocks}
\end{equation}
with overlap equations $(Y_i)_{22}=(Y_{i+1})_{11}$ and trace equation
\begin{equation}
 (Y_1)_{11}+\sum_{i=1}^{n-1}(Y_i)_{22}=1.                \label{eq:structural-path-local-trace}
\end{equation}
The objective is the sum of inner products with
$\left[\begin{smallmatrix}1&-1\\-1&1\end{smallmatrix}\right]$.
Restriction of a global PSD matrix is feasible in this formulation.
Conversely, consistent PSD edge matrices have a PSD completion because the
path is chordal: construct Gram vectors successively, adding one orthogonal
coordinate per new vertex.  Objective and trace use only specified entries,
so the reformulation is exact.  Its barrier Hessian has constant-size blocks;
elimination gives a tridiagonal normal matrix with one dense border.  A
tridiagonal factorization and one scalar Schur complement use $O(n)$ memory
and arithmetic per Newton solve.  The conservative short-step total is
$O(n^{3/2}\log(n/\epsilon))$ arithmetic and $O(n)$ memory.  This is a
concrete instance of the formulation dependence emphasized in the chordal
conversion literature \cite{zhang2024-complexity-of-chordal-conversion-for}.

Second, the tangent itself has the implicit resolvent representation
\eqref{eq:structural-path-tangent}.  More generally, accurate low-rank SDP
methods may use low-rank endpoint structure to precondition or reformulate a
dense direction rather than factorizing that direction
\cite{bellavia2021-a-relaxed-interior-point-method,
zhang2024-well-conditioned-primal-dual-interior}.  These are compatible with
\cref{thm:structural-rank-progress}.

Third, high approximate matrix rank does not imply a hard amplitude state.
Writing $a=\sqrt\mu\beta$, the path estimates give
\begin{equation}
 \norm{U/a-(I-uu^T)}_{\rm op}=O(n^{-1/2}),\qquad
 \norm{U/a-(I-uu^T)}_F=O(1).                             \label{eq:structural-path-state-close}
\end{equation}
After normalization, $\operatorname{vec}(U)$ is $O(n^{-1/2})$-close to
\begin{equation}
 \frac{\ket I-\ket{uu^T}}{\sqrt{n-1}}.                  \label{eq:structural-path-lcu-state}
\end{equation}
For $u=\mathbf1/\sqrt n$, $\ket I/\sqrt n$ is a maximally entangled index
state and $\ket{uu^T}$ is a product of uniform states.  A two-term LCU
prepares \eqref{eq:structural-path-lcu-state} with constant success
probability and $\widetilde O(\log n)$ gates.  The theorem therefore makes no
state-preparation lower-bound claim.

Boundary analyticity under strict complementarity is classical
\cite{halicka2002-analyticity-of-the-central-path}.  The rank law is a sharp
local obstruction, but cone reformulation, implicit resolvents, and compact
state-preparation circuits are first-class escapes rather than footnotes.

\subsection{Sparse support and sparse-system embeddings}
\label{sec:structural-sparse-results}

The next statements concern standard-form complementarity, not arbitrary
optimization algorithms.  At an exact LP center $XSe=\mu e$, let
$0\leq\sigma<1$.  An inexact centering direction satisfies
\begin{equation}
 S\Delta x+X\Delta s=-(1-\sigma)\mu e+r,\qquad
 \norm r_2\leq\eta\mu.                                   \label{eq:structural-inexact-centering}
\end{equation}
Let $\operatorname{psupp}(\Delta x,\Delta s)$ be the indices where at least
one member of the coordinate pair is nonzero.

\begin{theorem}[Complementarity support barrier]
\label{thm:structural-support-barrier}
For $0\leq\sigma<1$, every direction satisfying
\eqref{eq:structural-inexact-centering} obeys
\begin{equation}
 \boxed{|\operatorname{psupp}(\Delta x,\Delta s)|
 \geq n-\frac{\eta^2}{(1-\sigma)^2}.}                    \label{eq:structural-support-count}
\end{equation}
For a short step $1-\sigma=a/\sqrt n$ with $\eta<a$, this is
$[1-(\eta/a)^2]n$.  If a $k$-sparse latent vector $z$ is reconstructed as
$(\Delta x,\Delta s)=Bz$, and each column of $B$ touches at most $g$
complementarity pairs, then
\begin{equation}
 \boxed{gk\geq n-\frac{\eta^2}{(1-\sigma)^2}.}            \label{eq:structural-locality-product}
\end{equation}
\end{theorem}

\begin{proof}
Let $Z$ contain the indices at which both update coordinates vanish.  For
each $i\in Z$, \eqref{eq:structural-inexact-centering} forces
$r_i=(1-\sigma)\mu$.  Hence
\[
 |Z|(1-\sigma)^2\mu^2\leq\norm r_2^2\leq\eta^2\mu^2,
\]
which proves \eqref{eq:structural-support-count}.  The support of $Bz$ is
contained in the union of at most $k$ column pair-supports, of total size at
most $gk$, proving \eqref{eq:structural-locality-product}.
\end{proof}

The count is realized by a sparse, incompressible direction.  For $n\geq3$,
consider
\begin{equation}
 \min x_2\quad\text{subject to}\quad
 x_1+x_2=1,\quad x_{i-1}-x_i=0\ (3\leq i\leq n),\quad x\geq0.
                                                               \label{eq:structural-incompressible-lp}
\end{equation}
Its equality matrix has maximum row and column sparsity two (the endpoint
columns have only one nonzero), and its unique
optimum is $e_1$.  The feasible segment is
$x(t)=(1-t,t,\ldots,t)$, and, for $0<t<(n-1)/n$, the central parameter is
\begin{equation}
 \mu(t)=\frac{t(1-t)}{n-1-nt}.
\end{equation}
Indeed, eliminating the equalities gives the one-variable barrier
$t-\mu[\log(1-t)+(n-1)\log t]$; its stationarity equation is the displayed
formula.
At $t=1/2$, every nonzero central tangent is proportional to
$(-1,1,\ldots,1)$.  Its normalized best $k$-sparse approximation error is
exactly $\sqrt{1-k/n}$; a sparse optimum therefore does not make the central
direction compressible.

There is also no hidden geometric restriction on a positive direction
profile.  Given any $u\in\R_{++}^n$ with $\norm u_2=1$, choose $\mu_0>0$ so
that $u_i^{-2}>4\mu_0$ and set
\begin{equation}
 c_1=-\sqrt{u_1^{-2}-4\mu_0},\qquad
 c_i=\sqrt{u_i^{-2}-4\mu_0}\quad(i\geq2).                \label{eq:structural-universal-cost}
\end{equation}
For the bound-constrained QP
\begin{equation}
 \min_{x\geq0}\ \tfrac12\norm x_2^2+c^Tx,               \label{eq:structural-universal-qp}
\end{equation}
the unique optimum is the one-sparse vector $(-c_1)e_1$, with strict
complementarity.  Central stationarity gives $s=x+c$ and
\begin{equation}
 x_i(\mu)=\frac{-c_i+\sqrt{c_i^2+4\mu}}2,\qquad
 x_i'(\mu)=\frac1{\sqrt{c_i^2+4\mu}},
\end{equation}
so $x'(\mu_0)=u$.  The row-scaled complementarity Jacobian is
$\diag(2x(\mu_0)+c)=\diag(1/u_i)$, which is one-sparse and constant
conditioned for near-uniform $u$.  Combining this realization with coherent
tomography gives the established
$\widetilde\Omega(n/\epsilon)$ solution-unitary readout bound on the
near-uniform hard family \cite{joran2023-quantum-tomography-state-preparation-unitaries}.
It is a modular readout bound: an algorithm allowed to inspect the explicit
coefficients \eqref{eq:structural-universal-cost} can learn $u$ there.

The condition-one box-LP family behind the older version of this observation
is superseded in strength by the original-primal and bounded-data families
of \cref{sec:lp-prefix-rigidified,sec:lp-bounded-data}.  Its remaining lesson
is only that perfect reduced conditioning does not remove coefficient-query
or classical-output cost.  Likewise, the older coordinate-reset ledger is
superseded by the projective path theorem of \cref{sec:positive-lazy}; the
QRAM-write limitation is stated once in
\cref{sec:structural-geometric-locality}.

\subsubsection{A query-preserving OSS progress step}

The support theorem says what a physical update must touch.  A distinct
question is whether standard sparse-QLS lower bounds occur inside a valid
QIPM Newton solve.  They do.  Let $H\in\R^{n\times n}$ be real symmetric,
invertible, with $\norm H_2=1$ and bidirectional row/column-slot sparse
access.  Let $d$ be the declared slot bound; declared candidate slots may
contain zeros.  Let $\norm q_2=1$.  Fix $\theta\in(0,1)$, choose
\begin{equation}
 0<\delta\leq\min\{\theta,1/2\},\quad
 0\leq\alpha\leq\delta/4,\quad
 a=\delta q/\sqrt2,\quad D=\diag(a),\quad
 \sigma=1-\alpha/\sqrt n.                                \label{eq:structural-oss-parameters}
\end{equation}
Consider
\begin{equation}
 \min c^Tx\quad\text{subject to}\quad
 Ax=0,\ x\geq0,\qquad
 A=\begin{bmatrix}H&-H\end{bmatrix},\quad c=(e+a,e-a).    \label{eq:structural-oss-lp}
\end{equation}
At $x^0=e_{2n},y^0=0,s^0=c$, the complementarity average is one and the
$N_2$ residual norm is $\delta$, while the $N_\infty$ residual is at most
$\delta/\sqrt2$.  Thus the point belongs to both radius-$\theta$
neighborhoods.

\begin{theorem}[Feasible OSS embedding]
\label{thm:structural-oss-embedding}
With null-space basis $V=(I,I)^T$ and $r=1-\sigma$, the feasible OSS step is
\begin{equation}
 \underbrace{\begin{bmatrix}-H&I+D\\H&I-D\end{bmatrix}}_{M_{\rm OSS}}
 \binom{\Delta y}{\lambda}
 =\binom{-re-a}{-re+a}.                                  \label{eq:structural-oss-system}
\end{equation}
It has row sparsity at most $d+1$, column sparsity at most $2d$, and
$\cond(M_{\rm OSS})=\Theta_\delta(\cond(H))$.  Its exact solution is
\begin{equation}
 \lambda=-re,\qquad \Delta y=\sigma H^{-1}a.             \label{eq:structural-oss-solution}
\end{equation}
The normalized solution has probability at least $2/3$ in the $\Delta y$
block.  The physical step lands at
\begin{equation}
 x^1=\sigma e,\qquad s^1=(e+ra,e-ra),\qquad \mu^1=\sigma,
                                                               \label{eq:structural-oss-endpoint}
\end{equation}
with centrality-residual norm $\sigma r\delta$.  Every sparse query to the
LP or OSS data is simulated with $O(1)$ queries to $H$ and evaluations of the
known $q$.
\end{theorem}

\begin{proof}
Primal and dual feasibility at the start are immediate, and
$\norm{X^0s^0-e}_2=\norm{(a,-a)}_2=\delta$.  Substitution of
$\Delta x=V\lambda$ and $\Delta s=-A^T\Delta y$ into the centering equation
gives \eqref{eq:structural-oss-system}.  The orthogonal row transform
\[
 Q=2^{-1/2}\begin{bmatrix}-I&I\\I&I\end{bmatrix}
\]
gives
\begin{equation}
 QM_{\rm OSS}=\sqrt2\begin{bmatrix}H&-D\\0&I\end{bmatrix}.
\end{equation}
The triangular factor has constant norm and inverse
$\left[\begin{smallmatrix}H^{-1}&H^{-1}D\\0&I\end{smallmatrix}\right]$;
since $\norm D\leq\delta/\sqrt2$, its inverse norm is
$\Theta_\delta(\norm{H^{-1}})$.

Applying $Q$ to the right-hand side gives
$H\Delta y-D\lambda=a$ and $\lambda=-re$.  Since $De=a$, this proves
\eqref{eq:structural-oss-solution}.  Also
$\norm{\Delta y}_2^2\geq\sigma^2\delta^2/2$, whereas
$\norm\lambda_2^2=\alpha^2$, so the block probability is at least
\[
 \frac{\sigma^2\delta^2/2}{\sigma^2\delta^2/2+\alpha^2}>\frac23.
\]
Direct substitution proves \eqref{eq:structural-oss-endpoint} and gives
$X^1s^1-\mu^1e=\sigma r(a,-a)$.  All remaining entries and sparse positions
are fixed, proving the oracle simulation.
\end{proof}

After symmetric dilation, the theorem transfers the
$\Omega(\kappa\sqrt d)$ fixed-error and
$\Omega(\kappa\log(1/\epsilon))$ constant-sparsity QLS lower bounds of
Mori et al.\ to a strictly feasible short progress step
\cite{mori2026-sparsity-dependent-complexity-lower-bound}.  It is an
embedding of known QLS restrictions, not a new QLS technique or an
end-to-end LP lower bound.  Indeed, the physical $(\Delta x,\Delta s)$ above
is explicit; the hard block is $\Delta y$, and algorithms that bypass
accurate OSS-state preparation are outside the statement.
Here $d$ is the declared sparse-oracle slot bound, which may count padded
zero slots.  A symmetric dilation of a nonsymmetric source matrix requires
bidirectional row- and column-slot access.  Mori et al.'s clock family has
known successor/predecessor slot locations and reversible hidden updates, so
both directions are simulated with constant query overhead.

\subsubsection{What bounded treewidth preserves}

The same precision restriction survives at treewidth two.  Mori et al.'s
hard matrix has the form
\begin{equation}
 A=I-\rho B,                                               \label{eq:structural-clock-matrix}
\end{equation}
where $B$ is a clock permutation and every row and column has one identity
and one permutation position.  Form its symmetric dilation
\begin{equation}
 K_A=\begin{bmatrix}0&A^T\\A&0\end{bmatrix},\qquad
 \widehat b=\binom0b.                                    \label{eq:structural-treewidth-dilation}
\end{equation}

\begin{theorem}[Treewidth-two KKT precision]
\label{thm:structural-treewidth-two}
There are real symmetric, indefinite, two-sparse equality-KKT matrices $K_A$
whose support graphs have treewidth at most two and for which preparing an
$\epsilon$-approximation to the normalized solution requires
\begin{equation}
 \Omega\!\left(\cond(K_A)\log(1/\epsilon)\right)          \label{eq:structural-treewidth-qls}
\end{equation}
queries in the parameter range of the underlying QLS theorem.
\end{theorem}

\begin{proof}
One has
$K_A^{-1}\widehat b=(A^{-1}b,0)^T$, and the singular values of $K_A$ are
those of $A$, each repeated.  The identity positions of $A$ give one perfect
matching in the bipartite support graph of $K_A$, while the clock-permutation
positions give a second.  Their union is a disjoint union of paths and even
cycles, so its treewidth is at most two.  Successor and predecessor clock
queries simulate every sparse query with constant overhead.  The state and
condition number are preserved, so the QLS precision lower bound transfers.
Finally, $K_A$ is the equality-KKT matrix of $\min0$ subject to $Ax=b$.
\end{proof}

The two-sparse clock dilation of Wang--Zhang has the same low-treewidth
support phenomenon but proves a different query-depth statement and not the
$\log(1/\epsilon)$ precision factor
\cite{wang2024-tight-quantum-depth-lower-bound}.
This does not say that bounded-treewidth optimization is generically hard.
The stronger classical fact points the other way: supplied width-$\tau$
elimination gives $O(N\tau^2)$ scalar Cholesky work and $O(N\tau)$ storage
for an $N$-dimensional SPD Newton system.  A dense classical output already
costs $\Omega(N)$, so the dimension-only full-output speedup ceiling is
$O(\tau^2)$.  Nearly linear bounded-treewidth IPMs sharpen this baseline
\cite{dong2020-a-nearly-linear-time-algorithm,gu2022-a-faster-small-treewidth-sdp}.

There is nevertheless scalar LP value hardness at an exact chosen width.
For Boolean matrices $Z^i\in\{0,1\}^{c\times b}$, start from the pre-dual
program
\begin{equation}
 \begin{aligned}
  \max_{w,v^1,\ldots,v^a}\quad&\sum_{k=1}^c w_k\\
  \text{subject to}\quad&
  w-\sum_{i=1}^a Z^iv^i\leq0,\qquad
  e^Tv^i\leq1\ (1\leq i\leq a),\\
  &w,v^i\geq0.
 \end{aligned}                                            \label{eq:structural-apers-predual}
\end{equation}
Its value is
\begin{equation}
 \sum_{i=1}^a\max_{j\in[b]}\sum_{k=1}^c Z_{ijk}.          \label{eq:structural-apers-value}
\end{equation}
Add a nonnegative, zero-objective $z_i$ to every top inequality and bottom
inequality $i$, always with coefficient $+1$.

\begin{theorem}[LP value hardness at exact treewidth]
\label{thm:structural-exact-treewidth}
For integers $a,b,c\geq1$, the augmented LP has
$p=c+a(b+1)$ variables, $q=c+a$ constraints, and
constraint-intersection treewidth exactly $\tau=c$.  When $c\leq a$,
constant-additive approximation of its scalar value requires
\begin{equation}
 \Omega(a\sqrt b\,c)=\Omega(\tau\sqrt{pq})                \label{eq:structural-exact-treewidth-value}
\end{equation}
fixed-position quantum coefficient queries.  The corresponding randomized
bound is $\Omega(abc)=\Omega(p\tau)$.
\end{theorem}

\begin{proof}
For fixed $v^i$, each top constraint is tight at an optimum, so summing the
coordinates of $w$ separates the $a$ simplex problems in $v^i$ and gives
\eqref{eq:structural-apers-value}.  Setting $z=0$ extends every original
feasible point.  Conversely, deleting positive $z$ from an augmented point
only relaxes its inequalities, so the value is unchanged.

Let $C$ be the $c$ top-constraint vertices.  The column of $z_i$ makes
$C\cup\{i\}$ a clique, and no variable column contains two bottom vertices.
Thus the intersection graph is exactly $K_c\vee I_a$, with bags
$C\cup\{i\}$, and has treewidth $c$.  The dimensions follow by counting
$w,v,z$.  The published Boolean reduction for
\eqref{eq:structural-apers-value} gives
$\Omega(a\sqrt b\,c)$ quantum and $\Omega(abc)$ randomized fixed-position
queries.  Since
\[
 a\sqrt b\,c=\tau\sqrt{(q-\tau)(p-q)},\qquad
 abc=\tau(p-q),
\]
and $c\leq a$ makes $q-\tau=\Theta(q)$ and $p-q=\Theta(p)$, the displayed
bounds follow.
\end{proof}

This is a reparameterized structural corollary of an existing reduction
\cite{apers2026-quantum-speedups-linear-programming-interior}, not a new
adversary argument.  It applies to individual coefficient queries; a whole
column reveals $\Theta(\tau)$ coefficients, and a whole top row can reveal
far more.  The broad claim that treewidth was absent from QLS lower-bound
work would also be false in view of Khanpour--Talkington
\cite{khanpour2026-proving-limits-power-flow}.  The new restriction here is
only the treewidth-two KKT precision embedding.

\subsection{Geometric locality of central-point access}
\label{sec:structural-geometric-locality}

Sparse-query access suppresses where the input is stored.  On bounded-range
hardware, that geometry can itself be an obstruction.  Let hidden signs
$\varsigma_1,\ldots,\varsigma_N\in\{\pm1\}$ and introduce
$d_i=u_i-v_i$ with $u_i,v_i\geq0$.  For fixed $\alpha\in(0,1)$, consider
\begin{equation}
 \begin{aligned}
  \min\quad&\sum_{i=0}^N(u_i+v_i)+\alpha(u_N-v_N),\\
  \text{subject to}\quad&d_0=1,\\
  &d_i-\varsigma_i d_{i-1}=0\quad(1\leq i\leq N).
 \end{aligned}                                             \label{eq:structural-locality-lp}
\end{equation}
Every row has at most four nonzeros, every column at most two, and the
constraint-intersection graph is a path.  Put
$p_i=\prod_{j=1}^i\varsigma_j$ and $q_i=u_i+v_i$.  Feasibility gives
$d_i=p_i$.  Taking every $q_i=2$ is strictly primal feasible; with the
convention $A^Ty+s=c$, the choice $y=0$ gives $s=c>0$ and strict dual
feasibility.  The barrier restricted to the feasible affine space is
\begin{equation}
 \sum_{i=0}^Nq_i+\alpha p_N
 -\mu\sum_{i=0}^N\log\frac{q_i^2-1}{4}.                  \label{eq:structural-locality-barrier}
\end{equation}
Thus every central coordinate equals
\begin{equation}
 q(\mu)=\mu+\sqrt{\mu^2+1}.                              \label{eq:structural-locality-center}
\end{equation}
At $\mu=3/4$, $q=2$ and
$(u_i,v_i)=(3/2,1/2)$ or $(1/2,3/2)$ according to $p_i$.  Reading the final
cell with coordinate error below $1/2$ therefore reveals parity.

This information cost is invisible to reduced conditioning.  The public
vectors $V_i=(e_{u_i}+e_{v_i})/\sqrt2$ form an orthonormal null-space basis,
and for the full logarithmic-barrier Hessian $H_x=X^{-2}$,
\begin{equation}
 V^TH_xV
 =2\bigl[(q+1)^{-2}+(q-1)^{-2}\bigr]I.                   \label{eq:structural-locality-condition-one}
\end{equation}
It has condition number one for every $\mu>0$.

Place input register $\varsigma_i$ at a site $v_i$ of a hardware graph
$\Gamma$.  A depth-$T$, range-$r$ circuit has gates of hardware diameter at
most $r$ in each layer.  Its fixed initial workspace may contain arbitrary
input-independent entanglement.  A measurement at site $o$ is locally
parity-readable if it returns $p_N$ with probability at least $2/3$ for every
input.

\begin{theorem}[Central-oracle light cone]
\label{thm:structural-light-cone}
Let $R_\Gamma(o;v)=\max_{1\leq i\leq N}\operatorname{dist}_\Gamma(o,v_i)$
and let $b=\max_{w\in\Gamma}|\{i:v_i=w\}|$ be the maximum site occupancy.
Every such circuit constructing a locally readable final central cell obeys
\begin{equation}
 \boxed{T\geq\left\lceil\frac{R_\Gamma(o;v)}{r}\right\rceil.}
                                                          \label{eq:structural-light-cone}
\end{equation}
Consequently, when at most $b$ input registers share a site, a path
requires $T\geq(\lceil N/b\rceil-1)/(2r)$ for every choice of output site,
which is $\Omega(N/(br))$ once $N\ge2b$, and a $d$-dimensional
nearest-neighbor grid requires $(2rT+1)^d\geq\lceil N/b\rceil$, hence
$T\ge((N/b)^{1/d}-1)/(2r)$, which is $\Omega((N/b)^{1/d}/r)$ once
$N\ge2^db$.
Pairwise-distinct placement is the case $b=1$.
\end{theorem}

\begin{proof}
Pull a local output observable backwards through the circuit.  Its support
can grow by hardware distance at most $r$ per layer, so after $T$ layers it
is contained in $B_\Gamma(o,rT)$.  If some $v_j$ lies outside this ball, two
inputs differing only at $j$ have identical reduced input states on the
backward light cone.  Their workspace state is also identical, including
all of its input-independent entanglement.  The local output distributions
are therefore identical, although flipping $\varsigma_j$ flips $p_N$; one
distribution cannot return both answers with probability $2/3$.

Thus the ball $B_\Gamma(o,rT)$ contains every occupied site, proving
\eqref{eq:structural-light-cone}.  Under occupancy at most $b$, at least
$\lceil N/b\rceil$ distinct sites are occupied.  On a path a radius-$rT$
ball has at most $2rT+1$ sites, and on the grid at most $(2rT+1)^d$, which
proves the two consequences.
\end{proof}

This is the standard backward-light-cone argument
\cite{barak2022-classical-algorithms-and-quantum-limitations}; the claimed
contribution is only its conjunction with the strictly feasible,
condition-one central-point witness.  It is not a lower bound in the
all-to-all or free-QRAM query model, for a homogeneous self-dual embedding
with a universal start, for an answer left in a global observable whose
measurement is declared free, or when unboundedly many input registers may
share one site.  The same lower bound holds classically, so it
is an obstruction to a claimed quantum depth advantage, not a
quantum--classical separation.  Input-dependent pre-entanglement can already
contain the answer and merely moves its construction cost into advice.

A stronger per-iteration QRAM-write claim does not follow.  A lazy oracle may
store the input once and evaluate a changed weight from the queried index and
a scalar path parameter.  An $\Omega(N)$ refresh theorem would need a
materialized-array or cell-probe model, or a lower bound for the composed
evaluation oracle.  Numerical change in all coordinates is not enough.

\subsection{Synthesis}
\label{sec:structural-synthesis}

These tools charge different resources and should not be collapsed into one
informal claim that sparsity fails.  Symmetry charges invariant virtual work
even after equality rows are mixed.  Central mixtures trade sparse
one-bit-local input incidence against operator spread.  Complementarity
progress charges pair
support or NT aggregate rank.  Sparse-QLS hardness survives a legitimate OSS
step and even a treewidth-two KKT representation, while exact treewidth also
admits strong classical algorithms.  Physical locality charges communication
only when the input and output are locally materialized with bounded per-site
occupancy.

Their escape routes are equally structural.  A nonlinear output may not be
mixture stable; central matrices may have a large common-basis spectral
spread; a nonlocal reconstruction map may keep a latent state sparse; chordal
conversion or an implicit resolvent may avoid matrix-factor tomography; an
algorithm may bypass the OSS; and a global or lazy oracle may avoid local
materialization.  The results above identify where such an escape must enter.
They do not supply a universal quantum query or runtime lower bound.

\section{Corrections to Prior Work and Recorded Dead Ends}
\label{sec:corrections}

Corrections and negative results delimit what the positive theorems actually
establish.  The first result below corrects a norm estimate in a publicly
posted QCPM complexity derivation.  The next two results isolate representation and
initialization costs on sparse Hamiltonian and central-path constructions.
We then record false shortcuts that were tested and rejected.  None of these
statements is a general verdict against Hamiltonian optimization,
preconditioning, or state reuse.

\subsection{Correction to the QCPM simulation norm}
\label{sec:corrections-qcpm}

We consider arXiv:2311.03977v2 (16 October 2024), the quantum central-path
method (QCPM) of Augustino, Leng, Nannicini, Terlaky, and Wu
\cite{augustino2024-quantum-central-path-algorithm-linear}.  All formulas in
this subsection refer to that version.  Let $d=m+n+2$ be the dimension of its
self-dual embedding, and use its notation
\begin{equation}
 s(z)=Mz+q,\qquad F(z)=z\odot s(z),\qquad
 f_\mu(z)=\frac12\norm{F(z)-\mu e}_2^2.                 \label{eq:corr-qcpm-potential}
\end{equation}
The central-path Hamiltonian is
$H(\mu)=-(h^2/2)\nabla^2+f_\mu$, and Algorithm~1 sets
% QCPM v2 has an internal normalization inconsistency between Algorithm 1's
% h(t) and equation (23); the definition of a here deliberately follows
% Algorithm 1, as do all constants below.
\begin{equation}
 h(t)=a\mu(t)^2,
 \qquad a=\frac{\gamma^2}{2R_1C_{d,\delta}},            \label{eq:corr-qcpm-h}
\end{equation}
where $C_{d,\delta}=\sqrt d/2+(3/4)\log(2/\delta)$ is the
concentration factor displayed in Algorithm~1.  We assume $d\geq1$,
$a,\eta>0$, and $0<\varepsilon<1$.  The spatial domain $\Omega_D$ is
fixed and compact and contains $e$; statements involving $z(\mu)$ also
require that point to belong to $\Omega_D$.  The invoked
real-space simulation theorem measures a potential $V$ on a fixed domain
$\Omega_D$ by
\begin{equation}
 \norm V_{\infty,1}=\int\sup_{z\in\Omega_D}|V(z,\tau)|\,d\tau.             \label{eq:corr-qcpm-simulator-norm}
\end{equation}
The QCPM proof truncates to a fixed box containing the central-path
neighborhood and replaces the supremum in
\eqref{eq:corr-qcpm-simulator-norm} by an assumed bound
$\sup_{\Omega_D} f_\mu\leq K\gamma^2\mu^2$.  That bound is false.

\begin{proposition}[Global-supremum counterexample]
\label{prop:corr-qcpm-supremum}
If $\Omega_D$ contains the initial center $e=z(1)$, then, for every
$0<\mu\leq1$,
\begin{equation}
 \sup_{z\in\Omega_D}f_\mu(z)
 \geq f_\mu(e)=\frac d2(1-\mu)^2.                       \label{eq:corr-qcpm-counterexample}
\end{equation}
Consequently, at $\mu=\varepsilon\leq1/2$, the asserted bound would require
$K\geq d/(8\gamma^2\varepsilon^2)$.
\end{proposition}

\begin{proof}
The self-dual embedding is initialized so that $s(e)=e$.  Thus $F(e)=e$,
and \eqref{eq:corr-qcpm-potential} gives
\[
 f_\mu(e)=\tfrac12\norm{(1-\mu)e}_2^2
          =\tfrac d2(1-\mu)^2.
\]
Taking the supremum proves \eqref{eq:corr-qcpm-counterexample}.  When
$\varepsilon\leq1/2$, its right-hand side is at least $d/8$, which gives the
claim about $K$.
\end{proof}

Ground-state concentration controls state-weighted mass, not the operator
supremum on a fixed domain.  A scalar phase shift does not restore the
claimed scaling.  If $z(\mu)\in\Omega_D$, then $f_\mu(z(\mu))=0$, and
for every real scalar $c$ the triangle inequality gives
\[
 \sup_{z\in\Omega_D}|f_\mu(z)-c|
 \geq\frac12|f_\mu(e)-f_\mu(z(\mu))|
 =\frac d4(1-\mu)^2.
\]

There is a second, independent issue in the time dilation.  Let
$\mu:[0,1]\to[\varepsilon,1]$ be continuous.  After division
by $h(t)\mu(t)$, the evolution equation is
\begin{equation}
 i\eta\,\partial_t\Psi
 =\left[-\frac{\theta(t)}2\nabla^2
       +\frac{f_{\mu(t)}}{\mu(t)h(t)}\right]\Psi,
 \qquad \theta(t)=\frac{h(t)}{\mu(t)}=a\mu(t).           \label{eq:corr-qcpm-scaled-evolution}
\end{equation}
The linear substitution used in the QCPM proof treats $\theta$ as constant,
although it varies with $t$.  The clock that normalizes the kinetic term is
instead
\begin{equation}
 \tau(t)=\frac1\eta\int_0^t\theta(u)\,du
        =\frac a\eta\int_0^t\mu(u)\,du.                 \label{eq:corr-qcpm-clock}
\end{equation}
Indeed, $dt/d\tau=\eta/\theta(t)$, so
\eqref{eq:corr-qcpm-scaled-evolution} becomes
\begin{equation}
 i\partial_\tau\widetilde\Psi
 =\left[-\frac12\nabla^2
        +\frac{f_{\mu(t(\tau))}}{h(t(\tau))^2}\right]
   \widetilde\Psi.                                      \label{eq:corr-qcpm-corrected-evolution}
\end{equation}
Combining \eqref{eq:corr-qcpm-h}--\eqref{eq:corr-qcpm-clock} yields the
exact corrected identity
\begin{equation}
 \boxed{
 \Lambda:=\norm V_{\infty,1}
 =\frac1{a\eta}\int_0^1
   \frac{\sup_{z\in\Omega_D}f_{\mu(t)}(z)}{\mu(t)^3}\,dt.}
                                                               \label{eq:corr-qcpm-exact-norm}
\end{equation}
This identity does not use a lower-bound argument.  The clock correction by
itself works in the authors' favor: it lowers the norm weight from the
constant-$\theta$ argument's $\mu^{-4}$ to $\mu^{-3}$.  The
$\varepsilon^{-3}$ blow-up below is driven entirely by the false global
supremum bound.

The schedule in Algorithm~1 is
\begin{equation}
 g(t)=c_e^{-1}\int_0^t
 \exp\!\left[-\frac1{\tau(1-\tau)}\right]d\tau,
 \qquad
 \mu(t)=\varepsilon+(1-\varepsilon)(1-g(t)),            \label{eq:corr-qcpm-schedule}
\end{equation}
where $c_e$ normalizes $g(1)=1$.
The bump integrand is extended by zero at both endpoints.  It is
continuous, nonnegative, and positive on $(0,1)$, so $c_e>0$.

\begin{lemma}[Width of the flat endpoint]
\label{lem:corr-qcpm-endpoint}
For $0<\varepsilon\leq\min\{1/4,c_e\}$, put
$u_\varepsilon=1/(2\log(1/\varepsilon))$.  Then
$\mu(t)\leq2\varepsilon$ for every
$t\in[1-u_\varepsilon,1]$.
\end{lemma}

\begin{proof}
For $0<u<1$, substitute $v=1-\tau$ in the tail of
\eqref{eq:corr-qcpm-schedule}.  Since $(1-v)v\leq v$,
\begin{equation}
 1-g(1-u)
 \leq\frac1{c_e}\int_0^u e^{-1/v}\,dv
 \leq\frac{u}{c_e}e^{-1/u}.                             \label{eq:corr-qcpm-tail}
\end{equation}
At $u=u_\varepsilon$, the exponential equals $\varepsilon^2$.  The last
quantity in \eqref{eq:corr-qcpm-tail} is at most $\varepsilon$:
$\varepsilon\leq1/4$ implies $0<u_\varepsilon\leq1/2$, and
$u_\varepsilon\varepsilon\leq c_e$ follows from $\varepsilon\leq c_e$.
For $t\geq1-u_\varepsilon$,
monotonicity then gives
$\mu(t)\leq\varepsilon+(1-\varepsilon)\varepsilon\leq2\varepsilon$.
\end{proof}

\begin{theorem}[Corrected QCPM simulator-norm scale]
\label{thm:corr-qcpm-norm}
For the choices in Algorithm~1 and
$0<\varepsilon\leq\min\{1/4,c_e\}$,
\begin{equation}
 \boxed{
 \Lambda\geq
 \frac{d}{128a\eta\,\varepsilon^3\log(1/\varepsilon)}.} \label{eq:corr-qcpm-lower}
\end{equation}
If $\Omega_D=[0,D]^d$, a valid global upper certificate is
\begin{equation}
 \Lambda\leq\frac{F_D}{a\eta\varepsilon^3},
 \qquad
 F_D=\frac d2\left[
 D\bigl(D\norm M_\infty+\norm q_\infty\bigr)+1
 \right]^2.                                             \label{eq:corr-qcpm-upper}
\end{equation}
Here $\norm M_\infty$ is the maximum absolute row sum.
\end{theorem}

\begin{proof}
On the interval supplied by \cref{lem:corr-qcpm-endpoint}, take
$\varepsilon\leq1/4$.  Proposition~\ref{prop:corr-qcpm-supremum} gives
$\sup_{\Omega_D} f_{\mu(t)}\geq d/8$, while
$\mu(t)^{-3}\geq(2\varepsilon)^{-3}$.  The interval has length
$1/(2\log(1/\varepsilon))$.  Substitution in
\eqref{eq:corr-qcpm-exact-norm} proves \eqref{eq:corr-qcpm-lower}.

For the upper bound, every $z\in[0,D]^d$ satisfies
\[
 |s_i(z)|\leq D\norm M_\infty+\norm q_\infty.
\]
Hence $|F_i(z)-\mu|$ is at most
$D(D\norm M_\infty+\norm q_\infty)+1$ for $0<\mu\leq1$, and
$f_\mu(z)\leq F_D$.  Since $\mu(t)\geq\varepsilon$, the exact identity
\eqref{eq:corr-qcpm-exact-norm} gives \eqref{eq:corr-qcpm-upper}.
\end{proof}

Since $a^{-1}=2R_1C_{d,\delta}/\gamma^2$,
\cref{eq:corr-qcpm-lower} gives the simulator-norm lower bound
\begin{equation}
 \Omega\!\left(
  \frac{R_1dC_{d,\delta}}
       {\gamma^2\eta\varepsilon^3\log(1/\varepsilon)}
 \right).                                                \label{eq:corr-qcpm-complexity-scale}
\end{equation}
By contrast, the bound claimed in v2 is
\[
 \norm V_{\infty,1}\leq\frac{K\gamma^2}{\theta\eta},
 \qquad
 O\!\left(\frac{K\gamma^2}{a\eta\varepsilon}\right)
 =O\!\left(\frac{R_1C_{d,\delta}}{\eta\varepsilon}\right)
 \quad(K=O(1)).
\]
Thus $\gamma^2$ cancels against $1/a$.  Conditional on all other simulation
hypotheses, \eqref{eq:corr-qcpm-upper} supplies the
norm-based upper certificate
\[
 O\!\left(\frac{F_DR_1C_{d,\delta}}
 {\gamma^2\eta\varepsilon^3}\right)
\]
before polylogarithmic factors and potential-evaluation cost.

The calculation has a schedule-independent form.

\begin{corollary}[Norm--speed tradeoff]
\label{cor:corr-qcpm-tradeoff}
Let $\mu:[0,1]\to[\varepsilon,1]$ be $L$-Lipschitz, with
$\mu(0)=1$, $\mu(1)=\varepsilon$, and $0<L<\infty$.
In particular, a differentiable schedule with $|\dot\mu|\leq L$
satisfies this hypothesis by the mean value theorem.  Retain $h=a\mu^2$.  For
$\varepsilon\leq1/4$,
\begin{equation}
 \boxed{\Lambda\geq\frac{d}{64a\eta L\varepsilon^2}.}   \label{eq:corr-qcpm-tradeoff}
\end{equation}
\end{corollary}

\begin{proof}
The endpoint values imply $L\geq1-\varepsilon$, hence
$\varepsilon/L\leq1$.  On the terminal interval
$[1-\varepsilon/L,1]$, the Lipschitz bound gives
$\mu(t)\leq\varepsilon+L(1-t)\leq2\varepsilon$.  Throughout that interval,
\cref{prop:corr-qcpm-supremum} gives
$\sup_{\Omega_D} f_\mu\geq d/8$, and
$\mu^{-3}\geq(2\varepsilon)^{-3}$.  Equation
\eqref{eq:corr-qcpm-exact-norm} proves the result.
\end{proof}
No monotonicity assumption is needed for this bound.

Taking $L$ large weakens \eqref{eq:corr-qcpm-tradeoff}, but moves $L$ into
the derivative, smoothness, and adiabatic estimates.  The corollary alone
does not optimize that tradeoff, and it does not exclude a theorem that
rigorously exploits localization.

As a numerical check, adaptive quadrature gives
$c_e=0.007029858406609657$.  For
\[
 J(\varepsilon)=\varepsilon^3\log(1/\varepsilon)
 \int_0^1(1-\mu(t))^2\mu(t)^{-3}\,dt,
\]
the computed values are
\begin{center}
\begin{tabular}{c@{\qquad}c}
\toprule
$\varepsilon$ & $J(\varepsilon)$ \\
\midrule
$10^{-2}$ & $0.77120646$ \\
$10^{-3}$ & $0.89055595$ \\
$10^{-4}$ & $0.95036732$ \\
\bottomrule
\end{tabular}
\end{center}
The sampled values are consistent with a limit of one; no limit is
deduced from this computation.  Retaining only the witness
$f_\mu(e)=d(1-\mu)^2/2$ gives
\[
 \Lambda\geq\Lambda_e
 :=\frac{d}{2a\eta}\int_0^1
 (1-\mu(t))^2\mu(t)^{-3}\,dt
 \geq\frac{d}{128a\eta\varepsilon^3\log(1/\varepsilon)}
 \quad\bigl(0<\varepsilon\leq\min\{1/4,c_e\}\bigr).
\]
The proof certifies $J(\varepsilon)\geq1/64$ in this range and makes no
claim that this constant is sharp.  The quadrature is a sanity check on the proved scale, not
an estimate of the full supremum or part of the proof.

\paragraph{Formal verification of the norm correction.}
The companion Lean development in \path{formal/QipmFormal/QCPM/}
checks the finite-dimensional potential and initialization identity,
the compact-domain supremum bound, robustness under scalar shifts,
the corrected clock and its inverse, and the integral change of variables.
It also checks the actual normalized bump schedule, the explicit
endpoint threshold, both lower-bound constants, and the fixed-box upper
certificate.  The headline bounds use the spatial supremum of the stated
potential; their continuity and integrability are proved rather than
assumed as extra properties of that supremum.  The initialization row
equations $Me+q=e$ and inclusion of the required points in the fixed domain
are explicit hypotheses.  The source correspondence is to Algorithm~1
of v2, whose normalization differs from its equation~(23).
The verification does not formalize Schr\"odinger evolution, the external
simulation or adiabatic theorems, or an unconditional query lower bound.

The norm and clock errors invalidate the stated complexity derivation, not
the QCPM idea itself, and \eqref{eq:corr-qcpm-complexity-scale} is not an
unconditional query lower bound for every implementation of the evolution.
Chakrabarti et al.\ analyze missing error controls in the pseudo-spectral
real-space simulation
\cite{chakrabarti2025-speedups-quantum-dynamics}.  Gribling et al.\ cite
that analysis and separately identify the absence of a suitable directly
applicable infinite-dimensional adiabatic theorem
\cite{gribling2025-self-concordant-schrodinger-operators-spectral}.  We do
not claim those broader criticisms as new.  Even after correcting the norm,
the following obligations remain:
\begin{itemize}[leftmargin=*]
 \item The schedule in \eqref{eq:corr-qcpm-schedule} is nonconstant and
 $C^\infty$-flat, but it is not analytic through either endpoint.  An
 analytic function with its zero endpoint Taylor series would be locally
 constant.  The quoted exponential adiabatic theorem assumes analytic
 continuation through the closed interval.
 \item The adiabatic theorem must be applied to the generator in
 \eqref{eq:corr-qcpm-scaled-evolution}; the derivative discussion analyzes a
 differently scaled Hamiltonian and does not establish uniform derivative
 constants for time-dependent $h$.
 \item The exact potential \eqref{eq:corr-qcpm-potential} is generally
 quartic.  The displayed Gaussian state and oscillator gap concern its
 quadratic approximation.  A dimension-, $\mu$-, and input-uniform transfer
 to the exact ground state and spectral gap is not established.
 \item The proposition and algorithm use different choices of $h$; the
 proposition's choice has no $\gamma$ or $\delta$ dependence and therefore
 cannot by itself provide arbitrary neighborhood opening and failure
 probability.
 \item The initial exact ground state is assumed rather than prepared.
 Knowing its minimizer $e$ does not prepare the ground state of a quartic
 Hamiltonian.
 \item The real-space theorem assumes a smooth periodic potential and
 periodic boundary conditions.  Restriction of
 \eqref{eq:corr-qcpm-potential} to a box neither makes it periodic nor
 controls the boundary and truncation errors.
 \item The box size must retain its solution-radius dependence unless an
 input-size promise bounds the central path.
\end{itemize}
A localized or clipped-potential repair would need leakage, spurious-state,
uniform-gap, boundary, and comparison estimates; concentration alone does
not supply them.

\subsection{Sparse metric, dense inverse metric}
\label{sec:corrections-densification}

The proposal of Gribling, Apers, Nieuwboer, and Walter replaces the Euclidean
Laplacian by the Laplace--Beltrami operator of the Hessian metric
\cite{gribling2025-self-concordant-schrodinger-operators-spectral}.  If
$G(x)=\nabla^2\phi(x)$, its principal part is
\begin{equation}
 L\psi=\frac1{\sqrt{\det G}}
 \sum_{i,j}\partial_i\!\left(
 (G^{-1})_{ij}\sqrt{\det G}\,\partial_j\psi
 \right).                                                \label{eq:corr-laplace-beltrami}
\end{equation}
Sparse constraints do not make these standard-coordinate coefficients
entrywise sparse.

\begin{theorem}[Path-simplex inverse metric]
\label{thm:corr-path-simplex}
Let
\[
 \mathcal D_n=\{x\in\R^n:0<x_1<\cdots<x_n<1\}
\]
and use the logarithmic barrier
\begin{equation}
 \phi(x)=-\log x_1-\sum_{i=1}^{n-1}\log(x_{i+1}-x_i)
          -\log(1-x_n).                                  \label{eq:corr-path-barrier}
\end{equation}
Every inequality contains at most two variables, and the
constraint-intersection graph is a path.  At the analytic center
$\bar x_i=i/(n+1)$, the Hessian $G=\nabla^2\phi(\bar x)$ is tridiagonal, but
\begin{equation}
 (G^{-1})_{ij}
 =\frac{\min(i,j)(n+1-\max(i,j))}{(n+1)^3}>0             \label{eq:corr-path-green}
\end{equation}
for every $i,j$.  For
$i,j\in[\lceil(n+1)/4\rceil,\lfloor3(n+1)/4\rfloor]$,
these entries are $\Theta(1/n)$, so $\Theta(n^2)$ coefficients are of
comparable order.
\end{theorem}

\begin{proof}
Write the $n+1$ slacks as
$s_0=x_1$, $s_i=x_{i+1}-x_i$ for $1\leq i<n$, and
$s_n=1-x_n$.  At $\bar x$, every slack is $1/(n+1)$.  If $B$ is the signed
path-incidence Jacobian of this affine slack map, then
\[
 G=B^T\diag(s^{-2})B=(n+1)^2T_n,
\]
where $T_n$ has diagonal entries two and adjacent off-diagonal entries
minus one.  The discrete Green identity is
\[
 (T_n^{-1})_{ij}
 =\frac{\min(i,j)(n+1-\max(i,j))}{n+1}.
\]
Scaling gives \eqref{eq:corr-path-green}.  In the middle half, both factors
in its numerator are bounded above by $n$ and below by a fixed positive
multiple of $n$, proving the final assertion.
\end{proof}

Thus a direct standard-coordinate, nondivergence-form finite-difference
stencil for the principal symbol in \eqref{eq:corr-laplace-beltrami} has
$\Omega(n^2)$ mixed-derivative directions at the analytic-center grid point,
although the input has only $O(n)$ nonzeros.  This is a densification warning,
not a simulation or query lower bound.  Indeed, at an arbitrary interior
point put $r_k=s_k^2$, $R=\sum_{k=0}^nr_k$, and let $i\leq j$.  Weighted path
inversion gives
\begin{equation}
 (G(x)^{-1})_{ij}
 =\frac{(\sum_{k=0}^{i-1}r_k)(\sum_{k=j}^{n}r_k)}{R}.     \label{eq:corr-path-semiseparable}
\end{equation}
Every strict off-diagonal block is rank one, and prefix sums give an $O(n)$
matrix--vector apply.  The witness therefore rules out only the inference
from sparse Hessian support to an entrywise sparse inverse-metric tensor.
It does not rule out implicit, transform-based, or coordinate-adapted
simulation.

Abe and Nagai parameterize their Riemannian Hamiltonian simulation by the
kinetic-matrix sparsity and separately assume sparse-access oracles for that
matrix; they note that constructing those black-box oracles is generally
nontrivial \cite{abe2026-quantum-riemannian-hamiltonian-descent}.
Theorem~\ref{thm:corr-path-simplex} shows only that input-level $O(n)$
sparsity is lost: the direct stencil has $\Theta(n^2)$ terms.  On a tensor
grid with $q$ points per coordinate it still has at most $2n^2+1$ neighbors
per row, sparse relative to the $q^n$-dimensional grid.  Thus the theorem
does not rule out sparse-Hamiltonian access or refute the separately assumed
Abe--Nagai oracles.  Gribling et al.\ likewise leave the simulation and
qubit-realization questions open.

\subsection{A hidden-sign start-oracle separation}
\label{sec:corrections-hidden-sign}

We reuse, without reconstructing it, the path-sparse hidden-sign LP
\eqref{eq:structural-locality-lp} from
\cref{sec:structural-geometric-locality}.  Recall that
$p_0=1$, $p_i=\prod_{j=1}^i\varsigma_j$, and
$q_i=u_i+v_i$.

\begin{theorem}[Value, output, and central-start separation]
\label{thm:corr-hidden-sign}
For the family \eqref{eq:structural-locality-lp}:
\begin{enumerate}[leftmargin=*]
 \item It is strictly primal--dual feasible and has the unique optimum
 \begin{equation}
  \operatorname{OPT}=N+1+\alpha p_N.                    \label{eq:corr-hidden-opt}
 \end{equation}
 Estimating its shifted value $\operatorname{OPT}-(N+1)$ to additive error
 less than $\alpha$ requires $\Omega(N)$ bounded-error raw sparse queries.
 \item Every classical nonnegative output $\widehat x$ satisfying
 $\norm{A_\varsigma\widehat x-b}_1<1/2$ determines $p_N$ and therefore also
 requires $\Omega(N)$ queries.
 \item The primal central path has the public support-two orthonormal null
 basis stated in \cref{sec:structural-geometric-locality}, and its reduced
 log-barrier Hessian has condition number one for every $\mu>0$.
 \item At $\mu=3/4$, one lookup at the final cell of a reusable random-access
 value oracle for the exact primal central point reveals $p_N$.  If setup and
 that lookup use respectively $Q_{\rm setup}$ and $Q_{\rm lookup}$ raw
 queries, then
 $Q_{\rm setup}+Q_{\rm lookup}=\Omega(N)$.  In particular, query-free
 invocations after preprocessing force $Q_{\rm setup}=\Omega(N)$.
\end{enumerate}
\end{theorem}

\begin{proof}
The feasibility argument and central-path calculations are given throughout
\cref{sec:structural-geometric-locality}, with the displayed formulas in
\eqref{eq:structural-locality-barrier}--
\eqref{eq:structural-locality-condition-one}.  In particular, feasibility
fixes $d_i=p_i$ and leaves $q_i\geq1$, so the objective is
$\sum_iq_i+\alpha p_N$ and has the unique minimum
\eqref{eq:corr-hidden-opt}.  The shifted value reveals parity.

For an approximate output define
$r_0=\widehat d_0-1$ and
$r_i=\widehat d_i-\varsigma_i\widehat d_{i-1}$.  Telescoping gives
\begin{equation}
 \widehat d_N-p_N
 =\sum_{j=0}^N\left(\prod_{k=j+1}^N\varsigma_k\right)r_j,               \label{eq:corr-hidden-telescope}
\end{equation}
so $|\widehat d_N-p_N|<1/2$ under the asserted residual bound.  The sign of
$\widehat d_N=\widehat u_N-\widehat v_N$ is therefore $p_N$.

Finally, \eqref{eq:structural-locality-center} gives $q=2$ at $\mu=3/4$.
The final pair is $(3/2,1/2)$ or its swap according to $p_N$, so one
coordinate-oracle lookup determines parity even with coordinate error below
$1/2$.  Since quantum parity has query complexity $\Theta(N)$, each claimed
raw-query lower bound, including the combined setup--lookup bound, follows
from the local sign-oracle reduction used throughout \cref{sec:lp-hardness}.
\end{proof}

This theorem charges an original-form feasible start.  It does not apply to
infeasible-start or homogeneous self-dual methods, which can supply a
different universal start; nor does it say that the scalar reduced Newton
solve itself is hard.

\subsection{Dead ends and false shortcuts}
\label{sec:corrections-dead-ends}

\paragraph{First inverse versus filtering.}
A small value of
$\norm H\norm{H^{-1}f}/\norm f$ does not imply a fast DLS filtering solve.
Filtering depends on the second inverse through
$\norm H\norm{H^{-2}f}/\norm{H^{-1}f}$, and active--inactive coupling can
restore the full condition-number exponent.  The exact identities and sparse
sharpness pair are in
\cref{sec:beyond-kappa-coupling,sec:beyond-kappa-truncation}.

\paragraph{Sparse support versus inverse access.}
Sparse Hessian or KKT support does not imply a sparse inverse, a cheap block
encoding after elimination, or low circuit depth on bounded-range hardware.
The factor-normalization and inverse-compilation frontiers in
\cref{sec:preconditioning-whitening} prove the access cost; the bounded-range
depth claim requires the bounded-occupancy geometric input placement of
\cref{sec:structural-geometric-locality}.  The dense metric in
\cref{sec:corrections-densification} gives a separate coordinate-level
warning.  None of these statements says that every structured inverse is
expensive.

\paragraph{Condition one versus the surrounding maps.}
A condition-one reduced Newton matrix does not make the affine feasible
offset, coordinate gauge, transformed right-hand side, recovery map, or
requested output free.  The LP separations in
\cref{sec:lp-affine-separation,sec:lp-prefix-rigidified} place parity in
the affine offset while leaving the reduced geometry public.  The full-KKT
and preconditioner dichotomies in
\cref{sec:lp-full-kkt,sec:preconditioning-parity} show how input dependence
can relocate among these objects.

\paragraph{Exact Schur preconditioning.}
An exact Schur transformation can make the abstract product constantly
conditioned while relocating input dependence into construction, transformed
RHS loading, or recovery.  The complete ledger is proved in
\cref{sec:lp-full-kkt,sec:preconditioning-whitening}.  Supplying the
preconditioned product itself as a unit-cost oracle is a strictly stronger
input model, not a consequence of sparse access to the original system.

\paragraph{State output versus coordinate access.}
A normalized solution state is not a reusable coordinate oracle.  Its norm
and global phase are absent, while coherent diagonal construction also needs
controlled and adjoint access and pays a normalization factor.  The tight
state-to-diagonal product appears in
\cref{sec:preconditioning-state-interface}; the implicit-coordinate escape
route in \cref{sec:positive-compressed-dual} states the stronger interface it
actually constructs.

\paragraph{One-step bounds versus a fixed trajectory.}
Reapplying a one-step lower bound at every point of one fixed LP trajectory
is invalid without freshness, a memory restriction, an adaptivity/causality
restriction, or mandatory explicit output.  The fresh-oracle direct sum is
\cref{thm:prec-online-direct-sum}; the surrounding discussion shows how
caching a static input can erase later raw-query cost.

\paragraph{Residual bounds versus projective variation.}
Ordinary short-step neighborhood and relative OSS residual bounds do not
imply finite projective variation.  Section~\ref{sec:positive-lazy} gives a
sparse cycling counterexample and the additional hypotheses needed for
reuse.  The refined audit in \cref{sec:positive-face-repair} also corrects
an earlier version's proposed sharpness family: its convergent leakage and
transverse-forcing ledgers still allow linearly growing realized variation,
and a separate oscillating-displacement family shows that adding the
$O(\mu_t)$ tail bound on face leakage does not help either; finite realized
variation is therefore stated as an independent hypothesis rather than
derived from the ledgers.

\paragraph{Expanderization.}
Replacing a path by an expander does not evade the bounded-degree
local-linear mixture or cut/resistance obstruction.  Those no-go results are
\cref{sec:frontier-mixture,sec:frontier-electrical}.  In the favorable
well-spread regime, \cref{sec:positive-checkpoint-span} gives the matching
classical sampling shortcut; making rare columns influential instead moves
cost into normalization or one-time search.

\paragraph{Scalar gap versus density visibility.}
A constant scalar objective gap can coexist with a trace-normalized central
density that is only $O(1/G)$ from a public state.  Thus an output lower bound
must state its accuracy and normalization contract.  The zero-query boundary
and the mass--visibility inequalities are proved in
\cref{sec:condensation-visibility,sec:condensation-query-lower}; the dilution
example is \cref{rem:wta-dilution}.

\paragraph{Pairwise distance versus a common decoder.}
Pairwise distinguishability of states indexed by an unknown winner set does
not provide one public measurement that decodes every set.  The failed
collision shortcut loses a factor $1/k$.  Section~\ref{sec:state-conversion-model}
repairs this only by measuring a proposed label and verifying that label with
fresh input queries; the verifier, not pairwise trace distance alone, supplies
the common decision procedure.

\subsection{Disposition of remaining algorithmic candidates}
\label{sec:corrections-dispositions}

To state the factorized rectangular observation precisely, let $B_\mu$ have
full row rank, let $f=B_\mu h$ lie in $\operatorname{range}(B_\mu)$, and let
$\alpha_{B_\mu}\geq\norm{B_\mu}$ be its block-encoding normalization.  Define
\begin{equation}
 \rho_{B_\mu}(f)=\alpha_{B_\mu}
 \frac{\norm{(B_\mu B_\mu^T)^{-1}f}_2}
 {\sqrt{f^T(B_\mu B_\mu^T)^{-1}f}}.                    \label{eq:corr-rectangular-rho}
\end{equation}
This pseudoinverse parameter is distinct from the nonsingular filtering
parameter in \eqref{eq:dls-general-filtering-parameter}.  In the mixed
strictly-complementary exact-centering asymptotic, assume
$B_\mu B_\mu^T=\mu^{-1}C+\mu D+o(\mu)$,
$f=b\in U:=\operatorname{range}C$, and the generic active--inactive coupling
of \cref{sec:beyond-kappa-coupling}, so that
\[
 (B_\mu B_\mu^T)^{-1}b=\mu(p,-g)+o(\mu),
 \qquad b^T(B_\mu B_\mu^T)^{-1}b=\Theta(\mu).
\]
With the essential normalization
$\alpha_{B_\mu}=\Theta(\mu^{-1/2})$, substitution in
\eqref{eq:corr-rectangular-rho} gives the proved algebraic observation
$\rho_{B_\mu}(b)=\Theta(1)$.

Here $D_\mu=(XS^{-1})^{1/2}$, $B_\mu=AD_\mu$, and
$h=(XS)^{-1/2}r_c$ is the scaled complementarity right-hand side; thus
$f=B_\mu h$ is the row-space right-hand side denoted ``$Bh$''.  The two
scaled direction components are
\[
 u=\Pi_{\ker B_\mu}h,
 \qquad
 v=B_\mu^+f=B_\mu^+(B_\mu h),
\]
with $u+v=h$.  They recover the physical directions through
$\Delta x=D_\mu u$ and $\Delta s=D_\mu^{-1}v$.  The
constant value of \eqref{eq:corr-rectangular-rho} is not an end-to-end QIPM
speedup.  Six obligations remain: a DLS-quality rank-deficient rectangular
filter theorem with its cutoff and error law; preparation of $f=B_\mu h$;
coherent construction, norm control, and cancellation-safe physical recovery
of $u$ and $v$; inexact-neighborhood propagation; maintenance of the changing
diagonal oracles; and a complete output ledger with a strong classical
comparison.  Li's instance-dependent rectangular minimum-norm algorithm
already uses closely related second-inverse geometry
\cite{li2025-new-condition-number-application-multivariate-polynomial}, so
the primitive also substantially collides with prior work.

A further candidate, the augmented spectral Newton sketch, instead samples the
augmented matrix \([D_{\rm N}^{1/2}A^T,v_{\rm N}]\), where \(D_{\rm N}\) and
\(v_{\rm N}\) denote the sketch's Newton scaling and regression target.
Under its stated row-access assumptions,
a \((1\pm\varepsilon)\) embedding of the augmented Gram matrix gives a valid
least-squares embedding lemma at
\(\Otilde(\sqrt{nm}/\varepsilon)\) row-query scale; the infeasible-step
analysis, implicit update mechanism, and end-to-end comparison remain open,
and the direction overlaps the factorized rectangular candidate above.

State-to-diagonal lifting is a useful but known primitive, not a new result:
Rattew and Rebentrost give the amplitude-to-diagonal block encoding, with the
phase-aligned coherent access assumptions recorded in
\cref{sec:preconditioning-state-interface}
\cite{rattew2023-non-linear-transformations-of-quantum}.  Likewise, the
projective-plus-implicit expander candidate is closed by the dichotomy in
\cref{sec:positive-checkpoint-span}: bounded contributions admit a comparable
classical sketch, while rare influential columns incur normalization or
one-time search cost.

\section{Discussion and Open Problems}
\label{sec:discussion}

We asked whether QIPMs beat classical IPMs on sparse problems.  Instead of a
single yes or no, the results separate three levels, each with its own
obstructions and ways around them, that informal complexity claims often
combine: the geometry of the central path, the model of access to the
resulting linear-algebra objects, and the interface through which one
iteration produces the next.

\subsection{What the paper establishes}
\label{sec:discussion-verdict}

\paragraph{Geometry.}
For barrier-Newton methods, part of the condition-number growth is forced
by the optimization problem.  \Cref{thm:sublevel-chord-law} bounds the
condition number of the reduced barrier Hessian from below at a fixed
objective gap, for every self-concordant barrier, and
\cref{thm:canonical-achievability} shows that the log and log-det barriers
attain the resulting order up to instance-dependent constants.
\Cref{thm:width-spectrum-rigidity} extends this to a geometric description
of the whole spectrum.  In the stated minimax sense, changing the barrier
cannot in general remove the lower bound given by sublevel chords.

This is not the folklore claim that every central path has
\(\kappa=\Theta(\mu^{-2})\); the rate depends on the geometry of the
sublevel sets.  A positive-dimensional optimal face gives growth of order
\(g^{-2}\) in the objective gap \(g\); a curved unique optimum can give
growth of order \(g^{-1}\); at a unique LP vertex the condition number can
stay bounded; and an attained singularity degree can give fractional
exponents (\cref{cor:conditioning-regimes,cor:fractional-conditioning-laws}).
Nor does a large condition number alone determine solver cost: Krylov
iteration counts can be polylogarithmic in an arbitrarily large ratio
between two eigenvalue clusters
(\cref{prop:two-cluster-hessian,prop:two-cluster-cg}), and the Newton
right-hand side on the path can avoid the small eigenvalues
(\cref{prop:newton-rhs-alignment,cor:top-window-solve}).

\paragraph{Access models.}
The quantum cost depends on how the matrix is supplied, not only on its
spectrum.  The clustered comparison in \cref{thm:quantum-cluster-obstruction}
separates plain normalized block-encoding access from access to a factor
\(G\) with \(H=G^\dagger G\).  With plain access, the known worst-case query cost
\(\Thetatilde(\kappa)\) remains for two-point spectra in sufficiently large
dimension.  Norm-tight factor access admits a polynomial of square-root cost
but an end-to-end solve only under overlap assumptions; a generic
factor-access QLS solver of square-root cost would contradict the
plain-access lower bound, because the square-root factor of that hard
instance is easy to query.  By the exact coupling and
spectral-dimension laws
(\cref{thm:dls-coupling,thm:spectral-dimension-laws}), the distribution of the
right-hand side over the spectrum decides whether truncation or filtering
can improve on \(\kappa\).

The same holds for preconditioning.  A whitening factor may trade condition
number for a larger normalization; a congruence may trade the transformed
condition number for sensitivity when the physical solution is recovered; and
a perfect preconditioner that depends on the input may move the hard
information into setup or RHS preparation
(\cref{thm:prec-whitening-frontier,thm:prec-congruence-frontier,thm:prec-factor-rhs-hardness,thm:prec-parity-ledger}).
These results bound products of costs or the total cost of all phases.
They show not that preconditioning is useless, but that a small transformed
condition number is not by itself a complexity bound.

\paragraph{Interfaces.}
In our hard families the main obstruction is often loading or recovery
rather than inversion.  The prefix LP has a reduced Hessian of condition
number one along its central path, yet preparing a robust primal direction
in the original coordinates needs linearly many raw queries
(\cref{thm:lp-prefix-capstone}).  The SDP families distinguish a gap in the
optimal value (\cref{thm:sdp-intrinsic}), parity-hard primal and
complete-KKT solution states (\cref{thm:sdp-newton-state}), and a public Schur
solve for which loading the original coordinates obeys an accuracy--query
tradeoff (\cref{thm:sdp-schur,thm:sdp-accuracy-frontier}).  The canonical KKT
block encoding is exact at one raw query per call, yet preparing the root,
global, or complete-KKT state still needs linearly many calls to it, and
first constructing a block encoding of the tangent projector needs linearly
many raw queries (\cref{sec:sdp-intrinsic-access}).  Nor is a normalized state a
reusable coordinate oracle: producing the changing diagonal of the next Newton system
obeys the state-to-diagonal and state-to-next-solve product bounds in
\cref{thm:prec-state-diagonal,thm:prec-state-next-solve}.

In the condensation results, the winner mass controls both the condition
number of the reduced Hessian and the trace
distance from a public state with no winner
(\cref{thm:cond-mass-conditioning,thm:cond-visibility}).  The bounded-incidence
holonomy SDP has matching query bounds for the value and for a compact
output that names the winner (\cref{thm:cond-unique-or,thm:wta-search-upper}),
while the cost of preparing the central state changes across the phases of
condensation (\cref{sec:wta-central-state,eq:wta-preparation-phases}).  The
decoder that verifies sampled labels turns a trace-close mixed output into a
public decision procedure with one-sided error, and parity collapse gives
the query cost as a function of the allowed trace-distance error
(\cref{thm:sc-m1,thm:sc-small-success,cor:sc-s2}).  The lower bound holds for
a specified accuracy of the output state; a constant primal objective gap
alone would not imply it.

In short, conditioning is governed by sublevel geometry, quantum
performance by the access actually supplied, and end-to-end performance by
loading, recovery, updates, checks, and output between iterations.  Sparsity
can help at each level, but no single notion of sparsity controls all
three.

\subsection{What the paper does not establish}
\label{sec:discussion-nonclaims}

The paper proves neither an unconditional sparse-QIPM advantage nor an
unconditional lower bound.  The Apers--Gribling result for tall LPs remains
the only proved asymptotic quantum--classical separation, in row queries,
among the QIPM results considered here
\cite{apers2026-quantum-speedups-linear-programming-interior}; our results do
not weaken it.  Conversely, no parity or holonomy family rules out a method
with a different oracle, a weaker but still useful output, an
optimization primitive free of the condition number, or a structure-aware
preconditioner outside the stated class.

The barrier theorem concerns the reduced Hessian in an orthonormal basis of
the tangent space, not automatically an augmented KKT system, primal--dual
scaling, or a Hamiltonian central-path method.  The clustered QLS theorem
depends on the access model and has conditions on the dimension and on
overlaps.  The LP and SDP loading theorems count raw
coefficient queries; a supplied parity, prefix, tangent-projector, or
preconditioned-product oracle is a stronger input model.  The online direct
sum requires fresh independent input in each round, whereas a fixed instance
can be cached (\cref{thm:prec-online-direct-sum,sec:preconditioning-lazy}).
The physical light-cone bound concerns data and output placed locally with
bounded occupancy per site, not all-to-all query access
(\cref{thm:structural-light-cone}).

Likewise, the positive modules are components, not a complete algorithm.
Projective reuse of a rescaled Newton direction, as analyzed here, needs
finite realized variation and a fixed residual tolerance; equilibration of
the active range and face repair need known strict-complementarity tail margins and spectral gaps;
compressed dual output needs a retained dimension that is actually smaller
and must count the cost of its coordinate compiler; and the block-angular
module needs a projected barrier with a small parameter and a local
representation kept between iterations (\cref{sec:positive-synthesis}).  The
counterexamples in
\cref{sec:positive-lazy,sec:positive-face-repair,sec:positive-checkpoint-span}
show that the corresponding module fails when one of these conditions is
removed.

We propose the five accounting rules of \cref{sec:positive-modules} as a
standard for future positive claims, and add the comparison rules used
throughout the paper: match the residual
and precision requirements, count setup and reuse costs on both sides, and
compare with the best classical method for the same structure.  A query
advantage is not automatically a gate or time advantage, and a comparison of
two upper bounds is not a separation.

\subsection{Collected open problems}
\label{sec:discussion-open}

We collect every open-problem environment in
\cref{sec:conditioning,sec:spectra,sec:state-conversion,sec:lp-hardness,sec:sdp-hardness}
and the further open questions raised in the technical discussion, grouped
by the level at which new information is needed.

\subsubsection{Geometry and conditioning}

\begin{enumerate}[leftmargin=*,label=G\arabic*.]
 \item \textbf{Degenerate-vertex constants.}  Determine the limiting spectrum,
 not only its boundedness, of the canonical reduced log-barrier Hessian at a
 unique degenerate LP vertex from the geometry of the active set
 (\cref{open:degenerate-vertex-constants}).
 \item \textbf{Facial characterization of chord growth.}  Characterize the
 exact growth \(D(g)\) of objective-sublevel chords of an SDP from its facial
 data, and decide when the singularity-degree exponent is attained by the
 sublevel diameters
 (\cref{open:facial-chord-growth}).
 \item Relate the neighborhood \eqref{eq:cond-robust-residual} used by robust
 condensation, which bounds the blockwise stationarity residual in the
 barrier norm by \(\eta\) and the trace residual by \(\zeta\), to the
 algorithm-specific inexact Newton residuals of
 \cref{def:inexact-newton-residual}, whose forcing tolerance is a separate
 \(\eta\) (\cref{sec:condensation-discussion}).
\end{enumerate}

\subsubsection{Quantum access and state conversion}

\begin{enumerate}[leftmargin=*,label=A\arabic*.]
 \item \textbf{Exact subnormalization shift.}  Decide whether plain access to
 \(H\) permits a block encoding of \(I-H\) with unit subnormalization without
 paying the apparent cost of uniform amplification at the shifted spectral
 edge (\cref{open:exact-subnormalization-shift}).  This is the
 promise-sensitive form of the normalization question of Orsucci and Dunjko
 \cite{orsucci2021-on-solving-classes-of-positive}.
 \item \textbf{General inner predicate.}  Prove or refute the
 \(\sqrt\delta\,Q(f)\sqrt{G/k}\) small-success composition lower bound for a
 general inner predicate \(f\), beyond the proof by parity collapse
 (\cref{open:sc-general-composition}).
 \item Determine the minimax output radius and query curve when the public
 internal states of the losing blocks differ by \(\eta_{\rm int}\), including
 the regime below condensation (\cref{sec:state-conversion-open}).
 \item Resolve the state-conversion range \(\delta=o(k/G)\), where a public
 output with uniform labels may already satisfy part of the output requirement
 and the winner mass alone does not determine the boundary
 (\cref{sec:state-conversion-open}).
 \item \textbf{Intrinsically nonabelian word hardness.}  Prove a lower bound
 on local-symbol queries for distinguishing conjugacy classes under a
 promise on words whose letters do not commute, and transfer it to output of
 the conjugacy class in the group-holonomy SDP, instead of restricting to a
 cyclic subgroup
 (\cref{open:sdp-nonabelian}).
\end{enumerate}

\subsubsection{Hard families and classes of gadgets}

\begin{enumerate}[leftmargin=*,label=F\arabic*.]
 \item \textbf{Bounded-scale full-KKT parity amplifier.}  Construct, or rule
 out, a standard-form LP of linear size with bounded row and column degree;
 bounded coefficients, each depending on at most one input bit; public Newton
 right-hand sides that are simple to prepare; primal and slack starting
 values bounded above and below by positive constants; a complete Newton
 direction bounded coordinatewise; and a parity-hard primal component with
 constant mass in the state of that complete direction
 (\cref{open:lp-bounded-full-kkt}).
 \item Determine whether an arbitrary sparse equality LP with bounded range
 can support a nonlinear decoder of the normalized state with a fixed global
 residual tolerance without introducing nonlocal coefficient access
 (\cref{sec:frontier-synthesis}).
 \item Determine whether a nonorthogonal mixed-mode gadget can solve the same
 bounded-scale complete-direction problem while counting the cost of every
 input-dependent source and recovery map
 (\cref{sec:lp-beyond-pairs,sec:frontier-synthesis}).
\end{enumerate}

\subsubsection{Positive algorithmic results}

\begin{enumerate}[leftmargin=*,label=P\arabic*.]
 \item Exhibit one explicit family and access model that satisfy all six
 requirements of \cref{sec:positive-synthesis} simultaneously: stable direct
 normalization and inverse, stable forcing and RHS preparation, finite
 projective variation or a barrier with a small parameter, cheap static
 construction, useful compressed output, and a classical lower bound that
 holds even with randomized preprocessing.
 \item Find a structured access model that supplies a quantum operator with
 constant normalization without also supplying a classical sampling
 distribution with bounded leverage, and realize it on an explicit sparse LP
 family (\cref{sec:positive-checkpoint-span}).
 \item Construct a projected block-angular barrier with parameter
 \(\Otilde(k+\poly(p))\) whose local center and derivatives can be recomputed
 from one scenario block and the explicit border
 (\cref{sec:positive-block-angular}).
\end{enumerate}

The list omits directions that a theorem here already rules out: turning a
bounded-degree local-linear parity gadget into an
expander does not escape the mixture and resistance bounds, and ordinary
residual control in short steps does not imply finite projective variation
(\cref{sec:frontier-synthesis,sec:positive-lazy}).  A new proposal must drop a
premise of these theorems, not relabel the same construction.

\subsection{Outlook}
\label{sec:discussion-outlook}

The most useful next result would be small and complete: one structured
family, one explicit access model, one useful
output requirement, an accounting of all costs along the path, and a matching
classical lower bound or algorithmic comparison.  Problem~P1 is deliberately
hard because each module alone already has a classical or quantum
obstacle.  Stable spectral geometry is not enough if RHS preparation is
hard; cheap coherent access is not enough if a classical algorithm can
sample the same averages; and compressed output is not enough if the next
iterate must be written out in full.

Among nearer-term targets, Problem~A1 would show whether the separation
between plain and factor access for clustered QLS is an artifact of the
current construction or a property of the access models; Problem~A2 would
make the state-conversion curve a reusable composition principle rather than
a result about parity; Problem~A5 would separate matrix-valued holonomy from
its abelian restriction; and Problem~P3 would connect the border module to a
candidate end-to-end QIPM.

The methodological lesson is that a future sparse-QIPM advantage cannot be
claimed from one favorable parameter: it must hold for the geometry, the
access model, and the output interface at once.  This paper proves bounds
and counterexamples for each of these, with conditional modules where one
can be made cheap.  Solving Problem~P1 would give the end-to-end answer that
this paper does not claim.


\bibliographystyle{abbrvnat}
\bibliography{bibliography}

\end{document}
